\documentclass[10pt,oneside,openany]{book}
\usepackage[T1]{fontenc}
\usepackage[utf8]{inputenc}
\usepackage{babel}
\babelprovide[import,main]{spanish}
\usepackage{lmodern,amsmath,amssymb,amsthm,mathrsfs}
\usepackage{iftex}
\ifLuaTeX\usepackage{luatex85}\fi
\usepackage[all]{xy}
\usepackage{tikz-cd}
\usepackage{xstring,xparse,xcolor,etoolbox}
\usepackage[a4paper,left=18mm,right=18mm,top=22mm,bottom=23mm,marginparwidth=14mm,marginparsep=2mm]{geometry}
\usepackage{array,longtable,multicol,booktabs}
\usepackage[shortlabels]{enumitem}
\usepackage{marginnote,microtype,graphicx,adjustbox}
\usepackage[unicode=true,pdfencoding=auto,colorlinks=true,linkcolor=blue,citecolor=blue,urlcolor=blue,hypertexnames=false,hyperfootnotes=false,linktocpage=true]{hyperref}
\usepackage{bookmark}
\renewcommand{\thesection}{\arabic{section}}
\renewcommand{\thesubsection}{\thesection.\arabic{subsection}}
% Keep PDF destinations unique when the cumulative reader resets table counters.
\renewcommand{\theHtable}{\arabic{part}.\arabic{table}}
\setcounter{tocdepth}{2}
\setlength{\emergencystretch}{12em}
\tolerance=3000
\binoppenalty=200
\relpenalty=200
\hfuzz=0.1pt
\vfuzz=0.1pt
\raggedbottom
\allowdisplaybreaks[3]
\newtheoremstyle{egaenv}{}{}{\normalfont}{}{\bfseries}{.}{ }{\thmnote{(#3)}}
\newtheoremstyle{egastatement}{}{}{\itshape}{}{\bfseries}{. ---}{ }{\thmname{#1}\thmnote{ (#3)}}
\newtheoremstyle{egadescription}{}{}{\normalfont}{}{\bfseries}{. ---}{ }{\thmname{#1}\thmnote{ (#3)}}
\theoremstyle{egaenv}\newtheorem*{env}{---}
\theoremstyle{egastatement}
\newtheorem*{theorem}{Teorema}
\newtheorem*{proposition}{Proposición}
\newtheorem*{lemma}{Lema}
\newtheorem*{corollary}{Corolario}
\theoremstyle{egadescription}
\newtheorem*{definition}{Definición}
\newtheorem*{remark}{Observación}
\newtheorem*{remarks}{Observaciones}
\newtheorem*{example}{Ejemplo}
\newtheorem*{examples}{Ejemplos}
\newtheorem*{scholium}{Escolio}
\newtheorem*{notation}{Notación}
\newenvironment{namedtheorem}[2][]{\par\medskip\noindent\setmanualcurrentlabel{#1}\textbf{#2\ifstrempty{#1}{}{ (#1)}. ---}\ifstrequal{#2}{Observación}{\normalfont}{\ifstrequal{#2}{Observaciones}{\normalfont}{\itshape}}}{\par\medskip}
\newenvironment{namedlemma}[2][]{\par\medskip\noindent\setmanualcurrentlabel{#1}\textbf{#2\ifstrempty{#1}{}{ (#1)}. ---}\itshape}{\par\medskip}
\renewcommand{\proofname}{Demostración}
\renewcommand{\contentsname}{Índice general}
\renewcommand{\bibname}{Bibliografía}
\providecommand{\refname}{Bibliografía}
\newcommand{\oldpage}[2][]{\marginnote{\normalfont\scriptsize\textbf{#1}\,|\,#2}\ignorespaces}
\newcommand{\sectionbreak}{\begin{center}***\end{center}}
\newcommand{\erratum}[2][]{#2\marginnote{\scriptsize\textbf{Err}\textsubscript{#1}}}
\providecommand{\shHom}{\mathcal{H}\!\mathit{om}}
\providecommand{\shProj}{\mathcal{P}\!\mathit{roj}}
% Literal macro definition from the French EGA II notation index.
\providecommand{\egaShProj}{\mathcal{P}\!\mathit{roj}}
\providecommand{\shExt}{\mathcal{E}\!\mathit{xt}}
\providecommand{\shTor}{\mathcal{T}\!\mathit{or}}
\providecommand{\shGr}{\mathcal{G}\!\mathit{r}}
\providecommand{\shDer}{\mathcal{D}\!\mathit{er}}
\providecommand{\shDiff}{\mathcal{D}\!\mathit{iff}}
\providecommand{\shHomcont}{\mathcal{H}\!\mathit{om.cont}}
\providecommand{\shAut}{\mathcal{A}\!\mathit{ut}}
\providecommand{\Frac}{\operatorname{Frac}}
\providecommand{\unsure}[1]{#1}
\providecommand{\completelyunsure}{\textbf{??}}
\providecommand{\og}{\guillemotleft{}}
\providecommand{\fg}{\guillemotright{}}
\makeatletter
\newcommand{\setmanualcurrentlabel}[1]{\def\@currentlabel{#1}}
\newcommand{\EGARefNumber}[1]{\begingroup\StrSubstitute{#1}{-fr}{}[\egarefkey]\StrSubstitute{\egarefkey}{-es}{}[\egarefkey]\StrSubstitute{\egarefkey}{.fr}{}[\egarefkey]\StrBehind{\egarefkey}{.}\endgroup}
\newcommand{\EGAMaybeRefLink}[2]{\@ifundefined{r@#1}{#2}{\hyperref[#1]{#2}}}
\newcommand{\EGAReferenceAlias}[2]{\ifcsdef{r@#2}{\csletcs{r@#1}{r@#2}}{}}
\NewDocumentCommand{\sref}{o m o}{\textup{(\IfNoValueF{#1}{\textbf{#1},~}\EGAMaybeRefLink{#2}{\EGARefNumber{#2}}\IfNoValueF{#3}{,~#3})}}
\newcommand{\eref}[1]{\textup{(\EGAMaybeRefLink{#1}{\EGARefNumber{#1}})}}
\newcommand{\EGAReaderTarget}[2]{\phantomsection\def\@currentlabel{#2}\label{#1}}
\newcommand{\EGAReadLink}[2]{\EGAMaybeRefLink{#1}{#2}}
\newcommand{\EGAUnavailableHyperref}[2]{\EGAMaybeRefLink{#1}{#2}}
\newcommand{\sgaStableTarget}[1]{\hypertarget{#1}{}}
\newcommand{\sgaStableOrOriginalLink}[2]{\EGAMaybeRefLink{#1}{#2}}
\newcommand{\EGAIVFourExternalLabel}[1]{}
\newcommand{\EGAIVFourExternalCitation}[2]{\expandafter\def\csname b@#1\endcsname{#2}}
\makeatother
\newcommand{\EGAResetVolumeCounters}{\setcounter{section}{0}\setcounter{subsection}{0}\setcounter{equation}{0}\setcounter{footnote}{0}\setcounter{figure}{0}\setcounter{table}{0}}
\newcommand{\EGAVolumeTitlePage}[3]{\clearpage\EGAResetVolumeCounters\chapter*{#2\\#3}\addcontentsline{toc}{chapter}{#2: #3}}
\newcommand{\MR}[1]{MR #1}
\newcommand{\MRhref}[2]{\href{#1}{#2}}
\providecommand{\bysame}{\leavevmode\hbox to3em{\hrulefill}\,}
\pagestyle{plain}

% BEGIN authoritative French CC0 preamble-base.tex
% Notation of the EGA editions: the abbreviations the transcription uses for alphabets, symbols, operators and
% arrows. Written for these editions and dedicated to the public domain under CC0 1.0.
\makeatletter
% a command that may or may not exist already is (re)defined in one step
\newcommand{\ega@define}[2]{\providecommand{#1}{}\renewcommand{#1}{#2}}
\newcommand{\ega@operator}[2]{\providecommand{#1}{}\renewcommand{#1}{\operatorname{#2}}}

% alphabets: sheaves, bold letters, categories
\ega@define{\sh}{\mathcal}
\ega@define{\bb}{\mathbf}
\ega@define{\cat}{\mathtt}
\ega@define{\C}{\cat{C}}
\ega@define{\Set}{\cat{Set}}
\ega@define{\CHom}{\cat{Hom}}
\ega@define{\op}{^{\cat{op}}}
\ega@define{\GL}{\bb{GL}}
\ega@define{\bbGamma}{\boldsymbol{\Gamma}}

% letters and symbols
\renewcommand{\leq}{\leqslant}
\renewcommand{\geq}{\geqslant}
\ega@define{\setmin}{-}
\ega@define{\emp}{\varnothing}
\ega@define{\vphi}{\varphi}
\ega@define{\kres}{k}
\ega@define{\rad}{\mathfrak{r}}
\ega@define{\nilrad}{\mathfrak{N}}
\ega@define{\HH}{\mathrm{H}}
\ega@define{\CHH}{\check{\mathrm{H}}}
\ega@define{\RR}{\mathrm{R}}
\ega@define{\LL}{\mathrm{L}}
\ega@define{\red}{\mathrm{red}}
\providecommand{\dual}{}\renewcommand{\dual}[1]{{#1}^{\vee}}
% the tilde of the sheaf associated with a module, written after it: M\supertilde
\ega@define{\supertilde}{{\,\widetilde{\,}\,}}

% operators
\ega@operator{\rg}{rg}
\ega@operator{\gr}{gr}
\ega@operator{\Gr}{Gr}
\ega@operator{\Hom}{Hom}
\ega@operator{\Proj}{Proj}
\ega@operator{\Tor}{Tor}
\ega@operator{\Ext}{Ext}
\ega@operator{\Supp}{Supp}
\ega@operator{\Ker}{Ker}
\ega@operator{\Im}{Im}
\ega@operator{\Coker}{Coker}
\ega@operator{\Spec}{Spec}
\ega@operator{\Spm}{Spm}
\ega@operator{\Spf}{Spf}
\ega@operator{\grad}{grad}
\ega@operator{\dimc}{dimc}
\ega@operator{\codim}{codim}
\ega@operator{\id}{id}
\ega@operator{\Der}{Der}
\ega@operator{\Diff}{Diff}
\ega@operator{\Hyp}{Hyp}
\ega@operator{\Tr}{Tr}

% arrows: long in displays, short in text, scripts and subscripts
\let\ega@shortmapsto\mapsto
\renewcommand{\to}{\mathchoice{\longrightarrow}{\rightarrow}{\rightarrow}{\rightarrow}}
\providecommand{\from}{}\renewcommand{\from}{\mathchoice{\longleftarrow}{\leftarrow}{\leftarrow}{\leftarrow}}
\renewcommand{\mapsto}{\mathchoice{\longmapsto}{\ega@shortmapsto}{\ega@shortmapsto}{\ega@shortmapsto}}
\ega@define{\isoto}{\simeq}
\ega@define{\simto}{\xrightarrow{\sim}}

% wide accents
\renewcommand{\hat}[1]{\widehat{#1}}
\renewcommand{\tilde}[1]{\widetilde{#1}}
\makeatother

% END authoritative preamble-base.tex

\hypersetup{pdftitle={EGA III, primera parte},pdfauthor={Alexander Grothendieck y Jean Dieudonné}}
\begin{document}
\begin{titlepage}\centering\vspace*{20mm}
{\Huge\bfseries Elementos de geometría algebraica\par}\vspace{14mm}
{\LARGE EGA III, primera parte\par}\vspace{10mm}
{\large Alexander Grothendieck\\con la colaboración de Jean Dieudonné\par}
\vfill{\Large Edición española diplomática\par}
\vspace{8mm}{\large 25 de noviembre de 2026\par}\end{titlepage}
\chapter*{Créditos y derechos}
\addcontentsline{toc}{chapter}{Créditos y derechos}
\noindent\textbf{Trabajo de esta ejecución.} Traducción y preparación por GPT-6.0 Astra (OpenAI), modo Pro de ChatGPT, noviembre de 2026. Traducción producida por IA; correcciones producidas por IA. Cotejo francés heredado: Claude Opus 5.5, esfuerzo max, noviembre de 2026.

\subsection*{Créditos previos conservados}

Las traducciones están producidas por IA, y también las correcciones, la edición y la composición. Tradujeron
modelos de OpenAI Codex (GPT-5.6 Sol y GPT-6.1 Sol, esfuerzo Ultra); GPT-6 Astra (esfuerzo Ultra) revisó, corrigió
y montó las ediciones. En la
preparación de algunas secciones de EGA se usaron fuentes inglesas que incorporan la traducción humana del proyecto
ryankeleti/ega; para FGA, la traducción inglesa de Tim Hosgood sirvió de referencia (véase «Fuentes y créditos»).

Los lectores de EGA I (traducción de OpenAI Codex, GPT-6.1 Sol, esfuerzo de razonamiento ultra, agosto–octubre
de 2026) se publicaron el 15 de noviembre de 2026; la composición (nada sobresale del margen más de 2 pt) y la
publicación las hizo Claude Opus 5.5 con esfuerzo high.

Las ediciones corregidas y las traducciones diplomáticas, que siguen el texto impreso, son documentos separados; las
notas matemáticas también.

\subsection*{Fuentes y créditos del repositorio}

\begin{itemize}
\item EGA: Alexander Grothendieck, con la colaboración de Jean Dieudonné, \emph{Éléments de géométrie algébrique},
  \emph{Publications mathématiques de l'IHÉS} 4, 8, 11, 17, 20, 24, 28 y 32 (1960–1967), en
  \href{https://www.numdam.org/journals/PMIHES/}{NUMDAM}.
\item FGA: Alexander Grothendieck, \emph{Fondements de la géométrie algébrique}, exposiciones 149, 182, 190, 195, 212, 221,
  232 y 236 del Seminario Bourbaki, en \href{https://www.numdam.org/}{NUMDAM}.
\item EGA: las fuentes inglesas usadas para preparar EGA III-1 (capítulo 0, §§8–10 y parte del §11; capítulo III, §§1 y
  3), EGA IV-1 (capítulo 0, §14) y EGA IV-4 (§16 y el comienzo de los §§17 y 20) incorporan la traducción humana de
  \href{https://github.com/ryankeleti/ega/tree/3d7afe28c99d1887ccdb85acf2ca1afd93fb1971}{ryankeleti/ega} (versión
  3d7afe2), cuyos colaboradores son Mateo Carmona, Tim Holzschuh, Tim Hosgood, Ryan Keleti, Thiago Solovera \& Nery y
  Wayne Yang (\href{https://github.com/ryankeleti/ega/blob/3d7afe28c99d1887ccdb85acf2ca1afd93fb1971/README.md}{créditos por sección}).
\item FGA: la traducción inglesa de \href{https://github.com/thosgood/fga/tree/9976dd5f8e2ddcb388f0a07f3aa334f120d1019f}{Tim Hosgood}
  (versión 9976dd5), que agradece las correcciones de Steve Hnizdur y el apoyo técnico de Jon Sterling.
\item La traducción española no se atribuye a estos traductores. \href{../../tree/repo/ATTRIBUTION.md}{ATTRIBUTION.md} da las versiones exactas
  de las fuentes inglesas.
\end{itemize}

\subsection*{Fuentes y autoría}

\subsection*{FGA: traducción inglesa de referencia}

La obra original, \emph{Fondements de la géométrie algébrique}, es de Alexandre Grothendieck. La traducción inglesa de \href{https://github.com/thosgood/fga/tree/9976dd5f8e2ddcb388f0a07f3aa334f120d1019f}{Tim Hosgood} es una fuente de referencia de esta edición; su proyecto agradece las correcciones de Steve Hnizdur y el apoyo técnico de Jon Sterling. La versión 9976dd5f8e2ddcb388f0a07f3aa334f120d1019f identifica el texto inglés usado en la comparación; no establece la versión exacta incorporada al principio. La traducción española y las correcciones y la composición posteriores están producidas por IA y no se atribuyen a Tim Hosgood. Esta atribución no constituye permiso ni aval.

\subsection*{EGA: fuentes inglesas}

Las fuentes inglesas utilizadas para preparar las traducciones españolas de EGA incorporan la traducción humana de \href{https://github.com/ryankeleti/ega}{ryankeleti/ega}, versión \href{https://github.com/ryankeleti/ega/tree/3d7afe28c99d1887ccdb85acf2ca1afd93fb1971}{3d7afe28c99d1887ccdb85acf2ca1afd93fb1971}. El proyecto nombra como colaboradores a Mateo Carmona, Tim Holzschuh, Tim Hosgood, Ryan Keleti, Thiago Solovera \& Nery y Wayne Yang; sus \href{https://github.com/ryankeleti/ega/blob/3d7afe28c99d1887ccdb85acf2ca1afd93fb1971/README.md}{créditos por sección} y su \href{https://github.com/ryankeleti/ega/commits/3d7afe28c99d1887ccdb85acf2ca1afd93fb1971}{historial de contribuciones} precisan la autoría. Las secciones afectadas son:

\begin{itemize}
\item EGA III-1: los §§ 8–10 y parte del § 11 del capítulo 0, y los §§ 1 y 3 del capítulo III.
\item EGA IV-1, capítulo 0: el § 14.
\item EGA IV-4: el § 16 y las partes iniciales de los §§ 17 y 20 del capítulo IV.
\end{itemize}

La traducción inglesa y la estructura TeX heredadas conservan esa procedencia. La traducción española y las adiciones y correcciones producidas por IA no se atribuyen a esos colaboradores. La atribución no implica autorización ni aval.

\subsection*{Traducción española}

La traducción al español, las correcciones, la edición y la composición están producidas por IA. Tradujeron modelos de OpenAI Codex (GPT-5.6 Sol y GPT-6.1 Sol, esfuerzo Ultra); GPT-6 Astra (esfuerzo Ultra) revisó, corrigió y montó las ediciones. No se afirma revisión humana. La autoría de las obras originales y de las contribuciones humanas citadas se conserva.

\begin{itemize}
\item \href{../../tree/repo/editions/II\_CURRENT/FUENTES.md}{EGA II: alcance e identidad de las fuentes}.
\item \href{../../tree/repo/editions/III1\_CURRENT/FUENTES.md}{EGA III-1: alcance e identidad de las fuentes}.
\item \href{../../tree/repo/editions/IV1\_CURRENT/FUENTES.md}{EGA IV-1, capítulo 0: alcance e identidad de las fuentes}.
\item \href{../../tree/repo/editions/IV4/FUENTES.md}{EGA IV-4: alcance e identidad de las fuentes}.
\end{itemize}
\subsection*{Fuentes y contribuciones}

La obra original es de Alexandre Grothendieck, con la colaboración de Jean Dieudonné: Éléments de géométrie algébrique III, première partie, Publications Mathématiques de l’IHÉS 11 (1961).

Las fuentes inglesas utilizadas en la preparación de esta traducción española incorporan la traducción humana de \href{https://github.com/ryankeleti/ega}{ryankeleti/ega}, versión \href{https://github.com/ryankeleti/ega/tree/3d7afe28c99d1887ccdb85acf2ca1afd93fb1971}{3d7afe28c99d1887ccdb85acf2ca1afd93fb1971}, en los §§ 8–10 y parte del § 11 del capítulo 0, y los §§ 1 y 3 del capítulo III. El proyecto nombra a Mateo Carmona, Tim Holzschuh, Tim Hosgood, Ryan Keleti, Thiago Solovera \& Nery y Wayne Yang como colaboradores; sus \href{https://github.com/ryankeleti/ega/blob/3d7afe28c99d1887ccdb85acf2ca1afd93fb1971/README.md}{créditos por sección} y su \href{https://github.com/ryankeleti/ega/commits/3d7afe28c99d1887ccdb85acf2ca1afd93fb1971}{historial de contribuciones} precisan la autoría.

La traducción inglesa y la estructura TeX heredadas conservan esa procedencia. La traducción al español, las adiciones posteriores y las correcciones y composición producidas por IA se distinguen de esas contribuciones humanas. No se atribuye a los colaboradores de la traducción inglesa la traducción española ni el material posterior que no figura en su versión. La atribución no implica autorización ni aval.

\subsection*{Identidad de las fuentes inglesas}

| Módulo inglés | SHA-256 de la versión utilizada |
| --- | --- |
| ega0/ega0-8.tex | 437a0b257ad4cc4fb327fde803738555aec37b6199982602d17df2e73c4f8215 |
| ega0/ega0-9.tex | 6f7e0b828f658412d1547fcd2391f4441a06598f327245ea39b4e204071e1e13 |
| ega0/ega0-10.tex | 9f8095d85c88ffa611438b38b874a79710013ec9ab29b77a51d3a2b9e05ff943 |
| ega0/ega0-11.tex | 6d5461a5e290defb591e4df163b5305887dd4ddfc3b5b0897552d59d446be20f |
| ega0/ega0-12.tex | 5cc7fd5fbf9170860069538e64a1c5ab892c5207cc96eff928cb6a771327819e |
| ega0/ega0-13.tex | fbf53c246359db7743005178cf7a45e73687c8b2095fc767664f74b23be1f5d9 |
| ega3/ega3-1.tex | 357fe715c0c8c8857e9e079fed86ae11c3a3bcbfb8abdeb26f0196e6cca99370 |
| ega3/ega3-2.tex | 2a00783ed287ef8244057b8aa11c4b75b5fdf3414cbaf79e3dceea6de5fcea3f |
| ega3/ega3-3.tex | 6ea5607a06235f4af977deed8f30306576fd667e7af4664e195053b8b94c5651 |
| ega3/ega3-4.tex | 6839cce7e94f347a6cc6ae83938307c33d7093ebaf32d0f16a98d3a21c4868eb |
| ega3/ega3-5.tex | 049f86a740b5a38cdf36b6f7001d6e7a8b34d1975ed5cd990a75ba833d0173e8 |
\subsection*{Declaración de derechos de esta edición}
Los derechos sobre las obras originales pertenecen a sus autores o a los representantes o herederos de estos. Nuestras traducciones, correcciones y composición tipográfica fueron producidas por IA. No tienen autor humano y, según nuestro entender, no están protegidas por derechos de autor; si aun así existiera algún derecho nuestro, se dedica al dominio público bajo CC0 1.0. Todo este trabajo reproduce o traduce las obras originales, por lo que nada de lo que aquí se publica concede licencia ni autoriza reproducción, traducción u otro uso de ellas. Estas ediciones se publican con fines educativos y sin ánimo de lucro, y no reclamamos ningún derecho que no poseamos. Si usted es titular de derechos sobre alguna de estas obras y se opone a su publicación aquí, abra una incidencia en GitHub y la retiraremos.
\clearpage\tableofcontents\clearpage
\EGAReferenceAlias{0III}{chapter:0III-fr}
\EGAReferenceAlias{0III-fr}{chapter:0III-fr}
\EGAReferenceAlias{0III-es}{chapter:0III-fr}
\EGAReferenceAlias{0III.fr}{chapter:0III-fr}
\EGAReferenceAlias{0III.es}{chapter:0III-fr}
\EGAReferenceAlias{0III.eq}{chapter:0III-fr}
\EGAReferenceAlias{eq:0III}{chapter:0III-fr}
\EGAReferenceAlias{chapter:0III}{chapter:0III-fr}
\EGAReferenceAlias{chapter:0III-es}{chapter:0III-fr}
\EGAReferenceAlias{chapter:0III.fr}{chapter:0III-fr}
\EGAReferenceAlias{chapter:0III.es}{chapter:0III-fr}
\EGAReferenceAlias{chapter:0III.eq}{chapter:0III-fr}
\EGAReferenceAlias{eq:chapter:0III}{chapter:0III-fr}
\EGAReferenceAlias{0.8}{section:0.8}
\EGAReferenceAlias{0.8-fr}{section:0.8}
\EGAReferenceAlias{0.8-es}{section:0.8}
\EGAReferenceAlias{0.8.fr}{section:0.8}
\EGAReferenceAlias{0.8.es}{section:0.8}
\EGAReferenceAlias{0.8.eq}{section:0.8}
\EGAReferenceAlias{eq:0.8}{section:0.8}
\EGAReferenceAlias{section:0.8-es}{section:0.8}
\EGAReferenceAlias{section:0.8.fr}{section:0.8}
\EGAReferenceAlias{section:0.8.es}{section:0.8}
\EGAReferenceAlias{section:0.8.eq}{section:0.8}
\EGAReferenceAlias{eq:section:0.8}{section:0.8}
\EGAReferenceAlias{0.8.1}{subsection:0.8.1}
\EGAReferenceAlias{0.8.1-fr}{subsection:0.8.1}
\EGAReferenceAlias{0.8.1-es}{subsection:0.8.1}
\EGAReferenceAlias{0.8.1.fr}{subsection:0.8.1}
\EGAReferenceAlias{0.8.1.es}{subsection:0.8.1}
\EGAReferenceAlias{0.8.1.eq}{subsection:0.8.1}
\EGAReferenceAlias{eq:0.8.1}{subsection:0.8.1}
\EGAReferenceAlias{subsection:0.8.1-es}{subsection:0.8.1}
\EGAReferenceAlias{subsection:0.8.1.fr}{subsection:0.8.1}
\EGAReferenceAlias{subsection:0.8.1.es}{subsection:0.8.1}
\EGAReferenceAlias{subsection:0.8.1.eq}{subsection:0.8.1}
\EGAReferenceAlias{eq:subsection:0.8.1}{subsection:0.8.1}
\EGAReferenceAlias{0.8.1.1-es}{0.8.1.1}
\EGAReferenceAlias{0.8.1.1.fr}{0.8.1.1}
\EGAReferenceAlias{0.8.1.1.es}{0.8.1.1}
\EGAReferenceAlias{0.8.1.1.eq}{0.8.1.1}
\EGAReferenceAlias{eq:0.8.1.1}{0.8.1.1}
\EGAReferenceAlias{0.8.1.2-es}{0.8.1.2}
\EGAReferenceAlias{0.8.1.2.fr}{0.8.1.2}
\EGAReferenceAlias{0.8.1.2.es}{0.8.1.2}
\EGAReferenceAlias{0.8.1.2.eq}{0.8.1.2}
\EGAReferenceAlias{eq:0.8.1.2}{0.8.1.2}
\EGAReferenceAlias{0.8.1.3-es}{0.8.1.3}
\EGAReferenceAlias{0.8.1.3.fr}{0.8.1.3}
\EGAReferenceAlias{0.8.1.3.es}{0.8.1.3}
\EGAReferenceAlias{0.8.1.3.eq}{0.8.1.3}
\EGAReferenceAlias{eq:0.8.1.3}{0.8.1.3}
\EGAReferenceAlias{0.8.1.3.1}{0.8.1.3.1-fr}
\EGAReferenceAlias{0.8.1.3.1-es}{0.8.1.3.1-fr}
\EGAReferenceAlias{0.8.1.3.1.fr}{0.8.1.3.1-fr}
\EGAReferenceAlias{0.8.1.3.1.es}{0.8.1.3.1-fr}
\EGAReferenceAlias{0.8.1.3.1.eq}{0.8.1.3.1-fr}
\EGAReferenceAlias{eq:0.8.1.3.1}{0.8.1.3.1-fr}
\EGAReferenceAlias{0.8.1.4-es}{0.8.1.4}
\EGAReferenceAlias{0.8.1.4.fr}{0.8.1.4}
\EGAReferenceAlias{0.8.1.4.es}{0.8.1.4}
\EGAReferenceAlias{0.8.1.4.eq}{0.8.1.4}
\EGAReferenceAlias{eq:0.8.1.4}{0.8.1.4}
\EGAReferenceAlias{0.8.1.5-es}{0.8.1.5}
\EGAReferenceAlias{0.8.1.5.fr}{0.8.1.5}
\EGAReferenceAlias{0.8.1.5.es}{0.8.1.5}
\EGAReferenceAlias{0.8.1.5.eq}{0.8.1.5}
\EGAReferenceAlias{eq:0.8.1.5}{0.8.1.5}
\EGAReferenceAlias{0.8.1.6}{0.8.1.6-fr}
\EGAReferenceAlias{0.8.1.6-es}{0.8.1.6-fr}
\EGAReferenceAlias{0.8.1.6.fr}{0.8.1.6-fr}
\EGAReferenceAlias{0.8.1.6.es}{0.8.1.6-fr}
\EGAReferenceAlias{0.8.1.6.eq}{0.8.1.6-fr}
\EGAReferenceAlias{eq:0.8.1.6}{0.8.1.6-fr}
\EGAReferenceAlias{0.8.1.7-es}{0.8.1.7}
\EGAReferenceAlias{0.8.1.7.fr}{0.8.1.7}
\EGAReferenceAlias{0.8.1.7.es}{0.8.1.7}
\EGAReferenceAlias{0.8.1.7.eq}{0.8.1.7}
\EGAReferenceAlias{eq:0.8.1.7}{0.8.1.7}
\EGAReferenceAlias{0.8.1.8-es}{0.8.1.8}
\EGAReferenceAlias{0.8.1.8.fr}{0.8.1.8}
\EGAReferenceAlias{0.8.1.8.es}{0.8.1.8}
\EGAReferenceAlias{0.8.1.8.eq}{0.8.1.8}
\EGAReferenceAlias{eq:0.8.1.8}{0.8.1.8}
\EGAReferenceAlias{0.8.1.9-es}{0.8.1.9}
\EGAReferenceAlias{0.8.1.9.fr}{0.8.1.9}
\EGAReferenceAlias{0.8.1.9.es}{0.8.1.9}
\EGAReferenceAlias{0.8.1.9.eq}{0.8.1.9}
\EGAReferenceAlias{eq:0.8.1.9}{0.8.1.9}
\EGAReferenceAlias{0.8.1.10-es}{0.8.1.10}
\EGAReferenceAlias{0.8.1.10.fr}{0.8.1.10}
\EGAReferenceAlias{0.8.1.10.es}{0.8.1.10}
\EGAReferenceAlias{0.8.1.10.eq}{0.8.1.10}
\EGAReferenceAlias{eq:0.8.1.10}{0.8.1.10}
\EGAReferenceAlias{0.8.1.11-es}{0.8.1.11}
\EGAReferenceAlias{0.8.1.11.fr}{0.8.1.11}
\EGAReferenceAlias{0.8.1.11.es}{0.8.1.11}
\EGAReferenceAlias{0.8.1.11.eq}{0.8.1.11}
\EGAReferenceAlias{eq:0.8.1.11}{0.8.1.11}
\EGAReferenceAlias{0.8.2}{subsection:0.8.2}
\EGAReferenceAlias{0.8.2-fr}{subsection:0.8.2}
\EGAReferenceAlias{0.8.2-es}{subsection:0.8.2}
\EGAReferenceAlias{0.8.2.fr}{subsection:0.8.2}
\EGAReferenceAlias{0.8.2.es}{subsection:0.8.2}
\EGAReferenceAlias{0.8.2.eq}{subsection:0.8.2}
\EGAReferenceAlias{eq:0.8.2}{subsection:0.8.2}
\EGAReferenceAlias{subsection:0.8.2-es}{subsection:0.8.2}
\EGAReferenceAlias{subsection:0.8.2.fr}{subsection:0.8.2}
\EGAReferenceAlias{subsection:0.8.2.es}{subsection:0.8.2}
\EGAReferenceAlias{subsection:0.8.2.eq}{subsection:0.8.2}
\EGAReferenceAlias{eq:subsection:0.8.2}{subsection:0.8.2}
\EGAReferenceAlias{0.8.2.1-es}{0.8.2.1}
\EGAReferenceAlias{0.8.2.1.fr}{0.8.2.1}
\EGAReferenceAlias{0.8.2.1.es}{0.8.2.1}
\EGAReferenceAlias{0.8.2.1.eq}{0.8.2.1}
\EGAReferenceAlias{eq:0.8.2.1}{0.8.2.1}
\EGAReferenceAlias{0.8.2.2-es}{0.8.2.2}
\EGAReferenceAlias{0.8.2.2.fr}{0.8.2.2}
\EGAReferenceAlias{0.8.2.2.es}{0.8.2.2}
\EGAReferenceAlias{0.8.2.2.eq}{0.8.2.2}
\EGAReferenceAlias{eq:0.8.2.2}{0.8.2.2}
\EGAReferenceAlias{0.8.2.3-es}{0.8.2.3}
\EGAReferenceAlias{0.8.2.3.fr}{0.8.2.3}
\EGAReferenceAlias{0.8.2.3.es}{0.8.2.3}
\EGAReferenceAlias{0.8.2.3.eq}{0.8.2.3}
\EGAReferenceAlias{eq:0.8.2.3}{0.8.2.3}
\EGAReferenceAlias{0.8.2.4-es}{0.8.2.4}
\EGAReferenceAlias{0.8.2.4.fr}{0.8.2.4}
\EGAReferenceAlias{0.8.2.4.es}{0.8.2.4}
\EGAReferenceAlias{0.8.2.4.eq}{0.8.2.4}
\EGAReferenceAlias{eq:0.8.2.4}{0.8.2.4}
\EGAReferenceAlias{0.8.2.5-es}{0.8.2.5}
\EGAReferenceAlias{0.8.2.5.fr}{0.8.2.5}
\EGAReferenceAlias{0.8.2.5.es}{0.8.2.5}
\EGAReferenceAlias{0.8.2.5.eq}{0.8.2.5}
\EGAReferenceAlias{eq:0.8.2.5}{0.8.2.5}
\EGAReferenceAlias{0.8.2.6-es}{0.8.2.6}
\EGAReferenceAlias{0.8.2.6.fr}{0.8.2.6}
\EGAReferenceAlias{0.8.2.6.es}{0.8.2.6}
\EGAReferenceAlias{0.8.2.6.eq}{0.8.2.6}
\EGAReferenceAlias{eq:0.8.2.6}{0.8.2.6}
\EGAReferenceAlias{0.8.2.7}{0.8.2.7-fr}
\EGAReferenceAlias{0.8.2.7-es}{0.8.2.7-fr}
\EGAReferenceAlias{0.8.2.7.fr}{0.8.2.7-fr}
\EGAReferenceAlias{0.8.2.7.es}{0.8.2.7-fr}
\EGAReferenceAlias{0.8.2.7.eq}{0.8.2.7-fr}
\EGAReferenceAlias{eq:0.8.2.7}{0.8.2.7-fr}
\EGAReferenceAlias{0.8.2.8-es}{0.8.2.8}
\EGAReferenceAlias{0.8.2.8.fr}{0.8.2.8}
\EGAReferenceAlias{0.8.2.8.es}{0.8.2.8}
\EGAReferenceAlias{0.8.2.8.eq}{0.8.2.8}
\EGAReferenceAlias{eq:0.8.2.8}{0.8.2.8}
\EGAReferenceAlias{0.9}{section:0.9}
\EGAReferenceAlias{0.9-fr}{section:0.9}
\EGAReferenceAlias{0.9-es}{section:0.9}
\EGAReferenceAlias{0.9.fr}{section:0.9}
\EGAReferenceAlias{0.9.es}{section:0.9}
\EGAReferenceAlias{0.9.eq}{section:0.9}
\EGAReferenceAlias{eq:0.9}{section:0.9}
\EGAReferenceAlias{section:0.9-es}{section:0.9}
\EGAReferenceAlias{section:0.9.fr}{section:0.9}
\EGAReferenceAlias{section:0.9.es}{section:0.9}
\EGAReferenceAlias{section:0.9.eq}{section:0.9}
\EGAReferenceAlias{eq:section:0.9}{section:0.9}
\EGAReferenceAlias{0.9.1}{subsection:0.9.1}
\EGAReferenceAlias{0.9.1-fr}{subsection:0.9.1}
\EGAReferenceAlias{0.9.1-es}{subsection:0.9.1}
\EGAReferenceAlias{0.9.1.fr}{subsection:0.9.1}
\EGAReferenceAlias{0.9.1.es}{subsection:0.9.1}
\EGAReferenceAlias{0.9.1.eq}{subsection:0.9.1}
\EGAReferenceAlias{eq:0.9.1}{subsection:0.9.1}
\EGAReferenceAlias{subsection:0.9.1-es}{subsection:0.9.1}
\EGAReferenceAlias{subsection:0.9.1.fr}{subsection:0.9.1}
\EGAReferenceAlias{subsection:0.9.1.es}{subsection:0.9.1}
\EGAReferenceAlias{subsection:0.9.1.eq}{subsection:0.9.1}
\EGAReferenceAlias{eq:subsection:0.9.1}{subsection:0.9.1}
\EGAReferenceAlias{0.9.1.1-es}{0.9.1.1}
\EGAReferenceAlias{0.9.1.1.fr}{0.9.1.1}
\EGAReferenceAlias{0.9.1.1.es}{0.9.1.1}
\EGAReferenceAlias{0.9.1.1.eq}{0.9.1.1}
\EGAReferenceAlias{eq:0.9.1.1}{0.9.1.1}
\EGAReferenceAlias{0.9.1.2-es}{0.9.1.2}
\EGAReferenceAlias{0.9.1.2.fr}{0.9.1.2}
\EGAReferenceAlias{0.9.1.2.es}{0.9.1.2}
\EGAReferenceAlias{0.9.1.2.eq}{0.9.1.2}
\EGAReferenceAlias{eq:0.9.1.2}{0.9.1.2}
\EGAReferenceAlias{0.9.1.3-es}{0.9.1.3}
\EGAReferenceAlias{0.9.1.3.fr}{0.9.1.3}
\EGAReferenceAlias{0.9.1.3.es}{0.9.1.3}
\EGAReferenceAlias{0.9.1.3.eq}{0.9.1.3}
\EGAReferenceAlias{eq:0.9.1.3}{0.9.1.3}
\EGAReferenceAlias{0.9.1.4-es}{0.9.1.4}
\EGAReferenceAlias{0.9.1.4.fr}{0.9.1.4}
\EGAReferenceAlias{0.9.1.4.es}{0.9.1.4}
\EGAReferenceAlias{0.9.1.4.eq}{0.9.1.4}
\EGAReferenceAlias{eq:0.9.1.4}{0.9.1.4}
\EGAReferenceAlias{0.9.1.5-es}{0.9.1.5}
\EGAReferenceAlias{0.9.1.5.fr}{0.9.1.5}
\EGAReferenceAlias{0.9.1.5.es}{0.9.1.5}
\EGAReferenceAlias{0.9.1.5.eq}{0.9.1.5}
\EGAReferenceAlias{eq:0.9.1.5}{0.9.1.5}
\EGAReferenceAlias{0.9.1.6-es}{0.9.1.6}
\EGAReferenceAlias{0.9.1.6.fr}{0.9.1.6}
\EGAReferenceAlias{0.9.1.6.es}{0.9.1.6}
\EGAReferenceAlias{0.9.1.6.eq}{0.9.1.6}
\EGAReferenceAlias{eq:0.9.1.6}{0.9.1.6}
\EGAReferenceAlias{0.9.1.7-es}{0.9.1.7}
\EGAReferenceAlias{0.9.1.7.fr}{0.9.1.7}
\EGAReferenceAlias{0.9.1.7.es}{0.9.1.7}
\EGAReferenceAlias{0.9.1.7.eq}{0.9.1.7}
\EGAReferenceAlias{eq:0.9.1.7}{0.9.1.7}
\EGAReferenceAlias{0.9.1.8-es}{0.9.1.8}
\EGAReferenceAlias{0.9.1.8.fr}{0.9.1.8}
\EGAReferenceAlias{0.9.1.8.es}{0.9.1.8}
\EGAReferenceAlias{0.9.1.8.eq}{0.9.1.8}
\EGAReferenceAlias{eq:0.9.1.8}{0.9.1.8}
\EGAReferenceAlias{0.9.1.9-es}{0.9.1.9}
\EGAReferenceAlias{0.9.1.9.fr}{0.9.1.9}
\EGAReferenceAlias{0.9.1.9.es}{0.9.1.9}
\EGAReferenceAlias{0.9.1.9.eq}{0.9.1.9}
\EGAReferenceAlias{eq:0.9.1.9}{0.9.1.9}
\EGAReferenceAlias{0.9.1.10-es}{0.9.1.10}
\EGAReferenceAlias{0.9.1.10.fr}{0.9.1.10}
\EGAReferenceAlias{0.9.1.10.es}{0.9.1.10}
\EGAReferenceAlias{0.9.1.10.eq}{0.9.1.10}
\EGAReferenceAlias{eq:0.9.1.10}{0.9.1.10}
\EGAReferenceAlias{0.9.1.11-es}{0.9.1.11}
\EGAReferenceAlias{0.9.1.11.fr}{0.9.1.11}
\EGAReferenceAlias{0.9.1.11.es}{0.9.1.11}
\EGAReferenceAlias{0.9.1.11.eq}{0.9.1.11}
\EGAReferenceAlias{eq:0.9.1.11}{0.9.1.11}
\EGAReferenceAlias{0.9.1.12-es}{0.9.1.12}
\EGAReferenceAlias{0.9.1.12.fr}{0.9.1.12}
\EGAReferenceAlias{0.9.1.12.es}{0.9.1.12}
\EGAReferenceAlias{0.9.1.12.eq}{0.9.1.12}
\EGAReferenceAlias{eq:0.9.1.12}{0.9.1.12}
\EGAReferenceAlias{0.9.1.13-es}{0.9.1.13}
\EGAReferenceAlias{0.9.1.13.fr}{0.9.1.13}
\EGAReferenceAlias{0.9.1.13.es}{0.9.1.13}
\EGAReferenceAlias{0.9.1.13.eq}{0.9.1.13}
\EGAReferenceAlias{eq:0.9.1.13}{0.9.1.13}
\EGAReferenceAlias{0.9.2}{subsection:0.9.2}
\EGAReferenceAlias{0.9.2-fr}{subsection:0.9.2}
\EGAReferenceAlias{0.9.2-es}{subsection:0.9.2}
\EGAReferenceAlias{0.9.2.fr}{subsection:0.9.2}
\EGAReferenceAlias{0.9.2.es}{subsection:0.9.2}
\EGAReferenceAlias{0.9.2.eq}{subsection:0.9.2}
\EGAReferenceAlias{eq:0.9.2}{subsection:0.9.2}
\EGAReferenceAlias{subsection:0.9.2-es}{subsection:0.9.2}
\EGAReferenceAlias{subsection:0.9.2.fr}{subsection:0.9.2}
\EGAReferenceAlias{subsection:0.9.2.es}{subsection:0.9.2}
\EGAReferenceAlias{subsection:0.9.2.eq}{subsection:0.9.2}
\EGAReferenceAlias{eq:subsection:0.9.2}{subsection:0.9.2}
\EGAReferenceAlias{0.9.2.1-es}{0.9.2.1}
\EGAReferenceAlias{0.9.2.1.fr}{0.9.2.1}
\EGAReferenceAlias{0.9.2.1.es}{0.9.2.1}
\EGAReferenceAlias{0.9.2.1.eq}{0.9.2.1}
\EGAReferenceAlias{eq:0.9.2.1}{0.9.2.1}
\EGAReferenceAlias{0.9.2.2-es}{0.9.2.2}
\EGAReferenceAlias{0.9.2.2.fr}{0.9.2.2}
\EGAReferenceAlias{0.9.2.2.es}{0.9.2.2}
\EGAReferenceAlias{0.9.2.2.eq}{0.9.2.2}
\EGAReferenceAlias{eq:0.9.2.2}{0.9.2.2}
\EGAReferenceAlias{0.9.2.3-es}{0.9.2.3}
\EGAReferenceAlias{0.9.2.3.fr}{0.9.2.3}
\EGAReferenceAlias{0.9.2.3.es}{0.9.2.3}
\EGAReferenceAlias{0.9.2.3.eq}{0.9.2.3}
\EGAReferenceAlias{eq:0.9.2.3}{0.9.2.3}
\EGAReferenceAlias{0.9.2.4-es}{0.9.2.4}
\EGAReferenceAlias{0.9.2.4.fr}{0.9.2.4}
\EGAReferenceAlias{0.9.2.4.es}{0.9.2.4}
\EGAReferenceAlias{0.9.2.4.eq}{0.9.2.4}
\EGAReferenceAlias{eq:0.9.2.4}{0.9.2.4}
\EGAReferenceAlias{0.9.2.5-es}{0.9.2.5}
\EGAReferenceAlias{0.9.2.5.fr}{0.9.2.5}
\EGAReferenceAlias{0.9.2.5.es}{0.9.2.5}
\EGAReferenceAlias{0.9.2.5.eq}{0.9.2.5}
\EGAReferenceAlias{eq:0.9.2.5}{0.9.2.5}
\EGAReferenceAlias{0.9.2.6-es}{0.9.2.6}
\EGAReferenceAlias{0.9.2.6.fr}{0.9.2.6}
\EGAReferenceAlias{0.9.2.6.es}{0.9.2.6}
\EGAReferenceAlias{0.9.2.6.eq}{0.9.2.6}
\EGAReferenceAlias{eq:0.9.2.6}{0.9.2.6}
\EGAReferenceAlias{0.9.3}{subsection:0.9.3}
\EGAReferenceAlias{0.9.3-fr}{subsection:0.9.3}
\EGAReferenceAlias{0.9.3-es}{subsection:0.9.3}
\EGAReferenceAlias{0.9.3.fr}{subsection:0.9.3}
\EGAReferenceAlias{0.9.3.es}{subsection:0.9.3}
\EGAReferenceAlias{0.9.3.eq}{subsection:0.9.3}
\EGAReferenceAlias{eq:0.9.3}{subsection:0.9.3}
\EGAReferenceAlias{subsection:0.9.3-es}{subsection:0.9.3}
\EGAReferenceAlias{subsection:0.9.3.fr}{subsection:0.9.3}
\EGAReferenceAlias{subsection:0.9.3.es}{subsection:0.9.3}
\EGAReferenceAlias{subsection:0.9.3.eq}{subsection:0.9.3}
\EGAReferenceAlias{eq:subsection:0.9.3}{subsection:0.9.3}
\EGAReferenceAlias{0.9.3.1-es}{0.9.3.1}
\EGAReferenceAlias{0.9.3.1.fr}{0.9.3.1}
\EGAReferenceAlias{0.9.3.1.es}{0.9.3.1}
\EGAReferenceAlias{0.9.3.1.eq}{0.9.3.1}
\EGAReferenceAlias{eq:0.9.3.1}{0.9.3.1}
\EGAReferenceAlias{0.9.3.2-es}{0.9.3.2}
\EGAReferenceAlias{0.9.3.2.fr}{0.9.3.2}
\EGAReferenceAlias{0.9.3.2.es}{0.9.3.2}
\EGAReferenceAlias{0.9.3.2.eq}{0.9.3.2}
\EGAReferenceAlias{eq:0.9.3.2}{0.9.3.2}
\EGAReferenceAlias{0.9.3.3-es}{0.9.3.3}
\EGAReferenceAlias{0.9.3.3.fr}{0.9.3.3}
\EGAReferenceAlias{0.9.3.3.es}{0.9.3.3}
\EGAReferenceAlias{0.9.3.3.eq}{0.9.3.3}
\EGAReferenceAlias{eq:0.9.3.3}{0.9.3.3}
\EGAReferenceAlias{0.9.3.4-es}{0.9.3.4}
\EGAReferenceAlias{0.9.3.4.fr}{0.9.3.4}
\EGAReferenceAlias{0.9.3.4.es}{0.9.3.4}
\EGAReferenceAlias{0.9.3.4.eq}{0.9.3.4}
\EGAReferenceAlias{eq:0.9.3.4}{0.9.3.4}
\EGAReferenceAlias{0.10}{section:0.10}
\EGAReferenceAlias{0.10-fr}{section:0.10}
\EGAReferenceAlias{0.10-es}{section:0.10}
\EGAReferenceAlias{0.10.fr}{section:0.10}
\EGAReferenceAlias{0.10.es}{section:0.10}
\EGAReferenceAlias{0.10.eq}{section:0.10}
\EGAReferenceAlias{eq:0.10}{section:0.10}
\EGAReferenceAlias{section:0.10-es}{section:0.10}
\EGAReferenceAlias{section:0.10.fr}{section:0.10}
\EGAReferenceAlias{section:0.10.es}{section:0.10}
\EGAReferenceAlias{section:0.10.eq}{section:0.10}
\EGAReferenceAlias{eq:section:0.10}{section:0.10}
\EGAReferenceAlias{0.10.1}{subsection:0.10.1}
\EGAReferenceAlias{0.10.1-fr}{subsection:0.10.1}
\EGAReferenceAlias{0.10.1-es}{subsection:0.10.1}
\EGAReferenceAlias{0.10.1.fr}{subsection:0.10.1}
\EGAReferenceAlias{0.10.1.es}{subsection:0.10.1}
\EGAReferenceAlias{0.10.1.eq}{subsection:0.10.1}
\EGAReferenceAlias{eq:0.10.1}{subsection:0.10.1}
\EGAReferenceAlias{subsection:0.10.1-es}{subsection:0.10.1}
\EGAReferenceAlias{subsection:0.10.1.fr}{subsection:0.10.1}
\EGAReferenceAlias{subsection:0.10.1.es}{subsection:0.10.1}
\EGAReferenceAlias{subsection:0.10.1.eq}{subsection:0.10.1}
\EGAReferenceAlias{eq:subsection:0.10.1}{subsection:0.10.1}
\EGAReferenceAlias{0.10.1.1-es}{0.10.1.1}
\EGAReferenceAlias{0.10.1.1.fr}{0.10.1.1}
\EGAReferenceAlias{0.10.1.1.es}{0.10.1.1}
\EGAReferenceAlias{0.10.1.1.eq}{0.10.1.1}
\EGAReferenceAlias{eq:0.10.1.1}{0.10.1.1}
\EGAReferenceAlias{0.10.1.1.1}{0.10.1.1.1-fr}
\EGAReferenceAlias{0.10.1.1.1-es}{0.10.1.1.1-fr}
\EGAReferenceAlias{0.10.1.1.1.fr}{0.10.1.1.1-fr}
\EGAReferenceAlias{0.10.1.1.1.es}{0.10.1.1.1-fr}
\EGAReferenceAlias{0.10.1.1.1.eq}{0.10.1.1.1-fr}
\EGAReferenceAlias{eq:0.10.1.1.1}{0.10.1.1.1-fr}
\EGAReferenceAlias{0.10.1.1.2}{0.10.1.1.2-fr}
\EGAReferenceAlias{0.10.1.1.2-es}{0.10.1.1.2-fr}
\EGAReferenceAlias{0.10.1.1.2.fr}{0.10.1.1.2-fr}
\EGAReferenceAlias{0.10.1.1.2.es}{0.10.1.1.2-fr}
\EGAReferenceAlias{0.10.1.1.2.eq}{0.10.1.1.2-fr}
\EGAReferenceAlias{eq:0.10.1.1.2}{0.10.1.1.2-fr}
\EGAReferenceAlias{0.10.1.2-es}{0.10.1.2}
\EGAReferenceAlias{0.10.1.2.fr}{0.10.1.2}
\EGAReferenceAlias{0.10.1.2.es}{0.10.1.2}
\EGAReferenceAlias{0.10.1.2.eq}{0.10.1.2}
\EGAReferenceAlias{eq:0.10.1.2}{0.10.1.2}
\EGAReferenceAlias{0.10.1.3-es}{0.10.1.3}
\EGAReferenceAlias{0.10.1.3.fr}{0.10.1.3}
\EGAReferenceAlias{0.10.1.3.es}{0.10.1.3}
\EGAReferenceAlias{0.10.1.3.eq}{0.10.1.3}
\EGAReferenceAlias{eq:0.10.1.3}{0.10.1.3}
\EGAReferenceAlias{0.10.2}{subsection:0.10.2}
\EGAReferenceAlias{0.10.2-fr}{subsection:0.10.2}
\EGAReferenceAlias{0.10.2-es}{subsection:0.10.2}
\EGAReferenceAlias{0.10.2.fr}{subsection:0.10.2}
\EGAReferenceAlias{0.10.2.es}{subsection:0.10.2}
\EGAReferenceAlias{0.10.2.eq}{subsection:0.10.2}
\EGAReferenceAlias{eq:0.10.2}{subsection:0.10.2}
\EGAReferenceAlias{subsection:0.10.2-es}{subsection:0.10.2}
\EGAReferenceAlias{subsection:0.10.2.fr}{subsection:0.10.2}
\EGAReferenceAlias{subsection:0.10.2.es}{subsection:0.10.2}
\EGAReferenceAlias{subsection:0.10.2.eq}{subsection:0.10.2}
\EGAReferenceAlias{eq:subsection:0.10.2}{subsection:0.10.2}
\EGAReferenceAlias{0.10.2.1}{0.10.2.1-fr}
\EGAReferenceAlias{0.10.2.1-es}{0.10.2.1-fr}
\EGAReferenceAlias{0.10.2.1.fr}{0.10.2.1-fr}
\EGAReferenceAlias{0.10.2.1.es}{0.10.2.1-fr}
\EGAReferenceAlias{0.10.2.1.eq}{0.10.2.1-fr}
\EGAReferenceAlias{eq:0.10.2.1}{0.10.2.1-fr}
\EGAReferenceAlias{0.10.2.2-es}{0.10.2.2}
\EGAReferenceAlias{0.10.2.2.fr}{0.10.2.2}
\EGAReferenceAlias{0.10.2.2.es}{0.10.2.2}
\EGAReferenceAlias{0.10.2.2.eq}{0.10.2.2}
\EGAReferenceAlias{eq:0.10.2.2}{0.10.2.2}
\EGAReferenceAlias{0.10.2.3-es}{0.10.2.3}
\EGAReferenceAlias{0.10.2.3.fr}{0.10.2.3}
\EGAReferenceAlias{0.10.2.3.es}{0.10.2.3}
\EGAReferenceAlias{0.10.2.3.eq}{0.10.2.3}
\EGAReferenceAlias{eq:0.10.2.3}{0.10.2.3}
\EGAReferenceAlias{0.10.2.4-es}{0.10.2.4}
\EGAReferenceAlias{0.10.2.4.fr}{0.10.2.4}
\EGAReferenceAlias{0.10.2.4.es}{0.10.2.4}
\EGAReferenceAlias{0.10.2.4.eq}{0.10.2.4}
\EGAReferenceAlias{eq:0.10.2.4}{0.10.2.4}
\EGAReferenceAlias{0.10.2.5-es}{0.10.2.5}
\EGAReferenceAlias{0.10.2.5.fr}{0.10.2.5}
\EGAReferenceAlias{0.10.2.5.es}{0.10.2.5}
\EGAReferenceAlias{0.10.2.5.eq}{0.10.2.5}
\EGAReferenceAlias{eq:0.10.2.5}{0.10.2.5}
\EGAReferenceAlias{0.10.2.6-es}{0.10.2.6}
\EGAReferenceAlias{0.10.2.6.fr}{0.10.2.6}
\EGAReferenceAlias{0.10.2.6.es}{0.10.2.6}
\EGAReferenceAlias{0.10.2.6.eq}{0.10.2.6}
\EGAReferenceAlias{eq:0.10.2.6}{0.10.2.6}
\EGAReferenceAlias{0.10.2.7-es}{0.10.2.7}
\EGAReferenceAlias{0.10.2.7.fr}{0.10.2.7}
\EGAReferenceAlias{0.10.2.7.es}{0.10.2.7}
\EGAReferenceAlias{0.10.2.7.eq}{0.10.2.7}
\EGAReferenceAlias{eq:0.10.2.7}{0.10.2.7}
\EGAReferenceAlias{0.10.2.8-es}{0.10.2.8}
\EGAReferenceAlias{0.10.2.8.fr}{0.10.2.8}
\EGAReferenceAlias{0.10.2.8.es}{0.10.2.8}
\EGAReferenceAlias{0.10.2.8.eq}{0.10.2.8}
\EGAReferenceAlias{eq:0.10.2.8}{0.10.2.8}
\EGAReferenceAlias{0.10.2.8.1}{0.10.2.8.1-fr}
\EGAReferenceAlias{0.10.2.8.1-es}{0.10.2.8.1-fr}
\EGAReferenceAlias{0.10.2.8.1.fr}{0.10.2.8.1-fr}
\EGAReferenceAlias{0.10.2.8.1.es}{0.10.2.8.1-fr}
\EGAReferenceAlias{0.10.2.8.1.eq}{0.10.2.8.1-fr}
\EGAReferenceAlias{eq:0.10.2.8.1}{0.10.2.8.1-fr}
\EGAReferenceAlias{0.10.3}{subsection:0.10.3}
\EGAReferenceAlias{0.10.3-fr}{subsection:0.10.3}
\EGAReferenceAlias{0.10.3-es}{subsection:0.10.3}
\EGAReferenceAlias{0.10.3.fr}{subsection:0.10.3}
\EGAReferenceAlias{0.10.3.es}{subsection:0.10.3}
\EGAReferenceAlias{0.10.3.eq}{subsection:0.10.3}
\EGAReferenceAlias{eq:0.10.3}{subsection:0.10.3}
\EGAReferenceAlias{subsection:0.10.3-es}{subsection:0.10.3}
\EGAReferenceAlias{subsection:0.10.3.fr}{subsection:0.10.3}
\EGAReferenceAlias{subsection:0.10.3.es}{subsection:0.10.3}
\EGAReferenceAlias{subsection:0.10.3.eq}{subsection:0.10.3}
\EGAReferenceAlias{eq:subsection:0.10.3}{subsection:0.10.3}
\EGAReferenceAlias{0.10.3.1-es}{0.10.3.1}
\EGAReferenceAlias{0.10.3.1.fr}{0.10.3.1}
\EGAReferenceAlias{0.10.3.1.es}{0.10.3.1}
\EGAReferenceAlias{0.10.3.1.eq}{0.10.3.1}
\EGAReferenceAlias{eq:0.10.3.1}{0.10.3.1}
\EGAReferenceAlias{0.10.3.1.1-es}{0.10.3.1.1}
\EGAReferenceAlias{0.10.3.1.1.fr}{0.10.3.1.1}
\EGAReferenceAlias{0.10.3.1.1.es}{0.10.3.1.1}
\EGAReferenceAlias{0.10.3.1.1.eq}{0.10.3.1.1}
\EGAReferenceAlias{eq:0.10.3.1.1}{0.10.3.1.1}
\EGAReferenceAlias{0.10.3.1.2}{0.10.3.1.2-fr}
\EGAReferenceAlias{0.10.3.1.2-es}{0.10.3.1.2-fr}
\EGAReferenceAlias{0.10.3.1.2.fr}{0.10.3.1.2-fr}
\EGAReferenceAlias{0.10.3.1.2.es}{0.10.3.1.2-fr}
\EGAReferenceAlias{0.10.3.1.2.eq}{0.10.3.1.2-fr}
\EGAReferenceAlias{eq:0.10.3.1.2}{0.10.3.1.2-fr}
\EGAReferenceAlias{0.10.3.1.3-es}{0.10.3.1.3}
\EGAReferenceAlias{0.10.3.1.3.fr}{0.10.3.1.3}
\EGAReferenceAlias{0.10.3.1.3.es}{0.10.3.1.3}
\EGAReferenceAlias{0.10.3.1.3.eq}{0.10.3.1.3}
\EGAReferenceAlias{eq:0.10.3.1.3}{0.10.3.1.3}
\EGAReferenceAlias{0.10.3.1.4-es}{0.10.3.1.4}
\EGAReferenceAlias{0.10.3.1.4.fr}{0.10.3.1.4}
\EGAReferenceAlias{0.10.3.1.4.es}{0.10.3.1.4}
\EGAReferenceAlias{0.10.3.1.4.eq}{0.10.3.1.4}
\EGAReferenceAlias{eq:0.10.3.1.4}{0.10.3.1.4}
\EGAReferenceAlias{0.10.3.1.5}{0.10.3.1.5-fr}
\EGAReferenceAlias{0.10.3.1.5-es}{0.10.3.1.5-fr}
\EGAReferenceAlias{0.10.3.1.5.fr}{0.10.3.1.5-fr}
\EGAReferenceAlias{0.10.3.1.5.es}{0.10.3.1.5-fr}
\EGAReferenceAlias{0.10.3.1.5.eq}{0.10.3.1.5-fr}
\EGAReferenceAlias{eq:0.10.3.1.5}{0.10.3.1.5-fr}
\EGAReferenceAlias{0.10.3.2-es}{0.10.3.2}
\EGAReferenceAlias{0.10.3.2.fr}{0.10.3.2}
\EGAReferenceAlias{0.10.3.2.es}{0.10.3.2}
\EGAReferenceAlias{0.10.3.2.eq}{0.10.3.2}
\EGAReferenceAlias{eq:0.10.3.2}{0.10.3.2}
\EGAReferenceAlias{0.11}{section:0.11}
\EGAReferenceAlias{0.11-fr}{section:0.11}
\EGAReferenceAlias{0.11-es}{section:0.11}
\EGAReferenceAlias{0.11.fr}{section:0.11}
\EGAReferenceAlias{0.11.es}{section:0.11}
\EGAReferenceAlias{0.11.eq}{section:0.11}
\EGAReferenceAlias{eq:0.11}{section:0.11}
\EGAReferenceAlias{section:0.11-es}{section:0.11}
\EGAReferenceAlias{section:0.11.fr}{section:0.11}
\EGAReferenceAlias{section:0.11.es}{section:0.11}
\EGAReferenceAlias{section:0.11.eq}{section:0.11}
\EGAReferenceAlias{eq:section:0.11}{section:0.11}
\EGAReferenceAlias{0.11.1}{subsection:0.11.1}
\EGAReferenceAlias{0.11.1-fr}{subsection:0.11.1}
\EGAReferenceAlias{0.11.1-es}{subsection:0.11.1}
\EGAReferenceAlias{0.11.1.fr}{subsection:0.11.1}
\EGAReferenceAlias{0.11.1.es}{subsection:0.11.1}
\EGAReferenceAlias{0.11.1.eq}{subsection:0.11.1}
\EGAReferenceAlias{eq:0.11.1}{subsection:0.11.1}
\EGAReferenceAlias{subsection:0.11.1-es}{subsection:0.11.1}
\EGAReferenceAlias{subsection:0.11.1.fr}{subsection:0.11.1}
\EGAReferenceAlias{subsection:0.11.1.es}{subsection:0.11.1}
\EGAReferenceAlias{subsection:0.11.1.eq}{subsection:0.11.1}
\EGAReferenceAlias{eq:subsection:0.11.1}{subsection:0.11.1}
\EGAReferenceAlias{0.11.1.1}{0.11.1.1-fr}
\EGAReferenceAlias{0.11.1.1-es}{0.11.1.1-fr}
\EGAReferenceAlias{0.11.1.1.fr}{0.11.1.1-fr}
\EGAReferenceAlias{0.11.1.1.es}{0.11.1.1-fr}
\EGAReferenceAlias{0.11.1.1.eq}{0.11.1.1-fr}
\EGAReferenceAlias{eq:0.11.1.1}{0.11.1.1-fr}
\EGAReferenceAlias{0.11.1.1.1}{0.11.1.1.1-fr}
\EGAReferenceAlias{0.11.1.1.1-es}{0.11.1.1.1-fr}
\EGAReferenceAlias{0.11.1.1.1.fr}{0.11.1.1.1-fr}
\EGAReferenceAlias{0.11.1.1.1.es}{0.11.1.1.1-fr}
\EGAReferenceAlias{0.11.1.1.1.eq}{0.11.1.1.1-fr}
\EGAReferenceAlias{eq:0.11.1.1.1}{0.11.1.1.1-fr}
\EGAReferenceAlias{0.11.1.1.2}{0.11.1.1.2-fr}
\EGAReferenceAlias{0.11.1.1.2-es}{0.11.1.1.2-fr}
\EGAReferenceAlias{0.11.1.1.2.fr}{0.11.1.1.2-fr}
\EGAReferenceAlias{0.11.1.1.2.es}{0.11.1.1.2-fr}
\EGAReferenceAlias{0.11.1.1.2.eq}{0.11.1.1.2-fr}
\EGAReferenceAlias{eq:0.11.1.1.2}{0.11.1.1.2-fr}
\EGAReferenceAlias{0.11.1.1.3}{0.11.1.1.3-fr}
\EGAReferenceAlias{0.11.1.1.3-es}{0.11.1.1.3-fr}
\EGAReferenceAlias{0.11.1.1.3.fr}{0.11.1.1.3-fr}
\EGAReferenceAlias{0.11.1.1.3.es}{0.11.1.1.3-fr}
\EGAReferenceAlias{0.11.1.1.3.eq}{0.11.1.1.3-fr}
\EGAReferenceAlias{eq:0.11.1.1.3}{0.11.1.1.3-fr}
\EGAReferenceAlias{0.11.1.2-es}{0.11.1.2}
\EGAReferenceAlias{0.11.1.2.fr}{0.11.1.2}
\EGAReferenceAlias{0.11.1.2.es}{0.11.1.2}
\EGAReferenceAlias{0.11.1.2.eq}{0.11.1.2}
\EGAReferenceAlias{eq:0.11.1.2}{0.11.1.2}
\EGAReferenceAlias{0.11.1.3-es}{0.11.1.3}
\EGAReferenceAlias{0.11.1.3.fr}{0.11.1.3}
\EGAReferenceAlias{0.11.1.3.es}{0.11.1.3}
\EGAReferenceAlias{0.11.1.3.eq}{0.11.1.3}
\EGAReferenceAlias{eq:0.11.1.3}{0.11.1.3}
\EGAReferenceAlias{0.11.1.4-es}{0.11.1.4}
\EGAReferenceAlias{0.11.1.4.fr}{0.11.1.4}
\EGAReferenceAlias{0.11.1.4.es}{0.11.1.4}
\EGAReferenceAlias{0.11.1.4.eq}{0.11.1.4}
\EGAReferenceAlias{eq:0.11.1.4}{0.11.1.4}
\EGAReferenceAlias{0.11.1.5-es}{0.11.1.5}
\EGAReferenceAlias{0.11.1.5.fr}{0.11.1.5}
\EGAReferenceAlias{0.11.1.5.es}{0.11.1.5}
\EGAReferenceAlias{0.11.1.5.eq}{0.11.1.5}
\EGAReferenceAlias{eq:0.11.1.5}{0.11.1.5}
\EGAReferenceAlias{0.11.1.6}{0.11.1.6-fr}
\EGAReferenceAlias{0.11.1.6-es}{0.11.1.6-fr}
\EGAReferenceAlias{0.11.1.6.fr}{0.11.1.6-fr}
\EGAReferenceAlias{0.11.1.6.es}{0.11.1.6-fr}
\EGAReferenceAlias{0.11.1.6.eq}{0.11.1.6-fr}
\EGAReferenceAlias{eq:0.11.1.6}{0.11.1.6-fr}
\EGAReferenceAlias{0.11.1.7-es}{0.11.1.7}
\EGAReferenceAlias{0.11.1.7.fr}{0.11.1.7}
\EGAReferenceAlias{0.11.1.7.es}{0.11.1.7}
\EGAReferenceAlias{0.11.1.7.eq}{0.11.1.7}
\EGAReferenceAlias{eq:0.11.1.7}{0.11.1.7}
\EGAReferenceAlias{0.11.1.8-es}{0.11.1.8}
\EGAReferenceAlias{0.11.1.8.fr}{0.11.1.8}
\EGAReferenceAlias{0.11.1.8.es}{0.11.1.8}
\EGAReferenceAlias{0.11.1.8.eq}{0.11.1.8}
\EGAReferenceAlias{eq:0.11.1.8}{0.11.1.8}
\EGAReferenceAlias{0.11.1.9}{0.11.1.9-fr}
\EGAReferenceAlias{0.11.1.9-es}{0.11.1.9-fr}
\EGAReferenceAlias{0.11.1.9.fr}{0.11.1.9-fr}
\EGAReferenceAlias{0.11.1.9.es}{0.11.1.9-fr}
\EGAReferenceAlias{0.11.1.9.eq}{0.11.1.9-fr}
\EGAReferenceAlias{eq:0.11.1.9}{0.11.1.9-fr}
\EGAReferenceAlias{0.11.1.10-es}{0.11.1.10}
\EGAReferenceAlias{0.11.1.10.fr}{0.11.1.10}
\EGAReferenceAlias{0.11.1.10.es}{0.11.1.10}
\EGAReferenceAlias{0.11.1.10.eq}{0.11.1.10}
\EGAReferenceAlias{eq:0.11.1.10}{0.11.1.10}
\EGAReferenceAlias{0.11.2}{subsection:0.11.2}
\EGAReferenceAlias{0.11.2-fr}{subsection:0.11.2}
\EGAReferenceAlias{0.11.2-es}{subsection:0.11.2}
\EGAReferenceAlias{0.11.2.fr}{subsection:0.11.2}
\EGAReferenceAlias{0.11.2.es}{subsection:0.11.2}
\EGAReferenceAlias{0.11.2.eq}{subsection:0.11.2}
\EGAReferenceAlias{eq:0.11.2}{subsection:0.11.2}
\EGAReferenceAlias{subsection:0.11.2-es}{subsection:0.11.2}
\EGAReferenceAlias{subsection:0.11.2.fr}{subsection:0.11.2}
\EGAReferenceAlias{subsection:0.11.2.es}{subsection:0.11.2}
\EGAReferenceAlias{subsection:0.11.2.eq}{subsection:0.11.2}
\EGAReferenceAlias{eq:subsection:0.11.2}{subsection:0.11.2}
\EGAReferenceAlias{0.11.2.1-es}{0.11.2.1}
\EGAReferenceAlias{0.11.2.1.fr}{0.11.2.1}
\EGAReferenceAlias{0.11.2.1.es}{0.11.2.1}
\EGAReferenceAlias{0.11.2.1.eq}{0.11.2.1}
\EGAReferenceAlias{eq:0.11.2.1}{0.11.2.1}
\EGAReferenceAlias{0.11.2.2-es}{0.11.2.2}
\EGAReferenceAlias{0.11.2.2.fr}{0.11.2.2}
\EGAReferenceAlias{0.11.2.2.es}{0.11.2.2}
\EGAReferenceAlias{0.11.2.2.eq}{0.11.2.2}
\EGAReferenceAlias{eq:0.11.2.2}{0.11.2.2}
\EGAReferenceAlias{0.11.2.3-es}{0.11.2.3}
\EGAReferenceAlias{0.11.2.3.fr}{0.11.2.3}
\EGAReferenceAlias{0.11.2.3.es}{0.11.2.3}
\EGAReferenceAlias{0.11.2.3.eq}{0.11.2.3}
\EGAReferenceAlias{eq:0.11.2.3}{0.11.2.3}
\EGAReferenceAlias{0.11.2.4}{0.11.2.4-fr}
\EGAReferenceAlias{0.11.2.4-es}{0.11.2.4-fr}
\EGAReferenceAlias{0.11.2.4.fr}{0.11.2.4-fr}
\EGAReferenceAlias{0.11.2.4.es}{0.11.2.4-fr}
\EGAReferenceAlias{0.11.2.4.eq}{0.11.2.4-fr}
\EGAReferenceAlias{eq:0.11.2.4}{0.11.2.4-fr}
\EGAReferenceAlias{0.11.3}{subsection:0.11.3}
\EGAReferenceAlias{0.11.3-fr}{subsection:0.11.3}
\EGAReferenceAlias{0.11.3-es}{subsection:0.11.3}
\EGAReferenceAlias{0.11.3.fr}{subsection:0.11.3}
\EGAReferenceAlias{0.11.3.es}{subsection:0.11.3}
\EGAReferenceAlias{0.11.3.eq}{subsection:0.11.3}
\EGAReferenceAlias{eq:0.11.3}{subsection:0.11.3}
\EGAReferenceAlias{subsection:0.11.3-es}{subsection:0.11.3}
\EGAReferenceAlias{subsection:0.11.3.fr}{subsection:0.11.3}
\EGAReferenceAlias{subsection:0.11.3.es}{subsection:0.11.3}
\EGAReferenceAlias{subsection:0.11.3.eq}{subsection:0.11.3}
\EGAReferenceAlias{eq:subsection:0.11.3}{subsection:0.11.3}
\EGAReferenceAlias{0.11.3.1-es}{0.11.3.1}
\EGAReferenceAlias{0.11.3.1.fr}{0.11.3.1}
\EGAReferenceAlias{0.11.3.1.es}{0.11.3.1}
\EGAReferenceAlias{0.11.3.1.eq}{0.11.3.1}
\EGAReferenceAlias{eq:0.11.3.1}{0.11.3.1}
\EGAReferenceAlias{0.11.3.2-es}{0.11.3.2}
\EGAReferenceAlias{0.11.3.2.fr}{0.11.3.2}
\EGAReferenceAlias{0.11.3.2.es}{0.11.3.2}
\EGAReferenceAlias{0.11.3.2.eq}{0.11.3.2}
\EGAReferenceAlias{eq:0.11.3.2}{0.11.3.2}
\EGAReferenceAlias{0.11.3.2.1}{0.11.3.2.1-fr}
\EGAReferenceAlias{0.11.3.2.1-es}{0.11.3.2.1-fr}
\EGAReferenceAlias{0.11.3.2.1.fr}{0.11.3.2.1-fr}
\EGAReferenceAlias{0.11.3.2.1.es}{0.11.3.2.1-fr}
\EGAReferenceAlias{0.11.3.2.1.eq}{0.11.3.2.1-fr}
\EGAReferenceAlias{eq:0.11.3.2.1}{0.11.3.2.1-fr}
\EGAReferenceAlias{0.11.3.2.2}{0.11.3.2.2-fr}
\EGAReferenceAlias{0.11.3.2.2-es}{0.11.3.2.2-fr}
\EGAReferenceAlias{0.11.3.2.2.fr}{0.11.3.2.2-fr}
\EGAReferenceAlias{0.11.3.2.2.es}{0.11.3.2.2-fr}
\EGAReferenceAlias{0.11.3.2.2.eq}{0.11.3.2.2-fr}
\EGAReferenceAlias{eq:0.11.3.2.2}{0.11.3.2.2-fr}
\EGAReferenceAlias{0.11.3.3-es}{0.11.3.3}
\EGAReferenceAlias{0.11.3.3.fr}{0.11.3.3}
\EGAReferenceAlias{0.11.3.3.es}{0.11.3.3}
\EGAReferenceAlias{0.11.3.3.eq}{0.11.3.3}
\EGAReferenceAlias{eq:0.11.3.3}{0.11.3.3}
\EGAReferenceAlias{0.11.3.4}{0.11.3.4-fr}
\EGAReferenceAlias{0.11.3.4-es}{0.11.3.4-fr}
\EGAReferenceAlias{0.11.3.4.fr}{0.11.3.4-fr}
\EGAReferenceAlias{0.11.3.4.es}{0.11.3.4-fr}
\EGAReferenceAlias{0.11.3.4.eq}{0.11.3.4-fr}
\EGAReferenceAlias{eq:0.11.3.4}{0.11.3.4-fr}
\EGAReferenceAlias{0.11.3.4.1}{0.11.3.4.1-fr}
\EGAReferenceAlias{0.11.3.4.1-es}{0.11.3.4.1-fr}
\EGAReferenceAlias{0.11.3.4.1.fr}{0.11.3.4.1-fr}
\EGAReferenceAlias{0.11.3.4.1.es}{0.11.3.4.1-fr}
\EGAReferenceAlias{0.11.3.4.1.eq}{0.11.3.4.1-fr}
\EGAReferenceAlias{eq:0.11.3.4.1}{0.11.3.4.1-fr}
\EGAReferenceAlias{0.11.3.4.2}{0.11.3.4.2-fr}
\EGAReferenceAlias{0.11.3.4.2-es}{0.11.3.4.2-fr}
\EGAReferenceAlias{0.11.3.4.2.fr}{0.11.3.4.2-fr}
\EGAReferenceAlias{0.11.3.4.2.es}{0.11.3.4.2-fr}
\EGAReferenceAlias{0.11.3.4.2.eq}{0.11.3.4.2-fr}
\EGAReferenceAlias{eq:0.11.3.4.2}{0.11.3.4.2-fr}
\EGAReferenceAlias{0.11.3.5}{0.11.3.5-fr}
\EGAReferenceAlias{0.11.3.5-es}{0.11.3.5-fr}
\EGAReferenceAlias{0.11.3.5.fr}{0.11.3.5-fr}
\EGAReferenceAlias{0.11.3.5.es}{0.11.3.5-fr}
\EGAReferenceAlias{0.11.3.5.eq}{0.11.3.5-fr}
\EGAReferenceAlias{eq:0.11.3.5}{0.11.3.5-fr}
\EGAReferenceAlias{0.11.3.6-es}{0.11.3.6}
\EGAReferenceAlias{0.11.3.6.fr}{0.11.3.6}
\EGAReferenceAlias{0.11.3.6.es}{0.11.3.6}
\EGAReferenceAlias{0.11.3.6.eq}{0.11.3.6}
\EGAReferenceAlias{eq:0.11.3.6}{0.11.3.6}
\EGAReferenceAlias{0.11.4}{subsection:0.11.4}
\EGAReferenceAlias{0.11.4-fr}{subsection:0.11.4}
\EGAReferenceAlias{0.11.4-es}{subsection:0.11.4}
\EGAReferenceAlias{0.11.4.fr}{subsection:0.11.4}
\EGAReferenceAlias{0.11.4.es}{subsection:0.11.4}
\EGAReferenceAlias{0.11.4.eq}{subsection:0.11.4}
\EGAReferenceAlias{eq:0.11.4}{subsection:0.11.4}
\EGAReferenceAlias{subsection:0.11.4-es}{subsection:0.11.4}
\EGAReferenceAlias{subsection:0.11.4.fr}{subsection:0.11.4}
\EGAReferenceAlias{subsection:0.11.4.es}{subsection:0.11.4}
\EGAReferenceAlias{subsection:0.11.4.eq}{subsection:0.11.4}
\EGAReferenceAlias{eq:subsection:0.11.4}{subsection:0.11.4}
\EGAReferenceAlias{0.11.4.1-es}{0.11.4.1}
\EGAReferenceAlias{0.11.4.1.fr}{0.11.4.1}
\EGAReferenceAlias{0.11.4.1.es}{0.11.4.1}
\EGAReferenceAlias{0.11.4.1.eq}{0.11.4.1}
\EGAReferenceAlias{eq:0.11.4.1}{0.11.4.1}
\EGAReferenceAlias{0.11.4.2-es}{0.11.4.2}
\EGAReferenceAlias{0.11.4.2.fr}{0.11.4.2}
\EGAReferenceAlias{0.11.4.2.es}{0.11.4.2}
\EGAReferenceAlias{0.11.4.2.eq}{0.11.4.2}
\EGAReferenceAlias{eq:0.11.4.2}{0.11.4.2}
\EGAReferenceAlias{0.11.4.2.1}{0.11.4.2.1-fr}
\EGAReferenceAlias{0.11.4.2.1-es}{0.11.4.2.1-fr}
\EGAReferenceAlias{0.11.4.2.1.fr}{0.11.4.2.1-fr}
\EGAReferenceAlias{0.11.4.2.1.es}{0.11.4.2.1-fr}
\EGAReferenceAlias{0.11.4.2.1.eq}{0.11.4.2.1-fr}
\EGAReferenceAlias{eq:0.11.4.2.1}{0.11.4.2.1-fr}
\EGAReferenceAlias{0.11.4.2.2}{0.11.4.2.2-fr}
\EGAReferenceAlias{0.11.4.2.2-es}{0.11.4.2.2-fr}
\EGAReferenceAlias{0.11.4.2.2.fr}{0.11.4.2.2-fr}
\EGAReferenceAlias{0.11.4.2.2.es}{0.11.4.2.2-fr}
\EGAReferenceAlias{0.11.4.2.2.eq}{0.11.4.2.2-fr}
\EGAReferenceAlias{eq:0.11.4.2.2}{0.11.4.2.2-fr}
\EGAReferenceAlias{0.11.4.3-es}{0.11.4.3}
\EGAReferenceAlias{0.11.4.3.fr}{0.11.4.3}
\EGAReferenceAlias{0.11.4.3.es}{0.11.4.3}
\EGAReferenceAlias{0.11.4.3.eq}{0.11.4.3}
\EGAReferenceAlias{eq:0.11.4.3}{0.11.4.3}
\EGAReferenceAlias{0.11.4.3.1}{0.11.4.3.1-fr}
\EGAReferenceAlias{0.11.4.3.1-es}{0.11.4.3.1-fr}
\EGAReferenceAlias{0.11.4.3.1.fr}{0.11.4.3.1-fr}
\EGAReferenceAlias{0.11.4.3.1.es}{0.11.4.3.1-fr}
\EGAReferenceAlias{0.11.4.3.1.eq}{0.11.4.3.1-fr}
\EGAReferenceAlias{eq:0.11.4.3.1}{0.11.4.3.1-fr}
\EGAReferenceAlias{0.11.4.3.2}{0.11.4.3.2-fr}
\EGAReferenceAlias{0.11.4.3.2-es}{0.11.4.3.2-fr}
\EGAReferenceAlias{0.11.4.3.2.fr}{0.11.4.3.2-fr}
\EGAReferenceAlias{0.11.4.3.2.es}{0.11.4.3.2-fr}
\EGAReferenceAlias{0.11.4.3.2.eq}{0.11.4.3.2-fr}
\EGAReferenceAlias{eq:0.11.4.3.2}{0.11.4.3.2-fr}
\EGAReferenceAlias{0.11.4.4-es}{0.11.4.4}
\EGAReferenceAlias{0.11.4.4.fr}{0.11.4.4}
\EGAReferenceAlias{0.11.4.4.es}{0.11.4.4}
\EGAReferenceAlias{0.11.4.4.eq}{0.11.4.4}
\EGAReferenceAlias{eq:0.11.4.4}{0.11.4.4}
\EGAReferenceAlias{0.11.4.5}{0.11.4.5-fr}
\EGAReferenceAlias{0.11.4.5-es}{0.11.4.5-fr}
\EGAReferenceAlias{0.11.4.5.fr}{0.11.4.5-fr}
\EGAReferenceAlias{0.11.4.5.es}{0.11.4.5-fr}
\EGAReferenceAlias{0.11.4.5.eq}{0.11.4.5-fr}
\EGAReferenceAlias{eq:0.11.4.5}{0.11.4.5-fr}
\EGAReferenceAlias{0.11.4.5.1}{0.11.4.5.1-fr}
\EGAReferenceAlias{0.11.4.5.1-es}{0.11.4.5.1-fr}
\EGAReferenceAlias{0.11.4.5.1.fr}{0.11.4.5.1-fr}
\EGAReferenceAlias{0.11.4.5.1.es}{0.11.4.5.1-fr}
\EGAReferenceAlias{0.11.4.5.1.eq}{0.11.4.5.1-fr}
\EGAReferenceAlias{eq:0.11.4.5.1}{0.11.4.5.1-fr}
\EGAReferenceAlias{0.11.4.6-es}{0.11.4.6}
\EGAReferenceAlias{0.11.4.6.fr}{0.11.4.6}
\EGAReferenceAlias{0.11.4.6.es}{0.11.4.6}
\EGAReferenceAlias{0.11.4.6.eq}{0.11.4.6}
\EGAReferenceAlias{eq:0.11.4.6}{0.11.4.6}
\EGAReferenceAlias{0.11.4.6.1}{0.11.4.6.1-fr}
\EGAReferenceAlias{0.11.4.6.1-es}{0.11.4.6.1-fr}
\EGAReferenceAlias{0.11.4.6.1.fr}{0.11.4.6.1-fr}
\EGAReferenceAlias{0.11.4.6.1.es}{0.11.4.6.1-fr}
\EGAReferenceAlias{0.11.4.6.1.eq}{0.11.4.6.1-fr}
\EGAReferenceAlias{eq:0.11.4.6.1}{0.11.4.6.1-fr}
\EGAReferenceAlias{0.11.4.6.2}{0.11.4.6.2-fr}
\EGAReferenceAlias{0.11.4.6.2-es}{0.11.4.6.2-fr}
\EGAReferenceAlias{0.11.4.6.2.fr}{0.11.4.6.2-fr}
\EGAReferenceAlias{0.11.4.6.2.es}{0.11.4.6.2-fr}
\EGAReferenceAlias{0.11.4.6.2.eq}{0.11.4.6.2-fr}
\EGAReferenceAlias{eq:0.11.4.6.2}{0.11.4.6.2-fr}
\EGAReferenceAlias{0.11.4.7-es}{0.11.4.7}
\EGAReferenceAlias{0.11.4.7.fr}{0.11.4.7}
\EGAReferenceAlias{0.11.4.7.es}{0.11.4.7}
\EGAReferenceAlias{0.11.4.7.eq}{0.11.4.7}
\EGAReferenceAlias{eq:0.11.4.7}{0.11.4.7}
\EGAReferenceAlias{0.11.4.7.1}{0.11.4.7.1-fr}
\EGAReferenceAlias{0.11.4.7.1-es}{0.11.4.7.1-fr}
\EGAReferenceAlias{0.11.4.7.1.fr}{0.11.4.7.1-fr}
\EGAReferenceAlias{0.11.4.7.1.es}{0.11.4.7.1-fr}
\EGAReferenceAlias{0.11.4.7.1.eq}{0.11.4.7.1-fr}
\EGAReferenceAlias{eq:0.11.4.7.1}{0.11.4.7.1-fr}
\EGAReferenceAlias{0.11.5}{subsection:0.11.5}
\EGAReferenceAlias{0.11.5-fr}{subsection:0.11.5}
\EGAReferenceAlias{0.11.5-es}{subsection:0.11.5}
\EGAReferenceAlias{0.11.5.fr}{subsection:0.11.5}
\EGAReferenceAlias{0.11.5.es}{subsection:0.11.5}
\EGAReferenceAlias{0.11.5.eq}{subsection:0.11.5}
\EGAReferenceAlias{eq:0.11.5}{subsection:0.11.5}
\EGAReferenceAlias{subsection:0.11.5-es}{subsection:0.11.5}
\EGAReferenceAlias{subsection:0.11.5.fr}{subsection:0.11.5}
\EGAReferenceAlias{subsection:0.11.5.es}{subsection:0.11.5}
\EGAReferenceAlias{subsection:0.11.5.eq}{subsection:0.11.5}
\EGAReferenceAlias{eq:subsection:0.11.5}{subsection:0.11.5}
\EGAReferenceAlias{0.11.5.1}{0.11.5.1-fr}
\EGAReferenceAlias{0.11.5.1-es}{0.11.5.1-fr}
\EGAReferenceAlias{0.11.5.1.fr}{0.11.5.1-fr}
\EGAReferenceAlias{0.11.5.1.es}{0.11.5.1-fr}
\EGAReferenceAlias{0.11.5.1.eq}{0.11.5.1-fr}
\EGAReferenceAlias{eq:0.11.5.1}{0.11.5.1-fr}
\EGAReferenceAlias{0.11.5.2-es}{0.11.5.2}
\EGAReferenceAlias{0.11.5.2.fr}{0.11.5.2}
\EGAReferenceAlias{0.11.5.2.es}{0.11.5.2}
\EGAReferenceAlias{0.11.5.2.eq}{0.11.5.2}
\EGAReferenceAlias{eq:0.11.5.2}{0.11.5.2}
\EGAReferenceAlias{0.11.5.2.1-es}{0.11.5.2.1}
\EGAReferenceAlias{0.11.5.2.1.fr}{0.11.5.2.1}
\EGAReferenceAlias{0.11.5.2.1.es}{0.11.5.2.1}
\EGAReferenceAlias{0.11.5.2.1.eq}{0.11.5.2.1}
\EGAReferenceAlias{eq:0.11.5.2.1}{0.11.5.2.1}
\EGAReferenceAlias{0.11.5.2.2-es}{0.11.5.2.2}
\EGAReferenceAlias{0.11.5.2.2.fr}{0.11.5.2.2}
\EGAReferenceAlias{0.11.5.2.2.es}{0.11.5.2.2}
\EGAReferenceAlias{0.11.5.2.2.eq}{0.11.5.2.2}
\EGAReferenceAlias{eq:0.11.5.2.2}{0.11.5.2.2}
\EGAReferenceAlias{0.11.5.2.3}{0.11.5.2.3-fr}
\EGAReferenceAlias{0.11.5.2.3-es}{0.11.5.2.3-fr}
\EGAReferenceAlias{0.11.5.2.3.fr}{0.11.5.2.3-fr}
\EGAReferenceAlias{0.11.5.2.3.es}{0.11.5.2.3-fr}
\EGAReferenceAlias{0.11.5.2.3.eq}{0.11.5.2.3-fr}
\EGAReferenceAlias{eq:0.11.5.2.3}{0.11.5.2.3-fr}
\EGAReferenceAlias{0.11.5.2.4-es}{0.11.5.2.4}
\EGAReferenceAlias{0.11.5.2.4.fr}{0.11.5.2.4}
\EGAReferenceAlias{0.11.5.2.4.es}{0.11.5.2.4}
\EGAReferenceAlias{0.11.5.2.4.eq}{0.11.5.2.4}
\EGAReferenceAlias{eq:0.11.5.2.4}{0.11.5.2.4}
\EGAReferenceAlias{0.11.5.2.5-es}{0.11.5.2.5}
\EGAReferenceAlias{0.11.5.2.5.fr}{0.11.5.2.5}
\EGAReferenceAlias{0.11.5.2.5.es}{0.11.5.2.5}
\EGAReferenceAlias{0.11.5.2.5.eq}{0.11.5.2.5}
\EGAReferenceAlias{eq:0.11.5.2.5}{0.11.5.2.5}
\EGAReferenceAlias{0.11.5.3-es}{0.11.5.3}
\EGAReferenceAlias{0.11.5.3.fr}{0.11.5.3}
\EGAReferenceAlias{0.11.5.3.es}{0.11.5.3}
\EGAReferenceAlias{0.11.5.3.eq}{0.11.5.3}
\EGAReferenceAlias{eq:0.11.5.3}{0.11.5.3}
\EGAReferenceAlias{0.11.5.3.1}{0.11.5.3.1-fr}
\EGAReferenceAlias{0.11.5.3.1-es}{0.11.5.3.1-fr}
\EGAReferenceAlias{0.11.5.3.1.fr}{0.11.5.3.1-fr}
\EGAReferenceAlias{0.11.5.3.1.es}{0.11.5.3.1-fr}
\EGAReferenceAlias{0.11.5.3.1.eq}{0.11.5.3.1-fr}
\EGAReferenceAlias{eq:0.11.5.3.1}{0.11.5.3.1-fr}
\EGAReferenceAlias{0.11.5.4-es}{0.11.5.4}
\EGAReferenceAlias{0.11.5.4.fr}{0.11.5.4}
\EGAReferenceAlias{0.11.5.4.es}{0.11.5.4}
\EGAReferenceAlias{0.11.5.4.eq}{0.11.5.4}
\EGAReferenceAlias{eq:0.11.5.4}{0.11.5.4}
\EGAReferenceAlias{0.11.5.5}{0.11.5.5-fr}
\EGAReferenceAlias{0.11.5.5-es}{0.11.5.5-fr}
\EGAReferenceAlias{0.11.5.5.fr}{0.11.5.5-fr}
\EGAReferenceAlias{0.11.5.5.es}{0.11.5.5-fr}
\EGAReferenceAlias{0.11.5.5.eq}{0.11.5.5-fr}
\EGAReferenceAlias{eq:0.11.5.5}{0.11.5.5-fr}
\EGAReferenceAlias{0.11.5.6}{0.11.5.6-fr}
\EGAReferenceAlias{0.11.5.6-es}{0.11.5.6-fr}
\EGAReferenceAlias{0.11.5.6.fr}{0.11.5.6-fr}
\EGAReferenceAlias{0.11.5.6.es}{0.11.5.6-fr}
\EGAReferenceAlias{0.11.5.6.eq}{0.11.5.6-fr}
\EGAReferenceAlias{eq:0.11.5.6}{0.11.5.6-fr}
\EGAReferenceAlias{0.11.5.6.1}{0.11.5.6.1-fr}
\EGAReferenceAlias{0.11.5.6.1-es}{0.11.5.6.1-fr}
\EGAReferenceAlias{0.11.5.6.1.fr}{0.11.5.6.1-fr}
\EGAReferenceAlias{0.11.5.6.1.es}{0.11.5.6.1-fr}
\EGAReferenceAlias{0.11.5.6.1.eq}{0.11.5.6.1-fr}
\EGAReferenceAlias{eq:0.11.5.6.1}{0.11.5.6.1-fr}
\EGAReferenceAlias{0.11.5.6.2}{0.11.5.6.2-fr}
\EGAReferenceAlias{0.11.5.6.2-es}{0.11.5.6.2-fr}
\EGAReferenceAlias{0.11.5.6.2.fr}{0.11.5.6.2-fr}
\EGAReferenceAlias{0.11.5.6.2.es}{0.11.5.6.2-fr}
\EGAReferenceAlias{0.11.5.6.2.eq}{0.11.5.6.2-fr}
\EGAReferenceAlias{eq:0.11.5.6.2}{0.11.5.6.2-fr}
\EGAReferenceAlias{0.11.5.7-es}{0.11.5.7}
\EGAReferenceAlias{0.11.5.7.fr}{0.11.5.7}
\EGAReferenceAlias{0.11.5.7.es}{0.11.5.7}
\EGAReferenceAlias{0.11.5.7.eq}{0.11.5.7}
\EGAReferenceAlias{eq:0.11.5.7}{0.11.5.7}
\EGAReferenceAlias{0.11.6}{subsection:0.11.6}
\EGAReferenceAlias{0.11.6-fr}{subsection:0.11.6}
\EGAReferenceAlias{0.11.6-es}{subsection:0.11.6}
\EGAReferenceAlias{0.11.6.fr}{subsection:0.11.6}
\EGAReferenceAlias{0.11.6.es}{subsection:0.11.6}
\EGAReferenceAlias{0.11.6.eq}{subsection:0.11.6}
\EGAReferenceAlias{eq:0.11.6}{subsection:0.11.6}
\EGAReferenceAlias{subsection:0.11.6-es}{subsection:0.11.6}
\EGAReferenceAlias{subsection:0.11.6.fr}{subsection:0.11.6}
\EGAReferenceAlias{subsection:0.11.6.es}{subsection:0.11.6}
\EGAReferenceAlias{subsection:0.11.6.eq}{subsection:0.11.6}
\EGAReferenceAlias{eq:subsection:0.11.6}{subsection:0.11.6}
\EGAReferenceAlias{0.11.6.1-es}{0.11.6.1}
\EGAReferenceAlias{0.11.6.1.fr}{0.11.6.1}
\EGAReferenceAlias{0.11.6.1.es}{0.11.6.1}
\EGAReferenceAlias{0.11.6.1.eq}{0.11.6.1}
\EGAReferenceAlias{eq:0.11.6.1}{0.11.6.1}
\EGAReferenceAlias{0.11.6.2-es}{0.11.6.2}
\EGAReferenceAlias{0.11.6.2.fr}{0.11.6.2}
\EGAReferenceAlias{0.11.6.2.es}{0.11.6.2}
\EGAReferenceAlias{0.11.6.2.eq}{0.11.6.2}
\EGAReferenceAlias{eq:0.11.6.2}{0.11.6.2}
\EGAReferenceAlias{0.11.6.2.1}{0.11.6.2.1-fr}
\EGAReferenceAlias{0.11.6.2.1-es}{0.11.6.2.1-fr}
\EGAReferenceAlias{0.11.6.2.1.fr}{0.11.6.2.1-fr}
\EGAReferenceAlias{0.11.6.2.1.es}{0.11.6.2.1-fr}
\EGAReferenceAlias{0.11.6.2.1.eq}{0.11.6.2.1-fr}
\EGAReferenceAlias{eq:0.11.6.2.1}{0.11.6.2.1-fr}
\EGAReferenceAlias{0.11.6.2.2}{0.11.6.2.2-fr}
\EGAReferenceAlias{0.11.6.2.2-es}{0.11.6.2.2-fr}
\EGAReferenceAlias{0.11.6.2.2.fr}{0.11.6.2.2-fr}
\EGAReferenceAlias{0.11.6.2.2.es}{0.11.6.2.2-fr}
\EGAReferenceAlias{0.11.6.2.2.eq}{0.11.6.2.2-fr}
\EGAReferenceAlias{eq:0.11.6.2.2}{0.11.6.2.2-fr}
\EGAReferenceAlias{0.11.6.3-es}{0.11.6.3}
\EGAReferenceAlias{0.11.6.3.fr}{0.11.6.3}
\EGAReferenceAlias{0.11.6.3.es}{0.11.6.3}
\EGAReferenceAlias{0.11.6.3.eq}{0.11.6.3}
\EGAReferenceAlias{eq:0.11.6.3}{0.11.6.3}
\EGAReferenceAlias{0.11.6.3.1}{0.11.6.3.1-fr}
\EGAReferenceAlias{0.11.6.3.1-es}{0.11.6.3.1-fr}
\EGAReferenceAlias{0.11.6.3.1.fr}{0.11.6.3.1-fr}
\EGAReferenceAlias{0.11.6.3.1.es}{0.11.6.3.1-fr}
\EGAReferenceAlias{0.11.6.3.1.eq}{0.11.6.3.1-fr}
\EGAReferenceAlias{eq:0.11.6.3.1}{0.11.6.3.1-fr}
\EGAReferenceAlias{0.11.6.3.2}{0.11.6.3.2-fr}
\EGAReferenceAlias{0.11.6.3.2-es}{0.11.6.3.2-fr}
\EGAReferenceAlias{0.11.6.3.2.fr}{0.11.6.3.2-fr}
\EGAReferenceAlias{0.11.6.3.2.es}{0.11.6.3.2-fr}
\EGAReferenceAlias{0.11.6.3.2.eq}{0.11.6.3.2-fr}
\EGAReferenceAlias{eq:0.11.6.3.2}{0.11.6.3.2-fr}
\EGAReferenceAlias{0.11.6.4-es}{0.11.6.4}
\EGAReferenceAlias{0.11.6.4.fr}{0.11.6.4}
\EGAReferenceAlias{0.11.6.4.es}{0.11.6.4}
\EGAReferenceAlias{0.11.6.4.eq}{0.11.6.4}
\EGAReferenceAlias{eq:0.11.6.4}{0.11.6.4}
\EGAReferenceAlias{0.11.6.4.1}{0.11.6.4.1-fr}
\EGAReferenceAlias{0.11.6.4.1-es}{0.11.6.4.1-fr}
\EGAReferenceAlias{0.11.6.4.1.fr}{0.11.6.4.1-fr}
\EGAReferenceAlias{0.11.6.4.1.es}{0.11.6.4.1-fr}
\EGAReferenceAlias{0.11.6.4.1.eq}{0.11.6.4.1-fr}
\EGAReferenceAlias{eq:0.11.6.4.1}{0.11.6.4.1-fr}
\EGAReferenceAlias{0.11.6.4.2}{0.11.6.4.2-fr}
\EGAReferenceAlias{0.11.6.4.2-es}{0.11.6.4.2-fr}
\EGAReferenceAlias{0.11.6.4.2.fr}{0.11.6.4.2-fr}
\EGAReferenceAlias{0.11.6.4.2.es}{0.11.6.4.2-fr}
\EGAReferenceAlias{0.11.6.4.2.eq}{0.11.6.4.2-fr}
\EGAReferenceAlias{eq:0.11.6.4.2}{0.11.6.4.2-fr}
\EGAReferenceAlias{0.11.6.5-es}{0.11.6.5}
\EGAReferenceAlias{0.11.6.5.fr}{0.11.6.5}
\EGAReferenceAlias{0.11.6.5.es}{0.11.6.5}
\EGAReferenceAlias{0.11.6.5.eq}{0.11.6.5}
\EGAReferenceAlias{eq:0.11.6.5}{0.11.6.5}
\EGAReferenceAlias{0.11.6.5.1}{0.11.6.5.1-fr}
\EGAReferenceAlias{0.11.6.5.1-es}{0.11.6.5.1-fr}
\EGAReferenceAlias{0.11.6.5.1.fr}{0.11.6.5.1-fr}
\EGAReferenceAlias{0.11.6.5.1.es}{0.11.6.5.1-fr}
\EGAReferenceAlias{0.11.6.5.1.eq}{0.11.6.5.1-fr}
\EGAReferenceAlias{eq:0.11.6.5.1}{0.11.6.5.1-fr}
\EGAReferenceAlias{0.11.6.5.2}{0.11.6.5.2-fr}
\EGAReferenceAlias{0.11.6.5.2-es}{0.11.6.5.2-fr}
\EGAReferenceAlias{0.11.6.5.2.fr}{0.11.6.5.2-fr}
\EGAReferenceAlias{0.11.6.5.2.es}{0.11.6.5.2-fr}
\EGAReferenceAlias{0.11.6.5.2.eq}{0.11.6.5.2-fr}
\EGAReferenceAlias{eq:0.11.6.5.2}{0.11.6.5.2-fr}
\EGAReferenceAlias{0.11.7}{subsection:0.11.7}
\EGAReferenceAlias{0.11.7-fr}{subsection:0.11.7}
\EGAReferenceAlias{0.11.7-es}{subsection:0.11.7}
\EGAReferenceAlias{0.11.7.fr}{subsection:0.11.7}
\EGAReferenceAlias{0.11.7.es}{subsection:0.11.7}
\EGAReferenceAlias{0.11.7.eq}{subsection:0.11.7}
\EGAReferenceAlias{eq:0.11.7}{subsection:0.11.7}
\EGAReferenceAlias{subsection:0.11.7-es}{subsection:0.11.7}
\EGAReferenceAlias{subsection:0.11.7.fr}{subsection:0.11.7}
\EGAReferenceAlias{subsection:0.11.7.es}{subsection:0.11.7}
\EGAReferenceAlias{subsection:0.11.7.eq}{subsection:0.11.7}
\EGAReferenceAlias{eq:subsection:0.11.7}{subsection:0.11.7}
\EGAReferenceAlias{0.11.7.1-es}{0.11.7.1}
\EGAReferenceAlias{0.11.7.1.fr}{0.11.7.1}
\EGAReferenceAlias{0.11.7.1.es}{0.11.7.1}
\EGAReferenceAlias{0.11.7.1.eq}{0.11.7.1}
\EGAReferenceAlias{eq:0.11.7.1}{0.11.7.1}
\EGAReferenceAlias{0.11.7.2-es}{0.11.7.2}
\EGAReferenceAlias{0.11.7.2.fr}{0.11.7.2}
\EGAReferenceAlias{0.11.7.2.es}{0.11.7.2}
\EGAReferenceAlias{0.11.7.2.eq}{0.11.7.2}
\EGAReferenceAlias{eq:0.11.7.2}{0.11.7.2}
\EGAReferenceAlias{0.11.7.3-es}{0.11.7.3}
\EGAReferenceAlias{0.11.7.3.fr}{0.11.7.3}
\EGAReferenceAlias{0.11.7.3.es}{0.11.7.3}
\EGAReferenceAlias{0.11.7.3.eq}{0.11.7.3}
\EGAReferenceAlias{eq:0.11.7.3}{0.11.7.3}
\EGAReferenceAlias{0.11.8}{subsection:0.11.8}
\EGAReferenceAlias{0.11.8-fr}{subsection:0.11.8}
\EGAReferenceAlias{0.11.8-es}{subsection:0.11.8}
\EGAReferenceAlias{0.11.8.fr}{subsection:0.11.8}
\EGAReferenceAlias{0.11.8.es}{subsection:0.11.8}
\EGAReferenceAlias{0.11.8.eq}{subsection:0.11.8}
\EGAReferenceAlias{eq:0.11.8}{subsection:0.11.8}
\EGAReferenceAlias{subsection:0.11.8-es}{subsection:0.11.8}
\EGAReferenceAlias{subsection:0.11.8.fr}{subsection:0.11.8}
\EGAReferenceAlias{subsection:0.11.8.es}{subsection:0.11.8}
\EGAReferenceAlias{subsection:0.11.8.eq}{subsection:0.11.8}
\EGAReferenceAlias{eq:subsection:0.11.8}{subsection:0.11.8}
\EGAReferenceAlias{0.11.8.1}{0.11.8.1-fr}
\EGAReferenceAlias{0.11.8.1-es}{0.11.8.1-fr}
\EGAReferenceAlias{0.11.8.1.fr}{0.11.8.1-fr}
\EGAReferenceAlias{0.11.8.1.es}{0.11.8.1-fr}
\EGAReferenceAlias{0.11.8.1.eq}{0.11.8.1-fr}
\EGAReferenceAlias{eq:0.11.8.1}{0.11.8.1-fr}
\EGAReferenceAlias{0.11.8.2-es}{0.11.8.2}
\EGAReferenceAlias{0.11.8.2.fr}{0.11.8.2}
\EGAReferenceAlias{0.11.8.2.es}{0.11.8.2}
\EGAReferenceAlias{0.11.8.2.eq}{0.11.8.2}
\EGAReferenceAlias{eq:0.11.8.2}{0.11.8.2}
\EGAReferenceAlias{0.11.8.3-es}{0.11.8.3}
\EGAReferenceAlias{0.11.8.3.fr}{0.11.8.3}
\EGAReferenceAlias{0.11.8.3.es}{0.11.8.3}
\EGAReferenceAlias{0.11.8.3.eq}{0.11.8.3}
\EGAReferenceAlias{eq:0.11.8.3}{0.11.8.3}
\EGAReferenceAlias{0.11.8.3.1}{0.11.8.3.1-fr}
\EGAReferenceAlias{0.11.8.3.1-es}{0.11.8.3.1-fr}
\EGAReferenceAlias{0.11.8.3.1.fr}{0.11.8.3.1-fr}
\EGAReferenceAlias{0.11.8.3.1.es}{0.11.8.3.1-fr}
\EGAReferenceAlias{0.11.8.3.1.eq}{0.11.8.3.1-fr}
\EGAReferenceAlias{eq:0.11.8.3.1}{0.11.8.3.1-fr}
\EGAReferenceAlias{0.11.8.4-es}{0.11.8.4}
\EGAReferenceAlias{0.11.8.4.fr}{0.11.8.4}
\EGAReferenceAlias{0.11.8.4.es}{0.11.8.4}
\EGAReferenceAlias{0.11.8.4.eq}{0.11.8.4}
\EGAReferenceAlias{eq:0.11.8.4}{0.11.8.4}
\EGAReferenceAlias{0.11.8.5}{0.11.8.5-fr}
\EGAReferenceAlias{0.11.8.5-es}{0.11.8.5-fr}
\EGAReferenceAlias{0.11.8.5.fr}{0.11.8.5-fr}
\EGAReferenceAlias{0.11.8.5.es}{0.11.8.5-fr}
\EGAReferenceAlias{0.11.8.5.eq}{0.11.8.5-fr}
\EGAReferenceAlias{eq:0.11.8.5}{0.11.8.5-fr}
\EGAReferenceAlias{0.11.8.6-es}{0.11.8.6}
\EGAReferenceAlias{0.11.8.6.fr}{0.11.8.6}
\EGAReferenceAlias{0.11.8.6.es}{0.11.8.6}
\EGAReferenceAlias{0.11.8.6.eq}{0.11.8.6}
\EGAReferenceAlias{eq:0.11.8.6}{0.11.8.6}
\EGAReferenceAlias{0.11.8.7-es}{0.11.8.7}
\EGAReferenceAlias{0.11.8.7.fr}{0.11.8.7}
\EGAReferenceAlias{0.11.8.7.es}{0.11.8.7}
\EGAReferenceAlias{0.11.8.7.eq}{0.11.8.7}
\EGAReferenceAlias{eq:0.11.8.7}{0.11.8.7}
\EGAReferenceAlias{0.11.8.8-es}{0.11.8.8}
\EGAReferenceAlias{0.11.8.8.fr}{0.11.8.8}
\EGAReferenceAlias{0.11.8.8.es}{0.11.8.8}
\EGAReferenceAlias{0.11.8.8.eq}{0.11.8.8}
\EGAReferenceAlias{eq:0.11.8.8}{0.11.8.8}
\EGAReferenceAlias{0.11.8.9-es}{0.11.8.9}
\EGAReferenceAlias{0.11.8.9.fr}{0.11.8.9}
\EGAReferenceAlias{0.11.8.9.es}{0.11.8.9}
\EGAReferenceAlias{0.11.8.9.eq}{0.11.8.9}
\EGAReferenceAlias{eq:0.11.8.9}{0.11.8.9}
\EGAReferenceAlias{0.11.8.9.1}{0.11.8.9.1-fr}
\EGAReferenceAlias{0.11.8.9.1-es}{0.11.8.9.1-fr}
\EGAReferenceAlias{0.11.8.9.1.fr}{0.11.8.9.1-fr}
\EGAReferenceAlias{0.11.8.9.1.es}{0.11.8.9.1-fr}
\EGAReferenceAlias{0.11.8.9.1.eq}{0.11.8.9.1-fr}
\EGAReferenceAlias{eq:0.11.8.9.1}{0.11.8.9.1-fr}
\EGAReferenceAlias{0.11.8.10-es}{0.11.8.10}
\EGAReferenceAlias{0.11.8.10.fr}{0.11.8.10}
\EGAReferenceAlias{0.11.8.10.es}{0.11.8.10}
\EGAReferenceAlias{0.11.8.10.eq}{0.11.8.10}
\EGAReferenceAlias{eq:0.11.8.10}{0.11.8.10}
\EGAReferenceAlias{0.11.9}{subsection:0.11.9}
\EGAReferenceAlias{0.11.9-fr}{subsection:0.11.9}
\EGAReferenceAlias{0.11.9-es}{subsection:0.11.9}
\EGAReferenceAlias{0.11.9.fr}{subsection:0.11.9}
\EGAReferenceAlias{0.11.9.es}{subsection:0.11.9}
\EGAReferenceAlias{0.11.9.eq}{subsection:0.11.9}
\EGAReferenceAlias{eq:0.11.9}{subsection:0.11.9}
\EGAReferenceAlias{subsection:0.11.9-es}{subsection:0.11.9}
\EGAReferenceAlias{subsection:0.11.9.fr}{subsection:0.11.9}
\EGAReferenceAlias{subsection:0.11.9.es}{subsection:0.11.9}
\EGAReferenceAlias{subsection:0.11.9.eq}{subsection:0.11.9}
\EGAReferenceAlias{eq:subsection:0.11.9}{subsection:0.11.9}
\EGAReferenceAlias{0.11.9.1-es}{0.11.9.1}
\EGAReferenceAlias{0.11.9.1.fr}{0.11.9.1}
\EGAReferenceAlias{0.11.9.1.es}{0.11.9.1}
\EGAReferenceAlias{0.11.9.1.eq}{0.11.9.1}
\EGAReferenceAlias{eq:0.11.9.1}{0.11.9.1}
\EGAReferenceAlias{0.11.9.1.1}{0.11.9.1.1-fr}
\EGAReferenceAlias{0.11.9.1.1-es}{0.11.9.1.1-fr}
\EGAReferenceAlias{0.11.9.1.1.fr}{0.11.9.1.1-fr}
\EGAReferenceAlias{0.11.9.1.1.es}{0.11.9.1.1-fr}
\EGAReferenceAlias{0.11.9.1.1.eq}{0.11.9.1.1-fr}
\EGAReferenceAlias{eq:0.11.9.1.1}{0.11.9.1.1-fr}
\EGAReferenceAlias{0.11.9.1.2}{0.11.9.1.2-fr}
\EGAReferenceAlias{0.11.9.1.2-es}{0.11.9.1.2-fr}
\EGAReferenceAlias{0.11.9.1.2.fr}{0.11.9.1.2-fr}
\EGAReferenceAlias{0.11.9.1.2.es}{0.11.9.1.2-fr}
\EGAReferenceAlias{0.11.9.1.2.eq}{0.11.9.1.2-fr}
\EGAReferenceAlias{eq:0.11.9.1.2}{0.11.9.1.2-fr}
\EGAReferenceAlias{0.11.9.1.3}{0.11.9.1.3-fr}
\EGAReferenceAlias{0.11.9.1.3-es}{0.11.9.1.3-fr}
\EGAReferenceAlias{0.11.9.1.3.fr}{0.11.9.1.3-fr}
\EGAReferenceAlias{0.11.9.1.3.es}{0.11.9.1.3-fr}
\EGAReferenceAlias{0.11.9.1.3.eq}{0.11.9.1.3-fr}
\EGAReferenceAlias{eq:0.11.9.1.3}{0.11.9.1.3-fr}
\EGAReferenceAlias{0.11.9.1.4}{0.11.9.1.4-fr}
\EGAReferenceAlias{0.11.9.1.4-es}{0.11.9.1.4-fr}
\EGAReferenceAlias{0.11.9.1.4.fr}{0.11.9.1.4-fr}
\EGAReferenceAlias{0.11.9.1.4.es}{0.11.9.1.4-fr}
\EGAReferenceAlias{0.11.9.1.4.eq}{0.11.9.1.4-fr}
\EGAReferenceAlias{eq:0.11.9.1.4}{0.11.9.1.4-fr}
\EGAReferenceAlias{0.11.9.2}{0.11.9.2-fr}
\EGAReferenceAlias{0.11.9.2-es}{0.11.9.2-fr}
\EGAReferenceAlias{0.11.9.2.fr}{0.11.9.2-fr}
\EGAReferenceAlias{0.11.9.2.es}{0.11.9.2-fr}
\EGAReferenceAlias{0.11.9.2.eq}{0.11.9.2-fr}
\EGAReferenceAlias{eq:0.11.9.2}{0.11.9.2-fr}
\EGAReferenceAlias{0.11.9.3-es}{0.11.9.3}
\EGAReferenceAlias{0.11.9.3.fr}{0.11.9.3}
\EGAReferenceAlias{0.11.9.3.es}{0.11.9.3}
\EGAReferenceAlias{0.11.9.3.eq}{0.11.9.3}
\EGAReferenceAlias{eq:0.11.9.3}{0.11.9.3}
\EGAReferenceAlias{0.11.10}{subsection:0.11.10}
\EGAReferenceAlias{0.11.10-fr}{subsection:0.11.10}
\EGAReferenceAlias{0.11.10-es}{subsection:0.11.10}
\EGAReferenceAlias{0.11.10.fr}{subsection:0.11.10}
\EGAReferenceAlias{0.11.10.es}{subsection:0.11.10}
\EGAReferenceAlias{0.11.10.eq}{subsection:0.11.10}
\EGAReferenceAlias{eq:0.11.10}{subsection:0.11.10}
\EGAReferenceAlias{subsection:0.11.10-es}{subsection:0.11.10}
\EGAReferenceAlias{subsection:0.11.10.fr}{subsection:0.11.10}
\EGAReferenceAlias{subsection:0.11.10.es}{subsection:0.11.10}
\EGAReferenceAlias{subsection:0.11.10.eq}{subsection:0.11.10}
\EGAReferenceAlias{eq:subsection:0.11.10}{subsection:0.11.10}
\EGAReferenceAlias{0.11.10.1-es}{0.11.10.1}
\EGAReferenceAlias{0.11.10.1.fr}{0.11.10.1}
\EGAReferenceAlias{0.11.10.1.es}{0.11.10.1}
\EGAReferenceAlias{0.11.10.1.eq}{0.11.10.1}
\EGAReferenceAlias{eq:0.11.10.1}{0.11.10.1}
\EGAReferenceAlias{0.11.10.1.1}{0.11.10.1.1-fr}
\EGAReferenceAlias{0.11.10.1.1-es}{0.11.10.1.1-fr}
\EGAReferenceAlias{0.11.10.1.1.fr}{0.11.10.1.1-fr}
\EGAReferenceAlias{0.11.10.1.1.es}{0.11.10.1.1-fr}
\EGAReferenceAlias{0.11.10.1.1.eq}{0.11.10.1.1-fr}
\EGAReferenceAlias{eq:0.11.10.1.1}{0.11.10.1.1-fr}
\EGAReferenceAlias{0.11.10.1.2}{0.11.10.1.2-fr}
\EGAReferenceAlias{0.11.10.1.2-es}{0.11.10.1.2-fr}
\EGAReferenceAlias{0.11.10.1.2.fr}{0.11.10.1.2-fr}
\EGAReferenceAlias{0.11.10.1.2.es}{0.11.10.1.2-fr}
\EGAReferenceAlias{0.11.10.1.2.eq}{0.11.10.1.2-fr}
\EGAReferenceAlias{eq:0.11.10.1.2}{0.11.10.1.2-fr}
\EGAReferenceAlias{0.11.10.2-es}{0.11.10.2}
\EGAReferenceAlias{0.11.10.2.fr}{0.11.10.2}
\EGAReferenceAlias{0.11.10.2.es}{0.11.10.2}
\EGAReferenceAlias{0.11.10.2.eq}{0.11.10.2}
\EGAReferenceAlias{eq:0.11.10.2}{0.11.10.2}
\EGAReferenceAlias{0.11.10.3-es}{0.11.10.3}
\EGAReferenceAlias{0.11.10.3.fr}{0.11.10.3}
\EGAReferenceAlias{0.11.10.3.es}{0.11.10.3}
\EGAReferenceAlias{0.11.10.3.eq}{0.11.10.3}
\EGAReferenceAlias{eq:0.11.10.3}{0.11.10.3}
\EGAReferenceAlias{0.12}{section:0.12}
\EGAReferenceAlias{0.12-fr}{section:0.12}
\EGAReferenceAlias{0.12-es}{section:0.12}
\EGAReferenceAlias{0.12.fr}{section:0.12}
\EGAReferenceAlias{0.12.es}{section:0.12}
\EGAReferenceAlias{0.12.eq}{section:0.12}
\EGAReferenceAlias{eq:0.12}{section:0.12}
\EGAReferenceAlias{section:0.12-es}{section:0.12}
\EGAReferenceAlias{section:0.12.fr}{section:0.12}
\EGAReferenceAlias{section:0.12.es}{section:0.12}
\EGAReferenceAlias{section:0.12.eq}{section:0.12}
\EGAReferenceAlias{eq:section:0.12}{section:0.12}
\EGAReferenceAlias{0.12.1}{subsection:0.12.1}
\EGAReferenceAlias{0.12.1-fr}{subsection:0.12.1}
\EGAReferenceAlias{0.12.1-es}{subsection:0.12.1}
\EGAReferenceAlias{0.12.1.fr}{subsection:0.12.1}
\EGAReferenceAlias{0.12.1.es}{subsection:0.12.1}
\EGAReferenceAlias{0.12.1.eq}{subsection:0.12.1}
\EGAReferenceAlias{eq:0.12.1}{subsection:0.12.1}
\EGAReferenceAlias{subsection:0.12.1-es}{subsection:0.12.1}
\EGAReferenceAlias{subsection:0.12.1.fr}{subsection:0.12.1}
\EGAReferenceAlias{subsection:0.12.1.es}{subsection:0.12.1}
\EGAReferenceAlias{subsection:0.12.1.eq}{subsection:0.12.1}
\EGAReferenceAlias{eq:subsection:0.12.1}{subsection:0.12.1}
\EGAReferenceAlias{0.12.1.1-es}{0.12.1.1}
\EGAReferenceAlias{0.12.1.1.fr}{0.12.1.1}
\EGAReferenceAlias{0.12.1.1.es}{0.12.1.1}
\EGAReferenceAlias{0.12.1.1.eq}{0.12.1.1}
\EGAReferenceAlias{eq:0.12.1.1}{0.12.1.1}
\EGAReferenceAlias{0.12.1.2-es}{0.12.1.2}
\EGAReferenceAlias{0.12.1.2.fr}{0.12.1.2}
\EGAReferenceAlias{0.12.1.2.es}{0.12.1.2}
\EGAReferenceAlias{0.12.1.2.eq}{0.12.1.2}
\EGAReferenceAlias{eq:0.12.1.2}{0.12.1.2}
\EGAReferenceAlias{0.12.1.2.1}{0.12.1.2.1-fr}
\EGAReferenceAlias{0.12.1.2.1-es}{0.12.1.2.1-fr}
\EGAReferenceAlias{0.12.1.2.1.fr}{0.12.1.2.1-fr}
\EGAReferenceAlias{0.12.1.2.1.es}{0.12.1.2.1-fr}
\EGAReferenceAlias{0.12.1.2.1.eq}{0.12.1.2.1-fr}
\EGAReferenceAlias{eq:0.12.1.2.1}{0.12.1.2.1-fr}
\EGAReferenceAlias{0.12.1.3-es}{0.12.1.3}
\EGAReferenceAlias{0.12.1.3.fr}{0.12.1.3}
\EGAReferenceAlias{0.12.1.3.es}{0.12.1.3}
\EGAReferenceAlias{0.12.1.3.eq}{0.12.1.3}
\EGAReferenceAlias{eq:0.12.1.3}{0.12.1.3}
\EGAReferenceAlias{0.12.1.3.1}{0.12.1.3.1-fr}
\EGAReferenceAlias{0.12.1.3.1-es}{0.12.1.3.1-fr}
\EGAReferenceAlias{0.12.1.3.1.fr}{0.12.1.3.1-fr}
\EGAReferenceAlias{0.12.1.3.1.es}{0.12.1.3.1-fr}
\EGAReferenceAlias{0.12.1.3.1.eq}{0.12.1.3.1-fr}
\EGAReferenceAlias{eq:0.12.1.3.1}{0.12.1.3.1-fr}
\EGAReferenceAlias{0.12.1.3.2}{0.12.1.3.2-fr}
\EGAReferenceAlias{0.12.1.3.2-es}{0.12.1.3.2-fr}
\EGAReferenceAlias{0.12.1.3.2.fr}{0.12.1.3.2-fr}
\EGAReferenceAlias{0.12.1.3.2.es}{0.12.1.3.2-fr}
\EGAReferenceAlias{0.12.1.3.2.eq}{0.12.1.3.2-fr}
\EGAReferenceAlias{eq:0.12.1.3.2}{0.12.1.3.2-fr}
\EGAReferenceAlias{0.12.1.3.3}{0.12.1.3.3-fr}
\EGAReferenceAlias{0.12.1.3.3-es}{0.12.1.3.3-fr}
\EGAReferenceAlias{0.12.1.3.3.fr}{0.12.1.3.3-fr}
\EGAReferenceAlias{0.12.1.3.3.es}{0.12.1.3.3-fr}
\EGAReferenceAlias{0.12.1.3.3.eq}{0.12.1.3.3-fr}
\EGAReferenceAlias{eq:0.12.1.3.3}{0.12.1.3.3-fr}
\EGAReferenceAlias{0.12.1.3.4}{0.12.1.3.4-fr}
\EGAReferenceAlias{0.12.1.3.4-es}{0.12.1.3.4-fr}
\EGAReferenceAlias{0.12.1.3.4.fr}{0.12.1.3.4-fr}
\EGAReferenceAlias{0.12.1.3.4.es}{0.12.1.3.4-fr}
\EGAReferenceAlias{0.12.1.3.4.eq}{0.12.1.3.4-fr}
\EGAReferenceAlias{eq:0.12.1.3.4}{0.12.1.3.4-fr}
\EGAReferenceAlias{0.12.1.3.5}{0.12.1.3.5-fr}
\EGAReferenceAlias{0.12.1.3.5-es}{0.12.1.3.5-fr}
\EGAReferenceAlias{0.12.1.3.5.fr}{0.12.1.3.5-fr}
\EGAReferenceAlias{0.12.1.3.5.es}{0.12.1.3.5-fr}
\EGAReferenceAlias{0.12.1.3.5.eq}{0.12.1.3.5-fr}
\EGAReferenceAlias{eq:0.12.1.3.5}{0.12.1.3.5-fr}
\EGAReferenceAlias{0.12.1.4}{0.12.1.4-fr}
\EGAReferenceAlias{0.12.1.4-es}{0.12.1.4-fr}
\EGAReferenceAlias{0.12.1.4.fr}{0.12.1.4-fr}
\EGAReferenceAlias{0.12.1.4.es}{0.12.1.4-fr}
\EGAReferenceAlias{0.12.1.4.eq}{0.12.1.4-fr}
\EGAReferenceAlias{eq:0.12.1.4}{0.12.1.4-fr}
\EGAReferenceAlias{0.12.1.4.1}{0.12.1.4.1-fr}
\EGAReferenceAlias{0.12.1.4.1-es}{0.12.1.4.1-fr}
\EGAReferenceAlias{0.12.1.4.1.fr}{0.12.1.4.1-fr}
\EGAReferenceAlias{0.12.1.4.1.es}{0.12.1.4.1-fr}
\EGAReferenceAlias{0.12.1.4.1.eq}{0.12.1.4.1-fr}
\EGAReferenceAlias{eq:0.12.1.4.1}{0.12.1.4.1-fr}
\EGAReferenceAlias{0.12.1.4.2}{0.12.1.4.2-fr}
\EGAReferenceAlias{0.12.1.4.2-es}{0.12.1.4.2-fr}
\EGAReferenceAlias{0.12.1.4.2.fr}{0.12.1.4.2-fr}
\EGAReferenceAlias{0.12.1.4.2.es}{0.12.1.4.2-fr}
\EGAReferenceAlias{0.12.1.4.2.eq}{0.12.1.4.2-fr}
\EGAReferenceAlias{eq:0.12.1.4.2}{0.12.1.4.2-fr}
\EGAReferenceAlias{0.12.1.5}{0.12.1.5-fr}
\EGAReferenceAlias{0.12.1.5-es}{0.12.1.5-fr}
\EGAReferenceAlias{0.12.1.5.fr}{0.12.1.5-fr}
\EGAReferenceAlias{0.12.1.5.es}{0.12.1.5-fr}
\EGAReferenceAlias{0.12.1.5.eq}{0.12.1.5-fr}
\EGAReferenceAlias{eq:0.12.1.5}{0.12.1.5-fr}
\EGAReferenceAlias{0.12.1.5.1}{0.12.1.5.1-fr}
\EGAReferenceAlias{0.12.1.5.1-es}{0.12.1.5.1-fr}
\EGAReferenceAlias{0.12.1.5.1.fr}{0.12.1.5.1-fr}
\EGAReferenceAlias{0.12.1.5.1.es}{0.12.1.5.1-fr}
\EGAReferenceAlias{0.12.1.5.1.eq}{0.12.1.5.1-fr}
\EGAReferenceAlias{eq:0.12.1.5.1}{0.12.1.5.1-fr}
\EGAReferenceAlias{0.12.1.5.2}{0.12.1.5.2-fr}
\EGAReferenceAlias{0.12.1.5.2-es}{0.12.1.5.2-fr}
\EGAReferenceAlias{0.12.1.5.2.fr}{0.12.1.5.2-fr}
\EGAReferenceAlias{0.12.1.5.2.es}{0.12.1.5.2-fr}
\EGAReferenceAlias{0.12.1.5.2.eq}{0.12.1.5.2-fr}
\EGAReferenceAlias{eq:0.12.1.5.2}{0.12.1.5.2-fr}
\EGAReferenceAlias{0.12.1.5.3}{0.12.1.5.3-fr}
\EGAReferenceAlias{0.12.1.5.3-es}{0.12.1.5.3-fr}
\EGAReferenceAlias{0.12.1.5.3.fr}{0.12.1.5.3-fr}
\EGAReferenceAlias{0.12.1.5.3.es}{0.12.1.5.3-fr}
\EGAReferenceAlias{0.12.1.5.3.eq}{0.12.1.5.3-fr}
\EGAReferenceAlias{eq:0.12.1.5.3}{0.12.1.5.3-fr}
\EGAReferenceAlias{0.12.1.5.4}{0.12.1.5.4-fr}
\EGAReferenceAlias{0.12.1.5.4-es}{0.12.1.5.4-fr}
\EGAReferenceAlias{0.12.1.5.4.fr}{0.12.1.5.4-fr}
\EGAReferenceAlias{0.12.1.5.4.es}{0.12.1.5.4-fr}
\EGAReferenceAlias{0.12.1.5.4.eq}{0.12.1.5.4-fr}
\EGAReferenceAlias{eq:0.12.1.5.4}{0.12.1.5.4-fr}
\EGAReferenceAlias{0.12.1.5.5}{0.12.1.5.5-fr}
\EGAReferenceAlias{0.12.1.5.5-es}{0.12.1.5.5-fr}
\EGAReferenceAlias{0.12.1.5.5.fr}{0.12.1.5.5-fr}
\EGAReferenceAlias{0.12.1.5.5.es}{0.12.1.5.5-fr}
\EGAReferenceAlias{0.12.1.5.5.eq}{0.12.1.5.5-fr}
\EGAReferenceAlias{eq:0.12.1.5.5}{0.12.1.5.5-fr}
\EGAReferenceAlias{0.12.1.5.6}{0.12.1.5.6-fr}
\EGAReferenceAlias{0.12.1.5.6-es}{0.12.1.5.6-fr}
\EGAReferenceAlias{0.12.1.5.6.fr}{0.12.1.5.6-fr}
\EGAReferenceAlias{0.12.1.5.6.es}{0.12.1.5.6-fr}
\EGAReferenceAlias{0.12.1.5.6.eq}{0.12.1.5.6-fr}
\EGAReferenceAlias{eq:0.12.1.5.6}{0.12.1.5.6-fr}
\EGAReferenceAlias{0.12.1.6-es}{0.12.1.6}
\EGAReferenceAlias{0.12.1.6.fr}{0.12.1.6}
\EGAReferenceAlias{0.12.1.6.es}{0.12.1.6}
\EGAReferenceAlias{0.12.1.6.eq}{0.12.1.6}
\EGAReferenceAlias{eq:0.12.1.6}{0.12.1.6}
\EGAReferenceAlias{0.12.1.6.1}{0.12.1.6.1-fr}
\EGAReferenceAlias{0.12.1.6.1-es}{0.12.1.6.1-fr}
\EGAReferenceAlias{0.12.1.6.1.fr}{0.12.1.6.1-fr}
\EGAReferenceAlias{0.12.1.6.1.es}{0.12.1.6.1-fr}
\EGAReferenceAlias{0.12.1.6.1.eq}{0.12.1.6.1-fr}
\EGAReferenceAlias{eq:0.12.1.6.1}{0.12.1.6.1-fr}
\EGAReferenceAlias{0.12.1.6.2}{0.12.1.6.2-fr}
\EGAReferenceAlias{0.12.1.6.2-es}{0.12.1.6.2-fr}
\EGAReferenceAlias{0.12.1.6.2.fr}{0.12.1.6.2-fr}
\EGAReferenceAlias{0.12.1.6.2.es}{0.12.1.6.2-fr}
\EGAReferenceAlias{0.12.1.6.2.eq}{0.12.1.6.2-fr}
\EGAReferenceAlias{eq:0.12.1.6.2}{0.12.1.6.2-fr}
\EGAReferenceAlias{0.12.1.7-es}{0.12.1.7}
\EGAReferenceAlias{0.12.1.7.fr}{0.12.1.7}
\EGAReferenceAlias{0.12.1.7.es}{0.12.1.7}
\EGAReferenceAlias{0.12.1.7.eq}{0.12.1.7}
\EGAReferenceAlias{eq:0.12.1.7}{0.12.1.7}
\EGAReferenceAlias{0.12.1.7.1}{0.12.1.7.1-fr}
\EGAReferenceAlias{0.12.1.7.1-es}{0.12.1.7.1-fr}
\EGAReferenceAlias{0.12.1.7.1.fr}{0.12.1.7.1-fr}
\EGAReferenceAlias{0.12.1.7.1.es}{0.12.1.7.1-fr}
\EGAReferenceAlias{0.12.1.7.1.eq}{0.12.1.7.1-fr}
\EGAReferenceAlias{eq:0.12.1.7.1}{0.12.1.7.1-fr}
\EGAReferenceAlias{0.12.1.7.2}{0.12.1.7.2-fr}
\EGAReferenceAlias{0.12.1.7.2-es}{0.12.1.7.2-fr}
\EGAReferenceAlias{0.12.1.7.2.fr}{0.12.1.7.2-fr}
\EGAReferenceAlias{0.12.1.7.2.es}{0.12.1.7.2-fr}
\EGAReferenceAlias{0.12.1.7.2.eq}{0.12.1.7.2-fr}
\EGAReferenceAlias{eq:0.12.1.7.2}{0.12.1.7.2-fr}
\EGAReferenceAlias{0.12.1.7.3}{0.12.1.7.3-fr}
\EGAReferenceAlias{0.12.1.7.3-es}{0.12.1.7.3-fr}
\EGAReferenceAlias{0.12.1.7.3.fr}{0.12.1.7.3-fr}
\EGAReferenceAlias{0.12.1.7.3.es}{0.12.1.7.3-fr}
\EGAReferenceAlias{0.12.1.7.3.eq}{0.12.1.7.3-fr}
\EGAReferenceAlias{eq:0.12.1.7.3}{0.12.1.7.3-fr}
\EGAReferenceAlias{0.12.1.7.4}{0.12.1.7.4-fr}
\EGAReferenceAlias{0.12.1.7.4-es}{0.12.1.7.4-fr}
\EGAReferenceAlias{0.12.1.7.4.fr}{0.12.1.7.4-fr}
\EGAReferenceAlias{0.12.1.7.4.es}{0.12.1.7.4-fr}
\EGAReferenceAlias{0.12.1.7.4.eq}{0.12.1.7.4-fr}
\EGAReferenceAlias{eq:0.12.1.7.4}{0.12.1.7.4-fr}
\EGAReferenceAlias{0.12.1.7.5}{0.12.1.7.5-fr}
\EGAReferenceAlias{0.12.1.7.5-es}{0.12.1.7.5-fr}
\EGAReferenceAlias{0.12.1.7.5.fr}{0.12.1.7.5-fr}
\EGAReferenceAlias{0.12.1.7.5.es}{0.12.1.7.5-fr}
\EGAReferenceAlias{0.12.1.7.5.eq}{0.12.1.7.5-fr}
\EGAReferenceAlias{eq:0.12.1.7.5}{0.12.1.7.5-fr}
\EGAReferenceAlias{0.12.2}{subsection:0.12.2-fr}
\EGAReferenceAlias{0.12.2-fr}{subsection:0.12.2-fr}
\EGAReferenceAlias{0.12.2-es}{subsection:0.12.2-fr}
\EGAReferenceAlias{0.12.2.fr}{subsection:0.12.2-fr}
\EGAReferenceAlias{0.12.2.es}{subsection:0.12.2-fr}
\EGAReferenceAlias{0.12.2.eq}{subsection:0.12.2-fr}
\EGAReferenceAlias{eq:0.12.2}{subsection:0.12.2-fr}
\EGAReferenceAlias{subsection:0.12.2}{subsection:0.12.2-fr}
\EGAReferenceAlias{subsection:0.12.2-es}{subsection:0.12.2-fr}
\EGAReferenceAlias{subsection:0.12.2.fr}{subsection:0.12.2-fr}
\EGAReferenceAlias{subsection:0.12.2.es}{subsection:0.12.2-fr}
\EGAReferenceAlias{subsection:0.12.2.eq}{subsection:0.12.2-fr}
\EGAReferenceAlias{eq:subsection:0.12.2}{subsection:0.12.2-fr}
\EGAReferenceAlias{0.12.2.1}{0.12.2.1-fr}
\EGAReferenceAlias{0.12.2.1-es}{0.12.2.1-fr}
\EGAReferenceAlias{0.12.2.1.fr}{0.12.2.1-fr}
\EGAReferenceAlias{0.12.2.1.es}{0.12.2.1-fr}
\EGAReferenceAlias{0.12.2.1.eq}{0.12.2.1-fr}
\EGAReferenceAlias{eq:0.12.2.1}{0.12.2.1-fr}
\EGAReferenceAlias{0.12.2.2-es}{0.12.2.2}
\EGAReferenceAlias{0.12.2.2.fr}{0.12.2.2}
\EGAReferenceAlias{0.12.2.2.es}{0.12.2.2}
\EGAReferenceAlias{0.12.2.2.eq}{0.12.2.2}
\EGAReferenceAlias{eq:0.12.2.2}{0.12.2.2}
\EGAReferenceAlias{0.12.2.2.1}{0.12.2.2.1-fr}
\EGAReferenceAlias{0.12.2.2.1-es}{0.12.2.2.1-fr}
\EGAReferenceAlias{0.12.2.2.1.fr}{0.12.2.2.1-fr}
\EGAReferenceAlias{0.12.2.2.1.es}{0.12.2.2.1-fr}
\EGAReferenceAlias{0.12.2.2.1.eq}{0.12.2.2.1-fr}
\EGAReferenceAlias{eq:0.12.2.2.1}{0.12.2.2.1-fr}
\EGAReferenceAlias{0.12.2.3-es}{0.12.2.3}
\EGAReferenceAlias{0.12.2.3.fr}{0.12.2.3}
\EGAReferenceAlias{0.12.2.3.es}{0.12.2.3}
\EGAReferenceAlias{0.12.2.3.eq}{0.12.2.3}
\EGAReferenceAlias{eq:0.12.2.3}{0.12.2.3}
\EGAReferenceAlias{0.12.2.3.1}{0.12.2.3.1-fr}
\EGAReferenceAlias{0.12.2.3.1-es}{0.12.2.3.1-fr}
\EGAReferenceAlias{0.12.2.3.1.fr}{0.12.2.3.1-fr}
\EGAReferenceAlias{0.12.2.3.1.es}{0.12.2.3.1-fr}
\EGAReferenceAlias{0.12.2.3.1.eq}{0.12.2.3.1-fr}
\EGAReferenceAlias{eq:0.12.2.3.1}{0.12.2.3.1-fr}
\EGAReferenceAlias{0.12.2.3.2}{0.12.2.3.2-fr}
\EGAReferenceAlias{0.12.2.3.2-es}{0.12.2.3.2-fr}
\EGAReferenceAlias{0.12.2.3.2.fr}{0.12.2.3.2-fr}
\EGAReferenceAlias{0.12.2.3.2.es}{0.12.2.3.2-fr}
\EGAReferenceAlias{0.12.2.3.2.eq}{0.12.2.3.2-fr}
\EGAReferenceAlias{eq:0.12.2.3.2}{0.12.2.3.2-fr}
\EGAReferenceAlias{0.12.2.4-es}{0.12.2.4}
\EGAReferenceAlias{0.12.2.4.fr}{0.12.2.4}
\EGAReferenceAlias{0.12.2.4.es}{0.12.2.4}
\EGAReferenceAlias{0.12.2.4.eq}{0.12.2.4}
\EGAReferenceAlias{eq:0.12.2.4}{0.12.2.4}
\EGAReferenceAlias{0.12.2.4.1}{0.12.2.4.1-fr}
\EGAReferenceAlias{0.12.2.4.1-es}{0.12.2.4.1-fr}
\EGAReferenceAlias{0.12.2.4.1.fr}{0.12.2.4.1-fr}
\EGAReferenceAlias{0.12.2.4.1.es}{0.12.2.4.1-fr}
\EGAReferenceAlias{0.12.2.4.1.eq}{0.12.2.4.1-fr}
\EGAReferenceAlias{eq:0.12.2.4.1}{0.12.2.4.1-fr}
\EGAReferenceAlias{0.12.2.5-es}{0.12.2.5}
\EGAReferenceAlias{0.12.2.5.fr}{0.12.2.5}
\EGAReferenceAlias{0.12.2.5.es}{0.12.2.5}
\EGAReferenceAlias{0.12.2.5.eq}{0.12.2.5}
\EGAReferenceAlias{eq:0.12.2.5}{0.12.2.5}
\EGAReferenceAlias{0.12.2.5.1}{0.12.2.5.1-fr}
\EGAReferenceAlias{0.12.2.5.1-es}{0.12.2.5.1-fr}
\EGAReferenceAlias{0.12.2.5.1.fr}{0.12.2.5.1-fr}
\EGAReferenceAlias{0.12.2.5.1.es}{0.12.2.5.1-fr}
\EGAReferenceAlias{0.12.2.5.1.eq}{0.12.2.5.1-fr}
\EGAReferenceAlias{eq:0.12.2.5.1}{0.12.2.5.1-fr}
\EGAReferenceAlias{0.12.2.5.2}{0.12.2.5.2-fr}
\EGAReferenceAlias{0.12.2.5.2-es}{0.12.2.5.2-fr}
\EGAReferenceAlias{0.12.2.5.2.fr}{0.12.2.5.2-fr}
\EGAReferenceAlias{0.12.2.5.2.es}{0.12.2.5.2-fr}
\EGAReferenceAlias{0.12.2.5.2.eq}{0.12.2.5.2-fr}
\EGAReferenceAlias{eq:0.12.2.5.2}{0.12.2.5.2-fr}
\EGAReferenceAlias{0.12.2.5.3}{0.12.2.5.3-fr}
\EGAReferenceAlias{0.12.2.5.3-es}{0.12.2.5.3-fr}
\EGAReferenceAlias{0.12.2.5.3.fr}{0.12.2.5.3-fr}
\EGAReferenceAlias{0.12.2.5.3.es}{0.12.2.5.3-fr}
\EGAReferenceAlias{0.12.2.5.3.eq}{0.12.2.5.3-fr}
\EGAReferenceAlias{eq:0.12.2.5.3}{0.12.2.5.3-fr}
\EGAReferenceAlias{0.12.2.5.4}{0.12.2.5.4-fr}
\EGAReferenceAlias{0.12.2.5.4-es}{0.12.2.5.4-fr}
\EGAReferenceAlias{0.12.2.5.4.fr}{0.12.2.5.4-fr}
\EGAReferenceAlias{0.12.2.5.4.es}{0.12.2.5.4-fr}
\EGAReferenceAlias{0.12.2.5.4.eq}{0.12.2.5.4-fr}
\EGAReferenceAlias{eq:0.12.2.5.4}{0.12.2.5.4-fr}
\EGAReferenceAlias{0.12.2.6-es}{0.12.2.6}
\EGAReferenceAlias{0.12.2.6.fr}{0.12.2.6}
\EGAReferenceAlias{0.12.2.6.es}{0.12.2.6}
\EGAReferenceAlias{0.12.2.6.eq}{0.12.2.6}
\EGAReferenceAlias{eq:0.12.2.6}{0.12.2.6}
\EGAReferenceAlias{0.12.2.6.1}{0.12.2.6.1-fr}
\EGAReferenceAlias{0.12.2.6.1-es}{0.12.2.6.1-fr}
\EGAReferenceAlias{0.12.2.6.1.fr}{0.12.2.6.1-fr}
\EGAReferenceAlias{0.12.2.6.1.es}{0.12.2.6.1-fr}
\EGAReferenceAlias{0.12.2.6.1.eq}{0.12.2.6.1-fr}
\EGAReferenceAlias{eq:0.12.2.6.1}{0.12.2.6.1-fr}
\EGAReferenceAlias{0.12.2.6.2}{0.12.2.6.2-fr}
\EGAReferenceAlias{0.12.2.6.2-es}{0.12.2.6.2-fr}
\EGAReferenceAlias{0.12.2.6.2.fr}{0.12.2.6.2-fr}
\EGAReferenceAlias{0.12.2.6.2.es}{0.12.2.6.2-fr}
\EGAReferenceAlias{0.12.2.6.2.eq}{0.12.2.6.2-fr}
\EGAReferenceAlias{eq:0.12.2.6.2}{0.12.2.6.2-fr}
\EGAReferenceAlias{0.12.3}{subsection:0.12.3}
\EGAReferenceAlias{0.12.3-fr}{subsection:0.12.3}
\EGAReferenceAlias{0.12.3-es}{subsection:0.12.3}
\EGAReferenceAlias{0.12.3.fr}{subsection:0.12.3}
\EGAReferenceAlias{0.12.3.es}{subsection:0.12.3}
\EGAReferenceAlias{0.12.3.eq}{subsection:0.12.3}
\EGAReferenceAlias{eq:0.12.3}{subsection:0.12.3}
\EGAReferenceAlias{subsection:0.12.3-es}{subsection:0.12.3}
\EGAReferenceAlias{subsection:0.12.3.fr}{subsection:0.12.3}
\EGAReferenceAlias{subsection:0.12.3.es}{subsection:0.12.3}
\EGAReferenceAlias{subsection:0.12.3.eq}{subsection:0.12.3}
\EGAReferenceAlias{eq:subsection:0.12.3}{subsection:0.12.3}
\EGAReferenceAlias{0.12.3.1}{0.12.3.1-fr}
\EGAReferenceAlias{0.12.3.1-es}{0.12.3.1-fr}
\EGAReferenceAlias{0.12.3.1.fr}{0.12.3.1-fr}
\EGAReferenceAlias{0.12.3.1.es}{0.12.3.1-fr}
\EGAReferenceAlias{0.12.3.1.eq}{0.12.3.1-fr}
\EGAReferenceAlias{eq:0.12.3.1}{0.12.3.1-fr}
\EGAReferenceAlias{0.12.3.2-es}{0.12.3.2}
\EGAReferenceAlias{0.12.3.2.fr}{0.12.3.2}
\EGAReferenceAlias{0.12.3.2.es}{0.12.3.2}
\EGAReferenceAlias{0.12.3.2.eq}{0.12.3.2}
\EGAReferenceAlias{eq:0.12.3.2}{0.12.3.2}
\EGAReferenceAlias{0.12.3.3-es}{0.12.3.3}
\EGAReferenceAlias{0.12.3.3.fr}{0.12.3.3}
\EGAReferenceAlias{0.12.3.3.es}{0.12.3.3}
\EGAReferenceAlias{0.12.3.3.eq}{0.12.3.3}
\EGAReferenceAlias{eq:0.12.3.3}{0.12.3.3}
\EGAReferenceAlias{0.12.3.4-es}{0.12.3.4}
\EGAReferenceAlias{0.12.3.4.fr}{0.12.3.4}
\EGAReferenceAlias{0.12.3.4.es}{0.12.3.4}
\EGAReferenceAlias{0.12.3.4.eq}{0.12.3.4}
\EGAReferenceAlias{eq:0.12.3.4}{0.12.3.4}
\EGAReferenceAlias{0.12.3.4.1}{0.12.3.4.1-fr}
\EGAReferenceAlias{0.12.3.4.1-es}{0.12.3.4.1-fr}
\EGAReferenceAlias{0.12.3.4.1.fr}{0.12.3.4.1-fr}
\EGAReferenceAlias{0.12.3.4.1.es}{0.12.3.4.1-fr}
\EGAReferenceAlias{0.12.3.4.1.eq}{0.12.3.4.1-fr}
\EGAReferenceAlias{eq:0.12.3.4.1}{0.12.3.4.1-fr}
\EGAReferenceAlias{0.12.3.4.2}{0.12.3.4.2-fr}
\EGAReferenceAlias{0.12.3.4.2-es}{0.12.3.4.2-fr}
\EGAReferenceAlias{0.12.3.4.2.fr}{0.12.3.4.2-fr}
\EGAReferenceAlias{0.12.3.4.2.es}{0.12.3.4.2-fr}
\EGAReferenceAlias{0.12.3.4.2.eq}{0.12.3.4.2-fr}
\EGAReferenceAlias{eq:0.12.3.4.2}{0.12.3.4.2-fr}
\EGAReferenceAlias{0.12.3.4.3}{0.12.3.4.3-fr}
\EGAReferenceAlias{0.12.3.4.3-es}{0.12.3.4.3-fr}
\EGAReferenceAlias{0.12.3.4.3.fr}{0.12.3.4.3-fr}
\EGAReferenceAlias{0.12.3.4.3.es}{0.12.3.4.3-fr}
\EGAReferenceAlias{0.12.3.4.3.eq}{0.12.3.4.3-fr}
\EGAReferenceAlias{eq:0.12.3.4.3}{0.12.3.4.3-fr}
\EGAReferenceAlias{0.12.3.5-es}{0.12.3.5}
\EGAReferenceAlias{0.12.3.5.fr}{0.12.3.5}
\EGAReferenceAlias{0.12.3.5.es}{0.12.3.5}
\EGAReferenceAlias{0.12.3.5.eq}{0.12.3.5}
\EGAReferenceAlias{eq:0.12.3.5}{0.12.3.5}
\EGAReferenceAlias{0.12.4}{subsection:0.12.4}
\EGAReferenceAlias{0.12.4-fr}{subsection:0.12.4}
\EGAReferenceAlias{0.12.4-es}{subsection:0.12.4}
\EGAReferenceAlias{0.12.4.fr}{subsection:0.12.4}
\EGAReferenceAlias{0.12.4.es}{subsection:0.12.4}
\EGAReferenceAlias{0.12.4.eq}{subsection:0.12.4}
\EGAReferenceAlias{eq:0.12.4}{subsection:0.12.4}
\EGAReferenceAlias{subsection:0.12.4-es}{subsection:0.12.4}
\EGAReferenceAlias{subsection:0.12.4.fr}{subsection:0.12.4}
\EGAReferenceAlias{subsection:0.12.4.es}{subsection:0.12.4}
\EGAReferenceAlias{subsection:0.12.4.eq}{subsection:0.12.4}
\EGAReferenceAlias{eq:subsection:0.12.4}{subsection:0.12.4}
\EGAReferenceAlias{0.12.4.1-es}{0.12.4.1}
\EGAReferenceAlias{0.12.4.1.fr}{0.12.4.1}
\EGAReferenceAlias{0.12.4.1.es}{0.12.4.1}
\EGAReferenceAlias{0.12.4.1.eq}{0.12.4.1}
\EGAReferenceAlias{eq:0.12.4.1}{0.12.4.1}
\EGAReferenceAlias{0.12.4.1.1}{0.12.4.1.1-fr}
\EGAReferenceAlias{0.12.4.1.1-es}{0.12.4.1.1-fr}
\EGAReferenceAlias{0.12.4.1.1.fr}{0.12.4.1.1-fr}
\EGAReferenceAlias{0.12.4.1.1.es}{0.12.4.1.1-fr}
\EGAReferenceAlias{0.12.4.1.1.eq}{0.12.4.1.1-fr}
\EGAReferenceAlias{eq:0.12.4.1.1}{0.12.4.1.1-fr}
\EGAReferenceAlias{0.12.4.1.2}{0.12.4.1.2-fr}
\EGAReferenceAlias{0.12.4.1.2-es}{0.12.4.1.2-fr}
\EGAReferenceAlias{0.12.4.1.2.fr}{0.12.4.1.2-fr}
\EGAReferenceAlias{0.12.4.1.2.es}{0.12.4.1.2-fr}
\EGAReferenceAlias{0.12.4.1.2.eq}{0.12.4.1.2-fr}
\EGAReferenceAlias{eq:0.12.4.1.2}{0.12.4.1.2-fr}
\EGAReferenceAlias{0.12.4.2-es}{0.12.4.2}
\EGAReferenceAlias{0.12.4.2.fr}{0.12.4.2}
\EGAReferenceAlias{0.12.4.2.es}{0.12.4.2}
\EGAReferenceAlias{0.12.4.2.eq}{0.12.4.2}
\EGAReferenceAlias{eq:0.12.4.2}{0.12.4.2}
\EGAReferenceAlias{0.12.4.3-es}{0.12.4.3}
\EGAReferenceAlias{0.12.4.3.fr}{0.12.4.3}
\EGAReferenceAlias{0.12.4.3.es}{0.12.4.3}
\EGAReferenceAlias{0.12.4.3.eq}{0.12.4.3}
\EGAReferenceAlias{eq:0.12.4.3}{0.12.4.3}
\EGAReferenceAlias{0.12.4.4-es}{0.12.4.4}
\EGAReferenceAlias{0.12.4.4.fr}{0.12.4.4}
\EGAReferenceAlias{0.12.4.4.es}{0.12.4.4}
\EGAReferenceAlias{0.12.4.4.eq}{0.12.4.4}
\EGAReferenceAlias{eq:0.12.4.4}{0.12.4.4}
\EGAReferenceAlias{0.12.4.5}{0.12.4.5-fr}
\EGAReferenceAlias{0.12.4.5-es}{0.12.4.5-fr}
\EGAReferenceAlias{0.12.4.5.fr}{0.12.4.5-fr}
\EGAReferenceAlias{0.12.4.5.es}{0.12.4.5-fr}
\EGAReferenceAlias{0.12.4.5.eq}{0.12.4.5-fr}
\EGAReferenceAlias{eq:0.12.4.5}{0.12.4.5-fr}
\EGAReferenceAlias{0.12.4.6-es}{0.12.4.6}
\EGAReferenceAlias{0.12.4.6.fr}{0.12.4.6}
\EGAReferenceAlias{0.12.4.6.es}{0.12.4.6}
\EGAReferenceAlias{0.12.4.6.eq}{0.12.4.6}
\EGAReferenceAlias{eq:0.12.4.6}{0.12.4.6}
\EGAReferenceAlias{0.12.4.6.1}{0.12.4.6.1-fr}
\EGAReferenceAlias{0.12.4.6.1-es}{0.12.4.6.1-fr}
\EGAReferenceAlias{0.12.4.6.1.fr}{0.12.4.6.1-fr}
\EGAReferenceAlias{0.12.4.6.1.es}{0.12.4.6.1-fr}
\EGAReferenceAlias{0.12.4.6.1.eq}{0.12.4.6.1-fr}
\EGAReferenceAlias{eq:0.12.4.6.1}{0.12.4.6.1-fr}
\EGAReferenceAlias{0.12.4.7-es}{0.12.4.7}
\EGAReferenceAlias{0.12.4.7.fr}{0.12.4.7}
\EGAReferenceAlias{0.12.4.7.es}{0.12.4.7}
\EGAReferenceAlias{0.12.4.7.eq}{0.12.4.7}
\EGAReferenceAlias{eq:0.12.4.7}{0.12.4.7}
\EGAReferenceAlias{0.12.4.7.1}{0.12.4.7.1-fr}
\EGAReferenceAlias{0.12.4.7.1-es}{0.12.4.7.1-fr}
\EGAReferenceAlias{0.12.4.7.1.fr}{0.12.4.7.1-fr}
\EGAReferenceAlias{0.12.4.7.1.es}{0.12.4.7.1-fr}
\EGAReferenceAlias{0.12.4.7.1.eq}{0.12.4.7.1-fr}
\EGAReferenceAlias{eq:0.12.4.7.1}{0.12.4.7.1-fr}
\EGAReferenceAlias{0.12.4.8-es}{0.12.4.8}
\EGAReferenceAlias{0.12.4.8.fr}{0.12.4.8}
\EGAReferenceAlias{0.12.4.8.es}{0.12.4.8}
\EGAReferenceAlias{0.12.4.8.eq}{0.12.4.8}
\EGAReferenceAlias{eq:0.12.4.8}{0.12.4.8}
\EGAReferenceAlias{0.12.4.8.1}{0.12.4.8.1-fr}
\EGAReferenceAlias{0.12.4.8.1-es}{0.12.4.8.1-fr}
\EGAReferenceAlias{0.12.4.8.1.fr}{0.12.4.8.1-fr}
\EGAReferenceAlias{0.12.4.8.1.es}{0.12.4.8.1-fr}
\EGAReferenceAlias{0.12.4.8.1.eq}{0.12.4.8.1-fr}
\EGAReferenceAlias{eq:0.12.4.8.1}{0.12.4.8.1-fr}
\EGAReferenceAlias{0.12.4.8.2}{0.12.4.8.2-fr}
\EGAReferenceAlias{0.12.4.8.2-es}{0.12.4.8.2-fr}
\EGAReferenceAlias{0.12.4.8.2.fr}{0.12.4.8.2-fr}
\EGAReferenceAlias{0.12.4.8.2.es}{0.12.4.8.2-fr}
\EGAReferenceAlias{0.12.4.8.2.eq}{0.12.4.8.2-fr}
\EGAReferenceAlias{eq:0.12.4.8.2}{0.12.4.8.2-fr}
\EGAReferenceAlias{0.12.4.8.3}{0.12.4.8.3-fr}
\EGAReferenceAlias{0.12.4.8.3-es}{0.12.4.8.3-fr}
\EGAReferenceAlias{0.12.4.8.3.fr}{0.12.4.8.3-fr}
\EGAReferenceAlias{0.12.4.8.3.es}{0.12.4.8.3-fr}
\EGAReferenceAlias{0.12.4.8.3.eq}{0.12.4.8.3-fr}
\EGAReferenceAlias{eq:0.12.4.8.3}{0.12.4.8.3-fr}
\EGAReferenceAlias{0.13}{section:0.13}
\EGAReferenceAlias{0.13-fr}{section:0.13}
\EGAReferenceAlias{0.13-es}{section:0.13}
\EGAReferenceAlias{0.13.fr}{section:0.13}
\EGAReferenceAlias{0.13.es}{section:0.13}
\EGAReferenceAlias{0.13.eq}{section:0.13}
\EGAReferenceAlias{eq:0.13}{section:0.13}
\EGAReferenceAlias{section:0.13-es}{section:0.13}
\EGAReferenceAlias{section:0.13.fr}{section:0.13}
\EGAReferenceAlias{section:0.13.es}{section:0.13}
\EGAReferenceAlias{section:0.13.eq}{section:0.13}
\EGAReferenceAlias{eq:section:0.13}{section:0.13}
\EGAReferenceAlias{0.13.1}{subsection:0.13.1}
\EGAReferenceAlias{0.13.1-fr}{subsection:0.13.1}
\EGAReferenceAlias{0.13.1-es}{subsection:0.13.1}
\EGAReferenceAlias{0.13.1.fr}{subsection:0.13.1}
\EGAReferenceAlias{0.13.1.es}{subsection:0.13.1}
\EGAReferenceAlias{0.13.1.eq}{subsection:0.13.1}
\EGAReferenceAlias{eq:0.13.1}{subsection:0.13.1}
\EGAReferenceAlias{subsection:0.13.1-es}{subsection:0.13.1}
\EGAReferenceAlias{subsection:0.13.1.fr}{subsection:0.13.1}
\EGAReferenceAlias{subsection:0.13.1.es}{subsection:0.13.1}
\EGAReferenceAlias{subsection:0.13.1.eq}{subsection:0.13.1}
\EGAReferenceAlias{eq:subsection:0.13.1}{subsection:0.13.1}
\EGAReferenceAlias{0.13.1.1-es}{0.13.1.1}
\EGAReferenceAlias{0.13.1.1.fr}{0.13.1.1}
\EGAReferenceAlias{0.13.1.1.es}{0.13.1.1}
\EGAReferenceAlias{0.13.1.1.eq}{0.13.1.1}
\EGAReferenceAlias{eq:0.13.1.1}{0.13.1.1}
\EGAReferenceAlias{0.13.1.2-es}{0.13.1.2}
\EGAReferenceAlias{0.13.1.2.fr}{0.13.1.2}
\EGAReferenceAlias{0.13.1.2.es}{0.13.1.2}
\EGAReferenceAlias{0.13.1.2.eq}{0.13.1.2}
\EGAReferenceAlias{eq:0.13.1.2}{0.13.1.2}
\EGAReferenceAlias{0.13.1.3-es}{0.13.1.3}
\EGAReferenceAlias{0.13.1.3.fr}{0.13.1.3}
\EGAReferenceAlias{0.13.1.3.es}{0.13.1.3}
\EGAReferenceAlias{0.13.1.3.eq}{0.13.1.3}
\EGAReferenceAlias{eq:0.13.1.3}{0.13.1.3}
\EGAReferenceAlias{0.13.2}{subsection:0.13.2}
\EGAReferenceAlias{0.13.2-fr}{subsection:0.13.2}
\EGAReferenceAlias{0.13.2-es}{subsection:0.13.2}
\EGAReferenceAlias{0.13.2.fr}{subsection:0.13.2}
\EGAReferenceAlias{0.13.2.es}{subsection:0.13.2}
\EGAReferenceAlias{0.13.2.eq}{subsection:0.13.2}
\EGAReferenceAlias{eq:0.13.2}{subsection:0.13.2}
\EGAReferenceAlias{subsection:0.13.2-es}{subsection:0.13.2}
\EGAReferenceAlias{subsection:0.13.2.fr}{subsection:0.13.2}
\EGAReferenceAlias{subsection:0.13.2.es}{subsection:0.13.2}
\EGAReferenceAlias{subsection:0.13.2.eq}{subsection:0.13.2}
\EGAReferenceAlias{eq:subsection:0.13.2}{subsection:0.13.2}
\EGAReferenceAlias{0.13.2.1-es}{0.13.2.1}
\EGAReferenceAlias{0.13.2.1.fr}{0.13.2.1}
\EGAReferenceAlias{0.13.2.1.es}{0.13.2.1}
\EGAReferenceAlias{0.13.2.1.eq}{0.13.2.1}
\EGAReferenceAlias{eq:0.13.2.1}{0.13.2.1}
\EGAReferenceAlias{0.13.2.2-es}{0.13.2.2}
\EGAReferenceAlias{0.13.2.2.fr}{0.13.2.2}
\EGAReferenceAlias{0.13.2.2.es}{0.13.2.2}
\EGAReferenceAlias{0.13.2.2.eq}{0.13.2.2}
\EGAReferenceAlias{eq:0.13.2.2}{0.13.2.2}
\EGAReferenceAlias{0.13.2.3-es}{0.13.2.3}
\EGAReferenceAlias{0.13.2.3.fr}{0.13.2.3}
\EGAReferenceAlias{0.13.2.3.es}{0.13.2.3}
\EGAReferenceAlias{0.13.2.3.eq}{0.13.2.3}
\EGAReferenceAlias{eq:0.13.2.3}{0.13.2.3}
\EGAReferenceAlias{0.13.2.3.1}{0.13.2.3.1-fr}
\EGAReferenceAlias{0.13.2.3.1-es}{0.13.2.3.1-fr}
\EGAReferenceAlias{0.13.2.3.1.fr}{0.13.2.3.1-fr}
\EGAReferenceAlias{0.13.2.3.1.es}{0.13.2.3.1-fr}
\EGAReferenceAlias{0.13.2.3.1.eq}{0.13.2.3.1-fr}
\EGAReferenceAlias{eq:0.13.2.3.1}{0.13.2.3.1-fr}
\EGAReferenceAlias{0.13.2.4-es}{0.13.2.4}
\EGAReferenceAlias{0.13.2.4.fr}{0.13.2.4}
\EGAReferenceAlias{0.13.2.4.es}{0.13.2.4}
\EGAReferenceAlias{0.13.2.4.eq}{0.13.2.4}
\EGAReferenceAlias{eq:0.13.2.4}{0.13.2.4}
\EGAReferenceAlias{0.13.3}{subsection:0.13.3}
\EGAReferenceAlias{0.13.3-fr}{subsection:0.13.3}
\EGAReferenceAlias{0.13.3-es}{subsection:0.13.3}
\EGAReferenceAlias{0.13.3.fr}{subsection:0.13.3}
\EGAReferenceAlias{0.13.3.es}{subsection:0.13.3}
\EGAReferenceAlias{0.13.3.eq}{subsection:0.13.3}
\EGAReferenceAlias{eq:0.13.3}{subsection:0.13.3}
\EGAReferenceAlias{subsection:0.13.3-es}{subsection:0.13.3}
\EGAReferenceAlias{subsection:0.13.3.fr}{subsection:0.13.3}
\EGAReferenceAlias{subsection:0.13.3.es}{subsection:0.13.3}
\EGAReferenceAlias{subsection:0.13.3.eq}{subsection:0.13.3}
\EGAReferenceAlias{eq:subsection:0.13.3}{subsection:0.13.3}
\EGAReferenceAlias{0.13.3.1-es}{0.13.3.1}
\EGAReferenceAlias{0.13.3.1.fr}{0.13.3.1}
\EGAReferenceAlias{0.13.3.1.es}{0.13.3.1}
\EGAReferenceAlias{0.13.3.1.eq}{0.13.3.1}
\EGAReferenceAlias{eq:0.13.3.1}{0.13.3.1}
\EGAReferenceAlias{0.13.3.2-es}{0.13.3.2}
\EGAReferenceAlias{0.13.3.2.fr}{0.13.3.2}
\EGAReferenceAlias{0.13.3.2.es}{0.13.3.2}
\EGAReferenceAlias{0.13.3.2.eq}{0.13.3.2}
\EGAReferenceAlias{eq:0.13.3.2}{0.13.3.2}
\EGAReferenceAlias{0.13.4-es}{subsection:0.13.4}
\EGAReferenceAlias{0.13.4.fr}{subsection:0.13.4}
\EGAReferenceAlias{0.13.4.es}{subsection:0.13.4}
\EGAReferenceAlias{0.13.4.eq}{subsection:0.13.4}
\EGAReferenceAlias{eq:0.13.4}{subsection:0.13.4}
\EGAReferenceAlias{subsection:0.13.4-es}{subsection:0.13.4}
\EGAReferenceAlias{subsection:0.13.4.fr}{subsection:0.13.4}
\EGAReferenceAlias{subsection:0.13.4.es}{subsection:0.13.4}
\EGAReferenceAlias{subsection:0.13.4.eq}{subsection:0.13.4}
\EGAReferenceAlias{eq:subsection:0.13.4}{subsection:0.13.4}
\EGAReferenceAlias{0.13.4.1-es}{0.13.4.1}
\EGAReferenceAlias{0.13.4.1.fr}{0.13.4.1}
\EGAReferenceAlias{0.13.4.1.es}{0.13.4.1}
\EGAReferenceAlias{0.13.4.1.eq}{0.13.4.1}
\EGAReferenceAlias{eq:0.13.4.1}{0.13.4.1}
\EGAReferenceAlias{0.13.4.1.1}{0.13.4.1.1-fr}
\EGAReferenceAlias{0.13.4.1.1-es}{0.13.4.1.1-fr}
\EGAReferenceAlias{0.13.4.1.1.fr}{0.13.4.1.1-fr}
\EGAReferenceAlias{0.13.4.1.1.es}{0.13.4.1.1-fr}
\EGAReferenceAlias{0.13.4.1.1.eq}{0.13.4.1.1-fr}
\EGAReferenceAlias{eq:0.13.4.1.1}{0.13.4.1.1-fr}
\EGAReferenceAlias{0.13.4.1.2}{0.13.4.1.2-fr}
\EGAReferenceAlias{0.13.4.1.2-es}{0.13.4.1.2-fr}
\EGAReferenceAlias{0.13.4.1.2.fr}{0.13.4.1.2-fr}
\EGAReferenceAlias{0.13.4.1.2.es}{0.13.4.1.2-fr}
\EGAReferenceAlias{0.13.4.1.2.eq}{0.13.4.1.2-fr}
\EGAReferenceAlias{eq:0.13.4.1.2}{0.13.4.1.2-fr}
\EGAReferenceAlias{0.13.4.1.3}{0.13.4.1.3-fr}
\EGAReferenceAlias{0.13.4.1.3-es}{0.13.4.1.3-fr}
\EGAReferenceAlias{0.13.4.1.3.fr}{0.13.4.1.3-fr}
\EGAReferenceAlias{0.13.4.1.3.es}{0.13.4.1.3-fr}
\EGAReferenceAlias{0.13.4.1.3.eq}{0.13.4.1.3-fr}
\EGAReferenceAlias{eq:0.13.4.1.3}{0.13.4.1.3-fr}
\EGAReferenceAlias{0.13.4.2-es}{0.13.4.2}
\EGAReferenceAlias{0.13.4.2.fr}{0.13.4.2}
\EGAReferenceAlias{0.13.4.2.es}{0.13.4.2}
\EGAReferenceAlias{0.13.4.2.eq}{0.13.4.2}
\EGAReferenceAlias{eq:0.13.4.2}{0.13.4.2}
\EGAReferenceAlias{0.13.4.3-es}{0.13.4.3}
\EGAReferenceAlias{0.13.4.3.fr}{0.13.4.3}
\EGAReferenceAlias{0.13.4.3.es}{0.13.4.3}
\EGAReferenceAlias{0.13.4.3.eq}{0.13.4.3}
\EGAReferenceAlias{eq:0.13.4.3}{0.13.4.3}
\EGAReferenceAlias{0.13.4.3.1}{0.13.4.3.1-fr}
\EGAReferenceAlias{0.13.4.3.1-es}{0.13.4.3.1-fr}
\EGAReferenceAlias{0.13.4.3.1.fr}{0.13.4.3.1-fr}
\EGAReferenceAlias{0.13.4.3.1.es}{0.13.4.3.1-fr}
\EGAReferenceAlias{0.13.4.3.1.eq}{0.13.4.3.1-fr}
\EGAReferenceAlias{eq:0.13.4.3.1}{0.13.4.3.1-fr}
\EGAReferenceAlias{0.13.4.3.2}{0.13.4.3.2-fr}
\EGAReferenceAlias{0.13.4.3.2-es}{0.13.4.3.2-fr}
\EGAReferenceAlias{0.13.4.3.2.fr}{0.13.4.3.2-fr}
\EGAReferenceAlias{0.13.4.3.2.es}{0.13.4.3.2-fr}
\EGAReferenceAlias{0.13.4.3.2.eq}{0.13.4.3.2-fr}
\EGAReferenceAlias{eq:0.13.4.3.2}{0.13.4.3.2-fr}
\EGAReferenceAlias{0.13.4.4}{0.13.4.4-fr}
\EGAReferenceAlias{0.13.4.4-es}{0.13.4.4-fr}
\EGAReferenceAlias{0.13.4.4.fr}{0.13.4.4-fr}
\EGAReferenceAlias{0.13.4.4.es}{0.13.4.4-fr}
\EGAReferenceAlias{0.13.4.4.eq}{0.13.4.4-fr}
\EGAReferenceAlias{eq:0.13.4.4}{0.13.4.4-fr}
\EGAReferenceAlias{0.13.4.4.1}{0.13.4.4.1-fr}
\EGAReferenceAlias{0.13.4.4.1-es}{0.13.4.4.1-fr}
\EGAReferenceAlias{0.13.4.4.1.fr}{0.13.4.4.1-fr}
\EGAReferenceAlias{0.13.4.4.1.es}{0.13.4.4.1-fr}
\EGAReferenceAlias{0.13.4.4.1.eq}{0.13.4.4.1-fr}
\EGAReferenceAlias{eq:0.13.4.4.1}{0.13.4.4.1-fr}
\EGAReferenceAlias{0.13.4.5-es}{0.13.4.5}
\EGAReferenceAlias{0.13.4.5.fr}{0.13.4.5}
\EGAReferenceAlias{0.13.4.5.es}{0.13.4.5}
\EGAReferenceAlias{0.13.4.5.eq}{0.13.4.5}
\EGAReferenceAlias{eq:0.13.4.5}{0.13.4.5}
\EGAReferenceAlias{0.13.4.5.1}{0.13.4.5.1-fr}
\EGAReferenceAlias{0.13.4.5.1-es}{0.13.4.5.1-fr}
\EGAReferenceAlias{0.13.4.5.1.fr}{0.13.4.5.1-fr}
\EGAReferenceAlias{0.13.4.5.1.es}{0.13.4.5.1-fr}
\EGAReferenceAlias{0.13.4.5.1.eq}{0.13.4.5.1-fr}
\EGAReferenceAlias{eq:0.13.4.5.1}{0.13.4.5.1-fr}
\EGAReferenceAlias{0.13.5}{subsection:0.13.5}
\EGAReferenceAlias{0.13.5-fr}{subsection:0.13.5}
\EGAReferenceAlias{0.13.5-es}{subsection:0.13.5}
\EGAReferenceAlias{0.13.5.fr}{subsection:0.13.5}
\EGAReferenceAlias{0.13.5.es}{subsection:0.13.5}
\EGAReferenceAlias{0.13.5.eq}{subsection:0.13.5}
\EGAReferenceAlias{eq:0.13.5}{subsection:0.13.5}
\EGAReferenceAlias{subsection:0.13.5-es}{subsection:0.13.5}
\EGAReferenceAlias{subsection:0.13.5.fr}{subsection:0.13.5}
\EGAReferenceAlias{subsection:0.13.5.es}{subsection:0.13.5}
\EGAReferenceAlias{subsection:0.13.5.eq}{subsection:0.13.5}
\EGAReferenceAlias{eq:subsection:0.13.5}{subsection:0.13.5}
\EGAReferenceAlias{0.13.5.1-es}{0.13.5.1}
\EGAReferenceAlias{0.13.5.1.fr}{0.13.5.1}
\EGAReferenceAlias{0.13.5.1.es}{0.13.5.1}
\EGAReferenceAlias{0.13.5.1.eq}{0.13.5.1}
\EGAReferenceAlias{eq:0.13.5.1}{0.13.5.1}
\EGAReferenceAlias{0.13.5.2-es}{0.13.5.2}
\EGAReferenceAlias{0.13.5.2.fr}{0.13.5.2}
\EGAReferenceAlias{0.13.5.2.es}{0.13.5.2}
\EGAReferenceAlias{0.13.5.2.eq}{0.13.5.2}
\EGAReferenceAlias{eq:0.13.5.2}{0.13.5.2}
\EGAReferenceAlias{0.13.5.3-es}{0.13.5.3}
\EGAReferenceAlias{0.13.5.3.fr}{0.13.5.3}
\EGAReferenceAlias{0.13.5.3.es}{0.13.5.3}
\EGAReferenceAlias{0.13.5.3.eq}{0.13.5.3}
\EGAReferenceAlias{eq:0.13.5.3}{0.13.5.3}
\EGAReferenceAlias{0.13.5.4}{0.13.5.4-fr}
\EGAReferenceAlias{0.13.5.4-es}{0.13.5.4-fr}
\EGAReferenceAlias{0.13.5.4.fr}{0.13.5.4-fr}
\EGAReferenceAlias{0.13.5.4.es}{0.13.5.4-fr}
\EGAReferenceAlias{0.13.5.4.eq}{0.13.5.4-fr}
\EGAReferenceAlias{eq:0.13.5.4}{0.13.5.4-fr}
\EGAReferenceAlias{0.13.5.4.1}{0.13.5.4.1-fr}
\EGAReferenceAlias{0.13.5.4.1-es}{0.13.5.4.1-fr}
\EGAReferenceAlias{0.13.5.4.1.fr}{0.13.5.4.1-fr}
\EGAReferenceAlias{0.13.5.4.1.es}{0.13.5.4.1-fr}
\EGAReferenceAlias{0.13.5.4.1.eq}{0.13.5.4.1-fr}
\EGAReferenceAlias{eq:0.13.5.4.1}{0.13.5.4.1-fr}
\EGAReferenceAlias{0.13.5.5-es}{0.13.5.5}
\EGAReferenceAlias{0.13.5.5.fr}{0.13.5.5}
\EGAReferenceAlias{0.13.5.5.es}{0.13.5.5}
\EGAReferenceAlias{0.13.5.5.eq}{0.13.5.5}
\EGAReferenceAlias{eq:0.13.5.5}{0.13.5.5}
\EGAReferenceAlias{0.13.5.5.1}{0.13.5.5.1-fr}
\EGAReferenceAlias{0.13.5.5.1-es}{0.13.5.5.1-fr}
\EGAReferenceAlias{0.13.5.5.1.fr}{0.13.5.5.1-fr}
\EGAReferenceAlias{0.13.5.5.1.es}{0.13.5.5.1-fr}
\EGAReferenceAlias{0.13.5.5.1.eq}{0.13.5.5.1-fr}
\EGAReferenceAlias{eq:0.13.5.5.1}{0.13.5.5.1-fr}
\EGAReferenceAlias{0.13.5.5.2}{0.13.5.5.2-fr}
\EGAReferenceAlias{0.13.5.5.2-es}{0.13.5.5.2-fr}
\EGAReferenceAlias{0.13.5.5.2.fr}{0.13.5.5.2-fr}
\EGAReferenceAlias{0.13.5.5.2.es}{0.13.5.5.2-fr}
\EGAReferenceAlias{0.13.5.5.2.eq}{0.13.5.5.2-fr}
\EGAReferenceAlias{eq:0.13.5.5.2}{0.13.5.5.2-fr}
\EGAReferenceAlias{0.13.5.5.3}{0.13.5.5.3-fr}
\EGAReferenceAlias{0.13.5.5.3-es}{0.13.5.5.3-fr}
\EGAReferenceAlias{0.13.5.5.3.fr}{0.13.5.5.3-fr}
\EGAReferenceAlias{0.13.5.5.3.es}{0.13.5.5.3-fr}
\EGAReferenceAlias{0.13.5.5.3.eq}{0.13.5.5.3-fr}
\EGAReferenceAlias{eq:0.13.5.5.3}{0.13.5.5.3-fr}
\EGAReferenceAlias{0.13.5.6}{0.13.5.6-fr}
\EGAReferenceAlias{0.13.5.6-es}{0.13.5.6-fr}
\EGAReferenceAlias{0.13.5.6.fr}{0.13.5.6-fr}
\EGAReferenceAlias{0.13.5.6.es}{0.13.5.6-fr}
\EGAReferenceAlias{0.13.5.6.eq}{0.13.5.6-fr}
\EGAReferenceAlias{eq:0.13.5.6}{0.13.5.6-fr}
\EGAReferenceAlias{0.13.5.6.1}{0.13.5.6.1-fr}
\EGAReferenceAlias{0.13.5.6.1-es}{0.13.5.6.1-fr}
\EGAReferenceAlias{0.13.5.6.1.fr}{0.13.5.6.1-fr}
\EGAReferenceAlias{0.13.5.6.1.es}{0.13.5.6.1-fr}
\EGAReferenceAlias{0.13.5.6.1.eq}{0.13.5.6.1-fr}
\EGAReferenceAlias{eq:0.13.5.6.1}{0.13.5.6.1-fr}
\EGAReferenceAlias{0.13.5.7-es}{0.13.5.7}
\EGAReferenceAlias{0.13.5.7.fr}{0.13.5.7}
\EGAReferenceAlias{0.13.5.7.es}{0.13.5.7}
\EGAReferenceAlias{0.13.5.7.eq}{0.13.5.7}
\EGAReferenceAlias{eq:0.13.5.7}{0.13.5.7}
\EGAReferenceAlias{0.13.6}{subsection:0.13.6}
\EGAReferenceAlias{0.13.6-fr}{subsection:0.13.6}
\EGAReferenceAlias{0.13.6-es}{subsection:0.13.6}
\EGAReferenceAlias{0.13.6.fr}{subsection:0.13.6}
\EGAReferenceAlias{0.13.6.es}{subsection:0.13.6}
\EGAReferenceAlias{0.13.6.eq}{subsection:0.13.6}
\EGAReferenceAlias{eq:0.13.6}{subsection:0.13.6}
\EGAReferenceAlias{subsection:0.13.6-es}{subsection:0.13.6}
\EGAReferenceAlias{subsection:0.13.6.fr}{subsection:0.13.6}
\EGAReferenceAlias{subsection:0.13.6.es}{subsection:0.13.6}
\EGAReferenceAlias{subsection:0.13.6.eq}{subsection:0.13.6}
\EGAReferenceAlias{eq:subsection:0.13.6}{subsection:0.13.6}
\EGAReferenceAlias{0.13.6.1-es}{0.13.6.1}
\EGAReferenceAlias{0.13.6.1.fr}{0.13.6.1}
\EGAReferenceAlias{0.13.6.1.es}{0.13.6.1}
\EGAReferenceAlias{0.13.6.1.eq}{0.13.6.1}
\EGAReferenceAlias{eq:0.13.6.1}{0.13.6.1}
\EGAReferenceAlias{0.13.6.2-es}{0.13.6.2}
\EGAReferenceAlias{0.13.6.2.fr}{0.13.6.2}
\EGAReferenceAlias{0.13.6.2.es}{0.13.6.2}
\EGAReferenceAlias{0.13.6.2.eq}{0.13.6.2}
\EGAReferenceAlias{eq:0.13.6.2}{0.13.6.2}
\EGAReferenceAlias{0.13.6.3-es}{0.13.6.3}
\EGAReferenceAlias{0.13.6.3.fr}{0.13.6.3}
\EGAReferenceAlias{0.13.6.3.es}{0.13.6.3}
\EGAReferenceAlias{0.13.6.3.eq}{0.13.6.3}
\EGAReferenceAlias{eq:0.13.6.3}{0.13.6.3}
\EGAReferenceAlias{0.13.6.4-es}{0.13.6.4}
\EGAReferenceAlias{0.13.6.4.fr}{0.13.6.4}
\EGAReferenceAlias{0.13.6.4.es}{0.13.6.4}
\EGAReferenceAlias{0.13.6.4.eq}{0.13.6.4}
\EGAReferenceAlias{eq:0.13.6.4}{0.13.6.4}
\EGAReferenceAlias{0.13.6.4.1}{0.13.6.4.1-fr}
\EGAReferenceAlias{0.13.6.4.1-es}{0.13.6.4.1-fr}
\EGAReferenceAlias{0.13.6.4.1.fr}{0.13.6.4.1-fr}
\EGAReferenceAlias{0.13.6.4.1.es}{0.13.6.4.1-fr}
\EGAReferenceAlias{0.13.6.4.1.eq}{0.13.6.4.1-fr}
\EGAReferenceAlias{eq:0.13.6.4.1}{0.13.6.4.1-fr}
\EGAReferenceAlias{0.13.6.4.2}{0.13.6.4.2-fr}
\EGAReferenceAlias{0.13.6.4.2-es}{0.13.6.4.2-fr}
\EGAReferenceAlias{0.13.6.4.2.fr}{0.13.6.4.2-fr}
\EGAReferenceAlias{0.13.6.4.2.es}{0.13.6.4.2-fr}
\EGAReferenceAlias{0.13.6.4.2.eq}{0.13.6.4.2-fr}
\EGAReferenceAlias{eq:0.13.6.4.2}{0.13.6.4.2-fr}
\EGAReferenceAlias{0.13.6.5-es}{0.13.6.5}
\EGAReferenceAlias{0.13.6.5.fr}{0.13.6.5}
\EGAReferenceAlias{0.13.6.5.es}{0.13.6.5}
\EGAReferenceAlias{0.13.6.5.eq}{0.13.6.5}
\EGAReferenceAlias{eq:0.13.6.5}{0.13.6.5}
\EGAReferenceAlias{0.13.6.6-es}{0.13.6.6}
\EGAReferenceAlias{0.13.6.6.fr}{0.13.6.6}
\EGAReferenceAlias{0.13.6.6.es}{0.13.6.6}
\EGAReferenceAlias{0.13.6.6.eq}{0.13.6.6}
\EGAReferenceAlias{eq:0.13.6.6}{0.13.6.6}
\EGAReferenceAlias{0.13.6.6.1}{0.13.6.6.1-fr}
\EGAReferenceAlias{0.13.6.6.1-es}{0.13.6.6.1-fr}
\EGAReferenceAlias{0.13.6.6.1.fr}{0.13.6.6.1-fr}
\EGAReferenceAlias{0.13.6.6.1.es}{0.13.6.6.1-fr}
\EGAReferenceAlias{0.13.6.6.1.eq}{0.13.6.6.1-fr}
\EGAReferenceAlias{eq:0.13.6.6.1}{0.13.6.6.1-fr}
\EGAReferenceAlias{0.13.7-fr}{subsection:0.13.7}
\EGAReferenceAlias{0.13.7-es}{subsection:0.13.7}
\EGAReferenceAlias{0.13.7.fr}{subsection:0.13.7}
\EGAReferenceAlias{0.13.7.es}{subsection:0.13.7}
\EGAReferenceAlias{0.13.7.eq}{subsection:0.13.7}
\EGAReferenceAlias{eq:0.13.7}{subsection:0.13.7}
\EGAReferenceAlias{subsection:0.13.7-es}{subsection:0.13.7}
\EGAReferenceAlias{subsection:0.13.7.fr}{subsection:0.13.7}
\EGAReferenceAlias{subsection:0.13.7.es}{subsection:0.13.7}
\EGAReferenceAlias{subsection:0.13.7.eq}{subsection:0.13.7}
\EGAReferenceAlias{eq:subsection:0.13.7}{subsection:0.13.7}
\EGAReferenceAlias{0.13.7.1-es}{0.13.7.1}
\EGAReferenceAlias{0.13.7.1.fr}{0.13.7.1}
\EGAReferenceAlias{0.13.7.1.es}{0.13.7.1}
\EGAReferenceAlias{0.13.7.1.eq}{0.13.7.1}
\EGAReferenceAlias{eq:0.13.7.1}{0.13.7.1}
\EGAReferenceAlias{0.13.7.1.1}{0.13.7.1.1-fr}
\EGAReferenceAlias{0.13.7.1.1-es}{0.13.7.1.1-fr}
\EGAReferenceAlias{0.13.7.1.1.fr}{0.13.7.1.1-fr}
\EGAReferenceAlias{0.13.7.1.1.es}{0.13.7.1.1-fr}
\EGAReferenceAlias{0.13.7.1.1.eq}{0.13.7.1.1-fr}
\EGAReferenceAlias{eq:0.13.7.1.1}{0.13.7.1.1-fr}
\EGAReferenceAlias{0.13.7.2-es}{0.13.7.2}
\EGAReferenceAlias{0.13.7.2.fr}{0.13.7.2}
\EGAReferenceAlias{0.13.7.2.es}{0.13.7.2}
\EGAReferenceAlias{0.13.7.2.eq}{0.13.7.2}
\EGAReferenceAlias{eq:0.13.7.2}{0.13.7.2}
\EGAReferenceAlias{0.13.7.3-es}{0.13.7.3}
\EGAReferenceAlias{0.13.7.3.fr}{0.13.7.3}
\EGAReferenceAlias{0.13.7.3.es}{0.13.7.3}
\EGAReferenceAlias{0.13.7.3.eq}{0.13.7.3}
\EGAReferenceAlias{eq:0.13.7.3}{0.13.7.3}
\EGAReferenceAlias{0.13.7.3.1}{0.13.7.3.1-fr}
\EGAReferenceAlias{0.13.7.3.1-es}{0.13.7.3.1-fr}
\EGAReferenceAlias{0.13.7.3.1.fr}{0.13.7.3.1-fr}
\EGAReferenceAlias{0.13.7.3.1.es}{0.13.7.3.1-fr}
\EGAReferenceAlias{0.13.7.3.1.eq}{0.13.7.3.1-fr}
\EGAReferenceAlias{eq:0.13.7.3.1}{0.13.7.3.1-fr}
\EGAReferenceAlias{0.13.7.4-es}{0.13.7.4}
\EGAReferenceAlias{0.13.7.4.fr}{0.13.7.4}
\EGAReferenceAlias{0.13.7.4.es}{0.13.7.4}
\EGAReferenceAlias{0.13.7.4.eq}{0.13.7.4}
\EGAReferenceAlias{eq:0.13.7.4}{0.13.7.4}
\EGAReferenceAlias{0.13.7.4.1}{0.13.7.4.1-fr}
\EGAReferenceAlias{0.13.7.4.1-es}{0.13.7.4.1-fr}
\EGAReferenceAlias{0.13.7.4.1.fr}{0.13.7.4.1-fr}
\EGAReferenceAlias{0.13.7.4.1.es}{0.13.7.4.1-fr}
\EGAReferenceAlias{0.13.7.4.1.eq}{0.13.7.4.1-fr}
\EGAReferenceAlias{eq:0.13.7.4.1}{0.13.7.4.1-fr}
\EGAReferenceAlias{0.13.7.5-es}{0.13.7.5}
\EGAReferenceAlias{0.13.7.5.fr}{0.13.7.5}
\EGAReferenceAlias{0.13.7.5.es}{0.13.7.5}
\EGAReferenceAlias{0.13.7.5.eq}{0.13.7.5}
\EGAReferenceAlias{eq:0.13.7.5}{0.13.7.5}
\EGAReferenceAlias{0.13.7.5.1}{0.13.7.5.1-fr}
\EGAReferenceAlias{0.13.7.5.1-es}{0.13.7.5.1-fr}
\EGAReferenceAlias{0.13.7.5.1.fr}{0.13.7.5.1-fr}
\EGAReferenceAlias{0.13.7.5.1.es}{0.13.7.5.1-fr}
\EGAReferenceAlias{0.13.7.5.1.eq}{0.13.7.5.1-fr}
\EGAReferenceAlias{eq:0.13.7.5.1}{0.13.7.5.1-fr}
\EGAReferenceAlias{0.13.7.5.1b}{0.13.7.5.1b-fr}
\EGAReferenceAlias{0.13.7.5.1b-es}{0.13.7.5.1b-fr}
\EGAReferenceAlias{0.13.7.5.1b.fr}{0.13.7.5.1b-fr}
\EGAReferenceAlias{0.13.7.5.1b.es}{0.13.7.5.1b-fr}
\EGAReferenceAlias{0.13.7.5.1b.eq}{0.13.7.5.1b-fr}
\EGAReferenceAlias{eq:0.13.7.5.1b}{0.13.7.5.1b-fr}
\EGAReferenceAlias{0.13.7.6-es}{0.13.7.6}
\EGAReferenceAlias{0.13.7.6.fr}{0.13.7.6}
\EGAReferenceAlias{0.13.7.6.es}{0.13.7.6}
\EGAReferenceAlias{0.13.7.6.eq}{0.13.7.6}
\EGAReferenceAlias{eq:0.13.7.6}{0.13.7.6}
\EGAReferenceAlias{0.13.7.7-es}{0.13.7.7}
\EGAReferenceAlias{0.13.7.7.fr}{0.13.7.7}
\EGAReferenceAlias{0.13.7.7.es}{0.13.7.7}
\EGAReferenceAlias{0.13.7.7.eq}{0.13.7.7}
\EGAReferenceAlias{eq:0.13.7.7}{0.13.7.7}
\EGAReferenceAlias{0.13.7.7.1}{0.13.7.7.1-fr}
\EGAReferenceAlias{0.13.7.7.1-es}{0.13.7.7.1-fr}
\EGAReferenceAlias{0.13.7.7.1.fr}{0.13.7.7.1-fr}
\EGAReferenceAlias{0.13.7.7.1.es}{0.13.7.7.1-fr}
\EGAReferenceAlias{0.13.7.7.1.eq}{0.13.7.7.1-fr}
\EGAReferenceAlias{eq:0.13.7.7.1}{0.13.7.7.1-fr}
\EGAReferenceAlias{0.13.7.7.2-es}{0.13.7.7.2}
\EGAReferenceAlias{0.13.7.7.2.fr}{0.13.7.7.2}
\EGAReferenceAlias{0.13.7.7.2.es}{0.13.7.7.2}
\EGAReferenceAlias{0.13.7.7.2.eq}{0.13.7.7.2}
\EGAReferenceAlias{eq:0.13.7.7.2}{0.13.7.7.2}
\EGAReferenceAlias{0.13.7.7.3}{0.13.7.7.3-fr}
\EGAReferenceAlias{0.13.7.7.3-es}{0.13.7.7.3-fr}
\EGAReferenceAlias{0.13.7.7.3.fr}{0.13.7.7.3-fr}
\EGAReferenceAlias{0.13.7.7.3.es}{0.13.7.7.3-fr}
\EGAReferenceAlias{0.13.7.7.3.eq}{0.13.7.7.3-fr}
\EGAReferenceAlias{eq:0.13.7.7.3}{0.13.7.7.3-fr}
\EGAReferenceAlias{0.13.7.7.4-es}{0.13.7.7.4}
\EGAReferenceAlias{0.13.7.7.4.fr}{0.13.7.7.4}
\EGAReferenceAlias{0.13.7.7.4.es}{0.13.7.7.4}
\EGAReferenceAlias{0.13.7.7.4.eq}{0.13.7.7.4}
\EGAReferenceAlias{eq:0.13.7.7.4}{0.13.7.7.4}
\EGAReferenceAlias{0.13.7.8-es}{0.13.7.8}
\EGAReferenceAlias{0.13.7.8.fr}{0.13.7.8}
\EGAReferenceAlias{0.13.7.8.es}{0.13.7.8}
\EGAReferenceAlias{0.13.7.8.eq}{0.13.7.8}
\EGAReferenceAlias{eq:0.13.7.8}{0.13.7.8}
\EGAReferenceAlias{0.13.7.8.1}{0.13.7.8.1-fr}
\EGAReferenceAlias{0.13.7.8.1-es}{0.13.7.8.1-fr}
\EGAReferenceAlias{0.13.7.8.1.fr}{0.13.7.8.1-fr}
\EGAReferenceAlias{0.13.7.8.1.es}{0.13.7.8.1-fr}
\EGAReferenceAlias{0.13.7.8.1.eq}{0.13.7.8.1-fr}
\EGAReferenceAlias{eq:0.13.7.8.1}{0.13.7.8.1-fr}
\EGAReferenceAlias{III}{chapter:III-fr}
\EGAReferenceAlias{III-fr}{chapter:III-fr}
\EGAReferenceAlias{III-es}{chapter:III-fr}
\EGAReferenceAlias{III.fr}{chapter:III-fr}
\EGAReferenceAlias{III.es}{chapter:III-fr}
\EGAReferenceAlias{III.eq}{chapter:III-fr}
\EGAReferenceAlias{eq:III}{chapter:III-fr}
\EGAReferenceAlias{chapter:III}{chapter:III-fr}
\EGAReferenceAlias{chapter:III-es}{chapter:III-fr}
\EGAReferenceAlias{chapter:III.fr}{chapter:III-fr}
\EGAReferenceAlias{chapter:III.es}{chapter:III-fr}
\EGAReferenceAlias{chapter:III.eq}{chapter:III-fr}
\EGAReferenceAlias{eq:chapter:III}{chapter:III-fr}
\EGAReferenceAlias{III.1}{section:III.1}
\EGAReferenceAlias{III.1-fr}{section:III.1}
\EGAReferenceAlias{III.1-es}{section:III.1}
\EGAReferenceAlias{III.1.fr}{section:III.1}
\EGAReferenceAlias{III.1.es}{section:III.1}
\EGAReferenceAlias{III.1.eq}{section:III.1}
\EGAReferenceAlias{eq:III.1}{section:III.1}
\EGAReferenceAlias{section:III.1-es}{section:III.1}
\EGAReferenceAlias{section:III.1.fr}{section:III.1}
\EGAReferenceAlias{section:III.1.es}{section:III.1}
\EGAReferenceAlias{section:III.1.eq}{section:III.1}
\EGAReferenceAlias{eq:section:III.1}{section:III.1}
\EGAReferenceAlias{III.1.1}{subsection:III.1.1}
\EGAReferenceAlias{III.1.1-fr}{subsection:III.1.1}
\EGAReferenceAlias{III.1.1-es}{subsection:III.1.1}
\EGAReferenceAlias{III.1.1.fr}{subsection:III.1.1}
\EGAReferenceAlias{III.1.1.es}{subsection:III.1.1}
\EGAReferenceAlias{III.1.1.eq}{subsection:III.1.1}
\EGAReferenceAlias{eq:III.1.1}{subsection:III.1.1}
\EGAReferenceAlias{subsection:III.1.1-es}{subsection:III.1.1}
\EGAReferenceAlias{subsection:III.1.1.fr}{subsection:III.1.1}
\EGAReferenceAlias{subsection:III.1.1.es}{subsection:III.1.1}
\EGAReferenceAlias{subsection:III.1.1.eq}{subsection:III.1.1}
\EGAReferenceAlias{eq:subsection:III.1.1}{subsection:III.1.1}
\EGAReferenceAlias{III.1.1.1-es}{III.1.1.1}
\EGAReferenceAlias{III.1.1.1.fr}{III.1.1.1}
\EGAReferenceAlias{III.1.1.1.es}{III.1.1.1}
\EGAReferenceAlias{III.1.1.1.eq}{III.1.1.1}
\EGAReferenceAlias{eq:III.1.1.1}{III.1.1.1}
\EGAReferenceAlias{III.1.1.2-es}{III.1.1.2}
\EGAReferenceAlias{III.1.1.2.fr}{III.1.1.2}
\EGAReferenceAlias{III.1.1.2.es}{III.1.1.2}
\EGAReferenceAlias{III.1.1.2.eq}{III.1.1.2}
\EGAReferenceAlias{eq:III.1.1.2}{III.1.1.2}
\EGAReferenceAlias{III.1.1.2.1}{III.1.1.2.1-fr}
\EGAReferenceAlias{III.1.1.2.1-es}{III.1.1.2.1-fr}
\EGAReferenceAlias{III.1.1.2.1.fr}{III.1.1.2.1-fr}
\EGAReferenceAlias{III.1.1.2.1.es}{III.1.1.2.1-fr}
\EGAReferenceAlias{III.1.1.2.1.eq}{III.1.1.2.1-fr}
\EGAReferenceAlias{eq:III.1.1.2.1}{III.1.1.2.1-fr}
\EGAReferenceAlias{III.1.1.2.2}{III.1.1.2.2-fr}
\EGAReferenceAlias{III.1.1.2.2-es}{III.1.1.2.2-fr}
\EGAReferenceAlias{III.1.1.2.2.fr}{III.1.1.2.2-fr}
\EGAReferenceAlias{III.1.1.2.2.es}{III.1.1.2.2-fr}
\EGAReferenceAlias{III.1.1.2.2.eq}{III.1.1.2.2-fr}
\EGAReferenceAlias{eq:III.1.1.2.2}{III.1.1.2.2-fr}
\EGAReferenceAlias{III.1.1.2.3}{III.1.1.2.3-fr}
\EGAReferenceAlias{III.1.1.2.3-es}{III.1.1.2.3-fr}
\EGAReferenceAlias{III.1.1.2.3.fr}{III.1.1.2.3-fr}
\EGAReferenceAlias{III.1.1.2.3.es}{III.1.1.2.3-fr}
\EGAReferenceAlias{III.1.1.2.3.eq}{III.1.1.2.3-fr}
\EGAReferenceAlias{eq:III.1.1.2.3}{III.1.1.2.3-fr}
\EGAReferenceAlias{III.1.1.3}{III.1.1.3-fr}
\EGAReferenceAlias{III.1.1.3-es}{III.1.1.3-fr}
\EGAReferenceAlias{III.1.1.3.fr}{III.1.1.3-fr}
\EGAReferenceAlias{III.1.1.3.es}{III.1.1.3-fr}
\EGAReferenceAlias{III.1.1.3.eq}{III.1.1.3-fr}
\EGAReferenceAlias{eq:III.1.1.3}{III.1.1.3-fr}
\EGAReferenceAlias{III.1.1.3.1}{III.1.1.3.1-fr}
\EGAReferenceAlias{III.1.1.3.1-es}{III.1.1.3.1-fr}
\EGAReferenceAlias{III.1.1.3.1.fr}{III.1.1.3.1-fr}
\EGAReferenceAlias{III.1.1.3.1.es}{III.1.1.3.1-fr}
\EGAReferenceAlias{III.1.1.3.1.eq}{III.1.1.3.1-fr}
\EGAReferenceAlias{eq:III.1.1.3.1}{III.1.1.3.1-fr}
\EGAReferenceAlias{III.1.1.3.2}{III.1.1.3.2-fr}
\EGAReferenceAlias{III.1.1.3.2-es}{III.1.1.3.2-fr}
\EGAReferenceAlias{III.1.1.3.2.fr}{III.1.1.3.2-fr}
\EGAReferenceAlias{III.1.1.3.2.es}{III.1.1.3.2-fr}
\EGAReferenceAlias{III.1.1.3.2.eq}{III.1.1.3.2-fr}
\EGAReferenceAlias{eq:III.1.1.3.2}{III.1.1.3.2-fr}
\EGAReferenceAlias{III.1.1.3.3}{III.1.1.3.3-fr}
\EGAReferenceAlias{III.1.1.3.3-es}{III.1.1.3.3-fr}
\EGAReferenceAlias{III.1.1.3.3.fr}{III.1.1.3.3-fr}
\EGAReferenceAlias{III.1.1.3.3.es}{III.1.1.3.3-fr}
\EGAReferenceAlias{III.1.1.3.3.eq}{III.1.1.3.3-fr}
\EGAReferenceAlias{eq:III.1.1.3.3}{III.1.1.3.3-fr}
\EGAReferenceAlias{III.1.1.3.4}{III.1.1.3.4-fr}
\EGAReferenceAlias{III.1.1.3.4-es}{III.1.1.3.4-fr}
\EGAReferenceAlias{III.1.1.3.4.fr}{III.1.1.3.4-fr}
\EGAReferenceAlias{III.1.1.3.4.es}{III.1.1.3.4-fr}
\EGAReferenceAlias{III.1.1.3.4.eq}{III.1.1.3.4-fr}
\EGAReferenceAlias{eq:III.1.1.3.4}{III.1.1.3.4-fr}
\EGAReferenceAlias{III.1.1.3.5}{III.1.1.3.5-fr}
\EGAReferenceAlias{III.1.1.3.5-es}{III.1.1.3.5-fr}
\EGAReferenceAlias{III.1.1.3.5.fr}{III.1.1.3.5-fr}
\EGAReferenceAlias{III.1.1.3.5.es}{III.1.1.3.5-fr}
\EGAReferenceAlias{III.1.1.3.5.eq}{III.1.1.3.5-fr}
\EGAReferenceAlias{eq:III.1.1.3.5}{III.1.1.3.5-fr}
\EGAReferenceAlias{III.1.1.4-es}{III.1.1.4}
\EGAReferenceAlias{III.1.1.4.fr}{III.1.1.4}
\EGAReferenceAlias{III.1.1.4.es}{III.1.1.4}
\EGAReferenceAlias{III.1.1.4.eq}{III.1.1.4}
\EGAReferenceAlias{eq:III.1.1.4}{III.1.1.4}
\EGAReferenceAlias{III.1.1.4.1-es}{III.1.1.4.1}
\EGAReferenceAlias{III.1.1.4.1.fr}{III.1.1.4.1}
\EGAReferenceAlias{III.1.1.4.1.es}{III.1.1.4.1}
\EGAReferenceAlias{III.1.1.4.1.eq}{III.1.1.4.1}
\EGAReferenceAlias{eq:III.1.1.4.1}{III.1.1.4.1}
\EGAReferenceAlias{III.1.1.5-es}{III.1.1.5}
\EGAReferenceAlias{III.1.1.5.fr}{III.1.1.5}
\EGAReferenceAlias{III.1.1.5.es}{III.1.1.5}
\EGAReferenceAlias{III.1.1.5.eq}{III.1.1.5}
\EGAReferenceAlias{eq:III.1.1.5}{III.1.1.5}
\EGAReferenceAlias{III.1.1.5.1}{III.1.1.5.1-fr}
\EGAReferenceAlias{III.1.1.5.1-es}{III.1.1.5.1-fr}
\EGAReferenceAlias{III.1.1.5.1.fr}{III.1.1.5.1-fr}
\EGAReferenceAlias{III.1.1.5.1.es}{III.1.1.5.1-fr}
\EGAReferenceAlias{III.1.1.5.1.eq}{III.1.1.5.1-fr}
\EGAReferenceAlias{eq:III.1.1.5.1}{III.1.1.5.1-fr}
\EGAReferenceAlias{III.1.1.5.2}{III.1.1.5.2-fr}
\EGAReferenceAlias{III.1.1.5.2-es}{III.1.1.5.2-fr}
\EGAReferenceAlias{III.1.1.5.2.fr}{III.1.1.5.2-fr}
\EGAReferenceAlias{III.1.1.5.2.es}{III.1.1.5.2-fr}
\EGAReferenceAlias{III.1.1.5.2.eq}{III.1.1.5.2-fr}
\EGAReferenceAlias{eq:III.1.1.5.2}{III.1.1.5.2-fr}
\EGAReferenceAlias{III.1.1.6}{III.1.1.6-fr}
\EGAReferenceAlias{III.1.1.6-es}{III.1.1.6-fr}
\EGAReferenceAlias{III.1.1.6.fr}{III.1.1.6-fr}
\EGAReferenceAlias{III.1.1.6.es}{III.1.1.6-fr}
\EGAReferenceAlias{III.1.1.6.eq}{III.1.1.6-fr}
\EGAReferenceAlias{eq:III.1.1.6}{III.1.1.6-fr}
\EGAReferenceAlias{III.1.1.6.1}{III.1.1.6.1-fr}
\EGAReferenceAlias{III.1.1.6.1-es}{III.1.1.6.1-fr}
\EGAReferenceAlias{III.1.1.6.1.fr}{III.1.1.6.1-fr}
\EGAReferenceAlias{III.1.1.6.1.es}{III.1.1.6.1-fr}
\EGAReferenceAlias{III.1.1.6.1.eq}{III.1.1.6.1-fr}
\EGAReferenceAlias{eq:III.1.1.6.1}{III.1.1.6.1-fr}
\EGAReferenceAlias{III.1.1.6.2}{III.1.1.6.2-fr}
\EGAReferenceAlias{III.1.1.6.2-es}{III.1.1.6.2-fr}
\EGAReferenceAlias{III.1.1.6.2.fr}{III.1.1.6.2-fr}
\EGAReferenceAlias{III.1.1.6.2.es}{III.1.1.6.2-fr}
\EGAReferenceAlias{III.1.1.6.2.eq}{III.1.1.6.2-fr}
\EGAReferenceAlias{eq:III.1.1.6.2}{III.1.1.6.2-fr}
\EGAReferenceAlias{III.1.1.6.3}{III.1.1.6.3-fr}
\EGAReferenceAlias{III.1.1.6.3-es}{III.1.1.6.3-fr}
\EGAReferenceAlias{III.1.1.6.3.fr}{III.1.1.6.3-fr}
\EGAReferenceAlias{III.1.1.6.3.es}{III.1.1.6.3-fr}
\EGAReferenceAlias{III.1.1.6.3.eq}{III.1.1.6.3-fr}
\EGAReferenceAlias{eq:III.1.1.6.3}{III.1.1.6.3-fr}
\EGAReferenceAlias{III.1.1.6.4}{III.1.1.6.4-fr}
\EGAReferenceAlias{III.1.1.6.4-es}{III.1.1.6.4-fr}
\EGAReferenceAlias{III.1.1.6.4.fr}{III.1.1.6.4-fr}
\EGAReferenceAlias{III.1.1.6.4.es}{III.1.1.6.4-fr}
\EGAReferenceAlias{III.1.1.6.4.eq}{III.1.1.6.4-fr}
\EGAReferenceAlias{eq:III.1.1.6.4}{III.1.1.6.4-fr}
\EGAReferenceAlias{III.1.1.6.5}{III.1.1.6.5-fr}
\EGAReferenceAlias{III.1.1.6.5-es}{III.1.1.6.5-fr}
\EGAReferenceAlias{III.1.1.6.5.fr}{III.1.1.6.5-fr}
\EGAReferenceAlias{III.1.1.6.5.es}{III.1.1.6.5-fr}
\EGAReferenceAlias{III.1.1.6.5.eq}{III.1.1.6.5-fr}
\EGAReferenceAlias{eq:III.1.1.6.5}{III.1.1.6.5-fr}
\EGAReferenceAlias{III.1.1.7-es}{III.1.1.7}
\EGAReferenceAlias{III.1.1.7.fr}{III.1.1.7}
\EGAReferenceAlias{III.1.1.7.es}{III.1.1.7}
\EGAReferenceAlias{III.1.1.7.eq}{III.1.1.7}
\EGAReferenceAlias{eq:III.1.1.7}{III.1.1.7}
\EGAReferenceAlias{III.1.1.7.1}{III.1.1.7.1-fr}
\EGAReferenceAlias{III.1.1.7.1-es}{III.1.1.7.1-fr}
\EGAReferenceAlias{III.1.1.7.1.fr}{III.1.1.7.1-fr}
\EGAReferenceAlias{III.1.1.7.1.es}{III.1.1.7.1-fr}
\EGAReferenceAlias{III.1.1.7.1.eq}{III.1.1.7.1-fr}
\EGAReferenceAlias{eq:III.1.1.7.1}{III.1.1.7.1-fr}
\EGAReferenceAlias{III.1.1.7.2}{III.1.1.7.2-fr}
\EGAReferenceAlias{III.1.1.7.2-es}{III.1.1.7.2-fr}
\EGAReferenceAlias{III.1.1.7.2.fr}{III.1.1.7.2-fr}
\EGAReferenceAlias{III.1.1.7.2.es}{III.1.1.7.2-fr}
\EGAReferenceAlias{III.1.1.7.2.eq}{III.1.1.7.2-fr}
\EGAReferenceAlias{eq:III.1.1.7.2}{III.1.1.7.2-fr}
\EGAReferenceAlias{III.1.1.8-es}{III.1.1.8}
\EGAReferenceAlias{III.1.1.8.fr}{III.1.1.8}
\EGAReferenceAlias{III.1.1.8.es}{III.1.1.8}
\EGAReferenceAlias{III.1.1.8.eq}{III.1.1.8}
\EGAReferenceAlias{eq:III.1.1.8}{III.1.1.8}
\EGAReferenceAlias{III.1.1.9-es}{III.1.1.9}
\EGAReferenceAlias{III.1.1.9.fr}{III.1.1.9}
\EGAReferenceAlias{III.1.1.9.es}{III.1.1.9}
\EGAReferenceAlias{III.1.1.9.eq}{III.1.1.9}
\EGAReferenceAlias{eq:III.1.1.9}{III.1.1.9}
\EGAReferenceAlias{III.1.1.10-es}{III.1.1.10}
\EGAReferenceAlias{III.1.1.10.fr}{III.1.1.10}
\EGAReferenceAlias{III.1.1.10.es}{III.1.1.10}
\EGAReferenceAlias{III.1.1.10.eq}{III.1.1.10}
\EGAReferenceAlias{eq:III.1.1.10}{III.1.1.10}
\EGAReferenceAlias{III.1.2}{subsection:III.1.2}
\EGAReferenceAlias{III.1.2-fr}{subsection:III.1.2}
\EGAReferenceAlias{III.1.2-es}{subsection:III.1.2}
\EGAReferenceAlias{III.1.2.fr}{subsection:III.1.2}
\EGAReferenceAlias{III.1.2.es}{subsection:III.1.2}
\EGAReferenceAlias{III.1.2.eq}{subsection:III.1.2}
\EGAReferenceAlias{eq:III.1.2}{subsection:III.1.2}
\EGAReferenceAlias{subsection:III.1.2-es}{subsection:III.1.2}
\EGAReferenceAlias{subsection:III.1.2.fr}{subsection:III.1.2}
\EGAReferenceAlias{subsection:III.1.2.es}{subsection:III.1.2}
\EGAReferenceAlias{subsection:III.1.2.eq}{subsection:III.1.2}
\EGAReferenceAlias{eq:subsection:III.1.2}{subsection:III.1.2}
\EGAReferenceAlias{III.1.2.1-es}{III.1.2.1}
\EGAReferenceAlias{III.1.2.1.fr}{III.1.2.1}
\EGAReferenceAlias{III.1.2.1.es}{III.1.2.1}
\EGAReferenceAlias{III.1.2.1.eq}{III.1.2.1}
\EGAReferenceAlias{eq:III.1.2.1}{III.1.2.1}
\EGAReferenceAlias{III.1.2.2-es}{III.1.2.2}
\EGAReferenceAlias{III.1.2.2.fr}{III.1.2.2}
\EGAReferenceAlias{III.1.2.2.es}{III.1.2.2}
\EGAReferenceAlias{III.1.2.2.eq}{III.1.2.2}
\EGAReferenceAlias{eq:III.1.2.2}{III.1.2.2}
\EGAReferenceAlias{III.1.2.2.1}{III.1.2.2.1-fr}
\EGAReferenceAlias{III.1.2.2.1-es}{III.1.2.2.1-fr}
\EGAReferenceAlias{III.1.2.2.1.fr}{III.1.2.2.1-fr}
\EGAReferenceAlias{III.1.2.2.1.es}{III.1.2.2.1-fr}
\EGAReferenceAlias{III.1.2.2.1.eq}{III.1.2.2.1-fr}
\EGAReferenceAlias{eq:III.1.2.2.1}{III.1.2.2.1-fr}
\EGAReferenceAlias{III.1.2.2.2}{III.1.2.2.2-fr}
\EGAReferenceAlias{III.1.2.2.2-es}{III.1.2.2.2-fr}
\EGAReferenceAlias{III.1.2.2.2.fr}{III.1.2.2.2-fr}
\EGAReferenceAlias{III.1.2.2.2.es}{III.1.2.2.2-fr}
\EGAReferenceAlias{III.1.2.2.2.eq}{III.1.2.2.2-fr}
\EGAReferenceAlias{eq:III.1.2.2.2}{III.1.2.2.2-fr}
\EGAReferenceAlias{III.1.2.2.3}{III.1.2.2.3-fr}
\EGAReferenceAlias{III.1.2.2.3-es}{III.1.2.2.3-fr}
\EGAReferenceAlias{III.1.2.2.3.fr}{III.1.2.2.3-fr}
\EGAReferenceAlias{III.1.2.2.3.es}{III.1.2.2.3-fr}
\EGAReferenceAlias{III.1.2.2.3.eq}{III.1.2.2.3-fr}
\EGAReferenceAlias{eq:III.1.2.2.3}{III.1.2.2.3-fr}
\EGAReferenceAlias{III.1.2.3-es}{III.1.2.3}
\EGAReferenceAlias{III.1.2.3.fr}{III.1.2.3}
\EGAReferenceAlias{III.1.2.3.es}{III.1.2.3}
\EGAReferenceAlias{III.1.2.3.eq}{III.1.2.3}
\EGAReferenceAlias{eq:III.1.2.3}{III.1.2.3}
\EGAReferenceAlias{III.1.2.3.1}{III.1.2.3.1-fr}
\EGAReferenceAlias{III.1.2.3.1-es}{III.1.2.3.1-fr}
\EGAReferenceAlias{III.1.2.3.1.fr}{III.1.2.3.1-fr}
\EGAReferenceAlias{III.1.2.3.1.es}{III.1.2.3.1-fr}
\EGAReferenceAlias{III.1.2.3.1.eq}{III.1.2.3.1-fr}
\EGAReferenceAlias{eq:III.1.2.3.1}{III.1.2.3.1-fr}
\EGAReferenceAlias{III.1.2.3.2}{III.1.2.3.2-fr}
\EGAReferenceAlias{III.1.2.3.2-es}{III.1.2.3.2-fr}
\EGAReferenceAlias{III.1.2.3.2.fr}{III.1.2.3.2-fr}
\EGAReferenceAlias{III.1.2.3.2.es}{III.1.2.3.2-fr}
\EGAReferenceAlias{III.1.2.3.2.eq}{III.1.2.3.2-fr}
\EGAReferenceAlias{eq:III.1.2.3.2}{III.1.2.3.2-fr}
\EGAReferenceAlias{III.1.2.4-es}{III.1.2.4}
\EGAReferenceAlias{III.1.2.4.fr}{III.1.2.4}
\EGAReferenceAlias{III.1.2.4.es}{III.1.2.4}
\EGAReferenceAlias{III.1.2.4.eq}{III.1.2.4}
\EGAReferenceAlias{eq:III.1.2.4}{III.1.2.4}
\EGAReferenceAlias{III.1.2.5-es}{III.1.2.5}
\EGAReferenceAlias{III.1.2.5.fr}{III.1.2.5}
\EGAReferenceAlias{III.1.2.5.es}{III.1.2.5}
\EGAReferenceAlias{III.1.2.5.eq}{III.1.2.5}
\EGAReferenceAlias{eq:III.1.2.5}{III.1.2.5}
\EGAReferenceAlias{III.1.2.5.1}{III.1.2.5.1-fr}
\EGAReferenceAlias{III.1.2.5.1-es}{III.1.2.5.1-fr}
\EGAReferenceAlias{III.1.2.5.1.fr}{III.1.2.5.1-fr}
\EGAReferenceAlias{III.1.2.5.1.es}{III.1.2.5.1-fr}
\EGAReferenceAlias{III.1.2.5.1.eq}{III.1.2.5.1-fr}
\EGAReferenceAlias{eq:III.1.2.5.1}{III.1.2.5.1-fr}
\EGAReferenceAlias{III.1.3}{subsection:III.1.3}
\EGAReferenceAlias{III.1.3-fr}{subsection:III.1.3}
\EGAReferenceAlias{III.1.3-es}{subsection:III.1.3}
\EGAReferenceAlias{III.1.3.fr}{subsection:III.1.3}
\EGAReferenceAlias{III.1.3.es}{subsection:III.1.3}
\EGAReferenceAlias{III.1.3.eq}{subsection:III.1.3}
\EGAReferenceAlias{eq:III.1.3}{subsection:III.1.3}
\EGAReferenceAlias{subsection:III.1.3-es}{subsection:III.1.3}
\EGAReferenceAlias{subsection:III.1.3.fr}{subsection:III.1.3}
\EGAReferenceAlias{subsection:III.1.3.es}{subsection:III.1.3}
\EGAReferenceAlias{subsection:III.1.3.eq}{subsection:III.1.3}
\EGAReferenceAlias{eq:subsection:III.1.3}{subsection:III.1.3}
\EGAReferenceAlias{III.1.3.1-es}{III.1.3.1}
\EGAReferenceAlias{III.1.3.1.fr}{III.1.3.1}
\EGAReferenceAlias{III.1.3.1.es}{III.1.3.1}
\EGAReferenceAlias{III.1.3.1.eq}{III.1.3.1}
\EGAReferenceAlias{eq:III.1.3.1}{III.1.3.1}
\EGAReferenceAlias{III.1.3.2-es}{III.1.3.2}
\EGAReferenceAlias{III.1.3.2.fr}{III.1.3.2}
\EGAReferenceAlias{III.1.3.2.es}{III.1.3.2}
\EGAReferenceAlias{III.1.3.2.eq}{III.1.3.2}
\EGAReferenceAlias{eq:III.1.3.2}{III.1.3.2}
\EGAReferenceAlias{III.1.3.3-es}{III.1.3.3}
\EGAReferenceAlias{III.1.3.3.fr}{III.1.3.3}
\EGAReferenceAlias{III.1.3.3.es}{III.1.3.3}
\EGAReferenceAlias{III.1.3.3.eq}{III.1.3.3}
\EGAReferenceAlias{eq:III.1.3.3}{III.1.3.3}
\EGAReferenceAlias{III.1.3.4-es}{III.1.3.4}
\EGAReferenceAlias{III.1.3.4.fr}{III.1.3.4}
\EGAReferenceAlias{III.1.3.4.es}{III.1.3.4}
\EGAReferenceAlias{III.1.3.4.eq}{III.1.3.4}
\EGAReferenceAlias{eq:III.1.3.4}{III.1.3.4}
\EGAReferenceAlias{III.1.4}{subsection:III.1.4}
\EGAReferenceAlias{III.1.4-fr}{subsection:III.1.4}
\EGAReferenceAlias{III.1.4-es}{subsection:III.1.4}
\EGAReferenceAlias{III.1.4.fr}{subsection:III.1.4}
\EGAReferenceAlias{III.1.4.es}{subsection:III.1.4}
\EGAReferenceAlias{III.1.4.eq}{subsection:III.1.4}
\EGAReferenceAlias{eq:III.1.4}{subsection:III.1.4}
\EGAReferenceAlias{subsection:III.1.4-es}{subsection:III.1.4}
\EGAReferenceAlias{subsection:III.1.4.fr}{subsection:III.1.4}
\EGAReferenceAlias{subsection:III.1.4.es}{subsection:III.1.4}
\EGAReferenceAlias{subsection:III.1.4.eq}{subsection:III.1.4}
\EGAReferenceAlias{eq:subsection:III.1.4}{subsection:III.1.4}
\EGAReferenceAlias{III.1.4.1-es}{III.1.4.1}
\EGAReferenceAlias{III.1.4.1.fr}{III.1.4.1}
\EGAReferenceAlias{III.1.4.1.es}{III.1.4.1}
\EGAReferenceAlias{III.1.4.1.eq}{III.1.4.1}
\EGAReferenceAlias{eq:III.1.4.1}{III.1.4.1}
\EGAReferenceAlias{III.1.4.2-es}{III.1.4.2}
\EGAReferenceAlias{III.1.4.2.fr}{III.1.4.2}
\EGAReferenceAlias{III.1.4.2.es}{III.1.4.2}
\EGAReferenceAlias{III.1.4.2.eq}{III.1.4.2}
\EGAReferenceAlias{eq:III.1.4.2}{III.1.4.2}
\EGAReferenceAlias{III.1.4.3-es}{III.1.4.3}
\EGAReferenceAlias{III.1.4.3.fr}{III.1.4.3}
\EGAReferenceAlias{III.1.4.3.es}{III.1.4.3}
\EGAReferenceAlias{III.1.4.3.eq}{III.1.4.3}
\EGAReferenceAlias{eq:III.1.4.3}{III.1.4.3}
\EGAReferenceAlias{III.1.4.3.1}{III.1.4.3.1-fr}
\EGAReferenceAlias{III.1.4.3.1-es}{III.1.4.3.1-fr}
\EGAReferenceAlias{III.1.4.3.1.fr}{III.1.4.3.1-fr}
\EGAReferenceAlias{III.1.4.3.1.es}{III.1.4.3.1-fr}
\EGAReferenceAlias{III.1.4.3.1.eq}{III.1.4.3.1-fr}
\EGAReferenceAlias{eq:III.1.4.3.1}{III.1.4.3.1-fr}
\EGAReferenceAlias{III.1.4.3.2}{III.1.4.3.2-fr}
\EGAReferenceAlias{III.1.4.3.2-es}{III.1.4.3.2-fr}
\EGAReferenceAlias{III.1.4.3.2.fr}{III.1.4.3.2-fr}
\EGAReferenceAlias{III.1.4.3.2.es}{III.1.4.3.2-fr}
\EGAReferenceAlias{III.1.4.3.2.eq}{III.1.4.3.2-fr}
\EGAReferenceAlias{eq:III.1.4.3.2}{III.1.4.3.2-fr}
\EGAReferenceAlias{III.1.4.4-es}{III.1.4.4}
\EGAReferenceAlias{III.1.4.4.fr}{III.1.4.4}
\EGAReferenceAlias{III.1.4.4.es}{III.1.4.4}
\EGAReferenceAlias{III.1.4.4.eq}{III.1.4.4}
\EGAReferenceAlias{eq:III.1.4.4}{III.1.4.4}
\EGAReferenceAlias{III.1.4.4.1}{III.1.4.4.1-fr}
\EGAReferenceAlias{III.1.4.4.1-es}{III.1.4.4.1-fr}
\EGAReferenceAlias{III.1.4.4.1.fr}{III.1.4.4.1-fr}
\EGAReferenceAlias{III.1.4.4.1.es}{III.1.4.4.1-fr}
\EGAReferenceAlias{III.1.4.4.1.eq}{III.1.4.4.1-fr}
\EGAReferenceAlias{eq:III.1.4.4.1}{III.1.4.4.1-fr}
\EGAReferenceAlias{III.1.4.5-es}{III.1.4.5}
\EGAReferenceAlias{III.1.4.5.fr}{III.1.4.5}
\EGAReferenceAlias{III.1.4.5.es}{III.1.4.5}
\EGAReferenceAlias{III.1.4.5.eq}{III.1.4.5}
\EGAReferenceAlias{eq:III.1.4.5}{III.1.4.5}
\EGAReferenceAlias{III.1.4.5.1}{III.1.4.5.1-fr}
\EGAReferenceAlias{III.1.4.5.1-es}{III.1.4.5.1-fr}
\EGAReferenceAlias{III.1.4.5.1.fr}{III.1.4.5.1-fr}
\EGAReferenceAlias{III.1.4.5.1.es}{III.1.4.5.1-fr}
\EGAReferenceAlias{III.1.4.5.1.eq}{III.1.4.5.1-fr}
\EGAReferenceAlias{eq:III.1.4.5.1}{III.1.4.5.1-fr}
\EGAReferenceAlias{III.1.4.6-es}{III.1.4.6}
\EGAReferenceAlias{III.1.4.6.fr}{III.1.4.6}
\EGAReferenceAlias{III.1.4.6.es}{III.1.4.6}
\EGAReferenceAlias{III.1.4.6.eq}{III.1.4.6}
\EGAReferenceAlias{eq:III.1.4.6}{III.1.4.6}
\EGAReferenceAlias{III.1.4.7-es}{III.1.4.7}
\EGAReferenceAlias{III.1.4.7.fr}{III.1.4.7}
\EGAReferenceAlias{III.1.4.7.es}{III.1.4.7}
\EGAReferenceAlias{III.1.4.7.eq}{III.1.4.7}
\EGAReferenceAlias{eq:III.1.4.7}{III.1.4.7}
\EGAReferenceAlias{III.1.4.8-es}{III.1.4.8}
\EGAReferenceAlias{III.1.4.8.fr}{III.1.4.8}
\EGAReferenceAlias{III.1.4.8.es}{III.1.4.8}
\EGAReferenceAlias{III.1.4.8.eq}{III.1.4.8}
\EGAReferenceAlias{eq:III.1.4.8}{III.1.4.8}
\EGAReferenceAlias{III.1.4.9-es}{III.1.4.9}
\EGAReferenceAlias{III.1.4.9.fr}{III.1.4.9}
\EGAReferenceAlias{III.1.4.9.es}{III.1.4.9}
\EGAReferenceAlias{III.1.4.9.eq}{III.1.4.9}
\EGAReferenceAlias{eq:III.1.4.9}{III.1.4.9}
\EGAReferenceAlias{III.1.4.10-es}{III.1.4.10}
\EGAReferenceAlias{III.1.4.10.fr}{III.1.4.10}
\EGAReferenceAlias{III.1.4.10.es}{III.1.4.10}
\EGAReferenceAlias{III.1.4.10.eq}{III.1.4.10}
\EGAReferenceAlias{eq:III.1.4.10}{III.1.4.10}
\EGAReferenceAlias{III.1.4.10.1}{III.1.4.10.1-fr}
\EGAReferenceAlias{III.1.4.10.1-es}{III.1.4.10.1-fr}
\EGAReferenceAlias{III.1.4.10.1.fr}{III.1.4.10.1-fr}
\EGAReferenceAlias{III.1.4.10.1.es}{III.1.4.10.1-fr}
\EGAReferenceAlias{III.1.4.10.1.eq}{III.1.4.10.1-fr}
\EGAReferenceAlias{eq:III.1.4.10.1}{III.1.4.10.1-fr}
\EGAReferenceAlias{III.1.4.10.2}{III.1.4.10.2-fr}
\EGAReferenceAlias{III.1.4.10.2-es}{III.1.4.10.2-fr}
\EGAReferenceAlias{III.1.4.10.2.fr}{III.1.4.10.2-fr}
\EGAReferenceAlias{III.1.4.10.2.es}{III.1.4.10.2-fr}
\EGAReferenceAlias{III.1.4.10.2.eq}{III.1.4.10.2-fr}
\EGAReferenceAlias{eq:III.1.4.10.2}{III.1.4.10.2-fr}
\EGAReferenceAlias{III.1.4.10.3}{III.1.4.10.3-fr}
\EGAReferenceAlias{III.1.4.10.3-es}{III.1.4.10.3-fr}
\EGAReferenceAlias{III.1.4.10.3.fr}{III.1.4.10.3-fr}
\EGAReferenceAlias{III.1.4.10.3.es}{III.1.4.10.3-fr}
\EGAReferenceAlias{III.1.4.10.3.eq}{III.1.4.10.3-fr}
\EGAReferenceAlias{eq:III.1.4.10.3}{III.1.4.10.3-fr}
\EGAReferenceAlias{III.1.4.11-es}{III.1.4.11}
\EGAReferenceAlias{III.1.4.11.fr}{III.1.4.11}
\EGAReferenceAlias{III.1.4.11.es}{III.1.4.11}
\EGAReferenceAlias{III.1.4.11.eq}{III.1.4.11}
\EGAReferenceAlias{eq:III.1.4.11}{III.1.4.11}
\EGAReferenceAlias{III.1.4.11.1}{III.1.4.11.1-fr}
\EGAReferenceAlias{III.1.4.11.1-es}{III.1.4.11.1-fr}
\EGAReferenceAlias{III.1.4.11.1.fr}{III.1.4.11.1-fr}
\EGAReferenceAlias{III.1.4.11.1.es}{III.1.4.11.1-fr}
\EGAReferenceAlias{III.1.4.11.1.eq}{III.1.4.11.1-fr}
\EGAReferenceAlias{eq:III.1.4.11.1}{III.1.4.11.1-fr}
\EGAReferenceAlias{III.1.4.12}{III.1.4.12-fr}
\EGAReferenceAlias{III.1.4.12-es}{III.1.4.12-fr}
\EGAReferenceAlias{III.1.4.12.fr}{III.1.4.12-fr}
\EGAReferenceAlias{III.1.4.12.es}{III.1.4.12-fr}
\EGAReferenceAlias{III.1.4.12.eq}{III.1.4.12-fr}
\EGAReferenceAlias{eq:III.1.4.12}{III.1.4.12-fr}
\EGAReferenceAlias{III.1.4.13}{III.1.4.13-fr}
\EGAReferenceAlias{III.1.4.13-es}{III.1.4.13-fr}
\EGAReferenceAlias{III.1.4.13.fr}{III.1.4.13-fr}
\EGAReferenceAlias{III.1.4.13.es}{III.1.4.13-fr}
\EGAReferenceAlias{III.1.4.13.eq}{III.1.4.13-fr}
\EGAReferenceAlias{eq:III.1.4.13}{III.1.4.13-fr}
\EGAReferenceAlias{III.1.4.14-es}{III.1.4.14}
\EGAReferenceAlias{III.1.4.14.fr}{III.1.4.14}
\EGAReferenceAlias{III.1.4.14.es}{III.1.4.14}
\EGAReferenceAlias{III.1.4.14.eq}{III.1.4.14}
\EGAReferenceAlias{eq:III.1.4.14}{III.1.4.14}
\EGAReferenceAlias{III.1.4.15-es}{III.1.4.15}
\EGAReferenceAlias{III.1.4.15.fr}{III.1.4.15}
\EGAReferenceAlias{III.1.4.15.es}{III.1.4.15}
\EGAReferenceAlias{III.1.4.15.eq}{III.1.4.15}
\EGAReferenceAlias{eq:III.1.4.15}{III.1.4.15}
\EGAReferenceAlias{III.1.4.15.1}{III.1.4.15.1-fr}
\EGAReferenceAlias{III.1.4.15.1-es}{III.1.4.15.1-fr}
\EGAReferenceAlias{III.1.4.15.1.fr}{III.1.4.15.1-fr}
\EGAReferenceAlias{III.1.4.15.1.es}{III.1.4.15.1-fr}
\EGAReferenceAlias{III.1.4.15.1.eq}{III.1.4.15.1-fr}
\EGAReferenceAlias{eq:III.1.4.15.1}{III.1.4.15.1-fr}
\EGAReferenceAlias{III.1.4.15.2}{III.1.4.15.2-fr}
\EGAReferenceAlias{III.1.4.15.2-es}{III.1.4.15.2-fr}
\EGAReferenceAlias{III.1.4.15.2.fr}{III.1.4.15.2-fr}
\EGAReferenceAlias{III.1.4.15.2.es}{III.1.4.15.2-fr}
\EGAReferenceAlias{III.1.4.15.2.eq}{III.1.4.15.2-fr}
\EGAReferenceAlias{eq:III.1.4.15.2}{III.1.4.15.2-fr}
\EGAReferenceAlias{III.1.4.15.3}{III.1.4.15.3-fr}
\EGAReferenceAlias{III.1.4.15.3-es}{III.1.4.15.3-fr}
\EGAReferenceAlias{III.1.4.15.3.fr}{III.1.4.15.3-fr}
\EGAReferenceAlias{III.1.4.15.3.es}{III.1.4.15.3-fr}
\EGAReferenceAlias{III.1.4.15.3.eq}{III.1.4.15.3-fr}
\EGAReferenceAlias{eq:III.1.4.15.3}{III.1.4.15.3-fr}
\EGAReferenceAlias{III.1.4.15.4}{III.1.4.15.4-fr}
\EGAReferenceAlias{III.1.4.15.4-es}{III.1.4.15.4-fr}
\EGAReferenceAlias{III.1.4.15.4.fr}{III.1.4.15.4-fr}
\EGAReferenceAlias{III.1.4.15.4.es}{III.1.4.15.4-fr}
\EGAReferenceAlias{III.1.4.15.4.eq}{III.1.4.15.4-fr}
\EGAReferenceAlias{eq:III.1.4.15.4}{III.1.4.15.4-fr}
\EGAReferenceAlias{III.1.4.15.5}{III.1.4.15.5-fr}
\EGAReferenceAlias{III.1.4.15.5-es}{III.1.4.15.5-fr}
\EGAReferenceAlias{III.1.4.15.5.fr}{III.1.4.15.5-fr}
\EGAReferenceAlias{III.1.4.15.5.es}{III.1.4.15.5-fr}
\EGAReferenceAlias{III.1.4.15.5.eq}{III.1.4.15.5-fr}
\EGAReferenceAlias{eq:III.1.4.15.5}{III.1.4.15.5-fr}
\EGAReferenceAlias{III.1.4.16-es}{III.1.4.16}
\EGAReferenceAlias{III.1.4.16.fr}{III.1.4.16}
\EGAReferenceAlias{III.1.4.16.es}{III.1.4.16}
\EGAReferenceAlias{III.1.4.16.eq}{III.1.4.16}
\EGAReferenceAlias{eq:III.1.4.16}{III.1.4.16}
\EGAReferenceAlias{III.1.4.16.1}{III.1.4.16.1-fr}
\EGAReferenceAlias{III.1.4.16.1-es}{III.1.4.16.1-fr}
\EGAReferenceAlias{III.1.4.16.1.fr}{III.1.4.16.1-fr}
\EGAReferenceAlias{III.1.4.16.1.es}{III.1.4.16.1-fr}
\EGAReferenceAlias{III.1.4.16.1.eq}{III.1.4.16.1-fr}
\EGAReferenceAlias{eq:III.1.4.16.1}{III.1.4.16.1-fr}
\EGAReferenceAlias{III.1.4.17-es}{III.1.4.17}
\EGAReferenceAlias{III.1.4.17.fr}{III.1.4.17}
\EGAReferenceAlias{III.1.4.17.es}{III.1.4.17}
\EGAReferenceAlias{III.1.4.17.eq}{III.1.4.17}
\EGAReferenceAlias{eq:III.1.4.17}{III.1.4.17}
\EGAReferenceAlias{III.2}{section:III.2}
\EGAReferenceAlias{III.2-fr}{section:III.2}
\EGAReferenceAlias{III.2-es}{section:III.2}
\EGAReferenceAlias{III.2.fr}{section:III.2}
\EGAReferenceAlias{III.2.es}{section:III.2}
\EGAReferenceAlias{III.2.eq}{section:III.2}
\EGAReferenceAlias{eq:III.2}{section:III.2}
\EGAReferenceAlias{section:III.2-es}{section:III.2}
\EGAReferenceAlias{section:III.2.fr}{section:III.2}
\EGAReferenceAlias{section:III.2.es}{section:III.2}
\EGAReferenceAlias{section:III.2.eq}{section:III.2}
\EGAReferenceAlias{eq:section:III.2}{section:III.2}
\EGAReferenceAlias{III.2.1}{subsection:III.2.1}
\EGAReferenceAlias{III.2.1-fr}{subsection:III.2.1}
\EGAReferenceAlias{III.2.1-es}{subsection:III.2.1}
\EGAReferenceAlias{III.2.1.fr}{subsection:III.2.1}
\EGAReferenceAlias{III.2.1.es}{subsection:III.2.1}
\EGAReferenceAlias{III.2.1.eq}{subsection:III.2.1}
\EGAReferenceAlias{eq:III.2.1}{subsection:III.2.1}
\EGAReferenceAlias{subsection:III.2.1-es}{subsection:III.2.1}
\EGAReferenceAlias{subsection:III.2.1.fr}{subsection:III.2.1}
\EGAReferenceAlias{subsection:III.2.1.es}{subsection:III.2.1}
\EGAReferenceAlias{subsection:III.2.1.eq}{subsection:III.2.1}
\EGAReferenceAlias{eq:subsection:III.2.1}{subsection:III.2.1}
\EGAReferenceAlias{III.2.1.1-es}{III.2.1.1}
\EGAReferenceAlias{III.2.1.1.fr}{III.2.1.1}
\EGAReferenceAlias{III.2.1.1.es}{III.2.1.1}
\EGAReferenceAlias{III.2.1.1.eq}{III.2.1.1}
\EGAReferenceAlias{eq:III.2.1.1}{III.2.1.1}
\EGAReferenceAlias{III.2.1.1.1}{III.2.1.1.1-fr}
\EGAReferenceAlias{III.2.1.1.1-es}{III.2.1.1.1-fr}
\EGAReferenceAlias{III.2.1.1.1.fr}{III.2.1.1.1-fr}
\EGAReferenceAlias{III.2.1.1.1.es}{III.2.1.1.1-fr}
\EGAReferenceAlias{III.2.1.1.1.eq}{III.2.1.1.1-fr}
\EGAReferenceAlias{eq:III.2.1.1.1}{III.2.1.1.1-fr}
\EGAReferenceAlias{III.2.1.1.2}{III.2.1.1.2-fr}
\EGAReferenceAlias{III.2.1.1.2-es}{III.2.1.1.2-fr}
\EGAReferenceAlias{III.2.1.1.2.fr}{III.2.1.1.2-fr}
\EGAReferenceAlias{III.2.1.1.2.es}{III.2.1.1.2-fr}
\EGAReferenceAlias{III.2.1.1.2.eq}{III.2.1.1.2-fr}
\EGAReferenceAlias{eq:III.2.1.1.2}{III.2.1.1.2-fr}
\EGAReferenceAlias{III.2.1.1.3}{III.2.1.1.3-fr}
\EGAReferenceAlias{III.2.1.1.3-es}{III.2.1.1.3-fr}
\EGAReferenceAlias{III.2.1.1.3.fr}{III.2.1.1.3-fr}
\EGAReferenceAlias{III.2.1.1.3.es}{III.2.1.1.3-fr}
\EGAReferenceAlias{III.2.1.1.3.eq}{III.2.1.1.3-fr}
\EGAReferenceAlias{eq:III.2.1.1.3}{III.2.1.1.3-fr}
\EGAReferenceAlias{III.2.1.2-es}{III.2.1.2}
\EGAReferenceAlias{III.2.1.2.fr}{III.2.1.2}
\EGAReferenceAlias{III.2.1.2.es}{III.2.1.2}
\EGAReferenceAlias{III.2.1.2.eq}{III.2.1.2}
\EGAReferenceAlias{eq:III.2.1.2}{III.2.1.2}
\EGAReferenceAlias{III.2.1.2.1}{III.2.1.2.1-fr}
\EGAReferenceAlias{III.2.1.2.1-es}{III.2.1.2.1-fr}
\EGAReferenceAlias{III.2.1.2.1.fr}{III.2.1.2.1-fr}
\EGAReferenceAlias{III.2.1.2.1.es}{III.2.1.2.1-fr}
\EGAReferenceAlias{III.2.1.2.1.eq}{III.2.1.2.1-fr}
\EGAReferenceAlias{eq:III.2.1.2.1}{III.2.1.2.1-fr}
\EGAReferenceAlias{III.2.1.2.2}{III.2.1.2.2-fr}
\EGAReferenceAlias{III.2.1.2.2-es}{III.2.1.2.2-fr}
\EGAReferenceAlias{III.2.1.2.2.fr}{III.2.1.2.2-fr}
\EGAReferenceAlias{III.2.1.2.2.es}{III.2.1.2.2-fr}
\EGAReferenceAlias{III.2.1.2.2.eq}{III.2.1.2.2-fr}
\EGAReferenceAlias{eq:III.2.1.2.2}{III.2.1.2.2-fr}
\EGAReferenceAlias{III.2.1.2.3}{III.2.1.2.3-fr}
\EGAReferenceAlias{III.2.1.2.3-es}{III.2.1.2.3-fr}
\EGAReferenceAlias{III.2.1.2.3.fr}{III.2.1.2.3-fr}
\EGAReferenceAlias{III.2.1.2.3.es}{III.2.1.2.3-fr}
\EGAReferenceAlias{III.2.1.2.3.eq}{III.2.1.2.3-fr}
\EGAReferenceAlias{eq:III.2.1.2.3}{III.2.1.2.3-fr}
\EGAReferenceAlias{III.2.1.3-es}{III.2.1.3}
\EGAReferenceAlias{III.2.1.3.fr}{III.2.1.3}
\EGAReferenceAlias{III.2.1.3.es}{III.2.1.3}
\EGAReferenceAlias{III.2.1.3.eq}{III.2.1.3}
\EGAReferenceAlias{eq:III.2.1.3}{III.2.1.3}
\EGAReferenceAlias{III.2.1.3.1}{III.2.1.3.1-fr}
\EGAReferenceAlias{III.2.1.3.1-es}{III.2.1.3.1-fr}
\EGAReferenceAlias{III.2.1.3.1.fr}{III.2.1.3.1-fr}
\EGAReferenceAlias{III.2.1.3.1.es}{III.2.1.3.1-fr}
\EGAReferenceAlias{III.2.1.3.1.eq}{III.2.1.3.1-fr}
\EGAReferenceAlias{eq:III.2.1.3.1}{III.2.1.3.1-fr}
\EGAReferenceAlias{III.2.1.3.2}{III.2.1.3.2-fr}
\EGAReferenceAlias{III.2.1.3.2-es}{III.2.1.3.2-fr}
\EGAReferenceAlias{III.2.1.3.2.fr}{III.2.1.3.2-fr}
\EGAReferenceAlias{III.2.1.3.2.es}{III.2.1.3.2-fr}
\EGAReferenceAlias{III.2.1.3.2.eq}{III.2.1.3.2-fr}
\EGAReferenceAlias{eq:III.2.1.3.2}{III.2.1.3.2-fr}
\EGAReferenceAlias{III.2.1.4-es}{III.2.1.4}
\EGAReferenceAlias{III.2.1.4.fr}{III.2.1.4}
\EGAReferenceAlias{III.2.1.4.es}{III.2.1.4}
\EGAReferenceAlias{III.2.1.4.eq}{III.2.1.4}
\EGAReferenceAlias{eq:III.2.1.4}{III.2.1.4}
\EGAReferenceAlias{III.2.1.4.1}{III.2.1.4.1-fr}
\EGAReferenceAlias{III.2.1.4.1-es}{III.2.1.4.1-fr}
\EGAReferenceAlias{III.2.1.4.1.fr}{III.2.1.4.1-fr}
\EGAReferenceAlias{III.2.1.4.1.es}{III.2.1.4.1-fr}
\EGAReferenceAlias{III.2.1.4.1.eq}{III.2.1.4.1-fr}
\EGAReferenceAlias{eq:III.2.1.4.1}{III.2.1.4.1-fr}
\EGAReferenceAlias{III.2.1.4.2}{III.2.1.4.2-fr}
\EGAReferenceAlias{III.2.1.4.2-es}{III.2.1.4.2-fr}
\EGAReferenceAlias{III.2.1.4.2.fr}{III.2.1.4.2-fr}
\EGAReferenceAlias{III.2.1.4.2.es}{III.2.1.4.2-fr}
\EGAReferenceAlias{III.2.1.4.2.eq}{III.2.1.4.2-fr}
\EGAReferenceAlias{eq:III.2.1.4.2}{III.2.1.4.2-fr}
\EGAReferenceAlias{III.2.1.5-es}{III.2.1.5}
\EGAReferenceAlias{III.2.1.5.fr}{III.2.1.5}
\EGAReferenceAlias{III.2.1.5.es}{III.2.1.5}
\EGAReferenceAlias{III.2.1.5.eq}{III.2.1.5}
\EGAReferenceAlias{eq:III.2.1.5}{III.2.1.5}
\EGAReferenceAlias{III.2.1.5.1}{III.2.1.5.1-fr}
\EGAReferenceAlias{III.2.1.5.1-es}{III.2.1.5.1-fr}
\EGAReferenceAlias{III.2.1.5.1.fr}{III.2.1.5.1-fr}
\EGAReferenceAlias{III.2.1.5.1.es}{III.2.1.5.1-fr}
\EGAReferenceAlias{III.2.1.5.1.eq}{III.2.1.5.1-fr}
\EGAReferenceAlias{eq:III.2.1.5.1}{III.2.1.5.1-fr}
\EGAReferenceAlias{III.2.1.5.2}{III.2.1.5.2-fr}
\EGAReferenceAlias{III.2.1.5.2-es}{III.2.1.5.2-fr}
\EGAReferenceAlias{III.2.1.5.2.fr}{III.2.1.5.2-fr}
\EGAReferenceAlias{III.2.1.5.2.es}{III.2.1.5.2-fr}
\EGAReferenceAlias{III.2.1.5.2.eq}{III.2.1.5.2-fr}
\EGAReferenceAlias{eq:III.2.1.5.2}{III.2.1.5.2-fr}
\EGAReferenceAlias{III.2.1.6-es}{III.2.1.6}
\EGAReferenceAlias{III.2.1.6.fr}{III.2.1.6}
\EGAReferenceAlias{III.2.1.6.es}{III.2.1.6}
\EGAReferenceAlias{III.2.1.6.eq}{III.2.1.6}
\EGAReferenceAlias{eq:III.2.1.6}{III.2.1.6}
\EGAReferenceAlias{III.2.1.6.1}{III.2.1.6.1-fr}
\EGAReferenceAlias{III.2.1.6.1-es}{III.2.1.6.1-fr}
\EGAReferenceAlias{III.2.1.6.1.fr}{III.2.1.6.1-fr}
\EGAReferenceAlias{III.2.1.6.1.es}{III.2.1.6.1-fr}
\EGAReferenceAlias{III.2.1.6.1.eq}{III.2.1.6.1-fr}
\EGAReferenceAlias{eq:III.2.1.6.1}{III.2.1.6.1-fr}
\EGAReferenceAlias{III.2.1.7-es}{III.2.1.7}
\EGAReferenceAlias{III.2.1.7.fr}{III.2.1.7}
\EGAReferenceAlias{III.2.1.7.es}{III.2.1.7}
\EGAReferenceAlias{III.2.1.7.eq}{III.2.1.7}
\EGAReferenceAlias{eq:III.2.1.7}{III.2.1.7}
\EGAReferenceAlias{III.2.1.8-es}{III.2.1.8}
\EGAReferenceAlias{III.2.1.8.fr}{III.2.1.8}
\EGAReferenceAlias{III.2.1.8.es}{III.2.1.8}
\EGAReferenceAlias{III.2.1.8.eq}{III.2.1.8}
\EGAReferenceAlias{eq:III.2.1.8}{III.2.1.8}
\EGAReferenceAlias{III.2.1.8.1}{III.2.1.8.1-fr}
\EGAReferenceAlias{III.2.1.8.1-es}{III.2.1.8.1-fr}
\EGAReferenceAlias{III.2.1.8.1.fr}{III.2.1.8.1-fr}
\EGAReferenceAlias{III.2.1.8.1.es}{III.2.1.8.1-fr}
\EGAReferenceAlias{III.2.1.8.1.eq}{III.2.1.8.1-fr}
\EGAReferenceAlias{eq:III.2.1.8.1}{III.2.1.8.1-fr}
\EGAReferenceAlias{III.2.1.8.2}{III.2.1.8.2-fr}
\EGAReferenceAlias{III.2.1.8.2-es}{III.2.1.8.2-fr}
\EGAReferenceAlias{III.2.1.8.2.fr}{III.2.1.8.2-fr}
\EGAReferenceAlias{III.2.1.8.2.es}{III.2.1.8.2-fr}
\EGAReferenceAlias{III.2.1.8.2.eq}{III.2.1.8.2-fr}
\EGAReferenceAlias{eq:III.2.1.8.2}{III.2.1.8.2-fr}
\EGAReferenceAlias{III.2.1.9}{III.2.1.9-fr}
\EGAReferenceAlias{III.2.1.9-es}{III.2.1.9-fr}
\EGAReferenceAlias{III.2.1.9.fr}{III.2.1.9-fr}
\EGAReferenceAlias{III.2.1.9.es}{III.2.1.9-fr}
\EGAReferenceAlias{III.2.1.9.eq}{III.2.1.9-fr}
\EGAReferenceAlias{eq:III.2.1.9}{III.2.1.9-fr}
\EGAReferenceAlias{III.2.1.10-es}{III.2.1.10}
\EGAReferenceAlias{III.2.1.10.fr}{III.2.1.10}
\EGAReferenceAlias{III.2.1.10.es}{III.2.1.10}
\EGAReferenceAlias{III.2.1.10.eq}{III.2.1.10}
\EGAReferenceAlias{eq:III.2.1.10}{III.2.1.10}
\EGAReferenceAlias{III.2.1.10.1}{III.2.1.10.1-fr}
\EGAReferenceAlias{III.2.1.10.1-es}{III.2.1.10.1-fr}
\EGAReferenceAlias{III.2.1.10.1.fr}{III.2.1.10.1-fr}
\EGAReferenceAlias{III.2.1.10.1.es}{III.2.1.10.1-fr}
\EGAReferenceAlias{III.2.1.10.1.eq}{III.2.1.10.1-fr}
\EGAReferenceAlias{eq:III.2.1.10.1}{III.2.1.10.1-fr}
\EGAReferenceAlias{III.2.1.11-es}{III.2.1.11}
\EGAReferenceAlias{III.2.1.11.fr}{III.2.1.11}
\EGAReferenceAlias{III.2.1.11.es}{III.2.1.11}
\EGAReferenceAlias{III.2.1.11.eq}{III.2.1.11}
\EGAReferenceAlias{eq:III.2.1.11}{III.2.1.11}
\EGAReferenceAlias{III.2.1.11.1}{III.2.1.11.1-fr}
\EGAReferenceAlias{III.2.1.11.1-es}{III.2.1.11.1-fr}
\EGAReferenceAlias{III.2.1.11.1.fr}{III.2.1.11.1-fr}
\EGAReferenceAlias{III.2.1.11.1.es}{III.2.1.11.1-fr}
\EGAReferenceAlias{III.2.1.11.1.eq}{III.2.1.11.1-fr}
\EGAReferenceAlias{eq:III.2.1.11.1}{III.2.1.11.1-fr}
\EGAReferenceAlias{III.2.1.11.2}{III.2.1.11.2-fr}
\EGAReferenceAlias{III.2.1.11.2-es}{III.2.1.11.2-fr}
\EGAReferenceAlias{III.2.1.11.2.fr}{III.2.1.11.2-fr}
\EGAReferenceAlias{III.2.1.11.2.es}{III.2.1.11.2-fr}
\EGAReferenceAlias{III.2.1.11.2.eq}{III.2.1.11.2-fr}
\EGAReferenceAlias{eq:III.2.1.11.2}{III.2.1.11.2-fr}
\EGAReferenceAlias{III.2.1.12-es}{III.2.1.12}
\EGAReferenceAlias{III.2.1.12.fr}{III.2.1.12}
\EGAReferenceAlias{III.2.1.12.es}{III.2.1.12}
\EGAReferenceAlias{III.2.1.12.eq}{III.2.1.12}
\EGAReferenceAlias{eq:III.2.1.12}{III.2.1.12}
\EGAReferenceAlias{III.2.1.13-es}{III.2.1.13}
\EGAReferenceAlias{III.2.1.13.fr}{III.2.1.13}
\EGAReferenceAlias{III.2.1.13.es}{III.2.1.13}
\EGAReferenceAlias{III.2.1.13.eq}{III.2.1.13}
\EGAReferenceAlias{eq:III.2.1.13}{III.2.1.13}
\EGAReferenceAlias{III.2.1.14-es}{III.2.1.14}
\EGAReferenceAlias{III.2.1.14.fr}{III.2.1.14}
\EGAReferenceAlias{III.2.1.14.es}{III.2.1.14}
\EGAReferenceAlias{III.2.1.14.eq}{III.2.1.14}
\EGAReferenceAlias{eq:III.2.1.14}{III.2.1.14}
\EGAReferenceAlias{III.2.1.15-es}{III.2.1.15}
\EGAReferenceAlias{III.2.1.15.fr}{III.2.1.15}
\EGAReferenceAlias{III.2.1.15.es}{III.2.1.15}
\EGAReferenceAlias{III.2.1.15.eq}{III.2.1.15}
\EGAReferenceAlias{eq:III.2.1.15}{III.2.1.15}
\EGAReferenceAlias{III.2.1.16-es}{III.2.1.16}
\EGAReferenceAlias{III.2.1.16.fr}{III.2.1.16}
\EGAReferenceAlias{III.2.1.16.es}{III.2.1.16}
\EGAReferenceAlias{III.2.1.16.eq}{III.2.1.16}
\EGAReferenceAlias{eq:III.2.1.16}{III.2.1.16}
\EGAReferenceAlias{III.2.1.16.1}{III.2.1.16.1-fr}
\EGAReferenceAlias{III.2.1.16.1-es}{III.2.1.16.1-fr}
\EGAReferenceAlias{III.2.1.16.1.fr}{III.2.1.16.1-fr}
\EGAReferenceAlias{III.2.1.16.1.es}{III.2.1.16.1-fr}
\EGAReferenceAlias{III.2.1.16.1.eq}{III.2.1.16.1-fr}
\EGAReferenceAlias{eq:III.2.1.16.1}{III.2.1.16.1-fr}
\EGAReferenceAlias{III.2.1.16.2}{III.2.1.16.2-fr}
\EGAReferenceAlias{III.2.1.16.2-es}{III.2.1.16.2-fr}
\EGAReferenceAlias{III.2.1.16.2.fr}{III.2.1.16.2-fr}
\EGAReferenceAlias{III.2.1.16.2.es}{III.2.1.16.2-fr}
\EGAReferenceAlias{III.2.1.16.2.eq}{III.2.1.16.2-fr}
\EGAReferenceAlias{eq:III.2.1.16.2}{III.2.1.16.2-fr}
\EGAReferenceAlias{III.2.2}{subsection:III.2.2}
\EGAReferenceAlias{III.2.2-fr}{subsection:III.2.2}
\EGAReferenceAlias{III.2.2-es}{subsection:III.2.2}
\EGAReferenceAlias{III.2.2.fr}{subsection:III.2.2}
\EGAReferenceAlias{III.2.2.es}{subsection:III.2.2}
\EGAReferenceAlias{III.2.2.eq}{subsection:III.2.2}
\EGAReferenceAlias{eq:III.2.2}{subsection:III.2.2}
\EGAReferenceAlias{subsection:III.2.2-es}{subsection:III.2.2}
\EGAReferenceAlias{subsection:III.2.2.fr}{subsection:III.2.2}
\EGAReferenceAlias{subsection:III.2.2.es}{subsection:III.2.2}
\EGAReferenceAlias{subsection:III.2.2.eq}{subsection:III.2.2}
\EGAReferenceAlias{eq:subsection:III.2.2}{subsection:III.2.2}
\EGAReferenceAlias{III.2.2.1}{III.2.2.1-fr}
\EGAReferenceAlias{III.2.2.1-es}{III.2.2.1-fr}
\EGAReferenceAlias{III.2.2.1.fr}{III.2.2.1-fr}
\EGAReferenceAlias{III.2.2.1.es}{III.2.2.1-fr}
\EGAReferenceAlias{III.2.2.1.eq}{III.2.2.1-fr}
\EGAReferenceAlias{eq:III.2.2.1}{III.2.2.1-fr}
\EGAReferenceAlias{III.2.2.2}{III.2.2.2-fr}
\EGAReferenceAlias{III.2.2.2-es}{III.2.2.2-fr}
\EGAReferenceAlias{III.2.2.2.fr}{III.2.2.2-fr}
\EGAReferenceAlias{III.2.2.2.es}{III.2.2.2-fr}
\EGAReferenceAlias{III.2.2.2.eq}{III.2.2.2-fr}
\EGAReferenceAlias{eq:III.2.2.2}{III.2.2.2-fr}
\EGAReferenceAlias{III.2.2.3-es}{III.2.2.3}
\EGAReferenceAlias{III.2.2.3.fr}{III.2.2.3}
\EGAReferenceAlias{III.2.2.3.es}{III.2.2.3}
\EGAReferenceAlias{III.2.2.3.eq}{III.2.2.3}
\EGAReferenceAlias{eq:III.2.2.3}{III.2.2.3}
\EGAReferenceAlias{III.2.2.4-es}{III.2.2.4}
\EGAReferenceAlias{III.2.2.4.fr}{III.2.2.4}
\EGAReferenceAlias{III.2.2.4.es}{III.2.2.4}
\EGAReferenceAlias{III.2.2.4.eq}{III.2.2.4}
\EGAReferenceAlias{eq:III.2.2.4}{III.2.2.4}
\EGAReferenceAlias{III.2.2.5-es}{III.2.2.5}
\EGAReferenceAlias{III.2.2.5.fr}{III.2.2.5}
\EGAReferenceAlias{III.2.2.5.es}{III.2.2.5}
\EGAReferenceAlias{III.2.2.5.eq}{III.2.2.5}
\EGAReferenceAlias{eq:III.2.2.5}{III.2.2.5}
\EGAReferenceAlias{III.2.3}{subsection:III.2.3}
\EGAReferenceAlias{III.2.3-fr}{subsection:III.2.3}
\EGAReferenceAlias{III.2.3-es}{subsection:III.2.3}
\EGAReferenceAlias{III.2.3.fr}{subsection:III.2.3}
\EGAReferenceAlias{III.2.3.es}{subsection:III.2.3}
\EGAReferenceAlias{III.2.3.eq}{subsection:III.2.3}
\EGAReferenceAlias{eq:III.2.3}{subsection:III.2.3}
\EGAReferenceAlias{subsection:III.2.3-es}{subsection:III.2.3}
\EGAReferenceAlias{subsection:III.2.3.fr}{subsection:III.2.3}
\EGAReferenceAlias{subsection:III.2.3.es}{subsection:III.2.3}
\EGAReferenceAlias{subsection:III.2.3.eq}{subsection:III.2.3}
\EGAReferenceAlias{eq:subsection:III.2.3}{subsection:III.2.3}
\EGAReferenceAlias{III.2.3.1-es}{III.2.3.1}
\EGAReferenceAlias{III.2.3.1.fr}{III.2.3.1}
\EGAReferenceAlias{III.2.3.1.es}{III.2.3.1}
\EGAReferenceAlias{III.2.3.1.eq}{III.2.3.1}
\EGAReferenceAlias{eq:III.2.3.1}{III.2.3.1}
\EGAReferenceAlias{III.2.3.2-es}{III.2.3.2}
\EGAReferenceAlias{III.2.3.2.fr}{III.2.3.2}
\EGAReferenceAlias{III.2.3.2.es}{III.2.3.2}
\EGAReferenceAlias{III.2.3.2.eq}{III.2.3.2}
\EGAReferenceAlias{eq:III.2.3.2}{III.2.3.2}
\EGAReferenceAlias{III.2.3.3-es}{III.2.3.3}
\EGAReferenceAlias{III.2.3.3.fr}{III.2.3.3}
\EGAReferenceAlias{III.2.3.3.es}{III.2.3.3}
\EGAReferenceAlias{III.2.3.3.eq}{III.2.3.3}
\EGAReferenceAlias{eq:III.2.3.3}{III.2.3.3}
\EGAReferenceAlias{III.2.3.4}{III.2.3.4-fr}
\EGAReferenceAlias{III.2.3.4-es}{III.2.3.4-fr}
\EGAReferenceAlias{III.2.3.4.fr}{III.2.3.4-fr}
\EGAReferenceAlias{III.2.3.4.es}{III.2.3.4-fr}
\EGAReferenceAlias{III.2.3.4.eq}{III.2.3.4-fr}
\EGAReferenceAlias{eq:III.2.3.4}{III.2.3.4-fr}
\EGAReferenceAlias{III.2.3.4.1}{III.2.3.4.1-fr}
\EGAReferenceAlias{III.2.3.4.1-es}{III.2.3.4.1-fr}
\EGAReferenceAlias{III.2.3.4.1.fr}{III.2.3.4.1-fr}
\EGAReferenceAlias{III.2.3.4.1.es}{III.2.3.4.1-fr}
\EGAReferenceAlias{III.2.3.4.1.eq}{III.2.3.4.1-fr}
\EGAReferenceAlias{eq:III.2.3.4.1}{III.2.3.4.1-fr}
\EGAReferenceAlias{III.2.3.5.1}{III.2.3.5.1-fr}
\EGAReferenceAlias{III.2.3.5.1-es}{III.2.3.5.1-fr}
\EGAReferenceAlias{III.2.3.5.1.fr}{III.2.3.5.1-fr}
\EGAReferenceAlias{III.2.3.5.1.es}{III.2.3.5.1-fr}
\EGAReferenceAlias{III.2.3.5.1.eq}{III.2.3.5.1-fr}
\EGAReferenceAlias{eq:III.2.3.5.1}{III.2.3.5.1-fr}
\EGAReferenceAlias{III.2.3.5}{III.2.3.5-fr}
\EGAReferenceAlias{III.2.3.5-es}{III.2.3.5-fr}
\EGAReferenceAlias{III.2.3.5.fr}{III.2.3.5-fr}
\EGAReferenceAlias{III.2.3.5.es}{III.2.3.5-fr}
\EGAReferenceAlias{III.2.3.5.eq}{III.2.3.5-fr}
\EGAReferenceAlias{eq:III.2.3.5}{III.2.3.5-fr}
\EGAReferenceAlias{III.2.3.6-es}{III.2.3.6}
\EGAReferenceAlias{III.2.3.6.fr}{III.2.3.6}
\EGAReferenceAlias{III.2.3.6.es}{III.2.3.6}
\EGAReferenceAlias{III.2.3.6.eq}{III.2.3.6}
\EGAReferenceAlias{eq:III.2.3.6}{III.2.3.6}
\EGAReferenceAlias{III.2.3.7-es}{III.2.3.7}
\EGAReferenceAlias{III.2.3.7.fr}{III.2.3.7}
\EGAReferenceAlias{III.2.3.7.es}{III.2.3.7}
\EGAReferenceAlias{III.2.3.7.eq}{III.2.3.7}
\EGAReferenceAlias{eq:III.2.3.7}{III.2.3.7}
\EGAReferenceAlias{III.2.3.8-es}{III.2.3.8}
\EGAReferenceAlias{III.2.3.8.fr}{III.2.3.8}
\EGAReferenceAlias{III.2.3.8.es}{III.2.3.8}
\EGAReferenceAlias{III.2.3.8.eq}{III.2.3.8}
\EGAReferenceAlias{eq:III.2.3.8}{III.2.3.8}
\EGAReferenceAlias{III.2.4}{subsection:III.2.4}
\EGAReferenceAlias{III.2.4-fr}{subsection:III.2.4}
\EGAReferenceAlias{III.2.4-es}{subsection:III.2.4}
\EGAReferenceAlias{III.2.4.fr}{subsection:III.2.4}
\EGAReferenceAlias{III.2.4.es}{subsection:III.2.4}
\EGAReferenceAlias{III.2.4.eq}{subsection:III.2.4}
\EGAReferenceAlias{eq:III.2.4}{subsection:III.2.4}
\EGAReferenceAlias{subsection:III.2.4-es}{subsection:III.2.4}
\EGAReferenceAlias{subsection:III.2.4.fr}{subsection:III.2.4}
\EGAReferenceAlias{subsection:III.2.4.es}{subsection:III.2.4}
\EGAReferenceAlias{subsection:III.2.4.eq}{subsection:III.2.4}
\EGAReferenceAlias{eq:subsection:III.2.4}{subsection:III.2.4}
\EGAReferenceAlias{III.2.4.1}{III.2.4.1-fr}
\EGAReferenceAlias{III.2.4.1-es}{III.2.4.1-fr}
\EGAReferenceAlias{III.2.4.1.fr}{III.2.4.1-fr}
\EGAReferenceAlias{III.2.4.1.es}{III.2.4.1-fr}
\EGAReferenceAlias{III.2.4.1.eq}{III.2.4.1-fr}
\EGAReferenceAlias{eq:III.2.4.1}{III.2.4.1-fr}
\EGAReferenceAlias{III.2.4.1.1}{III.2.4.1.1-fr}
\EGAReferenceAlias{III.2.4.1.1-es}{III.2.4.1.1-fr}
\EGAReferenceAlias{III.2.4.1.1.fr}{III.2.4.1.1-fr}
\EGAReferenceAlias{III.2.4.1.1.es}{III.2.4.1.1-fr}
\EGAReferenceAlias{III.2.4.1.1.eq}{III.2.4.1.1-fr}
\EGAReferenceAlias{eq:III.2.4.1.1}{III.2.4.1.1-fr}
\EGAReferenceAlias{III.2.4.1.2}{III.2.4.1.2-fr}
\EGAReferenceAlias{III.2.4.1.2-es}{III.2.4.1.2-fr}
\EGAReferenceAlias{III.2.4.1.2.fr}{III.2.4.1.2-fr}
\EGAReferenceAlias{III.2.4.1.2.es}{III.2.4.1.2-fr}
\EGAReferenceAlias{III.2.4.1.2.eq}{III.2.4.1.2-fr}
\EGAReferenceAlias{eq:III.2.4.1.2}{III.2.4.1.2-fr}
\EGAReferenceAlias{III.2.4.2-es}{III.2.4.2}
\EGAReferenceAlias{III.2.4.2.fr}{III.2.4.2}
\EGAReferenceAlias{III.2.4.2.es}{III.2.4.2}
\EGAReferenceAlias{III.2.4.2.eq}{III.2.4.2}
\EGAReferenceAlias{eq:III.2.4.2}{III.2.4.2}
\EGAReferenceAlias{III.2.5}{subsection:III.2.5}
\EGAReferenceAlias{III.2.5-fr}{subsection:III.2.5}
\EGAReferenceAlias{III.2.5-es}{subsection:III.2.5}
\EGAReferenceAlias{III.2.5.fr}{subsection:III.2.5}
\EGAReferenceAlias{III.2.5.es}{subsection:III.2.5}
\EGAReferenceAlias{III.2.5.eq}{subsection:III.2.5}
\EGAReferenceAlias{eq:III.2.5}{subsection:III.2.5}
\EGAReferenceAlias{subsection:III.2.5-es}{subsection:III.2.5}
\EGAReferenceAlias{subsection:III.2.5.fr}{subsection:III.2.5}
\EGAReferenceAlias{subsection:III.2.5.es}{subsection:III.2.5}
\EGAReferenceAlias{subsection:III.2.5.eq}{subsection:III.2.5}
\EGAReferenceAlias{eq:subsection:III.2.5}{subsection:III.2.5}
\EGAReferenceAlias{III.2.5.1-es}{III.2.5.1}
\EGAReferenceAlias{III.2.5.1.fr}{III.2.5.1}
\EGAReferenceAlias{III.2.5.1.es}{III.2.5.1}
\EGAReferenceAlias{III.2.5.1.eq}{III.2.5.1}
\EGAReferenceAlias{eq:III.2.5.1}{III.2.5.1}
\EGAReferenceAlias{III.2.5.1.1}{III.2.5.1.1-fr}
\EGAReferenceAlias{III.2.5.1.1-es}{III.2.5.1.1-fr}
\EGAReferenceAlias{III.2.5.1.1.fr}{III.2.5.1.1-fr}
\EGAReferenceAlias{III.2.5.1.1.es}{III.2.5.1.1-fr}
\EGAReferenceAlias{III.2.5.1.1.eq}{III.2.5.1.1-fr}
\EGAReferenceAlias{eq:III.2.5.1.1}{III.2.5.1.1-fr}
\EGAReferenceAlias{III.2.5.2-es}{III.2.5.2}
\EGAReferenceAlias{III.2.5.2.fr}{III.2.5.2}
\EGAReferenceAlias{III.2.5.2.es}{III.2.5.2}
\EGAReferenceAlias{III.2.5.2.eq}{III.2.5.2}
\EGAReferenceAlias{eq:III.2.5.2}{III.2.5.2}
\EGAReferenceAlias{III.2.5.2.1}{III.2.5.2.1-fr}
\EGAReferenceAlias{III.2.5.2.1-es}{III.2.5.2.1-fr}
\EGAReferenceAlias{III.2.5.2.1.fr}{III.2.5.2.1-fr}
\EGAReferenceAlias{III.2.5.2.1.es}{III.2.5.2.1-fr}
\EGAReferenceAlias{III.2.5.2.1.eq}{III.2.5.2.1-fr}
\EGAReferenceAlias{eq:III.2.5.2.1}{III.2.5.2.1-fr}
\EGAReferenceAlias{III.2.5.3}{III.2.5.3-fr}
\EGAReferenceAlias{III.2.5.3-es}{III.2.5.3-fr}
\EGAReferenceAlias{III.2.5.3.fr}{III.2.5.3-fr}
\EGAReferenceAlias{III.2.5.3.es}{III.2.5.3-fr}
\EGAReferenceAlias{III.2.5.3.eq}{III.2.5.3-fr}
\EGAReferenceAlias{eq:III.2.5.3}{III.2.5.3-fr}
\EGAReferenceAlias{III.2.5.3.1}{III.2.5.3.1-fr}
\EGAReferenceAlias{III.2.5.3.1-es}{III.2.5.3.1-fr}
\EGAReferenceAlias{III.2.5.3.1.fr}{III.2.5.3.1-fr}
\EGAReferenceAlias{III.2.5.3.1.es}{III.2.5.3.1-fr}
\EGAReferenceAlias{III.2.5.3.1.eq}{III.2.5.3.1-fr}
\EGAReferenceAlias{eq:III.2.5.3.1}{III.2.5.3.1-fr}
\EGAReferenceAlias{III.2.5.4-es}{III.2.5.4}
\EGAReferenceAlias{III.2.5.4.fr}{III.2.5.4}
\EGAReferenceAlias{III.2.5.4.es}{III.2.5.4}
\EGAReferenceAlias{III.2.5.4.eq}{III.2.5.4}
\EGAReferenceAlias{eq:III.2.5.4}{III.2.5.4}
\EGAReferenceAlias{III.2.6}{subsection:III.2.6-fr}
\EGAReferenceAlias{III.2.6-fr}{subsection:III.2.6-fr}
\EGAReferenceAlias{III.2.6-es}{subsection:III.2.6-fr}
\EGAReferenceAlias{III.2.6.fr}{subsection:III.2.6-fr}
\EGAReferenceAlias{III.2.6.es}{subsection:III.2.6-fr}
\EGAReferenceAlias{III.2.6.eq}{subsection:III.2.6-fr}
\EGAReferenceAlias{eq:III.2.6}{subsection:III.2.6-fr}
\EGAReferenceAlias{subsection:III.2.6}{subsection:III.2.6-fr}
\EGAReferenceAlias{subsection:III.2.6-es}{subsection:III.2.6-fr}
\EGAReferenceAlias{subsection:III.2.6.fr}{subsection:III.2.6-fr}
\EGAReferenceAlias{subsection:III.2.6.es}{subsection:III.2.6-fr}
\EGAReferenceAlias{subsection:III.2.6.eq}{subsection:III.2.6-fr}
\EGAReferenceAlias{eq:subsection:III.2.6}{subsection:III.2.6-fr}
\EGAReferenceAlias{III.2.6.1}{III.2.6.1-fr}
\EGAReferenceAlias{III.2.6.1-es}{III.2.6.1-fr}
\EGAReferenceAlias{III.2.6.1.fr}{III.2.6.1-fr}
\EGAReferenceAlias{III.2.6.1.es}{III.2.6.1-fr}
\EGAReferenceAlias{III.2.6.1.eq}{III.2.6.1-fr}
\EGAReferenceAlias{eq:III.2.6.1}{III.2.6.1-fr}
\EGAReferenceAlias{III.2.6.2}{III.2.6.2-fr}
\EGAReferenceAlias{III.2.6.2-es}{III.2.6.2-fr}
\EGAReferenceAlias{III.2.6.2.fr}{III.2.6.2-fr}
\EGAReferenceAlias{III.2.6.2.es}{III.2.6.2-fr}
\EGAReferenceAlias{III.2.6.2.eq}{III.2.6.2-fr}
\EGAReferenceAlias{eq:III.2.6.2}{III.2.6.2-fr}
\EGAReferenceAlias{III.2.6.2.1}{III.2.6.2.1-fr}
\EGAReferenceAlias{III.2.6.2.1-es}{III.2.6.2.1-fr}
\EGAReferenceAlias{III.2.6.2.1.fr}{III.2.6.2.1-fr}
\EGAReferenceAlias{III.2.6.2.1.es}{III.2.6.2.1-fr}
\EGAReferenceAlias{III.2.6.2.1.eq}{III.2.6.2.1-fr}
\EGAReferenceAlias{eq:III.2.6.2.1}{III.2.6.2.1-fr}
\EGAReferenceAlias{III.2.6.2.2}{III.2.6.2.2-fr}
\EGAReferenceAlias{III.2.6.2.2-es}{III.2.6.2.2-fr}
\EGAReferenceAlias{III.2.6.2.2.fr}{III.2.6.2.2-fr}
\EGAReferenceAlias{III.2.6.2.2.es}{III.2.6.2.2-fr}
\EGAReferenceAlias{III.2.6.2.2.eq}{III.2.6.2.2-fr}
\EGAReferenceAlias{eq:III.2.6.2.2}{III.2.6.2.2-fr}
\EGAReferenceAlias{III.2.6.2.3}{III.2.6.2.3-fr}
\EGAReferenceAlias{III.2.6.2.3-es}{III.2.6.2.3-fr}
\EGAReferenceAlias{III.2.6.2.3.fr}{III.2.6.2.3-fr}
\EGAReferenceAlias{III.2.6.2.3.es}{III.2.6.2.3-fr}
\EGAReferenceAlias{III.2.6.2.3.eq}{III.2.6.2.3-fr}
\EGAReferenceAlias{eq:III.2.6.2.3}{III.2.6.2.3-fr}
\EGAReferenceAlias{III.2.6.2.4}{III.2.6.2.4-fr}
\EGAReferenceAlias{III.2.6.2.4-es}{III.2.6.2.4-fr}
\EGAReferenceAlias{III.2.6.2.4.fr}{III.2.6.2.4-fr}
\EGAReferenceAlias{III.2.6.2.4.es}{III.2.6.2.4-fr}
\EGAReferenceAlias{III.2.6.2.4.eq}{III.2.6.2.4-fr}
\EGAReferenceAlias{eq:III.2.6.2.4}{III.2.6.2.4-fr}
\EGAReferenceAlias{III.2.6.2.5}{III.2.6.2.5-fr}
\EGAReferenceAlias{III.2.6.2.5-es}{III.2.6.2.5-fr}
\EGAReferenceAlias{III.2.6.2.5.fr}{III.2.6.2.5-fr}
\EGAReferenceAlias{III.2.6.2.5.es}{III.2.6.2.5-fr}
\EGAReferenceAlias{III.2.6.2.5.eq}{III.2.6.2.5-fr}
\EGAReferenceAlias{eq:III.2.6.2.5}{III.2.6.2.5-fr}
\EGAReferenceAlias{III.2.6.2.6}{III.2.6.2.6-fr}
\EGAReferenceAlias{III.2.6.2.6-es}{III.2.6.2.6-fr}
\EGAReferenceAlias{III.2.6.2.6.fr}{III.2.6.2.6-fr}
\EGAReferenceAlias{III.2.6.2.6.es}{III.2.6.2.6-fr}
\EGAReferenceAlias{III.2.6.2.6.eq}{III.2.6.2.6-fr}
\EGAReferenceAlias{eq:III.2.6.2.6}{III.2.6.2.6-fr}
\EGAReferenceAlias{III.2.6.3-es}{III.2.6.3}
\EGAReferenceAlias{III.2.6.3.fr}{III.2.6.3}
\EGAReferenceAlias{III.2.6.3.es}{III.2.6.3}
\EGAReferenceAlias{III.2.6.3.eq}{III.2.6.3}
\EGAReferenceAlias{eq:III.2.6.3}{III.2.6.3}
\EGAReferenceAlias{III.3}{section:III.3}
\EGAReferenceAlias{III.3-fr}{section:III.3}
\EGAReferenceAlias{III.3-es}{section:III.3}
\EGAReferenceAlias{III.3.fr}{section:III.3}
\EGAReferenceAlias{III.3.es}{section:III.3}
\EGAReferenceAlias{III.3.eq}{section:III.3}
\EGAReferenceAlias{eq:III.3}{section:III.3}
\EGAReferenceAlias{section:III.3-es}{section:III.3}
\EGAReferenceAlias{section:III.3.fr}{section:III.3}
\EGAReferenceAlias{section:III.3.es}{section:III.3}
\EGAReferenceAlias{section:III.3.eq}{section:III.3}
\EGAReferenceAlias{eq:section:III.3}{section:III.3}
\EGAReferenceAlias{III.3.1}{subsection:III.3.1}
\EGAReferenceAlias{III.3.1-fr}{subsection:III.3.1}
\EGAReferenceAlias{III.3.1-es}{subsection:III.3.1}
\EGAReferenceAlias{III.3.1.fr}{subsection:III.3.1}
\EGAReferenceAlias{III.3.1.es}{subsection:III.3.1}
\EGAReferenceAlias{III.3.1.eq}{subsection:III.3.1}
\EGAReferenceAlias{eq:III.3.1}{subsection:III.3.1}
\EGAReferenceAlias{subsection:III.3.1-es}{subsection:III.3.1}
\EGAReferenceAlias{subsection:III.3.1.fr}{subsection:III.3.1}
\EGAReferenceAlias{subsection:III.3.1.es}{subsection:III.3.1}
\EGAReferenceAlias{subsection:III.3.1.eq}{subsection:III.3.1}
\EGAReferenceAlias{eq:subsection:III.3.1}{subsection:III.3.1}
\EGAReferenceAlias{III.3.1.1}{III.3.1.1-fr}
\EGAReferenceAlias{III.3.1.1-es}{III.3.1.1-fr}
\EGAReferenceAlias{III.3.1.1.fr}{III.3.1.1-fr}
\EGAReferenceAlias{III.3.1.1.es}{III.3.1.1-fr}
\EGAReferenceAlias{III.3.1.1.eq}{III.3.1.1-fr}
\EGAReferenceAlias{eq:III.3.1.1}{III.3.1.1-fr}
\EGAReferenceAlias{III.3.1.2}{III.3.1.2-fr}
\EGAReferenceAlias{III.3.1.2-es}{III.3.1.2-fr}
\EGAReferenceAlias{III.3.1.2.fr}{III.3.1.2-fr}
\EGAReferenceAlias{III.3.1.2.es}{III.3.1.2-fr}
\EGAReferenceAlias{III.3.1.2.eq}{III.3.1.2-fr}
\EGAReferenceAlias{eq:III.3.1.2}{III.3.1.2-fr}
\EGAReferenceAlias{III.3.1.3}{III.3.1.3-fr}
\EGAReferenceAlias{III.3.1.3-es}{III.3.1.3-fr}
\EGAReferenceAlias{III.3.1.3.fr}{III.3.1.3-fr}
\EGAReferenceAlias{III.3.1.3.es}{III.3.1.3-fr}
\EGAReferenceAlias{III.3.1.3.eq}{III.3.1.3-fr}
\EGAReferenceAlias{eq:III.3.1.3}{III.3.1.3-fr}
\EGAReferenceAlias{III.3.2}{subsection:III.3.2-fr}
\EGAReferenceAlias{III.3.2-fr}{subsection:III.3.2-fr}
\EGAReferenceAlias{III.3.2-es}{subsection:III.3.2-fr}
\EGAReferenceAlias{III.3.2.fr}{subsection:III.3.2-fr}
\EGAReferenceAlias{III.3.2.es}{subsection:III.3.2-fr}
\EGAReferenceAlias{III.3.2.eq}{subsection:III.3.2-fr}
\EGAReferenceAlias{eq:III.3.2}{subsection:III.3.2-fr}
\EGAReferenceAlias{subsection:III.3.2}{subsection:III.3.2-fr}
\EGAReferenceAlias{subsection:III.3.2-es}{subsection:III.3.2-fr}
\EGAReferenceAlias{subsection:III.3.2.fr}{subsection:III.3.2-fr}
\EGAReferenceAlias{subsection:III.3.2.es}{subsection:III.3.2-fr}
\EGAReferenceAlias{subsection:III.3.2.eq}{subsection:III.3.2-fr}
\EGAReferenceAlias{eq:subsection:III.3.2}{subsection:III.3.2-fr}
\EGAReferenceAlias{III.3.2.1}{III.3.2.1-fr}
\EGAReferenceAlias{III.3.2.1-es}{III.3.2.1-fr}
\EGAReferenceAlias{III.3.2.1.fr}{III.3.2.1-fr}
\EGAReferenceAlias{III.3.2.1.es}{III.3.2.1-fr}
\EGAReferenceAlias{III.3.2.1.eq}{III.3.2.1-fr}
\EGAReferenceAlias{eq:III.3.2.1}{III.3.2.1-fr}
\EGAReferenceAlias{III.3.2.2-es}{III.3.2.2}
\EGAReferenceAlias{III.3.2.2.fr}{III.3.2.2}
\EGAReferenceAlias{III.3.2.2.es}{III.3.2.2}
\EGAReferenceAlias{III.3.2.2.eq}{III.3.2.2}
\EGAReferenceAlias{eq:III.3.2.2}{III.3.2.2}
\EGAReferenceAlias{III.3.2.3-es}{III.3.2.3}
\EGAReferenceAlias{III.3.2.3.fr}{III.3.2.3}
\EGAReferenceAlias{III.3.2.3.es}{III.3.2.3}
\EGAReferenceAlias{III.3.2.3.eq}{III.3.2.3}
\EGAReferenceAlias{eq:III.3.2.3}{III.3.2.3}
\EGAReferenceAlias{III.3.2.4-es}{III.3.2.4}
\EGAReferenceAlias{III.3.2.4.fr}{III.3.2.4}
\EGAReferenceAlias{III.3.2.4.es}{III.3.2.4}
\EGAReferenceAlias{III.3.2.4.eq}{III.3.2.4}
\EGAReferenceAlias{eq:III.3.2.4}{III.3.2.4}
\EGAReferenceAlias{III.3.3}{subsection:III.3.3}
\EGAReferenceAlias{III.3.3-fr}{subsection:III.3.3}
\EGAReferenceAlias{III.3.3-es}{subsection:III.3.3}
\EGAReferenceAlias{III.3.3.fr}{subsection:III.3.3}
\EGAReferenceAlias{III.3.3.es}{subsection:III.3.3}
\EGAReferenceAlias{III.3.3.eq}{subsection:III.3.3}
\EGAReferenceAlias{eq:III.3.3}{subsection:III.3.3}
\EGAReferenceAlias{subsection:III.3.3-es}{subsection:III.3.3}
\EGAReferenceAlias{subsection:III.3.3.fr}{subsection:III.3.3}
\EGAReferenceAlias{subsection:III.3.3.es}{subsection:III.3.3}
\EGAReferenceAlias{subsection:III.3.3.eq}{subsection:III.3.3}
\EGAReferenceAlias{eq:subsection:III.3.3}{subsection:III.3.3}
\EGAReferenceAlias{III.3.3.1-es}{III.3.3.1}
\EGAReferenceAlias{III.3.3.1.fr}{III.3.3.1}
\EGAReferenceAlias{III.3.3.1.es}{III.3.3.1}
\EGAReferenceAlias{III.3.3.1.eq}{III.3.3.1}
\EGAReferenceAlias{eq:III.3.3.1}{III.3.3.1}
\EGAReferenceAlias{III.3.3.1.1}{III.3.3.1.1-fr}
\EGAReferenceAlias{III.3.3.1.1-es}{III.3.3.1.1-fr}
\EGAReferenceAlias{III.3.3.1.1.fr}{III.3.3.1.1-fr}
\EGAReferenceAlias{III.3.3.1.1.es}{III.3.3.1.1-fr}
\EGAReferenceAlias{III.3.3.1.1.eq}{III.3.3.1.1-fr}
\EGAReferenceAlias{eq:III.3.3.1.1}{III.3.3.1.1-fr}
\EGAReferenceAlias{III.3.3.2-es}{III.3.3.2}
\EGAReferenceAlias{III.3.3.2.fr}{III.3.3.2}
\EGAReferenceAlias{III.3.3.2.es}{III.3.3.2}
\EGAReferenceAlias{III.3.3.2.eq}{III.3.3.2}
\EGAReferenceAlias{eq:III.3.3.2}{III.3.3.2}
\EGAReferenceAlias{III.3.3.2.1}{III.3.3.2.1-fr}
\EGAReferenceAlias{III.3.3.2.1-es}{III.3.3.2.1-fr}
\EGAReferenceAlias{III.3.3.2.1.fr}{III.3.3.2.1-fr}
\EGAReferenceAlias{III.3.3.2.1.es}{III.3.3.2.1-fr}
\EGAReferenceAlias{III.3.3.2.1.eq}{III.3.3.2.1-fr}
\EGAReferenceAlias{eq:III.3.3.2.1}{III.3.3.2.1-fr}
\EGAReferenceAlias{III.3.4}{subsection:III.3.4}
\EGAReferenceAlias{III.3.4-fr}{subsection:III.3.4}
\EGAReferenceAlias{III.3.4-es}{subsection:III.3.4}
\EGAReferenceAlias{III.3.4.fr}{subsection:III.3.4}
\EGAReferenceAlias{III.3.4.es}{subsection:III.3.4}
\EGAReferenceAlias{III.3.4.eq}{subsection:III.3.4}
\EGAReferenceAlias{eq:III.3.4}{subsection:III.3.4}
\EGAReferenceAlias{subsection:III.3.4-es}{subsection:III.3.4}
\EGAReferenceAlias{subsection:III.3.4.fr}{subsection:III.3.4}
\EGAReferenceAlias{subsection:III.3.4.es}{subsection:III.3.4}
\EGAReferenceAlias{subsection:III.3.4.eq}{subsection:III.3.4}
\EGAReferenceAlias{eq:subsection:III.3.4}{subsection:III.3.4}
\EGAReferenceAlias{III.3.4.1-es}{III.3.4.1}
\EGAReferenceAlias{III.3.4.1.fr}{III.3.4.1}
\EGAReferenceAlias{III.3.4.1.es}{III.3.4.1}
\EGAReferenceAlias{III.3.4.1.eq}{III.3.4.1}
\EGAReferenceAlias{eq:III.3.4.1}{III.3.4.1}
\EGAReferenceAlias{III.3.4.2-es}{III.3.4.2}
\EGAReferenceAlias{III.3.4.2.fr}{III.3.4.2}
\EGAReferenceAlias{III.3.4.2.es}{III.3.4.2}
\EGAReferenceAlias{III.3.4.2.eq}{III.3.4.2}
\EGAReferenceAlias{eq:III.3.4.2}{III.3.4.2}
\EGAReferenceAlias{III.3.4.2.1}{III.3.4.2.1-fr}
\EGAReferenceAlias{III.3.4.2.1-es}{III.3.4.2.1-fr}
\EGAReferenceAlias{III.3.4.2.1.fr}{III.3.4.2.1-fr}
\EGAReferenceAlias{III.3.4.2.1.es}{III.3.4.2.1-fr}
\EGAReferenceAlias{III.3.4.2.1.eq}{III.3.4.2.1-fr}
\EGAReferenceAlias{eq:III.3.4.2.1}{III.3.4.2.1-fr}
\EGAReferenceAlias{III.3.4.2.2}{III.3.4.2.2-fr}
\EGAReferenceAlias{III.3.4.2.2-es}{III.3.4.2.2-fr}
\EGAReferenceAlias{III.3.4.2.2.fr}{III.3.4.2.2-fr}
\EGAReferenceAlias{III.3.4.2.2.es}{III.3.4.2.2-fr}
\EGAReferenceAlias{III.3.4.2.2.eq}{III.3.4.2.2-fr}
\EGAReferenceAlias{eq:III.3.4.2.2}{III.3.4.2.2-fr}
\EGAReferenceAlias{III.3.4.3-es}{III.3.4.3}
\EGAReferenceAlias{III.3.4.3.fr}{III.3.4.3}
\EGAReferenceAlias{III.3.4.3.es}{III.3.4.3}
\EGAReferenceAlias{III.3.4.3.eq}{III.3.4.3}
\EGAReferenceAlias{eq:III.3.4.3}{III.3.4.3}
\EGAReferenceAlias{III.3.4.4-es}{III.3.4.4}
\EGAReferenceAlias{III.3.4.4.fr}{III.3.4.4}
\EGAReferenceAlias{III.3.4.4.es}{III.3.4.4}
\EGAReferenceAlias{III.3.4.4.eq}{III.3.4.4}
\EGAReferenceAlias{eq:III.3.4.4}{III.3.4.4}
\EGAReferenceAlias{III.3.4.4.1}{III.3.4.4.1-fr}
\EGAReferenceAlias{III.3.4.4.1-es}{III.3.4.4.1-fr}
\EGAReferenceAlias{III.3.4.4.1.fr}{III.3.4.4.1-fr}
\EGAReferenceAlias{III.3.4.4.1.es}{III.3.4.4.1-fr}
\EGAReferenceAlias{III.3.4.4.1.eq}{III.3.4.4.1-fr}
\EGAReferenceAlias{eq:III.3.4.4.1}{III.3.4.4.1-fr}
\EGAReferenceAlias{III.3.4.4.2}{III.3.4.4.2-fr}
\EGAReferenceAlias{III.3.4.4.2-es}{III.3.4.4.2-fr}
\EGAReferenceAlias{III.3.4.4.2.fr}{III.3.4.4.2-fr}
\EGAReferenceAlias{III.3.4.4.2.es}{III.3.4.4.2-fr}
\EGAReferenceAlias{III.3.4.4.2.eq}{III.3.4.4.2-fr}
\EGAReferenceAlias{eq:III.3.4.4.2}{III.3.4.4.2-fr}
\EGAReferenceAlias{III.3.4.4.3}{III.3.4.4.3-fr}
\EGAReferenceAlias{III.3.4.4.3-es}{III.3.4.4.3-fr}
\EGAReferenceAlias{III.3.4.4.3.fr}{III.3.4.4.3-fr}
\EGAReferenceAlias{III.3.4.4.3.es}{III.3.4.4.3-fr}
\EGAReferenceAlias{III.3.4.4.3.eq}{III.3.4.4.3-fr}
\EGAReferenceAlias{eq:III.3.4.4.3}{III.3.4.4.3-fr}
\EGAReferenceAlias{III.3.4.4.4}{III.3.4.4.4-fr}
\EGAReferenceAlias{III.3.4.4.4-es}{III.3.4.4.4-fr}
\EGAReferenceAlias{III.3.4.4.4.fr}{III.3.4.4.4-fr}
\EGAReferenceAlias{III.3.4.4.4.es}{III.3.4.4.4-fr}
\EGAReferenceAlias{III.3.4.4.4.eq}{III.3.4.4.4-fr}
\EGAReferenceAlias{eq:III.3.4.4.4}{III.3.4.4.4-fr}
\EGAReferenceAlias{III.3.4.5-es}{III.3.4.5}
\EGAReferenceAlias{III.3.4.5.fr}{III.3.4.5}
\EGAReferenceAlias{III.3.4.5.es}{III.3.4.5}
\EGAReferenceAlias{III.3.4.5.eq}{III.3.4.5}
\EGAReferenceAlias{eq:III.3.4.5}{III.3.4.5}
\EGAReferenceAlias{III.3.4.5.1}{III.3.4.5.1-fr}
\EGAReferenceAlias{III.3.4.5.1-es}{III.3.4.5.1-fr}
\EGAReferenceAlias{III.3.4.5.1.fr}{III.3.4.5.1-fr}
\EGAReferenceAlias{III.3.4.5.1.es}{III.3.4.5.1-fr}
\EGAReferenceAlias{III.3.4.5.1.eq}{III.3.4.5.1-fr}
\EGAReferenceAlias{eq:III.3.4.5.1}{III.3.4.5.1-fr}
\EGAReferenceAlias{III.3.4.6-es}{III.3.4.6}
\EGAReferenceAlias{III.3.4.6.fr}{III.3.4.6}
\EGAReferenceAlias{III.3.4.6.es}{III.3.4.6}
\EGAReferenceAlias{III.3.4.6.eq}{III.3.4.6}
\EGAReferenceAlias{eq:III.3.4.6}{III.3.4.6}
\EGAReferenceAlias{III.3.4.6.1}{III.3.4.6.1-fr}
\EGAReferenceAlias{III.3.4.6.1-es}{III.3.4.6.1-fr}
\EGAReferenceAlias{III.3.4.6.1.fr}{III.3.4.6.1-fr}
\EGAReferenceAlias{III.3.4.6.1.es}{III.3.4.6.1-fr}
\EGAReferenceAlias{III.3.4.6.1.eq}{III.3.4.6.1-fr}
\EGAReferenceAlias{eq:III.3.4.6.1}{III.3.4.6.1-fr}
\EGAReferenceAlias{III.4}{section:III.4}
\EGAReferenceAlias{III.4-fr}{section:III.4}
\EGAReferenceAlias{III.4-es}{section:III.4}
\EGAReferenceAlias{III.4.fr}{section:III.4}
\EGAReferenceAlias{III.4.es}{section:III.4}
\EGAReferenceAlias{III.4.eq}{section:III.4}
\EGAReferenceAlias{eq:III.4}{section:III.4}
\EGAReferenceAlias{section:III.4-es}{section:III.4}
\EGAReferenceAlias{section:III.4.fr}{section:III.4}
\EGAReferenceAlias{section:III.4.es}{section:III.4}
\EGAReferenceAlias{section:III.4.eq}{section:III.4}
\EGAReferenceAlias{eq:section:III.4}{section:III.4}
\EGAReferenceAlias{III.4.1}{subsection:III.4.1}
\EGAReferenceAlias{III.4.1-fr}{subsection:III.4.1}
\EGAReferenceAlias{III.4.1-es}{subsection:III.4.1}
\EGAReferenceAlias{III.4.1.fr}{subsection:III.4.1}
\EGAReferenceAlias{III.4.1.es}{subsection:III.4.1}
\EGAReferenceAlias{III.4.1.eq}{subsection:III.4.1}
\EGAReferenceAlias{eq:III.4.1}{subsection:III.4.1}
\EGAReferenceAlias{subsection:III.4.1-es}{subsection:III.4.1}
\EGAReferenceAlias{subsection:III.4.1.fr}{subsection:III.4.1}
\EGAReferenceAlias{subsection:III.4.1.es}{subsection:III.4.1}
\EGAReferenceAlias{subsection:III.4.1.eq}{subsection:III.4.1}
\EGAReferenceAlias{eq:subsection:III.4.1}{subsection:III.4.1}
\EGAReferenceAlias{III.4.1.1-es}{III.4.1.1}
\EGAReferenceAlias{III.4.1.1.fr}{III.4.1.1}
\EGAReferenceAlias{III.4.1.1.es}{III.4.1.1}
\EGAReferenceAlias{III.4.1.1.eq}{III.4.1.1}
\EGAReferenceAlias{eq:III.4.1.1}{III.4.1.1}
\EGAReferenceAlias{III.4.1.2-es}{III.4.1.2}
\EGAReferenceAlias{III.4.1.2.fr}{III.4.1.2}
\EGAReferenceAlias{III.4.1.2.es}{III.4.1.2}
\EGAReferenceAlias{III.4.1.2.eq}{III.4.1.2}
\EGAReferenceAlias{eq:III.4.1.2}{III.4.1.2}
\EGAReferenceAlias{III.4.1.2.1}{III.4.1.2.1-fr}
\EGAReferenceAlias{III.4.1.2.1-es}{III.4.1.2.1-fr}
\EGAReferenceAlias{III.4.1.2.1.fr}{III.4.1.2.1-fr}
\EGAReferenceAlias{III.4.1.2.1.es}{III.4.1.2.1-fr}
\EGAReferenceAlias{III.4.1.2.1.eq}{III.4.1.2.1-fr}
\EGAReferenceAlias{eq:III.4.1.2.1}{III.4.1.2.1-fr}
\EGAReferenceAlias{III.4.1.2.2}{III.4.1.2.2-fr}
\EGAReferenceAlias{III.4.1.2.2-es}{III.4.1.2.2-fr}
\EGAReferenceAlias{III.4.1.2.2.fr}{III.4.1.2.2-fr}
\EGAReferenceAlias{III.4.1.2.2.es}{III.4.1.2.2-fr}
\EGAReferenceAlias{III.4.1.2.2.eq}{III.4.1.2.2-fr}
\EGAReferenceAlias{eq:III.4.1.2.2}{III.4.1.2.2-fr}
\EGAReferenceAlias{III.4.1.2.3}{III.4.1.2.3-fr}
\EGAReferenceAlias{III.4.1.2.3-es}{III.4.1.2.3-fr}
\EGAReferenceAlias{III.4.1.2.3.fr}{III.4.1.2.3-fr}
\EGAReferenceAlias{III.4.1.2.3.es}{III.4.1.2.3-fr}
\EGAReferenceAlias{III.4.1.2.3.eq}{III.4.1.2.3-fr}
\EGAReferenceAlias{eq:III.4.1.2.3}{III.4.1.2.3-fr}
\EGAReferenceAlias{III.4.1.3-es}{III.4.1.3}
\EGAReferenceAlias{III.4.1.3.fr}{III.4.1.3}
\EGAReferenceAlias{III.4.1.3.es}{III.4.1.3}
\EGAReferenceAlias{III.4.1.3.eq}{III.4.1.3}
\EGAReferenceAlias{eq:III.4.1.3}{III.4.1.3}
\EGAReferenceAlias{III.4.1.3.1}{III.4.1.3.1-fr}
\EGAReferenceAlias{III.4.1.3.1-es}{III.4.1.3.1-fr}
\EGAReferenceAlias{III.4.1.3.1.fr}{III.4.1.3.1-fr}
\EGAReferenceAlias{III.4.1.3.1.es}{III.4.1.3.1-fr}
\EGAReferenceAlias{III.4.1.3.1.eq}{III.4.1.3.1-fr}
\EGAReferenceAlias{eq:III.4.1.3.1}{III.4.1.3.1-fr}
\EGAReferenceAlias{III.4.1.3.2}{III.4.1.3.2-fr}
\EGAReferenceAlias{III.4.1.3.2-es}{III.4.1.3.2-fr}
\EGAReferenceAlias{III.4.1.3.2.fr}{III.4.1.3.2-fr}
\EGAReferenceAlias{III.4.1.3.2.es}{III.4.1.3.2-fr}
\EGAReferenceAlias{III.4.1.3.2.eq}{III.4.1.3.2-fr}
\EGAReferenceAlias{eq:III.4.1.3.2}{III.4.1.3.2-fr}
\EGAReferenceAlias{III.4.1.3.3}{III.4.1.3.3-fr}
\EGAReferenceAlias{III.4.1.3.3-es}{III.4.1.3.3-fr}
\EGAReferenceAlias{III.4.1.3.3.fr}{III.4.1.3.3-fr}
\EGAReferenceAlias{III.4.1.3.3.es}{III.4.1.3.3-fr}
\EGAReferenceAlias{III.4.1.3.3.eq}{III.4.1.3.3-fr}
\EGAReferenceAlias{eq:III.4.1.3.3}{III.4.1.3.3-fr}
\EGAReferenceAlias{III.4.1.3.4}{III.4.1.3.4-fr}
\EGAReferenceAlias{III.4.1.3.4-es}{III.4.1.3.4-fr}
\EGAReferenceAlias{III.4.1.3.4.fr}{III.4.1.3.4-fr}
\EGAReferenceAlias{III.4.1.3.4.es}{III.4.1.3.4-fr}
\EGAReferenceAlias{III.4.1.3.4.eq}{III.4.1.3.4-fr}
\EGAReferenceAlias{eq:III.4.1.3.4}{III.4.1.3.4-fr}
\EGAReferenceAlias{III.4.1.3.5}{III.4.1.3.5-fr}
\EGAReferenceAlias{III.4.1.3.5-es}{III.4.1.3.5-fr}
\EGAReferenceAlias{III.4.1.3.5.fr}{III.4.1.3.5-fr}
\EGAReferenceAlias{III.4.1.3.5.es}{III.4.1.3.5-fr}
\EGAReferenceAlias{III.4.1.3.5.eq}{III.4.1.3.5-fr}
\EGAReferenceAlias{eq:III.4.1.3.5}{III.4.1.3.5-fr}
\EGAReferenceAlias{III.4.1.3.6}{III.4.1.3.6-fr}
\EGAReferenceAlias{III.4.1.3.6-es}{III.4.1.3.6-fr}
\EGAReferenceAlias{III.4.1.3.6.fr}{III.4.1.3.6-fr}
\EGAReferenceAlias{III.4.1.3.6.es}{III.4.1.3.6-fr}
\EGAReferenceAlias{III.4.1.3.6.eq}{III.4.1.3.6-fr}
\EGAReferenceAlias{eq:III.4.1.3.6}{III.4.1.3.6-fr}
\EGAReferenceAlias{III.4.1.4-es}{III.4.1.4}
\EGAReferenceAlias{III.4.1.4.fr}{III.4.1.4}
\EGAReferenceAlias{III.4.1.4.es}{III.4.1.4}
\EGAReferenceAlias{III.4.1.4.eq}{III.4.1.4}
\EGAReferenceAlias{eq:III.4.1.4}{III.4.1.4}
\EGAReferenceAlias{III.4.1.4.1}{III.4.1.4.1-fr}
\EGAReferenceAlias{III.4.1.4.1-es}{III.4.1.4.1-fr}
\EGAReferenceAlias{III.4.1.4.1.fr}{III.4.1.4.1-fr}
\EGAReferenceAlias{III.4.1.4.1.es}{III.4.1.4.1-fr}
\EGAReferenceAlias{III.4.1.4.1.eq}{III.4.1.4.1-fr}
\EGAReferenceAlias{eq:III.4.1.4.1}{III.4.1.4.1-fr}
\EGAReferenceAlias{III.4.1.4.2}{III.4.1.4.2-fr}
\EGAReferenceAlias{III.4.1.4.2-es}{III.4.1.4.2-fr}
\EGAReferenceAlias{III.4.1.4.2.fr}{III.4.1.4.2-fr}
\EGAReferenceAlias{III.4.1.4.2.es}{III.4.1.4.2-fr}
\EGAReferenceAlias{III.4.1.4.2.eq}{III.4.1.4.2-fr}
\EGAReferenceAlias{eq:III.4.1.4.2}{III.4.1.4.2-fr}
\EGAReferenceAlias{III.4.1.4.3}{III.4.1.4.3-fr}
\EGAReferenceAlias{III.4.1.4.3-es}{III.4.1.4.3-fr}
\EGAReferenceAlias{III.4.1.4.3.fr}{III.4.1.4.3-fr}
\EGAReferenceAlias{III.4.1.4.3.es}{III.4.1.4.3-fr}
\EGAReferenceAlias{III.4.1.4.3.eq}{III.4.1.4.3-fr}
\EGAReferenceAlias{eq:III.4.1.4.3}{III.4.1.4.3-fr}
\EGAReferenceAlias{III.4.1.5-es}{III.4.1.5}
\EGAReferenceAlias{III.4.1.5.fr}{III.4.1.5}
\EGAReferenceAlias{III.4.1.5.es}{III.4.1.5}
\EGAReferenceAlias{III.4.1.5.eq}{III.4.1.5}
\EGAReferenceAlias{eq:III.4.1.5}{III.4.1.5}
\EGAReferenceAlias{III.4.1.6}{III.4.1.6-fr}
\EGAReferenceAlias{III.4.1.6-es}{III.4.1.6-fr}
\EGAReferenceAlias{III.4.1.6.fr}{III.4.1.6-fr}
\EGAReferenceAlias{III.4.1.6.es}{III.4.1.6-fr}
\EGAReferenceAlias{III.4.1.6.eq}{III.4.1.6-fr}
\EGAReferenceAlias{eq:III.4.1.6}{III.4.1.6-fr}
\EGAReferenceAlias{III.4.1.6.1}{III.4.1.6.1-fr}
\EGAReferenceAlias{III.4.1.6.1-es}{III.4.1.6.1-fr}
\EGAReferenceAlias{III.4.1.6.1.fr}{III.4.1.6.1-fr}
\EGAReferenceAlias{III.4.1.6.1.es}{III.4.1.6.1-fr}
\EGAReferenceAlias{III.4.1.6.1.eq}{III.4.1.6.1-fr}
\EGAReferenceAlias{eq:III.4.1.6.1}{III.4.1.6.1-fr}
\EGAReferenceAlias{III.4.1.7}{III.4.1.7-fr}
\EGAReferenceAlias{III.4.1.7-es}{III.4.1.7-fr}
\EGAReferenceAlias{III.4.1.7.fr}{III.4.1.7-fr}
\EGAReferenceAlias{III.4.1.7.es}{III.4.1.7-fr}
\EGAReferenceAlias{III.4.1.7.eq}{III.4.1.7-fr}
\EGAReferenceAlias{eq:III.4.1.7}{III.4.1.7-fr}
\EGAReferenceAlias{III.4.1.7.1}{III.4.1.7.1-fr}
\EGAReferenceAlias{III.4.1.7.1-es}{III.4.1.7.1-fr}
\EGAReferenceAlias{III.4.1.7.1.fr}{III.4.1.7.1-fr}
\EGAReferenceAlias{III.4.1.7.1.es}{III.4.1.7.1-fr}
\EGAReferenceAlias{III.4.1.7.1.eq}{III.4.1.7.1-fr}
\EGAReferenceAlias{eq:III.4.1.7.1}{III.4.1.7.1-fr}
\EGAReferenceAlias{III.4.1.7.2}{III.4.1.7.2-fr}
\EGAReferenceAlias{III.4.1.7.2-es}{III.4.1.7.2-fr}
\EGAReferenceAlias{III.4.1.7.2.fr}{III.4.1.7.2-fr}
\EGAReferenceAlias{III.4.1.7.2.es}{III.4.1.7.2-fr}
\EGAReferenceAlias{III.4.1.7.2.eq}{III.4.1.7.2-fr}
\EGAReferenceAlias{eq:III.4.1.7.2}{III.4.1.7.2-fr}
\EGAReferenceAlias{III.4.1.7.3}{III.4.1.7.3-fr}
\EGAReferenceAlias{III.4.1.7.3-es}{III.4.1.7.3-fr}
\EGAReferenceAlias{III.4.1.7.3.fr}{III.4.1.7.3-fr}
\EGAReferenceAlias{III.4.1.7.3.es}{III.4.1.7.3-fr}
\EGAReferenceAlias{III.4.1.7.3.eq}{III.4.1.7.3-fr}
\EGAReferenceAlias{eq:III.4.1.7.3}{III.4.1.7.3-fr}
\EGAReferenceAlias{III.4.1.7.4}{III.4.1.7.4-fr}
\EGAReferenceAlias{III.4.1.7.4-es}{III.4.1.7.4-fr}
\EGAReferenceAlias{III.4.1.7.4.fr}{III.4.1.7.4-fr}
\EGAReferenceAlias{III.4.1.7.4.es}{III.4.1.7.4-fr}
\EGAReferenceAlias{III.4.1.7.4.eq}{III.4.1.7.4-fr}
\EGAReferenceAlias{eq:III.4.1.7.4}{III.4.1.7.4-fr}
\EGAReferenceAlias{III.4.1.7.5}{III.4.1.7.5-fr}
\EGAReferenceAlias{III.4.1.7.5-es}{III.4.1.7.5-fr}
\EGAReferenceAlias{III.4.1.7.5.fr}{III.4.1.7.5-fr}
\EGAReferenceAlias{III.4.1.7.5.es}{III.4.1.7.5-fr}
\EGAReferenceAlias{III.4.1.7.5.eq}{III.4.1.7.5-fr}
\EGAReferenceAlias{eq:III.4.1.7.5}{III.4.1.7.5-fr}
\EGAReferenceAlias{III.4.1.7.6}{III.4.1.7.6-fr}
\EGAReferenceAlias{III.4.1.7.6-es}{III.4.1.7.6-fr}
\EGAReferenceAlias{III.4.1.7.6.fr}{III.4.1.7.6-fr}
\EGAReferenceAlias{III.4.1.7.6.es}{III.4.1.7.6-fr}
\EGAReferenceAlias{III.4.1.7.6.eq}{III.4.1.7.6-fr}
\EGAReferenceAlias{eq:III.4.1.7.6}{III.4.1.7.6-fr}
\EGAReferenceAlias{III.4.1.7.7}{III.4.1.7.7-fr}
\EGAReferenceAlias{III.4.1.7.7-es}{III.4.1.7.7-fr}
\EGAReferenceAlias{III.4.1.7.7.fr}{III.4.1.7.7-fr}
\EGAReferenceAlias{III.4.1.7.7.es}{III.4.1.7.7-fr}
\EGAReferenceAlias{III.4.1.7.7.eq}{III.4.1.7.7-fr}
\EGAReferenceAlias{eq:III.4.1.7.7}{III.4.1.7.7-fr}
\EGAReferenceAlias{III.4.1.7.8}{III.4.1.7.8-fr}
\EGAReferenceAlias{III.4.1.7.8-es}{III.4.1.7.8-fr}
\EGAReferenceAlias{III.4.1.7.8.fr}{III.4.1.7.8-fr}
\EGAReferenceAlias{III.4.1.7.8.es}{III.4.1.7.8-fr}
\EGAReferenceAlias{III.4.1.7.8.eq}{III.4.1.7.8-fr}
\EGAReferenceAlias{eq:III.4.1.7.8}{III.4.1.7.8-fr}
\EGAReferenceAlias{III.4.1.7.9}{III.4.1.7.9-fr}
\EGAReferenceAlias{III.4.1.7.9-es}{III.4.1.7.9-fr}
\EGAReferenceAlias{III.4.1.7.9.fr}{III.4.1.7.9-fr}
\EGAReferenceAlias{III.4.1.7.9.es}{III.4.1.7.9-fr}
\EGAReferenceAlias{III.4.1.7.9.eq}{III.4.1.7.9-fr}
\EGAReferenceAlias{eq:III.4.1.7.9}{III.4.1.7.9-fr}
\EGAReferenceAlias{III.4.1.7.10}{III.4.1.7.10-fr}
\EGAReferenceAlias{III.4.1.7.10-es}{III.4.1.7.10-fr}
\EGAReferenceAlias{III.4.1.7.10.fr}{III.4.1.7.10-fr}
\EGAReferenceAlias{III.4.1.7.10.es}{III.4.1.7.10-fr}
\EGAReferenceAlias{III.4.1.7.10.eq}{III.4.1.7.10-fr}
\EGAReferenceAlias{eq:III.4.1.7.10}{III.4.1.7.10-fr}
\EGAReferenceAlias{III.4.1.7.11}{III.4.1.7.11-fr}
\EGAReferenceAlias{III.4.1.7.11-es}{III.4.1.7.11-fr}
\EGAReferenceAlias{III.4.1.7.11.fr}{III.4.1.7.11-fr}
\EGAReferenceAlias{III.4.1.7.11.es}{III.4.1.7.11-fr}
\EGAReferenceAlias{III.4.1.7.11.eq}{III.4.1.7.11-fr}
\EGAReferenceAlias{eq:III.4.1.7.11}{III.4.1.7.11-fr}
\EGAReferenceAlias{III.4.1.7.12}{III.4.1.7.12-fr}
\EGAReferenceAlias{III.4.1.7.12-es}{III.4.1.7.12-fr}
\EGAReferenceAlias{III.4.1.7.12.fr}{III.4.1.7.12-fr}
\EGAReferenceAlias{III.4.1.7.12.es}{III.4.1.7.12-fr}
\EGAReferenceAlias{III.4.1.7.12.eq}{III.4.1.7.12-fr}
\EGAReferenceAlias{eq:III.4.1.7.12}{III.4.1.7.12-fr}
\EGAReferenceAlias{III.4.1.7.13}{III.4.1.7.13-fr}
\EGAReferenceAlias{III.4.1.7.13-es}{III.4.1.7.13-fr}
\EGAReferenceAlias{III.4.1.7.13.fr}{III.4.1.7.13-fr}
\EGAReferenceAlias{III.4.1.7.13.es}{III.4.1.7.13-fr}
\EGAReferenceAlias{III.4.1.7.13.eq}{III.4.1.7.13-fr}
\EGAReferenceAlias{eq:III.4.1.7.13}{III.4.1.7.13-fr}
\EGAReferenceAlias{III.4.1.7.14}{III.4.1.7.14-fr}
\EGAReferenceAlias{III.4.1.7.14-es}{III.4.1.7.14-fr}
\EGAReferenceAlias{III.4.1.7.14.fr}{III.4.1.7.14-fr}
\EGAReferenceAlias{III.4.1.7.14.es}{III.4.1.7.14-fr}
\EGAReferenceAlias{III.4.1.7.14.eq}{III.4.1.7.14-fr}
\EGAReferenceAlias{eq:III.4.1.7.14}{III.4.1.7.14-fr}
\EGAReferenceAlias{III.4.1.7.15}{III.4.1.7.15-fr}
\EGAReferenceAlias{III.4.1.7.15-es}{III.4.1.7.15-fr}
\EGAReferenceAlias{III.4.1.7.15.fr}{III.4.1.7.15-fr}
\EGAReferenceAlias{III.4.1.7.15.es}{III.4.1.7.15-fr}
\EGAReferenceAlias{III.4.1.7.15.eq}{III.4.1.7.15-fr}
\EGAReferenceAlias{eq:III.4.1.7.15}{III.4.1.7.15-fr}
\EGAReferenceAlias{III.4.1.7.16}{III.4.1.7.16-fr}
\EGAReferenceAlias{III.4.1.7.16-es}{III.4.1.7.16-fr}
\EGAReferenceAlias{III.4.1.7.16.fr}{III.4.1.7.16-fr}
\EGAReferenceAlias{III.4.1.7.16.es}{III.4.1.7.16-fr}
\EGAReferenceAlias{III.4.1.7.16.eq}{III.4.1.7.16-fr}
\EGAReferenceAlias{eq:III.4.1.7.16}{III.4.1.7.16-fr}
\EGAReferenceAlias{III.4.1.7.17}{III.4.1.7.17-fr}
\EGAReferenceAlias{III.4.1.7.17-es}{III.4.1.7.17-fr}
\EGAReferenceAlias{III.4.1.7.17.fr}{III.4.1.7.17-fr}
\EGAReferenceAlias{III.4.1.7.17.es}{III.4.1.7.17-fr}
\EGAReferenceAlias{III.4.1.7.17.eq}{III.4.1.7.17-fr}
\EGAReferenceAlias{eq:III.4.1.7.17}{III.4.1.7.17-fr}
\EGAReferenceAlias{III.4.1.7.18}{III.4.1.7.18-fr}
\EGAReferenceAlias{III.4.1.7.18-es}{III.4.1.7.18-fr}
\EGAReferenceAlias{III.4.1.7.18.fr}{III.4.1.7.18-fr}
\EGAReferenceAlias{III.4.1.7.18.es}{III.4.1.7.18-fr}
\EGAReferenceAlias{III.4.1.7.18.eq}{III.4.1.7.18-fr}
\EGAReferenceAlias{eq:III.4.1.7.18}{III.4.1.7.18-fr}
\EGAReferenceAlias{III.4.1.7.19}{III.4.1.7.19-fr}
\EGAReferenceAlias{III.4.1.7.19-es}{III.4.1.7.19-fr}
\EGAReferenceAlias{III.4.1.7.19.fr}{III.4.1.7.19-fr}
\EGAReferenceAlias{III.4.1.7.19.es}{III.4.1.7.19-fr}
\EGAReferenceAlias{III.4.1.7.19.eq}{III.4.1.7.19-fr}
\EGAReferenceAlias{eq:III.4.1.7.19}{III.4.1.7.19-fr}
\EGAReferenceAlias{III.4.1.7.20}{III.4.1.7.20-fr}
\EGAReferenceAlias{III.4.1.7.20-es}{III.4.1.7.20-fr}
\EGAReferenceAlias{III.4.1.7.20.fr}{III.4.1.7.20-fr}
\EGAReferenceAlias{III.4.1.7.20.es}{III.4.1.7.20-fr}
\EGAReferenceAlias{III.4.1.7.20.eq}{III.4.1.7.20-fr}
\EGAReferenceAlias{eq:III.4.1.7.20}{III.4.1.7.20-fr}
\EGAReferenceAlias{III.4.1.7.21}{III.4.1.7.21-fr}
\EGAReferenceAlias{III.4.1.7.21-es}{III.4.1.7.21-fr}
\EGAReferenceAlias{III.4.1.7.21.fr}{III.4.1.7.21-fr}
\EGAReferenceAlias{III.4.1.7.21.es}{III.4.1.7.21-fr}
\EGAReferenceAlias{III.4.1.7.21.eq}{III.4.1.7.21-fr}
\EGAReferenceAlias{eq:III.4.1.7.21}{III.4.1.7.21-fr}
\EGAReferenceAlias{III.4.1.7.22}{III.4.1.7.22-fr}
\EGAReferenceAlias{III.4.1.7.22-es}{III.4.1.7.22-fr}
\EGAReferenceAlias{III.4.1.7.22.fr}{III.4.1.7.22-fr}
\EGAReferenceAlias{III.4.1.7.22.es}{III.4.1.7.22-fr}
\EGAReferenceAlias{III.4.1.7.22.eq}{III.4.1.7.22-fr}
\EGAReferenceAlias{eq:III.4.1.7.22}{III.4.1.7.22-fr}
\EGAReferenceAlias{III.4.1.8}{III.4.1.8-fr}
\EGAReferenceAlias{III.4.1.8-es}{III.4.1.8-fr}
\EGAReferenceAlias{III.4.1.8.fr}{III.4.1.8-fr}
\EGAReferenceAlias{III.4.1.8.es}{III.4.1.8-fr}
\EGAReferenceAlias{III.4.1.8.eq}{III.4.1.8-fr}
\EGAReferenceAlias{eq:III.4.1.8}{III.4.1.8-fr}
\EGAReferenceAlias{III.4.1.9-es}{III.4.1.9}
\EGAReferenceAlias{III.4.1.9.fr}{III.4.1.9}
\EGAReferenceAlias{III.4.1.9.es}{III.4.1.9}
\EGAReferenceAlias{III.4.1.9.eq}{III.4.1.9}
\EGAReferenceAlias{eq:III.4.1.9}{III.4.1.9}
\EGAReferenceAlias{III.4.1.10-es}{III.4.1.10}
\EGAReferenceAlias{III.4.1.10.fr}{III.4.1.10}
\EGAReferenceAlias{III.4.1.10.es}{III.4.1.10}
\EGAReferenceAlias{III.4.1.10.eq}{III.4.1.10}
\EGAReferenceAlias{eq:III.4.1.10}{III.4.1.10}
\EGAReferenceAlias{III.4.2}{subsection:III.4.2}
\EGAReferenceAlias{III.4.2-fr}{subsection:III.4.2}
\EGAReferenceAlias{III.4.2-es}{subsection:III.4.2}
\EGAReferenceAlias{III.4.2.fr}{subsection:III.4.2}
\EGAReferenceAlias{III.4.2.es}{subsection:III.4.2}
\EGAReferenceAlias{III.4.2.eq}{subsection:III.4.2}
\EGAReferenceAlias{eq:III.4.2}{subsection:III.4.2}
\EGAReferenceAlias{subsection:III.4.2-es}{subsection:III.4.2}
\EGAReferenceAlias{subsection:III.4.2.fr}{subsection:III.4.2}
\EGAReferenceAlias{subsection:III.4.2.es}{subsection:III.4.2}
\EGAReferenceAlias{subsection:III.4.2.eq}{subsection:III.4.2}
\EGAReferenceAlias{eq:subsection:III.4.2}{subsection:III.4.2}
\EGAReferenceAlias{III.4.2.1}{III.4.2.1-fr}
\EGAReferenceAlias{III.4.2.1-es}{III.4.2.1-fr}
\EGAReferenceAlias{III.4.2.1.fr}{III.4.2.1-fr}
\EGAReferenceAlias{III.4.2.1.es}{III.4.2.1-fr}
\EGAReferenceAlias{III.4.2.1.eq}{III.4.2.1-fr}
\EGAReferenceAlias{eq:III.4.2.1}{III.4.2.1-fr}
\EGAReferenceAlias{III.4.2.1.1}{III.4.2.1.1-fr}
\EGAReferenceAlias{III.4.2.1.1-es}{III.4.2.1.1-fr}
\EGAReferenceAlias{III.4.2.1.1.fr}{III.4.2.1.1-fr}
\EGAReferenceAlias{III.4.2.1.1.es}{III.4.2.1.1-fr}
\EGAReferenceAlias{III.4.2.1.1.eq}{III.4.2.1.1-fr}
\EGAReferenceAlias{eq:III.4.2.1.1}{III.4.2.1.1-fr}
\EGAReferenceAlias{III.4.2.2-es}{III.4.2.2}
\EGAReferenceAlias{III.4.2.2.fr}{III.4.2.2}
\EGAReferenceAlias{III.4.2.2.es}{III.4.2.2}
\EGAReferenceAlias{III.4.2.2.eq}{III.4.2.2}
\EGAReferenceAlias{eq:III.4.2.2}{III.4.2.2}
\EGAReferenceAlias{III.4.2.3-es}{III.4.2.3}
\EGAReferenceAlias{III.4.2.3.fr}{III.4.2.3}
\EGAReferenceAlias{III.4.2.3.es}{III.4.2.3}
\EGAReferenceAlias{III.4.2.3.eq}{III.4.2.3}
\EGAReferenceAlias{eq:III.4.2.3}{III.4.2.3}
\EGAReferenceAlias{III.4.2.4-es}{III.4.2.4}
\EGAReferenceAlias{III.4.2.4.fr}{III.4.2.4}
\EGAReferenceAlias{III.4.2.4.es}{III.4.2.4}
\EGAReferenceAlias{III.4.2.4.eq}{III.4.2.4}
\EGAReferenceAlias{eq:III.4.2.4}{III.4.2.4}
\EGAReferenceAlias{III.4.3-es}{subsection:III.4.3}
\EGAReferenceAlias{III.4.3.fr}{subsection:III.4.3}
\EGAReferenceAlias{III.4.3.es}{subsection:III.4.3}
\EGAReferenceAlias{III.4.3.eq}{subsection:III.4.3}
\EGAReferenceAlias{eq:III.4.3}{subsection:III.4.3}
\EGAReferenceAlias{subsection:III.4.3-es}{subsection:III.4.3}
\EGAReferenceAlias{subsection:III.4.3.fr}{subsection:III.4.3}
\EGAReferenceAlias{subsection:III.4.3.es}{subsection:III.4.3}
\EGAReferenceAlias{subsection:III.4.3.eq}{subsection:III.4.3}
\EGAReferenceAlias{eq:subsection:III.4.3}{subsection:III.4.3}
\EGAReferenceAlias{III.4.3.1}{III.4.3.1-fr}
\EGAReferenceAlias{III.4.3.1-es}{III.4.3.1-fr}
\EGAReferenceAlias{III.4.3.1.fr}{III.4.3.1-fr}
\EGAReferenceAlias{III.4.3.1.es}{III.4.3.1-fr}
\EGAReferenceAlias{III.4.3.1.eq}{III.4.3.1-fr}
\EGAReferenceAlias{eq:III.4.3.1}{III.4.3.1-fr}
\EGAReferenceAlias{III.4.3.2}{III.4.3.2-fr}
\EGAReferenceAlias{III.4.3.2-es}{III.4.3.2-fr}
\EGAReferenceAlias{III.4.3.2.fr}{III.4.3.2-fr}
\EGAReferenceAlias{III.4.3.2.es}{III.4.3.2-fr}
\EGAReferenceAlias{III.4.3.2.eq}{III.4.3.2-fr}
\EGAReferenceAlias{eq:III.4.3.2}{III.4.3.2-fr}
\EGAReferenceAlias{III.4.3.3-es}{III.4.3.3}
\EGAReferenceAlias{III.4.3.3.fr}{III.4.3.3}
\EGAReferenceAlias{III.4.3.3.es}{III.4.3.3}
\EGAReferenceAlias{III.4.3.3.eq}{III.4.3.3}
\EGAReferenceAlias{eq:III.4.3.3}{III.4.3.3}
\EGAReferenceAlias{III.4.3.4-es}{III.4.3.4}
\EGAReferenceAlias{III.4.3.4.fr}{III.4.3.4}
\EGAReferenceAlias{III.4.3.4.es}{III.4.3.4}
\EGAReferenceAlias{III.4.3.4.eq}{III.4.3.4}
\EGAReferenceAlias{eq:III.4.3.4}{III.4.3.4}
\EGAReferenceAlias{III.4.3.5-es}{III.4.3.5}
\EGAReferenceAlias{III.4.3.5.fr}{III.4.3.5}
\EGAReferenceAlias{III.4.3.5.es}{III.4.3.5}
\EGAReferenceAlias{III.4.3.5.eq}{III.4.3.5}
\EGAReferenceAlias{eq:III.4.3.5}{III.4.3.5}
\EGAReferenceAlias{III.4.3.6-es}{III.4.3.6}
\EGAReferenceAlias{III.4.3.6.fr}{III.4.3.6}
\EGAReferenceAlias{III.4.3.6.es}{III.4.3.6}
\EGAReferenceAlias{III.4.3.6.eq}{III.4.3.6}
\EGAReferenceAlias{eq:III.4.3.6}{III.4.3.6}
\EGAReferenceAlias{III.4.3.7-es}{III.4.3.7}
\EGAReferenceAlias{III.4.3.7.fr}{III.4.3.7}
\EGAReferenceAlias{III.4.3.7.es}{III.4.3.7}
\EGAReferenceAlias{III.4.3.7.eq}{III.4.3.7}
\EGAReferenceAlias{eq:III.4.3.7}{III.4.3.7}
\EGAReferenceAlias{III.4.3.8-es}{III.4.3.8}
\EGAReferenceAlias{III.4.3.8.fr}{III.4.3.8}
\EGAReferenceAlias{III.4.3.8.es}{III.4.3.8}
\EGAReferenceAlias{III.4.3.8.eq}{III.4.3.8}
\EGAReferenceAlias{eq:III.4.3.8}{III.4.3.8}
\EGAReferenceAlias{III.4.3.9}{III.4.3.9-fr}
\EGAReferenceAlias{III.4.3.9-es}{III.4.3.9-fr}
\EGAReferenceAlias{III.4.3.9.fr}{III.4.3.9-fr}
\EGAReferenceAlias{III.4.3.9.es}{III.4.3.9-fr}
\EGAReferenceAlias{III.4.3.9.eq}{III.4.3.9-fr}
\EGAReferenceAlias{eq:III.4.3.9}{III.4.3.9-fr}
\EGAReferenceAlias{III.4.3.10-es}{III.4.3.10}
\EGAReferenceAlias{III.4.3.10.fr}{III.4.3.10}
\EGAReferenceAlias{III.4.3.10.es}{III.4.3.10}
\EGAReferenceAlias{III.4.3.10.eq}{III.4.3.10}
\EGAReferenceAlias{eq:III.4.3.10}{III.4.3.10}
\EGAReferenceAlias{III.4.3.11}{III.4.3.11-fr}
\EGAReferenceAlias{III.4.3.11-es}{III.4.3.11-fr}
\EGAReferenceAlias{III.4.3.11.fr}{III.4.3.11-fr}
\EGAReferenceAlias{III.4.3.11.es}{III.4.3.11-fr}
\EGAReferenceAlias{III.4.3.11.eq}{III.4.3.11-fr}
\EGAReferenceAlias{eq:III.4.3.11}{III.4.3.11-fr}
\EGAReferenceAlias{III.4.3.12-es}{III.4.3.12}
\EGAReferenceAlias{III.4.3.12.fr}{III.4.3.12}
\EGAReferenceAlias{III.4.3.12.es}{III.4.3.12}
\EGAReferenceAlias{III.4.3.12.eq}{III.4.3.12}
\EGAReferenceAlias{eq:III.4.3.12}{III.4.3.12}
\EGAReferenceAlias{III.4.3.13}{III.4.3.13-fr}
\EGAReferenceAlias{III.4.3.13-es}{III.4.3.13-fr}
\EGAReferenceAlias{III.4.3.13.fr}{III.4.3.13-fr}
\EGAReferenceAlias{III.4.3.13.es}{III.4.3.13-fr}
\EGAReferenceAlias{III.4.3.13.eq}{III.4.3.13-fr}
\EGAReferenceAlias{eq:III.4.3.13}{III.4.3.13-fr}
\EGAReferenceAlias{III.4.4}{subsection:III.4.4-fr}
\EGAReferenceAlias{III.4.4-fr}{subsection:III.4.4-fr}
\EGAReferenceAlias{III.4.4-es}{subsection:III.4.4-fr}
\EGAReferenceAlias{III.4.4.fr}{subsection:III.4.4-fr}
\EGAReferenceAlias{III.4.4.es}{subsection:III.4.4-fr}
\EGAReferenceAlias{III.4.4.eq}{subsection:III.4.4-fr}
\EGAReferenceAlias{eq:III.4.4}{subsection:III.4.4-fr}
\EGAReferenceAlias{subsection:III.4.4}{subsection:III.4.4-fr}
\EGAReferenceAlias{subsection:III.4.4-es}{subsection:III.4.4-fr}
\EGAReferenceAlias{subsection:III.4.4.fr}{subsection:III.4.4-fr}
\EGAReferenceAlias{subsection:III.4.4.es}{subsection:III.4.4-fr}
\EGAReferenceAlias{subsection:III.4.4.eq}{subsection:III.4.4-fr}
\EGAReferenceAlias{eq:subsection:III.4.4}{subsection:III.4.4-fr}
\EGAReferenceAlias{III.4.4.1-es}{III.4.4.1}
\EGAReferenceAlias{III.4.4.1.fr}{III.4.4.1}
\EGAReferenceAlias{III.4.4.1.es}{III.4.4.1}
\EGAReferenceAlias{III.4.4.1.eq}{III.4.4.1}
\EGAReferenceAlias{eq:III.4.4.1}{III.4.4.1}
\EGAReferenceAlias{III.4.4.2-es}{III.4.4.2}
\EGAReferenceAlias{III.4.4.2.fr}{III.4.4.2}
\EGAReferenceAlias{III.4.4.2.es}{III.4.4.2}
\EGAReferenceAlias{III.4.4.2.eq}{III.4.4.2}
\EGAReferenceAlias{eq:III.4.4.2}{III.4.4.2}
\EGAReferenceAlias{III.4.4.3}{III.4.4.3-fr}
\EGAReferenceAlias{III.4.4.3-es}{III.4.4.3-fr}
\EGAReferenceAlias{III.4.4.3.fr}{III.4.4.3-fr}
\EGAReferenceAlias{III.4.4.3.es}{III.4.4.3-fr}
\EGAReferenceAlias{III.4.4.3.eq}{III.4.4.3-fr}
\EGAReferenceAlias{eq:III.4.4.3}{III.4.4.3-fr}
\EGAReferenceAlias{III.4.4.4}{III.4.4.4-fr}
\EGAReferenceAlias{III.4.4.4-es}{III.4.4.4-fr}
\EGAReferenceAlias{III.4.4.4.fr}{III.4.4.4-fr}
\EGAReferenceAlias{III.4.4.4.es}{III.4.4.4-fr}
\EGAReferenceAlias{III.4.4.4.eq}{III.4.4.4-fr}
\EGAReferenceAlias{eq:III.4.4.4}{III.4.4.4-fr}
\EGAReferenceAlias{III.4.4.5-es}{III.4.4.5}
\EGAReferenceAlias{III.4.4.5.fr}{III.4.4.5}
\EGAReferenceAlias{III.4.4.5.es}{III.4.4.5}
\EGAReferenceAlias{III.4.4.5.eq}{III.4.4.5}
\EGAReferenceAlias{eq:III.4.4.5}{III.4.4.5}
\EGAReferenceAlias{III.4.4.6-es}{III.4.4.6}
\EGAReferenceAlias{III.4.4.6.fr}{III.4.4.6}
\EGAReferenceAlias{III.4.4.6.es}{III.4.4.6}
\EGAReferenceAlias{III.4.4.6.eq}{III.4.4.6}
\EGAReferenceAlias{eq:III.4.4.6}{III.4.4.6}
\EGAReferenceAlias{III.4.4.7-es}{III.4.4.7}
\EGAReferenceAlias{III.4.4.7.fr}{III.4.4.7}
\EGAReferenceAlias{III.4.4.7.es}{III.4.4.7}
\EGAReferenceAlias{III.4.4.7.eq}{III.4.4.7}
\EGAReferenceAlias{eq:III.4.4.7}{III.4.4.7}
\EGAReferenceAlias{III.4.4.8-es}{III.4.4.8}
\EGAReferenceAlias{III.4.4.8.fr}{III.4.4.8}
\EGAReferenceAlias{III.4.4.8.es}{III.4.4.8}
\EGAReferenceAlias{III.4.4.8.eq}{III.4.4.8}
\EGAReferenceAlias{eq:III.4.4.8}{III.4.4.8}
\EGAReferenceAlias{III.4.4.9-es}{III.4.4.9}
\EGAReferenceAlias{III.4.4.9.fr}{III.4.4.9}
\EGAReferenceAlias{III.4.4.9.es}{III.4.4.9}
\EGAReferenceAlias{III.4.4.9.eq}{III.4.4.9}
\EGAReferenceAlias{eq:III.4.4.9}{III.4.4.9}
\EGAReferenceAlias{III.4.4.10-es}{III.4.4.10}
\EGAReferenceAlias{III.4.4.10.fr}{III.4.4.10}
\EGAReferenceAlias{III.4.4.10.es}{III.4.4.10}
\EGAReferenceAlias{III.4.4.10.eq}{III.4.4.10}
\EGAReferenceAlias{eq:III.4.4.10}{III.4.4.10}
\EGAReferenceAlias{III.4.4.11-es}{III.4.4.11}
\EGAReferenceAlias{III.4.4.11.fr}{III.4.4.11}
\EGAReferenceAlias{III.4.4.11.es}{III.4.4.11}
\EGAReferenceAlias{III.4.4.11.eq}{III.4.4.11}
\EGAReferenceAlias{eq:III.4.4.11}{III.4.4.11}
\EGAReferenceAlias{III.4.4.12}{III.4.4.12-fr}
\EGAReferenceAlias{III.4.4.12-es}{III.4.4.12-fr}
\EGAReferenceAlias{III.4.4.12.fr}{III.4.4.12-fr}
\EGAReferenceAlias{III.4.4.12.es}{III.4.4.12-fr}
\EGAReferenceAlias{III.4.4.12.eq}{III.4.4.12-fr}
\EGAReferenceAlias{eq:III.4.4.12}{III.4.4.12-fr}
\EGAReferenceAlias{III.4.5}{subsection:III.4.5-fr}
\EGAReferenceAlias{III.4.5-fr}{subsection:III.4.5-fr}
\EGAReferenceAlias{III.4.5-es}{subsection:III.4.5-fr}
\EGAReferenceAlias{III.4.5.fr}{subsection:III.4.5-fr}
\EGAReferenceAlias{III.4.5.es}{subsection:III.4.5-fr}
\EGAReferenceAlias{III.4.5.eq}{subsection:III.4.5-fr}
\EGAReferenceAlias{eq:III.4.5}{subsection:III.4.5-fr}
\EGAReferenceAlias{subsection:III.4.5}{subsection:III.4.5-fr}
\EGAReferenceAlias{subsection:III.4.5-es}{subsection:III.4.5-fr}
\EGAReferenceAlias{subsection:III.4.5.fr}{subsection:III.4.5-fr}
\EGAReferenceAlias{subsection:III.4.5.es}{subsection:III.4.5-fr}
\EGAReferenceAlias{subsection:III.4.5.eq}{subsection:III.4.5-fr}
\EGAReferenceAlias{eq:subsection:III.4.5}{subsection:III.4.5-fr}
\EGAReferenceAlias{III.4.5.1}{III.4.5.1-fr}
\EGAReferenceAlias{III.4.5.1-es}{III.4.5.1-fr}
\EGAReferenceAlias{III.4.5.1.fr}{III.4.5.1-fr}
\EGAReferenceAlias{III.4.5.1.es}{III.4.5.1-fr}
\EGAReferenceAlias{III.4.5.1.eq}{III.4.5.1-fr}
\EGAReferenceAlias{eq:III.4.5.1}{III.4.5.1-fr}
\EGAReferenceAlias{III.4.5.2-es}{III.4.5.2}
\EGAReferenceAlias{III.4.5.2.fr}{III.4.5.2}
\EGAReferenceAlias{III.4.5.2.es}{III.4.5.2}
\EGAReferenceAlias{III.4.5.2.eq}{III.4.5.2}
\EGAReferenceAlias{eq:III.4.5.2}{III.4.5.2}
\EGAReferenceAlias{III.4.5.3-es}{III.4.5.3}
\EGAReferenceAlias{III.4.5.3.fr}{III.4.5.3}
\EGAReferenceAlias{III.4.5.3.es}{III.4.5.3}
\EGAReferenceAlias{III.4.5.3.eq}{III.4.5.3}
\EGAReferenceAlias{eq:III.4.5.3}{III.4.5.3}
\EGAReferenceAlias{III.4.5.3.1}{III.4.5.3.1-fr}
\EGAReferenceAlias{III.4.5.3.1-es}{III.4.5.3.1-fr}
\EGAReferenceAlias{III.4.5.3.1.fr}{III.4.5.3.1-fr}
\EGAReferenceAlias{III.4.5.3.1.es}{III.4.5.3.1-fr}
\EGAReferenceAlias{III.4.5.3.1.eq}{III.4.5.3.1-fr}
\EGAReferenceAlias{eq:III.4.5.3.1}{III.4.5.3.1-fr}
\EGAReferenceAlias{III.4.6}{subsection:III.4.6}
\EGAReferenceAlias{III.4.6-fr}{subsection:III.4.6}
\EGAReferenceAlias{III.4.6-es}{subsection:III.4.6}
\EGAReferenceAlias{III.4.6.fr}{subsection:III.4.6}
\EGAReferenceAlias{III.4.6.es}{subsection:III.4.6}
\EGAReferenceAlias{III.4.6.eq}{subsection:III.4.6}
\EGAReferenceAlias{eq:III.4.6}{subsection:III.4.6}
\EGAReferenceAlias{subsection:III.4.6-fr}{subsection:III.4.6}
\EGAReferenceAlias{subsection:III.4.6-es}{subsection:III.4.6}
\EGAReferenceAlias{subsection:III.4.6.fr}{subsection:III.4.6}
\EGAReferenceAlias{subsection:III.4.6.es}{subsection:III.4.6}
\EGAReferenceAlias{subsection:III.4.6.eq}{subsection:III.4.6}
\EGAReferenceAlias{eq:subsection:III.4.6}{subsection:III.4.6}
\EGAReferenceAlias{III.4.6.1}{III.4.6.1-fr}
\EGAReferenceAlias{III.4.6.1-es}{III.4.6.1-fr}
\EGAReferenceAlias{III.4.6.1.fr}{III.4.6.1-fr}
\EGAReferenceAlias{III.4.6.1.es}{III.4.6.1-fr}
\EGAReferenceAlias{III.4.6.1.eq}{III.4.6.1-fr}
\EGAReferenceAlias{eq:III.4.6.1}{III.4.6.1-fr}
\EGAReferenceAlias{III.4.6.1.1}{III.4.6.1.1-fr}
\EGAReferenceAlias{III.4.6.1.1-es}{III.4.6.1.1-fr}
\EGAReferenceAlias{III.4.6.1.1.fr}{III.4.6.1.1-fr}
\EGAReferenceAlias{III.4.6.1.1.es}{III.4.6.1.1-fr}
\EGAReferenceAlias{III.4.6.1.1.eq}{III.4.6.1.1-fr}
\EGAReferenceAlias{eq:III.4.6.1.1}{III.4.6.1.1-fr}
\EGAReferenceAlias{III.4.6.1.2}{III.4.6.1.2-fr}
\EGAReferenceAlias{III.4.6.1.2-es}{III.4.6.1.2-fr}
\EGAReferenceAlias{III.4.6.1.2.fr}{III.4.6.1.2-fr}
\EGAReferenceAlias{III.4.6.1.2.es}{III.4.6.1.2-fr}
\EGAReferenceAlias{III.4.6.1.2.eq}{III.4.6.1.2-fr}
\EGAReferenceAlias{eq:III.4.6.1.2}{III.4.6.1.2-fr}
\EGAReferenceAlias{III.4.6.1.3}{III.4.6.1.3-fr}
\EGAReferenceAlias{III.4.6.1.3-es}{III.4.6.1.3-fr}
\EGAReferenceAlias{III.4.6.1.3.fr}{III.4.6.1.3-fr}
\EGAReferenceAlias{III.4.6.1.3.es}{III.4.6.1.3-fr}
\EGAReferenceAlias{III.4.6.1.3.eq}{III.4.6.1.3-fr}
\EGAReferenceAlias{eq:III.4.6.1.3}{III.4.6.1.3-fr}
\EGAReferenceAlias{III.4.6.1.4}{III.4.6.1.4-fr}
\EGAReferenceAlias{III.4.6.1.4-es}{III.4.6.1.4-fr}
\EGAReferenceAlias{III.4.6.1.4.fr}{III.4.6.1.4-fr}
\EGAReferenceAlias{III.4.6.1.4.es}{III.4.6.1.4-fr}
\EGAReferenceAlias{III.4.6.1.4.eq}{III.4.6.1.4-fr}
\EGAReferenceAlias{eq:III.4.6.1.4}{III.4.6.1.4-fr}
\EGAReferenceAlias{III.4.6.2}{III.4.6.2-fr}
\EGAReferenceAlias{III.4.6.2-es}{III.4.6.2-fr}
\EGAReferenceAlias{III.4.6.2.fr}{III.4.6.2-fr}
\EGAReferenceAlias{III.4.6.2.es}{III.4.6.2-fr}
\EGAReferenceAlias{III.4.6.2.eq}{III.4.6.2-fr}
\EGAReferenceAlias{eq:III.4.6.2}{III.4.6.2-fr}
\EGAReferenceAlias{III.4.6.2.1}{III.4.6.2.1-fr}
\EGAReferenceAlias{III.4.6.2.1-es}{III.4.6.2.1-fr}
\EGAReferenceAlias{III.4.6.2.1.fr}{III.4.6.2.1-fr}
\EGAReferenceAlias{III.4.6.2.1.es}{III.4.6.2.1-fr}
\EGAReferenceAlias{III.4.6.2.1.eq}{III.4.6.2.1-fr}
\EGAReferenceAlias{eq:III.4.6.2.1}{III.4.6.2.1-fr}
\EGAReferenceAlias{III.4.6.3-es}{III.4.6.3}
\EGAReferenceAlias{III.4.6.3.fr}{III.4.6.3}
\EGAReferenceAlias{III.4.6.3.es}{III.4.6.3}
\EGAReferenceAlias{III.4.6.3.eq}{III.4.6.3}
\EGAReferenceAlias{eq:III.4.6.3}{III.4.6.3}
\EGAReferenceAlias{III.4.6.4-es}{III.4.6.4}
\EGAReferenceAlias{III.4.6.4.fr}{III.4.6.4}
\EGAReferenceAlias{III.4.6.4.es}{III.4.6.4}
\EGAReferenceAlias{III.4.6.4.eq}{III.4.6.4}
\EGAReferenceAlias{eq:III.4.6.4}{III.4.6.4}
\EGAReferenceAlias{III.4.6.4.1}{III.4.6.4.1-fr}
\EGAReferenceAlias{III.4.6.4.1-es}{III.4.6.4.1-fr}
\EGAReferenceAlias{III.4.6.4.1.fr}{III.4.6.4.1-fr}
\EGAReferenceAlias{III.4.6.4.1.es}{III.4.6.4.1-fr}
\EGAReferenceAlias{III.4.6.4.1.eq}{III.4.6.4.1-fr}
\EGAReferenceAlias{eq:III.4.6.4.1}{III.4.6.4.1-fr}
\EGAReferenceAlias{III.4.6.5-es}{III.4.6.5}
\EGAReferenceAlias{III.4.6.5.fr}{III.4.6.5}
\EGAReferenceAlias{III.4.6.5.es}{III.4.6.5}
\EGAReferenceAlias{III.4.6.5.eq}{III.4.6.5}
\EGAReferenceAlias{eq:III.4.6.5}{III.4.6.5}
\EGAReferenceAlias{III.4.6.6-es}{III.4.6.6}
\EGAReferenceAlias{III.4.6.6.fr}{III.4.6.6}
\EGAReferenceAlias{III.4.6.6.es}{III.4.6.6}
\EGAReferenceAlias{III.4.6.6.eq}{III.4.6.6}
\EGAReferenceAlias{eq:III.4.6.6}{III.4.6.6}
\EGAReferenceAlias{III.4.6.7}{III.4.6.7-fr}
\EGAReferenceAlias{III.4.6.7-es}{III.4.6.7-fr}
\EGAReferenceAlias{III.4.6.7.fr}{III.4.6.7-fr}
\EGAReferenceAlias{III.4.6.7.es}{III.4.6.7-fr}
\EGAReferenceAlias{III.4.6.7.eq}{III.4.6.7-fr}
\EGAReferenceAlias{eq:III.4.6.7}{III.4.6.7-fr}
\EGAReferenceAlias{III.4.6.7.1}{III.4.6.7.1-fr}
\EGAReferenceAlias{III.4.6.7.1-es}{III.4.6.7.1-fr}
\EGAReferenceAlias{III.4.6.7.1.fr}{III.4.6.7.1-fr}
\EGAReferenceAlias{III.4.6.7.1.es}{III.4.6.7.1-fr}
\EGAReferenceAlias{III.4.6.7.1.eq}{III.4.6.7.1-fr}
\EGAReferenceAlias{eq:III.4.6.7.1}{III.4.6.7.1-fr}
\EGAReferenceAlias{III.4.6.8}{III.4.6.8-fr}
\EGAReferenceAlias{III.4.6.8-es}{III.4.6.8-fr}
\EGAReferenceAlias{III.4.6.8.fr}{III.4.6.8-fr}
\EGAReferenceAlias{III.4.6.8.es}{III.4.6.8-fr}
\EGAReferenceAlias{III.4.6.8.eq}{III.4.6.8-fr}
\EGAReferenceAlias{eq:III.4.6.8}{III.4.6.8-fr}
\EGAReferenceAlias{III.4.7.1}{III.4.7.1-fr}
\EGAReferenceAlias{III.4.7.1-es}{III.4.7.1-fr}
\EGAReferenceAlias{III.4.7.1.fr}{III.4.7.1-fr}
\EGAReferenceAlias{III.4.7.1.es}{III.4.7.1-fr}
\EGAReferenceAlias{III.4.7.1.eq}{III.4.7.1-fr}
\EGAReferenceAlias{eq:III.4.7.1}{III.4.7.1-fr}
\EGAReferenceAlias{III.4.8.1-es}{III.4.8.1}
\EGAReferenceAlias{III.4.8.1.fr}{III.4.8.1}
\EGAReferenceAlias{III.4.8.1.es}{III.4.8.1}
\EGAReferenceAlias{III.4.8.1.eq}{III.4.8.1}
\EGAReferenceAlias{eq:III.4.8.1}{III.4.8.1}
\EGAReferenceAlias{III.4.8.2-es}{III.4.8.2}
\EGAReferenceAlias{III.4.8.2.fr}{III.4.8.2}
\EGAReferenceAlias{III.4.8.2.es}{III.4.8.2}
\EGAReferenceAlias{III.4.8.2.eq}{III.4.8.2}
\EGAReferenceAlias{eq:III.4.8.2}{III.4.8.2}
\EGAReferenceAlias{III.4.8.3-es}{III.4.8.3}
\EGAReferenceAlias{III.4.8.3.fr}{III.4.8.3}
\EGAReferenceAlias{III.4.8.3.es}{III.4.8.3}
\EGAReferenceAlias{III.4.8.3.eq}{III.4.8.3}
\EGAReferenceAlias{eq:III.4.8.3}{III.4.8.3}
\EGAReferenceAlias{III.4.8.4-es}{III.4.8.4}
\EGAReferenceAlias{III.4.8.4.fr}{III.4.8.4}
\EGAReferenceAlias{III.4.8.4.es}{III.4.8.4}
\EGAReferenceAlias{III.4.8.4.eq}{III.4.8.4}
\EGAReferenceAlias{eq:III.4.8.4}{III.4.8.4}
\EGAReferenceAlias{III.4.8.5-es}{III.4.8.5}
\EGAReferenceAlias{III.4.8.5.fr}{III.4.8.5}
\EGAReferenceAlias{III.4.8.5.es}{III.4.8.5}
\EGAReferenceAlias{III.4.8.5.eq}{III.4.8.5}
\EGAReferenceAlias{eq:III.4.8.5}{III.4.8.5}
\EGAReferenceAlias{III.4.8.6-es}{III.4.8.6}
\EGAReferenceAlias{III.4.8.6.fr}{III.4.8.6}
\EGAReferenceAlias{III.4.8.6.es}{III.4.8.6}
\EGAReferenceAlias{III.4.8.6.eq}{III.4.8.6}
\EGAReferenceAlias{eq:III.4.8.6}{III.4.8.6}
\EGAReferenceAlias{III.4.8.7-es}{III.4.8.7}
\EGAReferenceAlias{III.4.8.7.fr}{III.4.8.7}
\EGAReferenceAlias{III.4.8.7.es}{III.4.8.7}
\EGAReferenceAlias{III.4.8.7.eq}{III.4.8.7}
\EGAReferenceAlias{eq:III.4.8.7}{III.4.8.7}
\EGAReferenceAlias{III.4.8.8-es}{III.4.8.8}
\EGAReferenceAlias{III.4.8.8.fr}{III.4.8.8}
\EGAReferenceAlias{III.4.8.8.es}{III.4.8.8}
\EGAReferenceAlias{III.4.8.8.eq}{III.4.8.8}
\EGAReferenceAlias{eq:III.4.8.8}{III.4.8.8}
\EGAReferenceAlias{III.4.8.9-es}{III.4.8.9}
\EGAReferenceAlias{III.4.8.9.fr}{III.4.8.9}
\EGAReferenceAlias{III.4.8.9.es}{III.4.8.9}
\EGAReferenceAlias{III.4.8.9.eq}{III.4.8.9}
\EGAReferenceAlias{eq:III.4.8.9}{III.4.8.9}
\EGAReferenceAlias{III.4.8.10-es}{III.4.8.10}
\EGAReferenceAlias{III.4.8.10.fr}{III.4.8.10}
\EGAReferenceAlias{III.4.8.10.es}{III.4.8.10}
\EGAReferenceAlias{III.4.8.10.eq}{III.4.8.10}
\EGAReferenceAlias{eq:III.4.8.10}{III.4.8.10}
\EGAReferenceAlias{III.4.8.11-es}{III.4.8.11}
\EGAReferenceAlias{III.4.8.11.fr}{III.4.8.11}
\EGAReferenceAlias{III.4.8.11.es}{III.4.8.11}
\EGAReferenceAlias{III.4.8.11.eq}{III.4.8.11}
\EGAReferenceAlias{eq:III.4.8.11}{III.4.8.11}
\EGAReferenceAlias{III.5}{section:III.5}
\EGAReferenceAlias{III.5-fr}{section:III.5}
\EGAReferenceAlias{III.5-es}{section:III.5}
\EGAReferenceAlias{III.5.fr}{section:III.5}
\EGAReferenceAlias{III.5.es}{section:III.5}
\EGAReferenceAlias{III.5.eq}{section:III.5}
\EGAReferenceAlias{eq:III.5}{section:III.5}
\EGAReferenceAlias{section:III.5-es}{section:III.5}
\EGAReferenceAlias{section:III.5.fr}{section:III.5}
\EGAReferenceAlias{section:III.5.es}{section:III.5}
\EGAReferenceAlias{section:III.5.eq}{section:III.5}
\EGAReferenceAlias{eq:section:III.5}{section:III.5}
\EGAReferenceAlias{III.5.1}{subsection:III.5.1}
\EGAReferenceAlias{III.5.1-fr}{subsection:III.5.1}
\EGAReferenceAlias{III.5.1-es}{subsection:III.5.1}
\EGAReferenceAlias{III.5.1.fr}{subsection:III.5.1}
\EGAReferenceAlias{III.5.1.es}{subsection:III.5.1}
\EGAReferenceAlias{III.5.1.eq}{subsection:III.5.1}
\EGAReferenceAlias{eq:III.5.1}{subsection:III.5.1}
\EGAReferenceAlias{subsection:III.5.1-es}{subsection:III.5.1}
\EGAReferenceAlias{subsection:III.5.1.fr}{subsection:III.5.1}
\EGAReferenceAlias{subsection:III.5.1.es}{subsection:III.5.1}
\EGAReferenceAlias{subsection:III.5.1.eq}{subsection:III.5.1}
\EGAReferenceAlias{eq:subsection:III.5.1}{subsection:III.5.1}
\EGAReferenceAlias{III.5.1.1-es}{III.5.1.1}
\EGAReferenceAlias{III.5.1.1.fr}{III.5.1.1}
\EGAReferenceAlias{III.5.1.1.es}{III.5.1.1}
\EGAReferenceAlias{III.5.1.1.eq}{III.5.1.1}
\EGAReferenceAlias{eq:III.5.1.1}{III.5.1.1}
\EGAReferenceAlias{III.5.1.2-es}{III.5.1.2}
\EGAReferenceAlias{III.5.1.2.fr}{III.5.1.2}
\EGAReferenceAlias{III.5.1.2.es}{III.5.1.2}
\EGAReferenceAlias{III.5.1.2.eq}{III.5.1.2}
\EGAReferenceAlias{eq:III.5.1.2}{III.5.1.2}
\EGAReferenceAlias{III.5.1.3}{III.5.1.3-fr}
\EGAReferenceAlias{III.5.1.3-es}{III.5.1.3-fr}
\EGAReferenceAlias{III.5.1.3.fr}{III.5.1.3-fr}
\EGAReferenceAlias{III.5.1.3.es}{III.5.1.3-fr}
\EGAReferenceAlias{III.5.1.3.eq}{III.5.1.3-fr}
\EGAReferenceAlias{eq:III.5.1.3}{III.5.1.3-fr}
\EGAReferenceAlias{III.5.1.3.1}{III.5.1.3.1.eq-fr}
\EGAReferenceAlias{III.5.1.3.1-es}{III.5.1.3.1.eq-fr}
\EGAReferenceAlias{III.5.1.3.1.fr}{III.5.1.3.1.eq-fr}
\EGAReferenceAlias{III.5.1.3.1.es}{III.5.1.3.1.eq-fr}
\EGAReferenceAlias{III.5.1.3.1.eq}{III.5.1.3.1.eq-fr}
\EGAReferenceAlias{eq:III.5.1.3.1}{III.5.1.3.1.eq-fr}
\EGAReferenceAlias{III.5.1.3.1.eq-es}{III.5.1.3.1.eq-fr}
\EGAReferenceAlias{III.5.1.3.1.eq.fr}{III.5.1.3.1.eq-fr}
\EGAReferenceAlias{III.5.1.3.1.eq.es}{III.5.1.3.1.eq-fr}
\EGAReferenceAlias{III.5.1.3.1.eq.eq}{III.5.1.3.1.eq-fr}
\EGAReferenceAlias{eq:III.5.1.3.1.eq}{III.5.1.3.1.eq-fr}
\EGAReferenceAlias{III.5.1.4-es}{III.5.1.4}
\EGAReferenceAlias{III.5.1.4.fr}{III.5.1.4}
\EGAReferenceAlias{III.5.1.4.es}{III.5.1.4}
\EGAReferenceAlias{III.5.1.4.eq}{III.5.1.4}
\EGAReferenceAlias{eq:III.5.1.4}{III.5.1.4}
\EGAReferenceAlias{III.5.1.5-es}{III.5.1.5}
\EGAReferenceAlias{III.5.1.5.fr}{III.5.1.5}
\EGAReferenceAlias{III.5.1.5.es}{III.5.1.5}
\EGAReferenceAlias{III.5.1.5.eq}{III.5.1.5}
\EGAReferenceAlias{eq:III.5.1.5}{III.5.1.5}
\EGAReferenceAlias{III.5.1.6-es}{III.5.1.6}
\EGAReferenceAlias{III.5.1.6.fr}{III.5.1.6}
\EGAReferenceAlias{III.5.1.6.es}{III.5.1.6}
\EGAReferenceAlias{III.5.1.6.eq}{III.5.1.6}
\EGAReferenceAlias{eq:III.5.1.6}{III.5.1.6}
\EGAReferenceAlias{III.5.1.7-es}{III.5.1.7}
\EGAReferenceAlias{III.5.1.7.fr}{III.5.1.7}
\EGAReferenceAlias{III.5.1.7.es}{III.5.1.7}
\EGAReferenceAlias{III.5.1.7.eq}{III.5.1.7}
\EGAReferenceAlias{eq:III.5.1.7}{III.5.1.7}
\EGAReferenceAlias{III.5.1.8-es}{III.5.1.8}
\EGAReferenceAlias{III.5.1.8.fr}{III.5.1.8}
\EGAReferenceAlias{III.5.1.8.es}{III.5.1.8}
\EGAReferenceAlias{III.5.1.8.eq}{III.5.1.8}
\EGAReferenceAlias{eq:III.5.1.8}{III.5.1.8}
\EGAReferenceAlias{III.5.2}{subsection:III.5.2}
\EGAReferenceAlias{III.5.2-fr}{subsection:III.5.2}
\EGAReferenceAlias{III.5.2-es}{subsection:III.5.2}
\EGAReferenceAlias{III.5.2.fr}{subsection:III.5.2}
\EGAReferenceAlias{III.5.2.es}{subsection:III.5.2}
\EGAReferenceAlias{III.5.2.eq}{subsection:III.5.2}
\EGAReferenceAlias{eq:III.5.2}{subsection:III.5.2}
\EGAReferenceAlias{subsection:III.5.2-es}{subsection:III.5.2}
\EGAReferenceAlias{subsection:III.5.2.fr}{subsection:III.5.2}
\EGAReferenceAlias{subsection:III.5.2.es}{subsection:III.5.2}
\EGAReferenceAlias{subsection:III.5.2.eq}{subsection:III.5.2}
\EGAReferenceAlias{eq:subsection:III.5.2}{subsection:III.5.2}
\EGAReferenceAlias{III.5.2.1-es}{III.5.2.1}
\EGAReferenceAlias{III.5.2.1.fr}{III.5.2.1}
\EGAReferenceAlias{III.5.2.1.es}{III.5.2.1}
\EGAReferenceAlias{III.5.2.1.eq}{III.5.2.1}
\EGAReferenceAlias{eq:III.5.2.1}{III.5.2.1}
\EGAReferenceAlias{III.5.2.2-es}{III.5.2.2}
\EGAReferenceAlias{III.5.2.2.fr}{III.5.2.2}
\EGAReferenceAlias{III.5.2.2.es}{III.5.2.2}
\EGAReferenceAlias{III.5.2.2.eq}{III.5.2.2}
\EGAReferenceAlias{eq:III.5.2.2}{III.5.2.2}
\EGAReferenceAlias{III.5.2.3}{III.5.2.3-fr}
\EGAReferenceAlias{III.5.2.3-es}{III.5.2.3-fr}
\EGAReferenceAlias{III.5.2.3.fr}{III.5.2.3-fr}
\EGAReferenceAlias{III.5.2.3.es}{III.5.2.3-fr}
\EGAReferenceAlias{III.5.2.3.eq}{III.5.2.3-fr}
\EGAReferenceAlias{eq:III.5.2.3}{III.5.2.3-fr}
\EGAReferenceAlias{III.5.2.3.1}{III.5.2.3.1-fr}
\EGAReferenceAlias{III.5.2.3.1-es}{III.5.2.3.1-fr}
\EGAReferenceAlias{III.5.2.3.1.fr}{III.5.2.3.1-fr}
\EGAReferenceAlias{III.5.2.3.1.es}{III.5.2.3.1-fr}
\EGAReferenceAlias{III.5.2.3.1.eq}{III.5.2.3.1-fr}
\EGAReferenceAlias{eq:III.5.2.3.1}{III.5.2.3.1-fr}
\EGAReferenceAlias{III.5.2.3.2}{III.5.2.3.2-fr}
\EGAReferenceAlias{III.5.2.3.2-es}{III.5.2.3.2-fr}
\EGAReferenceAlias{III.5.2.3.2.fr}{III.5.2.3.2-fr}
\EGAReferenceAlias{III.5.2.3.2.es}{III.5.2.3.2-fr}
\EGAReferenceAlias{III.5.2.3.2.eq}{III.5.2.3.2-fr}
\EGAReferenceAlias{eq:III.5.2.3.2}{III.5.2.3.2-fr}
\EGAReferenceAlias{III.5.2.3.3}{III.5.2.3.3-fr}
\EGAReferenceAlias{III.5.2.3.3-es}{III.5.2.3.3-fr}
\EGAReferenceAlias{III.5.2.3.3.fr}{III.5.2.3.3-fr}
\EGAReferenceAlias{III.5.2.3.3.es}{III.5.2.3.3-fr}
\EGAReferenceAlias{III.5.2.3.3.eq}{III.5.2.3.3-fr}
\EGAReferenceAlias{eq:III.5.2.3.3}{III.5.2.3.3-fr}
\EGAReferenceAlias{III.5.2.3.4}{III.5.2.3.4-fr}
\EGAReferenceAlias{III.5.2.3.4-es}{III.5.2.3.4-fr}
\EGAReferenceAlias{III.5.2.3.4.fr}{III.5.2.3.4-fr}
\EGAReferenceAlias{III.5.2.3.4.es}{III.5.2.3.4-fr}
\EGAReferenceAlias{III.5.2.3.4.eq}{III.5.2.3.4-fr}
\EGAReferenceAlias{eq:III.5.2.3.4}{III.5.2.3.4-fr}
\EGAReferenceAlias{III.5.2.3.5}{III.5.2.3.5-fr}
\EGAReferenceAlias{III.5.2.3.5-es}{III.5.2.3.5-fr}
\EGAReferenceAlias{III.5.2.3.5.fr}{III.5.2.3.5-fr}
\EGAReferenceAlias{III.5.2.3.5.es}{III.5.2.3.5-fr}
\EGAReferenceAlias{III.5.2.3.5.eq}{III.5.2.3.5-fr}
\EGAReferenceAlias{eq:III.5.2.3.5}{III.5.2.3.5-fr}
\EGAReferenceAlias{III.5.2.4}{III.5.2.4-fr}
\EGAReferenceAlias{III.5.2.4-es}{III.5.2.4-fr}
\EGAReferenceAlias{III.5.2.4.fr}{III.5.2.4-fr}
\EGAReferenceAlias{III.5.2.4.es}{III.5.2.4-fr}
\EGAReferenceAlias{III.5.2.4.eq}{III.5.2.4-fr}
\EGAReferenceAlias{eq:III.5.2.4}{III.5.2.4-fr}
\EGAReferenceAlias{III.5.2.5-es}{III.5.2.5}
\EGAReferenceAlias{III.5.2.5.fr}{III.5.2.5}
\EGAReferenceAlias{III.5.2.5.es}{III.5.2.5}
\EGAReferenceAlias{III.5.2.5.eq}{III.5.2.5}
\EGAReferenceAlias{eq:III.5.2.5}{III.5.2.5}
\EGAReferenceAlias{III.5.2.6-es}{III.5.2.6}
\EGAReferenceAlias{III.5.2.6.fr}{III.5.2.6}
\EGAReferenceAlias{III.5.2.6.es}{III.5.2.6}
\EGAReferenceAlias{III.5.2.6.eq}{III.5.2.6}
\EGAReferenceAlias{eq:III.5.2.6}{III.5.2.6}
\EGAReferenceAlias{III.5.3}{subsection:III.5.3}
\EGAReferenceAlias{III.5.3-fr}{subsection:III.5.3}
\EGAReferenceAlias{III.5.3-es}{subsection:III.5.3}
\EGAReferenceAlias{III.5.3.fr}{subsection:III.5.3}
\EGAReferenceAlias{III.5.3.es}{subsection:III.5.3}
\EGAReferenceAlias{III.5.3.eq}{subsection:III.5.3}
\EGAReferenceAlias{eq:III.5.3}{subsection:III.5.3}
\EGAReferenceAlias{subsection:III.5.3-es}{subsection:III.5.3}
\EGAReferenceAlias{subsection:III.5.3.fr}{subsection:III.5.3}
\EGAReferenceAlias{subsection:III.5.3.es}{subsection:III.5.3}
\EGAReferenceAlias{subsection:III.5.3.eq}{subsection:III.5.3}
\EGAReferenceAlias{eq:subsection:III.5.3}{subsection:III.5.3}
\EGAReferenceAlias{III.5.3.1-es}{III.5.3.1}
\EGAReferenceAlias{III.5.3.1.fr}{III.5.3.1}
\EGAReferenceAlias{III.5.3.1.es}{III.5.3.1}
\EGAReferenceAlias{III.5.3.1.eq}{III.5.3.1}
\EGAReferenceAlias{eq:III.5.3.1}{III.5.3.1}
\EGAReferenceAlias{III.5.3.2-es}{III.5.3.2}
\EGAReferenceAlias{III.5.3.2.fr}{III.5.3.2}
\EGAReferenceAlias{III.5.3.2.es}{III.5.3.2}
\EGAReferenceAlias{III.5.3.2.eq}{III.5.3.2}
\EGAReferenceAlias{eq:III.5.3.2}{III.5.3.2}
\EGAReferenceAlias{III.5.3.3-es}{III.5.3.3}
\EGAReferenceAlias{III.5.3.3.fr}{III.5.3.3}
\EGAReferenceAlias{III.5.3.3.es}{III.5.3.3}
\EGAReferenceAlias{III.5.3.3.eq}{III.5.3.3}
\EGAReferenceAlias{eq:III.5.3.3}{III.5.3.3}
\EGAReferenceAlias{III.5.3.4}{III.5.3.4-fr}
\EGAReferenceAlias{III.5.3.4-es}{III.5.3.4-fr}
\EGAReferenceAlias{III.5.3.4.fr}{III.5.3.4-fr}
\EGAReferenceAlias{III.5.3.4.es}{III.5.3.4-fr}
\EGAReferenceAlias{III.5.3.4.eq}{III.5.3.4-fr}
\EGAReferenceAlias{eq:III.5.3.4}{III.5.3.4-fr}
\EGAReferenceAlias{III.5.3.5}{III.5.3.5-fr}
\EGAReferenceAlias{III.5.3.5-es}{III.5.3.5-fr}
\EGAReferenceAlias{III.5.3.5.fr}{III.5.3.5-fr}
\EGAReferenceAlias{III.5.3.5.es}{III.5.3.5-fr}
\EGAReferenceAlias{III.5.3.5.eq}{III.5.3.5-fr}
\EGAReferenceAlias{eq:III.5.3.5}{III.5.3.5-fr}
\EGAReferenceAlias{III.5.4}{subsection:III.5.4}
\EGAReferenceAlias{III.5.4-fr}{subsection:III.5.4}
\EGAReferenceAlias{III.5.4-es}{subsection:III.5.4}
\EGAReferenceAlias{III.5.4.fr}{subsection:III.5.4}
\EGAReferenceAlias{III.5.4.es}{subsection:III.5.4}
\EGAReferenceAlias{III.5.4.eq}{subsection:III.5.4}
\EGAReferenceAlias{eq:III.5.4}{subsection:III.5.4}
\EGAReferenceAlias{subsection:III.5.4-es}{subsection:III.5.4}
\EGAReferenceAlias{subsection:III.5.4.fr}{subsection:III.5.4}
\EGAReferenceAlias{subsection:III.5.4.es}{subsection:III.5.4}
\EGAReferenceAlias{subsection:III.5.4.eq}{subsection:III.5.4}
\EGAReferenceAlias{eq:subsection:III.5.4}{subsection:III.5.4}
\EGAReferenceAlias{III.5.4.1}{III.5.4.1-fr}
\EGAReferenceAlias{III.5.4.1-es}{III.5.4.1-fr}
\EGAReferenceAlias{III.5.4.1.fr}{III.5.4.1-fr}
\EGAReferenceAlias{III.5.4.1.es}{III.5.4.1-fr}
\EGAReferenceAlias{III.5.4.1.eq}{III.5.4.1-fr}
\EGAReferenceAlias{eq:III.5.4.1}{III.5.4.1-fr}
\EGAReferenceAlias{III.5.4.2}{III.5.4.2-fr}
\EGAReferenceAlias{III.5.4.2-es}{III.5.4.2-fr}
\EGAReferenceAlias{III.5.4.2.fr}{III.5.4.2-fr}
\EGAReferenceAlias{III.5.4.2.es}{III.5.4.2-fr}
\EGAReferenceAlias{III.5.4.2.eq}{III.5.4.2-fr}
\EGAReferenceAlias{eq:III.5.4.2}{III.5.4.2-fr}
\EGAReferenceAlias{III.5.4.3-es}{III.5.4.3}
\EGAReferenceAlias{III.5.4.3.fr}{III.5.4.3}
\EGAReferenceAlias{III.5.4.3.es}{III.5.4.3}
\EGAReferenceAlias{III.5.4.3.eq}{III.5.4.3}
\EGAReferenceAlias{eq:III.5.4.3}{III.5.4.3}
\EGAReferenceAlias{III.5.4.4}{III.5.4.4-fr}
\EGAReferenceAlias{III.5.4.4-es}{III.5.4.4-fr}
\EGAReferenceAlias{III.5.4.4.fr}{III.5.4.4-fr}
\EGAReferenceAlias{III.5.4.4.es}{III.5.4.4-fr}
\EGAReferenceAlias{III.5.4.4.eq}{III.5.4.4-fr}
\EGAReferenceAlias{eq:III.5.4.4}{III.5.4.4-fr}
\EGAReferenceAlias{III.5.4.5}{III.5.4.5-fr}
\EGAReferenceAlias{III.5.4.5-es}{III.5.4.5-fr}
\EGAReferenceAlias{III.5.4.5.fr}{III.5.4.5-fr}
\EGAReferenceAlias{III.5.4.5.es}{III.5.4.5-fr}
\EGAReferenceAlias{III.5.4.5.eq}{III.5.4.5-fr}
\EGAReferenceAlias{eq:III.5.4.5}{III.5.4.5-fr}
\EGAReferenceAlias{III.5.4.6-es}{III.5.4.6}
\EGAReferenceAlias{III.5.4.6.fr}{III.5.4.6}
\EGAReferenceAlias{III.5.4.6.es}{III.5.4.6}
\EGAReferenceAlias{III.5.4.6.eq}{III.5.4.6}
\EGAReferenceAlias{eq:III.5.4.6}{III.5.4.6}
\EGAReferenceAlias{III.5.5}{subsection:III.5.5}
\EGAReferenceAlias{III.5.5-fr}{subsection:III.5.5}
\EGAReferenceAlias{III.5.5-es}{subsection:III.5.5}
\EGAReferenceAlias{III.5.5.fr}{subsection:III.5.5}
\EGAReferenceAlias{III.5.5.es}{subsection:III.5.5}
\EGAReferenceAlias{III.5.5.eq}{subsection:III.5.5}
\EGAReferenceAlias{eq:III.5.5}{subsection:III.5.5}
\EGAReferenceAlias{subsection:III.5.5-es}{subsection:III.5.5}
\EGAReferenceAlias{subsection:III.5.5.fr}{subsection:III.5.5}
\EGAReferenceAlias{subsection:III.5.5.es}{subsection:III.5.5}
\EGAReferenceAlias{subsection:III.5.5.eq}{subsection:III.5.5}
\EGAReferenceAlias{eq:subsection:III.5.5}{subsection:III.5.5}
\EGAReferenceAlias{III.5.5.1}{III.5.5.1-fr}
\EGAReferenceAlias{III.5.5.1-es}{III.5.5.1-fr}
\EGAReferenceAlias{III.5.5.1.fr}{III.5.5.1-fr}
\EGAReferenceAlias{III.5.5.1.es}{III.5.5.1-fr}
\EGAReferenceAlias{III.5.5.1.eq}{III.5.5.1-fr}
\EGAReferenceAlias{eq:III.5.5.1}{III.5.5.1-fr}
\EGAReferenceAlias{III.5.5.2-es}{III.5.5.2}
\EGAReferenceAlias{III.5.5.2.fr}{III.5.5.2}
\EGAReferenceAlias{III.5.5.2.es}{III.5.5.2}
\EGAReferenceAlias{III.5.5.2.eq}{III.5.5.2}
\EGAReferenceAlias{eq:III.5.5.2}{III.5.5.2}




\title{Elementos de geometría algebraica\\
III. Estudio cohomológico de los haces coherentes\\
Primera parte}
\author{Alexandre Grothendieck y Jean Dieudonné}
\date{}
\maketitle

\begin{center}
\large Traducción española — Rama española diplomática — texto diplomático
\end{center}

\begin{quote}\small
Traducción, correcciones y composición producidas por IA: OpenAI Codex, GPT-5.6 Sol y GPT-6 Astra, esfuerzo Ultra.
\end{quote}

\clearpage
\thispagestyle{plain}
\markboth{PRELIMINARES}{PRELIMINARES}
\oldpage[0\textsubscript{III}]{5}
\phantomsection
\label{chapter:0III-fr}

\begin{center}
{\large CAPÍTULO 0 (\textit{continuación})\footnotemark\par}
\vspace{1.5em}
{\Large\bfseries PRELIMINARES\par}
\end{center}

\footnotetext{Para facilitar la búsqueda de las referencias, se remitirá
 en adelante a los parágrafos del capítulo 0 publicados con el capítulo I mediante
 el signo $\mathbf{0}_{\mathrm I}$.}

\setcounter{section}{7}


% BEGIN tree/repo/editions/III1_DIPLOMATIC/ega0/ega0-8.tex
\section{Funtores representables}
\label{section:0.8}
\label{section:0.8-fr}

\subsection{Funtores representables}
\label{subsection:0.8.1}
\label{subsection:0.8.1-fr}


\begin{env}[8.1.1]
\label{0.8.1.1}
\label{0.8.1.1-fr}
Denotaremos por $\boldsymbol{\mathit{Ens}}$ la categoría de conjuntos.
Sea $\boldsymbol{C}$ una categoría; para dos objetos $X$, $Y$ de $\boldsymbol{C}$, pondremos $h_X(Y)=\operatorname{Hom}(Y,X)$; para cada morfismo $u:Y\to Y'$ de $\boldsymbol{C}$, denotaremos por $h_X(u)$ la aplicación $v\mapsto vu$ de $\operatorname{Hom}(Y',X)$ en $\operatorname{Hom}(Y,X)$.
Es inmediato que, con estas definiciones, $h_X:\boldsymbol{C}\to\boldsymbol{\mathit{Ens}}$ es un \emph{funtor contravariante}, es decir, un objeto de la categoría $\boldsymbol{\mathit{Hom}}(\boldsymbol{C}^0,\boldsymbol{\mathit{Ens}})$ de funtores covariantes de la categoría $\boldsymbol{C}^0$ (la dual de la categoría $\boldsymbol{C}$) en la categoría $\boldsymbol{\mathit{Ens}}$ (T, 1.7, (\textit{d}) y \cite{III-29}).
\end{env}

\begin{env}[8.1.2]
\label{0.8.1.2}
\label{0.8.1.2-fr}
Sea ahora $w:X\to X'$ un morfismo de $\boldsymbol{C}$; para cada $Y\in\boldsymbol{C}$ y cada $v\in\operatorname{Hom}(Y,X)=h_X(Y)$,
se tiene $wv\in\operatorname{Hom}(Y,X')=h_{X'}(Y)$; designemos por $h_w(Y)$ la aplicación $v\mapsto wv$ de
$h_X(Y)$ en $h_{X'}(Y)$. Es inmediato que, para cada morfismo $u:Y\to Y'$ de $\boldsymbol{C}$, el
diagrama
\[
  \xymatrix{
    h_X(Y')\ar[r]^{h_X(u)}\ar[d]_{h_w(Y')} &
    h_X(Y)\ar[d]^{h_w(Y)}\\
    h_{X'}(Y')\ar[r]_{h_{X'}(u)} &
    h_{X'}(Y)
  }
\]
es conmutativo; en otros términos, $h_w$ es una \emph{transformación natural (o morfismo funtorial)
$h_X\to h_{X'}$} (T, 1.2), y también un morfismo de la categoría $\boldsymbol{\mathit{Hom}}(\boldsymbol{C}^0,\boldsymbol{\mathit{Ens}})$ (T, 1.7, (\textit{d})). Las
definiciones de $h_X$ y de $h_w$ constituyen, por tanto, la definición de un \emph{funtor covariante
canónico}
\[
  h:\boldsymbol{C}\to\boldsymbol{\mathit{Hom}}(\boldsymbol{C}^0,\boldsymbol{\mathit{Ens}}).
  \tag{8.1.2.1}
\]
\end{env}

\begin{env}[8.1.3]
\label{0.8.1.3}
\label{0.8.1.3-fr}
Sea $X$ un objeto de $\boldsymbol{C}$, y sea $F$ un funtor contravariante de $\boldsymbol{C}$ en $\boldsymbol{\mathit{Ens}}$
(un objeto de $\boldsymbol{\mathit{Hom}}(\boldsymbol{C}^0,\boldsymbol{\mathit{Ens}})$). Sea $g:h_X\to F$ una \emph{transformación natural}: para todo $Y\in\boldsymbol{C}$,
\oldpage[0\textsubscript{III}]{6}
$g(Y)$ es, pues, una aplicación $h_X(Y)\to F(Y)$ tal que, para cada morfismo $u:Y\to Y'$ de $\boldsymbol{C}$,
el diagrama
\[
  \xymatrix{
    h_X(Y')\ar[r]^{h_X(u)}\ar[d]_{g(Y')} &
    h_X(Y)\ar[d]^{g(Y)}\\
    F(Y')\ar[r]_{F(u)} & F(Y)
  }
  \tag{8.1.3.1}\label{0.8.1.3.1-fr}
\]
es conmutativo. En particular, se tiene una aplicación $g(X):h_X(X)=\operatorname{Hom}(X,X)\to F(X)$ y, por consiguiente, un elemento
\[
  \alpha(g)=(g(X))(1_X)\in F(X)
  \tag{8.1.3.2}
\]
y, en consecuencia, una aplicación canónica
\[
  \alpha:\operatorname{Hom}(h_X,F)\to F(X).
  \tag{8.1.3.3}
\]

Recíprocamente, consideremos un elemento $\xi\in F(X)$; para cada morfismo $v:Y\to X$ de $\boldsymbol{C}$, $F(v)$ es una
aplicación $F(X)\to F(Y)$; consideremos la aplicación
\[
  v\mapsto(F(v))(\xi)
  \tag{8.1.3.4}
\]
de $h_X(Y)$ en $F(Y)$; si se designa por $(\beta(\xi))(Y)$,
\[
  \beta(\xi):h_X\to F
  \tag{8.1.3.5}
\]
es una \emph{transformación natural}, pues para cada morfismo $u:Y\to Y'$ de $\boldsymbol{C}$,
$(F(vu))(\xi)=(F(v)\circ F(u))(\xi)$, lo cual hace conmutativo (\hyperref[0.8.1.3.1]{8.1.3.1}) para $g=\beta(\xi)$.
Se ha definido así una aplicación canónica
\[
  \beta:F(X)\to\operatorname{Hom}(h_X,F).
  \tag{8.1.3.6}
\]
\end{env}

\begin{proposition}[8.1.4]
\label{0.8.1.4}
\label{0.8.1.4-fr}
Las aplicaciones $\alpha$ y $\beta$ son biyecciones inversas entre sí.
\end{proposition}


Calculemos $\alpha(\beta(\xi))$ para $\xi\in F(X)$; para cada $Y\in\boldsymbol{C}$, $(\beta(\xi))(Y)$ es la aplicación
$g_1(Y):v\mapsto(F(v))(\xi)$ de $h_X(Y)$ en $F(Y)$. Se tiene, pues,
\[
  \alpha(\beta(\xi))=(g_1(X))(1_X)=(F(1_X))(\xi)
  =1_{F(X)}(\xi)=\xi.
\]
Calculemos ahora $\beta(\alpha(g))$ para $g\in\operatorname{Hom}(h_X,F)$; para cada $Y\in\boldsymbol{C}$, $(\beta(\alpha(g)))(Y)$
es la aplicación $v\mapsto(F(v))((g(X))(1_X))$; por la conmutatividad de (\hyperref[0.8.1.3.1]{8.1.3.1}), esta aplicación
no es otra que $v\mapsto(g(Y))((h_X(v))(1_X))=(g(Y))(v)$, por definición de $h_X(v)$; en otros términos,
es igual a $g(Y)$, lo que termina la demostración.


\begin{env}[8.1.5]
\label{0.8.1.5}
\label{0.8.1.5-fr}
Recordemos que una \emph{subcategoría} $\boldsymbol{C}'$ de una categoría $\boldsymbol{C}$ se define por la condición de que sus objetos
sean objetos de $\boldsymbol{C}$, y de que, si $X'$, $Y'$ son dos objetos de $\boldsymbol{C}'$, el conjunto
$\operatorname{Hom}_{\mathbf{C}'}(X',Y')$ de morfismos $X'\to Y'$ \emph{en $\boldsymbol{C}'$} sea un subconjunto del conjunto $\operatorname{Hom}_{\mathbf{C}}(X',Y')$ de
morfismos $X'\to Y'$ \emph{en $\boldsymbol{C}$}, siendo la aplicación canónica de ``composición de morfismos''
\[
  \operatorname{Hom}_{\mathbf{C}'}(X',Y')\times\operatorname{Hom}_{\mathbf{C}'}(Y',Z')
  \to\operatorname{Hom}_{\mathbf{C}'}(X',Z')
\]
\oldpage[0\textsubscript{III}]{7}
la restricción de la aplicación canónica
\[
  \operatorname{Hom}_{\mathbf{C}}(X',Y')\times\operatorname{Hom}_{\mathbf{C}}(Y',Z')
  \to\operatorname{Hom}_{\mathbf{C}}(X',Z').
\]

Se dice que $\boldsymbol{C}'$ es una subcategoría \emph{plena} de $\boldsymbol{C}$ si $\operatorname{Hom}_{\mathbf{C}'}(X',Y')=\operatorname{Hom}_{\mathbf{C}}(X',Y')$ para todo
par de objetos de $\boldsymbol{C}'$. La subcategoría $\boldsymbol{C}''$ de $\boldsymbol{C}$ formada por los objetos de $\boldsymbol{C}$ isomorfos a
objetos de $\boldsymbol{C}'$ es entonces también una subcategoría plena de $\boldsymbol{C}$, \emph{equivalente} (T, 1.2) a $\boldsymbol{C}'$, como se
comprueba fácilmente.

Un funtor covariante $F:\boldsymbol{C}_1\to\boldsymbol{C}_2$ se llama \emph{plenamente fiel} si, para todo par de objetos
$X_1$, $Y_1$ de $\boldsymbol{C}_1$, la aplicación $u\mapsto F(u)$ de $\operatorname{Hom}(X_1,Y_1)$ en $\operatorname{Hom}(F(X_1),F(Y_1))$ es
\emph{biyectiva}; esto implica que la subcategoría $F(\boldsymbol{C}_1)$ de $\boldsymbol{C}_2$ es \emph{plena}. Además,
si dos objetos $X_1$, $X'_1$ tienen la misma imagen $X_2$, existe un único isomorfismo
$u:X_1\to X'_1$ tal que $F(u)=1_{X_2}$. Para cada objeto $X_2$ de $F(\boldsymbol{C}_1)$, sea $G(X_2)$ uno de
los objetos $X_1$ de $\boldsymbol{C}_1$ tales que $F(X_1)=X_2$ ($G$ se define mediante el axioma de elección); para
cada morfismo $v:X_2\to Y_2$ de $F(\boldsymbol{C}_1)$, $G(v)$ será el único morfismo $u:G(X_2)\to G(Y_2)$ tal
que $F(u)=v$; $G$ es entonces un \emph{funtor} de $F(\boldsymbol{C}_1)$ en $\boldsymbol{C}_1$; $FG$ es el funtor identidad de
$F(\boldsymbol{C}_1)$, y lo anterior muestra que existe un isomorfismo de funtores $\varphi:\boldsymbol{\mathit{1}}_{\mathbf{C}_1}\to GF$
tal que $F$, $G$, $\varphi$ y la identidad $\boldsymbol{\mathit{1}}_{F(\mathbf{C}_1)}\to FG$ definen una \emph{equivalencia} entre
la categoría $\boldsymbol{C}_1$ y la subcategoría plena $F(\boldsymbol{C}_1)$ de $\boldsymbol{C}_2$ (T, 1.2).
\end{env}

\begin{env}[8.1.6]
\label{0.8.1.6-fr}
Apliquemos la prop. (\hyperref[0.8.1.4-fr]{8.1.4}) al caso en que el funtor $F$ es $h_{X'}$, siendo $X'$ un
objeto cualquiera de $\boldsymbol{C}$; la aplicación $\beta:\operatorname{Hom}(X,X')\to\operatorname{Hom}(h_X,h_{X'})$ no es otra que la aplicación
$w\mapsto h_w$ definida en \sref{0.8.1.2}; puesto que esta aplicación es \emph{biyectiva}, se ve,
con la terminología de \sref{0.8.1.5}, que:
\end{env}

\begin{proposition}[8.1.7]
\label{0.8.1.7}
\label{0.8.1.7-fr}
El funtor canónico $h:\boldsymbol{C}\to\boldsymbol{\mathit{Hom}}(\boldsymbol{C}^0,\boldsymbol{\mathit{Ens}})$ es plenamente fiel.
\end{proposition}

\begin{env}[8.1.8]
\label{0.8.1.8}
\label{0.8.1.8-fr}
Sea $F$ un funtor contravariante de $\boldsymbol{C}$ en $\boldsymbol{\mathit{Ens}}$; se dice que $F$ es \emph{representable} si
existe un objeto $X\in\boldsymbol{C}$ tal que $F$ es \emph{isomorfo} a $h_X$; de la
Proposición \sref{0.8.1.7} se sigue que los datos de un $X\in\boldsymbol{C}$ y de un isomorfismo de funtores
$g:h_X\to F$ determinan $X$ salvo isomorfismo único. La Proposición \sref{0.8.1.7} implica entonces
que $h$ define una \emph{equivalencia} entre $\boldsymbol{C}$ y la subcategoría plena de
$\boldsymbol{\mathit{Hom}}(\boldsymbol{C}^0,\boldsymbol{\mathit{Ens}})$ formada por los \emph{funtores contravariantes representables}. De la
Proposición \sref{0.8.1.4} se sigue que el dato de una transformación natural $g:h_X\to F$ es
equivalente al de un elemento $\xi\in F(X)$; decir que $g$ es un \emph{isomorfismo} equivale
a la siguiente condición sobre $\xi$: \emph{para todo objeto $Y$ de $\boldsymbol{C}$, la aplicación
$v\mapsto(F(v))(\xi)$ de $\operatorname{Hom}(Y,X)$ en $F(Y)$ es biyectiva}. Cuando $\xi$ satisface esta condición, se dirá que el par $(X,\xi)$ \emph{representa} al funtor representable $F$. Por abuso de lenguaje, se dirá también que el objeto $X\in\boldsymbol{C}$ representa a $F$ si existe un $\xi\in F(X)$ tal que
$(X,\xi)$ representa a $F$, es decir, si $h_X$ es isomorfo a $F$.

Sean $F$, $F'$ dos funtores contravariantes representables de $\boldsymbol{C}$ en $\boldsymbol{\mathit{Ens}}$, y sean $h_X\to F$ y
$h_{X'}\to F'$ dos isomorfismos de funtores. Se sigue entonces de \sref{0.8.1.6} que
existe una correspondencia biyectiva canónica entre $\operatorname{Hom}(X,X')$ y el conjunto $\operatorname{Hom}(F,F')$ de
transformaciones naturales $F\to F'$.
\end{env}

\begin{env}[8.1.9]
\label{0.8.1.9}
\label{0.8.1.9-fr}
\textit{Ejemplo I. Límites proyectivos}.
La noción de funtor contravariante representable comprende, en particular, la noción ``dual'' de la noción usual de ``solución de un problema universal''.
\oldpage[0\textsubscript{III}]{8}
Más generalmente, veremos que la noción de \emph{límite proyectivo} es un caso particular de la noción de funtor representable.
Recordemos que, en una categoría $\boldsymbol{C}$, se define un \emph{sistema proyectivo} mediante el dato de un conjunto preordenado $I$, de una familia $(A_\alpha)_{\alpha\in I}$ de objetos de $\boldsymbol{C}$ y, para cada par de índices $(\alpha,\beta)$ tal que $\alpha\leq\beta$, de un morfismo $u_{\alpha\beta}:A_\beta\to A_\alpha$.
Un \emph{límite proyectivo} de este sistema en $\boldsymbol{C}$ consiste en un objeto $B$ de $\boldsymbol{C}$ (denotado $\varprojlim A_\alpha$) y, para cada $\alpha\in I$, un morfismo $u_\alpha:B\to A_\alpha$ tales que: 1\textsuperscript{o}. $u_\alpha=u_{\alpha\beta}u_\beta$ para $\alpha\leq\beta$; 2\textsuperscript{o}. para todo objeto $X$ de $\boldsymbol{C}$ y toda familia $(v_\alpha)_{\alpha\in I}$ de morfismos $v_\alpha:X\to A_\alpha$ tal que $v_\alpha=u_{\alpha\beta}v_\beta$ para $\alpha\leq\beta$, existe un único morfismo $v:X\to B$ (denotado $\varprojlim v_\alpha$) tal que $v_\alpha=u_\alpha v$ para todo $\alpha\in I$ (T, 1.8).
Esto puede interpretarse de la manera siguiente: los $u_{\alpha\beta}$ definen canónicamente aplicaciones
\[
  \overline{u}_{\alpha\beta}:\operatorname{Hom}(X,A_\beta)
  \to\operatorname{Hom}(X,A_\alpha)
\]
que definen un \emph{sistema proyectivo} de conjuntos $(\operatorname{Hom}(X,A_\alpha),\overline{u}_{\alpha\beta})$, y $(v_\alpha)$ es, por definición, un elemento del conjunto $\varprojlim_\alpha\operatorname{Hom}(X,A_\alpha)$; está claro que $X\mapsto\varprojlim\operatorname{Hom}(X,A_\alpha)$ es un \emph{funtor contravariante} de $\boldsymbol{C}$ en $\boldsymbol{\mathit{Ens}}$, y la existencia del límite proyectivo $B$ equivale a decir que $(v_\alpha)\mapsto\varprojlim v_\alpha$ es un \emph{isomorfismo} de funtores en $X$
\[
  \varprojlim\operatorname{Hom}(X,A_\alpha)
  \xrightarrow{\sim}\operatorname{Hom}(X,B)
  \tag{8.1.9.1}
\]
es decir, que el funtor $X\mapsto\varprojlim\operatorname{Hom}(X,A_\alpha)$ es \emph{representable}.
\end{env}

\begin{env}[8.1.10]
\label{0.8.1.10}
\label{0.8.1.10-fr}
\textit{Ejemplo II. Objetos finales}.
Sea $\boldsymbol{C}$ una categoría, y sea $\{a\}$ un conjunto unitario.
Consideremos el funtor contravariante $F:\boldsymbol{C}\to\boldsymbol{\mathit{Ens}}$ que envía cada objeto $X$ de $\boldsymbol{C}$ al conjunto $\{a\}$, y cada morfismo $X\to X'$ de $\boldsymbol{C}$ a la única aplicación $\{a\}\to\{a\}$.
Decir que este funtor es \emph{representable} significa que existe un objeto $\boldsymbol{e}\in\boldsymbol{C}$ tal que, para todo $Y\in\boldsymbol{C}$, $\operatorname{Hom}(Y,\boldsymbol{e})=h_{\mathbf{e}}(Y)$ es un \emph{conjunto unitario}; se dice que $\boldsymbol{e}$ es un \emph{objeto final} de $\boldsymbol{C}$, y está claro que dos objetos finales de $\boldsymbol{C}$ son isomorfos (lo cual permite definir, en general mediante el axioma de elección, \emph{un} objeto final de $\boldsymbol{C}$, que se denota entonces por $\boldsymbol{e}_{\mathbf{C}}$).
Por ejemplo, en la categoría $\boldsymbol{\mathit{Ens}}$, los objetos finales son los conjuntos unitarios; en la categoría de \emph{álgebras aumentadas} sobre un cuerpo $K$ (donde los morfismos son los homomorfismos de álgebras compatibles con el aumento), $K$ es un objeto final; en la categoría de \emph{$S$-preesquemas} \sref[I]{I.2.5.1}, $S$ es un objeto final.
\end{env}

\begin{env}[8.1.11]
\label{0.8.1.11}
\label{0.8.1.11-fr}
Para dos objetos $X$ e $Y$ de una categoría $\boldsymbol{C}$, pongamos $h'_X(Y)=\operatorname{Hom}(X,Y)$ y, para todo morfismo $u:Y\to Y'$, sea $h_X(u)$ la aplicación $v\mapsto uv$ de $\operatorname{Hom}(X,Y)$ en $\operatorname{Hom}(X,Y')$; $h'_X$ es entonces un \emph{funtor covariante $\boldsymbol{C}\to\boldsymbol{\mathit{Ens}}$}, de modo que se deduce, como en \sref{0.8.1.2}, la definición de un funtor covariante canónico $h':\boldsymbol{C}^0\to\boldsymbol{\mathit{Hom}}(\boldsymbol{C},\boldsymbol{\mathit{Ens}})$; un funtor \emph{covariante} $F$ de $\boldsymbol{C}$ en $\boldsymbol{\mathit{Ens}}$, es decir, \emph{un} objeto de $\boldsymbol{\mathit{Hom}}(\boldsymbol{C},\boldsymbol{\mathit{Ens}})$, es entonces \emph{representable} si existe un objeto $X\in\boldsymbol{C}$ (necesariamente único salvo isomorfismo único) tal que $F$ es \emph{isomorfo a $h'_X$}; dejamos al lector el desarrollo de las nociones ``duales'' de las anteriores, que esta vez comprenden la noción de \emph{límite inductivo}, y en particular la noción usual de ``solución de un problema universal''.
\end{env}

\subsection{Estructuras algebraicas en categorías}
\label{subsection:0.8.2}
\label{subsection:0.8.2-fr}

\begin{env}[8.2.1]
\label{0.8.2.1}
\label{0.8.2.1-fr}
\oldpage[0\textsubscript{III}]{9}
Dados dos funtores contravariantes $F$ y $F'$ de $\boldsymbol{C}$ en $\boldsymbol{\mathit{Ens}}$, recordemos que, para todo objeto $Y\in\boldsymbol{C}$, se pone $(F\times F')(Y)=F(Y)\times F'(Y)$, y para todo morfismo $u:Y\to Y'$ de $\boldsymbol{C}$, se pone $(F\times F')(u)=F(u)\times F'(u)$, que es la aplicación $(t,t')\mapsto(F(u)(t),F'(u)(t'))$ de $F(Y')\times F'(Y')$ en $F(Y)\times F'(Y)$; $F\times F':\boldsymbol{C}\to\boldsymbol{\mathit{Ens}}$ es así un \emph{funtor contravariante} (que no es otro que el \emph{producto} de los objetos $F$ y $F'$ en la categoría $\boldsymbol{\mathit{Hom}}(\boldsymbol{C}^0,\boldsymbol{\mathit{Ens}})$).
Dado un objeto $X\in\boldsymbol{C}$, llamaremos \emph{ley de composición interna} sobre $X$ a una \emph{transformación natural}
\[
  \gamma_X:h_X\times h_X\to h_X.
  \tag{8.2.1.1}
\]

En otros términos (T, 1.2), para todo objeto $Y\in\boldsymbol{C}$, $\gamma_X(Y)$ es una aplicación $h_X(Y)\times h_X(Y)\to h_X(Y)$ (por tanto, por definición, una \emph{ley de composición interna} sobre el conjunto $h_X(Y)$) con la condición de que, para todo morfismo $u:Y\to Y'$ de $\boldsymbol{C}$, el diagrama
\[
  \xymatrix{
    h_X(Y')\times h_X(Y')
      \ar[rr]^{h_X(u)\times h_X(u)}\ar[dd]_{\gamma_X(Y')} &&
    h_X(Y)\times h_X(Y)\ar[dd]^{\gamma_X(Y)}\\\\
    h_X(Y')\ar[rr]_{h_X(u)} && h_X(Y)
  }
\]
sea conmutativo; esto implica que, para las leyes de composición $\gamma_X(Y)$ y $\gamma_X(Y')$, $h_X(u)$ es un \emph{homomorfismo} de $h_X(Y')$ en $h_X(Y)$.

De manera análoga, dados dos objetos $Z$ y $X$ de $\boldsymbol{C}$, se llama \emph{ley de composición externa sobre $X$, con $Z$ como dominio de operadores} a una transformación natural
\[
  \omega_{X,Z}:h_Z\times h_X\to h_X.
  \tag{8.2.1.2}
\]

Se ve como antes que, para todo $Y\in\boldsymbol{C}$, $\omega_{X,Z}(Y)$ es una ley de composición externa sobre $h_X(Y)$, con $h_Z(Y)$ como dominio de operadores, y tal que, para todo morfismo $u:Y\to Y'$, $h_X(u)$ y $h_Z(u)$ forman un \emph{dihomomorfismo} de $(h_Z(Y'),h_X(Y'))$ en $(h_Z(Y),h_X(Y))$.
\end{env}

\begin{env}[8.2.2]
\label{0.8.2.2}
\label{0.8.2.2-fr}
Sea $X'$ un segundo objeto de $\boldsymbol{C}$, y supongamos dada sobre $X'$ una ley de composición interna $\gamma_{X'}$; diremos que un morfismo
$w:X\to X'$ de $\boldsymbol{C}$ es un \emph{homomorfismo} para las leyes de composición si, para todo $Y\in\boldsymbol{C}$, $h_w(Y):h_X(Y)\to h_{X'}(Y)$ es un \emph{homomorfismo} para las leyes de composición $\gamma_X(Y)$ y $\gamma_{X'}(Y)$.
Si $X''$ es un tercer
\oldpage[0\textsubscript{III}]{10}
objeto de $\boldsymbol{C}$ provisto de una ley de composición interna $\gamma_{X''}$, y $w':X'\to X''$ es un morfismo de $\boldsymbol{C}$ que es un homomorfismo para $\gamma_{X'}$ y $\gamma_{X''}$, está claro que el morfismo $w'w:X\to X''$ es un homomorfismo para las leyes de composición $\gamma_X$ y $\gamma_{X''}$.
Un isomorfismo $w:X\xrightarrow{\sim}X'$ de $\boldsymbol{C}$ se llama \emph{isomorfismo para las leyes de composición} $\gamma_X$ y $\gamma_{X'}$ si $w$ es un homomorfismo para estas leyes de composición, y si su morfismo inverso $w^{-1}$ es un homomorfismo para las leyes de composición $\gamma_{X'}$ y $\gamma_X$.

Se definen de manera análoga los \emph{dihomomorfismos} para pares de objetos de $\boldsymbol{C}$ provistos de leyes de composición externa.
\end{env}

\begin{env}[8.2.3]
\label{0.8.2.3}
\label{0.8.2.3-fr}
Cuando una ley de composición interna $\gamma_X$ sobre un objeto $X\in \boldsymbol{C}$ es tal que $\gamma_X(Y)$ es una ley de \emph{grupo} sobre $h_X(Y)$ para \emph{todo} $Y\in \boldsymbol{C}$, se dice que $X$, provisto de esta ley, es un \emph{$\boldsymbol{C}$-grupo} o un \emph{objeto grupo de $\boldsymbol{C}$}.
Se definen análogamente los \emph{$\boldsymbol{C}$-anillos}, \emph{$\boldsymbol{C}$-módulos}, etc.
\end{env}

\begin{env}[8.2.4]
\label{0.8.2.4}
\label{0.8.2.4-fr}
Supongamos que el \emph{producto} $X\times X$ de un objeto $X\in \boldsymbol{C}$ por sí mismo existe en $\boldsymbol{C}$; por definición, se tiene entonces $h_{X\times X}=h_X\times h_X$ salvo isomorfismo canónico, pues es un caso particular del límite proyectivo \sref{0.8.1.9}; una ley de composición interna sobre $X$ puede, por tanto, considerarse como un morfismo funtorial $\gamma_X:h_{X\times X}\to h_X$, y determina así canónicamente \sref{0.8.1.6} un elemento $c_X\in\operatorname{Hom}(X\times X,X)$ tal que $h_{c_X}=\gamma_X$; en este caso, el dato de una ley de composición interna sobre $X$ equivale al dato de un morfismo $X\times X\to X$; cuando $\boldsymbol{C}$ es la categoría \textbf{\textit{Ens}}, se recupera la noción clásica de ley de composición interna sobre un conjunto. Se tiene un resultado análogo para una ley de composición externa cuando el producto $Z\times X$ existe en $\boldsymbol{C}$.
\end{env}

\begin{env}[8.2.5]
\label{0.8.2.5}
\label{0.8.2.5-fr}
Con la notación anterior, supongamos además que $X\times X\times X$ existe en $\boldsymbol{C}$; la caracterización del producto como objeto que representa un funtor \sref{0.8.1.9} implica la existencia de isomorfismos canónicos
\[
  (X\times X)\times X\xrightarrow{\sim}X\times X\times X
  \xrightarrow{\sim}X\times(X\times X)~;
\]
si se identifica canónicamente $X\times X\times X$ con $(X\times X)\times X$, la aplicación $\gamma_X(Y)\times 1_{h_X(Y)}$ se identifica con $h_{c_X\times 1_X}(Y)$ para todo $Y\in \boldsymbol{C}$.
En consecuencia, es equivalente decir que, para todo $Y\in \boldsymbol{C}$, la ley interna $\gamma_X(Y)$ es asociativa, o que el diagrama de aplicaciones
\[
  \xymatrix{
    h_X(Y)\times h_X(Y)\times h_X(Y)
      \ar[rr]^{\gamma_X(Y)\times 1}\ar[dd]_{1\times\gamma_X(Y)} &&
    h_X(Y)\times h_X(Y)\ar[dd]^{\gamma_X(Y)}\\\\
    h_X(Y)\times h_X(Y)\ar[rr]_{\gamma_X(Y)} && h_X(Y)
  }
\]
\oldpage[0\textsubscript{III}]{11}
es conmutativo, o que el diagrama de morfismos
\[
  \xymatrix{
    X\times X\times X
      \ar[rr]^{c_X\times 1_X}\ar[dd]_{1_X\times c_X} &&
    X\times X\ar[dd]^{c_X}\\\\
    X\times X\ar[rr]_{c_X} && X
  }
\]
es conmutativo.
\end{env}

\begin{env}[8.2.6]
\label{0.8.2.6}
\label{0.8.2.6-fr}
Bajo las hipótesis de \sref{0.8.2.5}, si se quiere expresar que, para todo $Y\in \boldsymbol{C}$, que la ley interna $\gamma_X(Y)$ es una ley de \emph{grupo}, es necesario, en primer lugar, que sea asociativa y, en segundo lugar, que exista una aplicación $\alpha_X(Y):h_X(Y)\to h_X(Y)$ con las propiedades del \emph{inverso} en un grupo; como para todo morfismo $u:Y\to Y'$ de $\boldsymbol{C}$, $h_X(u)$ debe ser un homomorfismo de grupos $h_X(Y')\to h_X(Y)$, se ve primero que $\alpha_X:h_X\to h_X$ debe ser una \emph{transformación natural}.
Por otra parte, se pueden expresar las propiedades características del inverso $s\leadsto s^{-1}$ de un grupo $G$ sin hacer intervenir el elemento identidad: basta comprobar \textbf{que} las dos aplicaciones compuestas
\[
  (s,t)\leadsto(s,s^{-1},t)\leadsto(s,s^{-1}t)\leadsto s(s^{-1}t)
\]
\[
  (s,t)\leadsto(s,s^{-1},t)\leadsto(s,ts^{-1})\leadsto(ts^{-1})s
\]
son iguales a la segunda proyección $(s,t)\leadsto t$ de $G\times G$ en $G$.
Por \sref{0.8.1.3}, se tiene $\alpha_X=h_{a_X}$, donde $a_X\in\operatorname{Hom}(X,X)$; la primera condición anterior expresa entonces que el morfismo compuesto
\[
  X\times X\xrightarrow{(1_X,a_X)\times 1_X}X\times X\times X
  \xrightarrow{1_X\times c_X}X\times X\xrightarrow{c_X}X
\]
es la segunda proyección $X\times X\to X$ en $\boldsymbol{C}$, y la segunda condición es análoga.
\end{env}

\begin{env}[8.2.7]
\label{0.8.2.7-fr}
Supongamos ahora que existe en $\boldsymbol{C}$ un \emph{objeto final} $\boldsymbol{e}$ (\hyperref[0.8.1.10-fr]{8.1.10}). Supongamos siempre que $\gamma_X(Y)$ es una ley de grupo sobre $h_X(Y)$ para todo $Y\in \boldsymbol{C}$, y denotemos por $\eta_X(Y)$ el elemento identidad de $\gamma_X(Y)$.
Como, para todo morfismo $u:Y\to Y'$ de $\boldsymbol{C}$, $h_X(u)$ es un homomorfismo de grupos, se tiene $\eta_X(Y)=(h_X(u))(\eta_X(Y'))$; tomando en particular $Y'=\boldsymbol{e}$, caso en el cual $u$ es el único elemento $\varepsilon$ de $\operatorname{Hom}(Y,\boldsymbol{e})$, se ve que el elemento $\eta_X(\boldsymbol{e})$ determina por completo $\eta_X(Y)$ para todo $Y\in \boldsymbol{C}$.
Pongamos $e_X=\eta_X(X)$, el elemento identidad del grupo $h_X(X)=\operatorname{Hom}(X,X)$; la conmutatividad del diagrama
\[
  \xymatrix{
    h_X(\boldsymbol{e})\ar[r]^{h_X(\varepsilon)}\ar[d]_{h_{e_X}(\boldsymbol{e})} &
    h_X(Y)\ar[d]^{h_{e_X}(Y)}\\
    h_X(\boldsymbol{e})\ar[r]_{h_X(\varepsilon)} & h_X(Y)
  }
\]
(cf. \sref{0.8.1.2}) muestra que, sobre el conjunto $h_X(Y)$, la aplicación $h_{e_X}(Y)$ no es otra que
\oldpage[0\textsubscript{III}]{12}
$s\leadsto\eta_X(Y)$, que envía todo elemento al elemento identidad.
Se comprueba entonces que el hecho de que $\eta_X(Y)$ sea el elemento identidad de $\gamma_X(Y)$ para todo $Y\in \boldsymbol{C}$ equivale a decir que el morfismo compuesto
\[
  X\xrightarrow{(1_X,1_X)}X\times X
  \xrightarrow{1_X\times e_X}X\times X\xrightarrow{c_X}X,
\]
y su análogo en el que se intercambian $1_X$ y $e_X$, son ambos \emph{iguales a} $1_X$.
\end{env}

\begin{env}[8.2.8]
\label{0.8.2.8}
\label{0.8.2.8-fr}
Naturalmente, se podrían ampliar con facilidad los ejemplos de estructuras algebraicas en categorías.
El ejemplo de los grupos se ha tratado con suficiente detalle, pero en adelante dejaremos por lo general al lector el desarrollo de las nociones análogas para los ejemplos de estructuras algebraicas que encontremos.
\end{env}

% END tree/repo/editions/III1_DIPLOMATIC/ega0/ega0-8.tex


% BEGIN tree/repo/editions/III1_DIPLOMATIC/ega0/ega0-9.tex
\section{Conjuntos construibles}
\label{section:0.9}
\label{section:0.9-fr}

\subsection{Conjuntos construibles}
\label{subsection:0.9.1}
\label{subsection:0.9.1-fr}

\begin{definition}[9.1.1]
\label{0.9.1.1}
\label{0.9.1.1-fr}
Se dice que una aplicación continua $f:X\to Y$ es cuasicompacta si, para todo subconjunto abierto cuasicompacto $U$ de $Y$, $f^{-1}(U)$ es cuasicompacto.
Se dice que un subconjunto $Z$ de un espacio topológico $X$ es retrocompacto en $X$ si la inyección canónica $Z\to X$ es cuasicompacta; en otros términos, si para todo subconjunto abierto cuasicompacto $U$ de $X$, $U\cap Z$ es cuasicompacto.
\end{definition}

Un subconjunto \emph{cerrado} de $X$ es retrocompacto en $X$, pero un subconjunto cuasicompacto de $X$ no es necesariamente retrocompacto en $X$.
Si $X$ es cuasicompacto, todo subconjunto abierto retrocompacto de $X$ es cuasicompacto.
Está claro que toda unión \emph{finita} de conjuntos retrocompactos en $X$ es retrocompacta en $X$, pues toda unión finita de conjuntos cuasicompactos es cuasicompacta.
Toda intersección finita de abiertos retrocompactos de $X$ es un abierto retrocompacto de $X$.
En un espacio \emph{localmente noetheriano} $X$, todo conjunto cuasicompacto es un subespacio noetheriano y, en consecuencia, \emph{todo subconjunto} de $X$ es retrocompacto en $X$.

\begin{definition}[9.1.2]
\label{0.9.1.2}
\label{0.9.1.2-fr}
Dado un espacio topológico $X$, se dice que un subconjunto de $X$ es construible si pertenece al menor conjunto de subconjuntos $\mathfrak{F}$ de $X$ que contiene todos los subconjuntos abiertos retrocompactos de $\mathfrak{F}$ y es estable por intersecciones finitas y complementarios \textup{(lo que implica que $\mathfrak{F}$ también es} estable por uniones finitas\textup{)}.
\end{definition}

\begin{proposition}[9.1.3]
\label{0.9.1.3}
\label{0.9.1.3-fr}
Para que un subconjunto de $X$ sea construible, es necesario y suficiente que sea una unión finita de conjuntos de la forma $U\cap\complement{V}$, donde $U$ y $V$ son abiertos retrocompactos de $X$.
\end{proposition}


Está claro que la condición es suficiente.
Para ver que es necesaria, consideremos el conjunto $\mathfrak{G}$ de las uniones finitas de conjuntos de la forma $U\cap\complement{V}$, donde $U$ y $V$ son abiertos retrocompactos de $X$; basta ver que el complementario de todo conjunto de $\mathfrak{G}$ pertenece a $\mathfrak{G}$.
Sea $Z=\bigcup_{i\in I}(U_i\cap\complement{V_i})$, donde $I$ es finito y los $U_i$ y $V_i$ son abiertos retrocompactos de $X$; se tiene
$\complement{Z}=\bigcap_{i\in I}(V_i\cup\complement{U_i})$, de modo que $Z$ es una unión finita de conjuntos que son intersecciones de cierto número de los $V_i$ y de cierto número de los $\complement{U_i}$, por tanto de la forma
\oldpage[0\textsubscript{III}]{13}
$V\cap\complement{U}$, donde $U$ es la unión de cierto número de los $U_i$ y $V$ es la intersección de cierto número de los $V_i$; pero se ha observado antes que las uniones e intersecciones finitas de abiertos retrocompactos de $X$ son abiertos retrocompactos de $X$, de donde se sigue la conclusión.


\begin{corollary}[9.1.4]
\label{0.9.1.4}
\label{0.9.1.4-fr}
Todo subconjunto construible de $X$ es retrocompacto en $X$.
\end{corollary}


Basta demostrar que, si $U$ y $V$ son abiertos retrocompactos de $X$, entonces $U\cap\complement{V}$ es retrocompacto en $X$; si $W$ es un abierto cuasicompacto de $X$, entonces $W\cap U\cap\complement{V}$ es cerrado en el espacio cuasicompacto $W\cap U$ y, por tanto, es cuasicompacto.


En particular:
\begin{corollary}[9.1.5]
\label{0.9.1.5}
\label{0.9.1.5-fr}
Para que un subconjunto abierto $U$ de $X$ sea construible, es necesario y suficiente que sea retrocompacto en $X$.
Para que un subconjunto cerrado $F$ de $X$ sea construible, es necesario y suficiente que el abierto $\complement{F}$ sea retrocompacto.
\end{corollary}

\begin{env}[9.1.6]
\label{0.9.1.6}
\label{0.9.1.6-fr}
Un caso importante es aquel en que todo subconjunto abierto cuasicompacto de $X$ es retrocompacto; en otros términos, aquel en que la intersección de dos subconjuntos abiertos cuasicompactos de $X$ es cuasicompacta (cf.~\sref[I]{I.5.5.6}).
Cuando $X$ es además cuasicompacto, esto implica que los subconjuntos abiertos retrocompactos de $X$ coinciden con los subconjuntos abiertos cuasicompactos de $X$, y que los subconjuntos construibles de $X$ son las uniones finitas de conjuntos de la forma $U\cap\complement{V}$, donde $U$ y $V$ son abiertos cuasicompactos.
\end{env}

\begin{corollary}[9.1.7]
\label{0.9.1.7}
\label{0.9.1.7-fr}
Para que un subconjunto de un espacio noetheriano $X$ sea construible, es necesario y suficiente que sea una unión finita de subconjuntos localmente cerrados de $X$.
\end{corollary}

\begin{proposition}[9.1.8]
\label{0.9.1.8}
\label{0.9.1.8-fr}
Sea $X$ un espacio topológico y $U$ un subconjunto abierto de $X$.
\begin{enumerate}
  \item[{\rm(i)}] Si $T$ es un subconjunto construible de $X$, entonces $T\cap U$ es un subconjunto construible de $U$.
  \item[{\rm(ii)}] Supongamos, además, que $U$ es retrocompacto en $X$. Para que un subconjunto $Z$ de $U$ sea construible en $X$, es necesario y suficiente que sea construible en $U$.
\end{enumerate}
\end{proposition}


\medskip\noindent
\begin{enumerate}
  \item[(i)] Mediante la Proposición \sref{0.9.1.3}, se reduce a demostrar que, si $T$ es un abierto retrocompacto de $X$, entonces $T\cap U$ es un abierto retrocompacto de $U$; en otros términos, para todo abierto cuasicompacto $W\subset U$, $T\cap U\cap W=T\cap W$ es cuasicompacto, lo que se sigue inmediatamente de la hipótesis.
  \item[(ii)] La condición es necesaria por (i), por lo que queda demostrar que es suficiente.
    Por la Proposición \sref{0.9.1.3}, basta considerar el caso en que $Z$ es un abierto retrocompacto \emph{en $U$}, pues entonces se seguirá que $U\setmin Z$ es construible en $X$, y si $Z$ y $Z'$ son dos abiertos retrocompactos \emph{en $U$}, entonces $Z\cap(U\setmin Z')$ será construible en $X$.
    Si $W$ es un abierto cuasicompacto de $X$ y $Z$ un abierto retrocompacto de $U$, se tiene $Z\cap W=Z\cap(W\cap U)$, y por hipótesis $W\cap U$ es un abierto cuasicompacto de $U$; por tanto, $W\cap Z$ es cuasicompacto y, en consecuencia, $Z$ es un abierto retrocompacto de $X$, y \emph{a fortiori} construible en $X$.
\end{enumerate}


\begin{corollary}[9.1.9]
\label{0.9.1.9}
\label{0.9.1.9-fr}
Sea $X$ un espacio topológico y $(U_i)_{i\in I}$ un recubrimiento finito de $X$ por abiertos retrocompactos de $X$.
Para que un subconjunto $Z$ de $X$ sea construible en $X$, es necesario y suficiente que para todo $i\in I$, $Z\cap U_i$ sea construible en $U_i$.
\end{corollary}

\begin{env}[9.1.10]
\label{0.9.1.10}
\label{0.9.1.10-fr}
Supongamos, en particular, que $X$ es cuasicompacto y que todo punto
\oldpage[0\textsubscript{III}]{14}
de $X$ admite un sistema fundamental de vecindades abiertas retrocompactas en $X$ (y \emph{a fortiori} cuasicompactas); entonces la condición para que un subconjunto $Z$ de $X$ sea construible en $X$ es de naturaleza \emph{local}; en otros términos, es necesario y suficiente que, para todo $x\in X$, exista una vecindad abierta $V$ de $x$ tal que $V\cap Z$ sea construible en $V$.
En efecto, si se satisface esta condición, entonces para todo $x\in X$ existe una vecindad abierta $V$ de $x$ que es \emph{retrocompacta en $X$} y tal que $V\cap Z$ es construible en $V$, por las hipótesis sobre $X$ y por la Proposición \sref{0.9.1.8}[i]; basta entonces recubrir $X$ por un número finito de estas vecindades y aplicar el Corolario \sref{0.9.1.9}.
\end{env}

\begin{definition}[9.1.11]
\label{0.9.1.11}
\label{0.9.1.11-fr}
Sea $X$ un espacio topológico.
Se dice que un subconjunto $T$ de $X$ es localmente construible en $X$ si, para todo $x\in X$, existe una vecindad abierta $V$ de $x$ tal que $T\cap V$ es construible en $V$.
\end{definition}

De la Proposición \sref{0.9.1.8}[i] se sigue que, si $V$ es tal que $V\cap T$ es construible en $V$, entonces para todo abierto $W\subset V$, $W\cap T$ es construible en $W$.
Si $T$ es localmente construible en $X$, entonces, para todo abierto $U$ de $X$, $T\cap U$ es localmente construible en $U$, en virtud de la observación anterior.
La misma observación muestra que el conjunto de subconjuntos localmente construibles de $X$ es estable por uniones finitas e intersecciones finitas; por otra parte, está claro que también es estable por paso al complementario.

\begin{proposition}[9.1.12]
\label{0.9.1.12}
\label{0.9.1.12-fr}
Sea $X$ un espacio topológico.
Todo conjunto construible en $X$ es localmente construible en $X$.
El recíproco es cierto si $X$ es cuasicompacto y si su topología admite una base formada por los conjuntos retrocompactos de $X$.
\end{proposition}


La primera afirmación se sigue de la Definición \sref{0.9.1.11}, y la segunda de \sref{0.9.1.10}.


\begin{corollary}[9.1.13]
\label{0.9.1.13}
\label{0.9.1.13-fr}
Sea $X$ un espacio topológico cuya topología admite una base formada por los conjuntos retrocompactos de $X$.
Entonces todo subconjunto $T$ localmente construible de $X$ es retrocompacto en $X$.
\end{corollary}


Sea $U$ un abierto cuasicompacto de $X$; $T\cap U$ es localmente construible en $U$, luego construible en $U$ por la Proposición~\sref{0.9.1.12} y, en consecuencia, cuasicompacto por el Corolario \sref{0.9.1.4}.


\subsection{Subconjuntos construibles de espacios noetherianos}
\label{subsection:0.9.2}
\label{subsection:0.9.2-fr}

\begin{env}[9.2.1]
\label{0.9.2.1}
\label{0.9.2.1-fr}
Se ha visto \sref{0.9.1.7} que, en un espacio noetheriano $X$, los subconjuntos construibles de $X$ son las \emph{uniones finitas de subconjuntos localmente cerrados} de $X$.

La imagen inversa de un conjunto construible de $X$ por una aplicación continua de un espacio noetheriano $X'$ en $X$ es construible en $X'$.
Si $Y$ es un subconjunto construible de un espacio noetheriano $X$, entonces los subconjuntos de $Y$ que son construibles como subespacios de $Y$ coinciden con los que son construibles como subespacios de $X$.
\end{env}

\begin{proposition}[9.2.2]
\label{0.9.2.2}
\label{0.9.2.2-fr}
Sea $X$ un espacio noetheriano irreducible, y sea $E$ un subconjunto construible de $X$.
Para que $E$ sea denso en todas partes en $X$, es necesario y suficiente que $E$ contenga un subconjunto abierto no vacío de $X$.
\end{proposition}


\oldpage[0\textsubscript{III}]{15}
La condición es evidentemente suficiente, pues todo abierto no vacío es denso en $X$.
Recíprocamente, sea $E=\bigcup_{\boldsymbol{i}=\mathbf{1}}^n(U_i\cap F_{\boldsymbol{i}})$ un subconjunto construible de $X$, siendo los $U_i$ abiertos no vacíos y los $F_i$ cerrados en $X$; se tiene entonces $\overline{E}\subset\bigcup_i F_i$.
En consecuencia, si $\overline{E}=X$, entonces $X$ es igual a uno de los $F_i$, luego $E\supset U_i$, lo que termina la demostración.


Cuando $X$ admite un punto genérico $x$
(\hyperref[0.2.1.2-fr]{$\mathbf{0}_{\mathrm I}$, 2.1.2}), la condición de (\hyperref[0.9.2.2-fr]{9.2.2}) equivale a la relación $x\in E$.

\begin{proposition}[9.2.3]
\label{0.9.2.3}
\label{0.9.2.3-fr}
Sea $X$ un espacio noetheriano.
Para que un subconjunto $E$ de $X$ sea construible, es necesario y suficiente que, para todo subconjunto cerrado irreducible $Y$ de $X$, $E\cap Y$ sea raro en $Y$ o contenga un subconjunto abierto no vacío de $Y$.
\end{proposition}


La necesidad de la condición se sigue de que $E\cap Y$ debe ser un subconjunto construible de $Y$ y de la Proposición \sref{0.9.2.2}, pues un subconjunto no denso de $Y$ es necesariamente raro en el espacio irreducible $Y$ (\hyperref[0.2.1.1-fr]{$\mathbf{0}_{\mathrm I}$, 2.1.1}). Para demostrar que la condición es suficiente, apliquemos el principio de inducción noetheriana (\hyperref[0.2.2.2-fr]{$\mathbf{0}_{\mathrm I}$, 2.2.2}) al conjunto
$\mathfrak{F}$ de subconjuntos cerrados $Y$ de $X$ tales que $Y\cap E$ es construible (respecto de $Y$ o respecto de $X$, lo que es equivalente): se puede suponer así que, para todo subconjunto cerrado $Y\ne X$ de $X$, $E\cap Y$ es construible.
Supongamos primero que $X$ no es irreducible, y sean $X_i$ ($1\leq i\leq m$) sus componentes irreducibles, necesariamente en número finito (\hyperref[0.2.2.5-fr]{$\mathbf{0}_{\mathrm I}$, 2.2.5}); por hipótesis, los
$E\cap X_i$ son construibles, luego también lo es su unión $E$.
Supongamos ahora que $X$ es irreducible; entonces, por hipótesis, si $E$ es raro, se tiene $\overline{E}\ne X$ y $E=E\cap\overline{E}$ es construible; si $E$ contiene un subconjunto abierto no vacío $U$ de $E$, es la unión de $U$ y $E\cap(X\setmin U)$; pero $X\setmin U$ es un conjunto cerrado distinto de $X$, por lo que $E\cap(X\setmin U)$ es construible; en consecuencia, $E$ mismo es construible, lo que termina la demostración.


\begin{corollary}[9.2.4]
\label{0.9.2.4}
\label{0.9.2.4-fr}
Sea $X$ un espacio noetheriano y $(E_\alpha)$ una familia filtrante creciente de subconjuntos construibles de $X$, tal que:
\begin{enumerate}
  \item[$1^\circ$] $X$ es la unión de la familia $(E_\alpha)$.
  \item[$2^\circ$] Todo subconjunto cerrado irreducible de $X$ está contenido en la adherencia de uno de los $E_\alpha$.
\end{enumerate}

Entonces existe un índice $\alpha$ tal que $X=E_\alpha$.

Cuando todo subconjunto cerrado irreducible de $X$ admite un punto genérico, puede suprimirse la hipótesis $2^\circ$.
\end{corollary}

Apliquemos el principio de inducción noetheriana
(\hyperref[0.2.2.2-fr]{$\mathbf{0}_{\mathrm I}$, 2.2.2}) al conjunto
$\mathfrak{M}$ de subconjuntos cerrados de $X$ contenidos en al menos uno de los $E_\alpha$; se puede suponer así que todo subconjunto cerrado $Y\ne X$ de $X$ está contenido en uno de los $E_\alpha$.
La proposición es evidente si $X$ no es irreducible, porque cada una de las componentes irreducibles $X_i$ de $X$ ($1\leq i\leq m$) está contenida en un $E_{\alpha_i}$, y existe un $E_\alpha$ que contiene todos los $E_{\alpha_i}$.
Supongamos ahora que $X$ es irreducible.
Por hipótesis, existe un $\beta$ tal que $X=\overline{E}_\beta$, de modo que \sref{0.9.2.2} $E_\beta$ contiene un subconjunto abierto no vacío $U$ de $X$.
Pero entonces el conjunto cerrado $X\setmin U$ está contenido en un $E_\gamma$, y basta tomar un $E_\alpha$ que contenga $E_\beta$ y $E_\gamma$.
Cuando todo subconjunto cerrado irreducible $Y$ de $X$
\oldpage[0\textsubscript{III}]{16}
admite un punto genérico $y$, existe $\alpha$ tal que $y\in E_\alpha$, de modo que $Y=\overline{\{y\}}\subset\overline{E}_\alpha$, y la condición $2^\circ$ es consecuencia de $1^\circ$.

\begin{proposition}[9.2.5]
\label{0.9.2.5}
\label{0.9.2.5-fr}
Sea $X$ un espacio noetheriano, sea $x$ un punto de $X$ y sea $E$ un subconjunto construible de $X$.
Para que $E$ sea una vecindad de $x$, es necesario y suficiente que, para todo subconjunto cerrado irreducible $Y$ de $X$ que contenga $x$, $E\cap Y$ sea denso en $Y$ (si existe un punto genérico $y$ de $Y$, esto significa también \textup{\sref{0.9.2.2}} que $y\in E$).
\end{proposition}


La condición es evidentemente necesaria; demostraremos que es suficiente.
Aplicando el principio de inducción noetheriana al conjunto $\mathfrak{M}$ de subconjuntos cerrados $Y$ de $X$ que contienen $x$ y tales que $E\cap Y$ es una vecindad de $x$ en $Y$, se puede suponer que todo subconjunto cerrado $Y\ne X$ de $X$ que contiene $x$ pertenece a $\mathfrak{M}$.
Si $X$ no es irreducible, cada una de las componentes irreducibles $X_i$ de $X$ que contienen $x$ es distinta de $X$, luego $E\cap X_i$ es una vecindad de $x$ respecto de $X_i$; en consecuencia, $E$ es una vecindad de $x$ en la unión de las componentes irreducibles de $X$ que contienen $x$, y como esta unión es una vecindad de $x$ en $X$, también lo es $E$.
Si $X$ es irreducible, $E$ es denso en $X$ por hipótesis, por lo que contiene un subconjunto abierto no vacío $U$ de $X$ \sref{0.9.2.2}; la proposición es entonces evidente si $x\in U$; en caso contrario, por hipótesis, $x$ es un punto interior de $E\cap(X\setmin U)$ relativo a $X\setmin U$, por tanto la adherencia en
$X$ de $X\setmin E$ no contiene $x$, y el complementario de esta adherencia es una vecindad de $x$ contenida en $E$, lo que termina la demostración.


\begin{corollary}[9.2.6]
\label{0.9.2.6}
\label{0.9.2.6-fr}
Sea $X$ un espacio noetheriano y sea $E$ un subconjunto de $X$.
Para que $E$ sea un abierto de $X$, es necesario y suficiente que, para todo subconjunto cerrado irreducible $Y$ de $X$ que interseca $E$, $E\cap Y$ contenga un subconjunto abierto no vacío de $Y$.
\end{corollary}


La condición es evidentemente necesaria; recíprocamente, si se satisface, implica que $E$ es construible por la Proposición \sref{0.9.2.3}.
Además, la Proposición \sref{0.9.2.5} muestra que $E$ es entonces una vecindad de cada uno de sus puntos, de donde se sigue la conclusión.


\subsection{Funciones construibles}
\label{subsection:0.9.3}
\label{subsection:0.9.3-fr}

\begin{definition}[9.3.1]
\label{0.9.3.1}
\label{0.9.3.1-fr}
Sea $h$ una aplicación de un espacio topológico $X$ en un conjunto $T$.
Se dice que $h$ es construible si $h^{-1}(t)$ es construible para todo $t\in T$, y vacío salvo para un número finito de valores de $t$; entonces, para todo subconjunto $S$ de $T$, $h^{-1}(S)$ es construible.
Se dice que $h$ es localmente construible si todo $x\in X$ posee una vecindad abierta $V$ tal que $h|V$ es construible.
\end{definition}

Toda función construible es localmente construible; el recíproco es cierto cuando $X$ es cuasicompacto y admite una base formada por los abiertos retrocompactos de $X$ (en particular, cuando $X$ es noetheriano).

\begin{proposition}[9.3.2]
\label{0.9.3.2}
\label{0.9.3.2-fr}
Sea $h$ una aplicación de un espacio noetheriano $X$ en un conjunto $T$.
Para que $h$ sea construible, es necesario y suficiente que, para todo subconjunto cerrado irreducible $Y$ de $X$, exista un subconjunto no vacío $U$ de $Y$, abierto respecto de $Y$, sobre el cual $h$ sea constante.
\end{proposition}


La condición es necesaria: en efecto, por hipótesis, $h$ solo toma en $Y$ un número finito de valores $t_i$, y cada uno de los conjuntos
$h^{-1}(t_i)\cap Y$ es construible en $Y$ \sref{0.9.2.1}; como no pueden ser todos subconjuntos raros del espacio $Y$, al menos uno de ellos contiene un abierto no vacío \sref{0.9.2.3}.

\oldpage[0\textsubscript{III}]{17}
Para ver que la condición es suficiente, apliquemos el principio de inducción noetheriana al conjunto $\mathfrak{M}$ de subconjuntos cerrados $Y$ de $X$ tales que la restricción $h|Y$ es construible; se puede suponer así que, para todo subconjunto cerrado $Y\ne X$ de $X$,
$h|Y$ es construible.
Si $X$ no es irreducible, la restricción de $h$ a cada una de las componentes irreducibles (en número finito) $X_i$ de $X$ es construible, y se sigue entonces inmediatamente de la Definición \sref{0.9.3.1} que $h$ es construible.
Si $X$ es irreducible, existe por hipótesis un subconjunto abierto no vacío $U$ de $X$ sobre el cual $h$ es constante; por otra parte, la restricción de $h$ a $X\setmin U$ es construible por hipótesis, y se sigue inmediatamente que $h$ es construible.


\begin{corollary}[9.3.3]
\label{0.9.3.3}
\label{0.9.3.3-fr}
Sea $X$ un espacio noetheriano en el que todo subconjunto cerrado irreducible admite un punto genérico.
Si $h$ es una aplicación de $X$ en un conjunto $T$ tal que, para todo $t\in T$, $h^{-1}(t)$ es construible, entonces $h$ es construible.
\end{corollary}


Si $Y$ es un subconjunto cerrado irreducible de $X$ e $y$ su punto genérico, entonces $Y\cap h^{-1}(h(y))$ es construible y contiene $y$; por tanto, \sref{0.9.2.2} este conjunto contiene un subconjunto abierto no vacío de $Y$, y basta aplicar la Proposición \sref{0.9.3.2}.


\begin{proposition}[9.3.4]
\label{0.9.3.4}
\label{0.9.3.4-fr}
Sea $X$ un espacio noetheriano en el que todo subconjunto cerrado irreducible admite un punto genérico, y sea $h$ una aplicación construible de $X$ en un conjunto ordenado.
Para que $h$ sea semicontinua superiormente sobre $X$, es necesario y suficiente que, para todo $x\in X$ y toda generalización \textup{(\hyperref[0.2.1.2-fr]{$\mathbf{0}_{\mathrm I}$, 2.1.2})} $x'$ de $x$, se tenga $h(x')\leq h(x)$.
\end{proposition}


La función $h$ solo toma un número finito de valores; por tanto, decir que es semicontinua superiormente significa que, para todo $x\in X$, el conjunto $E$ de los $y\in X$ tales que $h(y)\leq h(x)$ es una vecindad de $x$.
Por hipótesis, $E$ es un subconjunto construible de $X$; por otra parte, decir que un subconjunto cerrado irreducible $Y$ de $X$ contiene $x$ significa que su punto genérico $y$ es una generalización de $x$; la conclusión se sigue entonces de la Proposición \sref{0.9.2.5}.


% END tree/repo/editions/III1_DIPLOMATIC/ega0/ega0-9.tex


% BEGIN tree/repo/editions/III1_DIPLOMATIC/ega0/ega0-10.tex
\section{Complemento sobre los módulos planos}
\label{section:0.10}
\label{section:0.10-fr}

Para las demostraciones que falten en \hyperref[subsection:0.10.1]{(10.1)} y \hyperref[subsection:0.10.2]{(10.2)}, remitimos al lector a Bourbaki,~\emph{Alg.~comm.}, cap.~II y~III.

\subsection{Relaciones entre módulos planos y módulos libres}
\label{subsection:0.10.1}
\label{subsection:0.10.1-fr}

\begin{env}[10.1.1]
\label{0.10.1.1}
\label{0.10.1.1-fr}
Sea $A$ un anillo, $\mathfrak{J}$ un ideal de $A$ y $M$ un $A$-módulo;
para todo entero $p\geq0$, se tiene un homomorfismo canónico de $(A/\mathfrak{J})$-módulos
\[
\label{0.10.1.1.1-fr}
  \varphi_p:(M/\mathfrak{J}M)\otimes_{A/\mathfrak{J}}
  (\mathfrak{J}^p/\mathfrak{J}^{p+1})
  \longrightarrow \mathfrak{J}^pM/\mathfrak{J}^{p+1}M
  \tag{10.1.1.1}
\]
que es evidentemente \emph{sobreyectivo}.
Denotaremos por $\operatorname{gr}(A)=\bigoplus_{p\geq0}\mathfrak{J}^p/\mathfrak{J}^{p+1}$ el anillo graduado asociado a $A$ filtrado por los $\mathfrak{J}^p$, y por $\operatorname{gr}(M)=\bigoplus_{p\geq0}\mathfrak{J}^pM/\mathfrak{J}^{p+1}M$ el $\operatorname{gr}(A)$-módulo graduado asociado a $M$ filtrado por los $\mathfrak{J}^pM$;
se tiene entonces $\operatorname{gr}_p(A)=\mathfrak{J}^p/\mathfrak{J}^{p+1}$ y $\operatorname{gr}_p(M)=\mathfrak{J}^pM/\mathfrak{J}^{p+1}M$;
los $\varphi_p$ definen un homomorfismo \emph{sobreyectivo} de $\operatorname{gr}(A)$-módulos graduados
\[
\label{0.10.1.1.2-fr}
  \varphi:\operatorname{gr}_{\mathbf{0}}(M)\otimes_{\operatorname{gr}_0(A\boldsymbol{)}}
  \operatorname{gr}(A)\longrightarrow\operatorname{gr}(M).
  \tag{10.1.1.2}
\]
\end{env}

\oldpage[0\textsubscript{III}]{18}
\begin{env}[10.1.2]
\label{0.10.1.2}
\label{0.10.1.2-fr}
Supongamos que se satisface \emph{una} de las hipótesis siguientes:
\begin{enumerate}
  \item[(i)] $\mathfrak{J}$ es nilpotente;
  \item[(ii)] $A$ es noetheriano, $\mathfrak{J}$ está contenido en el radical de $A$ y $M$ es de tipo finito.
\end{enumerate}
Entonces las propiedades siguientes son equivalentes:
\begin{enumerate}
  \item[$a)$] $M$ es un $A$-módulo libre.
  \item[$b)$] $M/\mathfrak{J}M=M\otimes_A(A/\mathfrak{J})$ es un $(A/\mathfrak{J})$-módulo libre y $\operatorname{Tor}_1^A(M,A/\mathfrak{J})=0$.
  \item[$c)$] $M/\mathfrak{J}M$ es un $(A/\mathfrak{J})$-módulo libre y el homomorfismo canónico \sref{0.10.1.1.2} es inyectivo (y, por tanto, biyectivo).
\end{enumerate}
\end{env}

\begin{env}[10.1.3]
\label{0.10.1.3}
\label{0.10.1.3-fr}
Supongamos que $A/\mathfrak{J}$ es un \emph{cuerpo} (en otros términos, que $\mathfrak{J}$ sea maximal), y que se satisfaga una de las hipótesis (i), (ii) de (\hyperref[0.10.1.2-fr]{10.1.2}). Entonces las propiedades siguientes son equivalentes:
\begin{enumerate}
  \item[$a)$] $M$ es un $A$-módulo libre.
  \item[$b)$] $M$ es un $A$-módulo proyectivo.
  \item[$c)$] $M$ es un $A$-módulo plano.
  \item[$d)$] $\operatorname{Tor}_1^A(M,A/\mathfrak{J})=0$.
  \item[$e)$] El homomorfismo canónico (\hyperref[0.10.1.1.2-fr]{10.1.1.2}) es biyectivo.
\end{enumerate}
Este resultado se aplicará en particular en los dos casos siguientes:
\begin{enumerate}
  \item[(i)] $M$ es un módulo \emph{arbitrario} sobre un anillo local $A$ cuyo ideal maximal $\mathfrak{J}$ es \emph{nilpotente} (por ejemplo, un anillo artiniano local);
  \item[(ii)] $M$ es un módulo \emph{de tipo finito} sobre un anillo \emph{local noetheriano}.
\end{enumerate}
\end{env}

\subsection{Criterios locales de planitud}
\label{subsection:0.10.2}
\label{subsection:0.10.2-fr}

\begin{env}[10.2.1]
\label{0.10.2.1-fr}
Siendo las hipótesis y las notaciones las de (\hyperref[0.10.1.1-fr]{10.1.1}), consideremos las condiciones siguientes:
\begin{enumerate}
  \item[$a)$] $M$ es un $A$-módulo plano.
  \item[$b)$] $M/\mathfrak{J}M$ es un $(A/\mathfrak{J})$-módulo plano y $\operatorname{Tor}_1^A(M,A/\mathfrak{J})=0$.
  \item[$c)$] $M/\mathfrak{J}M$ es un $(A/\mathfrak{J})$-módulo plano y el homomorfismo canónico (\hyperref[0.10.1.1.2-fr]{10.1.1.2}) es biyectivo.
  \item[$d)$] Para todo $n\geq1$, $M/\mathfrak{J}^nM$ es un $(A/\mathfrak{J}^n)$-módulo plano.
\end{enumerate}

Entonces se tienen las implicaciones
\[
  a\Rightarrow b\Rightarrow c\Rightarrow d
\]
y, si $\mathfrak{J}$ es \emph{nilpotente}, las cuatro condiciones $a)$, $b)$, $c)$, $d)$ son \emph{equivalentes}. Lo mismo ocurre si $A$ es noetheriano y $M$ es \emph{idealmente separado}, es decir, si para todo ideal $\mathfrak{a}$ de $A$, el $A$-módulo $\mathfrak{a}\otimes_A M$ está \emph{separado} para la topología $\mathfrak{J}$-preádica.
\end{env}

\begin{env}[10.2.2]
\label{0.10.2.2}
\label{0.10.2.2-fr}
Sea $A$ un anillo noetheriano, $B$ una $A$-álgebra conmutativa noetheriana, $\mathfrak{J}$ un ideal de $A$ tal que $\mathfrak{J}B$ está contenido en el radical de $B$, y $M$ un $B$-\emph{módulo} de tipo finito.
Entonces, cuando $M$ se considera como $A$-\emph{módulo}, las condiciones $a)$, $b)$, $c)$, $d)$ de (\hyperref[0.10.2.1-fr]{10.2.1}) son \emph{equivalentes}.

\oldpage[0\textsubscript{III}]{19}
Este resultado se aplica sobre todo cuando $A$ y $B$ son anillos noetherianos \emph{locales} y el homomorfismo $A\to B$ es un homomorfismo \emph{local}.
Más precisamente, si $\mathfrak{J}$ es entonces el ideal \emph{maximal} de $A$, se puede, en las condiciones $b)$ y $c)$, suprimir la hipótesis de que $M/\mathfrak{J}M$ sea plano, que se satisface automáticamente, y la condición $d)$ significa que los módulos $M/\mathfrak{J}^nM$ son \emph{libres} sobre los $A/\mathfrak{J}^n$.
\end{env}

\begin{env}[10.2.3]
\label{0.10.2.3}
\label{0.10.2.3-fr}
Con las hipótesis sobre $A$, $B$, $\mathfrak{J}$ y $M$ del comienzo de \sref{0.10.2.2}, sea $\widehat A$ la completación separada de $A$ para la topología $\mathfrak{J}$-preádica, y $\widehat M$ la completación separada de $M$ para la topología $\mathfrak{J}B$-preádica.
Para que $M$ sea un $A$-módulo plano, es necesario y suficiente que $\widehat M$ sea un $\widehat A$-módulo plano.
\end{env}

\begin{env}[10.2.4]
\label{0.10.2.4}
\label{0.10.2.4-fr}
Sea $\rho:A\to B$ un homomorfismo local de anillos noetherianos locales, sea $k$ el cuerpo residual de $A$, y sean $M$ y $N$ dos $B$-módulos de tipo finito, suponiendo $N$ \emph{$A$-plano}.
Sea $u:M\to N$ un $B$-homomorfismo.
Entonces las condiciones siguientes son equivalentes.
\begin{enumerate}
  \item[(a)] $u$ es inyectivo y $\operatorname{Coker}(u)$ es un $A$-módulo plano.
  \item[(b)] $u\otimes1:M\otimes_A k\to N\otimes_A k$ es inyectivo.
\end{enumerate}
\end{env}

\begin{env}[10.2.5]
\label{0.10.2.5}
\label{0.10.2.5-fr}
Sean $\rho:A\to B$ y $\sigma:B\to C$ homomorfismos locales de anillos noetherianos locales, sea $k$ el cuerpo residual de $A$, y sea $M$ un $C$-módulo de tipo finito.
Se supone que $B$ es un $A$-módulo \emph{plano} Entonces las condiciones
siguientes son equivalentes:
\begin{enumerate}
  \item[$a)$] $M$ es un $B$-módulo plano.
  \item[$b)$] $M$ es un $A$-módulo plano y $M\otimes_A k$ es un $(B\otimes_A k)$-módulo plano.
\end{enumerate}
\end{env}

\begin{proposition}[10.2.6]
\label{0.10.2.6}
\label{0.10.2.6-fr}
Sean $A$ y $B$ anillos noetherianos locales, sea $\rho:A\to B$ un homomorfismo local, sea $\mathfrak{J}$ un ideal de $B$ contenido en el ideal maximal, y sea $M$ un $B$-módulo de tipo finito.
Supongamos que, para todo $n\geq0$, $M_n=M/\mathfrak{J}^{n+1}M$ es un $A$-módulo plano.
Entonces $M$ es un $A$-módulo plano.
\end{proposition}


Hay que probar que, para todo homomorfismo inyectivo $u:N'\to N$ de $A$-módulos de tipo finito, $v=1\otimes u:M\otimes_A N'\to M\otimes_A N$ es inyectivo.
Pero $M\otimes_A N'$ y $M\otimes_A N$ son $B$-módulos de tipo finito y, por tanto, están separados para la topología $\mathfrak{J}$-preádica (\hyperref[0.7.3.5-fr]{$\mathbf{0}_{\mathrm I}$, 7.3.5}); basta, pues, demostrar que el homomorfismo $\widehat v:(M\otimes_A N'){\,\widehat{\,}\,}\to(M\otimes_A N){\,\widehat{\,}\,}$ de los separados completados es inyectivo. Ahora bien, se tiene $\widehat v=\varprojlim v_n$, donde $v_n$ es el homomorfismo $1\otimes u:M_n\otimes_A N'\to M_n\otimes_A N$;
como, por hipótesis, $M_n$ es $A$-plano, $v_n$ es inyectivo para todo $n$, y también lo es $v$, pues el funtor $\varprojlim$ es exacto por la izquierda.


\begin{corollary}[10.2.7]
\label{0.10.2.7}
\label{0.10.2.7-fr}
Sea $A$ un anillo noetheriano, sea $B$ un anillo noetheriano local, sea $\rho:A\to B$ un homomorfismo, sea $f$ un elemento del ideal maximal de $B$, y sea $M$ un $B$-módulo de tipo finito.
Supongamos que la homotecia $f_M:x\to fx$ sobre $M$ es inyectiva y que $M/fM$ es un $A$-módulo plano.
Entonces $M$ es un $A$-módulo plano.
\end{corollary}

Pongamos $M_i:f^iM$ para $i\geq0$;
como $f_M$ es inyectiva, $M_i/M_{i+1}$ es isomorfo a $M/fM$ y, por tanto, $A$-plano para todo $i\geq0$;
la sucesión exacta
\[
  0\longrightarrow M_i/M_{i+1}\longrightarrow M/M_{i+1}
  \longrightarrow M/M_i\longrightarrow0
\]
nos da, por inducción sobre $i$, que $M/M_i$ es \emph{$A$-plano} para todo $i\geq0$ (\hyperref[0.6.1.2-fr]{$\mathbf{0}_{\mathrm I}$, 6.1.2}); se puede, pues, aplicar (\hyperref[0.10.2.6-fr]{10.2.6}). También se puede razonar directamente como sigue: para todo $A$-módulo $N$ de tipo finito, $M\otimes_A N$ es un $B$-módulo de tipo finito;
como $f$ pertenece al radical $\mathfrak n$ de $B$, la topología $(f)$-ádica sobre $M\otimes_A N$ es más fina que la
\oldpage[0\textsubscript{III}]{20}
topología $\mathfrak u$-ádica, y se sabe que esta última es \emph{separada}
(\hyperref[0.7.3.5-fr]{$\mathbf{0}_{\mathrm I}$, 7.3.5}). Por lo demás, como $M/M_i$ es $A$-plano, se tiene
\[
  f^i(M\otimes_A N)=\operatorname{Im}(M_i\otimes_A N\to M\otimes_A N)
  =\operatorname{Ker}(M\otimes_A N\to(M/M_i)\otimes_A N)
\]
(\hyperref[0.6.1.2-fr]{$\mathbf{0}_{\mathrm I}$, 6.1.2}). Sea entonces $N$ un $A$-módulo de tipo finito, y sea $N'$ un submódulo de $N$, con inyección canónica $j:N'\to N$; en el diagrama conmutativo
\[
\xymatrix{
  M\otimes_A N'\ar[r]\ar[d]_{1_M\otimes j} &
  (M/M_i)\otimes_A N'\ar[d]^{1_{M/M_i}\otimes j}\\
  M\otimes_A N\ar[r] & (M/M_i)\otimes_A N
}
\]
$1_{M/M_i}\otimes j$ es inyectivo, porque $M/M_i$ es $A$-plano; se concluye, pues, que
\[
  \operatorname{Ker}(M\otimes_A N'\to M\otimes_A N)
  \subset\operatorname{Ker}(M\otimes_A N'\to(M/M_i)\otimes_A N')
\]
para todo $i$; como la intersección de estos últimos núcleos es $0$, según se ha visto antes, también lo es el primer núcleo y, por tanto, $M$ es $A$-plano.


\begin{proposition}[10.2.8]
\label{0.10.2.8}
\label{0.10.2.8-fr}
Sea $A$ un anillo noetheriano reducido y sea $M$ un $A$-módulo de tipo finito.
Se supone que, para toda $A$-álgebra $B$ que sea un anillo de valoración discreta, $M\otimes_A B$ es un $B$-módulo plano (y, por tanto, libre \sref{0.10.1.3}).
Entonces $M$ es un $A$-módulo plano.
\end{proposition}


Se sabe que, para que $M$ sea plano, es necesario y suficiente que para todo ideal maximal $\mathfrak m$ de $A$, $M_{\mathfrak m}$ sea un $A_{\mathfrak m}$-módulo plano (\hyperref[0.6.3.3-fr]{$\mathbf{0}_{\mathrm I}$, 6.3.3}); se puede, pues, limitarse al caso en que $A$ es \emph{local} (\hyperref[0.1.2.8-fr]{$\mathbf{0}_{\mathrm I}$, 1.2.8}). Sean entonces $\mathfrak m$ el ideal maximal de $A$, sean $\mathfrak p_i$
$(1\leq i\leq r)$ sus ideales primos minimales, $k$ el cuerpo residual $A/\mathfrak m$.
Se sabe \sref[II]{II.7.1.7} que existe, para cada $i$, un anillo de valoración discreta $B_i$ que tiene el mismo cuerpo de fracciones $K_i$ que el anillo íntegro $A/\mathfrak p_i$, y que domina a este último. Pongamos $M_i=M\otimes_A B_i$.
Por hipótesis, $M_i$ es libre sobre $B_i$ y, denotando por $k_i$ el cuerpo residual de $B_i$
\[
\label{0.10.2.8.1-fr}
  \operatorname{rg}_{k_i}(M_i\otimes_{B_i}k_i)
  =\operatorname{rg}_{K_i}(M_i\otimes_{B_i}K_i)
  \tag{10.2.8.1}
\]

Pero está claro que el homomorfismo compuesto $A\to A/\mathfrak p_i\to B_i$ es local, de modo que $k$ es una extensión de $k$,
y se tiene
$M_i\otimes_{B_i}k_i=M\otimes_A k_i=(M\otimes_A k)\otimes_k k_i$, así como $M_i\otimes_{B_i}K_i=M\otimes_A K_i$.
La ecuación~\sref{0.10.2.8.1} implica, por tanto, que
\[
  \operatorname{rg}_k(M\otimes_A k)=\operatorname{rg}_{K_i}(M\otimes_A K_i)
  \qquad\text{para }1\leq i\leq r
\]
y, como $A$ es reducido, se sabe que esta condición implica que $M$ es un $A$-módulo \emph{libre} (Bourbaki, \emph{Alg. comm.}, cap.~II, \textsection~3, n\textsuperscript{o}~2, prop.~7).


\subsection{Existencia de extensiones planas de anillos locales}
\label{subsection:0.10.3}
\label{subsection:0.10.3-fr}

\begin{proposition}[10.3.1]
\label{0.10.3.1}
\label{0.10.3.1-fr}
Sea $A$ un anillo noetheriano local, con ideal maximal $\mathfrak J$ y cuerpo residual $k=A/\mathfrak J$.
Sea $K$ una extensión del cuerpo $k$.
Entonces existe un homomorfismo local de $A$ en un anillo noetheriano local $B$, tal que $B/\mathfrak J B$ es $k$-isomorfo a $K$ y tal que $B$ es un $A$-módulo plano.
\end{proposition}

El resto de esta sección está dedicado a demostrar esta proposición paso a paso.
\begin{env}[10.3.1.1]
\label{0.10.3.1.1}
\label{0.10.3.1.1-fr}
Supongamos primero que $K=k(T)$, donde $T$ es una indeterminada.
En el anillo de polinomios $A'=A[T]$, consideremos el ideal primo $\mathfrak J'=\mathfrak J A'$, formado por los
\oldpage[0\textsubscript{III}]{21}
polinomios cuyos coeficientes pertenecen al ideal $\mathfrak J$;
está claro que $A'/\mathfrak J'$ es canónicamente isomorfo a $k[T]$.
Mostraremos que el anillo de fracciones $B=A'_{\mathfrak J'}$ responde a la cuestión; es evidentemente un anillo local noetheriano cuyo ideal maximal $\mathfrak L=\mathfrak J B$.
Además, $B/\mathfrak L=(A'/\mathfrak J')_{\mathfrak J'}=(k[T])_{\mathfrak J'}$ es precisamente el cuerpo de fracciones $K$ de $k[T]$.
Finalmente, $B$ es un $A'$-módulo plano y $A'$ es un $A$-módulo libre, luego $B$ es un $A$-módulo plano (\hyperref[0.6.2.1-fr]{$\mathbf{0}_{\mathrm I}$, 6.2.1}).
\end{env}

\begin{env}[10.3.1.2]
\label{0.10.3.1.2-fr}
Supongamos después que $K=k(t)=k[t]$, donde $t$ es algebraico sobre $k$;
sea $f\in k[T]$ el polinomio minimal de $t$;
existe un polinomio mónico $F\in A[T]$ cuya imagen canónica en $k[T]$ es $f$. Pongamos también $A'=A[T]$, y sea $\mathfrak J'$ el ideal $\mathfrak J A'+(F)$ de $A'$.
Veremos que el anillo cociente $B=A'/(F)$ responde esta vez a la cuestión. En primer lugar, está claro que $B$ es un $A$-módulo \emph{libre} y, por tanto, plano.
El anillo $A'/\mathfrak J'$ es isomorfo a
\[
  (A'/\mathfrak J A')/((\mathfrak J A'+(F))/\mathfrak J A')
  =k[T]/(f)=K~;
\]
la imagen $\mathfrak L$ de $\mathfrak J'$ en $B$ es, por tanto, maximal, y se tiene evidentemente $\mathfrak L=\mathfrak J B$.
Finalmente, $B$ es un anillo semilocal, porque es un $A$-módulo de tipo finito (Bourbaki, \emph{Alg. comm.}, cap.~IV, \textsection~2, n\textsuperscript{o}~5, cor.~3 de la prop.~9), y sus ideales maximales están en correspondencia biyectiva con los de $B/\mathfrak J B$ (\cite[vol.~I, p.~259]{I-13});
los argumentos anteriores demuestran entonces que $B$ es un anillo local.
\end{env}

\begin{lemma}[10.3.1.3]
\label{0.10.3.1.3}
\label{0.10.3.1.3-fr}
Sea $(A_\lambda,f_{\mu\lambda})$ un sistema inductivo filtrante de anillos locales, tal que los $f_{\mu\lambda}$ son homomorfismos locales;
sea $\mathfrak m_\lambda$ el ideal maximal de $A_\lambda$, y sea $K_\lambda=A_\lambda/\mathfrak m_\lambda$.
Entonces $A'=\varinjlim A_\lambda$ es un anillo local, con ideal maximal $\mathfrak m'=\varinjlim\mathfrak m_\lambda$ y cuerpo residual $K=\varinjlim K_\lambda$.
Además, si $\mathfrak m_\mu=\mathfrak m_\lambda A_\mu$ para $\lambda<\mu$, se tiene $\mathfrak m'=\mathfrak m_\lambda A'$ para todo $\lambda$.
Si, además, para $\lambda<\mu$,
$A_\mu$ es un $A_\lambda$-módulo plano y todos los $A_\lambda$ son noetherianos, entonces $A'$ es noetheriano y es un $A_\lambda$-módulo plano para todo $\lambda$.
\end{lemma}


Como, por hipótesis, $f_{\mu\lambda}(\mathfrak m_\lambda)\subset\mathfrak m_\mu$ para $\lambda<\mu$, los $\mathfrak m_\lambda$ forman un sistema inductivo y su límite $\mathfrak m'$ es evidentemente un ideal de $A'$.
Además, si $x'\notin\mathfrak m'$, existe un $\lambda$ tal que $x'=f_\lambda(x_\lambda)$ para algún $x_\lambda\in A_\lambda$ (donde $f_\lambda:A_\lambda\to A'$ denota el homomorfismo canónico);
como $x'\notin\mathfrak m'$, necesariamente $x_\lambda\notin\mathfrak m_\lambda$, de modo que $x_\lambda$ admite un inverso $y_\lambda$ en $A_\lambda$, y $y'=f_\lambda(y_\lambda)$ es el inverso de $x'$ en $A'$, lo que demuestra que $A'$ es un anillo local con ideal maximal $\mathfrak m'$;
la afirmación sobre $K$ se sigue inmediatamente de que $\varinjlim$ es un funtor exacto.
La hipótesis $\mathfrak m_\mu=\mathfrak m_\lambda A_\mu$ implica que la aplicación canónica $\mathfrak m_\lambda\otimes_{A_\lambda}A_\mu\to\mathfrak m_\mu$ es sobreyectiva;
la igualdad $\mathfrak m'=\mathfrak m_\lambda A'$ se sigue entonces, de nuevo, de que el funtor $\varinjlim$ es exacto y conmuta con el producto tensorial.

Supongamos ahora que, para $\lambda<\mu$, se tenga $\mathfrak m_\mu=\mathfrak m_\lambda A_\mu$ y que $A_\mu$ es un $A_\lambda$-módulo plano.
Entonces $A'$ es un $A_\lambda$-módulo plano para todo $\lambda$, en virtud de (\hyperref[0.6.2.3-fr]{$\mathbf{0}_{\mathrm I}$, 6.2.3}); como $A'$ y $A_\lambda$ son anillos locales y $\mathfrak m'=\mathfrak m_\lambda A'$,
$A'$ es incluso un módulo \emph{fielmente plano} sobre $A_\lambda$-módulo \emph{fielmente plano} (\hyperref[0.6.6.2-fr]{$\mathbf{0}_{\mathrm I}$, 6.6.2}). Supongamos finalmente, además, que los $A_\lambda$ son \emph{noetherianos};
las topologías $\mathfrak m_\lambda$-preádicas son entonces separadas (\hyperref[0.7.3.5-fr]{$\mathbf{0}_{\mathrm I}$, 7.3.5}); mostremos que de ello se sigue primero que sobre $A'$ la topología $\mathfrak m'$-ádica es \emph{separada}. En efecto, si $x'\in A'$ pertenece a todos los $\mathfrak m'^{n}$ ($(n>0)$), entonces es la imagen de algún $x_\mu\in A_\mu$ para cierto índice $\mu$, y como la imagen inversa en $A_\mu$ de $\mathfrak m'^{n}=\mathfrak m_\mu^n A'$ es $\mathfrak m_\mu^n$
(\hyperref[0.6.6.1-fr]{$\mathbf{0}_{\mathrm I}$, 6.6.1}), $x_\mu$ pertenece a todos los $\mathfrak m_\mu^n$, luego $x_\mu=0$ por hipótesis y, por tanto, $x'=0$.
Denotemos por $\widehat A'$ la completación de $A'$ para la topología $\mathfrak m'$-ádica;
lo anterior muestra que se tiene $A'\subset\widehat A'$.
Mostraremos ahora que $\widehat A'$ es \emph{noetheriano} y \emph{$A_\lambda$-plano} para todo $\lambda$;
de esto,
\oldpage[0\textsubscript{III}]{22}
se seguirá que $\widehat A'$ es $A'$-plano (\hyperref[0.6.2.3-fr]{$\mathbf{0}_{\mathrm I}$, 6.2.3}), y como $\mathfrak m'\widehat A'\ne\widehat A'$,
$\widehat A'$ es un $A'$-módulo fielmente plano (\hyperref[0.6.6.2-fr]{$\mathbf{0}_{\mathrm I}$, 6.6.2}), de donde se concluirá finalmente que $A'$ es \emph{noetheriano} (\hyperref[0.6.5.2-fr]{$\mathbf{0}_{\mathrm I}$, 6.5.2}), lo que concluirá la demostración del lema.

Se tiene $\widehat A'=\varprojlim_n A'/\mathfrak m'^{n}$;
por el hecho de que $A'$ es $A_\lambda$-plano, se tiene
\[
  \mathfrak m'^{n}/\mathfrak m'^{n+1}
  =(\mathfrak m_\lambda^n/\mathfrak m_\lambda^{n+1})
    \otimes_{A_\lambda}A'
  =(\mathfrak m_\lambda^n/\mathfrak m_\lambda^{n+1})
    \otimes_{K_\lambda}(K_\lambda\otimes_{A_\lambda}A')
  =(\mathfrak m_\lambda^n/\mathfrak m_\lambda^{n+1})
    \otimes_{K_\lambda}K~;
\]
como $\mathfrak m_\lambda^n/\mathfrak m_\lambda^{n+1}$ es un $K_\lambda$-espacio vectorial de dimensión finita, $\mathfrak m'^{n}/\mathfrak m'^{n+1}$ es un $K$-espacio vectorial de dimensión finita para todo $n\geq0$. Resulta, pues, de (\hyperref[0.7.2.12-fr]{$\mathbf{0}_{\mathrm I}$, 7.2.12}) y (\hyperref[0.7.2.8-fr]{$\mathbf{0}_{\mathrm I}$, 7.2.8}) que $\widehat A'$ es \emph{noetheriano}.
Se sabe además que el ideal maximal de $\widehat A'$ es $\mathfrak m'\widehat A'$, y que $\widehat A'/\mathfrak m'^{n}\widehat A'$ es isomorfo a $A'/\mathfrak m'^{n}$;
como $A'/\mathfrak m'^{n}=(A_\lambda/\mathfrak m_\lambda^n)
\otimes_{A_\lambda}A'$,
$A'/\mathfrak m'^{n}$ es un $(A_\lambda/\mathfrak m_\lambda^n)$-módulo plano (\hyperref[0.6.2.1-fr]{$\mathbf{0}_{\mathrm I}$, 6.2.1}); el criterio (\hyperref[0.10.2.2-fr]{10.2.2}) es, pues, aplicable a la $A_\lambda$-álgebra noetheriana $\widehat A'$, y muestra que $\widehat A'$ es \emph{$A_\lambda$-plano}.


\begin{env}[10.3.1.4]
\label{0.10.3.1.4}
\label{0.10.3.1.4-fr}
Tratemos ahora el caso general.
Existe un ordinal $\gamma$ y, para todo ordinal $\lambda\leq\gamma$, un subcuerpo $k_\lambda$ de $K$ que contiene a $k$, tales que: $1^\circ$ Para todo $\lambda<\gamma$,
$k_{\lambda+1}$ es una extensión de $k_\lambda$ generada por un solo elemento; $2^\circ$ Para todo ordinal $\mu$, $k_\mu=\bigcup_{\lambda<\mu}k_\lambda$;
$3^\circ$ $K=k_\gamma$.
En efecto, basta considerar una biyección $\xi\to t_\xi$ del conjunto de ordinales $\xi\leq\beta$ (para algún $\beta$ apropiado) en $K$, y definir $k_\lambda$ por inducción transfinita (para $\lambda\leq\beta$) como la unión de los $k_\mu$ para $\mu<\lambda$ si $\lambda$ no tiene predecesor, y, si $\lambda=\nu+1$, como $k_\nu(t_\xi)$, donde $\xi$ es el menor ordinal tal que $t_\xi\notin k_\nu$;
$\gamma$ es entonces, por definición, el menor ordinal $\leq\beta$ tal que $k_\gamma=K$.

Teniendo esto presente, definiremos, por inducción transfinita, una familia de anillos noetherianos locales $A_\lambda$ para $\lambda\leq\gamma$, y homomorfismos locales $f_{\mu\lambda}:A_\lambda\to A_\mu$ para $\lambda\leq\mu$, que satisfagan las condiciones siguientes:
\begin{enumerate}
  \item[(i)] $(A_\lambda,f_{\mu\lambda})$ es un sistema inductivo y $A_0=A$;
  \item[(ii)] para todo $\lambda$, se tiene un $k$-isomorfismo $A_\lambda/\mathfrak J A_\lambda\xrightarrow{\sim}k_\lambda$;
  \item[(iii)] para $\lambda\leq\mu$,
$A_\mu$ es un $A_\lambda$-módulo plano.
\end{enumerate}

Supongamos, pues, definidos los $A_\lambda$ y los $f_{\mu\lambda}$ para $\lambda<\mu<\xi$, y supongamos primero que $\xi=\zeta+1$, de modo que $k_\xi=k_\zeta(t)$.
Si $t$ es trascendente sobre $k_\zeta$, se define $A_\xi$, siguiendo el procedimiento de \sref{0.10.3.1.1}, como $(A_\zeta[t])_{\mathfrak J A_\zeta[t]}$;
la aplicación canónica es $f_{\zeta\xi}$ y, para $\lambda<\zeta$, se toma $f_{\xi\lambda}=f_{\xi\zeta}\circ f_{\zeta\lambda}$;
la verificación de las condiciones~(i) a (iii) es entonces inmediata, en virtud de lo demostrado en \sref{0.10.3.1.1}.
Supongamos ahora que $t$ es algebraico, y sea $h$ su polinomio minimal en $k_\zeta[T]$,
$H$ un polinomio mónico de $A_\zeta[T]$ cuya imagen en $k_\zeta[T]$ es $h$;
se toma entonces $A_\xi$ igual a $A_\zeta[T]/(H)$, definiendo los $f_{\xi\lambda}$ como antes;
la verificación de las condiciones~(i) a (iii) se sigue entonces de lo demostrado en \sref{0.10.3.1.2}.

Supongamos ahora que $\xi$ no tiene predecesor;
se toma entonces $A_\xi$ como límite inductivo del sistema inductivo de anillos locales $(A_\lambda,f_{\mu\lambda})$ para $\lambda<\xi$;
$f_{\xi\lambda}$ como la aplicación canónica para $\lambda<\xi$.
El hecho de que $A_\xi$ sea local y noetheriano, de que los $f_{\xi\lambda}$ sean homomorfismos locales y de que las condiciones~(i) a (iii) se satisfagan para $\lambda\leq\xi$
\oldpage[0\textsubscript{III}]{23}
se sigue entonces de la hipótesis de inducción y del Lema~\sref{0.10.3.1.3}.
Con esta construcción, está claro que el anillo $B=A_\gamma$ satisface el enunciado de (\hyperref[0.10.3.1-fr]{10.3.1}).
\end{env}

Se observará que, en virtud de (\hyperref[0.10.2.1-fr]{10.2.1, $c)$}), se tiene un isomorfismo canónico
\[
\label{0.10.3.1.5-fr}
  \operatorname{gr}(A)\otimes_k K\xrightarrow{\sim}\operatorname{gr}(B).
  \tag{10.3.1.5}
\]

También se puede reemplazar $B$ por su completación $\mathfrak J B$-ádica $\widehat B$ sin modificar las conclusiones de \sref{0.10.3.1}, porque $\widehat B$ es un $B$-módulo plano (\hyperref[0.7.3.3-fr]{$\mathbf{0}_{\mathrm I}$, 7.3.3}), por tanto un $A$-módulo plano (\hyperref[0.6.2.1-fr]{$\mathbf{0}_{\mathrm I}$, 6.2.1}).

Se ha demostrado además el

\begin{corollary}[10.3.2]
\label{0.10.3.2}
\label{0.10.3.2-fr}
Si $K$ es una extensión de grado finito, se puede suponer que $B$ es una $A$-álgebra finita.
\end{corollary}

% END tree/repo/editions/III1_DIPLOMATIC/ega0/ega0-10.tex


% BEGIN tree/repo/editions/III1_DIPLOMATIC/ega0/ega0-11.tex
\section{Suplemento sobre álgebra homológica}
\label{section:0.11}
\label{section:0.11-fr}

\subsection{Repaso de las sucesiones espectrales}
\label{subsection:0.11.1}
\label{subsection:0.11.1-fr}

\begin{env}[11.1.1]
\label{0.11.1.1-fr}
Utilizaremos en lo que sigue una noción de sucesión espectral más general que
la definida en (T, 2.4); conservando las notaciones de (T, 2.4),
llamaremos \emph{sucesión espectral} en una categoría abeliana $\boldsymbol{C}$ a un
sistema E formado por los elementos siguientes:
\begin{enumerate}
  \item[$a)$] Una familia $(E_r^{pq})$ de objetos de $\boldsymbol{C}$ definidos para
    $p\in\mathbb Z$, $q\in\mathbb Z$ y $r\geq2$.
  \item[$b)$] Una familia de morfismos
    $d_r^{pq}:E_r^{pq}\to E_r^{p+r,q-r+1}$ tales que
    $d_r^{p+r,q-r+1}d_r^{pq}=0$. Se pone
    $Z_{r+1}(E_r^{pq})=\operatorname{Ker}(d_r^{pq})$,
    $B_{r+1}(E_r^{pq})=\operatorname{Im}(d_r^{p+r,q-r+1})$, de modo que
    \[
      B_{r+1}(E_r^{pq})\subset Z_{r+1}(E_r^{pq})\subset E_r^{pq}.
    \]
  \item[$c)$] Una familia de isomorfismos
    $\alpha_r^{pq}:Z_{r+1}(E_r^{pq})/B_{r+1}(E_r^{pq})
    \xrightarrow{\sim}E_{r+1}^{pq}$.
\end{enumerate}

Se definen entonces para $k\geq r+1$, por inducción sobre $k$, los subobjetos
$B_k(E_r^{pq})$ y $Z_k(E_r^{pq})$ como imágenes inversas, por el morfismo
canónico $E_r^{pq}\to E_r^{pq}/B_{r+1}(E_r^{pq})$ de los subobjetos de este
cociente identificados por $\alpha_r^{pq}$ con los subobjetos $B_k(E_{r+1}^{pq})$ y
$Z_k(E_{r+1}^{pq})$ respectivamente. Está claro que se tiene entonces, salvo
isomorfismo,
\[
\label{0.11.1.1.1-fr}
  Z_k(E_r^{pq})/B_k(E_r^{pq})=E_k^{pq}
  \qquad\text{para }k\geq r+1
  \tag{11.1.1.1}
\]
y, si se pone además $B_r(E_r^{pq})=0$ y $Z_r(E_r^{pq})=E_r^{pq}$, se tienen
las relaciones de inclusión
\[
\label{0.11.1.1.2-fr}
  0=B_r(E_r^{pq})\subset B_{r+1}(E_r^{pq})\subset B_{r+2}(E_r^{pq})
  \subset\cdots\subset Z_{r+2}(E_r^{pq})\subset Z_{r+1}(E_r^{pq})
  \subset Z_r(E_r^{pq})=E_r^{pq}.
  \tag{11.1.1.2}
\]

Los otros elementos del dato de E son entonces:
\begin{enumerate}
  \item[$d)$] Dos subobjetos $B_\infty(E_2^{pq})$ y
    $Z_\infty(E_2^{pq})$ de $E_2^{pq}$ tales que se tenga
    $B_\infty(E_2^{pq})\subset Z_\infty(E_2^{pq})$ y, para todo $k\geq2$,
    \[
      B_k(E_2^{pq})\subset B_\infty(E_2^{pq})
      \qquad\text{y}\qquad
      Z_\infty(E_2^{pq})\subset Z_k(E_2^{pq}).
    \]
    Se pone
    \[
    \label{0.11.1.1.3-fr}
      E_\infty^{pq}=Z_\infty(E_2^{pq})/B_\infty(E_2^{pq}).
      \tag{11.1.1.3}
    \]
  \item[$e)$]
    \oldpage[0\textsubscript{III}]{24}
    Una familia $(E^n)$ de objetos de $\boldsymbol{C}$, cada uno provisto de una
    \emph{filtración decreciente} $(F^p(E^n))_{p\in\mathbb Z}$. Se designa
    como de costumbre por $\operatorname{gr}(E^n)$ el objeto graduado asociado al
    objeto filtrado $E^n$, suma directa de los
    $\operatorname{gr}_p(E^n)=F^p(E^n)/F^{p+1}(E^n)$.
  \item[$f)$] Para todo par $(p,q)\in\mathbb Z\times\mathbb Z$, un
    isomorfismo
    $\beta^{pq}:E_\infty^{pq}\xrightarrow{\sim}
    \operatorname{gr}_p(E^{p+q})$.
\end{enumerate}

La familia $(E^n)$, sin las filtraciones, se llama el \emph{término límite}
de la sucesión espectral E.

Supongamos que la categoría $\boldsymbol{C}$ admite sumas directas infinitas,
o que para todo $r\geq2$ y todo $n\in\mathbb Z$ los pares $(p,q)$ tales
que $p+q=n$ y $E_r^{pq}\ne0$ son en número finito (basta que sea así
para $r=2$). Entonces los
\[
  E_r^{(n)}=\sum_{p+q=n}E_r^{pq}
\]
están definidos, y designando por $d_r^{(n)}$ el morfismo
$E_r^{(n)}\to E_r^{(n+1)}$ cuya restricción a $E_r^{pq}$ es $d_r^{pq}$
para cada par $(p,q)$ tal que $p+q=n$, se tiene
$d_r^{(n+1)}\circ d_r^{(n)}=0$, en otros términos $(E_r^{(n)})_{n\in\mathbb Z}$
es un \emph{complejo $E_r^{(\bullet)}$} en $\boldsymbol{C}$, con operador de
derivación de grado $+1$, y resulta de \emph{c)} que
$\mathrm H^n(E_n^{(\bullet)})$ es \emph{isomorfo a} $E_{r+1}^{(n)}$ para todo
$r\geq2$.
\end{env}

\begin{env}[11.1.2]
\label{0.11.1.2}
\label{0.11.1.2-fr}
Un \emph{morfismo $u:E\to E'$} de una sucesión espectral E en una sucesión espectral $E'=(E_r^{\prime pq},E^{\prime n})$ consiste en sistemas de morfismos $u_r^{pq}:E_r^{pq}\to E_r^{\prime pq}$ y $u^n:E^n\to E^{\prime n}$, siendo los $u^n$ compatibles con las filtraciones de $E^n$ y $E^{\prime n}$, y tales que los diagramas
\[
  \xymatrix{
    E_r^{pq}\ar[r]^-{d_r^{pq}}\ar[d]_{u_r^{pq}} &
    E_r^{p+r,q-r+1}\ar[d]^{u_r^{p+r,q-r+1}}\\
    E_r^{\prime pq}\ar[r]_{d_r^{\prime pq}} &
    E_r^{\prime p+r,q-r+1}
  }
\]
sean conmutativos; además, pasando a los cocientes, $u_r^{pq}$ da un morfismo \[
  \overline u_r^{pq}:
  Z_{r+1}(E_r^{pq})/B_{r+1}(E_r^{pq})\longrightarrow
  Z_{r+1}(E_r^{\prime pq})/B_{r+1}(E_r^{\prime pq})
\] y debe tenerse $\alpha_r^{\prime pq}\circ\overline u_r^{pq}
=u_{r+1}^{pq}\circ\alpha_r^{pq}$; finalmente, debe tenerse $u_2^{pq}(B_\infty(E_2^{pq}))\subset B_\infty(E_2^{\prime pq})$ y $u_2^{pq}(Z_\infty(E_2^{pq}))\subset Z_\infty(E_2^{\prime pq})$; pasando a los cocientes, $u_2^{pq}$ da entonces un morfismo $u_\infty^{\prime pq}:E_\infty^{pq}\to E_\infty^{\prime pq}$, y el diagrama
\[
  \xymatrix{
    E_\infty^{pq}\ar[r]^-{u_\infty^{\prime pq}}\ar[d]_{\beta^{pq}} &
    E_\infty^{\prime pq}\ar[d]^{\beta^{\prime pq}}\\
    \operatorname{gr}_{\boldsymbol{p}}(E^{p+q})
      \ar[r]_{\operatorname{gr}_{\boldsymbol{p}}(u^{p+q})} &
    \operatorname{gr}_{\boldsymbol{p}}(E^{\prime\boldsymbol{p}+q})
  }
\]
debe ser conmutativo.

Las definiciones anteriores muestran, por inducción sobre $r$, que si los $u_2^{pq}$ son \emph{isomorfismos}, también lo son los $u_r^{pq}$ para $r\geq2$; \emph{si además se sabe} que $u_2^{pq}(B_\infty(E_2^{pq}))=B_\infty(E_2^{\prime pq})$ y $u_2^{pq}(Z_\infty(E_2^{pq}))=Z_\infty(E_2^{\prime pq})$ y que los $u^n$ son \emph{isomorfismos}, entonces se puede concluir que $u$ es un \emph{isomorfismo}.
\end{env}

\begin{env}[11.1.3]
\label{0.11.1.3}
\label{0.11.1.3-fr}
\oldpage[0\textsubscript{III}]{25}
Recordemos que si $(F^p(X))_{p\in\mathbb Z}$ es una \emph{filtración} (decreciente) de un objeto $X\in\boldsymbol{C}$, se dice que esta filtración es \emph{separada} si $\inf(F^p(X))=0$, \emph{discreta} si existe un $p$ tal que $F^p(X)=0$, \emph{exhaustiva} (o \emph{coseparada}) si $\sup(F^p(X))=X$, y \emph{codiscreta} si existe un $p$ tal que $F^p(X)=X$.

Diremos que una sucesión espectral $E=(E_r^{pq},E^n)$ es \emph{débilmente convergente} si se tiene $B_\infty(E_2^{pq})=\sup_k(B_k(E_2^{pq}))$ y $Z_\infty(E_2^{pq})=\inf_k(Z_k(E_2^{pq}))$ (en otras palabras, los objetos $B_\infty(E_2^{pq})$ y $Z_\infty(E_2^{pq})$ están determinados por los datos \emph{a)} a \emph{c)} de la sucesión espectral E). Diremos que la sucesión espectral E es \emph{regular} si es débilmente convergente y si además:
\begin{enumerate}
  \item[$1^\circ$] Para todo par $(p,q)$, la sucesión decreciente $(Z_k(E_2^{pq}))_{k\geq2}$ es \emph{estable}; la hipótesis de que E es débilmente convergente implica entonces
$Z_\infty(E_2^{pq})=Z_k(E_2^{pq})$ para $k$ suficientemente grande (dependiendo de $p$ y $q$).
  \item[$2^\circ$] Para todo $n$, la filtración $(F^p(E^n))_{p\in\mathbb Z}$ de $E^n$ es \emph{discreta} y \emph{exhaustiva}.
\end{enumerate}

Se dice que la sucesión espectral E es \emph{corregular} si es débilmente convergente y si además:
\begin{enumerate}
  \item[$3^\circ$] Para todo par $(p,q)$, la sucesión creciente $(B_k(E_2^{pq}))_{k\geq2}$ es \emph{estable}, lo que implica que $B_\infty(E_2^{pq})=B_k(E_2^{pq})$ y, en consecuencia, $E_\infty^{pq}=\inf E_k^{pq}$.
  \item[$4^\circ$] Para todo $n$, la filtración de $E^n$ es \emph{codiscreta}.
\end{enumerate}

Finalmente, se dice que E es \emph{birregular} si es a la vez regular y corregular, es decir, si se tienen las condiciones siguientes:
\begin{enumerate}
  \item[\textit{a)}] Para todo par $(p,q)$, las sucesiones $(B_k(E_2^{pq}))_{k\geq2}$ y $(Z_k(E_2^{pq}))_{k\geq2}$ son \emph{estables} y se tiene $B_\infty(E_2^{pq})=B_k(E_2^{pq})$ y $Z_\infty(E_2^{pq})=Z_k(E_2^{pq})$ para $k$ suficientemente grande (lo que implica que $E_\infty^{pq}=E_k^{pq}$).
  \item[\textit{(b)}] Para todo $n$, la filtración $(F^p(E^n))_{p\in\mathbb Z}$ es \emph{discreta} y \emph{codiscreta} (también se dice \emph{finita}).
\end{enumerate}

Las sucesiones espectrales definidas en (T, 2.4) son, por tanto, sucesiones espectrales birregulares.
\end{env}

\begin{env}[11.1.4]
\label{0.11.1.4}
\label{0.11.1.4-fr}
Supongamos que en la categoría $\boldsymbol{C}$ existen los límites inductivos filtrantes y que el funtor $\varinjlim$ es \emph{exacto} (lo que equivale a decir que se satisface el axioma (AB~5) de (T, 1.5) (cf.~T, 1.8)).
La condición de que la filtración $(F^p(X))_{p\in\mathbb Z}$ de un objeto $X\in\boldsymbol{C}$ sea exhaustiva se expresa entonces por $\varinjlim_{p\to-\infty}F^p(X)=X$. Si una sucesión espectral E es débilmente convergente, se tiene
$B_\infty(E_2^{pq})=\varinjlim_{k\to\infty}B_k(E_2^{pq})$; si además $u:E\to E'$ es un morfismo de E en una sucesión espectral débilmente convergente $E'$ en $\boldsymbol{C}$, entonces se tiene $u_2^{pq}(B_\infty(E_2^{pq}))=B_\infty(E_2^{\prime pq})$, por la exactitud de $\varinjlim$.
Además:
\end{env}

\begin{proposition}[11.1.5]
\label{0.11.1.5}
\label{0.11.1.5-fr}
Sea $\mathbf{C}$ una categoría abeliana en la que los límites
inductivos filtrantes son exactos, \textup{E}, $E'$ dos sucesiones espectrales regulares en $\mathbf{C}$, y sea $u:E\to E'$ un morfismo de sucesiones espectrales.
Si los $u_2^{pq}$ son isomorfismos, entonces $u$ también lo es.
\end{proposition}

\begin{proof}
Ya sabemos por \sref{0.11.1.2} que los $u_r^{pq}$ son isomorfismos y que
\[
  u_2^{pq}(B_\infty(E_2^{pq}))=B_\infty(E_2^{\prime pq})~;
\]
\oldpage[0\textsubscript{III}]{26}
la hipótesis de que E y $E'$ son regulares implica también que $u_2^{pq}(Z_\infty(E_2^{pq}))=Z_\infty(E_2^{\prime pq})$ y, como $u_2^{pq}$ es un isomorfismo, también lo es $u_\infty^{pq}$; se concluye así que $\operatorname{gr}_p(u^{p+q})$ es asimismo un isomorfismo.
Pero, como las filtraciones de los $E^n$ y los $E^{\prime n}$ son discretas y exhaustivas, esto implica que los $u^n$ son también isomorfismos (Bourbaki, \emph{Alg. comm.}, cap.~III, \textsection2, n\textsuperscript{o}8, th.~1).
\end{proof}

\begin{env}[11.1.6]
\label{0.11.1.6-fr}
Resulta de (\hyperref[0.11.1.1.2-fr]{11.1.1.2}) y de la definición (\hyperref[0.11.1.1.3-fr]{11.1.1.3}) que si, en una sucesión espectral E, se tiene
$E_r^{pq}=0$, entonces $E_k^{pq}=0$ para $k\geq r$ y $E_\infty^{pq}=0$.
Se dice que una sucesión espectral \emph{degenera} si existe un entero $r\geq2$ y, para todo entero $n\in\mathbb Z$, un entero $q(n)$ tal que $E_r^{n-q,q}=0$ \emph{para todo} $q\ne q(n)$.
De la observación anterior se deduce primero que también se tiene $E_k^{n-q,q}=0$ para $k\geq r$ (incluido $k=\infty$) y $q\ne q(n)$.
Además, la definición de $E_{r+1}^{pq}$ muestra que se tiene $E_{r+1}^{n-q(n),q(n)}=E_r^{n-q(n),q(n)}$; si E es \emph{débilmente convergente}, se tiene, pues, también
$E_\infty^{n-q(n),q(n)}=E_r^{n-q(n),q(n)}$; en otras palabras, para todo $n\in\mathbb Z$, $\operatorname{gr}_p(E^n)=0$ para $p\ne q(n)$ y $\operatorname{gr}_{q(n)}(E^n)=E_r^{n-q(n),q(n)}$.
Si además la filtración de $E^n$ es \emph{discreta} y \emph{exhaustiva}, la sucesión E es \emph{regular} y se tiene
$E^n=E_r^{n-q(n),q(n)}$ salvo isomorfismo.
\end{env}

\begin{env}[11.1.7]
\label{0.11.1.7}
\label{0.11.1.7-fr}
Supongamos que en la categoría $\boldsymbol{C}$ existen los límites inductivos filtrantes y son exactos, y sea $(E_\lambda,u_{\mu\lambda})$ un sistema inductivo (sobre un conjunto filtrante de índices) de sucesiones espectrales en $\boldsymbol{C}$.
Entonces el \emph{límite inductivo} de este sistema inductivo existe en la categoría aditiva de las sucesiones espectrales de objetos de $\boldsymbol{C}$: para verlo, basta definir $E_r^{pq}$, $d_r^{pq}$, $\alpha_r^{pq}$, $B_\infty(E_2^{pq})$, $Z_\infty(E_2^{pq})$, $E^n$, $F^p(E^n)$ y $\beta^{pq}$ como los límites inductivos respectivos de los $E_{r,\lambda}^{pq}$, $d_{r,\lambda}^{pq}$, $\alpha_{r,\lambda}^{pq}$, $B_\infty(E_{2,\lambda}^{pq})$, $Z_\infty(E_{2,\lambda}^{pq})$, $E_\lambda^n$, $F^p(E_\lambda^n)$ y $\beta_\lambda^{pq}$; la verificación de las condiciones de \sref{0.11.1.1} se deduce de la exactitud del funtor $\varinjlim$ sobre $\boldsymbol{C}$.
\end{env}

\begin{remark}[11.1.8]
\label{0.11.1.8}
\label{0.11.1.8-fr}
Supongamos que la categoría $\boldsymbol{C}$ es la categoría de los $A$-módulos sobre un anillo $A$ \emph{noetheriano} (resp. un anillo $A$).
Las definiciones de \sref{0.11.1.1} muestran entonces que si, para un $r$ dado, los $E_r^{pq}$ son $A$-módulos \emph{de tipo finito} (resp. \emph{de longitud finita}), también lo son todos los módulos $E_s^{pq}$ para $s\geq r$ y, por tanto, también $E_\infty^{pq}$.
Si además la filtración del término límite $(E^n)$ es \emph{discreta} y \emph{codiscreta} para todo $n$, se concluye que cada $E_n$ es también un $A$-módulo \emph{de tipo finito} (resp. \emph{de longitud finita}).
\end{remark}

\begin{env}[11.1.9]
\label{0.11.1.9-fr}
Habremos de considerar condiciones que aseguren que una sucesión espectral E sea birregular de manera «uniforme» en $p+q=n$.
Se utilizará entonces el lema siguiente:
\end{env}

\begin{lemma}[11.1.10]
\label{0.11.1.10}
\label{0.11.1.10-fr}
Sea $(E_r^{pq})$ una familia de objetos de $\mathbf{C}$ relacionados por los datos \emph{(a)}, \emph{(b)} y \emph{(c)} de \textup{\sref{0.11.1.1}}.
Para un entero fijo $n$, las propiedades siguientes son equivalentes:
\begin{enumerate}
  \item[{\rm(a)}] Existe un entero $r(n)$ tal que, para $r\geq r(n)$, $p+q=n$ o $p+q=n-1$, los morfismos $d_r^{pq}$ son nulos.
  \item[{\rm(b)}] Existe un entero $r(n)$ tal que, para $p+q=n$ o $p+q=n+1$, se tiene $B_r(E_2^{pq})=B_s(E_2^{pq})$ para $s\geq r\geq r(n)$.
  \item[{\rm(c)}] Existe un entero $r(n)$ tal que, para $p+q=n$ o $p+q=n-1$, se tiene $Z_r(E_2^{pq})=Z_s(E_2^{pq})$ para $s\geq r\geq r(n)$.
  \item[{\rm(d)}] Existe un entero $r(n)$ tal que, para $p+q=n$, se tiene $B_r(E_2^{pq})=B_s(E_2^{pq})$ y $Z_r(E_2^{pq})=Z_s(E_2^{pq})$ para $s\geq r\geq r(n)$.
\end{enumerate}
\end{lemma}

\begin{proof}
\oldpage[0\textsubscript{III}]{27}
Según las condiciones \emph{(a), (b) y (c)} de \sref{0.11.1.1}, decir que $Z_{r+1}(E_2^{pq})=Z_r(E_2^{pq})$ equivale a decir que $d_r^{pq}=0$, y decir que $B_r(E_2^{p+r,q-r+1})=B_{r+1}(E_2^{p+r,q-r+1})$ equivale a decir que $d_r^{pq}=0$; el lema se deduce inmediatamente de esta observación.
\end{proof}

\subsection{La sucesión espectral de un complejo filtrado}
\label{subsection:0.11.2}
\label{subsection:0.11.2-fr}

\begin{env}[11.2.1]
\label{0.11.2.1}
\label{0.11.2.1-fr}
Sea $\boldsymbol{C}$ una categoría abeliana; convenimos en denotar mediante notaciones como $K^\bullet$ los \emph{complejos $(K^i)_{i\in\mathbb Z}$} de objetos de $\boldsymbol{C}$ cuyo diferencial es de grado $+1$, y mediante notaciones como $K_\bullet$ los complejos $(K_i)_{i\in\mathbb Z}$ de objetos de $\boldsymbol{C}$ cuyo diferencial es de grado $-1$.
A todo complejo $K^\bullet=(K^i)$ cuyo diferencial $d$ es de grado $+1$ se puede asociarle un complejo $K'_\bullet=(K'_i)$ poniendo $K'_i=K^{-i}$, siendo el diferencial $K'_i\to K'_{i-1}$ el operador $d:K^{-i}\to K^{-i+1}$; y \emph{viceversa}, lo que permitirá, según las circunstancias, considerar uno u otro de los tipos de complejos y traducir todo resultado de uno de los tipos en resultados para el otro.
Denotaremos análogamente mediante notaciones como $K^{\bullet\bullet}=(K^{ij})$ (resp. $K_{\bullet\bullet}=(K_{ij})$) los \emph{bicomplejos} (o \emph{complejos dobles}) de objetos de $\boldsymbol{C}$ en los que los \emph{dos} diferenciales son de grado $+1$ (resp. $-1$); todavía se puede pasar de un tipo al otro cambiando los signos de los índices, y se tienen notaciones y observaciones análogas para los multicomplejos cualesquiera.
Las notaciones $K^\bullet$ y $K_\bullet$ se utilizará también para \emph{objetos graduados} de $\boldsymbol{C}$, de tipo $\mathbb Z$, que no son necesariamente complejos (o que se pueden considerar como tales para operadores de derivación \emph{nulos}); por ejemplo, escribiremos
$\mathrm H^\bullet(K^\bullet)=(\mathrm H^i(K^\bullet))_{i\in\mathbb Z}$ para la \emph{cohomología} de un complejo $K^\bullet$ cuyo diferencial es de grado $+1$, y $\mathrm H_\bullet(K_\bullet)=(\mathrm H_i(K_\bullet))_{i\in\mathbb Z}$ para la \emph{homología} de un complejo $K_\bullet$ cuyo diferencial es de grado $-1$; cuando se pasa de $K^\bullet$ a $K'_\bullet$ por el procedimiento descrito anteriormente, se tiene $\mathrm H_i(K'_\bullet)=\mathrm H^{-i}(K^\bullet)$.

Recordemos en este caso que, para un complejo $K^\bullet$ (resp. $K_\bullet$), escribiremos en general $Z^i(K^\bullet)=\operatorname{Ker}(K^i\to K^{i+1})$ (``objeto de cociclos'') y $B^i(K^\bullet)=\operatorname{Im}(K^{i-1}\to K^i)$ (``objeto de cobordes'') (resp. $Z_i(K_\bullet)=\operatorname{Ker}(K_i\to K_{i-1})$ (``objeto de ciclos'') y $B_i(K_\bullet)=\operatorname{Im}(K_{i+1}\to K_i))$ (``objeto de bordes'')), de modo que $\mathrm H^i(K^\bullet)=Z^i(K^\bullet)/B^i(K^\bullet)$ (resp. $\mathrm H_i(K_\bullet)=Z_i(K_\bullet)/B_i(K_\bullet)$).

Si $K^\bullet=(K^i)$ (resp. $K_\bullet=(K_i)$) es un complejo en $\boldsymbol{C}$ y $T:\boldsymbol{C}\to\boldsymbol{C}'$ un funtor de $\boldsymbol{C}$ en una categoría abeliana $\boldsymbol{C}'$, denotaremos por $T(K^\bullet)$ (resp. $T(K_\bullet)$) el complejo $(T(K^i))$ (resp. $(T(K_i))$) en $\boldsymbol{C}'$.

No repasaremos la definición de los \emph{$\partial$-funtores} (T, 2.1), salvo para observar que diremos \emph{también} $\partial$-funtor en lugar de $\partial^*$-funtor cuando el morfismo $\partial$ disminuye el grado en una unidad, aclarando el contexto este punto si puede haber confusión.

Finalmente, diremos que un \emph{objeto graduado $(A_i)_{i\in\mathbb Z}$} de $\boldsymbol{C}$ está \emph{acotado inferiormente} (resp. \emph{superiormente}) si existe un $i_0$ tal que $A_i=0$ para $i<i_0$ (resp. $i>i_0$).
\end{env}

\begin{env}[11.2.2]
\label{0.11.2.2}
\label{0.11.2.2-fr}
Sea $K^\bullet$ un complejo en $\boldsymbol{C}$ cuyo diferencial $d$ es de grado $+1$, y supongamos que está provisto de una \emph{filtración $F(K^\bullet)=(F^p(K^\bullet))_{p\in\mathbb Z}$} formada por subobjetos \emph{graduados}
\oldpage[0\textsubscript{III}]{28}
de $K^\bullet$, es decir, $F^p(K^\bullet)=(K^i\cap F^p(K^\bullet))_{i\in\mathbb Z}$; se supone además que $d(F^p(K^\bullet))\subset F^p(K^\bullet)$ para todo $p\in\mathbb Z$.
Recordemos rápidamente cómo se define \emph{funtorialmente} una sucesión espectral $E(K^\bullet)$ a partir de $K^\bullet$ (M,~XV,~4 y G,~I,~4.3).
Para $r\geq2$, el morfismo canónico $F^p(K^\bullet)/F^{p+r}(K^\bullet)\to
F^p(K^\bullet)/F^{p+1}(K^\bullet)$ define un morfismo en cohomología
\[
  \mathrm H^{p+q}(F^p(K^\bullet)/F^{p+r}(K^\bullet))
  \longrightarrow
  \mathrm H^{p+q}(F^p(K^\bullet)/F^{p+1}(K^\bullet))
\]

Se denota por $Z_r^{pq}(K^\bullet)$ la imagen de este morfismo.
Análogamente, de la sucesión exacta
\[
  0\to F^p(K^\bullet)/F^{p+1}(K^\bullet)
  \to F^{p-r+1}(K^\bullet)/F^{p+1}(K^\bullet)
  \to F^{p-r+1}(K^\bullet)/F^p(K^\bullet)\to0
\]
se deduce de la sucesión exacta de cohomología un morfismo
\[
  \mathrm H^{p+q-1}(F^{p-r+1}(K^\bullet)/F^p(K^\bullet))
  \longrightarrow
  \mathrm H^{p+q}(F^p(K^\bullet)/F^{p+1}(K^\bullet))
\]
y se denota por $B_r^{pq}(K^\bullet)$ la imagen de este morfismo; se demuestra que $B_r^{pq}(K^\bullet)\subset Z_r^{pq}(K^\bullet)$ y se toma $E_r^{pq}(K^\bullet)=Z_r^{pq}(K^\bullet)/B_r^{pq}(K^\bullet)$; no precisaremos la definición de los $d_r^{pq}$ ni de los $\alpha_r^{pq}$. Se observará aquí que todos los $Z_r^{pq}(K^\bullet)$ y $B_p^{pq}(K^\bullet)$, para $p$ y $q$ fijos, son subobjetos del mismo objeto $\mathrm H^{p+q}(F^p(K^\bullet)/F^{p+1}(K^\bullet))$, que se denota por $Z_1^{pq}(K^\bullet)$; se pone $B_1^{pq}(K^\bullet)=0$, de modo que las definiciones anteriores de $Z_r^{pq}(K^\bullet)$ y $B_r^{pq}(K^\bullet)$ valen también para $r=1$; se pone asimismo $E_1^{pq}(K^\bullet)=Z_1^{pq}(K^\bullet)$. Se definen también los $d_1^{pq}$ y $\alpha_1^{pq}$ de modo que se satisfagan las condiciones de \sref{0.11.1.1} para $r=1$.
Por otra parte, se definen los subobjetos $Z_\infty^{pq}(K^\bullet)$ como la imagen del morfismo
\[
  \mathrm H^{p+q}(F^p(K^\bullet))\longrightarrow
  \mathrm H^{p+q}(F^p(K^\bullet)/F^{p+1}(K^\bullet))
  =E_1^{pq}(K^\bullet)
\]
y $B_\infty^{pq}(K^\bullet)$ como la imagen del morfismo
\[
  \mathrm H^{p+q-1}(K^\bullet/F^p(K^\bullet))\longrightarrow
  \mathrm H^{p+q}(F^p(K^\bullet)/F^{p+1}(K^\bullet))
  =E_1^{pq}(K^\bullet)
\] deducido como arriba de una sucesión exacta de cohomología. Se toma como $Z_\infty(E_2^{pq}(K^\bullet))$ y $B_\infty(E_2^{pq}(K^\bullet))$ como las imágenes canónicas en $E_2^{pq}(K^\bullet)$ de $Z_\infty^{pq}(K^\bullet)$ y $B_\infty^{pq}(K^\bullet)$.

Finalmente, se denota por $F^p(\mathrm H^n(K^\bullet))$ la imagen en $\mathrm H^n(K^\bullet)$ del morfismo $\mathrm H^n(F^p(K^\bullet))\to\mathrm H^n(K^\bullet)$ inducido por la inyección canónica $F^p(K^\bullet)\to K^\bullet$; por la sucesión exacta de cohomología, este es también el núcleo del morfismo $\mathrm H^n(K^\bullet)\to\mathrm H^n(K^\bullet/F^p(K^\bullet))$.
Esto define una filtración sobre $E^n(K^\bullet)=\mathrm H^n(K^\bullet)$; no daremos aquí la definición de los isomorfismos $\beta^{pq}$.
\end{env}

\begin{env}[11.2.3]
\label{0.11.2.3}
\label{0.11.2.3-fr}
El carácter \emph{funtorial} de $E(K^\bullet)$ se entiende del modo siguiente: dados dos complejos \emph{filtrados} $K^\bullet$ y $K^{\prime\bullet}$ en $\boldsymbol{C}$ y un morfismo de complejos $u:K^\bullet\to K^{\prime\bullet}$ \emph{compatible con las filtraciones}, se inducen de manera evidente los morfismos $u_r^{pq}$ (para $r\geq1$) y $u^n$, y se demuestra que estos morfismos son compatibles con los $d_r^{pq}$, $\alpha_r^{pq}$ y $\beta^{pq}$ en el sentido de \sref{0.11.1.2}, y definen así un morfismo bien definido $E(u):E(K^\bullet)\to E(K^{\prime\bullet})$ de sucesiones espectrales.
Además, se demuestra que si $u$ y $v$ son morfismos $K^\bullet\to K^{\prime\bullet}$ del tipo anterior, \emph{homótopos de orden $\leq k$}, entonces $u_r^{pq}=v_r^{pq}$ para $r>k$ y $u^n=v^n$ para todo $n$ (M,~XV,~3.1).
\end{env}

\begin{env}[11.2.4]
\label{0.11.2.4-fr}
\oldpage[0\textsubscript{III}]{29}
Supongamos que en $\boldsymbol{C}$ los límites inductivos filtrantes sean exactos. Entonces, si la filtración $(F^p(K^\bullet))$ de $K^\bullet$ es \emph{exhaustiva}, también lo es la filtración $(F^p(\mathrm H^n(K^\bullet)))$ para todo $n$, pues por hipótesis se tiene $K^\bullet=\varinjlim_{p\to-\infty}F^p(K^\bullet)$ y la hipótesis sobre $\boldsymbol{C}$ implica que la cohomología conmuta con los límites inductivos.
Además, por la misma razón, se tiene $B_\infty(E_2^{pq}(K^\bullet))=\sup_k B_k(E_2^{pq}(K^\bullet))$.
Se dice que la filtración $(F^p(K^\bullet))$ de $K^\bullet$ es \emph{regular} si para todo $n$ existe un entero $u(n)$ tal que $\mathrm H^n(F^p(K^\bullet))=0$ para $p>u(n)$.
Este es, en particular, el caso cuando la filtración de $K^\bullet$ es \emph{discreta}.
Cuando la filtración de $K^\bullet$ es regular y exhaustiva, y los límites inductivos filtrantes son exactos en $\boldsymbol{C}$, se tiene (M,~XV,~4) que la sucesión espectral $E(K^\bullet)$ es \emph{regular}.
\end{env}

\subsection{Las sucesiones espectrales de un bicomplejo}
\label{subsection:0.11.3}
\label{subsection:0.11.3-fr}

\begin{env}[11.3.1]
\label{0.11.3.1}
\label{0.11.3.1-fr}
En cuanto a las convenciones para los bicomplejos, seguimos las de~(T,~2.4) antes que las de~(M), suponiéndose así que los dos diferenciales $d'$ y $d''$ (de grado~$+1$) de un tal bicomplejo $K^{\bullet\bullet}=(K^{ij})$ se suponen, pues, \emph{permutables}. Supongamos satisfecha \emph{una} de las dos condiciones siguientes: $1^\circ$ Las \emph{sumas directas infinitas} existen en $\boldsymbol{C}$; $2^\circ$ Para todo $n\in\mathbb Z$, solo hay un número \emph{finito} de pares $(p,q)$ tales que $p+q=n$ y $K^{pq}\ne0$.
 
Entonces el bicomplejo $K^{\bullet\bullet}$ define un \emph{complejo} (simple) $(K^{\prime n})_{n\in\mathbb Z}$ con $K^{\prime n}=\sum_{i+j=n}K^{ij}$, viniendo dado el diferencial $d$ (de grado~$+1$) de este complejo por $dx=d'x+(-1)^i d''x$ para $x\in K^{ij}$.
Cuando hablemos más adelante del \emph{complejo (simple) definido por un bicomplejo $K^{\bullet\bullet}$}, se sobreentenderá siempre que se satisface una de las condiciones anteriores.
Se adoptarán convenciones análogas para los multicomplejos.

Se denota por $K^{i,\bullet}$ (resp.~$K^{\bullet,j}$) el complejo simple $(K^{ij})_{j\in\mathbb Z}$ (resp.~$(K^{ij})_{i\in\mathbb Z}$), y por $Z_{\mathrm{II}}^p(K^{i,\bullet})$, $B_{\mathrm{II}}^p(K^{i,\bullet})$, $\mathrm H_{\mathrm{II}}^p(K^{i,\bullet})$ (resp.~$Z_{\mathrm I}^p(K^{\bullet,j})$, $B_{\mathrm I}^p(K^{\bullet,j})$, $\mathrm H_{\mathrm I}^p(K^{\bullet,j})$) sus objetos de cociclos, de cobordes y de cohomología de grado $p$, respectivamente;
el diferencial $d':K^{i,\bullet}\to K^{i+1,\bullet}$ es un morfismo de complejos, que da así un operador sobre los cociclos, los cobordes y la cohomología,
\[
\begin{aligned}
  d'&:Z_{\mathrm{II}}^p(K^{i,\bullet})
      \longrightarrow Z_{\mathrm{II}}^p(K^{i+1,\bullet})\\
  d'&:B_{\mathrm{II}}^p(K^{i,\bullet})
      \longrightarrow B_{\mathrm{II}}^p(K^{i+1,\bullet})\\
  d'&:\mathrm H_{\mathrm{II}}^p(K^{i,\bullet})
      \longrightarrow \mathrm H_{\mathrm{II}}^p(K^{i+1,\bullet})
\end{aligned}
\]
y está claro que, para estos operadores, $(Z_{\mathrm{II}}^p(K^{i,\bullet}))_{i\in\mathbb Z}$, $(B_{\mathrm{II}}^p(K^{i,\bullet}))_{i\in\mathbb Z}$ y $(\mathrm H_{\mathrm{II}}^p(K^{i,\bullet}))_{i\in\mathbb Z}$ son complejos;
denotaremos el complejo $(\mathrm H_{\mathrm{II}}^p(K^{i,\bullet}))_{i\in\mathbb Z}$ por $\mathrm H_{\mathrm{II}}^p(K^{\bullet\bullet})$, y sus objetos de cociclos, cobordes y cohomología de grado $q$ por $Z_{\mathrm I}^q(\mathrm H_{\mathrm{II}}^p(K^{\bullet\bullet}))$, $B_{\mathrm I}^q(\mathrm H_{\mathrm{II}}^p(K^{\bullet\bullet}))$ y $\mathrm H_{\mathrm I}^q(\mathrm H_{\mathrm{II}}^p(K^{\bullet\bullet}))$.
Se definen análogamente los complejos $\mathrm H_{\mathrm I}^p(K^{\bullet\bullet})$ y sus objetos de cohomología $\mathrm H_{\mathrm{II}}^q(\mathrm H_{\mathrm I}^p(K^{\bullet\bullet}))$.
Recordemos, no obstante, que $\mathrm H^n(K^{\bullet\bullet})$ denota el objeto de grado $n$ de la cohomología del complejo \emph{(simple)} definido por $K^{\bullet\bullet}$.
\end{env}

\begin{env}[11.3.2]
\label{0.11.3.2}
\label{0.11.3.2-fr}
Sobre el complejo definido por un bicomplejo $K^{\bullet\bullet}$ pueden considerarse dos filtraciones canónicas, $(F_{\mathrm I}^p(K^{\bullet\bullet}))$ y $(F_{\mathrm{II}}^p(K^{\bullet\bullet}))$, dadas por
\[
\label{0.11.3.2.1-fr}
  F_{\mathrm I}^p(K^{\bullet\bullet})
  =\left(\sum_{i+j=n,\,i\geq p}K^{ij}\right)_{n\in\mathbb Z},
  \qquad
  F_{\mathrm{II}}^p(K^{\bullet\bullet})
  =\left(\sum_{i+j=n,\,j\geq p}K^{ij}\right)_{n\in\mathbb Z}
  \tag{11.3.2.1}
\]
\oldpage[0\textsubscript{III}]{30}
que, por definición, son subobjetos graduados del complejo (simple) definido por $K^{\bullet\bullet}$ y hacen así de este complejo un complejo filtrado;
además, está claro que estas filtraciones son \emph{exhaustivas} y \emph{separadas}.

A cada una de estas filtraciones le corresponde una sucesión espectral~\sref{0.11.2.2};
denotaremos por $'E(K^{\bullet\bullet})$ y $''E(K^{\bullet\bullet})$ las sucesiones espectrales correspondientes a $(F_{\mathrm I}^p(K^{\bullet\bullet}))$ y $(F_{\mathrm{II}}^p(K^{\bullet\bullet}))$, respectivamente, llamadas \emph{sucesiones espectrales del bicomplejo $K^{\bullet\bullet}$}, y que tienen ambas por término límite la cohomología $(\mathrm H^n(K^{\bullet\bullet}))$.
Se demuestra además~(M,~XV,~6) que se tiene
\[
\label{0.11.3.2.2-fr}
  'E_2^{pq}(K^{\bullet\bullet})
  =\mathrm H_{\mathrm I}^p(\mathrm H_{\mathrm{II}}^p(K^{\bullet\bullet})),
  \qquad
  ''E_2^{pq}(K^{\bullet\bullet})
  =\mathrm H_{\mathrm{II}}^p(\mathrm H_{\mathrm I}^p(K^{\bullet\bullet})).
  \tag{11.3.2.2}
\]

Todo morfismo $u:K^{\bullet\bullet}\to K^{\prime\bullet\bullet}$ de bicomplejos es \emph{ipso facto} compatible con las filtraciones del mismo tipo de $K^{\bullet\bullet}$ y $K^{\prime\bullet\bullet}$, y define así un morfismo para cada una de las dos sucesiones espectrales;
además, dos morfismos \emph{homótopos} definen una homotopía \emph{de orden $\leq1$} de los complejos (simples) filtrados correspondientes, y por tanto el \emph{mismo} morfismo para cada una de las dos sucesiones espectrales~(M,~XV,~6.1).
\end{env}

\begin{proposition}[11.3.3]
\label{0.11.3.3}
\label{0.11.3.3-fr}
Sea $K^{\bullet\bullet}=(K^{ij})$ un bicomplejo en una categoría abeliana $\boldsymbol{C}$.
\begin{enumerate}
  \item[{\rm(i)}] Si existen $i_0$ y $j_0$ tales que $K^{ij}=0$ para $i<i_0$ o $j<j_0$ \textup{(resp.}~$i>i_0$ o $j>j_0$), entonces las dos sucesiones espectrales $'E(K^{\bullet\bullet})$ y $''E(K^{\bullet\bullet})$ son birregulares.
  \item[{\rm(ii)}] Si existen $i_0$ e $i_1$ tales que $K^{ij}=0$ para $i<i_0$ o $i>i_1$ \textup{(resp.}~si existen $j_0$ y $j_1$ tales que $K^{ij}=0$ para $j<j_0$ o $j>j_1$), entonces las dos sucesiones espectrales $'E(K^{\bullet\bullet})$ y $''E(K^{\bullet\bullet})$ son birregulares.
\end{enumerate}
Supongamos además que en $\boldsymbol{C}$ los límites inductivos filtrantes existen y son exactos. Entonces:
\begin{enumerate}
\item[\textup{(iii)}] Si existe $i_0$ tal que $K^{ij}=0$ para $i>i_0$ \textup{(resp.}~si existe $j_0$ tal que $K^{ij}=0$ para $j<j_0$), entonces la sucesión espectral $'E(K^{\bullet\bullet})$ es regular.
  \item[{\rm(iv)}] Si existe $i_0$ tal que $K^{ij}=0$ para $i<i_0$ \textup{(resp.}~si existe $j_0$ tal que $K^{ij}=0$ para $j>j_0$), entonces la sucesión espectral $''E(K^{\bullet\bullet})$ es regular.
\end{enumerate}
\end{proposition}

\begin{proof}
La proposición se deduce inmediatamente de las definiciones~\sref{0.11.1.3} y de~\sref{0.11.2.4}, así como de las observaciones siguientes relativas a la filtración $F_{\mathrm I}$ (y de las observaciones análogas que se deducen para $F_{\mathrm{II}}$ intercambiando los papeles de los dos índices en $K^{\bullet\bullet}$):
\begin{enumerate}
  \item[{$1^\circ$}] Si existe $i_0$ tal que $K^{ij}=0$ para $i>i_0$, entonces la filtración $F_{\mathrm I}(K^{\bullet\bullet})$ es \emph{discreta}.
  \item[{$2^\circ$}] Si existe $i_0$ tal que $K^{ij}=0$ para $i<i_0$, entonces la filtración $F_{\mathrm I}(K^{\bullet\bullet})$ es \emph{codiscreta}.
    Se deduce inmediatamente que ocurre lo mismo para la filtración correspondiente $F_{\mathrm I}(\mathrm H^n(K^{\bullet\bullet}))$ para todo $n$.
    Además, la definición de los $B_r^{pq}$ correspondientes a la filtración $F_{\mathrm I}(K^{\bullet\bullet})$ \sref{0.11.2.2} muestra que, para todo par $(p,q)$, la sucesión $(B_r^{pq})_{r\geq2}$ es estable.
  \item[{$3^\circ$}] Si existe $j_0$ tal que $K^{ij}=0$ para $j<j_0$, entonces se tiene
    \[
      F_{\mathrm I}^{p+r}(K^{\bullet\bullet})
      \cap\left(\sum_{i+j=n}K^{ij}\right)=0
    \]
    desde que $p+r+j_0>n$, luego $Z_r^{pq}=Z_\infty(E_2^{pq})$ para $r>q-j_0+1$.
    Por otra parte, $\mathrm H^n(F_{\mathrm I}^p(K^{\bullet\bullet}))=0$ para $p>n-j_0+1$.
  \item[{$4^\circ$}] Si existe $j_0$ tal que $K^{ij}=0$ para $j>j_0$, entonces se tiene
    \[
      F_{\mathrm I}^{p-r+1}(K^{\bullet\bullet})
      \cap\left(\sum_{i+j=n}K^{ij}\right)=\sum_{i+j=n}K^{ij}
    \]
    \oldpage[0\textsubscript{III}]{31}
    desde que $p-r+1+j_0<n$, luego $B_r^{pq}=B_\infty(E_2^{pq})$ para $r<j_0-q+1$.
    Por otra parte, $\mathrm H^n(F_{\mathrm I}^p(K^{\bullet\bullet}))
    =\mathrm H^n(K^{\bullet\bullet})$ para $p+j_0<n-1$.
\end{enumerate}
\end{proof}

\begin{env}[11.3.4]
\label{0.11.3.4-fr}
Supongamos que el bicomplejo $K^{\bullet\bullet}=(K^{ij})$ sea tal que $K^{ij}=0$ para $i<0$ o $j<0$. Se sabe que se puede entonces definir para todo $p\in\mathbb Z$ un «homomorfismo de borde» canónico
\[
\label{0.11.3.4.1-fr}
  'E_2^{p0}(K^{\bullet\bullet})\longrightarrow
  \mathrm H^p(K^{\bullet\bullet})
  \tag{11.3.4.1}
\]
(M, XV, 6). Recordemos brevemente que esto se debe, por una parte, a que se tiene en la sucesión espectral $'E(K^{\bullet\bullet})$, $Z_r^{p0}=Z_{\mathrm I}^p(Z_{\mathrm{II}}^0(K^{\bullet\bullet}))$ para $2\leq r\leq+\infty$, y por otra al hecho de que $\mathrm H^p(F_{\mathrm I}^{p+1}(K^{\bullet\bullet}))=0$, de modo que el isomorfismo $\beta^{p0}:'E_\infty^{p0}\xrightarrow{\sim}
\mathrm H^p(F_{\mathrm I}^p)/\mathrm H^p(F_{\mathrm I}^{p+1})$ da un homomorfismo $'E_\infty^{p0}\to\mathrm H^p(F_{\mathrm I}^p(K^{\bullet\bullet}))
\to\mathrm H^p(K^{\bullet\bullet})$; la igualdad de todos los $Z_r^{p0}$ permite entonces definir homomorfismos canónicos $'E_r^{p0}\to'E_s^{p0}$ para $r\leq s$, y en particular un homomorfismo $'E_2^{p0}\to'E_\infty^{p0}$, de donde, por composición, el homomorfismo de borde $'E_2^{p0}\to\mathrm H^p(K^{\bullet\bullet})$; además, se verifica inmediatamente que, a la clase mód. $B_2^{p0}$ de un elemento $z\in Z_{\mathrm{II}}^0(K^{\bullet\bullet})\subset K^{p0}$ tal que $d'z=0$, el homomorfismo de borde así definido hace corresponder en $'E_\infty^{p0}$ la clase de $z$ mód. $B_\infty^{p0}$, y después a esta última la clase de cohomología de $z$ en $\mathrm H^p(K^{\bullet\bullet})$. Se ve, pues, finalmente que el \emph{homomorfismo de borde} (\hyperref[0.11.3.4.1-fr]{11.3.4.1}) \emph{proviene, por paso a la cohomología, de la inyección canónica} $Z_{\mathrm{II}}^0(K^{\bullet\bullet})\to K^{\bullet\bullet}$ (donde $K^{\bullet\bullet}$ se considera como complejo simple). Naturalmente, se interpreta del mismo modo el homomorfismo de borde
\[
\label{0.11.3.4.2-fr}
  ''E_2^{p0}(K^{\bullet\bullet})\longrightarrow
  \mathrm H^p(K^{\bullet\bullet})
  \tag{11.3.4.2}
\]
como proveniente de la inyección canónica $Z_{\mathrm I}^0(K^{\bullet\bullet})\to K^{\bullet\bullet}$.
\end{env}

\begin{env}[11.3.5]
\label{0.11.3.5-fr}
Sea ahora $K_{\bullet\bullet}=(K_{ij})$ un bicomplejo de
$\boldsymbol{C}$ cuyos dos operadores de derivación son de grado $-1$. Escribiremos
entonces $K_{i,\bullet}$ (resp. $K_{\bullet,j}$) para el complejo simple
$(K_{ij})_{j\in\mathbb Z}$ (resp. $(K_{ij})_{i\in\mathbb Z}$),
$\mathrm H_p^{\mathrm{II}}(K_{i,\bullet})$ (resp.
$\mathrm H_p^{\mathrm I}(K_{\bullet,j})$) para su objeto de homología de grado $p$,
$\mathrm H_p^{\mathrm{II}}(K_{\bullet\bullet})$ (resp.
$\mathrm H_p^{\mathrm I}(K_{\bullet\bullet})$) para el complejo
$(\mathrm H_p^{\mathrm{II}}(K_{i,\bullet}))_{i\in\mathbb Z}$ (resp.
$(\mathrm H_p^{\mathrm I}(K_{\bullet,j}))_{j\in\mathbb Z}$),
$\mathrm H_q^{\mathrm I}(\mathrm H_p^{\mathrm{II}}(K_{\bullet\bullet}))$
(resp.
$\mathrm H_q^{\mathrm{II}}(\mathrm H_p^{\mathrm I}(K_{\bullet\bullet}))$)
para su objeto de homología de grado $q$; notaciones análogas para los objetos de ciclos
y objetos de bordes; finalmente, $\mathrm H_n(K_{\bullet\bullet})$ designará
(cuando exista) el objeto de homología de grado $n$ del complejo simple (con
operador de derivación de grado $-1$) definido por $K_{\bullet\bullet}$.

Sea $K^{\prime\bullet\bullet}=(K^{\prime ij})$ con
$K^{\prime ij}=K_{-i,-j}$ el bicomplejo con operadores de derivación de grados
$+1$ asociado a $K_{\bullet\bullet}$; por definición, las
\emph{sucesiones espectrales} de $K_{\bullet\bullet}$ son las de
$K^{\prime\bullet\bullet}$, que se escriben $'E(K_{\bullet\bullet})$ y
$''E(K_{\bullet\bullet})$, donde sin embargo se cambian las notaciones, poniendo
\[
  'E_{pq}^r(K_{\bullet\bullet})
  ='E_r^{-p,-q}(K^{\prime\bullet\bullet}),\qquad
  ''E_{pq}^r(K_{\bullet\bullet})
  =''E_r^{-p,-q}(K^{\prime\bullet\bullet})
\]
para $2\leq r\leq\infty$. Con estas notaciones, se tiene
\[
  'E_{pq}^2(K_{\bullet\bullet})
  =\mathrm H_p^{\mathrm I}(\mathrm H_q^{\mathrm{II}}(K_{\bullet\bullet})),
  \qquad
  ''E_{pq}^2(K_{\bullet\bullet})
  =\mathrm H_p^{\mathrm{II}}(\mathrm H_p^{\mathrm I}(K_{\bullet\bullet})).
\]
Para evitar errores de signo, será en general preferible, para las
relaciones entre estas sucesiones espectrales y su término límite, volver al
complejo $K^{\prime\bullet\bullet}$. Señalemos no obstante los criterios
correspondientes a (\hyperref[0.11.3.3-fr]{11.3.3}):
\end{env}

\begin{env}[11.3.6]
\label{0.11.3.6}
\label{0.11.3.6-fr}
Las sucesiones espectrales $'E(K_{\bullet\bullet})$ y $''E(K_{\bullet\bullet})$ son \emph{birregulares} en los casos siguientes:
\begin{enumerate}
  \item[\textit{a)}] Existen $i_0$ y $j_0$ tales que $K_{ij}=0$ para $i>i_0$ o $j>j_0$ (resp. para $i<i_0$
    \oldpage[0\textsubscript{III}]{32}
    o $j<j_0$);
  \item[\textit{b)}] Existen $i_0$ e $i_1$ tales que $K_{ij}=0$ para $i<i_0$ y $i>i_1$;
  \item[\textit{c)}] Existen $j_0$ y $j_1$ tales que $K_{ij}=0$ para $j<j_0$ y $j>j_1$.
\end{enumerate}

La sucesión $'E(K_{\bullet\bullet})$ es \emph{regular} si existe $i_0$ tal que $K_{ij}=0$ para $i<i_0$, o si existe $j_0$ tal que $K_{ij}=0$ para $j>j_0$.

La sucesión $''E(K_{\bullet\bullet})$ es \emph{regular} si existe $i_0$ tal que $K_{ij}=0$ para $i>i_0$, o si existe $j_0$ tal que $K_{ij}=0$ para $j<j_0$.
\end{env}


\subsection{Hipercohomología de un funtor con respecto a un complejo $K^\bullet$}
\label{subsection:0.11.4}
\label{subsection:0.11.4-fr}

\begin{env}[11.4.1]
\label{0.11.4.1}
\label{0.11.4.1-fr}
Sea $\boldsymbol{C}$ una categoría abeliana.
Recordemos que se ha definido una \emph{resolución por la derecha} (o \emph{resolución cohomológica}) de un objeto $A$ de $\boldsymbol{C}$ como un complejo de objetos de $\boldsymbol{C}$, cuyo operador diferencial es de grado $+1$,
\[
  0\longrightarrow L^0\longrightarrow L^1\longrightarrow L^2
  \ldots
\]
provisto de un morfismo $\varepsilon:A\to L^0$, llamado la \emph{aumentación} de la resolución, que puede considerarse como un morfismo de complejos
\[
\xymatrix{
  0\ar[r] & A\ar[r]\ar[d]_{\varepsilon} & 0\ar[r]\ar[d]
    & 0\ar[r]\ar[d] & \ldots\\
  0\ar[r] & L^0\ar[r] & L^1\ar[r] & L^2\ar[r] & \ldots
})
\]
tal que la sucesión
\[
  0\longrightarrow A\xrightarrow{\varepsilon}L^0\longrightarrow L^1
  \longrightarrow\ldots
\]
sea \emph{exacta}.
Análogamente, una \emph{resolución por la izquierda} (o \emph{resolución homológica}) de $A$ es un complejo $0\leftarrow L_0\leftarrow L_1\leftarrow\ldots$ de objetos de $\boldsymbol{C}$, cuyo operador diferencial es de grado~$-1$, provisto de una aumentación $\varepsilon:L_0\to A$, de modo que la sucesión
\[
  0\longleftarrow A\xleftarrow{\varepsilon}L_0\longleftarrow L_1
  \longleftarrow\ldots
\]
sea \emph{exacta}.

Cuando la resolución por la derecha $(L_i)_{i\geq0}$ de un objeto $A$ es tal que $L_i=0$ para $i\geq n+1$, se dice que esta resolución es \emph{de longitud $\leq n$}.
Se define análogamente una resolución por la izquierda de longitud $\leq n$.
Una resolución que es de longitud $\leq n$ para algún entero $n$ se dice \emph{finita}.

Una resolución de $A$ se llama \emph{proyectiva} (resp. \emph{inyectiva}) si los objetos de $\boldsymbol{C}$ distintos de $A$ que la componen son \emph{proyectivos} (resp. \emph{inyectivos}).
Cuando $\boldsymbol{C}$ es la categoría de los módulos (a izquierda, por ejemplo) sobre un anillo, se dice que una resolución de $A$ es \emph{plana} (resp. \emph{libre}) cuando los módulos distintos de $A$ que la componen son \emph{planos} (resp. \emph{libres}).
\end{env}

\begin{env}[11.4.2]
\label{0.11.4.2}
\label{0.11.4.2-fr}
Sea $K^\bullet=(K^i)_{i\in\mathbb Z}$ un complejo de objetos de $\boldsymbol{C}$, cuyo operador diferencial es de grado~$+1$.
Se define una \emph{resolución de Cartan--Eilenberg por la derecha} de $K^\bullet$ como un par formado por un bicomplejo $L^{\bullet\bullet}=(L^{ij})$ con operadores diferenciales de grado~$+1$ y tal que $L^{ij}=0$ para $j<0$, y un morfismo de complejos simples $\varepsilon:K^\bullet\to L^{\bullet,0}$, de modo que se cumplan las condiciones siguientes:
\oldpage[0\textsubscript{III}]{33}
\begin{enumerate}
  \item[(i)] Para cada índice $i$, las sucesiones
    \[
    \begin{aligned}
      0\longrightarrow{}&K^i\xrightarrow{\varepsilon}L^{i0}
        \longrightarrow L^{i1}\longrightarrow\ldots\\
      0\longrightarrow{}&\mathrm B^i(K^\bullet)\xrightarrow{\varepsilon}
        \mathrm B_{\mathrm I}^i(L^{\bullet,0})
        \longrightarrow\mathrm B_{\mathrm I}^i(L^{\bullet,1})
        \longrightarrow\ldots\\
      0\longrightarrow{}&\mathrm Z^i(K^\bullet)\xrightarrow{\varepsilon}
        \mathrm Z_{\mathrm I}^i(L^{\bullet,0})
        \longrightarrow\mathrm Z_{\mathrm I}^i(L^{\bullet,1})
        \longrightarrow\ldots\\
      0\longrightarrow{}&\mathrm H^i(K^\bullet)\xrightarrow{\varepsilon}
        \mathrm H_{\mathrm I}^i(L^{\bullet,0})
        \longrightarrow\mathrm H_{\mathrm I}^i(L^{\bullet,1})
        \longrightarrow\ldots
    \end{aligned}
    \]
    son exactas, es decir, $(L^{i,\bullet})$, $(\mathrm B_{\mathrm I}^i(L^{\bullet\bullet}))$, $(\mathrm Z_{\mathrm I}^i(L^{\bullet\bullet}))$ y
$(\mathrm H_{\mathrm I}^i(L^{\bullet\bullet}))$ son, respectivamente, \emph{resoluciones} de $K^i$, $\mathrm B^i(K^\bullet)$, $\mathrm Z^i(K^\bullet)$ y $\mathrm H^i(K^\bullet)$.
  \item[(ii)] Para cada $j$, el complejo simple $L^{\bullet,j}$ es \emph{escindido} o, en otras palabras, las sucesiones exactas
    \[
      \label{0.11.4.2.1-fr}
      0\longrightarrow\mathrm B_{\mathrm I}^i(L^{\bullet,j})
      \longrightarrow\mathrm Z_{\mathrm I}^i(L^{\bullet,j})
      \longrightarrow\mathrm H_{\mathrm I}^i(L^{\bullet,j})
      \longrightarrow0
      \tag{11.4.2.1}
    \]
    \[
      \label{0.11.4.2.2-fr}
      0\longrightarrow\mathrm Z_{\mathrm I}^i(L^{\bullet,j})
      \longrightarrow L^{ij}
      \longrightarrow\mathrm B_{\mathrm I}^{i+1}(L^{\bullet,j})
      \longrightarrow0
      \tag{11.4.2.2}
    \]
    son \emph{escindidas}.
\end{enumerate}

Se demuestra (M,~XVII,~1.2) que si todo objeto de $\boldsymbol{C}$ es subobjeto de un objeto inyectivo, entonces todo complejo $K^\bullet$ de $\boldsymbol{C}$ admite una resolución de Cartan--Eilenberg \emph{inyectiva}, es decir, formada por objetos inyectivos $L^{ij}$ (la condición (ii) anterior implica entonces que los $\mathrm B_{\mathrm I}^i(L^{\bullet,j})$, $\mathrm Z_{\mathrm I}^i(L^{\bullet,j})$ y
$\mathrm H_{\mathrm I}^i(L^{\bullet,j})$ están también formados por objetos inyectivos).
Además, para todo morfismo $f:K^\bullet\to K'^\bullet$ de complejos de $\boldsymbol{C}$,
toda resolución de Cartan--Eilenberg $L^{\bullet\bullet}$ de $K^\bullet$ y toda resolución de Cartan--Eilenberg \emph{inyectiva} $L'^{\bullet\bullet}$ de $K'^\bullet$, existe un morfismo de bicomplejos $F:L^{\bullet\bullet}\to L'^{\bullet\bullet}$ compatible con $f$ y las aumentaciones, y si $f$ y $g$ son morfismos homótopos de $K^\bullet$ en $K'^\bullet$, entonces los morfismos correspondientes de $L^{\bullet\bullet}$ en $L'^{\bullet\bullet}$ son también homótopos \emph{(loc. cit.)}.

Cuando $K^\bullet$ está \emph{acotado inferiormente} (resp. \emph{superiormente}), puede tomarse $L^{\bullet\bullet}$ tal que $L^{ij}=0$ para $i<i_0$ (resp. $i>i_0$) si $K^i=0$ para $i<i_0$ (resp. $i>i_0$) (M,~XVII,~1.3).

Supongamos, por otra parte, que existe un entero $n$ tal que todo objeto de $\boldsymbol{C}$ admite una \emph{resolución inyectiva de longitud $\leq n$};
entonces puede suponerse que $L^{ij}=0$ para $j>n$ (M,~XVII,~1.4).
\end{env}


\begin{env}[11.4.3]
\label{0.11.4.3}
\label{0.11.4.3-fr}
Sea ahora $T$ un \emph{funtor covariante aditivo} de $\boldsymbol{C}$ en una categoría abeliana $\boldsymbol{C}'$.
Dado un complejo $K^\bullet$ de $\boldsymbol{C}$ y una resolución de Cartan--Eilenberg \emph{inyectiva} $L^{\bullet\bullet}$ de $K^\bullet$, supongamos que existe el complejo (simple) definido por el bicomplejo $T(L^{\bullet\bullet})$ existe
(cf. \hyperref[0.11.3.1-fr]{11.3.1}); entonces las dos sucesiones espectrales
$'E(T(L^{\bullet\bullet}))$ y $''E(T(L^{\bullet\bullet}))$ de este bicomplejo se llaman las \emph{sucesiones espectrales de hipercohomología} de $T$ con respecto a $K^\bullet$; en virtud de (\hyperref[0.11.4.2-fr]{11.4.2}) y (\hyperref[0.11.3.2-fr]{11.3.2}), efectivamente solo dependen de $K^\bullet$ y no de la elección de la resolución de Cartan--Eilenberg inyectiva $L^{\bullet\bullet}$;
además, dependen \emph{funtorialmente} de $K^\bullet$.
Tienen el mismo término límite $\mathrm H^\bullet(T(L^{\bullet\bullet}))$, llamado también la \emph{hipercohomología} de $T$ con respecto a $K^\bullet$, y denotado por $\mathbf R^\bullet T(K^\bullet)$.
Se demuestra que los términos $E_2$ de las dos sucesiones espectrales anteriores vienen dados por
\[
  \label{0.11.4.3.1-fr}
  'E_2^{pq}=\mathrm H^{\boldsymbol{p}}(\mathrm R^qT(K^\bullet))
  \tag{11.4.3.1}
\]
\[
  \label{0.11.4.3.2-fr}
  ''E_2^{pq}=\mathrm R^{\boldsymbol{p}}T(\mathrm H^q(K^\bullet))
  \tag{11.4.3.2}
\]
\oldpage[0\textsubscript{III}]{34}
donde $\mathrm R^pT$ denota, como de costumbre, el \emph{funtor derivado} de grado $p$ de $T$ para $p\in\mathbb Z$;
$\mathrm R^qT(K^\bullet)$ denota el complejo $(\mathrm R^qT(K^i))_{i\in\mathbb Z}$.
\phantomsection\label{0.11.4.4}
\label{0.11.4.4-fr}
Salvo mención explícita en contrario, supondremos en lo sucesivo que todo objeto de $\boldsymbol{C}$ es subobjeto de un objeto inyectivo de $\boldsymbol{C}$, de modo que existen resoluciones de Cartan--Eilenberg inyectivas para todo complejo de $\boldsymbol{C}$.
Como $L^{ij}=0$ para $j<0$, los criterios de \sref{0.11.3.3} muestran que las \emph{dos} sucesiones espectrales de hipercohomología de $T$ con respecto a $K^\bullet$ existen y
son \emph{birregulares} en cada uno de los dos casos siguientes:
\begin{enumerate}
  \item[$1^\circ$] $K^\bullet$ está acotado inferiormente;
  \item[$2^\circ$] Todo objeto de $\boldsymbol{C}$ admite una resolución inyectiva
de longitud a lo sumo igual a un entero $n$ independiente del objeto considerado.
\end{enumerate}
En efecto, puede suponerse, en el primer caso, por \sref{0.11.4.2}, que existe $i_0$ tal que $L^{ij}=0$ para $i<i_0$, y, en el segundo caso, que existe $j_1$ tal que $L^{ij}=0$ para $j>j_1$;
en cada uno de estos dos casos, está además claro que, para cada $n$ dado, solo existe un número finito de pares $(i,j)$ tales que $L^{ij}\neq0$ e $i+j=n$, lo que demuestra nuestras aserciones.

Si suponemos que los límites inductivos filtrantes existen en $\boldsymbol{C}'$ y son exactos (lo que implica, en particular, la existencia de sumas directas infinitas en $\boldsymbol{C}'$), entonces existe el complejo definido por el bicomplejo $T(L^{\bullet\bullet})$, y el criterio \sref{0.11.3.3} muestra que la sucesión $'E(T(L^{\bullet\bullet}))$ es siempre \emph{regular}.

Si $K^\bullet$ es un complejo tal que todos los $K^i$ son nulos salvo un único $K^{i_0}$, entonces
$\mathbf R^nT(K^\bullet)$ es isomorfo a $\mathrm R^{n-i_0}T(K^{i_0})$, como se deduce inmediatamente de las definiciones tomando una resolución de Cartan--Eilenberg $L^{\bullet\bullet}$ tal que $L^{ij}=0$ para $i\neq i_0$.

Si $K^\bullet$ y $K'^\bullet$ son dos complejos de $\boldsymbol{C}$, y $f$ y $g$ son morfismos homótopos de $K^\bullet$ en $K'^\bullet$, entonces los morfismos $\mathbf R^\bullet T(K^\bullet)\to
\mathbf R^\bullet T(K'^\bullet)$ inducidos por $f$ y $g$ son idénticos, y lo mismo ocurre con los morfismos de las sucesiones espectrales de hipercohomología.
\end{env}

\begin{proposition}[11.4.5]
\label{0.11.4.5-fr}
Supongamos que en $\boldsymbol{C}'$ los límites inductivos filtrantes existen y son exactos. Si $\mathrm R^nT(K^i)=0$ para todo $n>0$ y todo $i\in\mathbb Z$, se tienen isomorfismos funtoriales
\[
  \label{0.11.4.5.1-fr}
  \mathbf R^iT(K^\bullet)\xrightarrow{\sim}
  \mathrm H^i(T(K^\bullet))\qquad(i\in\mathbb Z).
  \tag{11.4.5.1}
\]
\end{proposition}

\begin{proof}
En efecto, los únicos términos $E_2$ no nulos de la primera sucesión espectral \sref{0.11.4.3.1} son entonces $'E_2^{p0}=\mathrm H^p(T(K^\bullet))$;
en otras palabras, esta sucesión es \emph{degenerada};
como es regular \sref{0.11.4.4}, la conclusión se deduce de \sref{0.11.1.6}.
\end{proof}

\begin{env}[11.4.6]
\label{0.11.4.6}
\label{0.11.4.6-fr}
Consideremos ahora, por ejemplo, un \emph{bifuntor} covariante $(M,N)\to T(M,N)$ de $\boldsymbol{C}\times\boldsymbol{C}'$ en $\boldsymbol{C}''$, donde $\boldsymbol{C}$, $\boldsymbol{C}'$ y $\boldsymbol{C}''$ son tres categorías abelianas;
se supone, por simplicidad, que $T$ es aditivo en cada uno de sus argumentos y, además, que todo objeto de $\boldsymbol{C}$ y todo objeto de $\boldsymbol{C}'$ son subobjetos de un objeto inyectivo,
y que los límites inductivos filtrantes existen en $\boldsymbol{C}''$ y son exactos.
Se define entonces la \emph{hipercohomología} de $T$ con respecto a los complejos $K^\bullet$ y $K'^\bullet$ de $\boldsymbol{C}$ y $\boldsymbol{C}'$, respectivamente, con operadores diferenciales de grado~$+1$, reemplazando $K^\bullet$ (resp. $K'^\bullet$) por una resolución de Cartan--Eilenberg inyectiva $L^{\bullet\bullet}$ (resp. $L'^{\bullet\bullet}$);
entonces $T(L^{\bullet\bullet},L'^{\bullet\bullet})$ es un cuadricomplejo de $\boldsymbol{C}''$, que se considera como un \emph{bicomplejo} de $\boldsymbol{C}''$, siendo los grados de $T(L^{ij},L'^{hk})$ los enteros $i+h$ y $j+k$.
La \emph{hipercohomología} de $T$ con respecto a $K^\bullet$ y $K'^\bullet$ es, por definición, la cohomología $\mathrm H^\bullet(T(L^{\bullet\bullet},L'^{\bullet\bullet}))$ de este bicomplejo (en otras palabras, la del complejo simple asociado) y se denota por $\mathbf R^\bullet T(K^\bullet,K'^\bullet)$;
\oldpage[0\textsubscript{III}]{35}
es el término límite de dos sucesiones espectrales cuyos términos $E_2$ vienen dados por
\[
  \label{0.11.4.6.1-fr}
  'E_2^{pq}=\mathrm H^p(\mathrm R^qT(K^\bullet,K'^\bullet))
  \tag{11.4.6.1}
\]
\[
  \label{0.11.4.6.2-fr}
  ''E_2^{pq}=\sum_{q'+q''=q}
  \mathrm R^pT(\mathrm H^{q'}(K^\bullet),\mathrm H^{q''}(K'^\bullet))
  \tag{11.4.6.2}
\]
(cf.~M,~XVII,~2).

Aquí $\mathrm R^qT(K^\bullet,K'^\bullet)$ es el bicomplejo $(\mathrm R^qT(K^i,K'^j))_{(i,j)\in\mathbb Z\times\mathbb Z}$, y el segundo miembro de \sref{0.11.4.6.1} es su cohomología cuando se lo considera como complejo simple.

Además, la primera sucesión espectral es siempre regular, y las dos sucesiones espectrales son birregulares cuando existe $n$ tal que todo objeto de $\boldsymbol{C}$ y todo objeto de $\boldsymbol{C}'$ admiten una resolución inyectiva de longitud $\leq n$, o cuando $K^\bullet$ \emph{y} $K'^\bullet$ están acotados inferiormente;
en este último caso, puede omitirse además la hipótesis de que existan los límites inductivos en $\boldsymbol{C}$ y $\boldsymbol{C}'$.

Si $K_1^\bullet$ y $K_1'^\bullet$ son otros dos complejos de $\boldsymbol{C}$ y $\boldsymbol{C}'$, respectivamente, $f$ y $g$ son morfismos \emph{homótopos} de $K^\bullet$ en $K_1^\bullet$,
y $f'$ y $g'$ son morfismos \emph{homótopos} de $K'^\bullet$ en $K_1'^\bullet$, entonces los morfismos $\mathbf R^\bullet T(K^\bullet,K'_\bullet)\to
\mathbf R^\bullet T(K_1^\bullet,K_1'^\bullet)$ inducidos por $f$ y $f'$, por una parte, y por $g$ y $g'$, por otra, son idénticos, y lo mismo ocurre con los morfismos de las sucesiones espectrales de hipercohomología.

Se generaliza fácilmente a todo multifuntor covariante aditivo.
\end{env}

\begin{proposition}[11.4.7]
\label{0.11.4.7}
\label{0.11.4.7-fr}
Supongamos que, para todo objeto inyectivo $I$ de $\boldsymbol{C}$ \textup{(resp.} $I'$ de $\boldsymbol{C}'$), $A'\leadsto T(I,A')$ \textup{(resp.} $A\leadsto T(A,I')$) es un funtor exacto.
Entonces, con la notación de \sref{0.11.4.6}, se tiene isomorfismos canónicos
\[
  \label{0.11.4.7.1-fr}
  \mathbf R^\bullet T(K^\bullet,K'^\bullet)
  \xrightarrow{\sim}\mathrm H^\bullet(T(L^{\bullet\bullet},K'^\bullet))
  \xrightarrow{\sim}\mathrm H^\bullet(T(K^\bullet,L'^{\bullet\bullet}))
  \tag{11.4.7.1}
\]
donde los dos últimos términos son la cohomología de los complejos simples definidos por los tricomplejos $T(L^{\bullet\bullet},K'^\bullet)$ y $T(K^\bullet,L'^{\bullet\bullet})$, respectivamente.
\end{proposition}

\begin{proof}
Definamos, por ejemplo, el primero de estos isomorfismos.
El cuadricomplejo $T(L^{\bullet\bullet},L'^{\bullet\bullet}))$ puede considerarse como un bicomplejo, donde los grados de $T(L^{ij},L'^{hk})$ son los enteros $i+j+h$ y $k$.
Como, para cada $h$, $L'^{h,\bullet}$ es una \emph{resolución} de $K'^h$, se tiene, para este bicomplejo, en virtud de las hipótesis sobre $T$, $\mathrm H_{\mathrm{II}}^q(T(L^{\bullet\bullet},L'^{\bullet\bullet}))=0$ para $q\neq0$, y $\mathrm H_{\mathrm{II}}^0(T(L^{\bullet\bullet},L'^{\bullet\bullet}))
=T(L^{\bullet\bullet},K'^\bullet)$;
la primera sucesión espectral de este bicomplejo es, pues, \emph{degenerada};
como $L'^{hk}=0$ para $k<0$, esta sucesión es también regular \sref{0.11.3.3}, y la conclusión se deduce entonces de \sref{0.11.1.6}.
\end{proof}

Se tienen resultados análogos para un multifuntor covariante de un número cualquiera $n$ de argumentos: en el cálculo de la hipercohomología no es necesario reemplazar \emph{todos} los complejos
por una resolución de Cartan--Eilenberg, sino únicamente $n-1$ de ellos, siempre que, al fijar $n-1$ argumentos arbitrarios tomándolos como objetos \emph{inyectivos}, el funtor covariante en el argumento restante sea \emph{exacto}.


\subsection{Paso al límite inductivo en la hipercohomología}
\label{subsection:0.11.5}
\label{subsection:0.11.5-fr}

\begin{lemma}[11.5.1]
\label{0.11.5.1-fr}
Sea $K^\bullet=(K^i)_{i\in\mathbb Z}$ un complejo de $\boldsymbol{C}$, y para todo entero $r\in\mathbb Z$, sea $K_{(r)}^\bullet$ el complejo tal que $K_{(r)}^i=0$ para $i<r$, $K_{(r)}^i=K^i$ para $i\geq r$. Sea $T$ un funtor covariante aditivo
\oldpage[0\textsubscript{III}]{36}
de $\boldsymbol{C}$ en $\boldsymbol{C}'$, que conmuta con los límites inductivos (se supone que los límites inductivos filtrantes existen y son exactos en $\boldsymbol{C}$ y $\boldsymbol{C}'$). Entonces $\mathbf R^\bullet T(K^\bullet)$ es isomorfo al límite inductivo $\varinjlim\mathbf R^\bullet T(K_{(r)}^\bullet)$ cuando $r$ tiende a $-\infty$.
\end{lemma}

La construcción de una resolución inyectiva de Cartan-Eilenberg de $K^\bullet$ se hace eligiendo \emph{arbitrariamente}, para cada $i$, una resolución inyectiva $(X_{\mathrm B}^{ij})_{j\geq0}$ de $\mathrm B^i(K^\bullet)$ y una resolución inyectiva $(X_{\mathrm H}^{ij})_{j\geq0}$ de $\mathrm H^i(K^\bullet)$; hecho esto, el método de construcción muestra que la resolución inyectiva $(L^{ij})_{j\geq0}$ de $K^i$ y los operadores de derivación $L^{i,j}\to L^{i+1,j}$ \emph{solo dependen de las resoluciones} $(X_{\mathrm B}^{i,\bullet})$, $(X_{\mathrm H}^{i,\bullet})$ \emph{y} $(X_{\mathrm B}^{i+1,\bullet})$ (M, XVII, 1.2). Ahora bien, está claro que se tiene $\mathrm B^i(K_{(r)}^\bullet)=\mathrm B^i(K^\bullet)$ y $\mathrm H^i(K_{(r)}^\bullet)=\mathrm H^i(K^\bullet)$ para $i\geq r+1$. Por otra parte, para cada $r$, se tiene una inyección canónica $\varphi_{r-1,r}^\bullet:K_{(r)}^\bullet\to K_{(r-1)}^\bullet$, siendo $\varphi_{r-1,r}^i$ la identidad para $i\neq r-1$. La observación precedente muestra que, si $L^{\bullet\bullet}=(L^{ij})$ es una resolución inyectiva de Cartan-Eilenberg de $K^\bullet$, se puede para cada $r$ definir una resolución inyectiva de Cartan-Eilenberg $L_{(r)}^{\bullet\bullet}=(L_{(r)}^{ij})$ de $K_{(r)}^\bullet$ tal que $L_{(r)}^{ij}=0$ para $i<r$ y $L_{(r)}^{ij}=L^{ij}$ para $i\geq r+1$. Por otra parte, se puede definir un morfismo de bicomplejos $\Phi_{r-1,r}^{\bullet\bullet}:L_{(r)}^{\bullet\bullet}\to
L_{(r-1)}^{\bullet\bullet}$ correspondiente a $\varphi_{r-1,r}^\bullet$, y el método de definición de este morfismo (\emph{loc. cit.}) muestra además que puede construirse de modo que $\Phi_{r-1,r}^{ij}$ sea la identidad para $i\neq r$ y $i\neq r-1$. Se ha definido así un sistema inductivo $(L_{(r)}^{\bullet\bullet})$ de bicomplejos de $\boldsymbol{C}$, del que $L^{\bullet\bullet}$ es evidentemente el límite inductivo cuando $r$ tiende a $-\infty$; por la permutabilidad de las sumas directas y de los límites inductivos, el complejo simple asociado a $L^{\bullet\bullet}$ es también límite inductivo del complejo simple asociado a $L_{(r)}^{\bullet\bullet}$. Como $T$ conmuta por hipótesis con los límites inductivos y lo mismo ocurre con la cohomología (por la exactitud del funtor $\varinjlim$), se tiene en efecto $\mathrm H^\bullet(T(L^{\bullet\bullet}))=\varinjlim\mathrm H^\bullet(T(L_{(r)}^{\bullet\bullet}))$ salvo isomorfismo.

El lema (\hyperref[0.11.5.1-fr]{11.5.1}) permite extender a complejos $K^\bullet$ cualesquiera, por paso al límite inductivo, propiedades válidas para los complejos \emph{acotados inferiormente}. Como primer ejemplo, demostraremos:

\begin{proposition}[11.5.2]
\label{0.11.5.2}
\label{0.11.5.2-fr}
Bajo las hipótesis de \sref{0.11.5.1} relativas a $\boldsymbol{C}$, $\boldsymbol{C}'$ y $T$, el funtor $\mathbf R^\bullet T(K^\bullet)$ es un funtor cohomológico sobre la categoría abeliana de los complejos de $\boldsymbol{C}$.
\end{proposition}

\begin{proof}
Mostremos primero que basta tratar los complejos \emph{acotados inferiormente}.
Dada una sucesión exacta
$0\longrightarrow K'^\bullet\longrightarrow K^\bullet\longrightarrow K''^\bullet\longrightarrow0$
de complejos, para todo $r$ induce una sucesión exacta
$0\longrightarrow K_{(r)}'^\bullet\longrightarrow K_{(r)}^\bullet\longrightarrow K_{(r)}''^\bullet\longrightarrow0$,
Por hipótesis, esta última da una sucesión exacta
\[
  \ldots\longrightarrow\mathbf R^nT(K_{(r)}'^\bullet)
  \longrightarrow\mathbf R^nT(K_{(r)}^\bullet)
  \longrightarrow\mathbf R^nT(K_{(r)}''^\bullet)
  \xrightarrow{\partial}\mathbf R^{n+1}T(K_{(r)}'^\bullet)
  \longrightarrow\ldots
\]
Estas sucesiones exactas forman un sistema inductivo.
El Lema~\sref{0.11.5.1} y la exactitud de $\varinjlim$ dan, por tanto, una sucesión exacta
\[
  \ldots\longrightarrow\mathbf R^nT(K'^\bullet)
  \longrightarrow\mathbf R^nT(K^\bullet)
  \longrightarrow\mathbf R^nT(K''^\bullet)
  \xrightarrow{\partial}\mathbf R^{n+1}T(K'^\bullet)
  \longrightarrow\ldots
\]
\end{proof}

Para tratar los complejos acotados inferiormente, basta considerar aquellos para los cuales $K^i=0$ cuando $i<0$; forman evidentemente una \emph{categoría abeliana} $\boldsymbol{K}$.

\begin{lemma}[11.5.2.1]
\label{0.11.5.2.1}
\label{0.11.5.2.1-fr}
En $\boldsymbol{K}$, sea $\mathfrak J$ la clase de los complejos $Q^\bullet=(Q^i)_{i\geq0}$ que tienen las
\oldpage[0\textsubscript{III}]{37}
propiedades siguientes: $1^\circ$ Todo $Q^i$ es un objeto inyectivo de $\boldsymbol{C}$; $2^\circ$ Para todo $i\geq0$, $\mathrm Z^i(Q^\bullet)=\mathrm B^i(Q^\bullet)$, y $\mathrm Z^i(Q^\bullet)$ es un sumando directo de $Q^i$.
Entonces:
\begin{enumerate}
  \item[{\rm(i)}] todo $Q^\bullet\in\mathfrak J$ es un objeto inyectivo de $\boldsymbol{K}$;
  \item[{\rm(ii)}] todo objeto de $\boldsymbol{K}$ es isomorfo a un subcomplejo de un complejo perteneciente a $\mathfrak J$.
\end{enumerate}
\end{lemma}

\begin{proof}
{\rm(i)}
Sea $A^\bullet=(A^i)$ un objeto de $\boldsymbol{K}$, sea $A'^\bullet=(A'^i)$ un subobjeto de $A^\bullet$, sea $Q^\bullet=(Q^i)$ perteneciente a $\mathfrak J$, y supongamos dado un morfismo
$f=(f^i):A'^\bullet\to Q^\bullet$, que se trata de prolongar a un morfismo $g=(g^i):A^\bullet\to Q^\bullet$.
Por simplicidad, utilizamos el lenguaje de la categoría de módulos (cf.~[27]).

Identifiquemos $Q$ con $\mathrm B^i(Q^\bullet)\oplus\mathrm B^{i+1}(Q^\bullet)$.
Procedemos por inducción sobre $i$.
Supongamos que $g^j$ ha sido definido para $j<i$, compatiblemente con los diferenciales $d^j:A^j\to A^{j+1}$ y $d'^j:Q^j\to Q^{j+1}$ para $j<i-1$ y tales además que: $1^\circ$ $g^{i-1}(\mathrm Z^{i-1}(A^\bullet))\subset
\mathrm Z^{i-1}(Q^\bullet)$; $2^\circ$ Si se pone $C^j=(d^j)^{-1}(A'^{j+1})$ para todo $j$, entonces $d'^{i-1}\circ g^{i-1}$ coincide con $f^i\circ d^{i-1}$ sobre $C^{i-1}$.
Componiendo $f^i:A'^i\to Q^i$ con las proyecciones, se obtienen dos morfismos $f'^i:A'^i\to\mathrm B^i(Q^\bullet)$ y $f''^i:A'^i\to\mathrm B^{i+1}(Q^\bullet)$. Como $d'^{i-1}\circ g^{i-1}$ envía $A^{i-1}$ en $\mathrm B^i(Q^\bullet)$ y se anula sobre $\mathrm Z^{i-1}(A^\bullet)$, define un morfismo
$h^i:\mathrm B^i(A^\bullet)\to\mathrm B^i(Q^\bullet)$, y puesto que $d'^{i-1}\circ g^{i-1}$ coincide con $f^i\circ d^{i-1}$ sobre $C^{i-1}$, $h^i$ coincide con $f'_i$ sobre $\mathrm B^i(A^\bullet)\cap A'^i$.
Como $\mathrm B^i(Q^\bullet)$ es un sumando directo de $Q^i$, es inyectivo; existe, por tanto, un morfismo
$g'^i:A^i\to\mathrm B^i(Q^\bullet)$
que coincide con $h^i$ sobre $\mathrm B^i(A^\bullet)$ y con $f'^i$ sobre $A'^i$.

Por otra parte, consideremos el morfismo
$f''^{i+1}\circ d^i:C^i\to\mathrm B^{i+1}(Q^\bullet)$,
que se anula sobre $\mathrm Z^i(A^\bullet)$.
Como $\mathrm B^{i+1}(Q^\bullet)$ es inyectivo, existe un morfismo
$g''^i:A^i\to\mathrm B^{i+1}(Q^\bullet)$
que coincide con $f''^{i+1}\circ d^i$ sobre $C^i$ y con $0$ en $\mathrm Z^i(A^\bullet)$.
Tomando $g^i=g'^i+g''^i$ puede continuarse la inducción.

\smallskip
\noindent{\rm(ii)}
Para sumergir $A^\bullet=(A^i)$ en un complejo perteneciente a $\mathfrak J$, se toma para cada $i\geq1$ un objeto inyectivo $Q'^i$ de $\boldsymbol{C}$ y una inyección $f'^i:A^i\to Q'^i$. Pongamos entonces $Q^i=0$ para $i<0$, $Q^0=Q'^1$ y $Q^i=Q'^i\oplus Q'^{i+1}$ para $i\geq1$, con el operador de derivación evidente. Pongamos
$f^i=f'^i+(f'^{i+1}\circ d^i)$ para todo $i$ (con $f'^i=0$ para $i\leq0$);
Es inmediato que $f^i$ es inyectivo para todo $i$, y que los $f^i$ conmutan con los diferenciales.
\end{proof}

\begin{corollary}[11.5.2.2]
\label{0.11.5.2.2}
\label{0.11.5.2.2-fr}
Todo objeto $K^\bullet$ de $\boldsymbol{K}$ admite una resolución por la derecha formada por objetos de $\mathfrak J$.
Si $L^{\bullet\bullet}$ es una tal resolución, entonces para toda resolución $L'^{\bullet\bullet}$ de $K^\bullet$ formada por objetos de $\boldsymbol{K}$ existe un morfismo de bicomplejos
$F:L'^{\bullet\bullet}\to L^{\bullet\bullet}$ compatible con las aumentaciones, y dos tales morfismos $F$, $F'$ son homótopos.
\end{corollary}

\begin{proof}
Esto no es sino (M, V, 1.1 \emph{a)}) aplicado a la categoría abeliana
$\boldsymbol{K}$.
\end{proof}

\begin{env}[11.5.2.3]
\label{0.11.5.2.3-fr}
Establecidos estos preliminares, consideremos una resolución de Cartan-Eilenberg inyectiva $L'^{\bullet\bullet}$ una resolución de Cartan--Eilenberg inyectiva de $K^\bullet$ y sea $L^{\bullet\bullet}$ una resolución de $K^\bullet$ formada por objetos de $\mathfrak J$.
Mostremos que existe un isomorfismo
$\mathrm H^\bullet(T(L'^{\bullet\bullet}))\xrightarrow{\sim}\mathrm H^\bullet(T(L^{\bullet\bullet}))$.
El Corolario~\sref{0.11.5.2.2} da un morfismo de bicomplejos
$T(L'^{\bullet\bullet})\to T(L^{\bullet\bullet})$,
y, por tanto, un morfismo
$'E(T(L'^{\bullet\bullet}))\to{}'E(T(L^{\bullet\bullet}))$
entre sus primeras sucesiones espectrales.
Por \sref{0.11.3.3}, estas sucesiones espectrales son \emph{regulares}, de modo que por \sref{0.11.1.5} basta comprobar que este morfismo es un isomorfismo sobre los términos $E_2$, o también que
$\mathrm H_{\mathrm{II}}^q(T(L'^{i,\bullet}))$ es igual a
$\mathrm R^qT(K^i)$; como $L'^{i,\bullet}$ es una resolución por la derecha de $K^i$, esto reduce la cuestión al lema siguiente.
\end{env}

\oldpage[0\textsubscript{III}]{38}
\begin{lemma}[11.5.2.4]
\label{0.11.5.2.4}
\label{0.11.5.2.4-fr}
Sea $L^\bullet=(L^i)_{i\in\mathbb Z}$ una resolución por la derecha de un objeto $A$ de $\boldsymbol{C}$ tal que $\mathrm R^nT(L^i)=0$ para todo $i\in\mathbb Z$ y todo $n>0$.
Entonces
$\mathrm H^\bullet T(L^\bullet)=\mathrm R^\bullet T(A)$.
\end{lemma}

\begin{proof}
Este es un caso particular de (T, 2.5.2).
\end{proof}

\begin{env}[11.5.2.5]
\label{0.11.5.2.5}
\label{0.11.5.2.5-fr}
La demostración de \sref{0.11.5.2} se deduce ahora inmediatamente, puesto que $L'^{\bullet\bullet}$ es una resolución inyectiva de $K^\bullet$ \emph{en la categoría abeliana $\boldsymbol{K}$}.
En otras palabras, el funtor
$K^\bullet\leadsto\mathrm H^\bullet(T(L'^{\bullet\bullet}))$
es precisamente el funtor cohomológico \emph{derivado por la derecha} de $T$ en la categoría $\boldsymbol{K}$ (T, 2.3).
\end{env}

\begin{proposition}[11.5.3]
\label{0.11.5.3}
\label{0.11.5.3-fr}
Bajo las hipótesis de \sref{0.11.5.1} relativas a $\boldsymbol{C}$, $\boldsymbol{C}'$ y $T$, sea $L^{\bullet\bullet}=(L^{ij})$ un bicomplejo tal que $L^{ij}=0$ para $j<0$ y tal que, para todo $i$, $L^{i,\bullet}$ es una resolución de $K^i$.
Supongamos además que $\mathrm R^nT(L^{ij})=0$ para todo par $(i,j)$ y todo $n>0$.
Entonces existe un isomorfismo funtorial
\[
  \label{0.11.5.3.1-fr}
  \mathbf R^\bullet T(K^\bullet)\xrightarrow{\sim}
  \mathrm H^\bullet(T(L^{\bullet\bullet}))).
  \tag{11.5.3.1}
\]
\end{proposition}

\begin{proof}
Sea $L_{(r)}^{\bullet\bullet}=(L_{(r)}^{ij})$ el bicomplejo tal que $L_{(r)}^{ij}=0$ para $i<r$, $L_{(r)}^{ij}=L^{ij}$ para $i\geq r$; es inmediato que $L^{\bullet\bullet}$ es el límite inductivo de los $L_{(r)}^{\bullet\bullet}$ cuando $r$ tiende a $-\infty$.
Por la hipótesis sobre $T$ y por \sref{0.11.5.1}, basta, por tanto, demostrar la proposición cuando $K^\bullet$ está acotado inferiormente, por ejemplo cuando $K^i=0$ para $i<0$, con $L^{ij}=0$ para $i<0$.
Sea $L'^{\bullet\bullet}=(L'^{ij})$ una resolución por la derecha de $K^\bullet$ formada por objetos de $\mathfrak J$, como en \sref{0.11.5.2.2}.
Existe un morfismo de bicomplejos
$L^{\bullet\bullet}\to L'^{\bullet\bullet}$
compatible con las aumentaciones y, por tanto, un morfismo
$'E(T(L^{\bullet\bullet}))\to{}'E(T(L'^{\bullet\bullet}))$
entre las primeras sucesiones espectrales.
El Lema~\sref{0.11.5.2.4} muestra, exactamente como en \sref{0.11.5.2.3}, que este morfismo es un isomorfismo, lo que demuestra la aserción.
\end{proof}

\begin{remark}[11.5.4]
\label{0.11.5.4}
\label{0.11.5.4-fr}
Los razonamientos anteriores muestran que las conclusiones de \sref{0.11.5.2} y \sref{0.11.5.3} se cumplen en la categoría de los complejos $K^\bullet$ \emph{acotados inferiormente} sin suponer que $T$ conmuta con los límites inductivos filtrantes.
Además, si se considera únicamente la categoría $\boldsymbol{K}$ de los complejos $K^\bullet$ tales que $K^i=0$ para $i<0$, la caracterización de $\mathbf R^\bullet T(K^\bullet)$ como el sistema de funtores derivados por la derecha de $T$ en $\boldsymbol{K}$ muestra que este funtor cohomológico es \emph{universal} (T, 2.3).

Otro caso en el que \sref{0.11.5.2} se cumple sin hipótesis adicional alguna sobre $T$ es aquel en que existe un entero $m>0$ tal que todo objeto de $\boldsymbol{C}$ admita una resolución inyectiva \emph{de longitud} $\leq m$. En efecto, en la demostración de (\hyperref[0.11.5.1-fr]{11.5.1}), todas las resoluciones inyectivas de objetos de $\boldsymbol{C}$ que intervienen pueden tomarse de longitud $\leq m$, de donde resulta inmediatamente que los términos de grado total $n$ en $T(L_{(r)}^{\bullet\bullet})$ son \emph{iguales a los de} $T(L^{\bullet\bullet})$ y en número \emph{finito}, en cuanto $r$ es suficientemente grande, lo que
implica que para todo $n$,
$\mathrm H^n(T(L^{\bullet\bullet}))=
\mathrm H^n(T(L_{(r)}^{\bullet\bullet}))$
para $r$ suficientemente grande.
Con la notación de \sref{0.11.5.2}, se tiene análogamente
$\mathbf R^nT(K_{(r)}^\bullet)=\mathbf R^nT(K^\bullet)$
para $r$ suficientemente grande (dependiendo de $n$), y análogamente para $K'^\bullet$ y $K''^\bullet$, lo que demuestra la conclusión.
Del mismo modo, \sref{0.11.5.3} se cumple sin condición adicional sobre $T$ cuando $\boldsymbol{C}$ satisface la hipótesis anterior y se toman las resoluciones $L^{i,\bullet}$ son de longitud $\leq m$.
\end{remark}

\begin{env}[11.5.5]
\label{0.11.5.5-fr}
El resultado de (\hyperref[0.11.5.2-fr]{11.5.2}) se generaliza a los multifuntores covariantes. Consideremos por ejemplo la situación de (\hyperref[0.11.4.6-fr]{11.4.6}), donde se supone que en $\boldsymbol{C}$, $\boldsymbol{C}'$ y $\boldsymbol{C}''$
\oldpage[0\textsubscript{III}]{39}
los límites inductivos filtrantes existen y son exactos, y que $T$ conmuta con los límites inductivos. Entonces $\mathbf R^\bullet T(K^\bullet,K'^\bullet)$ es un bifuntor \emph{cohomológico} en cada uno de los complejos $K^\bullet$, $K'^\bullet$; para verlo, se reduce, como en (\hyperref[0.11.5.2-fr]{11.5.2}), al caso en que $K^\bullet$ y $K'^\bullet$ están acotados inferiormente; tomando entonces para $K^\bullet$ y $K'^\bullet$ resoluciones inyectivas del tipo descrito en (\hyperref[0.11.5.2.2-fr]{11.5.2.2}), se reduce a la propiedad general demostrada en (M, V, 4.1).
\end{env}

\begin{env}[11.5.6]
\label{0.11.5.6-fr}
Del mismo modo, los resultados de (\hyperref[0.11.4.7-fr]{11.4.7}) y (\hyperref[0.11.5.3-fr]{11.5.3}) se generalizan como sigue. Supongamos (bajo las hipótesis de (\hyperref[0.11.5.5-fr]{11.5.5})) que se tienen dos bicomplejos $L^{\bullet\bullet}=(L^{ij})$, $L'^{\bullet\bullet}=(L'^{ij})$ tales que $L^{ij}=0$ y $L'^{ij}=0$ para $j<0$, tales que, para todo $i$, $L^{i,\bullet}$ es una resolución de $K^i$ y $L'^{i,\bullet}$ es una resolución de $K'^i$, y tales que
$\mathrm R^nT(L^{ij},L'^{hk})=0$
para $n>0$ y todo sistema de índices $(i,j,h,k)$.
Entonces existe un isomorfismo, funtorial en $K^\bullet$ y $K'^\bullet$,
\[
  \label{0.11.5.6.1-fr}
  \mathbf R^\bullet T(K^\bullet,K'^\bullet)\xrightarrow{\sim}
  \mathrm H^\bullet(T(L^{\bullet\bullet},L'^{\bullet\bullet}))
  \tag{11.5.6.1}
\]
Esto se demuestra como en \sref{0.11.5.3}, reduciéndose al caso en que $K^\bullet$ y $K'^\bullet$ están acotados inferiormente.

Supongamos además que, para todo par $(i,j)$ y todo par $(h,k)$, los funtores $A\leadsto T(A,L'^{hk})$ y $A'\leadsto T(L^{ij},A')$ sean exactos en $\boldsymbol{C}$ y $\boldsymbol{C}'$, respectivamente.
Entonces existen también isomorfismos funtoriales
\[
  \label{0.11.5.6.2-fr}
  \mathbf R^\bullet T(K^\bullet,K'^\bullet)
  \sim\mathrm H^\bullet(T(L^{\bullet\bullet},K'^\bullet))
  \xrightarrow{\sim}\mathrm H^\bullet(T(K^\bullet,L'^{\bullet\bullet})).
  \tag{11.5.6.2}
\]
La demostración es análoga a la de \sref{0.11.4.7}.
\end{env}

\begin{env}[11.5.7]
\label{0.11.5.7}
\label{0.11.5.7-fr}
Los resultados de \sref{0.11.5.5} y \sref{0.11.5.6} se cumplen también sin suponer que $T$ conmuta con los límites inductivos, con tal de restringirse a los complejos $K^\bullet$, $K'^\bullet$ acotados inferiormente, o de suponer que todo objeto de $\boldsymbol{C}$ (resp. de $\boldsymbol{C}'$) admite una resolución inyectiva de longitud acotada y, en \sref{0.11.5.6}, que los segundos grados de $L^{\bullet\bullet}$ y $L'^{\bullet\bullet}$ están acotados superiormente.
\end{env}

\subsection{Hiperhomología de un funtor con respecto a un complejo $K_\bullet$}
\label{subsection:0.11.6}
\label{subsection:0.11.6-fr}

\begin{env}[11.6.1]
\label{0.11.6.1}
\label{0.11.6.1-fr}
Sea $\boldsymbol{C}$ una categoría abeliana y sea $K_\bullet=(K_i)_{i\in\mathbb Z}$ un complejo de objetos de $\boldsymbol{C}$ cuyo diferencial es de grado $-1$.
Una \emph{resolución de Cartan--Eilenberg por la izquierda} de $K_\bullet$ consiste en un bicomplejo $L_{\bullet\bullet}=(L_{ij})$, cuyos diferenciales son de grado $-1$ y para el cual $L_{ij}=0$ cuando $j<0$, junto con un morfismo de complejos simples
$\varepsilon:L_{\bullet,0}\to K_\bullet$,
tal que se satisfacen las condiciones obtenidas de las de \sref{0.11.4.2} ``invirtiendo las flechas''.
Si todo objeto de $\boldsymbol{C}$ es cociente de un objeto proyectivo, entonces todo complejo $K_\bullet$ de $\boldsymbol{C}$ admite una \emph{resolución de Cartan--Eilenberg proyectiva}, es decir, formada por objetos proyectivos $L_{ij}$; posee propiedades funtoriales análogas a las de \sref{0.11.4.2}.
Además, si $K_\bullet$ está acotado inferiormente (resp. superiormente), puede exigirse $L_{ij}=0$ para $i<i_0$ (resp. $i>i_0$) siempre que $K_i=0$ para $i<i_0$ (resp. $i>i_0$).
Si todo objeto de $\boldsymbol{C}$ admite una resolución proyectiva de longitud $\leq n$, se puede suponer que $L_{ij}=0$ para $j>n$.
\end{env}

\begin{env}[11.6.2]
\label{0.11.6.2}
\label{0.11.6.2-fr}
\oldpage[0\textsubscript{III}]{40}
Supongamos que $T$ es un \emph{funtor covariante aditivo} de $\boldsymbol{C}$ en una categoría abeliana $\boldsymbol{C}'$.
La \emph{hiperhomología} $\mathbf{L}_\bullet T(K_\bullet)$ y las \emph{sucesiones espectrales de hiperhomología} de $T$ con respecto a un complejo $K_\bullet$ de $\boldsymbol{C}$ (cuando existen) se definen a partir de \sref{0.11.4.3}, de nuevo ``invirtiendo las flechas''.
Los términos $E^2$ de las dos sucesiones espectrales así obtenidas son
\[
\label{0.11.6.2.1-fr}
\tag{11.6.2.1}
{}'E^2_{pq}=H_p(L_qT(K_\bullet))
\]
y
\[
\label{0.11.6.2.2-fr}
\tag{11.6.2.2}
{}''E^2_{pq}=L_pT(H_q(K_\bullet))
\]
donde $L_pT$ denota el funtor derivado de grado $p$ de $T$ para $p\geq 0$, y $0$ para $p<0$, mientras que $L_qT(K_\bullet)$ denota el complejo $(L_qT(K_i))_{i\in\mathbb Z}$.

Las propiedades de la hiperhomología no se deducen todas por simple «inversión de las flechas» de las de la hipercohomología (a menos que se hagan hipótesis suplementarias del tipo AB~5*) de T, 1.5 sobre la categoría $\boldsymbol{C}'$) a causa de las condiciones de \emph{regularidad} de las dos sucesiones espectrales precedentes, a las que esta vez hay que aplicar los criterios de (\hyperref[0.11.3.4-fr]{11.3.4}). Estos últimos muestran que, cuando se supone que en $\boldsymbol{C}'$ los límites inductivos filtrantes existen y son exactos, entonces el complejo definido por el bicomplejo $T(L_{\bullet\bullet})$ existe, y la \emph{segunda} sucesión espectral $\,{}''E(T(L_{\bullet\bullet}))$ es esta vez \emph{regular}. Si se supone, o bien que $K_\bullet$ está acotado inferiormente, o bien que existe un entero $n$ tal que todo objeto de $\boldsymbol{C}$ admita una resolución proyectiva de longitud $\leq n$, entonces las \emph{dos} sucesiones espectrales de hiperhomología existen (sin hipótesis sobre $\boldsymbol{C}'$) y son \emph{birregulares}.
\end{env}

\begin{proposition}[11.6.3]
\label{0.11.6.3}
\label{0.11.6.3-fr}
Sean $\boldsymbol{C}$ y $\boldsymbol{C}'$ dos categorías abelianas, y sea $T$ un funtor covariante aditivo de $\boldsymbol{C}$ en $\boldsymbol{C}'$.
Entonces:
\begin{enumerate}
  \item[{\rm(i)}] {La hiperhomología $\mathbf{L}_\bullet T(K_\bullet)$ es un funtor homológico sobre la categoría abeliana de los complejos de $\boldsymbol{C}$ acotados inferiormente.}

  \item[{\rm(ii)}] {Sea $K_\bullet$ un complejo de $\boldsymbol{C}$ acotado inferiormente.
  Si $L_nT(K_i)=0$ para todo $n>0$ y todo $i\in\mathbb Z$, existen isomorfismos funtoriales}
  \[
\label{0.11.6.3.1-fr}
\tag{11.6.3.1}
\mathbf{L}_iT(K_\bullet)\xrightarrow{\sim} H_i(T(K_\bullet))\qquad(i\in\mathbb Z).
\]

  \item[{\rm(iii)}] {Sea $K_\bullet$ un complejo de $\boldsymbol{C}$ acotado inferiormente.
  Sea $L_{\bullet\bullet}=(L_{ij})$ un bicomplejo tal que $L_{ij}=0$ para $j<0$ y, para todo $i$, $L_{i,\bullet}$ es una resolución de $K_i$.
  Supongamos además que $L_nT(L_{ij})=0$ para todo par $(i,j)$ y todo $n>0$.
  Entonces existe un isomorfismo funtorial}
  \[
\label{0.11.6.3.2-fr}
\tag{11.6.3.2}
\mathbf{L}_\bullet T(K_\bullet)\xrightarrow{\sim} H_\bullet(T(L_{\bullet\bullet}))
\]
\end{enumerate}

\upshape Las demostraciones proceden como las de \sref{0.11.5.2}, \sref{0.11.4.5} y \sref{0.11.5.3} en el caso de complejos acotados inferiormente.
Dejamos al lector los detalles de estos razonamientos.
\end{proposition}

\begin{env}[11.6.4]
\label{0.11.6.4}
\label{0.11.6.4-fr}
Se tienen resultados enteramente análogos para los \emph{multifuntores} covariantes aditivos en cada argumento.
Por ejemplo, para un bifuntor $T$, las dos sucesiones espectrales de hiperhomología tienen términos $E^2$
\oldpage[0\textsubscript{III}]{41}
\[
\label{0.11.6.4.1-fr}
\tag{11.6.4.1}
{}'E^2_{pq}=H_p(L_qT(K_\bullet,K'_\bullet))
\]
y
\[
\label{0.11.6.4.2-fr}
\tag{11.6.4.2}
{}''E^2_{pq}=
\sum_{q'+q''=q}L_p(H_{q'}(K_\bullet),H_{q''}(K'_\bullet)).
\]
También aquí es la \emph{segunda} sucesión espectral la que es regular.
Las dos sucesiones son birregulares cuando $K_\bullet$ y $K'_\bullet$ están acotados inferiormente, o cuando los objetos de las categorías abelianas consideradas admiten resoluciones proyectivas de una longitud fija.

Además:
\end{env}

\begin{proposition}[11.6.5]
\label{0.11.6.5}
\label{0.11.6.5-fr}
Sean $\boldsymbol{C}$, $\boldsymbol{C}'$ y $\boldsymbol{C}''$ tres categorías abelianas, y sea $T$ un bifuntor covariante biaditivo de $\boldsymbol{C}\times\boldsymbol{C}'$ en $\boldsymbol{C}''$.
\begin{enumerate}
  \item[{\rm(i)}] {$\mathbf{L}_\bullet T(K_\bullet,K'_\bullet)$ es un bifuntor homológico en cada uno de los complejos acotados inferiormente $K_\bullet$, $K'_\bullet$, formados respectivamente por objetos de $\boldsymbol{C}$ y de $\boldsymbol{C}'$.}

  \item[{\rm(ii)}] {Supongamos que $K_\bullet$ y $K'_\bullet$ están acotados inferiormente.
  Sean $L_{\bullet\bullet}=(L_{ij})$ y $L'_{\bullet\bullet}=(L'_{ij})$ dos bicomplejos tales que $L_{ij}=0$ y $L'_{ij}=0$ para $j<0$ y tales que, para todo $i$, $L_{i,\bullet}$ es una resolución de $K_i$ y $L'_{i,\bullet}$ es una resolución de $K'_i$, y finalmente que $L_nT(L_{ij},L'_{hk})=0$ para todo $n>0$ y todo sistema de índices $(i,j,h,k)$.
  Entonces existe un isomorfismo funtorial}
  \[
\label{0.11.6.5.1-fr}
\tag{11.6.5.1}
\mathbf{L}_\bullet T(K_\bullet,K'_\bullet)
\xrightarrow{\sim} H_\bullet(T(L_{\bullet\bullet},L'_{\bullet\bullet})).
\]

  \item[{\rm(iii)}] {Supongamos además que, para todo par $(i,j)$ y todo par $(h,k)$, los funtores $A\leadsto T(A,L'_{hk})$ y $A'\leadsto T(L_{ij},A')$ sean exactos en $\boldsymbol{C}$ y $\boldsymbol{C}'$, respectivamente.
  Entonces existen isomorfismos funtoriales}
  \[
\label{0.11.6.5.2-fr}
\tag{11.6.5.2}
\mathbf{L}_\bullet T(K_\bullet,K'_\bullet)
\xrightarrow{\sim} H_\bullet(T(L_{\bullet\bullet},K'_\bullet))
\xrightarrow{\sim} H_\bullet(T(K_\bullet,L'_{\bullet\bullet})).
\]
\end{enumerate}

\upshape Las demostraciones son análogas a las de \sref{0.11.5.5} y \sref{0.11.5.6}.
\end{proposition}

\subsection{Hiperhomología de un funtor con respecto a un bicomplejo $K_{\bullet\bullet}$}
\label{subsection:0.11.7}
\label{subsection:0.11.7-fr}

\begin{env}[11.7.1]
\label{0.11.7.1}
\label{0.11.7.1-fr}
Sea $\boldsymbol{C}$ una categoría abeliana en la que todo objeto es cociente de un
objeto proyectivo.
Consideremos un bicomplejo $K_{\bullet\bullet}=(K_{ij})$ de objetos de $\boldsymbol{C}$ cuyos
dos grados están \emph{acotados inferiormente}.
Se puede limitar siempre al caso en que $K_{ij}=0$ para $i<0$ o
$j<0$, y es lo que haremos en adelante. Se puede considerar $K_{\bullet\bullet}$ como un complejo ordinario formado por los
objetos
$K_{i,\bullet}=(K_{ij})_{j\geq 0}$
de la categoría abeliana $\boldsymbol{K}$ de los complejos no negativos de objetos de
$\boldsymbol{C}$.
Del Lema~\sref{0.11.5.2.1} (o, más precisamente, del lema
``dual'' obtenido ``invirtiendo las flechas'') y de (M,~V,~2.2) se deduce
que $K_{\bullet\bullet}$ admite una \emph{resolución de Cartan--Eilenberg
proyectiva} en la categoría $\boldsymbol{K}$.
Una tal resolución es un \emph{tricomplejo}
$M_{\bullet\bullet\bullet}=(M_{ijk})$ en $\boldsymbol{C}$, con grados todos $\geq 0$, formado por objetos proyectivos, tal que para todo $i$, $M_{i,\bullet\bullet}$, $\mathrm B_i^{\mathrm I}(M_{\bullet\bullet\bullet})$, $\mathrm Z_i^{\mathrm I}(M_{\bullet\bullet\bullet})$, $H_i^{\mathrm I}(M_{\bullet\bullet\bullet})$ constituyen \emph{resoluciones proyectivas} de $K_{i,\bullet}$, $\mathrm B_i^{\mathrm I}(K_{\bullet\bullet})$, $\mathrm Z_i^{\mathrm I}(K_{\bullet\bullet})$, $H_i^{\mathrm I}(K_{\bullet\bullet})$ respectivamente en la categoría $\boldsymbol{K}$; en particular, para todo par $(i,j)$, $M_{ij,\bullet}$ es una \emph{resolución
proyectiva} de $K_{ij}$ en $\boldsymbol{C}$.
\end{env}

\begin{proposition}[11.7.2]
\label{0.11.7.2}
\label{0.11.7.2-fr}
Sea $T$ un funtor covariante aditivo de $\boldsymbol{C}$ en una categoría
abeliana $\boldsymbol{C}'$.
Con la notación de \sref{0.11.7.1}, la homología
$H_\bullet(T(M_{\bullet\bullet\bullet}))$
del complejo simple asociado al tricomplejo
$T(M_{\bullet\bullet\bullet})$ es canónicamente isomorfa a la
hiperhomología $\mathbf{L}_\bullet T(K_{\bullet\bullet})$ del
\oldpage[0\textsubscript{III}]{42}
complejo simple asociado a $K_{\bullet\bullet}$
\sref{0.11.6.2}.
Es también el término límite común de seis sucesiones espectrales birregulares denotadas
$\mathop{}^{(t)}E$, donde $t=a,b,a',b',c$ o $d$, cuyos
términos $E^2$ están dados por las fórmulas
\[
\begin{aligned}
{}^{(a)}E^2_{pq}&=L_pT(H_q(K_{\bullet\bullet}))&
{}^{(b)}E^2_{pq}&=H_p(\mathbf{L}_q^{\mathrm{II}}T(K_{\bullet\bullet}))\\
{}^{(a')}E^2_{pq}&=\mathbf{L}_pT(H_q^{\mathrm I}(K_{\bullet\bullet}))&
{}^{(b')}E^2_{pq}&=H_p(L_qT(K_{\bullet\bullet}))\\
{}^{(c)}E^2_{pq}&=\mathbf{L}_pT(H_q^{\mathrm{II}}(K_{\bullet\bullet}))&
{}^{(d)}E^2_{pq}&=H_p(\mathbf{L}_q^{\mathrm I}T(K_{\bullet\bullet}))
\end{aligned}
\]
\end{proposition}

(Se recuerda que utilizamos la notación $F(A_\bullet)$ para el complejo cuyos objetos son los $F(A_i)$ para todo complejo $A_\bullet=(A_i)$; por ejemplo $\mathbf{L}_q^{\mathrm{II}}T(K_{\bullet\bullet})$ denota el complejo
$(\mathbf{L}_q^{\mathrm{II}}T(K_{i,\bullet}))_{i\geq0}$,
donde $\mathbf{L}_q^{\mathrm{II}}T(K_{i,\bullet})$ es la hiperhomología de índice
$q$ del funtor $T$ relativa al complejo simple $K_{i,\bullet}$.)

\begin{proof}
Designemos por $L_\bullet$ el complejo simple asociado a
$K_{\bullet\bullet}$, de modo que
$L_i=\bigoplus_{r+s=i}K_{rs}$,
y pongamos
$N_{ij}=\bigoplus_{r+s=i}M_{rsj}$;
Está claro que, para todo $i$, $N_{i,\bullet}$ es una \emph{resolución proyectiva}
de $L_i$ en $\boldsymbol{C}$.
De \sref{0.11.6.3} y \sref{0.11.6.4} se deduce que existe un
isomorfismo funtorial
$
\mathbf{L}_\bullet T(L_\bullet)\xrightarrow{\sim} H_\bullet(T(N_{\bullet\bullet}))$;
Como el complejo simple asociado al bicomplejo
$T(N_{\bullet\bullet})$ es también el complejo simple asociado al
tricomplejo $T(M_{\bullet\bullet\bullet})$, esto demuestra la primera aserción.

Además, $\mathbf{L}_\bullet T(L_\bullet)$ es el término límite de las dos
sucesiones espectrales de hiperhomología \sref{0.11.6.2} de $T$ relativas al
complejo simple $L_\bullet$.
Estas son precisamente $\mathop{}^{(b')}E$ y $\mathop{}^{(a)}E$,
respectivamente.

Consideremos ahora $M_{\bullet\bullet\bullet}$ como un bicomplejo
$U_{\bullet\bullet}$, donde
$U_{ij}=\bigoplus_{r+s=j}M_{irs}$;
La homología $H_\bullet(T(M_{\bullet\bullet\bullet}))$ es también la
término límite de las dos sucesiones espectrales del bicomplejo
$T(U_{\bullet\bullet})$.
Para todo $i$, el bicomplejo $M_{i,\bullet\bullet}$ satisface las
hipótesis de \sref{0.11.6.3} relativas al complejo simple
$K_{i,\bullet}$.
Por consiguiente,
$H_q^{\mathrm{II}}(T(U_{i,\bullet}))=
\mathbf{L}_q^{\mathrm{II}}T(K_{i,\bullet}))$,
y la primera sucesión espectral de $T(U_{\bullet\bullet})$ es precisamente
$\mathop{}^{(b)}E$.

Por otra parte, para todo $r$, $M_{\bullet,r,\bullet}$ es una
resolución de Cartan--Eilenberg del complejo simple $K_{\bullet,r}$.
El cálculo efectuado en (M,~XV,~2) muestra que
$H_q^{\mathrm I}(T(M_{\bullet,rs}))=T(H_q^{\mathrm I}(M_{\bullet,rs}))$,
y, por tanto,
$H_q^{\mathrm I}(T(U_{\bullet,j}))=\bigoplus_{r+s=j}T(H_q^{\mathrm I}(M_{\bullet,rs}))$;
En otras palabras, el complejo simple
$H_q^{\mathrm I}(T(U_{\bullet\bullet}))$ es el complejo simple
asociado al bicomplejo
$T(H_q^{\mathrm I}(M_{\bullet\bullet\bullet})$.
Para todo $q$, $H_q^{\mathrm I}(M_{\bullet\bullet\bullet})$ es una resolución \emph{proyectiva} del complejo simple $H_q^{\mathrm I}(K_{\bullet\bullet})$ en la categoría $\boldsymbol{K}$; aplicando (\hyperref[0.11.6.3-fr]{11.6.3}), se ve que se tiene
\[
H_p^{\mathrm{II}}
\bigl(T(H_q^{\mathrm I}(M_{\bullet\bullet\bullet}))\bigr)
=\mathbf{L}_pT(H_q^{\mathrm I}(K_{\bullet\bullet})),
\]
y, por tanto, la sucesión espectral $\mathop{}^{(a')}E$.
Finalmente, las sucesiones espectrales $\mathop{}^{(c)}E$ y
$\mathop{}^{(d)}E$ se obtienen intercambiando los papeles de los
índices $i$ y $j$ en la definición del tricomplejo
$M_{\bullet\bullet\bullet}$ y aplicando los razonamientos anteriores al
tricomplejo resultante.
\end{proof}

Se llama a $\mathbf{L}_\bullet T(K_{\bullet\bullet})$ la \emph{hiperhomología} de
$T$ relativa al \emph{bicomplejo} $K_{\bullet\bullet}$.

\begin{remarks}[11.7.3]
\label{0.11.7.3}
\label{0.11.7.3-fr}
\emph{Observaciones}.
\begin{enumerate}
  \item[{\rm(i)}]
  De \sref{0.11.6.3} se deduce que
  $\mathbf{L}_\bullet T(K_{\bullet\bullet})$ es un \emph{funtor homológico} sobre la
  categoría de los bicomplejos $K_{\bullet\bullet}$ en $\boldsymbol{C}$ que están
  acotados inferiormente en cada grado.

\oldpage[0\textsubscript{III}]{43}
  \item[{\rm(ii)}]
  Sea $M_{\bullet\bullet\bullet}$ un tricomplejo en $\boldsymbol{C}$ tal que,
  para todo par $(r,s)$, $M_{rs,\bullet}$ sea una \emph{resolución} de $K_{rs}$,
  y tal que
  $L_nT(M_{ijk})=0$
  para toda terna $(i,j,k)$ y todo $n>0$.
  Entonces existe un isomorfismo
  $
\mathbf{L}_\bullet T(K_{\bullet\bullet})
\xrightarrow{\sim} H_\bullet(T(M_{\bullet\bullet\bullet}))$;
  En efecto, con la notación utilizada en la demostración de \sref{0.11.7.2},
  $N_{i,\bullet}$ es una resolución de $L_i$ tal que
  $L_nT(N_{ij})=0$ para todo $n>0$ y todo par $(i,j)$, por lo que
  basta aplicar \sref{0.11.6.3},~(iii).

  \item[{\rm(iii)}]
  Los resultados de \sref{0.11.7.2} se generalizan inmediatamente a los
  multifuntores covariantes.
  Por ejemplo, sea $\boldsymbol{C}'$ una segunda categoría abeliana en la que todo
  objeto sea cociente de un objeto proyectivo, sea
  $K'_{\bullet\bullet}$ un bicomplejo en $\boldsymbol{C}'$ cuyos dos grados estén
  acotados inferiormente, y sea $T$ un bifuntor covariante aditivo de
  $\boldsymbol{C}\times\boldsymbol{C}'$ en una categoría abeliana $\boldsymbol{C}''$.
  Si $L_\bullet$ y $L'_\bullet$ son los complejos simples asociados a
  $K_{\bullet\bullet}$ y $K'_{\bullet\bullet}$, respectivamente, se define
  la hiperhomología de $T$ relativa a los dos bicomplejos
  $K_{\bullet\bullet}$ y $K'_{\bullet\bullet}$ como
  $\mathbf{L}_\bullet T(L_\bullet,L'_\bullet)$;
  Aplicando \sref{0.11.6.4} y \sref{0.11.6.5}, se obtienen, como en
  \sref{0.11.7.2}, seis sucesiones espectrales que convergen a esta hiperhomología;
  dejamos al lector su formulación explícita.
\end{enumerate}
\end{remarks}

\subsection{Suplemento sobre la cohomología de los complejos simpliciales}
\label{subsection:0.11.8}
\label{subsection:0.11.8-fr}

\begin{env}[11.8.1]
\label{0.11.8.1-fr}
Sean $A$ un conjunto finito, $\Sigma(A)$ el conjunto de las sucesiones finitas $\sigma=(\alpha_0,\ldots,\alpha_h)$ de elementos de $A$ («símplices» de $A$); se pone $|\sigma|=\{\alpha_0,\ldots,\alpha_h\}$; se recuerda que el complejo de cadenas $\mathrm C_\bullet(A)$ es el grupo abeliano libre graduado engendrado por los elementos de $\Sigma(A)$, siendo $(\alpha_0,\ldots,\alpha_h)$ de \emph{grado} $h$, con una diferencial definida por $d(\alpha_0,\ldots,\alpha_h)=\sum_{j=0}^{h}(-1)^j(\alpha_0,\ldots,\widehat{\alpha}_j,\ldots,\alpha_h)$. El subgrupo $\mathrm D_\bullet(A)$ de $\mathrm C_\bullet(A)$ engendrado por las cadenas $\sigma=(\alpha_0,\ldots,\alpha_h)$ para las que dos de los $\alpha_i$ son iguales, y por las cadenas $\pi(\sigma)-\varepsilon_\pi.\sigma$, donde para toda permutación $\pi$, $\pi(\sigma)=(\alpha_{\pi(0)},\ldots,\alpha_{\pi(h)})$ y $\varepsilon_\pi$ es el signo de $\pi$, es un subcomplejo de $\mathrm C_\bullet(A)$ cuyos elementos se llaman cadenas \emph{degeneradas}; se pone $\mathrm L_\bullet(A)=\mathrm C_\bullet(A)/\mathrm D_\bullet(A)$, y se tiene un homomorfismo natural de complejos $p:\mathrm C_\bullet(A)\longrightarrow\mathrm L_\bullet(A)$. Por otra parte, se define un homomorfismo de complejos $j:\mathrm L_\bullet(A)\longrightarrow\mathrm C_\bullet(A)$ del modo siguiente: se ordena totalmente $A$; a la clase mód. $\mathrm D_\bullet(A)$ de un símplex $\sigma=(\alpha_0,\ldots,\alpha_h)$, se hace corresponder $0$ si dos de los $\alpha_i$ son iguales, y la sucesión $(\beta_0,\ldots,\beta_h)$ de los $\alpha_i$ ordenados \emph{en orden creciente} en caso contrario. Está claro que $p\circ
j$ es la \emph{identidad} de $\mathrm L_\bullet(A)$.
\end{env}

\begin{env}[11.8.2]
\label{0.11.8.2}
\label{0.11.8.2-fr}
Sea $B$ un segundo conjunto finito.
Si $d'$ y $d''$ son los diferenciales de
$\mathrm C_\bullet(A)$ y $\mathrm C_\bullet(B)$, el complejo producto
tensorial
$\mathrm C_\bullet(A)\otimes\mathrm C_\bullet(B)$
puede considerarse como el grupo abeliano libre engendrado por los elementos de
$\Sigma(A)\times\Sigma(B)$, con la diferencial
$d(\sigma,\tau)=(d'\sigma,\tau)+(-1)^{h+1}(\sigma,d''\tau)$ si $\operatorname{Card}|\sigma|=h+1$.

Los homomorfismos naturales
$\mathrm C_\bullet(A)\to\mathrm L_\bullet(A)$ y
$\mathrm C_\bullet(B)\to\mathrm L_\bullet(B)$ definen un homomorfismo
$p:\mathrm C_\bullet(A)\otimes\mathrm C_\bullet(B)\longrightarrow\mathrm L_\bullet(A)\otimes\mathrm L_\bullet(B)$,
siendo este último producto tensorial isomorfo a
$(\mathrm C_\bullet(A)\otimes\mathrm C_\bullet(B))/(\mathrm C_\bullet(A)\otimes\mathrm D_\bullet(B)+\mathrm D_\bullet(A)\otimes\mathrm C_\bullet(B))$.
Análogamente, utilizando los homomorfismos
$\mathrm L_\bullet(A)\to\mathrm C_\bullet(A)$ y
$\mathrm L_\bullet(B)\to\mathrm C_\bullet(B)$ definidos en
\sref{0.11.8.1}, con órdenes totales sobre $A$ y $B$, se obtiene un
homomorfismo
$j:\mathrm L_\bullet(A)\otimes\mathrm L_\bullet(B)\longrightarrow\mathrm C_\bullet(A)\otimes\mathrm C_\bullet(B)$
tal que $p\circ j$ es la identidad.
\end{env}

\oldpage[0\textsubscript{III}]{44}
Con esta notación:
\begin{proposition}[11.8.3]
\label{0.11.8.3}
\label{0.11.8.3-fr}
Existe una homotopía
$h:\mathrm C_\bullet(A)\otimes\mathrm C_\bullet(B)\longrightarrow\mathrm C_\bullet(A)\otimes\mathrm C_\bullet(B)$
tal que $h(\sigma,\tau)$ es una combinación lineal de pares de símplices
$(\sigma_i,\tau_i)$ que satisfacen
$|\sigma_i|\subset|\sigma|$ y $|\tau_i|\subset|\tau|$, y tal que,
si $f=j\circ p$, entonces
\[
\label{0.11.8.3.1-fr}
\tag{11.8.3.1}
f-1=h\circ d+d\circ h.
\]
\end{proposition}

\begin{proof}
Basta definir $h$ sobre cada par $(\sigma,\tau)$ de símplices, por
inducción sobre la suma de los grados de $\sigma$ y $\tau$, pues se puede tomar $h=0$ cuando esta suma es $0$. Sea
$\omega=f(\sigma,\tau)-(\sigma,\tau)-h(d(\sigma,\tau))$;
Por la hipótesis de inducción y la definición de $d$, se tiene
$\omega\in\mathrm C_\bullet(|\sigma|)\otimes\mathrm C_\bullet(|\tau|)$.
Además,
\[
d\omega=f(d(\sigma,\tau))-d(\sigma,\tau)-d(h(d(\sigma,\tau)))
=h(d(d(\sigma,\tau)))=0
\]
en virtud de (\hyperref[0.11.8.3.1-fr]{11.8.3.1}) y de la hipótesis de inducción. Ahora bien, se tiene $H_q(\mathrm C_\bullet(A))=0$ para $q>0$ (G, I, 3.7.4), por tanto también $H_q(\mathrm C_\bullet(A)\otimes\mathrm C_\bullet(B))=0$ para $q>0$, en virtud de la fórmula de Künneth (G, I, 5.5.2). Aplicando esta observación al reemplazar $A$ reemplazado por $|\sigma|$ y $B$ por
$|\tau|$, se ve que existe un elemento $\omega'$ de $\mathrm C_\bullet(|\sigma|)\otimes\mathrm C_\bullet(|\tau|)$ tal que $\omega=d\omega'$.
Tomar $h(\sigma,\tau)=\omega'$ verifica \eqref{0.11.8.3.1} para el par
$(\sigma,\tau)$ y permite continuar la inducción.
\end{proof}

\begin{env}[11.8.4]
\label{0.11.8.4}
\label{0.11.8.4-fr}
Con la notación de \sref{0.11.8.2}, escribamos
$(\sigma,\tau)\leq(\sigma',\tau')$
si $|\sigma|\subset|\sigma'|$ y $|\tau|\subset|\tau'|$.
Un \emph{sistema de coeficientes} $\mathscr S$ sobre
$\Sigma(A)\times\Sigma(B)$ consiste en una familia
$(\Gamma_{\sigma,\tau})$ de grupos abelianos, donde
$\Gamma_{\sigma,\tau}$ depende únicamente de los conjuntos $|\sigma|$ y
$|\tau|$, y una familia de homomorfismos $\Gamma_{\sigma,\tau}\longrightarrow\Gamma_{\sigma',\tau'}$, para $(\sigma,\tau)\leq(\sigma',\tau')$, que forman un \emph{sistema inductivo} para esta relación de preorden. Se define entonces un \emph{complejo de cocadenas} $\mathrm C^\bullet(A,B;\mathscr S)$ como el conjunto de las familias $\lambda=(\lambda(\sigma,\tau))$, donde $(\sigma,\tau)$ recorre $\Sigma(A)\times\Sigma(B)$, con $\lambda(\sigma,\tau)\in\Gamma_{\sigma,\tau}$ para todo par $(\sigma,\tau)$. La diferencial está dada del modo siguiente: si
$d(\sigma,\tau)=\sum_i\pm(\sigma_i,\tau_i)$,
entonces $|\sigma_i|\subset|\sigma|$ y $|\tau_i|\subset|\tau|$ para todo
$i$, y se pone
\[
d\lambda(\sigma,\tau)=\sum_i\pm\lambda_i(\sigma_i,\tau_i),
\]
donde $\lambda_i(\sigma_i,\tau_i)$ denota la imagen canónica de
$\lambda(\sigma_i,\tau_i)$ en $\Gamma_{\sigma,\tau}$. Diremos que una cocadena $\lambda\in\mathrm C^\bullet(A,B;\mathscr S)$ se llama
\emph{bialternada} si $\lambda(\sigma,\tau)=0$ siempre que $\sigma$
o $\tau$ tiene dos términos iguales, y si se tiene $\lambda(\pi(\sigma),\tau)=\varepsilon_\pi\lambda(\sigma,\tau)$ y $\lambda(\sigma,\pi'(\tau))=\varepsilon_{\pi'}\lambda(\sigma,\tau)$ para permutaciones cualesquiera $\pi$ y $\pi'$ de los índices.
Estas cocadenas engendran manifiestamente un \emph{subcomplejo}
$\mathrm L^\bullet(A,B;\mathscr S)$ de
$\mathrm C^\bullet(A,B;\mathscr S)$.
\end{env}

\begin{proposition}[11.8.5]
\label{0.11.8.5-fr}
{La inyección canónica
$\mathrm L^\bullet(A,B;\mathscr S)\longrightarrow\mathrm C^\bullet(A,B;\mathscr S)$
define un isomorfismo para la cohomología de estos dos complejos.}
\end{proposition}

\begin{proof}
Observemos que si $p$ y $j$ tienen el sentido definido en (\hyperref[0.11.8.2-fr]{11.8.2}), las aplicaciones ${}^tp:\lambda\to\lambda\circ p$ y ${}^tj:\lambda\to\lambda\circ j$ están definidas en $\mathrm L^\bullet(A,B;\mathscr S)$ y
$\mathrm C^\bullet(A,B;\mathscr S)$, respectivamente; la primera es precisamente
la inyección canónica.
Como ${}^tj\circ{}^tp$ es la identidad, solo queda demostrar
que ${}^tp\circ{}^tj$ es homótopa a la identidad.
Por \sref{0.11.8.3}, la aplicación
${}^th:\lambda\to\lambda\circ h$
está \emph{definida} sobre $\mathrm C^\bullet(A,B;\mathscr S)$, de modo que al trasponer
\eqref{0.11.8.3.1} se obtiene el resultado deseado.
\end{proof}

\begin{env}[11.8.6]
\label{0.11.8.6}
\label{0.11.8.6-fr}
\oldpage[0\textsubscript{III}]{45}
La Proposición~\sref{0.11.8.5} reduce el cálculo de la cohomología de
$\mathrm L^\bullet(A,B;\mathscr S)$ al de
$\mathrm C^\bullet(A,B;\mathscr S)$. Recordemos por otra parte que esta última es, en virtud del teorema de Eilenberg-Zilber (G, I, 3.10.2), canónicamente isomorfa a la cohomología del complejo de cocadenas definido como sigue: se forma el complejo de cadenas $\mathrm P_\bullet(A,B)$ de combinaciones lineales de pares
$(\sigma,\tau)\in\Sigma(A)\times\Sigma(B)$ para los cuales $\sigma$ y
$\tau$ tengan \emph{el mismo grado}; la diferencial de este complejo
viene dada por
$d:(\sigma,\tau)\leadsto\sum_{j,k}(-1)^{j+k}(\sigma_j,\tau_k)$ si
$d\sigma=\sum_j(-1)^j\sigma_j$ y $d\tau=\sum_k(-1)^k\tau_k$;
se tienen entonces dos homomorfismos canónicos de complejos
\[
f:\mathrm P_\bullet(A,B)\longrightarrow
\mathrm C_\bullet(A)\otimes\mathrm C_\bullet(B),\qquad
g:\mathrm C_\bullet(A)\otimes\mathrm C_\bullet(B)
\longrightarrow\mathrm P_\bullet(A,B),
\] y se demuestra (\emph{loc. cit.}) que hay \emph{homotopías} $h$ y $h'$ tales que
\[
f\circ g-1=d\circ h+h\circ d\qquad\text{y}\qquad
g\circ f-1=d\circ h'+h'\circ d.
\]
Además, se tiene $f(\sigma,\tau)\in\mathrm C_\bullet(|\sigma|)
\otimes\mathrm C_\bullet(|\tau|)$ y $g(\sigma,\tau)\in\mathrm P_\bullet(|\sigma|,|\tau|)$ y las homotopías $h$, $h'$ \emph{pueden tomarse de modo que} $h(\sigma,\tau)\in\mathrm C_\bullet(|\sigma|)
\otimes\mathrm C_\bullet(|\tau|)$ \emph{y} $h'(\sigma,\tau)\in\mathrm P_\bullet(|\sigma|,|\tau|)$. Este punto proviene de que la definición de $h(\sigma,\tau)$ y $h'(\sigma,\tau)$ pueden definirse por
\emph{inducción} sobre la suma de los grados de $\sigma$ y $\tau$, y de que los
$H^q(\mathrm C_\bullet(|\sigma|)\otimes\mathrm C_\bullet(|\tau|))$ y
$H^q(\mathrm P_\bullet(|\sigma|,|\tau|))$
son \emph{nulos} para $q>0$ (\emph{loc. cit.}); se razona entonces como en (\hyperref[0.11.8.3-fr]{11.8.3}) y resulta la conclusión. Se define entonces $\mathrm P^\bullet(A,B;\mathscr S)$ como el conjunto de las familias $\lambda=(\lambda(\sigma,\tau))$ donde $(\sigma,\tau)$ recorre los pares cuyos términos tienen el mismo grado, con $\lambda(\sigma,\tau)\in\Gamma_{\sigma,\tau}$, y como se tiene $d\sigma=\sum_j(-1)^j\sigma_j\in\mathrm C_\bullet(|\sigma|)$ y $d\tau=\sum_k(-1)^k\tau_k\in\mathrm C_\bullet(|\tau|)$,
\[
d\lambda(\sigma,\tau)=
\sum(-1)^{j+k}\lambda(\sigma_j,\tau_k)
\]
está definida y da el diferencial de
$\mathrm P^\bullet(A,B;\mathscr S)$. Dicho esto, las aplicaciones ${}^tf:\lambda\to\lambda\circ f$, ${}^tg:\lambda\to\lambda\circ g$, ${}^th:\lambda\to\lambda\circ h$ y ${}^th':\lambda\to\lambda\circ h'$ están todas \emph{definidas} en virtud de las observaciones precedentes;
${}^tf\circ{}^tg$ y ${}^tg\circ{}^tf$ son homótopas a la
identidad, dando el isomorfismo requerido entre la cohomología de
$\mathrm C^\bullet(A,B;\mathscr S)$ y la de
$\mathrm P^\bullet(A,B;\mathscr S)$.
\end{env}

\begin{env}[11.8.7]
\label{0.11.8.7}
\label{0.11.8.7-fr}
\emph{Observación}.
El razonamiento de \sref{0.11.8.3}, aplicado a
$\mathrm C_\bullet(A)$ y $\mathrm L_\bullet(A)$, muestra que si
$j$ y $p$ se definen como en \sref{0.11.8.1}, entonces
$f=j\circ p$ satisface de nuevo \eqref{0.11.8.3.1}, con
$|h(\sigma)|\subset|\sigma|$.
Como en \sref{0.11.8.5}, se obtiene un isomorfismo de la cohomología de
$\mathrm L^\bullet(A;\mathscr S)$ en la de
$\mathrm C^\bullet(A;\mathscr S)$, definiéndose los dos complejos de manera
evidente.
Este es el resultado cuya demostración se esboza en (G, I, 3.8.1).
\end{env}

\begin{env}[11.8.8]
\label{0.11.8.8}
\label{0.11.8.8-fr}
Volvamos ahora a la notación y las hipótesis de \sref{0.11.8.2}, y sea
$\mathscr S^\bullet=(\mathscr S^k)$
un \emph{complejo} \emph{de sistemas de coeficientes} sobre
$\Sigma(A)\times\Sigma(B)$.
Para cada $(\sigma,\tau)$, los grupos $\Gamma^k_{\sigma,\tau}$
forman, por tanto, un complejo de grupos abelianos $(k\in\mathbb Z)$, y se tienen los diagramas conmutativos
\[
\xymatrix{
\Gamma^k_{\sigma,\tau}\ar[r]\ar[d] &
\Gamma^{k+1}_{\sigma,\tau}\ar[d]\\
\Gamma^k_{\sigma',\tau'}\ar[r] &
\Gamma^{k+1}_{\sigma',\tau'}
}
\]
\oldpage[0\textsubscript{III}]{46}
para $(\sigma,\tau)\leq(\sigma',\tau')$. Entonces se verifica inmediatamente que
$\mathrm C^\bullet(A,B;\mathscr S^\bullet)
=\bigl(\mathrm C^h(A,B;\mathscr S^k)\bigr)_{(h,k)\in\mathbb Z\times\mathbb Z}$
es un \emph{bicomplejo} de grupos abelianos, y que
$\mathrm L^\bullet(A,B;\mathscr S^\bullet)
=\bigl(\mathrm L^h(A,B;\mathscr S^k)\bigr)$
es un subbicomplejo.
\end{env}

\begin{proposition}[11.8.9]
\label{0.11.8.9}
\label{0.11.8.9-fr}
La inyección canónica
$\mathrm L^\bullet(A,B;\mathscr S^\bullet)\to
\mathrm C^\bullet(A,B;\mathscr S^\bullet)$ define un isomorfismo para la cohomología de estos dos bicomplejos.
\end{proposition}

\begin{proof}
Pongamos $\mathrm C^{\bullet\bullet}=
\mathrm C^\bullet(A,B;\mathscr S^\bullet)$ y $\mathrm L^{\bullet\bullet}=
\mathrm L^\bullet(A,B;\mathscr S^\bullet)$ para simplificar, y observemos que, puesto que $\mathrm C^{hk}=\mathrm L^{hk}=0$ para $h<0$, las segundas
sucesiones espectrales de estos bicomplejos son \emph{regulares}
\sref{0.11.3.3}.
El homomorfismo
$\mathrm L^{\bullet\bullet}\to\mathrm C^{\bullet\bullet}$ da, por tanto,
un morfismo de sucesiones espectrales
${}''E(\mathrm L^{\bullet\bullet})\to
{}''E(\mathrm C^{\bullet\bullet})$
que, sobre los términos $E_2$, es
\[
\label{0.11.8.9.1-fr}
\tag{11.8.9.1}
H^p_{\mathrm{II}}(H^q_{\mathrm I}(\mathrm L^{\bullet\bullet}))
\longrightarrow
H^p_{\mathrm{II}}(H^q_{\mathrm I}(\mathrm C^{\bullet\bullet})).
\]
Sin embargo, para todo $k\in\mathbb Z$, \sref{0.11.8.3} muestra que
$H^q_{\mathrm I}(\mathrm L^{\bullet,k})\to
H^q_{\mathrm I}(\mathrm C^{\bullet,k})$
es \emph{biyectivo}.
La conclusión se deduce de \sref{0.11.1.5}.
\end{proof}

\begin{env}[11.8.10]
\label{0.11.8.10}
\label{0.11.8.10-fr}
Análogamente, con la notación de \sref{0.11.8.6}, existe un
homomorfismo canónico de bicomplejos
$\mathrm C^\bullet(A,B;\mathscr S^\bullet)\to
\mathrm P^\bullet(A,B;\mathscr S^\bullet)$
con la notación evidente.
El razonamiento de \sref{0.11.8.9}, utilizando ahora \sref{0.11.8.6}, muestra que
este homomorfismo induce asimismo un \emph{isomorfismo} en cohomología.
\end{env}

\subsection{Un lema sobre los complejos de tipo finito}
\label{subsection:0.11.9}
\label{subsection:0.11.9-fr}

\begin{proposition}[11.9.1]
\label{0.11.9.1}
\label{0.11.9.1-fr}
Sea $\boldsymbol{C}$ una categoría abeliana, y sean $\boldsymbol{K}'$ y $\boldsymbol{K}''$ subconjuntos
del conjunto de objetos de $\boldsymbol{C}$, con $\boldsymbol{K}''\subset\boldsymbol{K}'$, que satisfacen las
condiciones siguientes:
\begin{enumerate}
  \item[{\rm(i)}]
  Para todo $A'\in\boldsymbol{K}'$ y todo epimorfismo $u:A\to A'$ en
  $\boldsymbol{C}$, existen un objeto $B\in\boldsymbol{K}''$ y un morfismo $v:B\to A$ tales que $uv$ sea un epimorfismo.
\item[{\rm(ii)}] Para todo par de objetos $A\in\boldsymbol{K}'$, $B\in\boldsymbol{K}'$ y todo epimorfismo $u:A\to B$,
  el objeto $\operatorname{Ker}(u)$ pertenece a $\boldsymbol{K}'$.

  \item[{\rm(iii)}]
  El producto de dos objetos de $\boldsymbol{K}''$ pertenece a $\boldsymbol{K}''$.
\end{enumerate}

Sea $P_\bullet=(P_i)_{i\in\mathbb Z}$ un complejo en $\boldsymbol{C}$ tal
que $H_i(P_\bullet)\in\boldsymbol{K}'$ para todo $i$, y que exista $d$ para el cual $H_i(P_\bullet)=0$ cuando $i<d$.
Entonces existen un complejo $Q_\bullet=(Q_i)_{i\in\mathbb Z}$ en $\boldsymbol{C}$
tal que $Q_i\in\boldsymbol{K}''$ para todo $i$ y $Q_i=0$ para $i<d$, y un
homomorfismo de complejos
$u:Q_\bullet\to P_\bullet$
cuyo morfismo inducido
$H_\bullet(Q_\bullet)\to H_\bullet(P_\bullet)$
es un isomorfismo.
\end{proposition}

\begin{proof}
Demostremos primero la consecuencia siguiente de la condición~(i).

\medskip\noindent
(i \emph{bis})
\emph{Sea $u:C\to B$ un epimorfismo en $\boldsymbol{C}$, sea $A$ un
objeto de $\boldsymbol{K}'$, y sea $v:A\to B$ un morfismo en $\boldsymbol{C}$.
Entonces existen un objeto $D\in\boldsymbol{K}''$, un epimorfismo $u':D\to A$ y un
morfismo $v':D\to C$ tales que el diagrama}
\[
\xymatrix{
D\ar[r]^{u'}\ar[d]_{v'} & A\ar[d]^{v}\\
C\ar[r]_{u} & B
}
\tag{11.9.1.1}\label{0.11.9.1.1-fr}
\]
\emph{es conmutativo.}

Formemos el producto fibrado $C\times_B A$ en $\boldsymbol{C}$ y las proyecciones canónicas $p:C\times_B A\to C$, $q:C\times_B A\to A$, que hacen conmutativo el diagrama
\oldpage[0\textsubscript{III}]{47}
\[
\xymatrix{
C\times_B A\ar[r]^{q}\ar[d]_{\boldsymbol{p}} & A\ar[d]^{\boldsymbol{v}}\\
C\ar[r]_{\boldsymbol{u}} & B
}
\tag{11.9.1.2}\label{0.11.9.1.2-fr}
\]
es conmutativo.
El conúcleo de $q$ es el cociente de $A$ por $v^{-1}(u(C))$
([27, p.~1-12]).
Como $u$ es un epimorfismo, $u(C)=B$ y
$v^{-1}(u(C))=A$; por tanto, $q$ es un epimorfismo; basta entonces aplicar (i) al epimorfismo $q:C\times_B A\to A$, se obtienen un objeto
$D\in\boldsymbol{K}''$ y un morfismo $w:D\to C\times_B A$ tales que $qw$ sea un epimorfismo; se tomará $u'=qw$, $v'=pw$.

Dicho esto, para demostrar la proposición, procedamos por inducción. Supongamos, para un $i\geq d-1$, que para todo $j\leq i$, los objetos $Q_j$, los morfismos $d_j:Q_j\to Q_{j-1}$ y los morfismos $u_j:Q_j\to P_j$ de modo que $Q_j=0$ para $j<d$, que $d_{j-1}\circ d_j=0$ y $d_j\circ u_j=u_{j-1}\circ d_j$ para $j\leq i$; además, suponemos satisfechas las condiciones siguientes:
\begin{enumerate}
\item[$(\mathrm I_i)$]
  $Q_j\in\boldsymbol{K}''$ para $j\leq i$, y
  $B_j(Q_\bullet)\in\boldsymbol{K}'$ para $j<i$.

  \item[$(\mathrm{II}_i)$]
  Para $j<i$, el homomorfismo
  $H_j(Q_\bullet)\to H_j(P_\bullet)$
  inducido por $(u_k)_{k\leq i}$ es un \emph{isomorfismo}.

  \item[$(\mathrm{III}_i)$]
  El compuesto
  $v_i:Z_i(Q_\bullet)\to Z_i(P_\bullet)
\to H_i(P_\bullet)$
(donde la flecha de la izquierda es la restricción de $u_i$ y la flecha de la derecha el morfismo canónico) es un \emph{epimorfismo}.
\end{enumerate}

Observemos que, según (ii), $Z_i(Q_\bullet)$, núcleo del epimorfismo $Q_i\to B_{i-1}(Q_\bullet)$, pertenece a $\boldsymbol{K}'$ en virtud de la hipótesis $(\mathrm I_i)$. Se deduce además de (ii) que
$N_i=\operatorname{Ker}(v_i)$
pertenece a $\boldsymbol{K}'$, habida cuenta de la hipótesis $(\mathrm{III}_i)$. En virtud de (i \emph{bis}), existe un $Q'_{i+1}\in\boldsymbol{K}''$, un epimorfismo
$d'_{i+1}:Q'_{i+1}\to N_i$
y un morfismo $u'_{i+1}:Q'_{i+1}\to P_{i+1}$ tales que
\[
\xymatrix{
Q'_{i+1}\ar[r]^{d'_{i+1}}\ar[d]_{u'_{i+1}} &
N_i\ar[d]^{u_i}\\
P_{i+1}\ar[r]_{d_{i+1}} & B_i(P_\bullet)
}
\tag{11.9.1.3}\label{0.11.9.1.3-fr}
\]
es conmutativo.

El morfismo canónico
$Z_{i+1}(P_\bullet)\to H_{i+1}(P_\bullet)$
es un epimorfismo, y $H_{i+1}(P_\bullet)\in\boldsymbol{K}'$.
La condición~(i) proporciona, por tanto, un objeto $Q''_{i+1}\in\boldsymbol{K}''$ y un
morfismo
$u''_{i+1}:Q''_{i+1}\to Z_{i+1}(P_\bullet)$
tal que el compuesto
$Q''_{i+1}\to Z_{i+1}(P_\bullet)\to H_{i+1}(P_\bullet)$
es un epimorfismo.
Si $d''_{i+1}:Q''_{i+1}\to N_i$ igual a $0$, el diagrama
\[
\xymatrix{
Q''_{i+1}\ar[r]^{d''_{i+1}}\ar[d]_{u''_{i+1}} &
N_i\ar[d]^{u_i}\\
Z_{i+1}(P_\bullet)\ar[r]_{d_{i+1}} & P_i
}
\tag{11.9.1.4}\label{0.11.9.1.4-fr}
\]
\oldpage[0\textsubscript{III}]{48}
es conmutativo, siendo la flecha horizontal inferior $0$. Tomemos entonces $Q_{i+1}=Q'_{i+1}\times Q''_{i+1}$, que pertenece a $\boldsymbol{K}''$ en virtud de (iii), y
$d_{i+1}=d'_{i+1}+d''_{i+1}$, $u_{i+1}=u'_{i+1}+u''_{i+1}$.
Como $d_{i+1}(Q_{i+1})=d'_{i+1}(Q_{i+1})=N_i\subset Z_i(Q_\bullet)$, se tiene
$d_{i+1}\circ d_i=0$ y, con las notaciones usuales,
$B_i(Q_\bullet)=N_i$; esto demuestra
$(\mathrm I_{i+1})$.
La conmutatividad de \eqref{0.11.9.1.3} y
\eqref{0.11.9.1.4} da
$d_{i+1}\circ u_{i+1}=u_i\circ d_{i+1}$.
Por la definición de $N_i$, el morfismo
$H_i(Q_\bullet)=Z_i(Q_\bullet)/N_i\to H_i(P_\bullet)$
inducido por los $u_k$ ($k\leq i+1$) es el morfismo deducido de $v_i$ por paso a los cocientes, por tanto es un isomorfismo puesto que $v_i$ es un epimorfismo, de donde $(\mathrm{II}_{i+1})$. Finalmente, se tiene $Q''_{i+1}\subset Z_{i+1}(Q_\bullet)$ por definición; la elección de $u''_{i+1}$ muestra que
$v_{i+1}:Z_{i+1}(Q_\bullet)\to Z_{i+1}(P_\bullet)
\to H_{i+1}(P_\bullet)$
es un epimorfismo, porque su restricción a $Q''_{i+1}$ ya lo
es.
Así se cumple $(\mathrm{III}_{i+1})$, la inducción continúa y se deduce la
proposición.
\end{proof}

\begin{corollary}[11.9.2]
\label{0.11.9.2-fr}
Sean $A$ un anillo noetheriano (no necesariamente conmutativo), $P_\bullet=(P_i)_{i\in\mathbb Z}$ un complejo de $A$-módulos por la derecha. Se supone que los $H_i(P_\bullet)$ son $A$-módulos de tipo finito y que existe $d$ tal que $H_i(P_\bullet)=0$ para $i<d$. Entonces existe un complejo $Q_\bullet=(Q_i)_{i\in\mathbb Z}$ formado por $A$-módulos por la derecha libres de rango finito, tal que $Q_i=0$ para $i<d$, y un homomorfismo $u:Q_\bullet\to P_\bullet$ de complejos, tal que el homomorfismo $H_\bullet(Q_\bullet)\to H_\bullet(P_\bullet)$ correspondiente a $u$ sea biyectivo.
\end{corollary}

\begin{proof}
Se aplica (\hyperref[0.11.9.1-fr]{11.9.1}) tomando para $\boldsymbol{C}$ la categoría de los $A$-módulos por la derecha, para $\boldsymbol{K}'$ (resp. $\boldsymbol{K}''$) el conjunto de los $A$-módulos de tipo finito (resp. el conjunto de los $A$-módulos libres de rango finito); la verificación de las condiciones (i), (ii) y (iii) de (\hyperref[0.11.9.1-fr]{11.9.1}) es inmediata, habida cuenta de la hipótesis de que $A$ es noetheriano.
\end{proof}

\begin{env}[11.9.3]
\label{0.11.9.3}
\label{0.11.9.3-fr}
\emph{Observaciones}.
\begin{enumerate}
  \item[{\rm(i)}]
  Bajo las hipótesis de \sref{0.11.9.2}, supongamos además que los
  $P_i$ son $A$-módulos \emph{planos} a derecha.
  Entonces, para todo $A$-módulo a izquierda $M$, el homomorfismo de complejos
  $u\otimes 1:Q_\bullet\otimes_A M\to P_\bullet\otimes_A M$
  induce de nuevo un \emph{isomorfismo}
  $H_\bullet(Q_\bullet\otimes_A M)\xrightarrow{\sim}
H_\bullet(P_\bullet\otimes_A M)$ de la homología como lo veremos en el cap.~III.
\item[{\rm(ii)}] La conclusión de (\hyperref[0.11.9.2-fr]{11.9.2}) ya no es necesariamente exacta cuando no se supone $A$ \emph{noetheriano}; en efecto, al aplicarlo a un complejo reducido a $0$ salvo en un solo término, se concluiría que todo $A$-módulo por la izquierda finitamente engendrado admite una \emph{resolución} por
  módulos libres finitamente engendrados, lo que es falso en general
(cf. Bourbaki, \emph{Alg. comm.}, cap.~I, \S~2, ejerc.~6).

  Sin embargo, es posible reemplazar la hipótesis de que $A$ es
  noetheriano por la suposición de que los $H_i(P_\bullet)$ admiten
  $\infty$-presentaciones finitas (cf.\ Capítulo~IV).
\end{enumerate}
\end{env}

\subsection{Característica de Euler--Poincar\'e de un complejo de módulos de longitud finita}
\label{subsection:0.11.10}
\label{subsection:0.11.10-fr}

\begin{env}[11.10.1]
\label{0.11.10.1}
\label{0.11.10.1-fr}
Sea $A$ un anillo, no necesariamente conmutativo, y sea
\[
\label{0.11.10.1.1-fr}
\tag{11.10.1.1}
M^\bullet:0\longrightarrow M^0\longrightarrow M^1
\longrightarrow\ldots\longrightarrow M^n\longrightarrow0
\]
un \emph{complejo} de $A$-módulos a izquierda de \emph{longitud finita}.
La \emph{característica de Euler--Poincar\'e} de este complejo es el
entero
\oldpage[0\textsubscript{III}]{49}
\[
\label{0.11.10.1.2-fr}
\tag{11.10.1.2}
\chi(M^\bullet)=\sum_{i=0}^{n}(-1)^i\operatorname{long}M^i.
\]
\end{env}

\begin{proposition}[11.10.2]
\label{0.11.10.2}
\label{0.11.10.2-fr}
Para todo complejo finito $M^\bullet$ de $A$-módulos a izquierda de longitud
finita,
$\chi(M^\bullet)=\chi(H^\bullet(M^\bullet))$
donde $H^\bullet(M^\bullet)$ se considera como un complejo con diferencial
nulo.
En particular, si la sucesión \eqref{0.11.10.1.1} es exacta, entonces
$\chi(M^\bullet)=0$.
\end{proposition}

\begin{proof}
Pongamos para abreviar $B^i=B^i(M^\bullet)$, $Z^i=Z^i(M^\bullet)$, $H^i=H^i(M^\bullet)=Z^i/B^i$; los $B^i$, $Z^i$ y $H^i$ tienen todos
longitud finita.
Las sucesiones exactas
\[
0\longrightarrow B^i\longrightarrow Z^i\longrightarrow H^i
\longrightarrow0
\]
y
\[
0\longrightarrow Z^i\longrightarrow M^i\longrightarrow B^{i+1}
\longrightarrow0
\]
dan
\[
\operatorname{long}(Z^i)=\operatorname{long}(H^i)+
\operatorname{long}(B^i)
\]
y
\[
\operatorname{long}(M^i)=\operatorname{long}(Z^i)+
\operatorname{long}(B^{i+1})
\]
Por consiguiente,
\[
\operatorname{long}(M^i)-\operatorname{long}(H^i)=
\operatorname{long}(B^{i+1})+\operatorname{long}(B^i)
\]
Multipliquemos esta igualdad por $(-1)^i$ y hagamos la suma de las relaciones obtenidas para $0\leq i\leq n$.
Como $B^0=B^{n+1}=0$, se deduce la igualdad deseada.
\end{proof}

\begin{corollary}[11.10.3]
\label{0.11.10.3}
\label{0.11.10.3-fr}
Sea $E=(E_r^{pq})$ una sucesión espectral en la
categoría de los módulos sobre un anillo $A$.
Se supone que los $E_2^{pq}$ son $A$-módulos de longitud finita y que solo hay un número finito de pares $(p,q)$ tales que $E_2^{pq}\neq0$. Entonces las características de Euler-Poincaré
$\chi(E_r^{(\bullet)})$
de todos los complejos
$E_r^{(\bullet)}=(E_r^{(n)})_{n\in\mathbb Z}$
de \sref{0.11.1.1} son iguales.
Si, además, $E$ es débilmente convergente y si se pone $E_\infty^{(n)}=\bigoplus_{p+q=n}E_\infty^{pq}$ para todo $n\in\mathbb Z$, se tiene también $\chi(E_\infty^{(\bullet)})=\chi(E_2^{(\bullet)})$,
donde
$E_\infty^{(\bullet)}=(E_\infty^{(n)})_{n\in\mathbb Z}$
se considera como un complejo con diferencial nulo.
\end{corollary}

\begin{proof}
Observemos primero que si $E_2^{pq}=0$, entonces
$E_r^{pq}=0$ para $2\leq r\leq+\infty$.
Así, todo complejo $E_r^{(\bullet)}$ es finito y está formado por
$A$-módulos de longitud finita; la relación $\chi(E_r^{(\bullet)})=\chi(E_{r+1}^{(\bullet)})$ para todo $r$ finito resulta, pues, de (\hyperref[0.11.10.2-fr]{11.10.2}) y de la isomorfía entre $H^\bullet(E_r^{(\bullet)})$ y $E_{r+1}^{(\bullet)}$ (como complejos con derivación trivial). La hipótesis de que los $E_r^{pq}$ tienen longitud finita, para cada par
$(p,q)$, las sucesiones
$(B_k(E_2^{pq}))_{k\geq2}$ y $(Z_k(E_2^{pq})_{k\geq2}$ son estacionarias;
la hipótesis de que $E$ es débilmente convergente y de que $E_2^{pq}=0$ salvo para un número finito de pares $(p,q)$, implica, por tanto,
que algún entero $r\geq2$ satisface
$E_\infty^{pq}=E_r^{pq}$
para todo par $(p,q)$.
Se deduce la aserción relativa a
$\chi(E_\infty^{(\bullet)})$.
\end{proof}

% END tree/repo/editions/III1_DIPLOMATIC/ega0/ega0-11.tex


% BEGIN tree/repo/editions/III1_DIPLOMATIC/ega0/ega0-12.tex
\section{Suplemento sobre la cohomología de haces}
\label{section:0.12}
\label{section:0.12-fr}

\subsection{Cohomología de haces de módulos sobre espacios anillados}
\label{subsection:0.12.1}
\label{subsection:0.12.1-fr}

\begin{env}[12.1.1]
\label{0.12.1.1}
\label{0.12.1.1-fr}

Sea $(X,\mathcal O_X)$ un espacio anillado.
Recordemos que, para todo $\mathcal O_X$-módulo $\mathcal F$, se define la
\emph{cohomología} $H^\bullet(X,\mathcal F)$, que es un \emph{funtor cohomológico universal}
(T, 2.2) de la categoría $\boldsymbol{C}(X)$ de
$\mathcal O_X$-módulos en la categoría de grupos abelianos; es el funtor derivado
del funtor exacto por la izquierda
$\mathcal F\to\Gamma(X,\mathcal F)$.
El funtor $H^\bullet(X,\mathcal F)$ es isomorfo
\oldpage[0\textsubscript{III}]{50}
a la restricción a la categoría $\boldsymbol{C}(X)$ del funtor cohomológico
definido del mismo modo sobre la categoría de haces de grupos abelianos
en $X$ (G, II, 7.2.1).
\end{env}

\begin{env}[12.1.2]
\label{0.12.1.2}
\label{0.12.1.2-fr}
Pongamos $A=\Gamma(X,\mathcal O_X)$.
Como todo elemento de $A$ define un endomorfismo del grupo abeliano
$\Gamma(X,\mathcal F)$, define por funtorialidad un endomorfismo del
$\partial$-funtor $H^\bullet(X,\mathcal F)$.
Estos endomorfismos dotan a cada $H^p(X,\mathcal F)$ de una
estructura de $A$-\emph{módulo}, y el operador de conexión $\partial$ es
$A$-lineal. Además, para dos enteros positivos cualesquiera $p$, $q$, y dos $\mathcal O_X$-módulos $\mathcal F$, $\mathcal G$, se tiene un homomorfismo de $A$-módulos, llamado \emph{producto cup},
\[
\label{0.12.1.2.1-fr}
\tag{12.1.2.1}
H^p(X,\mathcal F)\otimes_A H^q(X,\mathcal G)
\longrightarrow H^{p+q}(X,\mathcal F\otimes_{\mathcal O_X}\mathcal G)
\]
(G, II, 6.6).
Estos homomorfismos hacen de la suma directa $S$ de los
$H^p(X,\mathcal O_X)$, para $p\geq0$, una $A$-\emph{álgebra graduada
anticomutativa}, y de la suma directa de los $H^p(X,\mathcal F)$ un
$S$-\emph{módulo graduado}.

Para todo recubrimiento abierto $\mathfrak U$ de $X$, denotaremos siempre
por $C^\bullet(\mathfrak U,\mathcal F)$ (contrariamente a
(G, II, 5.1)) el complejo de cocadenas alternadas del nervio de
$\mathfrak U$ con valores en el sistema de coeficientes
$\Gamma(U_\sigma,\mathcal F)$.
Está claro que
$C^\bullet(\mathfrak U,\mathcal F)$ es un
$A$-\emph{módulo graduado}, de modo que los grupos de cohomología
$H^p(\mathfrak U,\mathcal F)$ de este complejo están dotados de
una estructura de $A$-\emph{módulo}.
Además, las aplicaciones canónicas
$H^p(\mathfrak U,\mathcal F)\to H^p(X,\mathcal F)$
(G, II, 5.4) son $A$-\emph{lineales}.
\end{env}

\begin{env}[12.1.3]
\label{0.12.1.3}
\label{0.12.1.3-fr}
Sea $(X',\mathcal O_{X'})$ un segundo espacio anillado, y sea
$f=(\psi,\theta)$ un morfismo de $X'$ en $X$.
Pongamos $A'=\Gamma(X',\mathcal O_{X'})$; el $\psi$-morfismo $\theta$
define canónicamente un homomorfismo de anillos $A\to A'$.
Sea $\mathcal F$ un $\mathcal O_X$-módulo y $\mathcal F'$ un
$\mathcal O_{X'}$-módulo.
Para todo \emph{$f$-morfismo} $u:\mathcal F\to\mathcal F'$
\sref[0\textsubscript{I}]{0.4.4.1}, veremos que se puede definir,
para todo $p\geq0$, un dihomomorfismo
\[
\label{0.12.1.3.1-fr}
\tag{12.1.3.1}
u_p:H^p(X,\mathcal F)\longrightarrow H^p(X',\mathcal F').
\]

En efecto, como $\psi^*$ es exacto en la categoría de haces de grupos
abelianos sobre $X$, el funtor
$\mathcal F\leadsto H^\bullet(X',\psi^*(\mathcal F))$
es un $\partial$-funtor sobre dicha categoría, y se sabe que existe un
homomorfismo canónico de $\partial$-funtores
\[
\label{0.12.1.3.2-fr}
\tag{12.1.3.2}
H^\bullet(X,\mathcal F)\longrightarrow H^\bullet(X',\psi^*(\mathcal F))
\]
determinado únicamente por la condición de reducirse al homomorfismo canónico $\Gamma(X,\mathcal F)\to\Gamma(X',\psi^*(\mathcal F))$
(T, 3.2.2).
Además, todo elemento de $A$ determina un endomorfismo $\mu$ de
$\Gamma(X,\mathcal F)$ y un endomorfismo $\mu'$ de
$\Gamma(X',\psi^*(\mathcal F))$ tal que el diagrama
\[
\label{0.12.1.3.3-fr}
\tag{12.1.3.3}
\xymatrix{
\Gamma(X,\mathcal F)\ar[r]\ar[d]_\mu
&\Gamma(X',\psi^*(\mathcal F))\ar[d]^{\mu'}\\
\Gamma(X,\mathcal F)\ar[r]
&\Gamma(X',\psi^*(\mathcal F))
}
\]
es conmutativo.
\oldpage[0\textsubscript{III}]{51}
Por la propiedad de unicidad para prolongar morfismos de funtores
cohomológicos universales (T, 2.2), se obtienen prolongaciones únicas de
$\mu$ y $\mu'$ a la cohomología que hacen conmutativos los diagramas
análogos a \sref{0.12.1.3.3}.
Así, \sref{0.12.1.3.2} es un homomorfismo de $A$-\emph{módulos}.

Observemos ahora que
$f^*(\mathcal F)=\psi^*(\mathcal F)
\otimes_{\psi^*(\mathcal O_X)}\mathcal O_{X'}$,
y, por tanto, que existe un dihomomorfismo canónico
$\psi^*(\mathcal F)\to f^*(\mathcal F)$
del $\psi^*(\mathcal O_X)$-módulo $\psi^*(\mathcal F)$ en el
$\mathcal O_{X'}$-módulo $f^*(\mathcal F)$.
Por funtorialidad se obtiene, pues, un dihomomorfismo funtorial
\[
\label{0.12.1.3.4-fr}
\tag{12.1.3.4}
H^p(X',\psi^*(\mathcal F))\longrightarrow H^p(X',f^*(\mathcal F))
\]
siendo $A$ y $A'$ los anillos correspondientes.
Componiendo este dihomomorfismo con \sref{0.12.1.3.2}, se obtiene un
dihomomorfismo canónico, \emph{funtorial} en $\mathcal F$,
\[
\label{0.12.1.3.5-fr}
\tag{12.1.3.5}
\theta_p:H^p(X,\mathcal F)\longrightarrow H^p(X',f^*(\mathcal F)).
\]

Finalmente, por funtorialidad, el homomorfismo
$u^\sharp:f^*(\mathcal F)\to\mathcal F'$ da un homomorfismo de
$A'$-módulos
$H^p(X',f^*(\mathcal F))\to H^p(X',\mathcal F')$,
que, compuesto con \sref{0.12.1.3.5}, da
\sref{0.12.1.3.1}.

Sea $f'=(\psi',\theta'):X''\to X'$ un segundo morfismo de espacios
anillados, y sea $f''=f\circ f'$ el morfismo compuesto.
Teniendo en cuenta la compatibilidad del funtor $\psi'^*$ con los
productos tensoriales \sref[0\textsubscript{I}]{0.4.3.3}, se comprueba
inmediatamente que el compuesto del dihomomorfismo
$H^p(X',f^*(\mathcal F))\to H^p(X'',f'^*(f^*(\mathcal F)))$
con \sref{0.12.1.3.5} es el dihomomorfismo correspondiente
$H^p(X,\mathcal F)\to H^p(X'',f''^*(\mathcal F))$.
\end{env}

\begin{env}[12.1.4]
\label{0.12.1.4-fr}
Una definición directa del homomorfismo (\hyperref[0.12.1.3.2-fr]{12.1.3.2}) puede obtenerse de la manera siguiente: se considera una resolución inyectiva $\mathcal L^\bullet=(\mathcal L^i)$ de $\mathcal F$ constituida por haces de
grupos abelianos sobre $X$.
Como el funtor $\psi^*$ es exacto,
$\psi^*(\mathcal L^\bullet)$ es una \emph{resolución} de $\psi^*(\mathcal F)$
constituida por haces sobre $X'$.
Si $\mathcal L'^\bullet=(\mathcal L'^i)$ es una resolución inyectiva de
$\psi^*(\mathcal F)$ en la categoría de haces de grupos abelianos sobre
$X'$, existe un morfismo
$\psi^*(\mathcal L^\bullet)\to\mathcal L'^\bullet$
de complejos de haces de grupos abelianos, compatible con los
aumentos (M, V, 1.1.\emph{a)}, y determinado salvo homotopía.
Se obtienen así homomorfismos
\[
\Gamma(X,\mathcal L^\bullet)\longrightarrow
\Gamma(X',\psi^*(\mathcal L^\bullet))\longrightarrow
\Gamma(X',\mathcal L'^\bullet)
\]
de complejos de grupos abelianos.
Su compuesto, al pasar a la cohomología, da un morfismo de
$\partial$-funtores
$H^\bullet(X,\mathcal F)\to H^\bullet(X',\psi^*(\mathcal F))$; como coincide con (\hyperref[0.12.1.3.2-fr]{12.1.3.2}) en grado 0, es idéntico a él (T, 2.2).

Consideremos ahora un \emph{recubrimiento abierto}
$\mathfrak U=(U_\alpha)$ de $X$, y sea
$\mathfrak U'=(f^{-1}(U_\alpha))$ su recubrimiento abierto imagen inversa
de $X'$, imagen inversa de $\mathfrak U$. Los homomorfismos canónicos $\Gamma(V,\mathcal F)\to\Gamma(f^{-1}(V),f^*(\mathcal F))$ para todo abierto $V$ de $X$ definen inmediatamente (cf. G, II, 5.1) un homomorfismo de complejos
$C^\bullet(\mathfrak U,\mathcal F)\to
C^\bullet(\mathfrak U',f^*(\mathcal F))$,
y, por tanto, homomorfismos canónicos
\[
\label{0.12.1.4.1-fr}
\tag{12.1.4.1}
\theta_p:H^p(\mathfrak U,\mathcal F)
\longrightarrow H^p(\mathfrak U',f^*(\mathcal F)).
\]

\oldpage[0\textsubscript{III}]{52}
Además, existen diagramas conmutativos
\[
\label{0.12.1.4.2-fr}
\tag{12.1.4.2}
\xymatrix{
H^p(\mathfrak U,\mathcal F)\ar[r]^{\theta_p}\ar[d]
&H^p(\mathfrak U',f^*(\mathcal F))\ar[d]\\
H^p(X,\mathcal F)\ar[r]_{\theta_p}
&H^p(X',f^*(\mathcal F))
}
\]
donde las flechas verticales son los homomorfismos canónicos de
(G, II, 5.2).
Para demostrar la conmutatividad de \sref{0.12.1.4.2}, consideremos el
complejo de haces de cocadenas alternadas de $\mathcal F$ relativo a
$\mathfrak U$,
$\mathcal C^\bullet(\mathfrak U,\mathcal F)$, para el cual
$\Gamma(X,\mathcal C^\bullet(\mathfrak U,\mathcal F))=
C^\bullet(\mathfrak U,\mathcal F)$
(G, II, 5.2).
Los homomorfismos canónicos
$\Gamma(V,\mathcal F)\to
\Gamma(f^{-1}(V),\psi^*(\mathcal F))$
definen entonces un $\psi$-morfismo
$\mathcal C^\bullet(\mathfrak U,\mathcal F)\to
\mathcal C^\bullet(\mathfrak U',\psi^*(\mathcal F))$,
y, con la notación precedente, existe un diagrama conmutativo
\[
\xymatrix{
\Gamma(X,\mathcal C^\bullet(\mathfrak U,\mathcal F))\ar[r]\ar[d]
&\Gamma(X',\mathcal C^\bullet(\mathfrak U',\psi^*(\mathcal F)))\ar[d]\\
\Gamma(X,\mathcal L^\bullet)\ar[r]
&\Gamma(X',\psi^*(\mathcal L^\bullet))\ar[r]
&\Gamma(X',\mathcal L'^\bullet)
}
\]
Al pasar a la cohomología se obtienen diagramas conmutativos
\[
\xymatrix{
H^p(\mathfrak U,\mathcal F)\ar[r]\ar[d]
&H^p(\mathfrak U',\psi^*(\mathcal F))\ar[d]\\
H^p(X,\mathcal F)\ar[r]
&H^p(X',\psi^*(\mathcal F))
}
\]
donde las flechas verticales son los homomorfismos canónicos de
(G, II, 5.2).
Solo queda combinar estos diagramas con los diagramas conmutativos
\[
\xymatrix{
H^p(\mathfrak U',\psi^*(\mathcal F))\ar[r]\ar[d]
&H^p(\mathfrak U',f^*(\mathcal F))\ar[d]\\
H^p(X',\psi^*(\mathcal F))\ar[r]
&H^p(X',f^*(\mathcal F))
}
\]
\oldpage[0\textsubscript{III}]{53}
que provienen del homomorfismo
$\psi^*(\mathcal F)\to f^*(\mathcal F)$ y de la funtorialidad de los
homomorfismos canónicos de (G, II, 5.2), para obtener la conmutatividad
de \sref{0.12.1.4.2}.

Se observará que si
$A=\Gamma(X,\mathcal O_X)$ y $A'=\Gamma(X',\mathcal O_{X'})$, entonces
(12.1.4.3) es un dihomomorfismo de módulos correspondiente a
los anillos $A$ y $A'$.
Existe para \sref{0.12.1.4.1} una propiedad de transitividad respecto
de un compuesto de dos morfismos, análoga a la de
\sref{0.12.1.3.5}.
Finalmente, en las definiciones precedentes, en lugar de una resolución
inyectiva $\mathcal L^\bullet$ de $\mathcal F$, se habría podido partir igualmente de una resolución tal que $H^p(X,\mathcal L^i)=0$ para todo $i$ y todo $p>0$ (G, II, 4.7.1).
\end{env}

\begin{env}[12.1.5]
\label{0.12.1.5-fr}
Sean $\mathcal F$, $\mathcal G$, $\mathcal K$ tres $\mathcal O_X$-módulos, y consideremos
un $\mathcal O_X$-homomorfismo
$u:\mathcal F\otimes_{\mathcal O_X}\mathcal G\to\mathcal K$,
que da, en cohomología, los homomorfismos
\[
\label{0.12.1.5.1-fr}
\tag{12.1.5.1}
H^p(X,\mathcal F)\otimes_A H^q(X,\mathcal G)
\longrightarrow H^{p+q}(X,\mathcal K)
\]
deducidos del producto cup \sref{0.12.1.2.1}.
Demostraremos que, bajo las hipótesis y con la notación de
\sref{0.12.1.3}, los diagramas
\[
\label{0.12.1.5.2-fr}
\tag{12.1.5.2}
\xymatrix@C=10pt{
H^p(X,\mathcal F)\otimes_A H^q(X,\mathcal G)\ar[r]\ar[d]
&H^{p+q}(X,\mathcal K)\ar[d]\\
H^p(X',f^*(\mathcal F))\otimes_{A'}H^q(X',f^*(\mathcal G))\ar[r]
&H^{p+q}(X',f^*(\mathcal K))
}
\]
son conmutativos, donde las flechas verticales provienen de los homomorfismos canónicos (\hyperref[0.12.1.3.5-fr]{12.1.3.5}). Para ello, recordemos que (\hyperref[0.12.1.5.1-fr]{12.1.5.1}) puede obtenerse partiendo de las resoluciones \emph{canónicas} (G, II, 4.3) $\mathcal L^\bullet$, $\mathcal M^\bullet$, $\mathcal N^\bullet$ de $\mathcal F$, $\mathcal G$, $\mathcal K$ respectivamente, que están formadas por $\mathcal O_X$-módulos, y de la aplicación lineal
$\mathcal L^\bullet\otimes_{\mathcal O_X}\mathcal M^\bullet
\to\mathcal N^\bullet$
de complejos de $\mathcal O_X$-módulos correspondiente a $u$.
Esta aplicación proporciona un homomorfismo de complejos de $A$-módulos
$\Gamma(X,\mathcal L^\bullet)\otimes_A\Gamma(X,\mathcal M^\bullet)
\to\Gamma(X,\mathcal N^\bullet)$,
y, por tanto, al pasar a la cohomología, homomorfismos
$H^p(\Gamma(X,\mathcal L^\bullet))\otimes_A
H^q(\Gamma(X,\mathcal M^\bullet))
\to H^{p+q}(\Gamma(X,\mathcal N^\bullet))$
(G, II, 6.6).
Existe claramente un diagrama conmutativo
\[
\label{0.12.1.5.3-fr}
\tag{12.1.5.3}
\xymatrix@C=8pt{
\Gamma(X,\mathcal L^\bullet)\otimes_A\Gamma(X,\mathcal M^\bullet)
\ar[r]\ar[d]
&\Gamma(X,\mathcal N^\bullet)\ar[d]\\
\Gamma(X',\psi^*(\mathcal L^\bullet))
\otimes_{\Gamma(X',\psi^*(\mathcal O_X))}
\Gamma(X',\psi^*(\mathcal M^\bullet))\ar[r]
&\Gamma(X',\psi^*(\mathcal N^\bullet))
}
\]
\oldpage[0\textsubscript{III}]{54}
Al pasar a la cohomología se obtienen diagramas conmutativos
\[
\label{0.12.1.5.4-fr}
\tag{12.1.5.4}
\xymatrix@C=7pt{
H^p(X,\mathcal F)\otimes_A H^p(X,\mathcal G)\ar[r]\ar[d]
&H^{p+q}(X,\mathcal K)\ar[d]\\
H^p(\Gamma(X',\psi^*(\mathcal L^\bullet)))
\otimes_{\Gamma(X',\psi^*(\mathcal O_X))}
H^q(\Gamma(X',\psi^*(\mathcal M^\bullet)))\ar[r]
&H^{p+q}(\Gamma(X',\psi^*(\mathcal N^\bullet))).
}
\]
Pero como
$\psi^*(\mathcal L^\bullet)$,
$\psi^*(\mathcal M^\bullet)$, y
$\psi^*(\mathcal N^\bullet)$
son resoluciones de
$\psi^*(\mathcal F)$,
$\psi^*(\mathcal G)$, y
$\psi^*(\mathcal K)$ respectivamente, existe un diagrama conmutativo
(G, II, 6.6.1)
\[
\label{0.12.1.5.5-fr}
\tag{12.1.5.5}
\xymatrix@C=7pt{
H^p(\Gamma(X',\psi^*(\mathcal L^\bullet)))
\otimes_{\Gamma(X',\psi^*(\mathcal O_X))}
H^q(\Gamma(X',\psi^*(\mathcal M^\bullet)))\ar[r]\ar[d]
&H^{p+q}(\Gamma(X',\psi^*(\mathcal N^\bullet)))\ar[d]\\
H^p(X',\psi^*(\mathcal F))
\otimes_{\Gamma(X',\psi^*(\mathcal O_X))}
H^q(X',\psi^*(\mathcal G))\ar[r]
&H^{p+q}(X',\psi^*(\mathcal K))
}
\]
Finalmente, por funtorialidad, existe un diagrama conmutativo
\[
\label{0.12.1.5.6-fr}
\tag{12.1.5.6}
\xymatrix@C=7pt{
H^p(X',\psi^*(\mathcal F))
\otimes_{\Gamma(X',\psi^*(\mathcal O_X))}
H^q(X',\psi^*(\mathcal G))\ar[r]\ar[d]
&H^{p+q}(X',\psi^*(\mathcal K)).\ar[d]\\
H^p(X',f^*(\mathcal F))\otimes_{A'}H^q(X',f^*(\mathcal G))\ar[r]
&H^{p+q}(X',f^*(\mathcal K))
}
\]
Combinando los tres diagramas \sref{0.12.1.5.4},
\sref{0.12.1.5.5} y \sref{0.12.1.5.6}, se obtiene el diagrama
conmutativo deseado \sref{0.12.1.5.2}.
\end{env}

\begin{remark}[12.1.6]
\label{0.12.1.6}
\label{0.12.1.6-fr}
\oldpage[0\textsubscript{III}]{55}
Con la notación de \sref{0.12.1.3}, supongamos que existe un
diagrama conmutativo
\[
\label{0.12.1.6.1-fr}
\tag{12.1.6.1}
\xymatrix{
0\ar[r]&\mathcal F\ar[r]^r\ar[d]_u
&\mathcal G\ar[r]^s\ar[d]_v
&\mathcal H\ar[r]\ar[d]_w&0\\
0\ar[r]&\mathcal F'\ar[r]_{r'}
&\mathcal G'\ar[r]_{s'}
&\mathcal H'\ar[r]&0
}
\]
donde $r$, $s$ son homomorfismos de $\mathcal O_X$-módulos, $r'$, $s'$ homomorfismos de $\mathcal O_{X'}$-módulos, $u$, $v$, $w$ $f$-morfismos, y las filas son exactas.
Se sigue que el diagrama
\[
\label{0.12.1.6.2-fr}
\tag{12.1.6.2}
\xymatrix@C=6pt{
\ldots\ar[r]&H^p(X,\mathcal F)\ar[r]\ar[d]_{u_p}
&H^p(X,\mathcal G)\ar[r]\ar[d]_{v_p}
&H^p(X,\mathcal H)\ar[r]^\partial\ar[d]_{w_p}
&H^{p+1}(X,\mathcal F)\ar[r]\ar[d]_{u_{p+1}}&\ldots\\
\ldots\ar[r]&H^p(X',\mathcal F')\ar[r]
&H^p(X',\mathcal G')\ar[r]
&H^p(X',\mathcal H')\ar[r]_\partial
&H^{p+1}(X',\mathcal F')\ar[r]&\ldots
}
\]
es conmutativo.
En efecto, \sref{0.12.1.6.1} se factoriza como
\[
\xymatrix@R=14pt{
0\ar[r]&\mathcal F\ar[r]\ar[d]
&\mathcal G\ar[r]\ar[d]
&\mathcal H\ar[r]\ar[d]&0\\
0\ar[r]&\psi^*(\mathcal F)\ar[r]\ar[d]
&\psi^*(\mathcal G)\ar[r]\ar[d]
&\psi^*(\mathcal H)\ar[r]\ar[d]&0\\
0\ar[r]&\mathcal F'\ar[r]
&\mathcal G'\ar[r]
&\mathcal H'\ar[r]&0
}
\]
donde la fila intermedia es exacta
\sref[0\textsubscript{I}]{0.3.7.2}.
Basta entonces utilizar el hecho de que \sref{0.12.1.3.2} es un
homomorfismo de $\partial$-funtores y que los
$H^p(X',\mathcal F')$ forman un $\partial$-funtor en $\mathcal F'$.
\end{remark}

\begin{env}[12.1.7]
\label{0.12.1.7}
\label{0.12.1.7-fr}
Con las hipótesis y la notación de \sref{0.12.1.3}, consideremos ahora el
caso
$\mathcal F=f_*(\mathcal F')=\psi_*(\mathcal F')$; vamos a ver que el dihomomorfismo definido en (\hyperref[0.12.1.3-fr]{12.1.3}) \[
\label{0.12.1.7.1-fr}
\tag{12.1.7.1}
H^p(X,f_*(\mathcal F'))\longrightarrow H^p(X',\mathcal F')
\]
puede obtenerse (salvo un automorfismo de $H^p(X',\mathcal F')$) como un
homomorfismo de borde de una sucesión espectral para el funtor compuesto
$\mathcal F'\leadsto\Gamma(X',\psi_*(\mathcal F'))$ (T, 2.4). La descripción del homomorfismo (\hyperref[0.12.1.7.1-fr]{12.1.7.1}) dada en (\hyperref[0.12.1.4-fr]{12.1.4}) muestra aquí que se puede obtener este homomorfismo de la manera siguiente: se consideran resoluciones inyectivas $\mathcal L^\bullet$ y
$\mathcal L'^\bullet$ de $\psi_*(\mathcal F')$ y $\mathcal F'$,
respectivamente, y después un homomorfismo de complejos
$v:\psi^*(\mathcal L^\bullet)\to\mathcal L'^\bullet$
``por encima'' del homomorfismo canónico
$\psi^*(\psi_*(\mathcal F'))\to\mathcal F'$;
\oldpage[0\textsubscript{III}]{56}
se observa después que se tiene $\Gamma(X',\mathcal L'^\bullet)=\Gamma(X,\psi_*(\mathcal L'^\bullet))$
y que el homomorfismo compuesto
\[
\Gamma(X,\mathcal L^\bullet)\longrightarrow
\Gamma(X',\psi^*(\mathcal L^\bullet))
\xrightarrow{\Gamma(v)}\Gamma(X',\mathcal L'^\bullet)
\]
no es otro que
\[
\label{0.12.1.7.2-fr}
\tag{12.1.7.2}
\Gamma(v^\flat):\Gamma(X,\mathcal L^\bullet)
\longrightarrow\Gamma(X,\psi_*(\mathcal L'^\bullet))
\]
\sref[0\textsubscript{I}]{0.3.7.1}; y \sref{0.12.1.7.1} se
obtiene pasando a la cohomología en \sref{0.12.1.7.2}.

Por otra parte, las sucesiones espectrales para el funtor compuesto
$\mathcal F'\leadsto\Gamma(X,\psi_*(\mathcal F'))$
se obtienen tomando una resolución inyectiva de Cartan--Eilenberg
$\mathcal M^{\bullet\bullet}=(\mathcal M^{ij})$
del complejo $\psi_*(\mathcal L'^\bullet)$ en la categoría de
haces de grupos abelianos sobre $X$.
Las sucesiones espectrales en cuestión son las del bicomplejo
$\Gamma(X,\mathcal M^{\bullet\bullet})$; son birregulares, pues
$\mathcal M^{ij}=0$ para $i<0$ o $j<0$.
Ahora bien, la primera sucesión espectral de este bicomplejo degenera, pues
los haces $\psi_*(\mathcal L'^i)$ son flácidos (G, II, 3.1.1), por tanto $H_{\mathrm{II}}^q(\Gamma(X,\mathcal M^{i,\bullet}))
=H^q(\psi_*(\mathcal L'^i))=0$ para $q>0$ (G, II, 4.4.3); se tienen, pues, homomorfismos de borde \emph{biyectivos} (\hyperref[0.11.1.6-fr]{11.1.6})
\[
\label{0.12.1.7.3-fr}
\tag{12.1.7.3}
'E_2^{i0}=H^i(H_{\mathrm{II}}^0(\Gamma(X,\mathcal M^{\bullet\bullet})))
\longrightarrow H^i(\Gamma(X,\mathcal M^{\bullet\bullet}))
\]
y se sabe por \sref{0.11.3.4} que este homomorfismo proviene, al pasar
a la cohomología, del aumento
\[
\label{0.12.1.7.4-fr}
\tag{12.1.7.4}
\Gamma(X,\psi_*(\mathcal L'^\bullet))
\longrightarrow\Gamma(X,\mathcal M^{\bullet\bullet})
\]
que a su vez proviene del aumento
$\eta:\psi_*(\mathcal L'^\bullet)\to\mathcal M^{\bullet\bullet}$.

Para la segunda sucesión espectral, por otra parte, existen homomorfismos
de borde
\[
\label{0.12.1.7.5-fr}
\tag{12.1.7.5}
''E_2^{i0}=H^i(H_{\mathrm I}^0(\Gamma(X,\mathcal M^{\bullet\bullet}))
\longrightarrow H^i(\Gamma(X,\mathcal M^{\bullet\bullet}))
\]
que provienen, por \sref{0.11.3.4}, al pasar a la cohomología, del
homomorfismo de complejos
$Z_{\mathrm I}^0(\Gamma(X,\mathcal M^{\bullet\bullet}))
\to\Gamma(X,\mathcal M^{\bullet\bullet})$.
Como $\psi_*$ es exacto por la izquierda, la sucesión
\[
0\longrightarrow\psi_*(\mathcal F')\longrightarrow
\psi_*(\mathcal L'^0)\longrightarrow\psi_*(\mathcal L'^1)
\]
es exacta; por definición de una resolución de Cartan-Eilenberg (\hyperref[0.11.4.2-fr]{11.4.2}), se puede, pues, tomar $B_{\mathrm I}^0(\mathcal M^{\bullet\bullet})=0$, $Z_{\mathrm I}^0(\mathcal M^{\bullet\bullet})=\mathcal L^\bullet$; como el diagrama
\[
\xymatrix@C=50pt@R=24pt{
\mathcal L^0\ar[r]^{i^0}&\mathcal M^{00}\\
\psi_*(\mathcal F')\ar[u]^\varepsilon
\ar[r]_{\varepsilon'}\ar[ur]^{\varepsilon''}
&\psi_*(\mathcal L'^0)\ar[u]_{\eta^0}
}
\]
es conmutativo, la inyección de complejos
$i:\mathcal L^\bullet\to\mathcal M^{\bullet\bullet}$ es compatible con los
aumentos $\varepsilon$ y $\varepsilon''$.
Existen, pues, dos homomorfismos de complejos \textbf{de} $\mathcal L^\bullet$
en $\mathcal M^{\bullet\bullet}$,
\[
\xymatrix@C=52pt@R=22pt{
\mathcal L^\bullet\ar[r]^i\ar[dr]_{v^\flat}
&\mathcal M^{\bullet\bullet}\\
&\psi_*(\mathcal L'^\bullet)\ar[u]_\eta
}
\]
\oldpage[0\textsubscript{III}]{57}
compatibles con los aumentos $\varepsilon$ y
$\varepsilon''$.
Como $\mathcal L^\bullet$ es una resolución inyectiva y
$\mathcal M^{\bullet\bullet}$ está constituido por haces inyectivos, se sigue
de (M, V, 1.1 \textit{a)}) que estos dos homomorfismos son homotópicos.
Lo mismo vale, por tanto, para los dos homomorfismos correspondientes
$\Gamma(X,\mathcal L^\bullet)\to\Gamma(X,\mathcal M^{\bullet\bullet})$, y al pasar a la cohomología, se obtiene por tanto el \emph{mismo} homomorfismo; en otros términos, se ha mostrado que el homomorfismo de borde (\hyperref[0.12.1.7.5-fr]{12.1.7.5}), que se escribe $H^p(X,\psi_*(\mathcal F'))\to H^p(\Gamma(X,\mathcal M^{\bullet\bullet}))$
es el compuesto de \sref{0.12.1.7.1} y
\sref{0.12.1.7.3}, estando este último escrito
$H^p(X',\mathcal F')\to H^p(\Gamma(X,\mathcal M^{\bullet\bullet}))$ y que se ha visto que es un \emph{isomorfismo}; de donde nuestra aserción.
\end{env}

\subsection{Imágenes directas superiores}
\label{subsection:0.12.2-fr}

\begin{env}[12.2.1]
\label{0.12.2.1-fr}
Sean $(X,\mathcal O_X)$, $(Y,\mathcal O_Y)$ dos espacios anillados,
$f=(\psi,\theta)$ un morfismo de $X$ en $Y$, que define el funtor
\emph{imagen directa}
$f_*:\boldsymbol{C}(X)\longrightarrow\boldsymbol{C}(Y)$,
idéntico además a la restricción a $\boldsymbol{C}(X)$ del funtor $\psi_*$ definido
en la categoría de los haces de grupos abelianos sobre $X$. Este último
funtor es aditivo y exacto por la izquierda, y como todo haz de grupos
abelianos sobre $X$ es isomorfo a un subhaz de un haz \emph{inyectivo} de
grupos abelianos, se definen los \emph{funtores derivados} derechos
$\mathcal F\leadsto R^p\psi_*(\mathcal F)$ del funtor $\psi_*$; los
$R^p\psi_*(\mathcal F)$ son haces de grupos abelianos sobre $Y$, y
los $R^p\psi_*$ forman un \emph{funtor cohomológico universal} (T, 2.3).

Además, el haz $R^p\psi_*(\mathcal F)$ es el haz asociado al
prehaz $V\leadsto H^p(f^{-1}(V),\mathcal F)$ (T, 3.7.2). Si ahora
se supone que $\mathcal F$ es un $\mathcal O_X$-\emph{Módulo},
$H^p(f^{-1}(V),\mathcal F)$ está naturalmente provisto de una estructura de
$\Gamma(f^{-1}(V),\mathcal O_X)$-módulo, luego de
$\Gamma(V,\psi_*(\mathcal O_X))$-módulo, y el dato del homomorfismo
$\theta:\mathcal O_Y\to\psi_*(\mathcal O_X)$ permite deducir una
estructura de $\Gamma(V,\mathcal O_Y)$-módulo. Para las estructuras así
definidas, está claro que la restricción de un abierto $V$ a un abierto
$V'\subset V$ define un dihomomorfismo, y esto permite por tanto definir sobre
cada uno de los $R^p\psi_*(\mathcal F)$ una estructura de $\mathcal O_Y$-\emph{Módulo};
es este $\mathcal O_Y$-Módulo el que denotaremos por $R^p f_*(\mathcal F)$,
quedando así $R^p f_*$ definido como un funtor aditivo de $\boldsymbol{C}(X)$ en
$\boldsymbol{C}(Y)$. Además, los $R^p f_*$ forman un \emph{$\partial$-funtor}, pues si
$0\longrightarrow\mathcal F'\longrightarrow\mathcal F
\longrightarrow\mathcal F''\longrightarrow0$
es una sucesión exacta de $\mathcal O_X$-Módulos, la descripción de los
$R^p\psi_*$ y de la estructura de $\mathcal O_Y$-Módulo sobre
$R^p\psi_*(\mathcal F)$ dada arriba muestra inmediatamente que
el homomorfismo
$\partial:R^p\psi_*(\mathcal F'')\longrightarrow
R^{p+1}\psi_*(\mathcal F')$
es en este caso un homomorfismo de $\mathcal O_Y$-\emph{Módulos}. Finalmente, los
$R^p f_*$ se identifican con los funtores \emph{derivados} derechos de $f_*$: en efecto,
todo $\mathcal O_X$-Módulo admite una resolución \emph{inyectiva} formada por
$\mathcal O_X$-Módulos, y como tal resolución está formada por haces
\emph{flácidos} de grupos abelianos (G, II, 7.1), puede servir para calcular los
$R^p\psi_*(\mathcal F)$, puesto que $R^n\psi_*(\mathcal G)=0$ para $n\geq1$ y
para todo haz \emph{flácido} $\mathcal G$ (T, 2.4.1, Observación 3, y cor. de la
prop. 3.3.2). Se concluye pues que los $R^p f_*$ forman un \emph{funtor
cohomológico universal de} $\boldsymbol{C}(X)$ en $\boldsymbol{C}(Y)$ (T, 2.3).
\end{env}

\begin{env}[12.2.2]
\label{0.12.2.2}
\label{0.12.2.2-fr}
Sean $\mathcal F$ y $\mathcal G$ dos $\mathcal O_X$-módulos.
Con la notación de \sref{0.12.2.1}, para todo subconjunto abierto $V$ de
$Y$ existe el homomorfismo de producto cup \sref{0.12.1.2.1}
\[
H^p(f^{-1}(V),\mathcal F)
\otimes_{\Gamma(f^{-1}(V),\mathcal O_X)}
H^q(f^{-1}(V),\mathcal G)
\longrightarrow
H^{p+q}(f^{-1}(V),\mathcal F\otimes_{\mathcal O_X}\mathcal G)
\]
\oldpage[0\textsubscript{III}]{58}
Se sigue inmediatamente de la definición del producto cup
(G, II, 6.6) que estos homomorfismos conmutan con la restricción de
$V$ a un subespacio abierto $V'$ de $V$. Por otra parte, se tiene un homomorfismo de anillos
\[
\Gamma(V,\mathcal O_Y)\longrightarrow\Gamma(V,\psi_*(\mathcal O_X))
=\Gamma(f^{-1}(V),\mathcal O_X)
\]
proveniente de $\theta$, de donde un homomorfismo canónico de productos tensoriales
\[
\begin{aligned}
&H^p(f^{-1}(V),\mathcal F)\otimes_{\Gamma(V,\mathcal O_Y)}
H^q(f^{-1}(V),\mathcal G)\\
&\qquad\longrightarrow
H^p(f^{-1}(V),\mathcal F)
\otimes_{\Gamma(f^{-1}(V),\mathcal O_X)}
H^q(f^{-1}(V),\mathcal G)
\end{aligned}
\]
que es asimismo compatible con la restricción de $V'$ a $V$.
Por composición se obtiene así un homomorfismo de
$\Gamma(V,\mathcal O_Y)$-módulos.  Para los haces asociados a los
prehaces en cuestión, define un homomorfismo canónico,
funtorial en $\mathcal F$ y $\mathcal G$,
\[
\label{0.12.2.2.1-fr}
\tag{12.2.2.1}
R^p f_*(\mathcal F)\otimes_{\mathcal O_Y}R^q f_*(\mathcal G)
\longrightarrow
R^{p+q}f_*(\mathcal F\otimes_{\mathcal O_X}\mathcal G).
\]
Para $p=q=0$, este homomorfismo se reduce a
\sref[0\textsubscript{I}]{0.4.2.2.1}.
\end{env}

\begin{proposition}[12.2.3]
\label{0.12.2.3}
\label{0.12.2.3-fr}
Para todo $\mathcal O_X$-módulo $\mathcal F$ y todo $\mathcal O_Y$-módulo
$\mathcal L$ localmente libre de rango finito, existen isomorfismos canónicos
funtoriales
\[
\label{0.12.2.3.1-fr}
\tag{12.2.3.1}
R^p f_*(\mathcal F)\otimes_{\mathcal O_Y}\mathcal L
\xrightarrow{\sim}
R^p f_*(\mathcal F\otimes_{\mathcal O_X}f^*(\mathcal L)).
\]
\end{proposition}

\begin{proof}
El homomorfismo \sref{0.12.2.3.1} se obtiene componiendo el
siguiente caso particular de \sref{0.12.2.2.1},
\[
\label{0.12.2.3.2-fr}
\tag{12.2.3.2}
R^p f_*(\mathcal F)\otimes_{\mathcal O_Y}f_*(f^*(\mathcal L))
\longrightarrow
R^p f_*(\mathcal F\otimes_{\mathcal O_X}f^*(\mathcal L))
\]
con el homomorfismo del primer miembro de
\sref{0.12.2.3.1} al de \sref{0.12.2.3.2} inducido por el
homomorfismo canónico \sref[0\textsubscript{I}]{0.4.4.3.2}.
Para comprobar que \sref{0.12.2.3.1} es un isomorfismo cuando $\mathcal L$ es
localmente libre, puede reducirse inmediatamente al caso
$\mathcal L=\mathcal O_Y$, pues la cuestión es local sobre $Y$ y los
funtores en cuestión son aditivos en $\mathcal L$.
En ese caso, en vista de la definición de \sref{0.12.2.2.1}, la
proposición se reduce a comprobar que el homomorfismo correspondiente
de prehaces es biyectivo; esto se sigue inmediatamente de
$f^*(\mathcal O_Y)=\mathcal O_X$.
\end{proof}

\begin{env}[12.2.4]
\label{0.12.2.4}
\label{0.12.2.4-fr}
Sea $(Z,\mathcal O_Z)$ un tercer espacio anillado, y sea
$g:Y\to Z$ un morfismo de espacios anillados.
Se sabe (G, II, 7.1 y 3.1.1) que, para todo
$\mathcal O_X$-módulo inyectivo $\mathcal G$, $f_*(\mathcal G)$ es un haz flácido de grupos abelianos, y por consiguiente (\hyperref[0.12.2.1-fr]{12.2.1}) se tiene $R^p g_*(f_*(\mathcal G))=0$ para todo $p>0$. De ello se sigue (T, 2.4.1) que la \emph{sucesión espectral de Leray} de los funtores compuestos es aplicable al funtor compuesto $g_*f_*$: hay una sucesión espectral \emph{birregular} cuyo término límite es el funtor $R^\bullet h_*$, donde $h=g\circ f$, y cuyo término $E_2$ es
\[
\label{0.12.2.4.1-fr}
\tag{12.2.4.1}
E_2^{pq}=R^p g_*(R^q f_*(\mathcal F)).
\]
\end{env}

\begin{env}[12.2.5]
\label{0.12.2.5}
\label{0.12.2.5-fr}
Bajo las condiciones de \sref{0.12.2.4}, definiremos directamente
homomorfismos canónicos de $\mathcal O_Z$-módulos
\[
\label{0.12.2.5.1-fr}
\tag{12.2.5.1}
R^n g_*(f_*(\mathcal F))\longrightarrow R^n h_*(\mathcal F)
\]
y
\[
\label{0.12.2.5.2-fr}
\tag{12.2.5.2}
R^n h_*(\mathcal F)\longrightarrow g_*(R^n f_*(\mathcal F))
\]
\oldpage[0\textsubscript{III}]{59}
que pueden identificarse con los ``homomorfismos de borde'' de la sucesión
espectral de Leray (cf.\ \sref{0.12.1.7}).
Basta trabajar con los prehaces a los que están asociados los haces de
imágenes directas superiores \sref{0.12.2.1}.
Para ello, sea $W$ un subconjunto abierto cualquiera de $Z$ y su imagen inversa $g^{-1}(W)$ en $Y$.  Existe un
dihomomorfismo canónico
\[
\label{0.12.2.5.3-fr}
\tag{12.2.5.3}
H^n(g^{-1}(W),f_*(\mathcal F))\longrightarrow
H^n(f^{-1}(g^{-1}(W)),f^*(f_*(\mathcal F)))
\]
siendo los anillos correspondientes
$\Gamma(g^{-1}(W),\mathcal O_Y)$ y
$\Gamma(h^{-1}(W),\mathcal O_X)$.
Por otra parte, el homomorfismo canónico
\sref[0\textsubscript{I}]{0.4.4.3.3} da, por funtorialidad,
homomorfismos canónicos
\[
\label{0.12.2.5.4-fr}
\tag{12.2.5.4}
H^n(h^{-1}(W),f^*(f_*(\mathcal F)))\longrightarrow
H^n(h^{-1}(W),\mathcal F)
\]
que son homomorfismos de
$\Gamma(h^{-1}(W),\mathcal O_X)$-módulos.
Teniendo en cuenta el homomorfismo de anillos
$\Gamma(W,\mathcal O_Z)\to\Gamma(g^{-1}(W),\mathcal O_Y)$, se ve que, componiendo (\hyperref[0.12.2.5.4-fr]{12.2.5.4}) y (\hyperref[0.12.2.5.3-fr]{12.2.5.3}), se obtiene un homomorfismo de prehaces, que proporciona el homomorfismo de haces (\hyperref[0.12.2.5.1-fr]{12.2.5.1}).

La definición de (\hyperref[0.12.2.5.2-fr]{12.2.5.2}) es aún más simple; por definición, $R^n h_*(\mathcal F)$ está asociado al prehaz
$W\leadsto H^n(f^{-1}(g^{-1}(W)),\mathcal F)$
mientras que $R^n f_*(\mathcal F)$ está asociado al prehaz
$V\leadsto H^n(f^{-1}(V),\mathcal F)$;
Existe, por tanto, un homomorfismo canónico
\[
H^n(f^{-1}(g^{-1}(W)),\mathcal F)\longrightarrow
\Gamma(g^{-1}(W),R^n f_*(\mathcal F)),
\]
\textbf{y} es inmediato que estos homomorfismos definen un homomorfismo de
prehaces, que a su vez define \sref{0.12.2.5.2}.
\end{env}

\begin{env}[12.2.6]
\label{0.12.2.6}
\label{0.12.2.6-fr}
Bajo las hipótesis de \sref{0.12.2.4}, sean $\mathcal F$, $\mathcal G$
y $\mathcal K$ tres $\mathcal O_X$-módulos, y sea
$u:\mathcal F\otimes\mathcal G\to\mathcal K$
un $\mathcal O_X$-homomorfismo.
Entonces los diagramas siguientes son conmutativos:
\[
\xymatrix@C=10pt@R=24pt{
R^p g_*(f_*(\mathcal F))\otimes_{\mathcal O_Z}R^q g_*(f_*(\mathcal G))
\ar[r]\ar[d]&R^{p+q}g_*(f_*(\mathcal K))\ar[d]\\
R^p h_*(\mathcal F)\otimes_{\mathcal O_Z}R^q h_*(\mathcal G)
\ar[r]&R^{p+q}h_*(\mathcal K)
}
\tag{12.2.6.1}\label{0.12.2.6.1-fr}
\]
\textbf{y}
\[
\xymatrix@C=10pt@R=24pt{
R^p h_*(\mathcal F)\otimes_{\mathcal O_Z}R^q h_*(\mathcal G)
\ar[r]\ar[d]&R^{p+q}h_*(\mathcal K)\ar[d]\\
g_*(R^p f_*(\mathcal F))\otimes_{\mathcal O_Z}g_*(R^q f_*(\mathcal G))
\ar[r]&g_*(R^{p+q}f_*(\mathcal K))
}
\tag{12.2.6.2}\label{0.12.2.6.2-fr}
\]
\oldpage[0\textsubscript{III}]{60}
Aquí las flechas horizontales provienen de \sref{0.12.2.2.1} (la última
combinada con \sref[0\textsubscript{I}]{0.4.2.2.1}), y las flechas
verticales de \sref{0.12.2.5.1} y \sref{0.12.2.5.2}, respectivamente.

En efecto, basta verificar la aserción para los homomorfismos de prehaces
correspondientes.  Volviendo a las definiciones de
\sref{0.12.2.2} y \sref{0.12.2.5}, se reduce inmediatamente, para
\sref{0.12.2.6.1}, a los diagramas conmutativos
\sref{0.12.1.5.2}; la verificación es aún más sencilla para
\sref{0.12.2.6.2}.
\end{env}

\subsection{Suplementos sobre los funtores Ext de haces}
\label{subsection:0.12.3}
\label{subsection:0.12.3-fr}

\begin{env}[12.3.1]
\label{0.12.3.1-fr}
Consideremos un espacio anillado $(X,\mathcal O_X)$; no volveremos sobre la definición y las principales propiedades de los bifuntores $\operatorname{Ext}_{\mathcal O_X}^p(X;\mathcal F,\mathcal G)$ de la categoría de los $\mathcal O_X$-módulos en la de los $\Gamma(X,\mathcal O_X)$-módulos, y $\mathcal{E}xt_{\mathcal O_X}^p(\mathcal F,\mathcal G)$ de la categoría de los $\mathcal O_X$-módulos en sí misma, ni sobre la sucesión espectral birregular $E(\mathcal F,\mathcal G)$ que los relaciona (T, 4.2 y G, II, 7.3).
\end{env}

\begin{env}[12.3.2]
\label{0.12.3.2}
\label{0.12.3.2-fr}
Al igual que en (M, XIV, 1), se define la noción de una
\emph{extensión} de un $\mathcal O_X$-módulo $\mathcal F$ por un
$\mathcal O_X$-módulo $\mathcal G$ y la ley de composición sobre las clases de
extensiones equivalentes: los argumentos utilizados para los módulos se
trasladan manifiestamente a una categoría abeliana arbitraria.
La segunda demostración de (M, XIV, 1.1), que solo utiliza la existencia de
inmersiones en objetos inyectivos, sigue siendo, por tanto, válida en la
categoría de $\mathcal O_X$-módulos.  Muestra que
$\operatorname{Ext}_{\mathcal O_X}^1(X;\mathcal F,\mathcal G)$
se identifica canónicamente con el \emph{grupo abeliano de clases de
extensiones de $\mathcal F$ por $\mathcal G$}.
\end{env}

\begin{proposition}[12.3.3]
\label{0.12.3.3}
\label{0.12.3.3-fr}
Sea $(X,\mathcal O_X)$ un espacio anillado cuyo haz de anillos
$\mathcal O_X$ es coherente.
Entonces, para todo par de $\mathcal O_X$-módulos coherentes $\mathcal F$ y
$\mathcal G$ y todo $p\geq0$,
$\mathcal{E}xt_{\mathcal O_X}^p(\mathcal F,\mathcal G)$
es un $\mathcal O_X$-Módulo coherente.
\end{proposition}

\begin{proof}
Observemos que los $\mathcal{E}xt_{\mathcal O_X}^p(\mathcal F,\mathcal G)$ forman un funtor cohomológico
contravariante en $\mathcal F$.
Como $\mathcal F$ es coherente, para todo $p$ y todo punto $x\in X$
existen una vecindad abierta $U$ de $X$ y una sucesión exacta
de $(\mathcal O_X|U)$-módulos
\[
0\longrightarrow\mathcal R\longrightarrow\mathcal L_{p-1}
\longrightarrow\ldots\longrightarrow\mathcal L_0
\longrightarrow\mathcal F|U\longrightarrow0
\]
donde cada $\mathcal L_i$ ($0\leq i\leq p\setmin 1$) es isomorfo a
$\mathcal O_X^{n_i}|U$, y $\mathcal R$ es coherente.
Esto se sigue por inducción sobre $p$ de
\sref[0\textsubscript{I}]{0.5.3.2} y
\sref[0\textsubscript{I}]{0.5.3.4}, pues $\mathcal O_X$ es coherente. Observemos ahora que, para $p\geq1$,
$\operatorname{Ext}_{\mathcal O_X|U}^p(\mathcal L|U,\mathcal G|U)=0$
para todo $\mathcal O_X$-módulo $\mathcal L$ tal que $\mathcal L|U$ sea
isomorfo a $\mathcal O_X^n|U$ (T, 4.2.3).
El argumento de (M, V, 7.2) se aplica, por tanto, al funtor cohomológico
contravariante
$\mathcal F\leadsto
\mathrm{Hom}_{\mathcal O_X|U}(\mathcal F|U,\mathcal G|U)$
y da una sucesión exacta
\[
\mathcal{H}om_{\mathcal O_X|U}(\mathcal L_{p-1},\mathcal G|U)
\longrightarrow
\mathcal{H}om_{\mathcal O_X|U}(\mathcal R,\mathcal G|U)
\longrightarrow
\mathcal{E}xt_{\mathcal O_X|U}^p(\mathcal F|U,\mathcal G|U)
\longrightarrow0
\] y como los dos primeros términos de esta sucesión son $(\mathcal O_X|U)$-módulos coherentes
\sref[0\textsubscript{I}]{0.5.3.5}; por tanto, también lo es el tercero
\sref[0\textsubscript{I}]{0.5.3.4}.
\end{proof}

\oldpage[0\textsubscript{III}]{61}
\begin{proposition}[12.3.4]
\label{0.12.3.4}
\label{0.12.3.4-fr}
Sea $f:X\to Y$ un morfismo \textup{plano} de espacios anillados, y sean
$\mathcal F$ y $\mathcal G$ dos $\mathcal O_Y$-módulos.
\begin{enumerate}
  \item[\rm(i)] Existe un homomorfismo de bifuntores cohomológicos
  \[
\label{0.12.3.4.1-fr}
\tag{12.3.4.1}
f^*(\mathcal{E}xt_{\mathcal O_Y}^{\bullet}(\mathcal F,\mathcal G))
\xrightarrow{\sim}
\mathcal{E}xt_{\mathcal O_X}^{\bullet}(f^*(\mathcal F),f^*(\mathcal G))
\]
  que en grado $0$ se reduce al homomorfismo canónico
  \sref[0\textsubscript{I}]{0.4.4.6}.

  \item[\rm(ii)] Existe un morfismo canónico de sucesiones espectrales
  \[
\label{0.12.3.4.2-fr}
\tag{12.3.4.2}
E(\mathcal F,\mathcal G)\longrightarrow
E(f^*(\mathcal F),f^*(\mathcal G))
\]
  cuya aplicación sobre los términos $E_2$ viene dada por los homomorfismos
  \[
\label{0.12.3.4.3-fr}
\tag{12.3.4.3}
H^p(Y,\mathcal{E}xt_{\mathcal O_Y}^q(\mathcal F,\mathcal G))
\longrightarrow
H^p(X,\mathcal{E}xt_{\mathcal O_X}^q(f^*(\mathcal F),f^*(\mathcal G)))
\]
  deducidos de \sref{0.12.3.4.1} y \sref{0.12.1.3.1}.
\end{enumerate}
\end{proposition}

\begin{proof}
\begin{enumerate}
  \item[\rm(i)]
  Como $f^*$ es exacto sobre la categoría de $\mathcal O_Y$-módulos (\hyperref[0.6.7.2-fr]{\textbf{0}\textsubscript{I}, 6.7.2}), los funtores $\mathcal G\leadsto f^*(\mathcal{H}om_{\mathcal O_Y}(\mathcal F,\mathcal G))$ y $\mathcal G\leadsto\mathcal{H}om_{\mathcal O_X}(f^*(\mathcal F),f^*(\mathcal G))$ son exactos por la izquierda; se deduce canónicamente de (\hyperref[0.4.4.6-fr]{\textbf{0}\textsubscript{I}, 4.4.6}) un homomorfismo de sus funtores derivados. Para calcular estos últimos, se toma una resolución inyectiva $\mathcal L^\bullet=(\mathcal L^i)$ de $\mathcal G$.  Se tienen entonces
  homomorfismos
  $\mathcal H^p(f^*(\mathcal{H}om_{\mathcal O_Y}(\mathcal F,\mathcal L^\bullet)))\longrightarrow\mathcal H^p(\mathcal{H}om_{\mathcal O_X}(f^*(\mathcal F),f^*(\mathcal L^\bullet)))$
  entre los haces de cohomología de complejos.
  Por la exactitud de $f^*$, se tiene $\mathcal H^p(f^*(\mathcal{H}om_{\mathcal O_Y}(\mathcal F,\mathcal L^\bullet)))=f^*(\mathcal H^p(\mathcal{H}om_{\mathcal O_Y}(\mathcal F,\mathcal L^\bullet)))=f^*(\mathcal{E}xt_{\mathcal O_Y}^p(\mathcal F,\mathcal G))$ por definición. Por otra parte, la exactitud de $f^*$ implica también que
  $f^*(\mathcal L^\bullet)$ es una \emph{resolución} de $f^*(\mathcal G)$.
  Si $\mathcal L'^\bullet=(\mathcal L'^{i})$ es una resolución inyectiva
  de $f^*(\mathcal G)$ en la categoría de $\mathcal O_X$-módulos, existe un
  homomorfismo de complejos
  $f^*(\mathcal L^\bullet)\to\mathcal L'^\bullet$,
  determinado salvo homotopía, y por tanto un homomorfismo bien definido en
  cohomología.  Su compuesto con el homomorfismo precedente es
  \sref{0.12.3.4.1}.

  \item[\rm(ii)]
  Con la notación precedente, existe un homomorfismo de complejos de
  $\mathcal O_X$-módulos
  $f^*(\mathcal{H}om_{\mathcal O_Y}(\mathcal F,\mathcal L^\bullet))\longrightarrow\mathcal{H}om_{\mathcal O_X}(f^*(\mathcal F),\mathcal L'^\bullet)$.
  Sea $\mathcal M^{\bullet\bullet}$ una resolución inyectiva de
  Cartan--Eilenberg de
  $\mathcal{H}om_{\mathcal O_Y}(\mathcal F,\mathcal L^\bullet)$ en la categoría de
  $\mathcal O_Y$-módulos.  Como $f^*$ es exacto,
  $f^*(\mathcal M^{\bullet\bullet})$ es una resolución de Cartan--Eilenberg de
  $f^*(\mathcal{H}om_{\mathcal O_Y}(\mathcal F,\mathcal L^\bullet))$.
  Si $\mathcal M'^{\bullet\bullet}$ es una resolución inyectiva de
  Cartan--Eilenberg de
  $\mathcal{H}om_{\mathcal O_X}(f^*(\mathcal F),\mathcal L'^\bullet))$, existe
  \sref{0.11.4.2} un homomorfismo, determinado salvo homotopía,
  $f^*(\mathcal M^{\bullet\bullet})\longrightarrow\mathcal M'^{\bullet\bullet}$
  compatible con el homomorfismo precedente; en otras palabras, existe
  un $f$-morfismo
  $\mathcal M^{\bullet\bullet}\to
\mathcal M'^{\bullet\bullet}$ de
  de bicomplejos de haces.
  Induce un dihomomorfismo
  $\Gamma(Y,\mathcal M^{\bullet\bullet})\to
\Gamma(X,\mathcal M'^{\bullet\bullet})$
  de bicomplejos de módulos, determinado salvo homotopía, y por tanto
  un morfismo bien definido de sucesiones espectrales
  \sref{0.11.3.2}.  Este es precisamente el morfismo deseado
  \sref{0.12.3.4.2}, y la descripción
  \sref{0.12.3.4.3} se sigue inmediatamente de las definiciones.
\end{enumerate}
\end{proof}

\begin{proposition}[12.3.5]
\label{0.12.3.5}
\label{0.12.3.5-fr}
Bajo las hipótesis de \sref{0.12.3.4}, supongamos además que el
haz de anillos $\mathcal O_Y$ es coherente.
Entonces, para todo $\mathcal O_Y$-módulo coherente $\mathcal F$, los
homomorfismos canónicos \sref{0.12.3.4.1} son biyectivos.
\end{proposition}

\begin{proof}
\oldpage[0\textsubscript{III}]{62}
Como la cuestión es local sobre $Y$, se puede suponer que existe una
sucesión exacta
$0\longrightarrow\mathcal R\longrightarrow\mathcal O_Y^n\longrightarrow\mathcal F\longrightarrow0$,
Entonces $\mathcal R$ es también un $\mathcal O_Y$-módulo coherente
\sref[0\textsubscript{I}]{0.5.3.4}.
Para demostrar que los homomorfismos
\[
f^*(\mathcal{E}xt_{\mathcal O_Y}^p(\mathcal F,\mathcal G))
\longrightarrow
\mathcal{E}xt_{\mathcal O_X}^p(f^*(\mathcal F),f^*(\mathcal G))
\]
son biyectivos, razonamos por inducción sobre $p$.
Para $p=0$, la aserción es
\sref[0\textsubscript{I}]{0.6.7.6.1}.
Existe un diagrama conmutativo con filas exactas
\[
\hbox{\tiny$
\xymatrix@C=2pt@R=24pt{
f^*(\mathcal{E}xt_{\mathcal O_Y}^{p-1}(\mathcal O_Y^n,\mathcal G))
\ar[r]\ar[d]&
f^*(\mathcal{E}xt_{\mathcal O_Y}^{p-1}(\mathcal R,\mathcal G))
\ar[r]^\partial\ar[d]&
f^*(\mathcal{E}xt_{\mathcal O_Y}^{p}(\mathcal F,\mathcal G))
\ar[r]\ar[d]&
f^*(\mathcal{E}xt_{\mathcal O_Y}^{p}(\mathcal O_Y^n,\mathcal G))
\ar[d]\\
\mathcal{E}xt_{\mathcal O_X}^{p-1}(\mathcal O_X^n,f^*(\mathcal G))
\ar[r]&
\mathcal{E}xt_{\mathcal O_X}^{p-1}(f^*(\mathcal R),f^*(\mathcal G))
\ar[r]^\partial&
\mathcal{E}xt_{\mathcal O_X}^{p}(f^*(\mathcal F),f^*(\mathcal G))
\ar[r]&
\mathcal{E}xt_{\mathcal O_X}^{p}(\mathcal O_X^n,f^*(\mathcal G))
}
$}
\] puesto que $f^*(\mathcal O_Y)=\mathcal O_X$; como $f^*$ es exacto, las dos filas son exactas. Además, se tiene $\mathcal{E}xt_{\mathcal O_Y}^p(\mathcal O_Y^n,\mathcal G)=0$ para todo $p>0$ y análogamente $\mathcal{E}xt_{\mathcal O_X}^p(\mathcal O_X^n,f^*(\mathcal G))=0$ para todo $p>0$ (T, 4.2.3). Habida cuenta de la hipótesis de inducción, las dos primeras flechas verticales del diagrama precedente son isomorfismos, y los términos de la derecha son 0, por tanto
$f^*(\mathcal{E}xt_{\mathcal O_Y}^p(\mathcal F,\mathcal G))\longrightarrow\mathcal{E}xt_{\mathcal O_X}^p(f^*(\mathcal F),f^*(\mathcal G))$
es un isomorfismo.
\end{proof}

\subsection{Hipercohomología del funtor imagen directa}
\label{subsection:0.12.4}
\label{subsection:0.12.4-fr}

\begin{env}[12.4.1]
\label{0.12.4.1}
\label{0.12.4.1-fr}
Sean $(X,\mathcal O_X)$ y $(Y,\mathcal O_Y)$ dos espacios anillados, y sea
$f:X\to Y$ un morfismo de espacios anillados.
Se puede formar la hipercohomología de $f_*$ respecto de un complejo
\emph{arbitrario} de $\mathcal O_X$-módulos
$\mathcal K^\bullet=(\mathcal K^i)_{i\in\mathbb Z}$
(\hyperref[0.11.4.4-fr]{11.4.4}), pues en la categoría abeliana de los $\mathcal O_Y$-Módulos los límites inductivos filtrantes existen y son exactos (T, 3.1.1). Los $\mathcal O_Y$-módulos de hipercohomología
$\mathbf{R}^p f_*(\mathcal K^\bullet)$ se denotarán también
$\mathscr H^p(f,\mathcal K^\bullet)$, o simplemente
$\mathscr H_f^p(\mathcal K^\bullet)$.
Recordemos que $\mathscr H^\bullet(f,\mathcal K^\bullet)$ es la cohomología del
bicomplejo de $\mathcal O_Y$-módulos $f_*(\mathcal L^{\bullet\bullet})$,
donde $\mathcal L^{\bullet\bullet}$ es una resolución inyectiva de
Cartan--Eilenberg de $\mathcal K^\bullet$ en la categoría de
$\mathcal O_X$-módulos; $\mathscr H^\bullet(f,\mathcal K^\bullet)$ es el término límite de dos sucesiones espectrales ${}'\mathcal E(f,\mathcal K^\bullet)$ y ${}''\mathcal E(f,\mathcal K^\bullet)$ cuyos términos $E_2$
vienen dados por
\begin{equation}
\label{0.12.4.1.1-fr}
\tag{12.4.1.1}
{}'\mathcal E_2^{pq}=\mathcal H^p(\mathcal H^q(f,\mathcal K^\bullet))
\end{equation}
y
\begin{equation}
\label{0.12.4.1.2-fr}
\tag{12.4.1.2}
{}''\mathcal E_2^{pq}=\mathcal H^p(f,\mathcal H^q(\mathcal K^\bullet))
\quad(=R^p f_*(\mathcal H^q(\mathcal K^\bullet)))
\end{equation}
En estas fórmulas se ha adoptado la notación general $T(A^\bullet)$ para el transformado de un complejo por un funtor (\hyperref[0.11.2.1-fr]{11.2.1}), y se escribe
$\mathcal H^p(f,\mathcal F)$ en lugar de $R^p f_*(\mathcal F)$ para un $\mathcal O_X$-módulo $\mathcal F$. Recordemos además que la sucesión ${}'\mathcal E(f,\mathcal K^\bullet)$ es siempre \emph{regular}; las dos sucesiones espectrales ${}'\mathcal E(f,\mathcal K^\bullet)$ y ${}''\mathcal E(f,\mathcal K^\bullet)$ son \emph{birregulares} cuando $\mathcal K^\bullet$ está
acotado inferiormente,
\oldpage[0\textsubscript{III}]{63}
o cuando existe un entero $m$ tal que todo $\mathcal O_X$-módulo
admite una resolución flácida de longitud $\leq m$ \sref{0.11.4.4}.
\end{env}

\begin{env}[12.4.2]
\label{0.12.4.2}
\label{0.12.4.2-fr}
Denotaremos análogamente por $\mathbf H^\bullet(X,\mathcal K^\bullet)$ la
hipercohomología del funtor $\Gamma$ respecto de un complejo
$\mathcal K^\bullet$ de $\mathcal O_X$-módulos.
Así, los $\mathbf H^p(X,\mathcal K^\bullet)$ son
$\Gamma(X,\mathcal O_X)$-módulos.
Además, $\mathbf H^\bullet(X,\mathcal K^\bullet)$ puede considerarse como un
caso particular de $\mathscr H^\bullet(f,\mathcal K^\bullet)$, donde $f$ es un
morfismo de $(X,\mathcal O_X)$ a un espacio anillado constituido por un punto
dotado del anillo $\Gamma(X,\mathcal O_X)$.

Para todo subconjunto abierto $V$ de $X$, escribiremos
$\mathbf H^\bullet(V,\mathcal K^\bullet)$ en lugar de
$\mathbf H^\bullet(V,\mathcal K^\bullet|V)$.
\end{env}

\begin{proposition}[12.4.3]
\label{0.12.4.3}
\label{0.12.4.3-fr}
Para todo entero $p\in\mathbb Z$, el $\mathcal O_Y$-módulo
$\mathscr H^p(f,\mathcal K^\bullet)$ es canónicamente isomorfo al haz
asociado al prehaz
$U\leadsto\mathbf H^p(f^{-1}(U),\mathcal K^\bullet)$
sobre $Y$.
\end{proposition}

\begin{proof}
Con la notación de \sref{0.12.4.1}, el haz de cohomología
$\mathcal H^p(f_*(\mathcal L^{\bullet\bullet}))$ está asociado al
prehaz
$U\leadsto H^p(\Gamma(U,f_*(\mathcal L^{\bullet\bullet})))=H^p(\Gamma(f^{-1}(U),\mathcal L^{\bullet\bullet}))$.
Pero está claro que
$\mathcal L^{\bullet\bullet}|f^{-1}(U)$ es una resolución inyectiva de
Cartan--Eilenberg de $\mathcal K^\bullet|f^{-1}(U)$ (T, 3.1.3).
Por consiguiente,
$H^p(\Gamma(f^{-1}(U),\mathcal L^{\bullet\bullet}))=\mathbf H^p(f^{-1}(U),\mathcal K^\bullet)$
por definición.
\end{proof}

\begin{proposition}[12.4.4]
\label{0.12.4.4}
\label{0.12.4.4-fr}
La hipercohomología $\mathscr H^\bullet(f,\mathcal K^\bullet)$ es un
funtor cohomológico de $\mathcal K^\bullet$ en cada uno de los casos
siguientes:
\begin{enumerate}
  \item[\rm{a)}] $\mathcal K^\bullet$ varía en la categoría de complejos
    acotados inferiormente;
  \item[\rm{b)}] existe un entero $m$ tal que todo
    $\mathcal O_X$-módulo admite una resolución flácida de longitud $\leq m$;
  \item[\rm{c)}] $X$ es un espacio noetheriano.
\end{enumerate}
\end{proposition}

\begin{proof}
Los casos~\emph{a)} y~\emph{b)} son casos particulares de \sref{0.11.5.4}.
El caso~\emph{c)}, por otra parte, se sigue de \sref{0.11.5.2}, pues en
este caso el funtor $f_*$ conmuta con los límites inductivos
(G, II, 3.10.1).
\end{proof}

\begin{env}[12.4.5]
\label{0.12.4.5-fr}
Consideremos ahora un recubrimiento abierto $\mathfrak U=(U_\alpha)$ un recubrimiento abierto de $X$.
Para todo complejo de \emph{prehaces}
$\mathcal K^\bullet=(\mathcal K^j)$ sobre $X$, el bicomplejo $C^\bullet(\mathfrak U,\mathcal K^\bullet)$ cuya
componente de bigrado $(i,j)$ es
$C^i(\mathfrak U,\mathcal K^j)$, el grupo de
$i$-cocadenas alternadas del nervio de $\mathfrak U$ con valores
en $\mathcal K^j$ (G, II, 5.1).
Llamaremos a la cohomología de este bicomplejo la
\emph{hipercohomología} del recubrimiento $\mathfrak U$ con
coeficientes en $\mathcal K^\bullet$, y la denotaremos por
$\mathbf H^\bullet(\mathfrak U,\mathcal K^\bullet)=H^\bullet(C^\bullet(\mathfrak U,\mathcal K^\bullet))$.
La sucesión espectral de Leray de un recubrimiento
(T, 3.8.1 y G, II, 5.9.1) se generaliza a la hipercohomología como sigue.
\end{env}

\begin{proposition}[12.4.6]
\label{0.12.4.6}
\label{0.12.4.6-fr}
Sea $\mathcal K^\bullet=(\mathcal K^i)$ un complejo de
$\mathcal O_X$-módulos.
Existe un funtor espectral regular de $\mathcal K^\bullet$ con
término límite $\mathbf H^\bullet(X,\mathcal K^\bullet)$ y término $E_2$
\begin{equation}
\label{0.12.4.6.1-fr}
\tag{12.4.6.1}
\mathbf{E}_2^{pq}=\mathbf H^p(\mathfrak U,h^q(\mathcal K^\bullet))
\end{equation}
donde $h^q(\mathcal K^\bullet)$ denota el complejo de prehaces
$V\leadsto H^q(V,\mathcal K^\bullet)$
sobre $X$.
Esta sucesión espectral es birregular si $\mathcal K^\bullet$ está acotado inferiormente.
\end{proposition}

\begin{proof}
Consideremos una resolución inyectiva de Cartan-Eilenberg $\mathcal L^{\bullet\bullet}$ una resolución inyectiva de
Cartan--Eilenberg de $\mathcal K^\bullet$, y el tricomplejo $C^\bullet(\mathfrak U,\mathcal L^{\bullet\bullet})
=(C^i(\mathfrak U,\mathcal L^{jk}))$
Consideremos primero este tricomplejo como un bicomplejo en los grados $i$ y
$j+k$.
Como $i$ solo toma valores $\geq0$, la segunda sucesión espectral de este bicomplejo es regular (\hyperref[0.11.3.3-fr]{11.3.3}) y degenerada, pues se tiene $\mathbf H^q(\mathfrak U,\mathcal L^{jk})=0$ para todo $q>0$, siendo los $\mathcal O_X$-módulos $\mathcal L^{jk}$ haces flácidos
\oldpage[0\textsubscript{III}]{64}
(G, II, 5.2.3).
Se sigue que existe un isomorfismo canónico
$H^n(C^\bullet(\mathfrak U,\mathcal L^{\bullet\bullet}))\xrightarrow{\sim}H^n(\Gamma(X,\mathcal L^{\bullet\bullet}))$
\sref{0.11.1.6}, en virtud de (G, II, 5.2.2), y por tanto, por la
definición \sref{0.12.4.2}, un isomorfismo
$H^n(C^\bullet(\mathfrak U,\mathcal L^{\bullet\bullet}))\xrightarrow{\sim}\mathbf H^n(X,\mathcal K^\bullet)$.

Consideremos ahora
$C^\bullet(U,L^{\bullet\bullet})$
como un bicomplejo en los grados $i+j$ y $k$.
Como $k$ solo toma valores $\geq0$, la primera sucesión espectral de este bicomplejo es siempre regular; es birregular si $\mathcal L^{jk}=0$ para $j<j_0$, es decir, cuando
$\mathcal K^\bullet$ está acotado inferiormente \sref{0.11.3.3}.
Esta es la sucesión espectral deseada.
En efecto, para todo $j$, $\mathcal L^{j,\bullet}$ es una resolución
inyectiva de $\mathcal K^j$; por consiguiente,
$H^q(C^i(\mathfrak U,\mathcal L^{j,\bullet}))$
es precisamente el complejo de cocadenas
$C^i(\mathfrak U,h^q(\mathcal K^j))$.
Esto demuestra la proposición.
\end{proof}

\begin{corollary}[12.4.7]
\label{0.12.4.7}
\label{0.12.4.7-fr}
Si, para todo símplex $\sigma$ del nervio de $\mathfrak U$ y
todo entero $i$, se tiene $H^q(U_\sigma,\mathcal K^i)=0$ para $q>0$, entonces se tiene un isomorfismo canónico
\begin{equation}
\label{0.12.4.7.1-fr}
\tag{12.4.7.1}
\mathbf H^\bullet(\mathfrak U,\mathcal K^\bullet)
\xrightarrow{\sim}\mathbf H^\bullet(X,\mathcal K^\bullet).
\end{equation}
\end{corollary}

\begin{proof}
La hipótesis implica que
$C^\bullet(\mathfrak U,h^q(\mathcal K^j))=0$ para $q>0$, y por tanto que
$E_2^{pq}=0$ para $q>0$.
Como la sucesión espectral \sref{0.12.4.6.1} es degenerada y
regular, la conclusión se sigue de la definición
\sref{0.12.4.5} de
$\mathbf H^\bullet(\mathfrak U,\mathcal K^\bullet)$ y de
\sref{0.11.1.6}.
\end{proof}

\begin{env}[12.4.8]
\label{0.12.4.8}
\label{0.12.4.8-fr}
Sea $(X',\mathcal O_{X'})$ un segundo espacio anillado, y sea
$f=(\psi,\theta)$ un morfismo de $X'$ en $X$.
Por el mismo método que en \sref{0.12.1.3} y \sref{0.12.1.4}, se
define un dihomomorfismo para la hipercohomología de un complejo
$\mathcal K^\bullet$ de $\mathcal O_X$-módulos:
\begin{equation}
\label{0.12.4.8.1-fr}
\tag{12.4.8.1}
\mathbf H^p(X,\mathcal K^\bullet)
\longrightarrow\mathbf H^p(X',f^*(\mathcal K^\bullet)).
\end{equation}
Partiendo de una resolución inyectiva de Cartan--Eilenberg
$\mathcal L^{\bullet\bullet}$ de $\mathcal K^\bullet$, y como $\psi^*$ es exacto, $\psi^*(\mathcal L^{\bullet\bullet})$ es una
resolución de Cartan--Eilenberg de $\psi^*(\mathcal K^\bullet)$ en la
categoría de $\psi^*(\mathcal O_X)$-módulos.
Existe, por tanto, un morfismo
$\psi^*(\mathcal L^{\bullet\bullet})\to
\mathcal L'^{\bullet\bullet}$,
donde $\mathcal L'^{\bullet\bullet}$ es una resolución inyectiva de
Cartan--Eilenberg de $\psi^*(\mathcal K^\bullet)$.
Induce un morfismo en cohomología
$\mathbf H^\bullet(X,\mathcal K^\bullet)\longrightarrow\mathbf H^\bullet(X',\psi^*(\mathcal K^\bullet))$
Componiéndolo con el morfismo inducido por funtorialidad a partir de
$\psi^*(\mathcal K^\bullet)\to f^*(\mathcal K^\bullet)$, se obtiene el morfismo (\hyperref[0.12.4.8.1-fr]{12.4.8.1}) buscado.

Partiendo de (\hyperref[0.12.4.8.1-fr]{12.4.8.1}) y de (\hyperref[0.12.4.3-fr]{12.4.3}), se puede entonces, razonando como en (\hyperref[0.12.2.5-fr]{12.2.5}), definir, para dos morfismos $f:X\to Y$, $g:Y\to Z$ de espacios anillados, homomorfismos para la hipercohomología de un complejo $\mathcal K^\bullet$ de $\mathcal O_X$-módulos:
\begin{equation}
\label{0.12.4.8.2-fr}
\tag{12.4.8.2}
\mathscr H^n(g,f_*(\mathcal K^\bullet))
\longrightarrow\mathscr H^n(h,\mathcal K^\bullet)
\end{equation}
y
\begin{equation}
\label{0.12.4.8.3-fr}
\tag{12.4.8.3}
\mathscr H^n(h,\mathcal K^\bullet)
\longrightarrow g_*(\mathscr H^n(f,\mathcal K^\bullet)).
\end{equation}
Dejamos al lector los detalles de las definiciones.
\end{env}

% END tree/repo/editions/III1_DIPLOMATIC/ega0/ega0-12.tex


% BEGIN tree/repo/editions/III1_DIPLOMATIC/ega0/ega0-13.tex
\section{Límites proyectivos en álgebra homológica}
\label{section:0.13}
\label{section:0.13-fr}

\subsection{La condición de Mittag-Leffler}
\label{subsection:0.13.1}
\label{subsection:0.13.1-fr}

\begin{env}[13.1.1]
\label{0.13.1.1}
\label{0.13.1.1-fr}
Sea $\boldsymbol{C}$ una categoría abeliana en la que existan los productos
infinitos (axioma AB~3*) de T, 1.5).
Entonces existe el ínfimo de una familia de subobjetos de un objeto de
$\boldsymbol{C}$, y todo sistema proyectivo de objetos de $\boldsymbol{C}$ posee un
límite proyectivo; este último es un funtor exacto por la izquierda del
sistema proyectivo en cuestión (T, 1.8).
\oldpage[0\textsubscript{III}]{65}
Sea $(A_\alpha,f_{\alpha\beta})$ un sistema proyectivo de objetos de
$\boldsymbol{C}$ cuyo conjunto de índices $I$ es \emph{filtrante} a la derecha; sean $A=\varprojlim A_\alpha$
y, para cada $\alpha\in I$, sea
$f_\alpha:A\to A_\alpha$ el morfismo canónico.
Para todo $\alpha\in I$, los subobjetos
$f_{\alpha\beta}(A_\beta)$, para $\alpha\leq\beta$, forman una familia
filtrante decreciente de subobjetos de $A_\alpha$.
El subobjeto
$A'_\alpha=\inf_{\beta\geq\alpha}f_{\alpha\beta}(A_\beta)$
se llama el subobjeto de las \emph{imágenes universales} en $A_\alpha$; está claro que $f_\alpha(A)\subset A'_\alpha$ y $f_{\alpha\beta}(A'_\beta)\subset A'_\alpha$ para $\alpha\leq\beta$; por tanto
$(A'_\alpha,f_{\alpha\beta}|A'_\beta)$ es un sistema proyectivo y
$A=\varprojlim A'_\alpha$.
\end{env}

\begin{env}[13.1.2]
\label{0.13.1.2}
\label{0.13.1.2-fr}
Para un sistema proyectivo $(A_\alpha,f_{\alpha\beta})$ en $\boldsymbol{C}$,
se llama \emph{condición de Mittag-Leffler} a la siguiente:
\begin{itemize}
  \item[(ML)]
  \emph{Para todo índice $\alpha$, existe un $\beta\geq\alpha$ tal
  que, para todo $\gamma\geq\beta$,}
  $f_{\alpha\gamma}(A_\gamma)=f_{\alpha\beta}(A_\beta)$.
\end{itemize}

Está claro que la condición (ML) se cumple si todos los
$f_{\alpha\beta}$ son \emph{epimorfismos}, la condición (ML) se satisface. Recíprocamente, si se satisface (ML), y si para todo $\alpha\in I$ sea $A'_\alpha$ el subobjeto de las imágenes
universales en $A_\alpha$, la restricción de $f_{\alpha\beta}$ a $A'_\beta$ es un \emph{epimorfismo}
$A'_\beta\to A'_\alpha$ para $\alpha\leq\beta$: en efecto, si $\gamma\geq\beta$ es tal que $f_{\beta\delta}(A_\delta)=f_{\beta\gamma}(A_\gamma)$ para $\delta\geq\gamma$, se tiene
$A'_\beta=f_{\beta\gamma}(A_\gamma)$ y esto implica, por otra parte, $f_{\alpha\delta}(A_\delta)=f_{\alpha\gamma}(A_\gamma)$ para $\delta\geq\gamma$, por tanto $A'_\alpha=f_{\alpha\gamma}(A_\gamma)=f_{\alpha\beta}(A'_\beta)$.

Observemos también que la condición (ML) se satisface cuando los objetos $A_\alpha$ son
\emph{artinianos} en $\boldsymbol{C}$, es decir, cuando toda familia de
subobjetos de $A_\alpha$ posee un elemento minimal.
En efecto, un elemento minimal de la familia filtrante decreciente
$(f_{\alpha\beta}(A_\beta))$ de subobjetos de $A_\alpha$
es necesariamente el menor de esos subobjetos.
\end{env}

\begin{remark}[13.1.3]
\label{0.13.1.3}
\label{0.13.1.3-fr}
\emph{Observación} (13.1.3).
La condición (ML) puede formularse también cuando $\boldsymbol{C}$ es, por
ejemplo, la categoría de conjuntos.
Se puede definir de nuevo el subconjunto de las imágenes universales en
$A_\alpha$, y las observaciones hechas en \sref{0.13.1.1} y
\sref{0.13.1.2} siguen siendo válidas.
\end{remark}

\subsection{La condición de Mittag-Leffler para grupos abelianos}
\label{subsection:0.13.2}
\label{subsection:0.13.2-fr}

\begin{proposition}[13.2.1]
\label{0.13.2.1}
\label{0.13.2.1-fr}
Sea
\[
0\longrightarrow A_\alpha\xrightarrow{u_\alpha}B_\alpha
\xrightarrow{v_\alpha}C_\alpha\longrightarrow0
\]
una sucesión exacta de sistemas proyectivos de grupos abelianos, relativos
al mismo conjunto de índices filtrante $I$.
\begin{enumerate}
  \item[{\rm(i)}]
  Si $(B_\alpha)$ satisface {\rm(ML)}, entonces también $(C_\alpha)$.
  \item[{\rm(ii)}]
  Si $(A_\alpha)$ y $(C_\alpha)$ satisfacen {\rm(ML)}, entonces también
  $(B_\alpha)$.
\end{enumerate}
\end{proposition}

\begin{proof}
Sean $(f_{\alpha\beta})$, $(g_{\alpha\beta})$ y
$(h_{\alpha\beta})$ los sistemas de homomorfismos que definen
$(A_\alpha)$, $(B_\alpha)$ y $(C_\alpha)$ respectivamente.

(i) Supongamos que $g_{\alpha\beta}(B_\beta)=g_{\alpha\lambda}(B_\lambda)$ para $\lambda\geq\beta$; como $v_\beta$ y $v_\lambda$ son sobreyectivos,
$h_{\alpha\beta}(C_\beta)=v_\alpha(g_{\alpha\beta}(B_\beta))=v_\alpha(g_{\alpha\lambda}(B_\lambda))=h_{\alpha\lambda}(C_\lambda)$
para $\lambda\geq\beta$.

(ii) Sea $\alpha\in I$, y sea $\beta\geq\alpha$ un índice tal que para $\lambda\geq\beta$ se tenga $f_{\alpha\beta}(A_\beta)=f_{\alpha\lambda}(A_\lambda)$; sea, por otra parte, $\gamma\geq\beta$ un índice tal que, para $\lambda\geq\gamma$, se tiene
$h_{\beta\gamma}(C_\gamma)=h_{\beta\lambda}(C_\lambda)$. Sea entonces
$v_\alpha$ un elemento de $g_{\alpha\gamma}(B_\gamma)$; se tiene pues
$y_\alpha=g_{\alpha\gamma}(y_\gamma)$ con $y_\gamma\in B_\gamma$,
y pongamos $y_\beta=g_{\beta\gamma}(y_\gamma)$.
Entonces
$v_\beta(y_\beta)=h_{\beta\gamma}(v_\gamma(y_\gamma))$.
Para todo $\lambda\geq\gamma$, existe, por hipótesis, un
$y_\lambda\in B_\lambda$ tal que
$h_{\beta\gamma}(v_\gamma(y_\gamma))
=h_{\beta\lambda}(v_\lambda(y_\lambda))
=v_\beta(g_{\beta\lambda}(y_\lambda))$,
De aquí que
$v_\beta(y_\beta\setmin g_{\beta\lambda}(y_\lambda))=0$, y por tanto
existe un $x_\beta\in A_\beta$
\oldpage[0\textsubscript{III}]{66}
tal que
$y_\beta=g_{\beta\lambda}(y_\lambda)+u_\beta(x_\beta)$.
Se sigue que
$y_\alpha=g_{\alpha\lambda}(y_\lambda)+u_\alpha(f_{\alpha\beta}(x_\beta))$
Como $\lambda\geq\beta$, existe un $x_\lambda\in A_\lambda$ tal
que
$f_{\alpha\beta}(x_\beta)=f_{\alpha\lambda}(x_\lambda)$,
Por consiguiente,
$y_\alpha=g_{\alpha\lambda}(y_\lambda+u_\lambda(x_\lambda))\in g_{\alpha\lambda}(B_\lambda)$,
lo que demuestra la aserción.
\end{proof}

\begin{proposition}[13.2.2]
\label{0.13.2.2}
\label{0.13.2.2-fr}
Sea $I$ un conjunto ordenado filtrante que posee un subconjunto
cofinal \textup{numerable}, y sea
\[
0\longrightarrow A_\alpha\xrightarrow{u_\alpha}B_\alpha
\xrightarrow{v_\alpha}C_\alpha\longrightarrow0
\]
una sucesión exacta de sistemas proyectivos de grupos abelianos indexados
por $I$.
Si $(A_\alpha)$ satisface {\rm(ML)}, entonces la sucesión
\[
0\longrightarrow\varprojlim A_\alpha
\longrightarrow\varprojlim B_\alpha
\longrightarrow\varprojlim C_\alpha\longrightarrow0
\]
es exacta.
\end{proposition}

\begin{proof}
Solo queda demostrar que el homomorfismo
$v=\varprojlim v_\alpha:\varprojlim B_\alpha\to\varprojlim C_\alpha$ es sobreyectivo. Sea $z=(z_\alpha)$ un elemento de $\varprojlim C_\alpha$, y pongamos
$E_\alpha=v_\alpha^{-1}(z_\alpha)$
Los $E_\alpha$ forman manifiestamente un sistema proyectivo de conjuntos no
vacíos mediante las restricciones de los homomorfismos
$g_{\alpha\beta}:B_\beta\to B_\alpha$.

Demostremos que este sistema proyectivo satisface (ML).
Identificando $A_\alpha$ con un subgrupo de $B_\alpha$ por medio de
$u_\alpha$, para todo $\alpha\in I$, un índice
$\beta\geq\alpha$ tal que $g_{\alpha\beta}(A_\beta)=g_{\alpha\lambda}(A_\lambda)$ para $\lambda\geq\beta$; mostremos que se tiene también $g_{\alpha\beta}(E_\beta)=g_{\alpha\lambda}(E_\lambda)$ para $\lambda\geq\beta$. En efecto, tomemos un $y_\lambda\in E_\lambda$ y pongamos $y_\beta=g_{\beta\lambda}(y_\lambda)$, $y_\alpha=g_{\alpha\lambda}(y_\lambda)$; sea $y'_\alpha\in g_{\alpha\beta}(E_\beta)$, de modo que
$y'_\alpha=g_{\alpha\beta}(y'_\beta)$ para algún
$y'_\beta\in E_\beta$.
Entonces $y'_\beta\setmin y_\beta=x_\beta\in A_\beta$, y por hipótesis existe
un $x_\lambda\in A_\lambda$ tal que
$g_{\alpha\beta}(x_\beta)=g_{\alpha\lambda}(x_\lambda)$
Así,
\[
y'_\alpha=g_{\alpha\beta}(y_\beta)+g_{\alpha\beta}(x_\beta)
=g_{\alpha\lambda}(y_\lambda)+g_{\alpha\lambda}(x_\lambda)
=g_{\alpha\lambda}(y_\lambda+x_\lambda)
\in g_{\alpha\lambda}(E_\lambda),
\] lo que demuestra nuestra aserción. Dicho esto, se sabe (Bourbaki, \emph{Top. gén.}, cap.~II, 3\textsuperscript{a} ed., §~3, teorema~1) que bajo las hipótesis hechas sobre $I$, un sistema proyectivo de conjuntos no vacíos que satisface (ML)
posee un límite proyectivo no vacío.
Existe, por tanto, un punto
$y=(y_\alpha)\in\varprojlim E_\alpha$.
Como $v_\alpha(y_\alpha)=z_\alpha$ para todo $\alpha$, se tiene
$z=v(y)$.
\end{proof}

\begin{proposition}[13.2.3]
\label{0.13.2.3}
\label{0.13.2.3-fr}
Bajo las hipótesis sobre $I$ de \sref{0.13.2.2}, sea
$(K_\alpha^\bullet)_{\alpha\in I}$ un sistema proyectivo
de complejos de grupos abelianos
$K_\alpha^\bullet=(K_\alpha^n)_{n\in\mathbb Z}$
cuyas diferenciales tienen grado $+1$.
Para todo $n$, existe un homomorfismo canónico funtorial
\begin{equation}
\label{0.13.2.3.1-fr}
\tag{13.2.3.1}
h_n:H^n(\varprojlim K_\alpha^\bullet)
\longrightarrow\varprojlim H^n(K_\alpha^\bullet).
\end{equation}
Si, para todo grado $n$, el sistema proyectivo de grupos abelianos
$(K_\alpha^n)_{\alpha\in I}$ satisface {\rm(ML)},
entonces todo $h_n$ son sobreyectivos. Si además, para un grado $n$, el sistema proyectivo $(H^{n-1}(K_\alpha^\bullet))_{\alpha\in I}$
satisface {\rm(ML)}, entonces $h_n$ es biyectivo.
\end{proposition}

\begin{proof}
Pongamos, para todo $n$, $K^n=\varprojlim K_\alpha^n$; la definición de los homomorfismos $h_n$ están definidos por los diagramas conmutativos
\[
\xymatrix{
\ldots\ar[r]&K^{n-1}\ar[r]\ar[d]&K^n\ar[r]\ar[d]
&K^{n+1}\ar[r]\ar[d]&\ldots\\
\ldots\ar[r]&K_\alpha^{n-1}\ar[r]&K_\alpha^n\ar[r]
&K_\alpha^{n+1}\ar[r]&\ldots
}
\]
\oldpage[0\textsubscript{III}]{67}
siendo las diferenciales en $K^\bullet$ los límites proyectivos
de las diferenciales correspondientes en los
$K_\alpha^\bullet$.
Consideremos las sucesiones exactas
\[
(*_n)\qquad
0\longrightarrow \mathrm B^n(K_\alpha^\bullet)
\longrightarrow \mathrm Z^n(K_\alpha^\bullet)
\longrightarrow H^n(K_\alpha^\bullet)\longrightarrow0
\]
y
\[
(**_n)\qquad
0\longrightarrow \mathrm Z^{n-1}(K_\alpha^\bullet)
\longrightarrow K_\alpha^{n-1}
\longrightarrow \mathrm B^n(K_\alpha^\bullet)\longrightarrow0
\]
La hipótesis y la Proposición~\sref{0.13.2.1}~(i) muestran que el
sistema proyectivo
$(\mathrm B^n(K_\alpha^\bullet))_{\alpha\in I}$
satisface (ML) para todo $n$.
Se sigue de \sref{0.13.2.2} que
\[
(***_n)\qquad
0\longrightarrow\varprojlim_\alpha\mathrm B^n(K_\alpha^\bullet)
\longrightarrow\varprojlim_\alpha\mathrm Z^n(K_\alpha^\bullet)
\longrightarrow\varprojlim_\alpha H^n(K_\alpha^\bullet)
\longrightarrow0
\]
es exacta.
Ahora bien,
$\varprojlim\mathrm B^n(K_\alpha^\bullet)$
se identifica con un subgrupo de $K^{n+1}$ que contiene
$\mathrm B^n(K^\bullet)$, mientras que
$\varprojlim\mathrm Z^n(K_\alpha^\bullet)$
se identifica con un subgrupo de $\mathrm Z^n(K^\bullet)$.
Por consiguiente, $h_n$ es sobreyectiva. Si ahora se supone además que el sistema proyectivo $(H^{n-1}(K_\alpha^\bullet))_{\alpha\in I}$
satisface (ML).
Las sucesiones exactas $(*_{n-1})$ y la
Proposición~\sref{0.13.2.1}~(ii) muestran que
$(\mathrm Z^{n-1}(K_\alpha^\bullet))_{\alpha\in I}$
satisface (ML).
Aplicando \sref{0.13.2.2} a las sucesiones exactas $(**_n)$ se obtiene una
sucesión exacta
\[
0\longrightarrow
\varprojlim_\alpha\mathrm Z^{n-1}(K_\alpha^\bullet)
\longrightarrow K^{n-1}\xrightarrow{u}
\varprojlim_\alpha\mathrm B^n(K_\alpha^\bullet)
\longrightarrow0
\]
Como
$\varprojlim\mathrm B^n(K_\alpha^\bullet)
\supset\mathrm B^n(K^\bullet)$, y que el compuesto de la inyección $\varprojlim\mathrm B^n(K_\alpha^\bullet)\to K^n$ y de $u$ es el operador de derivación $K^{n-1}\to K^n$, la sobreyectividad de
$u$ implica
$\varprojlim\mathrm B^n(K_\alpha^\bullet)
=\mathrm B^n(K^\bullet)$,
Así, $h_n$ es inyectivo.
\end{proof}

\begin{env}[13.2.4]
\label{0.13.2.4}
\label{0.13.2.4-fr}
\emph{Observaciones} (13.2.4).
\begin{enumerate}
  \item[{\rm(i)}]
  El argumento de \sref{0.13.2.2}
  (cf.\ Bourbaki, \emph{loc.\ cit.}) muestra que la conclusión de
  esa proposición sigue siendo válida si se supone únicamente que los
  $A_\alpha$ pueden dotarse de estructuras de espacios metrizables completos
  en los que las traslaciones son homeomorfismos, que las aplicaciones
  $f_{\alpha\beta}:A_\beta\to A_\alpha$ que definen el sistema proyectivo $(A_\alpha)$ son \emph{uniformemente continuas} para las distancias consideradas, y finalmente que el sistema $(A_\alpha)$ satisface la condición
  \begin{itemize}
    \item[(ML$'$)]
    \emph{Para todo índice $\alpha$, existe un $\beta\geq\alpha$ tal
    que, para todo $\gamma\geq\beta$,
    $f_{\alpha\gamma}(A_\gamma)$ es denso en
    $f_{\alpha\beta}(A_\beta)$.}
  \end{itemize}

  Esto da un suplemento análogo a \sref{0.13.2.3}.
  Supongamos que $K_\alpha^n=0$ para $n<0$ y todo $\alpha$; supongamos además que $(K_\alpha^n)_{\alpha\in I}$ satisface (ML) para
  $n\geq0$; y que
  $A_\alpha=H^0(K_\alpha^\bullet)$ puede dotarse de
  estructuras de espacio métrico que tengan las propiedades recién enunciadas.
  Entonces las conclusiones de \sref{0.13.2.3} no cambian para
  $n\geq2$, y $h_1$ es además \emph{biyectivo}.
  En efecto, el argumento de \sref{0.13.2.2} muestra también que
  $(\mathrm B^1(K_\alpha^\bullet))_{\alpha\in I}$
  satisface (ML), que la sucesión $(***_1)$ es exacta, y finalmente, en virtud de lo que precede, que $\varprojlim\mathrm B^1(K_\alpha^\bullet)=\mathrm B^1(K^\bullet)$.
  Esto demuestra, entre otras cosas, las aserciones de (T, 3.10.2).

  \item[{\rm(ii)}]
  Se pueden introducir los funtores derivados por la izquierda
  $\varinjlim^{(n)}$ del funtor $\varinjlim$, y obtener enunciados más
  completos que los precedentes~\cite{III-28}.
\end{enumerate}
\end{env}

\subsection{Aplicación: cohomología de un límite proyectivo de haces}
\label{subsection:0.13.3}
\label{subsection:0.13.3-fr}

\begin{proposition}[13.3.1]
\label{0.13.3.1}
\label{0.13.3.1-fr}
\oldpage[0\textsubscript{III}]{68}
Sea $X$ un espacio topológico, sea
$(\mathcal F_k)_{k\in\mathbb N}$ un sistema proyectivo de haces de
grupos abelianos sobre $X$, y sea $\mathcal F=\varprojlim\mathcal F_k$.
Se suponen verificadas las condiciones siguientes:
\begin{enumerate}
  \item[{\rm(i)}]
  Existe una base $\mathfrak B$ de la topología de $X$ tal que,
  para todo $U\in\mathfrak B$ y todo $i\geq0$, el sistema
  proyectivo
  $(H^i(U,\mathcal F_k))_{k\in\mathbb N}$ satisface \textup{(ML)}.

  \item[{\rm(ii)}]
  Para todo $x\in X$ y todo $i>0$, se tiene
  $\varinjlim_U\bigl(\varprojlim_k H^i(U,\mathcal F_k)\bigr)=0$
  donde $U$ recorre las vecindades de $x$ pertenecientes a
  $\mathfrak B$.

  \item[{\rm(iii)}]
  Los homomorfismos
  $u_{hk}:\mathcal F_k\to\mathcal F_h$ $(h\leq k)$ que definen el
  sistema proyectivo $(\mathcal F_k)$ son sobreyectivos.
\end{enumerate}
Bajo estas condiciones, para todo $i>0$, el homomorfismo canónico
\[
h_i:H^i(X,\mathcal F)\longrightarrow
\varprojlim_k H^i(X,\mathcal F_k)
\]
es sobreyectivo. Si, además, para un valor dado de $i$ el sistema
proyectivo
$(H^{i-1}(X,\mathcal F_k))_{k\in\mathbb N}$ satisface \textup{(ML)}, entonces $h_i$ es
biyectivo.
\end{proposition}

\begin{proof}
\emph{a)}
Supongamos primero que los $\mathcal F_k$, así como los \emph{núcleos}
$\mathcal N_{hk}$ de los $u_{hk}$, son \emph{flácidos}; mostraremos entonces que
la condición~{\rm(iii)} del enunciado implica
$H^i(X,\mathcal F)=0$ para $i>0$.
Basta demostrar que, para \emph{todo} subconjunto abierto $U$ de $X$ y todo
recubrimiento $\mathfrak U$ de $U$ por subconjuntos abiertos de $U$, se tiene $H^i(\mathfrak U,\mathcal F)=0$ para $i>0$. De ello resultará en efecto, en primer lugar, que para la cohomología de Čech se tiene
$\check H^i(U,\mathcal F)=0$ para todo $i>0$, y por tanto, por
(G, II, 5.9.2) aplicado al conjunto de todos los subconjuntos abiertos de
$X$, que $H^i(X,\mathcal F)=0$ para todo $i>0$.

Como los $\mathcal F_k$ son flácidos, se tiene
$H^i(\mathfrak U,\mathcal F_k)=0$ para todo $i>0$ (G, II, 5.2.3); consideremos para cada $k$ el complejo $C^\bullet(\mathfrak U,\mathcal F_k)$ de cocadenas
alternadas sobre el nervio del recubrimiento $\mathfrak U$ (G, II, 5.1), que forman evidentemente un sistema proyectivo de complejos de grupos abelianos. Mostremos que todas las aplicaciones $C^\bullet(\mathfrak U,\mathcal F_k)\to
C^\bullet(\mathfrak U,\mathcal F_h)$ $(h\leq k)$ son \emph{sobreyectivas}. Evidentemente, basta, por definición, demostrar que para todo abierto $V$
de $X$, la aplicación
$\Gamma(V,\mathcal F_k)\to\Gamma(V,\mathcal F_h)$ es
es sobreyectiva. Pero la sucesión
$0\to\mathcal N_{hk}\to\mathcal F_k\to\mathcal F_h\to0$
es exacta por hipótesis y da, por tanto, la sucesión exacta de
cohomología
\[
\Gamma(V,\mathcal F_k)\longrightarrow\Gamma(V,\mathcal F_h)
\longrightarrow H^1(V,\mathcal N_{hk})=0
\]
pues $\mathcal N_{hk}$ es flácido.

El sistema proyectivo
$(C^\bullet(\mathfrak U,\mathcal F_k))_{k\in\mathbb N}$
satisface, por tanto, (ML), y lo mismo vale para
$(H^i(\mathfrak U,\mathcal F_k))_{k\in\mathbb N}$ para todo $i\geq0$:
para $i>0$, y como $H^0(\mathfrak U,\mathcal F_k)=\Gamma(U,\mathcal F_k)$
(G, II, 5.2.2), la condición (ML) se cumple también para $i=0$ por lo que precede. Se puede, pues, aplicar (\hyperref[0.13.2.3-fr]{13.2.3}), que muestra que $H^i(\mathfrak U,\mathcal F)=\varprojlim_k H^i(\mathfrak U,\mathcal F_k)=0$ para todo $i>0$.

\emph{b)} Pasemos al caso general, y consideremos para cada $k\in\mathbb N$, la \emph{resolución canónica} $\mathcal C^\bullet(X,\mathcal F_k)
=(\mathcal C^i(X,\mathcal F_k))_{i\geq0}$
de $\mathcal F_k$ por haces flácidos (G, II, 4.3).
Para todo $i\geq0$, está claro que
$(\mathcal C^i(X,\mathcal F_k))_{k\in\mathbb N}$ es un sistema proyectivo de
haces flácidos; demostremos que satisface las condiciones de la
parte~\emph{a)}.
\oldpage[0\textsubscript{III}]{69}
En efecto, si $\mathcal N_{hk}$ denota el núcleo de $u_{hk}$ para $h\leq k$,
entonces la sucesión
$0\to\mathcal N_{hk}\to\mathcal F_k\to\mathcal F_h\to0$
es exacta por~{\rm(iii)}, y nuestra aserción se sigue de la exactitud
del funtor
$\mathcal A\leadsto\mathcal C^i(X,\mathcal A)$ es exacto (G, II, 4.3). Sea $\mathcal G^i=\varprojlim_k\mathcal C^i(X,\mathcal F_k)$
Por la parte~\emph{a)}, se tiene
$H^j(X,\mathcal G^i)=0$ para $j>0$ e $i\geq0$.
Demostraremos que
$\mathcal G^\bullet=(\mathcal G^i)_{i\geq0}$ es una \emph{resolución} del haz
$\mathcal F$. Como $H^j(X,\mathcal G^\bullet)=0$ para $j>0$, la cohomología
$H^\bullet(X,\mathcal F)$ será entonces igual a
$H^\bullet(\Gamma(X,\mathcal G^\bullet))$ (G, II, 4.7.1). Está claro que, por paso al límite proyectivo, se deducen de las sucesiones exactas \[
0\longrightarrow\mathcal F_k
\longrightarrow\mathcal C^0(X,\mathcal F_k)
\longrightarrow\mathcal C^1(X,\mathcal F_k)\longrightarrow\ldots
\]
se obtiene manifiestamente un complejo de haces de grupos abelianos
\[
0\longrightarrow\mathcal F\longrightarrow\mathcal G^0
\longrightarrow\mathcal G^1\longrightarrow\ldots
\]
Para demostrar nuestra aserción, queda establecer que
$\mathcal H^i(\mathcal G^\bullet)=0$ para $i>0$.
Este haz está generado por el prehaz
$U\leadsto H^i(\Gamma(U,\mathcal G^\bullet))$
(G, II, 4.1).
Ahora bien, el complejo $\Gamma(U,\mathcal G^\bullet)$ es el límite proyectivo
del sistema proyectivo de complejos de grupos abelianos
$(\Gamma(U,\mathcal C^\bullet(X,\mathcal F_k)))_{k\in\mathbb N}$
(\textbf{0}\textsubscript{I}, \EGAUnavailableHyperref{0.3.2.6}{3.2.6}).
Se ha visto en la parte~\emph{a)} que, para todo $i\geq0$, las aplicaciones $\Gamma(U,\mathcal C^i(X,\mathcal F_k))\to
\Gamma(U,\mathcal C^i(X,\mathcal F_h))$ $(h\leq k)$ son \emph{sobreyectivas}; por otra parte, se tiene
$H^i(U,\mathcal F_k)=
H^i(\Gamma(U,\mathcal C^\bullet(X,\mathcal F_k)))$,
pues la resolución canónica $\mathcal C^\bullet(U,\mathcal F_k|U)$ es
inducida sobre $U$ por $\mathcal C^\bullet(X,\mathcal F_k)$.
Por la hipótesis~{\rm(i)}, para todo $U\in\mathfrak B$, se puede aplicar (\hyperref[0.13.2.3-fr]{13.2.3}) al sistema proyectivo de complejos $(\Gamma(U,\mathcal C^\bullet(X,\mathcal F_k)))_{k\in\mathbb N}$, y se tiene, pues, $H^i(\Gamma(U,\mathcal G^\bullet))=\varprojlim_k H^i(U,\mathcal F_k)$ para todo $i\geq0$. La hipótesis (ii) demuestra entonces, por definición, que los haces $\mathcal H^i(\mathcal G^\bullet)$ se anulan para $i>0$.

Así, para todo $i\geq0$, $H^i(X,\mathcal F)=H^i(\Gamma(X,\mathcal G^\bullet))$ y \[
\Gamma(X,\mathcal G^\bullet)
=\varprojlim_k\Gamma(X,\mathcal C^\bullet(X,\mathcal F_k)).
\] Se acaba de observar que las aplicaciones $\Gamma(X,\mathcal C^i(X,\mathcal F_k))\to
\Gamma(X,\mathcal C^i(X,\mathcal F_h))$ $(h\leq k)$ son todas sobreyectivas; la conclusión resulta, pues, de nuevo de (\hyperref[0.13.2.3-fr]{13.2.3}).
\end{proof}

\begin{env}[13.3.2]
\label{0.13.3.2}
\label{0.13.3.2-fr}
\emph{Observaciones} (13.3.2).
\begin{enumerate}
  \item[{\rm(i)}]
  El enunciado de \sref{0.13.3.1} solo tiene interés para $i>0$,
  pues, para $i=0$, $h_i$ es siempre un isomorfismo sin hipótesis
  alguna (\textbf{0}\textsubscript{I}, \EGAUnavailableHyperref{0.3.2.6}{3.2.6}).

  \item[{\rm(ii)}]
  Las condiciones~{\rm(i)} y~{\rm(ii)} de \sref{0.13.3.1} se satisfacen,
  en particular, si
  $H^i(U,\mathcal F_k)=0$ para todo $k$, todo $i>0$ y todo
  $U\in\mathfrak B$, y si, para $U\in\mathfrak B$, las aplicaciones
  $\Gamma(U,\mathcal F_k)\to\Gamma(U,\mathcal F_h)$
  son \emph{sobreyectivas}. Este es el caso más frecuente en el que se aplica
  \sref{0.13.3.1}.
\end{enumerate}
\end{env}

\subsection{Condición de Mittag-Leffler y objetos graduados asociados a sistemas proyectivos}
\label{subsection:0.13.4}
\label{subsection:0.13.4-fr}
\label{0.13.4}
\label{0.13.4-fr}

\begin{env}[13.4.1]
\label{0.13.4.1}
\label{0.13.4.1-fr}
Sea
$\mathbf A=(A_k,u_{kh})_{k\in\mathbb Z}$ un sistema proyectivo en una
categoría abeliana $\boldsymbol{C}$.
Decimos que está \emph{acotado inferiormente} si existe un $k_0$ tal que
$A_k=0$ para $k<k_0$. Definiremos sobre cada $A_k$ una \emph{filtración} $(\mathrm F^p(A_k))_{p\in\mathbb Z}$ por
\begin{equation}
\label{0.13.4.1.1-fr}
\tag{13.4.1.1}
\begin{cases}
\mathrm F^p(A_k)=\operatorname{Ker}(A_k\longrightarrow A_{p-1})&\text{para }p\leq k+1\\
\mathrm F^p(A_k)=0&\text{para }p\geq k+1
\end{cases}
\end{equation}

\oldpage[0\textsubscript{III}]{70}
Por hipótesis,
$\mathrm F^{k_0}(A_k)=A_k$ y
$\mathrm F^{k+1}(A_k)=0$; en otras palabras, la filtración considerada
es finita (\sref{0.11.1.3}).
Los objetos graduados asociados a esta filtración son, por tanto,
\[
\operatorname{gr}^p(A_k)=
\operatorname{Ker}(A_k\longrightarrow A_{p-1})
/\operatorname{Ker}(A_k\longrightarrow A_p)
\]
Por consiguiente, $\operatorname{gr}^p(A_k)$ es isomorfo a la imagen por
$A_k\to A_p$ de
$\operatorname{Ker}(A_k\to A_{p-1})$.
Por la transitividad de los morfismos que definen un sistema proyectivo,
se tiene, por tanto,
\begin{equation}
\label{0.13.4.1.2-fr}
\tag{13.4.1.2}
\operatorname{gr}^p(A_k)=
\operatorname{Ker}(A_p\longrightarrow A_{p-1})
\cap\operatorname{Im}(A_k\longrightarrow A_p)
\end{equation}
Pero \sref{0.13.4.1.1} da
$\operatorname{Ker}(A_p\to A_{p-1})=\operatorname{gr}^p(A_p)$, y por tanto
\begin{equation}
\label{0.13.4.1.3-fr}
\tag{13.4.1.3}
\operatorname{gr}^p(A_k)=
\operatorname{gr}^p(A_p)\cap\operatorname{Im}(A_k\longrightarrow A_p).
\end{equation}
Las definiciones precedentes muestran además que se tiene para $k\leq h$,
\[
u_{kh}(\mathrm F^p(A_h))\subset\mathrm F^p(A_k)
\]
y por consiguiente que los $\operatorname{gr}^p(u_{kh})$ definen un sistema proyectivo
$(\operatorname{gr}^p(A_k))_{k\in\mathbb Z}$ para todo $p\in\mathbb Z$.
\end{env}

\begin{env}[13.4.2]
\label{0.13.4.2}
\label{0.13.4.2-fr}
Decimos que el sistema proyectivo $\mathbf A$ es
\emph{esencialmente constante} si los morfismos
$A_{k+1}\to A_k$ son isomorfismos para todo $k$
suficientemente grande.
Decimos que $\mathbf A$ es \emph{estricto} si los morfismos
$A_i\to A_j$ $(j\leq i)$ son epimorfismos.

Cuando $\mathbf A$ es estricto, se sigue de \sref{0.13.4.1.3} que, para
$p\leq k\leq h$, el morfismo canónico
$\operatorname{gr}^p(A_h)\to\operatorname{gr}^p(A_k)$
es un isomorfismo; en otras palabras, el sistema proyectivo
$(\operatorname{gr}^p(A_k))_{k\in\mathbb Z}$ es esencialmente constante.
La sucesión de objetos $\operatorname{gr}^p(A_p)$ (identificados con $\varprojlim_k\operatorname{gr}^p(A_k)$ para todo $p$) se denota entonces
$\operatorname{gr}^\bullet(\mathbf A)$ y se llama el
\emph{objeto graduado asociado al sistema proyectivo estricto}
$\mathbf A=(A_k)$. Si se supone ahora que el sistema proyectivo $\mathbf A$, acotado
inferiormente, satisface (ML).
Por \sref{0.13.1.2}, el sistema proyectivo
$\mathbf A'=(A'_k)$ de objetos de ``imagen universal'' es estricto; también está
acotado inferiormente.
El objeto graduado $\operatorname{gr}^\bullet(\mathbf A')$ asociado a
$\mathbf A'$ se llama asimismo el objeto graduado
\emph{asociado a} $\mathbf A$ y se denota
$\operatorname{gr}^\bullet(\mathbf A)$.
\end{env}

\begin{proposition}[13.4.3]
\label{0.13.4.3}
\label{0.13.4.3-fr}
Sea $\mathbf A=(A_k)_{k\in\mathbb Z}$ un sistema proyectivo acotado
inferiormente en una categoría abeliana $\boldsymbol{C}$.
Las condiciones siguientes son equivalentes:
\begin{enumerate}
  \item[{\rm a)}] $\mathbf A$ satisface \textup{(ML)}.
  \item[{\rm b)}] Para todo $p\in\mathbb Z$, el sistema proyectivo
  $(\operatorname{gr}^p(A_k))_{k\in\mathbb Z}$ es esencialmente constante.
\end{enumerate}
Además, cuando se cumplen estas condiciones, existe, para todo
$p\in\mathbb Z$, un isomorfismo canónico
\begin{equation}
\label{0.13.4.3.1-fr}
\tag{13.4.3.1}
\operatorname{gr}^p(\mathbf A)\xrightarrow{\sim}
\varprojlim_k\operatorname{gr}^p(A_k).
\end{equation}
\end{proposition}

\begin{proof}
Se sigue inmediatamente de \sref{0.13.4.1.2} que \emph{a)} implica
\emph{b)}. La misma fórmula, aplicada al sistema proyectivo
$\mathbf A'$ con la notación de \sref{0.13.4.2}, da
\sref{0.13.4.3.1} por definición.

Para $k\leq h$, pongamos
$A_{kh}=\operatorname{Im}(A_h\to A_k)$
Si $k\leq h\leq j$, entonces
$A_{kj}\subset A_{kh}\subset A_k$.
Dotemos a $A_{kh}$ de la filtración inducida por
$(\mathrm F^p(A_k))$.
La transitividad de los morfismos que definen $\mathbf A$ muestra
inmediatamente que esta es también la filtración cociente de
$(\mathrm F^p(A_h))$; por consiguiente,
\begin{equation}
\label{0.13.4.3.2-fr}
\tag{13.4.3.2}
\operatorname{gr}^p(A_{kh})=
\operatorname{Im}\bigl(\operatorname{gr}^p(A_h)\longrightarrow
\operatorname{gr}^p(A_k)\bigr).
\end{equation}

\oldpage[0\textsubscript{III}]{71}
Supongamos ahora que se cumple \emph{b)}.
Para todo $p\in\mathbb Z$ y todo $k\geq p$, existe un entero
$\mathrm L(p,k)$ tal que el miembro derecho de
\sref{0.13.4.3.2} es constante para $h\geq\mathrm L(p,k)$.
Como $\operatorname{gr}^p(A_k)=0$ para $p<k_0$, para $k$ fijo solo
un número finito de los enteros $\mathrm L(p,k)$ puede ser no nulo cuando
$p$ recorre los enteros $\leq k$.
Sea $\mathrm L(k)=m$ el mayor de estos enteros.
Para todo $h\geq m$, se tiene
$A_{kh}\subset A_{km}$, y, por la definición de $m$, la inyección
canónica
$A_{kh}\to A_{km}$ induce un isomorfismo
$\operatorname{gr}^\bullet(A_{kh})\xrightarrow{\sim}
\operatorname{gr}^\bullet(A_{km})$
Como las filtraciones son finitas, esta inyección es ella misma biyectiva
(Bourbaki, \emph{Alg.\ comm.}, cap.~III, \textsection2,
n\textsuperscript{o}~8, teorema~1).
Esto demuestra que $\mathbf A$ satisface (ML).
\end{proof}

\begin{env}[13.4.4]
\label{0.13.4.4-fr}
Supongamos que en $\boldsymbol{C}$ el límite proyectivo $A=\varprojlim A_k$ existe. En las definiciones de (\hyperref[0.13.4.1-fr]{13.4.1}), se puede entonces reemplazar $A_k$
por $A$, y la filtración así definida sobre $A$ sigue satisfaciendo
\begin{equation}
\label{0.13.4.4.1-fr}
\tag{13.4.4.1}
\operatorname{gr}^p(A)=\operatorname{gr}^p(A_p)
\cap\operatorname{Im}(A\longrightarrow A_p).
\end{equation}
\end{env}

\begin{corollary}[13.4.5]
\label{0.13.4.5}
\label{0.13.4.5-fr}
Supongamos que $\boldsymbol{C}$ es la categoría de grupos abelianos.
Si el sistema proyectivo $\mathbf A$ satisface \textup{(ML)} y
$A=\varprojlim_k A_k$, entonces para todo $p\in\mathbb Z$ existe un
isomorfismo canónico
\begin{equation}
\label{0.13.4.5.1-fr}
\tag{13.4.5.1}
\operatorname{gr}^p(A)\xrightarrow{\sim}
\varprojlim_k\operatorname{gr}^p(A_k).
\end{equation}
\end{corollary}

\begin{proof}
En efecto,
$\operatorname{Im}(A_k\to A_p)=\operatorname{Im}(A\to A_p)$
para todo $k$ suficientemente grande
(Bourbaki, \emph{Top.\ g\'{e}n.}, cap.~II,
3\textsuperscript{e}~\'{e}d., \textsection3,
n\textsuperscript{o}~5, teorema~1), y la conclusión se sigue de
\sref{0.13.4.1.3} y \sref{0.13.4.4.1}.
\end{proof}

\subsection{Límites proyectivos de sucesiones espectrales de complejos filtrados}
\label{subsection:0.13.5}
\label{subsection:0.13.5-fr}

\begin{env}[13.5.1]
\label{0.13.5.1}
\label{0.13.5.1-fr}
Sea $\boldsymbol{C}$ una categoría abeliana y sea $X^\bullet$ un complejo
de objetos de $\boldsymbol{C}$ dotado de una filtración
$(\mathrm F^p(X^\bullet))_{p\in\mathbb Z}$ tal que
$\mathrm F^{p_0}(X^\bullet)=X^\bullet$ para algún índice $p_0$. Consideremos para cada $k\in\mathbb Z$ el complejo $X_k^\bullet=X^\bullet/\mathrm F^{k+1}(X^\bullet)$; está provisto canónicamente de la filtración formada por los $\mathrm F^p(X_k^\bullet)=
\mathrm F^p(X^\bullet)/\mathrm F^{k+1}(X^\bullet)$ para $p\leq k$ y por los $\mathrm F^p(X_k^\bullet)=0$ para $p\geq k+1$.
Además, existen morfismos canónicos
$X_{k+1}^\bullet\to X_k^\bullet$, que hacen de
$\mathbf X^\bullet=(X_k^\bullet)_{k\in\mathbb Z}$ un \emph{sistema proyectivo
de complejos filtrados de objetos de $\boldsymbol{C}$}.
Se observará que este sistema proyectivo es \emph{estricto} y que
$X_k^\bullet=0$ para $k<p_0$.
\end{env}

\begin{env}[13.5.2]
\label{0.13.5.2}
\label{0.13.5.2-fr}
Más generalmente, consideremos un sistema proyectivo estricto acotado
inferiormente
$\mathbf X^\bullet=(X_k^\bullet)_{k\in\mathbb Z}$ de complejos de objetos de
$\boldsymbol{C}$, y dotemos cada $X_k^\bullet$ de la filtración definida en
\sref{0.13.4.1} (aplicada en la categoría abeliana de complejos acotados
inferiormente de objetos de $\boldsymbol{C}$).
Los morfismos
$X_k^\bullet\to X_p^\bullet$ ($(p\leq k)$) se convierten entonces
en morfismos de complejos \emph{filtrados} con filtraciones finitas.
La funtorialidad de las sucesiones espectrales de complejos filtrados
(\sref{0.11.2.3}) muestra que los morfismos de transición de
$\mathbf X^\bullet$ inducen morfismos que hacen de
$\mathrm E(\mathbf X^\bullet)=(\mathrm E(X_k^\bullet))$
un \emph{sistema proyectivo de sucesiones espectrales}.
\end{env}

\begin{lemma}[13.5.3]
\label{0.13.5.3}
\label{0.13.5.3-fr}
Supongamos que el sistema proyectivo
$\mathbf X^\bullet=(X_k^\bullet)_{k\in\mathbb Z}$ de complejos filtrados se
obtiene como en \sref{0.13.5.2}. Entonces:
\begin{enumerate}
  \item[{\rm a)}] Para $r\geq p-p_0$, se tiene
  $\mathrm B_r^{pq}(X_k^\bullet)=\mathrm B_\infty^{pq}(X_k^\bullet)$
  para todo $k\in\mathbb Z$.
  \item[{\rm b)}] Para $k+1\leq p+r$, se tiene
  $\mathrm Z_r^{pq}(X_k^\bullet)=\mathrm Z_\infty^{pq}(X_k^\bullet)$.
  \item[{\rm c)}] Para $k+1\geq p+r$, los morfismos $\mathrm Z_r^{pq}(X_h^\bullet)\to\mathrm Z_r^{pq}(X_k^\bullet)$ y $\mathrm B_r^{pq}(X_h^\bullet)\to\mathrm B_r^{pq}(X_k^\bullet)$ son isomorfismos para todo $h\geq k$.
\end{enumerate}
\end{lemma}

\oldpage[0\textsubscript{III}]{72}
Estas tres propiedades se siguen inmediatamente de las definiciones de
\sref{0.11.2.2}, teniendo en cuenta que
$\mathrm F^{p-r+1}(X_k^\bullet)=X_k^\bullet$ cuando $p-r<p_0$.

\begin{env}[13.5.4]
\label{0.13.5.4-fr}
Supongamos satisfechas las hipótesis de (\hyperref[0.13.5.3-fr]{13.5.3}). Entonces, para $p,q,r$ fijos ($r$ finito), los sistemas proyectivos $(\mathrm Z_r^{pq}(X_k^\bullet))_{k\in\mathbb Z}$, $(\mathrm B_r^{pq}(X_k^\bullet))_{k\in\mathbb Z}$, $(\mathrm E_r^{pq}(X_k^\bullet))_{k\in\mathbb Z}$ son \emph{esencialmente constantes}; se denotará por $\mathrm Z_r^{pq}(\mathbf X^\bullet)$,
$\mathrm B_r^{pq}(\mathbf X^\bullet)$, y
$\mathrm E_r^{pq}(\mathbf X^\bullet)=
\mathrm Z_r^{pq}(\mathbf X^\bullet)/
\mathrm B_r^{pq}(\mathbf X^\bullet)$
Los objetos $\mathrm Z_r^{pq}(\mathbf X^\bullet)$ y
$\mathrm B_r^{pq}(\mathbf X^\bullet)$ se identifican canónicamente con
subobjetos de $\mathrm E_1^{pq}(\mathbf X^\bullet)$.
La definición de los $d_r^{pq}$ (M, XV, 1) muestra que estos
morfismos, tomados respecto de los $X_k^\bullet$, son también esencialmente
constantes, y definen, por tanto, morfismos
\begin{equation}
\label{0.13.5.4.1-fr}
\tag{13.5.4.1}
d_r^{pq}:\mathrm E_r^{pq}(\mathbf X^\bullet)
\longrightarrow\mathrm E_r^{p+r,q-r+1}(\mathbf X^\bullet)
\end{equation}
tal que
$d_r^{p+r,q-r+1}\circ d_r^{pq}=0$.
Además, existen isomorfismos canónicos de
$\operatorname{Ker}(d_r^{pq})$ sobre
$\mathrm Z_{r+1}^{pq}(\mathbf X^\bullet)/
\mathrm B_r^{pq}(\mathbf X^\bullet)$
y de $\operatorname{Im}(d_r^{pq})$ sobre
$\mathrm B_{r+1}^{p+r,q-r+1}(\mathbf X^\bullet)/
\mathrm B_r^{p+r,q-r+1}(\mathbf X^\bullet)$.
\end{env}

\begin{lemma}[13.5.5]
\label{0.13.5.5}
\label{0.13.5.5-fr}
Bajo las hipótesis de \sref{0.13.5.3}, para
$s\geq r>p-p_0$ existe un monomorfismo canónico
\begin{equation}
\label{0.13.5.5.1-fr}
\tag{13.5.5.1}
i:\mathrm E_s^{pq}(\mathbf X^\bullet)
\longrightarrow\mathrm E_r^{pq}(\mathbf X^\bullet)
\end{equation}
y un isomorfismo canónico
\begin{equation}
\label{0.13.5.5.2-fr}
\tag{13.5.5.2}
j_r:\mathrm E_r^{pq}(\mathbf X^\bullet)
\xrightarrow{\sim}\mathrm E_\infty^{pq}(X_{p+r-1}^\bullet)
\end{equation}
tal que el diagrama
\begin{equation}
\label{0.13.5.5.3-fr}
\tag{13.5.5.3}
\xymatrix{
\mathrm E_s^{pq}(\mathbf X^\bullet)\ar[r]^{js}\ar[d]_i
&\mathrm E_\infty^{pq}(X_{p+s-1}^\bullet)\ar[d]\\
\mathrm E_r^{pq}(\mathbf X^\bullet)\ar[r]_{jr}
&\mathrm E_\infty^{pq}(X_{p+r-1}^\bullet)
}
\end{equation}
es conmutativo, donde la flecha vertical derecha está inducida por el
morfismo
$X_{p+s-1}^\bullet\to X_{p+r-1}^\bullet$.
\end{lemma}

\begin{proof}
La existencia de $i$ proviene de que $\mathrm B_r^{pq}(X_k^\bullet)=\mathrm B_\infty^{pq}(X_k^\bullet)$ para $r>p-p_0$ (\hyperref[0.13.5.3-fr]{13.5.3}, \emph{a)}); se tiene
$\mathrm Z_r^{pq}(X_k^\bullet)=\mathrm Z_\infty^{pq}(X_k^\bullet)$ para $k+1\leq p+r$ (\hyperref[0.13.5.3-fr]{13.5.3}, \emph{b)}), de donde, en particular,
$\mathrm Z_r^{pq}(X_{p+r-1}^\bullet)=
\mathrm Z_\infty^{pq}(X_{p+r-1}^\bullet)$
Por otra parte,
$\mathrm Z_r^{pq}(X_{p+r-1}^\bullet)$ y
$\mathrm B_r^{pq}(X_{p+r-1}^\bullet)$ se identifican canónicamente con
$\mathrm Z_r^{pq}(\mathbf X^\bullet)$ y
$\mathrm B_r^{pq}(\mathbf X^\bullet)$ por
\sref{0.13.5.3},~\emph{c)}.
Esto demuestra la existencia de $j_r$ y la conmutatividad de
\sref{0.13.5.5.1}.
\end{proof}

\begin{corollary}[13.5.6]
\label{0.13.5.6-fr}
Bajo las hipótesis de \textup{(\hyperref[0.13.5.3-fr]{13.5.3})}, si uno de los límites proyectivos $\varprojlim_r\mathrm E_r^{pq}(\mathbf X^\bullet)$, $\varprojlim_k\mathrm E_\infty^{pq}(X_k^\bullet)$ existe, lo mismo ocurre con el otro, y se tiene un isomorfismo canónico
\begin{equation}
\label{0.13.5.6.1-fr}
\tag{13.5.6.1}
j_\infty:\varprojlim_r\mathrm E_r^{pq}(\mathbf X^\bullet)
\xrightarrow{\sim}
\varprojlim_k\mathrm E_\infty^{pq}(X_k^\bullet).
\end{equation}
Además, el sistema proyectivo
$(\mathrm E_r^{pq}(\mathbf X^\bullet))_{r\in\mathbb Z}$ es esencialmente
constante \textup{(\sref{0.13.4.2})} si y solo si el sistema proyectivo
$(\mathrm E_\infty^{pq}(X_k^\bullet))_{k\in\mathbb Z}$ es esencialmente
constante.
\end{corollary}

\begin{env}[13.5.7]
\label{0.13.5.7}
\label{0.13.5.7-fr}
Se denota por
$\mathrm B_\infty^{pq}(\mathbf X^\bullet)$ y
$\mathrm Z_\infty^{pq}(\mathbf X^\bullet)$ los subobjetos de
$\mathrm E_1^{pq}(\mathbf X^\bullet)$ iguales respectivamente a
$\mathrm B_r^{pq}(\mathbf X^\bullet)$ para $r>p-p_0$ y a
$\inf_r\mathrm Z_r^{pq}(\mathbf X^\bullet)$, cuando este último existe.
Así,
$\varprojlim_r\mathrm E_r^{pq}(\mathbf X^\bullet)$
\oldpage[0\textsubscript{III}]{73}
se identifica canónicamente con
$\mathrm E_\infty^{pq}(\mathbf X^\bullet)=
\mathrm Z_\infty^{pq}(\mathbf X^\bullet)/
\mathrm B_\infty^{pq}(\mathbf X^\bullet)$.
Se observará que los objetos
$\mathrm Z_r^{pq}(\mathbf X^\bullet)$,
$\mathrm B_r^{pq}(\mathbf X^\bullet)$,
$\mathrm E_r^{pq}(\mathbf X^\bullet)$
($(1\leq r\leq+\infty)$), y los morfismos $d_r^{pq}$ dependen
\emph{funtorialmente} del sistema proyectivo $\mathbf X^\bullet$, con las
restricciones de \sref{0.13.5.5}, y los morfismos definidos en
\sref{0.13.5.5} y \sref{0.13.5.6} son funtoriales.
\end{env}

\subsection{Sucesión espectral de un funtor relativa a un objeto dotado de una filtración finita}
\label{subsection:0.13.6}
\label{subsection:0.13.6-fr}

\begin{env}[13.6.1]
\label{0.13.6.1}
\label{0.13.6.1-fr}
Sean $\boldsymbol{C}$ y $\boldsymbol{C}'$ dos categorías abelianas, y sea
$\mathrm T:\boldsymbol{C}\to \boldsymbol{C}'$ un funtor covariante
aditivo.
Supongamos que todo objeto de $\boldsymbol{C}$ es isomorfo a un subobjeto de
un objeto inyectivo, de modo que existan los funtores derivados por la derecha
$\mathrm R^p\mathrm T$ ($(p\geq0)$).
\end{env}

\begin{lemma}[13.6.2]
\label{0.13.6.2}
\label{0.13.6.2-fr}
Sea $A$ un objeto de $\boldsymbol{C}$ dotado de una filtración \textup{finita}
$(\mathrm F^i(A))_{i\in\mathbb Z}$.
Existe una resolución inyectiva
$X^\bullet=(X^j)_{j\geq0}$ de $A$, dotada de una filtración finita
$(\mathrm F^i(X^\bullet))_{i\in\mathbb Z}$ tal que la relación $\mathrm F^i(A)=A$ \textup{(resp. $\mathrm F^i(A)=0$)} implique $\mathrm F^i(X^\bullet)=X^\bullet$ \textup{(resp. $\mathrm F^i(X^\bullet)=0$)} y que, para todo $i\in\mathbb Z$,
$\mathrm F^i(X^\bullet)$ es una resolución inyectiva de
$\mathrm F^i(A)$.
\end{lemma}

\begin{proof}
Sea $p$ (resp. $q>p$) el mayor índice tal que
$\mathrm F^i(A)=A$ (resp. el menor índice tal que
$\mathrm F^i(A)=0$).
Se razona por inducción sobre $q-p$, siendo inmediato el lema cuando
$q-p=1$. Habiendo formado una resolución inyectiva ${X''}^\bullet$ de
$A/\mathrm F^{q-1}(A)$ con las propiedades requeridas, se considera la sucesión exacta $0\to\mathrm F^{q-1}(A)\to A\to A/\mathrm F^{q-1}(A)\to0$, se toma una resolución inyectiva ${X'}^\bullet$ de
$\mathrm F^{q-1}(A)$, luego se determina una resolución inyectiva $X^\bullet$ de $A$ de modo que se obtenga una sucesión exacta
$0\to {X'}^\bullet\to X^\bullet\to {X''}^\bullet\to0$
compatible con la sucesión precedente (M, V, 2.2).
Está claro que $X^\bullet$ tiene las propiedades requeridas.
\end{proof}

\begin{corollary}[13.6.3]
\label{0.13.6.3}
\label{0.13.6.3-fr}
Sea $B$ un segundo objeto de $\boldsymbol{C}$, dotado de una filtración
finita $(\mathrm F^i(B))_{i\in\mathbb Z}$, sea $s$ un entero,
y sea $u:A\to B$ un morfismo tal que
$u(\mathrm F^i(A))\subset\mathrm F^{i+s}(B)$
para todo $i\in\mathbb Z$.
Si $Y^\bullet=(Y^j)_{j\geq0}$ es una resolución inyectiva de $B$
dotada de una filtración
$(\mathrm F^i(Y^\bullet))_{i\in\mathbb Z}$ que tenga las propiedades
enunciadas en \sref{0.13.6.2}, entonces existe un morfismo
$v:X^\bullet\to Y^\bullet$, compatible con $u$, tal
que
$v(\mathrm F^i(X^\bullet))\subset\mathrm F^{i+s}(Y^\bullet)$
para todo $i\in\mathbb Z$.
Además, cualesquiera dos de tales morfismos $v$ y $v'$ son homotópicos.
\end{corollary}

Esto se sigue inmediatamente, por inducción sobre $q-p$, de la construcción
precedente y de (M, V, 2.3).

\begin{env}[13.6.4]
\label{0.13.6.4}
\label{0.13.6.4-fr}
Bajo las hipótesis de \sref{0.13.6.2}, consideremos el complejo
$\mathrm T(X^\bullet)$ en $\boldsymbol{C}'$.
Está filtrado por los complejos
$\mathrm T(\mathrm F^i(X^\bullet))$, pues
$\mathrm F^i(X^\bullet)$ es un sumando directo de $X^\bullet$.
Se sigue de \sref{0.13.6.3} que la \emph{sucesión espectral} de
este complejo filtrado depende, salvo isomorfismo, únicamente del
objeto filtrado $A$.
Su término límite es la cohomología $\mathrm R^\bullet\mathrm T(A)$, dotada
de la filtración
\begin{equation}
\label{0.13.6.4.1-fr}
\tag{13.6.4.1}
\begin{split}
\mathrm F^p(\mathrm R^n\mathrm T(A))
&=\operatorname{Im}\bigl(\mathrm R^n\mathrm T(\mathrm F^p(A))
\longrightarrow\mathrm R^n\mathrm T(A)\bigr)\\
&=\operatorname{Ker}\bigl(\mathrm R^n\mathrm T(A)
\longrightarrow\mathrm R^n\mathrm T(A/\mathrm F^p(A))\bigr)
\end{split}
\end{equation}
por \sref{0.11.2.2}, y su término $\mathrm E_1$ es
\begin{equation}
\label{0.13.6.4.2-fr}
\tag{13.6.4.2}
\mathrm E_1^{pq}=\mathrm R^{p+q}\mathrm T(\operatorname{gr}^p(A))
\end{equation}
donde $\operatorname{gr}^p(A)$ designa, como de costumbre, $\mathrm F^p(A)/\mathrm F^{p+1}(A)$. Está claro, por (\hyperref[0.11.2.2-fr]{11.2.2}), que la filtración del término límite es \emph{finita}, y que para $p$ y $q$ dados, las sucesiones de los
\oldpage[0\textsubscript{III}]{74}
$\mathrm B_r^{pq}(A)=\mathrm B_r^{pq}(\mathrm T(X^\bullet))$ y $\mathrm Z_r^{pq}(A)=\mathrm Z_r^{pq}(\mathrm T(X^\bullet))$ son \emph{estacionarias}, por tanto la sucesión espectral precedente es \emph{birregular} (\hyperref[0.11.1.3-fr]{11.1.3}). Denotaremos esta sucesión por $\mathrm E(A)=(\mathrm E_r^{pq}(A))$ y la llamaremos la
\emph{sucesión espectral del funtor $\mathrm T$ relativa al
objeto filtrado $A$}.
\end{env}

\begin{env}[13.6.5]
\label{0.13.6.5}
\label{0.13.6.5-fr}
Supongamos ahora que se cumplen las hipótesis de \sref{0.13.6.3}, y
conservemos su notación.
Como $\mathrm F^i(X^\bullet)$ (resp.
$\mathrm F^i(Y^\bullet)$) es un sumando directo de $X^\bullet$
(resp. $Y^\bullet$), se tiene
$(\mathrm T v)(\mathrm T(\mathrm F^i(X^\bullet)))
\subset\mathrm T(\mathrm F^{i+s}(Y^\bullet))$
para todo $i\in\mathbb Z$.
Las definiciones de \sref{0.11.2.2} muestran entonces que, para
$1\leq r\leq+\infty$, $\mathrm T v$ define morfismos
$\mathrm B_r^{pq}(\mathrm T(X^\bullet))\to
\mathrm B_r^{p+s,q-s}(\mathrm T(Y^\bullet))$
y
$\mathrm Z_r^{pq}(\mathrm T(X^\bullet))\to
\mathrm Z_r^{p+s,q-s}(\mathrm T(Y^\bullet))$,
y, por tanto, un morfismo
\[
w_r:\mathrm E_r^{pq}(A)\to\mathrm E_r^{p+s,q-s}(B)~;
\]
Análogamente, sobre los términos límite se tienen morfismos
$u_n:\mathrm R^n\mathrm T(A)\to\mathrm R^n\mathrm T(B)$
tales que
$u_n(\mathrm F^p(\mathrm R^n\mathrm T(A)))\subset
\mathrm F^{p+s}(\mathrm R^n\mathrm T(B))$.
La definición de los $d_r^{pq}$ (M, XV, 1) muestra también que los
diagramas
\[
\xymatrix{
\mathrm E_r^{pq}(A)\ar[r]^{d_r^{pq}}\ar[d]_{w_r}
&\mathrm E_r^{p+r,q-r+1}(A)\ar[d]^{w_r}\\
\mathrm E_r^{p+s,q-s}(B)\ar[r]_{d_r^{p+s,q-s}}
&\mathrm E_r^{p+r+s,q-r-s+1}(B)
}
\]
son conmutativos.
Se sigue que existe un diagrama conmutativo análogo para los
isomorfismos $\alpha_r^{pq}$, cuya explicitación dejamos al lector.
Finalmente (\emph{loc.\ cit.}), se tienen también diagramas conmutativos
sobre los términos límite:
\[
\xymatrix{
\mathrm E_\infty^{pq}(A)\ar[r]^{\beta^{pq}}\ar[d]_{w_\infty}
&\operatorname{gr}^p(\mathrm R^{p+q}\mathrm T(A))
  \ar[d]^{u_{p+q}}\\
\mathrm E_\infty^{p+s,q-s}(B)\ar[r]_{\beta^{p+s,q-s}}
&\operatorname{gr}^{p+s}(\mathrm R^{p+q}\mathrm T(B))
}
\]
\end{env}

\begin{env}[13.6.6]
\label{0.13.6.6}
\label{0.13.6.6-fr}
Supongamos, en particular, que existe un anillo $S$, dotado de una
\emph{filtración} $(\mathrm F^i(S))_{i\in\mathbb Z}$, y un homomorfismo
de anillos
\begin{equation}
\label{0.13.6.6.1-fr}
\tag{13.6.6.1}
h:S\to\operatorname{Hom}_{\boldsymbol{C}}(A,A)
\end{equation}
\oldpage[0\textsubscript{III}]{75}
tal que, para todo $t\in\mathrm F^j(S)$,
$h_t(\mathrm F^i(A))\subset\mathrm F^{i+j}(A)$ para todo par $i$, $j$. Diremos para abreviar que $A$ está dotado de la estructura de un
\emph{$S$--$\boldsymbol{C}$-módulo filtrado} sobre el anillo filtrado $S$.

Al pasar a los objetos graduados asociados, todo $h_t$, para
$t\in\mathrm F^j(S)$, define un endomorfismo \emph{graduado}
$\overline h_t$ de $\operatorname{gr}^\bullet(A)$, homogéneo de grado $j$; además, este morfismo solo depende de la clase de $t$ en $\operatorname{gr}^j(S)$, y se define así un homomorfismo de \emph{anillos graduados} \[
\overline h:\operatorname{gr}^\bullet(S)\to
\operatorname{Hom}_{\boldsymbol{C}}^g(\operatorname{gr}^\bullet(A),
\operatorname{gr}^\bullet(A))
\] donde el segundo miembro es el anillo de los endomorfismos \emph{graduados} de $\operatorname{gr}^\bullet(A)$.
Decimos que $\operatorname{gr}^\bullet(A)$ está dotado de la estructura
de un \emph{$\operatorname{gr}^\bullet(S)$--$\boldsymbol{C}$-módulo graduado}.

Se sigue de \sref{0.13.6.5} que, para
$1\leq r\leq+\infty$, todo $\overline t\in\operatorname{gr}^j(S)$ define canónicamente en los objetos bigraduados $(\mathrm B_r^{pq}(A))_{(p,q)\in\mathbb Z\times\mathbb Z}$, $(\mathrm Z_r^{pq}(A))_{(p,q)\in\mathbb Z\times\mathbb Z}$ y $\mathrm E_r(A)=(\mathrm E_r^{pq}(A))_{(p,q)\in\mathbb Z\times\mathbb Z}$
endomorfismos bigraduados de grado $(s,-s)$.
Sobre $\mathrm E_r(A)$, para $r$ finito, este endomorfismo conmuta con
el endomorfismo bigraduado definido por los $d_r^{pq}$.
Estos endomorfismos satisfacen las condiciones usuales de asociatividad y
distributividad respecto de la adición en
$\operatorname{gr}^\bullet(S)$ y en los objetos bigraduados en cuestión.
Para abreviar, decimos que estos últimos son
\emph{$\operatorname{gr}^\bullet(S)$--$\boldsymbol{C}'$-módulos bigraduados};
es inmediato que los $\alpha_r^{pq}$ definen un \emph{isomorfismo} para
este tipo de estructura.

Para todo entero $n$, sea $\mathrm B_r^{(n)}(A)$
(resp. $\mathrm Z_r^{(n)}(A)$ y
$\mathrm E_r^{(n)}(A)$) el subobjeto graduado de $\mathrm B_r^{\bullet\bullet}(A)$ (resp. $\mathrm Z_r^{\bullet\bullet}(A)$, $\mathrm E_r^{\bullet\bullet}(A)$) formado por los $\mathrm B_r^{pq}(A)$ (resp. $\mathrm Z_r^{pq}(A)$, $\mathrm E_r^{pq}(A)$) tales que $p+q=n$ ($1\leq r\leq+\infty$).
Es inmediato que estos son
\emph{$\operatorname{gr}^\bullet(S)$--$\boldsymbol{C}'$-módulos graduados}.

Finalmente, todo $\overline t\in\operatorname{gr}^j(S)$ define, para todo
$n$ un endomorfismo graduado de grado $j$ en el objeto graduado $\operatorname{gr}^\bullet(\mathrm R^n\mathrm T(A))$, que adquiere así la
estructura de un
$\operatorname{gr}^\bullet(S)$--\emph{$\boldsymbol{C}'$-módulo graduado}; los $\beta^{pq}$ (para $p+q=n$) definen un \emph{isomorfismo} de
$\mathrm E^{(n)}(A)$ sobre
$\operatorname{gr}^\bullet(\mathrm R^n\mathrm T(A))$ para esta especie
de estructura.

Se observará que cuando $\boldsymbol{C}'$ es la categoría de \emph{grupos abelianos}, las
estructuras de $S$--$\boldsymbol{C}'$-módulo (resp. de
$\operatorname{gr}^\bullet(S)$--$\boldsymbol{C}'$-módulo graduado o bigraduado) son
simplemente las estructuras usuales de $S$-\emph{módulo} (resp. de
$\operatorname{gr}^\bullet(S)$-módulo graduado o bigraduado).
\end{env}

\subsection{Funtores derivados de un límite proyectivo de argumentos}
\label{subsection:0.13.7}
\label{subsection:0.13.7-fr}
\label{0.13.7}

\begin{env}[13.7.1]
\label{0.13.7.1}
\label{0.13.7.1-fr}
Sean $\boldsymbol{C}$ y $\boldsymbol{C}'$ dos categorías abelianas, suponiéndose $\boldsymbol{C}$ tal que todo objeto de $\boldsymbol{C}$ es isomorfo a un subobjeto de un
objeto inyectivo.
Sea
$\mathrm T:\boldsymbol{C}\to \boldsymbol{C}'$
un funtor covariante aditivo.
Consideremos un \emph{sistema proyectivo estricto}
$\mathbf A=(A_k)_{k\in\mathbb Z}$ en $\boldsymbol{C}$, acotado inferiormente; más
precisamente, supongamos que $A_k=0$ para $k<k_0$. Asociamos canónicamente a este sistema una filtración $(\mathrm F^p(A_k))_{k\in\mathbb Z}$ sobre cada $A_k$ mediante las fórmulas (\hyperref[0.13.4.1.1-fr]{13.4.1.1}), y como se trata de un sistema proyectivo \emph{estricto}, los morfismos canónicos
\begin{equation}
\label{0.13.7.1.1-fr}
\tag{13.7.1.1}
\mathrm F^i(A_h)/\mathrm F^j(A_h)\to
\mathrm F^i(A_k)/\mathrm F^j(A_k)\qquad(h\geq k)
\end{equation}
son \emph{isomorfismos} para $i\leq j\leq k+1$ son \emph{isomorfismos}. Recordemos además que se tiene $\mathrm F^{k_0}(A_k)=A_k$ y
$\mathrm F^{k+1}(A_k)=0$ para todo $k$.
\end{env}

\begin{env}[13.7.2]
\label{0.13.7.2}
\label{0.13.7.2-fr}
Para cada $k$ una \emph{resolución inyectiva} $X_k^\bullet=(X_k^j)_{j\geq0}$ de $A_k$ que tenga las
propiedades de \sref{0.13.6.2}.
Los morfismos canónicos $A_{k+1}\to A_k$ nos permiten, por
\sref{0.13.6.3}, definir para cada $k$ un morfismo de complejos
$X_{k+1}^\bullet\to X_k^\bullet$ compatibles
\oldpage[0\textsubscript{III}]{76}
con las filtraciones, y que hacen de
$\mathbf X^\bullet=(X_k^\bullet)_{k\in\mathbb Z}$ un \emph{sistema proyectivo} de complejos. Se puede suponer además que este sistema proyectivo es \emph{estricto}. Para ello, se observa que, en virtud del isomorfismo (\hyperref[0.13.7.1.1-fr]{13.7.1.1}), $A_k$ es isomorfo a $A_{k+1}/\mathrm F^{k+1}(A_{k+1})$; se puede, pues, tomar, en la construcción de $X_{k+1}^\bullet$, la resolución inyectiva de $A_{k+1}/\mathrm F^{k+1}(A_{k+1})$ \emph{igual} a
$X_k^\bullet$.
Se sigue entonces de (M, V, 2.3) que el morfismo de complejos
$X_{k+1}^\bullet\to X_k^\bullet$ puede
construirse de modo que, al pasar a los cocientes, induzca un
\emph{isomorfismo}
$X_{k+1}^\bullet/\mathrm F^{k+1}(X_{k+1}^\bullet)
\xrightarrow{\sim}X_k^\bullet$
que respeta las filtraciones; esta es precisamente la condición de
\sref{0.13.5.1}.
\end{env}

\begin{env}[13.7.3]
\label{0.13.7.3}
\label{0.13.7.3-fr}
Por construcción, el sistema proyectivo
$\mathrm T(\mathbf X^\bullet)$ de complejos
$\mathrm T(X_k^\bullet)$ satisface las hipótesis de
\sref{0.13.5.3}.
Los resultados de \sref{0.13.5.4}, \sref{0.13.5.5} y
\sref{0.13.5.6} se aplican, por tanto, a las sucesiones espectrales
$\mathrm E(\mathrm T(X_k^\bullet))=\mathrm E(A_k)$; escribiremos $\mathrm E_r^{pq}(\mathbf A)$ en lugar de
$\mathrm E_r^{pq}(\mathrm T(\mathbf X^\bullet))$ para $1\leq r\leq+\infty$ (cf. (\hyperref[0.13.5.7-fr]{13.5.7}) para $r=+\infty$) y lo mismo para las notaciones análogas. Se observará en particular que se tiene \begin{equation}
\label{0.13.7.3.1-fr}
\tag{13.7.3.1}
\mathrm E_1^{pq}(\mathbf A)=
\mathrm R^{p+q}\mathrm T(\operatorname{gr}^p(\mathbf A))
\end{equation}
por \sref{0.13.6.4.2} y el hecho de que el sistema proyectivo
$(\operatorname{gr}^p(A_k))$ es esencialmente constante.
Estos resultados, junto con \sref{0.13.4.3}, dan la proposición
siguiente, demostrada por primera vez por Shih Weishu mediante un método
diferente (inédito).
\end{env}

\begin{proposition}[13.7.4]
\label{0.13.7.4}
\label{0.13.7.4-fr}
(Shih).
Sea $n$ un entero.
Las dos condiciones siguientes son equivalentes:
\begin{enumerate}
  \item[{\rm a)}] Para todo par $(p,q)$ tal que $p+q=n$, el
  sistema proyectivo
  $(\mathrm E_r^{pq}(\mathbf A))_{r\geq2}$ es esencialmente constante.
  \item[{\rm b)}] El sistema proyectivo
  $\mathrm R^n\mathrm T(\mathbf A)=
(\mathrm R^n\mathrm T(A_k))_{k\in\mathbb Z}$ satisface \textup{(ML)}.
\end{enumerate}
Además, cuando se cumplen estas condiciones, se tiene un isomorfismo canónico
\begin{equation}
\label{0.13.7.4.1-fr}
\tag{13.7.4.1}
\operatorname{gr}^p(\mathrm R^n\mathrm T(\mathbf A))
\xrightarrow{\sim}\mathrm E_\infty^{p,n-p}(\mathbf A)
\qquad\textit{pour tout $p\in\mathbb Z$.}
\end{equation}
\end{proposition}
\begin{proof}
En efecto, en virtud de (\hyperref[0.13.5.6-fr]{13.5.6}), la condición \emph{a)} equivale a decir que el sistema proyectivo $(\mathrm E_\infty^{pq}(A_k))_{k\in\mathbb Z}$ es esencialmente constante para $p+q=n$, y por otra parte
$\operatorname{gr}^p(\mathrm R^n(\mathrm T(A_k))$ es canónicamente isomorfo a
$\mathrm E_\infty^{p,n-p}(A_k)$.
Se sigue de \sref{0.13.4.3} que {\rm a)} y {\rm b)} son
equivalentes, y \sref{0.13.7.4.1} no es sino
\sref{0.13.5.6.1} en el caso considerado.
\end{proof}

\begin{corollary}[13.7.5]
\label{0.13.7.5}
\label{0.13.7.5-fr}
Sea $\boldsymbol{\mathcal F}=(\mathcal F_k)_{k\in\mathbb N}$ un sistema proyectivo de
haces de grupos abelianos que satisface las condiciones {\rm(i)}, {\rm(ii)}
y {\rm(iii)} de \sref{0.13.3.1}, y sea
$\mathcal F=\varprojlim\mathcal F_k$.
Supongamos que, para el funtor
$\mathcal G\leadsto\Gamma(X,\mathcal G)$, el sistema proyectivo $(\mathrm E_r^{pq}(\boldsymbol{\mathcal F}))_{r\in\mathbb Z}$ sea esencialmente constante para todo par $(p,q)$ tal que $p+q=n$ o $p+q=n+1$. Consideremos sobre $H^{n+1}(X,\mathcal F)$ la filtración definida por $\mathrm F^p(H^{n+1}(X,\mathcal F))=
\operatorname{Ker}(H^{n+1}(X,\mathcal F)\to
H^{n+1}(X,\mathcal F_{p-1}))$.
Entonces existe un isomorfismo canónico
\begin{equation}
\label{0.13.7.5.1-fr}
\tag{13.7.5.1}
\operatorname{gr}^p(H^{n+1}(X,\mathcal F))
\xrightarrow{\sim}\mathrm E_\infty^{p,n-p+1}(\boldsymbol{\mathcal F})
\qquad\textit{pour tout $p\in\mathbb Z$.}
\end{equation}
\end{corollary}

\begin{proof}
Resulta de (\hyperref[0.13.7.4-fr]{13.7.4}) aplicado al caso en que $\boldsymbol{C}$ la categoría de haces de
grupos abelianos sobre $X$, como $\boldsymbol{C}'$ la categoría de grupos abelianos,
y $\mathrm T=\Gamma$, que se tiene un isomorfismo
\oldpage[0\textsubscript{III}]{77}
canónico
$\operatorname{gr}^p(\mathrm R^{n+1}\Gamma(\boldsymbol{\mathcal F}))
\xrightarrow{\sim}\mathrm E_\infty^{p,n-p+1}(\boldsymbol{\mathcal F})$ para todo $p\in\mathbb Z$. Por otra parte, como, en virtud de (\hyperref[0.13.7.4-fr]{13.7.4}), el sistema proyectivo $(H^n(X,\mathcal F_k))_{k\in\mathbb Z}$ satisface {\rm(ML)}; por tanto
\sref{0.13.3.1} da un isomorfismo canónico
\begin{equation}
\label{0.13.7.5.1b-fr}
\tag{13.7.5.1}
H^{n+1}(X,\mathcal F)\xrightarrow{\sim}
\varprojlim_k H^{n+1}(X,\mathcal F_k).
\end{equation}

Como el sistema proyectivo
$\mathrm R^{n+1}\Gamma(\boldsymbol{\mathcal F})$ satisface {\rm(ML)} por
\sref{0.13.7.4}, existe un isomorfismo canónico
$\operatorname{gr}^p(\mathrm R^{n+1}\Gamma(\boldsymbol{\mathcal F}))
\xrightarrow{\sim}
\varprojlim_k\operatorname{gr}^p(H^{n+1}(X,\mathcal F_k))$
por \sref{0.13.4.3}, y un isomorfismo canónico
$\varprojlim_k\operatorname{gr}^p(H^{n+1}(X,\mathcal F_k))
\xrightarrow{\sim}
\operatorname{gr}^p(\varprojlim_k H^{n+1}(X,\mathcal F_k))$
(\hyperref[0.13.4.5-fr]{13.4.5}). Todo se reduce pues a ver que
el isomorfismo (\hyperref[0.13.7.5.1-fr]{13.7.5.1}) es compatible con
las filtraciones de los dos miembros; pero esto resulta inmediatamente de las
definiciones y de la conmutatividad del diagrama
\[
\xymatrix{
H^{n+1}(X,\mathcal F)\ar[rr]^{\sim}\ar[dr]
&&\varprojlim_k H^{n+1}(X,\mathcal F_k)\ar[dl]\\
&H^{n+1}(X,\mathcal F_{p-1})&
}
\]
para todo $p$.
\end{proof}

\begin{env}[13.7.6]
\label{0.13.7.6}
\label{0.13.7.6-fr}
Sea $S$ un anillo dotado de una filtración
$(\mathrm F^i(S))_{i\in\mathbb Z}$ tal que
$\mathrm F^0(S)=S$ (y por tanto
$\operatorname{gr}^i(S)=0$ para $i<0$).
Supongamos que cada $A_k$, con la filtración definida en
\sref{0.13.7.1}, es un
$S$--$\boldsymbol{C}$-\emph{módulo filtrado} en el sentido de \sref{0.13.6.6}, y que los
morfismos $A_h\to A_k$ para $k\leq h$ son morfismos de
$S$--$\boldsymbol{C}$-módulos filtrados.
Para abreviar, diremos entonces que $\mathbf A$ es un
\emph{sistema proyectivo de $S$--$\boldsymbol{C}$-módulos filtrados}. Entonces es inmediato que los morfismos $\mathrm B_r^{pq}(A_h)\to\mathrm B_r^{pq}(A_k)$ y $\mathrm Z_r^{pq}(A_h)\to\mathrm Z_r^{pq}(A_k)$ para $k\leq h$ y
$1\leq r\leq+\infty$, son morfismos para las estructuras de $\operatorname{gr}^\bullet(S)$-$\boldsymbol{C}'$-\emph{módulo bigraduado} (\hyperref[0.13.6.5-fr]{13.6.5}), y que las familias $(\mathrm Z_r^{pq}(\mathbf A))$, $(\mathrm B_r^{pq}(\mathbf A))$ y $(\mathrm E_r^{pq}(\mathbf A))$ son $\operatorname{gr}^\bullet(S)$--$\boldsymbol{C}'$-\emph{módulos bigraduados}
(\sref{0.13.6.5}), y que, para $r$ finito, siendo los dos primeros submódulos de $(\mathrm E_1^{pq}(\mathbf A))$. Se denotará también por $\mathrm Z_r^{(n)}(\mathbf A)$,
$\mathrm B_r^{(n)}(\mathbf A)$ y
$\mathrm E_r^{(n)}(\mathbf A)$ para los respectivos subobjetos graduados
obtenidos conservando únicamente los términos con $p+q=n$; estos son
$\operatorname{gr}^\bullet(S)$--$\boldsymbol{C}'$-\emph{módulos graduados}.

Cuando el sistema proyectivo
$(\mathrm E_r^{pq}(\mathbf A))_{r\in\mathbb Z}$ es esencialmente constante,
$(\mathrm E_\infty^{pq}(\mathbf A))$ es asimismo un
$\operatorname{gr}^\bullet(S)$--$\boldsymbol{C}'$-módulo bigraduado, y cada
$\mathrm E_\infty^{(n)}(\mathbf A)$ es un
$\operatorname{gr}^\bullet(S)$--$\boldsymbol{C}'$-módulo graduado. Además, los $\beta^{p,n-p}:\mathrm E_\infty^{p,n-p}(A_k)\xrightarrow{\sim}
\operatorname{gr}^p(\mathrm R^n\mathrm T(A_k))$ constituyen para cada $k$ un isomorfismo para la estructura de $\operatorname{gr}^\bullet(S)$--$\boldsymbol{C}'$-módulos graduados de
$\mathrm E_\infty^{(n)}(A_k)$ sobre
$\operatorname{gr}^\bullet(\mathrm R^n\mathrm T(A_k))$.
Bajo las hipótesis precedentes, el isomorfismo
$\beta^{p,n-p}:\mathrm E_\infty^{(n)}(\mathbf A)\xrightarrow{\sim}
\varprojlim_k\operatorname{gr}^\bullet(\mathrm R^n\mathrm T(A_k))$
respeta, por tanto, también estas estructuras.
Lo mismo vale manifiestamente para el isomorfismo canónico
$\operatorname{gr}^\bullet(\mathrm R^n\mathrm T(\mathbf A)
\xrightarrow{\sim}
\varprojlim_k\operatorname{gr}^\bullet(\mathrm R^n\mathrm T(A_k))$,
por consiguiente, los isomorfismos \sref{0.13.7.4.1} forman un
isomorfismo de
$\operatorname{gr}^\bullet(S)$--$\boldsymbol{C}'$-\emph{módulos graduados}.
\end{env}

\begin{proposition}[13.7.7]
\label{0.13.7.7}
\label{0.13.7.7-fr}
Sea $S$ un anillo noetheriano $\mathfrak J$-ádico.
Supongamos que $\boldsymbol{C}$ es una categoría abeliana en la que todo objeto
es isomorfo a un subobjeto de un objeto inyectivo, y sea
$\mathrm T$ un funtor covariante aditivo de $\boldsymbol{C}$ en la
categoría de grupos abelianos.
Sea $\mathbf A=(A_k)_{k\in\mathbb Z}$ un
\oldpage[0\textsubscript{III}]{78}
sistema proyectivo estricto de $S$--$\boldsymbol{C}$-módulos filtrados, para la
filtración $\mathfrak J$-ádica sobre $S$, acotado inferiormente.
Para un entero fijo $n$ dado, se satisfaga la condición siguiente:
\begin{itemize}
\item[$(\mathrm F_n)$]
  \emph{El $\operatorname{gr}^\bullet(S)$-módulo graduado}
  $\mathrm E_1^{(m)}(\mathbf A)=
(\mathrm R^m\mathrm T(\operatorname{gr}^p(\mathbf A)))_{p\in\mathbb Z}$
  \emph{de \sref{0.13.7.3.1} es de tipo finito para
  $m=n$ y $m=n+1$.}
\end{itemize}
Entonces:
\begin{enumerate}
  \item[{\rm(i)}] Los sistemas proyectivos
  $(\mathrm R^n\mathrm T(A_k))_{k\in\mathbb Z}$ y
  $(\mathrm R^{n+1}\mathrm T(A_k))_{k\in\mathbb Z}$ satisfacen {\rm(ML)}.
  \item[{\rm(ii)}] Si
  $\mathrm R'^{\,n}\mathrm T(\mathbf A)=
\varprojlim_k\mathrm R^n\mathrm T(A_k)$,
  entonces $\mathrm R'^{\,n}\mathrm T(\mathbf A)$ es un $S$-módulo de tipo finito.
\item[\rm(iii)] La filtración definida por $\mathrm F^p(\mathrm R'^{\,n}\mathrm T(\mathbf A))=
\operatorname{Ker}(\mathrm R'^{\,n}\mathrm T(\mathbf A)
\to\mathrm R^n\mathrm T(A_{p-1}))$ $(p\in\mathbb Z)$ sobre $\mathrm R'^{\,n}\mathrm T(\mathbf A)$ es $\mathfrak J$-buena; es decir,
  $\mathfrak J\mathrm F^p(\mathrm R'^{\,n}\mathrm T(\mathbf A))
\subset\mathrm F^{p+1}(\mathrm R'^{\,n}(\mathrm T(\mathbf A))$
  para todo $p$, con igualdad para todo $p$ suficientemente grande.
  En particular, la topología sobre
  $\mathrm R'^{\,n}\mathrm T(\mathbf A)$ definida por esta filtración es la
  topología $\mathfrak J$-ádica.
  \item[{\rm(iv)}] El sistema proyectivo
  $(\mathrm E_r^{pq}(\mathbf A))_{r\in\mathbb Z}$ es esencialmente constante siempre que
  $p+q=n$ o $p+q=n+1$.
  Así, $\mathrm E_\infty^{pq}(\mathbf A)$ está, pues, definido \textup{(\hyperref[0.13.5.7-fr]{13.5.7})} y se tiene un isomorfismo canónico de $\operatorname{gr}^\bullet(S)$-módulos graduados
\begin{equation}
\label{0.13.7.7.1-fr}
\tag{13.7.7.1}
\operatorname{gr}^p(\mathrm R'^{\,n}\mathrm T(\mathbf A))
\xrightarrow{\sim}\mathrm E_\infty^{p,n-p}(\mathbf A)
\quad(p\in\mathbb Z).
\end{equation}
\end{enumerate}
\end{proposition}

\begin{proof}
Se observará que el isomorfismo (\hyperref[0.13.7.7.1-fr]{13.7.7.1}) permitirá denotar por $\mathrm R^n\mathrm T(\mathbf A)$ para el límite proyectivo
$\mathrm R'^{\,n}\mathrm T(\mathbf A)$ del sistema proyectivo
$\mathrm R^n\mathrm T(\mathbf A)$.

El anillo graduado $\operatorname{gr}^\bullet(S)$ es noetheriano
(Bourbaki, \emph{Alg.\ comm.}, cap.~III, \textsection2,
n\textsuperscript{o}~9, cor.~5 del teorema~2).
Por tanto, la sucesión creciente de
$\operatorname{gr}^\bullet(S)$-submódulos graduados
$\mathrm B_r^{(m)}(\mathbf A)$ de
$\mathrm E_1^{(m)}(\mathbf A)$ es \emph{estacionaria} para $m=n$ y
$m=n+1$.
La condición \emph{(b)} de \sref{0.11.1.10} queda, por tanto, satisfecha.
Se sigue que la condición \emph{(a)} de \sref{0.13.7.4} se cumple para
$n$ y para $n+1$, lo que demuestra {\rm(i)}.

Además, los isomorfismos \sref{0.13.7.4.1}, junto con las
observaciones de \sref{0.13.7.6}, muestran que
$\operatorname{gr}^\bullet(\mathrm R^n\mathrm T(\mathbf A))$ es un
$\operatorname{gr}^\bullet(S)$-módulo graduado isomorfo a
$\mathrm E_\infty^{(n)}(\mathbf A)=
\mathrm Z_\infty^{(n)}(\mathbf A)/\mathrm B_\infty^{(n)}(\mathbf A)$
Como $\mathrm Z_\infty^{(n)}(\mathbf A)$ es un submódulo de
$\mathrm E_1^{(n)}(\mathbf A)$, es de \emph{tipo finito}, y también lo es
$\mathrm E_\infty^{(n)}(\mathbf A)$.
Finalmente, para la filtración
$(\mathrm F^p(\mathrm R'^{\,n}\mathrm T(\mathbf A)))$, resulta de (\hyperref[0.13.4.5-fr]{13.4.5}) que $\operatorname{gr}^\bullet(\mathrm R^n\mathrm T(\mathbf A))$ y $\operatorname{gr}^\bullet(\mathrm R'^{\,n}\mathrm T(\mathbf A))$ son $\operatorname{gr}^\bullet(S)$-módulos graduados isomorfos.
Esto demuestra {\rm(iv)}.
Las aserciones {\rm(ii)} y {\rm(iii)} se seguirán de estos resultados
y del lema siguiente.
\end{proof}

\begin{lemma}[13.7.7.2]
\label{0.13.7.7.2}
\label{0.13.7.7.2-fr}
Sea $S$ un anillo noetheriano $\mathfrak J$-ádico, y sea $M$
un $S$-módulo dotado de una filtración codiscreta
$(\mathrm F^p(M))_{p\in\mathbb Z}$ tal que
$\mathfrak J\mathrm F^p(M)\subset\mathrm F^{p+1}(M)$
Así, $M$ es un módulo filtrado sobre $S$ dotado de su
filtración $\mathfrak J$-ádica.
Supongamos además que $M$ es separado para la topología definida por
$(\mathrm F^p(M))$.
Entonces las condiciones siguientes son equivalentes:
\begin{enumerate}
  \item[{\rm a)}] $M$ es un $S$-módulo de tipo finito y
  $(\mathrm F^p(M))$ es una filtración $\mathfrak J$-buena.
  \item[{\rm b)}] $\operatorname{gr}^\bullet(M)$ es un
  $\operatorname{gr}^\bullet(S)$-módulo de tipo finito.
  \item[{\rm c)}] Los $\operatorname{gr}^p(M)$ son $S$-módulos de
  tipo finito y, para todo $p$ suficientemente grande, los homomorfismos
  canónicos
  \begin{equation}
\label{0.13.7.7.3-fr}
\tag{13.7.7.3}
\mathfrak J\otimes_S\operatorname{gr}^p(M)\longrightarrow
\operatorname{gr}^{p+1}(M)
\end{equation}
  \oldpage[0\textsubscript{III}]{79}
  son sobreyectivos.
  Están inducidos por
  $\mathfrak J\otimes_S\mathrm F^p(M)\to\mathrm F^{p+1}(M)$, teniendo en cuenta el hecho
  de que la imagen del homomorfismo compuesto
  $\mathfrak J\otimes_S\mathrm F^{p+1}(M)\to
\mathfrak J\otimes_S\mathrm F^p(M)\to\mathrm F^{p+1}(M)$
  es
  $\mathfrak J\mathrm F^{p+1}(M)\subset\mathrm F^{p+2}(M)$.
\end{enumerate}

\textup{Para la demostración, véase Bourbaki, \emph{Alg.\ comm.}, cap.~III,
\textsection3, n\textsuperscript{o}~1, prop.~3.}
\end{lemma}

\begin{env}[13.7.7.4]
\label{0.13.7.7.4}
\label{0.13.7.7.4-fr}
Para aplicar el lema \sref{0.13.7.7.2}, queda observar que la
topología sobre $\mathrm R'^{\,n}\mathrm T(\mathbf A)$ por la filtración considerada hace de $\mathrm R'^{\,n}\mathrm T(\mathbf A)$ un $S$-módulo separado y completo:
esta es la topología de límite proyectivo de los grupos discretos
$\mathrm R^n\mathrm T(A_k)$.
Por otra parte, si $A_k=0$ para $k<k_0$, entonces también
$\mathrm R^n\mathrm T(A_k)=0$ para $k<k_0$; por tanto
$\mathrm F^{k_0}(\mathrm R'^{\,n}\mathrm T(\mathbf A))=
\mathrm R'^{\,n}\mathrm T(\mathbf A)$,
Así, todas las hipótesis del lema quedan efectivamente satisfechas.
\end{env}

\begin{corollary}[13.7.8]
\label{0.13.7.8}
\label{0.13.7.8-fr}
Si se cumple la hipótesis $(\mathrm F_n)$, entonces para todo $f\in S$
existe un isomorfismo canónico
\begin{equation}
\label{0.13.7.8.1-fr}
\tag{13.7.8.1}
\varprojlim_k\bigl((\mathrm R^n\mathrm T(A_k))_f\bigr)
\xrightarrow{\sim}\mathrm R^n\mathrm T(\mathbf A)\otimes_S S_{\{f\}}.
\end{equation}
\end{corollary}

\begin{proof}
En efecto, $\mathrm R^n\mathrm T(\mathbf A)$ es un $S$-módulo de tipo finito,
y $S_{\{f\}}$ es una $S$-álgebra ádica noetheriana
(\sref[0\textsubscript{I}]{0.7.6.11}), a saber, la completación separada
de $S_f$ para la topología $\mathfrak J$-preádica
(\sref[0\textsubscript{I}]{0.7.6.2}).
Se sigue de \sref[0\textsubscript{I}]{0.7.7.8} y
\sref[0\textsubscript{I}]{0.7.7.1} que
$\mathrm R^n\mathrm T(\mathbf A)\otimes_S S_{\{f\}}$
es isomorfo a la completación separada de
$\mathrm R^n\mathrm T(\mathbf A)\otimes_S S_f$ para la
topología $\mathfrak J$-preádica; un sistema fundamental de vecindades de $0$ para esta topología es
$(\mathfrak J^p\mathrm R^n\mathrm T(\mathbf A))\otimes_S S_f$,
los grupos
$\mathrm F^p(\mathrm R^n\mathrm T(\mathbf A))\otimes_S S_f$
forman, por tanto, también un sistema tal.
Este último grupo es el núcleo de la aplicación canónica
$(\mathrm R^n\mathrm T(\mathbf A))_f\to
(\mathrm R^n\mathrm T(A_{p-1}))_f$
Por consiguiente, el grupo separado asociado a
$\mathrm R^n\mathrm T(\mathbf A)\otimes_S S_f$ se identifica con un subgrupo
$G$ de
$\varprojlim_k((\mathrm R^n\mathrm T(A_k))_f)$.
Pero el sistema proyectivo
$((\mathrm R^n\mathrm T(A_k))_f)$ satisface manifiestamente {\rm(ML)}, y la
imagen de $(\mathrm R^n\mathrm T(\mathbf A))_f$ en cada
$(\mathrm R^n\mathrm T(A_k))_f$ es igual a la imagen común de los
$(\mathrm R^n\mathrm T(A_h))_f$ para todo $h\geq k$ suficientemente grande.
Se sigue inmediatamente que $G$ es \emph{denso en todas partes} en
$\varprojlim_k((\mathrm R^n\mathrm T(A_k))_f)$.
Como este último grupo es separado y completo, queda demostrado el
corolario.
\end{proof}

% END tree/repo/editions/III1_DIPLOMATIC/ega0/ega0-13.tex

\begin{flushright}
\textit{(Continuará.)}
\end{flushright}

\clearpage
\oldpage[0\textsubscript{III}]{80}
\thispagestyle{empty}
\null
\clearpage
\thispagestyle{plain}
\markboth{ESTUDIO COHOMOLÓGICO DE LOS HACES COHERENTES}{CAPÍTULO III}
\oldpage[III]{81}
\phantomsection
\label{chapter:III-fr}

\begin{center}
{\large CAPÍTULO III\par}
\vspace{1.5em}
{\Large\bfseries ESTUDIO COHOMOLÓGICO\par}
{\Large\bfseries DE LOS HACES COHERENTES\par}
\vspace{1.8em}
{\bfseries Sumario\par}
\end{center}

\medskip
\begin{tabular}{@{}r@{\ }p{0.82\textwidth}@{}}
\S~1. & Cohomología de los esquemas afines.\\
\S~2. & Estudio cohomológico de los morfismos proyectivos.\\
\S~3. & El teorema de finitud para los morfismos propios.\\
\S~4. & El teorema fundamental de los morfismos propios; aplicaciones.\\
\S~5. & Un teorema de existencia de haces algebraicos coherentes.\\
\S~6. & Funtores Tor y Ext locales y globales; fórmula de Künneth.\\
\S~7. & Estudio del cambio de base en los funtores cohomológicos covariantes de Módulos.\\
\S~8. & El teorema de dualidad sobre los fibrados proyectivos.\\
\S~9. & Cohomología relativa y cohomología local; dualidad local.\\
\S~10. & Relaciones entre cohomología proyectiva y cohomología local. Técnica de compleción formal a lo largo de un divisor.\\
\S~11. & Grupos de Picard globales y locales\footnotemark.
\end{tabular}

\medskip
Este capítulo da los teoremas fundamentales sobre la cohomología de los haces algebraicos coherentes, con excepción de los que se deducen de la teoría de residuos (teoremas de dualidad), que serán objeto de un capítulo posterior. Entre los primeros, hay esencialmente seis teoremas fundamentales, que constituyen el objeto de los seis primeros parágrafos del presente capítulo. Estos resultados serán en lo sucesivo herramientas esenciales, incluso en cuestiones que no son de naturaleza propiamente cohomológica; el lector verá sus primeros ejemplos ya en el \S~4. El \S~7 da resultados de naturaleza más técnica, pero de uso constante en las aplicaciones. Finalmente, en los \S\S~8 a 11, desarrollaremos ciertos resultados, ligados a la dualidad de los haces coherentes, particularmente importantes para las aplicaciones, y que pueden exponerse antes de la teoría general de residuos.

\footnotetext{El capítulo IV no depende de los \S\S~8 a 11, y probablemente se publicará antes que estos últimos.}

\oldpage[III]{82}
El contenido de los \S\S~1 y 2 se debe a J.-P. Serre, y el lector constatará que solo hemos tenido que seguir la exposición de (FAC). Los \S\S~8 y 9 están igualmente inspirados por (FAC) (siendo, sin embargo, menos evidentes las transposiciones exigidas por los diferentes contextos). Finalmente, como hemos dicho en la Introducción, el \S~4 debe considerarse como la formulación, en lenguaje moderno, del «teorema de invariancia» fundamental de la «teoría de las funciones holomorfas» de Zariski.

Señalemos finalmente que los resultados del n\textsuperscript{o}~3.4 (y las proposiciones preliminares de (\textbf{0}, 13.4 a 13.7)) no se utilizarán en lo que sigue del cap.~III y pueden, por tanto, omitirse en una primera lectura.

\setcounter{section}{0}


% BEGIN tree/repo/editions/III1_DIPLOMATIC/ega3/ega3-1.tex
\section{Cohomología de los esquemas afines}
\label{section:III.1}
\label{section:III.1-fr}

\subsection{Repaso del complejo del álgebra exterior}
\label{subsection:III.1.1}
\label{subsection:III.1.1-fr}

\begin{env}[1.1.1]
\label{III.1.1.1}
\label{III.1.1.1-fr}
Sea $A$ un anillo y sea $\mathbf f=(f_i)_{1\leq i\leq r}$ un sistema de $r$ elementos de $A$.
El \emph{complejo del álgebra exterior $K_\bullet(\mathbf f)$} correspondiente a $\mathbf f$ es un complejo de cadenas (G, I, 2.2) definido del modo siguiente: el $A$-módulo graduado $K_\bullet(\mathbf f)$ es igual al \emph{álgebra exterior $\bigwedge(A^r)$}, graduada de la manera usual, y la aplicación borde es la \emph{multiplicación interior $i_{\mathbf f}$} por $\mathbf f$ considerado como elemento del dual $(A^r)\check{}$; se recuerda que $i_{\mathbf f}$ es una \emph{antiderivación} de grado $-1$ de $\bigwedge(A^r)$, y que, si $(\mathbf e_i)_{1\leq i\leq r}$ es la base canónica de $A^r$, se tiene $i_{\mathbf f}(\mathbf e_i)=f_i$; la verificación de la condición $i_{\mathbf f}\circ i_{\mathbf f}=0$ es inmediata.

Una definición equivalente es la siguiente: para cada $i$, se considera un complejo de cadenas $K_\bullet(f_i)$ definido por $K_0(f_i)=K_1(f_i)=A$, $K_n(f_i)=0$ para $n\neq0,1$; la aplicación borde queda definida por la condición de que $d_1:A\to A$ sea la \emph{multiplicación por $f_i$}.
Se toma entonces como $K_\bullet(\mathbf f)$ el \emph{producto tensorial $K_\bullet(f_1)\otimes K_\bullet(f_2)\otimes\ldots\otimes K_\bullet(f_r)$} (G, I, 2.7), provisto de su grado total; la verificación del isomorfismo de este complejo con el definido anteriormente es inmediata.
\end{env}

\begin{env}[1.1.2]
\label{III.1.1.2}
\label{III.1.1.2-fr}
Para todo $A$-módulo $M$, se define el \emph{complejo de cadenas}
\begin{equation}
\label{III.1.1.2.1-fr}
\tag{1.1.2.1}
K_\bullet(\mathbf f,M)=K_\bullet(\mathbf f)\otimes_A M
\end{equation}
y el \emph{complejo de cocadenas} (G, I, 2.2)
\begin{equation}
\label{III.1.1.2.2-fr}
\tag{1.1.2.2}
K^\bullet(\mathbf f,M)=\operatorname{Hom}_A(K_\bullet(\mathbf f),M).
\end{equation}

Si $g$ es una $k$-cocadena de este último complejo y se pone
\[
g(i_1,\ldots,i_k)=g(\mathbf e_{i_1}\wedge\ldots\wedge\mathbf e_{i_k}),
\]
entonces $g$ se identifica con una aplicación \emph{alternada} de $[1,r]^k$ en $M$, y de las definiciones anteriores se deduce que
\begin{equation}
\label{III.1.1.2.3-fr}
\tag{1.1.2.3}
d^kg(i_1,i_2,\ldots,i_{k+1})=
\sum_{h=1}^{k+1}(-1)^{h-1}f_{i_h}
g(i_1,\ldots,\widehat{i}_h,\ldots,i_{k+1}).
\end{equation}
\end{env}

\oldpage[III]{83}

\begin{env}[1.1.3]
\label{III.1.1.3-fr}
De los complejos precedentes se deducen, como de costumbre, los \emph{$A$-módulos
de homología y de cohomología} (G, I, 2.2)
\begin{align}
\label{III.1.1.3.1-fr}
\tag{1.1.3.1}
H_\bullet(\mathbf f,M)&=H_\bullet(K_\bullet(\mathbf f,M))\\
\label{III.1.1.3.2-fr}
\tag{1.1.3.2}
H^\bullet(\mathbf f,M)&=H^\bullet(K^\bullet(\mathbf f,M)).
\end{align}
Se define además un $A$-\emph{isomorfismo}
$K_\bullet(\mathbf f,M)\xrightarrow{\sim}K^\bullet(\mathbf f,M)$ haciendo
corresponder a toda cadena
$z=\sum(\mathbf e_{i_1}\wedge\ldots\wedge\mathbf e_{i_k})\otimes
z_{i_1\ldots i_k}$ la cocadena $g_z$ tal que
$g_z(j_1,\ldots,j_{r-k})=\varepsilon z_{i_1\ldots i_k}$, donde
$(j_h)_{1\leq h\leq r-k}$ es la sucesión estrictamente creciente complementaria
de la sucesión estrictamente creciente $(i_h)_{1\leq h\leq k}$ en $[1,r]$ y
$\varepsilon=(-1)^\nu$, siendo $\nu$ el número de inversiones de la permutación
$i_1,\ldots,i_k,j_1,\ldots,j_{r-k}$ de $[1,r]$. Se verifica que
$g_{dz}=d(g_z)$, lo que da un isomorfismo
\begin{equation}
\label{III.1.1.3.3-fr}
\tag{1.1.3.3}
H^i(\mathbf f,M)\xrightarrow{\sim}H_{r-i}(\mathbf f,M)
\qquad\text{para }0\leq i\leq r.
\end{equation}
En este capítulo tendremos que considerar sobre todo los módulos de cohomología
$H^\bullet(\mathbf f,M)$.

Para un $\mathbf f$ dado, es inmediato (G, I, 2.1) que
$M\leadsto H^\bullet(\mathbf f,M)$ es un \emph{funtor cohomológico}
(T, II, 2.1), de la categoría de los $A$-módulos en la de los $A$-módulos
graduados, nulo para los grados $<0$ y $>r$. Además, se tiene
\begin{equation}
\label{III.1.1.3.4-fr}
\tag{1.1.3.4}
H^0(\mathbf f,M)=\operatorname{Hom}_A(A/(\mathbf f),M)
\end{equation}
designando por $(\mathbf f)$ el ideal de $A$ engendrado por
$f_1,\ldots,f_r$; esto resulta inmediatamente de
(\hyperref[III.1.1.2.3-fr]{1.1.2.3}), y está claro que
$H^0(\mathbf f,M)$ se identifica con el submódulo de $M$ \emph{anulado por
$(\mathbf f)$}. Análogamente, se tiene según
(\hyperref[III.1.1.2.3-fr]{1.1.2.3})
\begin{equation}
\label{III.1.1.3.5-fr}
\tag{1.1.3.5}
H^r(\mathbf f,M)=M/\Bigl(\sum_{i=1}^r f_iM\Bigr)
=(A/(\mathbf f))\otimes_A M.
\end{equation}
Utilizaremos el siguiente resultado conocido, del que recordaremos una
demostración para ser completos:
\end{env}

\begin{proposition}[1.1.4]
\label{III.1.1.4}
\label{III.1.1.4-fr}
Sean $A$ un anillo, $\mathbf f=(f_i)_{1\leq i\leq r}$ una familia finita de elementos de $A$ y $M$ un $A$-módulo.
Si, para $1\leq i\leq r$, la aplicación $z\leadsto f_i.z$ sobre $M_{i-1}=M/(f_1M+\ldots+f_{i-1}M)$ es inyectiva, entonces $H^i(\mathbf f,M)=0$ para $i\neq r$.
\end{proposition}

\begin{proof}
Basta probar que $H_i(\mathbf f,M)=0$ para todo $i>0$, en virtud de (\hyperref[III.1.1.3.3]{1.1.3.3}).
Razonemos por inducción sobre $r$, siendo trivial el caso $r=0$ siendo trivial. Pongamos $\mathbf f'=(f_i)_{1\leq i\leq r-1}$; esta familia satisface las condiciones del enunciado, de modo que, si se pone $L_\bullet=K_\bullet(\mathbf f',M)$, por hipótesis se tiene $H_i(L_\bullet)=0$ para $i>0$, y $H_0(L_\bullet)=M_{r-1}$ en virtud de (\hyperref[III.1.1.3.3]{1.1.3.3}) y (\hyperref[III.1.1.3.5]{1.1.3.5}).
Para abreviar, pongamos $K_\bullet=K_\bullet(f_r)=K_0\oplus K_1$, con $K_0=K_1=A$ y $d_1:K_1\to K_0$ igual a la multiplicación por $f_r$; por definición \sref{III.1.1.1}, se tiene $K_\bullet(\mathbf f,M)=K_\bullet\otimes_A L_\bullet$.
Se dispone del siguiente lema:

\begin{lemma}[1.1.4.1]
\label{III.1.1.4.1}
\label{III.1.1.4.1-fr}
Sea $K_\bullet$ un complejo de cadenas de $A$-módulos libres, nulo salvo en las dimensiones $0$ y $1$.
Para todo complejo de cadenas $L_\bullet$ de $A$-módulos se tiene una sucesión exacta
\[
0\longrightarrow H_0(K_\bullet\otimes H_p(L_\bullet))
\longrightarrow H_p(K_\bullet\otimes L_\bullet)
\longrightarrow H_1(K_\bullet\otimes H_{p-1}(L_\bullet))
\longrightarrow0
\]
para todo índice $p$.
\end{lemma}

\oldpage[III]{84}
Este es un caso particular de una sucesión exacta de términos de grado bajo de la sucesión espectral de K\"unneth (M, XVII, 5.2 \emph{(a)} y G, I, 5.5.2); puede demostrarse directamente como sigue.
Consideremos $K_0$ y $K_1$ como complejos de cadenas (nulos en las dimensiones $\neq0$ y $\neq1$, respectivamente); se tiene entonces una sucesión exacta de complejos
\[
0\longrightarrow K_0\otimes L_\bullet
\longrightarrow K_\bullet\otimes L_\bullet
\longrightarrow K_1\otimes L_\bullet\longrightarrow0
\] a la cual podemos aplicar la sucesión exacta de homología \[
\ldots\longrightarrow H_{p+1}(K_1\otimes L_\bullet)
\xrightarrow{\partial}H_p(K_0\otimes L_\bullet)
\longrightarrow H_p(K_\bullet\otimes L_\bullet)
\longrightarrow H_p(K_1\otimes L_\bullet)
\xrightarrow{\partial}H_{p-1}(K_0\otimes L_\bullet)
\longrightarrow\ldots
\]
Pero es evidente que $H_p(K_0\otimes L_\bullet)=K_0\otimes H_p(L_\bullet)$ y $H_p(K_1\otimes L_\bullet)=K_1\otimes H_{v-1}(L_\bullet)$ para todo $p$; además, se comprueba inmediatamente que el operador $\partial:K_1\otimes H_p(L_\bullet)\to K_0\otimes H_p(L_\bullet)$ no es otro que $d_1\otimes1$; el lema se deduce así de la sucesión exacta anterior y de la definición de $H_0(K\otimes H_p(L_\bullet))$ y $H_1(K_\bullet\otimes H_{p-1}(L_\bullet))$.

Establecido el lema, el final de la demostración de la Proposición \sref{III.1.1.4} es inmediato: la hipótesis de inducción y el Lema \sref{III.1.1.4.1} dan $H_p(K_\bullet\otimes L_\bullet)=0$ para $p\geq2$; además, si se prueba que $H_1(K_\bullet,H_0(L_\bullet))=0$, se deduce también del Lema \sref{III.1.1.4.1} que $H_1(K_\bullet\otimes L_\bullet)=0$; pero, por definición, $H_1(K_\bullet,H_0(L_\bullet))$ no es otro que el núcleo de la aplicación $z\leadsto f_r.z$ sobre $M_{r-1}$, y, como por hipótesis este núcleo es nulo, queda terminada la demostración.
\end{proof}

\begin{env}[1.1.5]
\label{III.1.1.5}
\label{III.1.1.5-fr}
Sea $\mathbf g=(g_i)_{1\leq i\leq r}$ una segunda sucesión de $r$ elementos de $A$, y pongamos $\mathbf f\mathbf g=(f_ig_i)_{1\leq i\leq r}$.
Puede definirse un homomorfismo canónico de complejos
\begin{equation}
\label{III.1.1.5.1-fr}
\tag{1.1.5.1}
\varphi_{\mathbf g}:K_\bullet(\mathbf f\mathbf g)\longrightarrow
K_\bullet(\mathbf f)
\end{equation}
como la extensión canónica al álgebra exterior $\bigwedge(A^r)$ de la aplicación $A$-lineal $(x_1,\ldots,x_r)\leadsto(g_1x_1,\ldots,g_rx_r)$ de $A^r$ en sí mismo.
Para ver que se trata de un homomorfismo de complejos, basta observar, en general, que, si $u:E\to F$ es una aplicación $A$-lineal, $\mathbf x\in\check F$ y $\mathbf y={}^tu(\mathbf x)\in\check E$, entonces se tiene la fórmula
\begin{equation}
\label{III.1.1.5.2-fr}
\tag{1.1.5.2}
(\bigwedge u)\circ i_{\mathbf y}=i_{\mathbf x}\circ(\bigwedge u)~;
\end{equation}
en efecto, los dos miembros son antiderivaciones de $\bigwedge F$, y basta verificar que coinciden en $F$, lo cual se deduce inmediatamente de las definiciones.

Cuando se identifica $K_\bullet(\mathbf f)$ con el producto tensorial de los $K_\bullet(f_i)$ \sref{III.1.1.1}, $\varphi_{\mathbf g}$ es el producto tensorial de los $\varphi_{g_i}$, donde $\varphi_{g_i}$ es la identidad en grado $0$ y la multiplicación por $g_i$ en grado $1$.
\end{env}

\begin{env}[1.1.6]
\label{III.1.1.6-fr}
En particular, para todo par de enteros $m,n$ tales que $0\leq n\leq m$, se tienen homomorfismos de complejos
\begin{equation}
\label{III.1.1.6.1-fr}
\tag{1.1.6.1}
\varphi_{\mathbf f^{m-n}}:K_\bullet(\mathbf f^m)\longrightarrow
K_\bullet(\mathbf f^n)
\end{equation}
y por consiguiente homomorfismos
\begin{align}
\label{III.1.1.6.2-fr}
\tag{1.1.6.2}
\varphi_{\mathbf f^{m-n}} &:K^\bullet(\mathbf f^n,M)\longrightarrow
K^\bullet(\mathbf f^m,M)\\
\label{III.1.1.6.3-fr}
\tag{1.1.6.3}
\varphi_{\mathbf f^{m-n}} &:H^\bullet(\mathbf f^n,M)\longrightarrow
H^\bullet(\mathbf f^m,M).
\end{align}

\oldpage[III]{85}
Estos últimos satisfacen evidentemente la condición de transitividad
$\varphi_{\mathbf f^{m-p}}=
\varphi_{\mathbf f^{m-n}}\circ\varphi_{\mathbf f^{n-p}}$ para $p\leq n\leq m$; definen, por tanto, dos \emph{sistemas inductivos} de $A$-módulos; se pondrá
\begin{align}
\label{III.1.1.6.4-fr}
\tag{1.1.6.4}
C^\bullet((\mathbf f),M)&=\varinjlim_n K^\bullet(\mathbf f^n,M)\\
\label{III.1.1.6.5-fr}
\tag{1.1.6.5}
H^\bullet((\mathbf f),M)&=H^\bullet(C^\bullet((\mathbf f),M))
=\varinjlim_n H^\bullet(\mathbf f^n,M)
\end{align}
proviniendo la última igualdad de que el paso al límite inductivo conmuta con el funtor $H^\bullet$ (G, I, 2.1).
Se verá más adelante \sref{III.1.4.3} que $H^\bullet((\mathbf f),M)$ depende en realidad solo del \emph{ideal $(\mathbf f)$} de $A$ (e incluso solo de la topología $(\mathbf f)$-preádica sobre $A$), lo que justifica las notaciones.

Es claro que $M\leadsto C^\bullet((\mathbf f),M)$ es un funtor exacto $A$-lineal y que $M\leadsto H^\bullet((\mathbf f),M)$ es un funtor cohomológico.
\end{env}

\begin{env}[1.1.7]
\label{III.1.1.7}
\label{III.1.1.7-fr}
Pónganse $\mathbf f=(f_i)\in A^r$ y $\mathbf g=(g_i)\in A^r$; designemos por $e_{\mathbf g}$ la multiplicación por la izquierda por el vector $\mathbf g\in A^r$ sobre el álgebra exterior $\bigwedge(A^r)$; se sabe que se tiene la \emph{fórmula de homotopía}
\begin{equation}
\label{III.1.1.7.1-fr}
\tag{1.1.7.1}
i_{\mathbf f}e_{\mathbf g}+e_{\mathbf g}i_{\mathbf f}
=\langle\mathbf g,\mathbf f\rangle 1
\end{equation}
en el $A$-módulo $A^r$ ($1$ denota el automorfismo identidad de $A^r$); esta relación implica también que \emph{en el complejo $K_\bullet(\mathbf f)$} se tiene
\begin{equation}
\label{III.1.1.7.2-fr}
\tag{1.1.7.2}
de_{\mathbf g}+e_{\mathbf g}d=\langle\mathbf g,\mathbf f\rangle 1.
\end{equation}

Si el ideal $(\mathbf f)$ es igual a $A$, existe $\mathbf g\in A^r$ tal que $\langle\mathbf g,\mathbf f\rangle=\sum_{i=1}^r g_if_i=1$.
Por consiguiente (G, I, 2.4):
\end{env}

\begin{proposition}[1.1.8]
\label{III.1.1.8}
\label{III.1.1.8-fr}
Supongamos que el ideal $(\mathbf f)$ generado por los $f_i$ es igual a $A$.
Entonces el complejo $K_\bullet(\mathbf f)$ es homotópicamente trivial, al igual que los complejos $K_\bullet(\mathbf f,M)$ y $K^\bullet(\mathbf f,M)$ para todo $A$-módulo $M$.
\end{proposition}

\begin{corollary}[1.1.9]
\label{III.1.1.9}
\label{III.1.1.9-fr}
Si $(\mathbf f)=A$, entonces $H^\bullet(\mathbf f,M)=0$ y $H^\bullet((\mathbf f),M)=0$ para todo $A$-módulo $M$.
\end{corollary}

\begin{proof}
En efecto, se tiene entonces $(\mathbf f^n)=A$ para todo $n$.
\end{proof}

\begin{remark}[1.1.10]
\label{III.1.1.10}
\label{III.1.1.10-fr}
Con las mismas notaciones anteriores, sean $X=\operatorname{Spec}(A)$ y sea $Y$ el subesquema cerrado de $X$ definido por el ideal $(\mathbf f)$. Probaremos en el \S~9 que $H^\bullet((\mathbf f),M)$ es isomorfo a la cohomología $H_Y^\bullet(X,\widetilde M)$ correspondiente al antifiltro $\Phi$ de subconjuntos cerrados de $Y$ (T, 3.2). Mostraremos también que la prop. (\hyperref[III.1.2.3-fr]{1.2.3}) aplicada a $X$ y a $\mathcal F=\widetilde M$, es un caso particular de una sucesión exacta de cohomología
\[
\ldots\longrightarrow H_Y^p(X,\mathcal F)\longrightarrow H^p(X,\mathcal F)
\longrightarrow H^p(X\setmin Y,\mathcal F)\longrightarrow H_Y^{p+1}(X,\mathcal F)
\longrightarrow\ldots
\]
\end{remark}

\subsection{Cohomología de \v Cech de un recubrimiento abierto}
\label{subsection:III.1.2}
\label{subsection:III.1.2-fr}

\begin{notation}[1.2.1]
\label{III.1.2.1}
\label{III.1.2.1-fr}
En esta sección se denota por:
\begin{itemize}
  \item $X$ un preesquema;
  \item $\mathcal F$ un $\mathcal O_X$-módulo cuasicoherente;
\oldpage[III]{86}
  \item $A=\Gamma(X,\mathcal O_X)$, $M=\Gamma(X,\mathcal F)$;
  \item $\mathbf f=(f_i)_{1\leq i\leq r}$ un sistema finito de elementos de $A$;
  \item $U_i=X_{f_i}$, el abierto \sref[0]{0.5.5.2} de los $x\in X$ tales que $f_i(x)\neq0$;
  \item $U=\bigcup_{i=1}^r U_i$;
  \item $\mathfrak U$ el recubrimiento $(U_i)_{1\leq i\leq r}$ de $U$.
\end{itemize}
\end{notation}

\begin{env}[1.2.2]
\label{III.1.2.2}
\label{III.1.2.2-fr}
Supongamos que $X$ es un preesquema cuyo espacio subyacente es \emph{noetheriano}, o bien un \emph{esquema} cuyo espacio subyacente es \emph{cuasicompacto}.
Se sabe entonces \sref[I]{I.9.3.3} que $\Gamma(U_i,\mathcal F)=M_{f_i}$. Pondremos \[
U_{i_0i_1\ldots i_p}=\bigcap_{k=0}^p U_{i_k}
=X_{f_{i_0}f_{i_1}\ldots f_{i_p}}
\]
\sref[0]{0.5.5.3}; se tiene, por tanto, también
\begin{equation}
\label{III.1.2.2.1-fr}
\tag{1.2.2.1}
\Gamma(U_{i_0i_1\ldots i_p},\mathcal F)=M_{f_{i_0}f_{i_1}\ldots f_{i_p}}
\end{equation}

Por \sref[0]{0.1.6.1}, $M_{f_{i_0}f_{i_1}\ldots f_{i_p}}$ se identifica con el límite inductivo $\varinjlim_n M_{i_0i_1\ldots i_p}^{(n)}$, donde el sistema inductivo está formado por los $M_{i_0i_1\ldots i_p}^{(n)}=M$, siendo los homomorfismos $\varphi_{nm}:M_{i_0i_1\ldots i_p}^{(m)}\to
M_{i_0i_1\ldots i_p}^{(n)}$ la multiplicación por $(f_{i_0}f_{i_1}\ldots f_{i_p})^{n-m}$ para $m\leq n$. Designemos por $C_n^p(M)$ el conjunto de aplicaciones \emph{alternadas} de $[1,r]^{p+1}$ en $M$ (para todo $n$); estos $A$-módulos forman también un sistema inductivo respecto de los $\varphi_{nm}$.
Si $C^p(\mathfrak U,\mathcal F)$ es el grupo de las $p$-cocadenas \emph{alternadas} de \v Cech relativas al recubrimiento $\mathfrak U$, con coeficientes en $\mathcal F$ (G, II, 5.1), de lo anterior se deduce que puede escribirse
\begin{equation}
\label{III.1.2.2.2-fr}
\tag{1.2.2.2}
C^p(\mathfrak U,\mathcal F)=\varinjlim_n C_n^p(M).
\end{equation}

Con las notaciones de \sref{III.1.1.2}, $C_n^p(M)$ se identifica con $K^{p+1}(\mathbf f^n,M)$, y la aplicación $\varphi_{nm}$ se identifica con la aplicación $\varphi_{\mathbf f^{n-m}}$ definida en \sref{III.1.1.6}.
Se tiene así, para todo $p\geq0$, un isomorfismo canónico funtorial en $\mathcal F$
\begin{equation}
\label{III.1.2.2.3-fr}
\tag{1.2.2.3}
C^p(\mathfrak U,\mathcal F)\xrightarrow{\sim}C^{p+1}((\mathbf f),M).
\end{equation}

Además, la fórmula (\hyperref[III.1.1.2.3]{1.1.2.3}) y la definición de la cohomología de un recubrimiento (G, II, 5.1) muestran que los isomorfismos (\hyperref[III.1.2.2.3]{1.2.2.3}) son compatibles con las aplicaciones coborde.
\end{env}

\begin{proposition}[1.2.3]
\label{III.1.2.3}
\label{III.1.2.3-fr}
Si $X$ es un preesquema cuyo espacio subyacente es noetheriano, o un esquema cuyo espacio subyacente es cuasicompacto, existe un isomorfismo canónico funtorial en $\mathcal F$
\begin{equation}
\label{III.1.2.3.1-fr}
\tag{1.2.3.1}
H^p(\mathfrak U,\mathcal F)\xrightarrow{\sim}H^{p+1}((\mathbf f),M)
\qquad\text{para todo }p\geq1.
\end{equation}

Además, se tiene una sucesión exacta funtorial en $\mathcal F$
\begin{equation}
\label{III.1.2.3.2-fr}
\tag{1.2.3.2}
0\longrightarrow H^0((\mathbf f),M)\longrightarrow M
\longrightarrow H^0(\mathfrak U,\mathcal F)
\longrightarrow H^1((\mathbf f),M)\longrightarrow0.
\end{equation}
\end{proposition}

\oldpage[III]{87}
\begin{proof}
Los isomorfismos (\hyperref[III.1.2.3.1]{1.2.3.1}) son consecuencias inmediatas de lo visto en \sref{III.1.2.2}.
Por otra parte, se tiene $C^0(\mathfrak U,\mathcal F)=C^1((\mathbf f),M)$; por consiguiente, $H^0(\mathfrak U,\mathcal F)$ se identifica con el subgrupo de los $1$-cociclos de $C^1((\mathbf f),M)$; como $M=C^0((\mathbf f),M)$, la sucesión exacta (\hyperref[III.1.2.3.2]{1.2.3.2}) no es otra que la dada por la definición de los grupos de cohomología $H^0((\mathbf f),M)$ y $H^1((\mathbf f),M)$.
\end{proof}

\begin{corollary}[1.2.4]
\label{III.1.2.4}
\label{III.1.2.4-fr}
Supongamos que los $X_{f_i}$ son cuasicompactos y que existen $g_i\in\Gamma(U,\mathcal F)$ tales que
$\sum_i g_i(f_i\vert U)=1\vert U$. Entonces, para todo
$(\mathcal O_X\vert U)$-Módulo cuasicoherente $\mathcal G$, se tiene $H^p(\mathfrak U,\mathcal G)=0$ para $p>0$; si además $U=X$, el homomorfismo canónico \textup{(\hyperref[III.1.2.3.2]{1.2.3.2})} $M\to H^0(\mathfrak U,\mathcal F)$ es biyectivo.
\end{corollary}

\begin{proof}
Como, por hipótesis, los $U_i=X_{f_i}$ son cuasicompactos, también lo es $U$, y puede reducirse al caso $U=X$; la hipótesis implica entonces
$H^p((\mathbf f),M)=0$ para todo $p\geq0$ \sref{III.1.1.9}.
El corolario se deduce inmediatamente de (\hyperref[III.1.2.3.1]{1.2.3.1}) y (\hyperref[III.1.2.3.2]{1.2.3.2}).
\end{proof}

Se observará que, puesto que $H^0(\mathfrak U,\mathcal F)=H^0(U,\mathcal F)$ (G, II, 5.2.2), se ha demostrado de nuevo \sref[I]{I.1.3.7} como caso particular.

\begin{remark}[1.2.5]
\label{III.1.2.5}
\label{III.1.2.5-fr}
Supongamos que $X$ es un \emph{esquema afín}; entonces los $U_i=X_{f_i}=D(f_i)$ son abiertos afines, al igual que los $U_{i_0i_1\ldots i_p}$ (pero $U$ no es necesariamente afín).
En este caso, los funtores $\Gamma(X,\mathcal F)$ y $\Gamma(U_{i_0i_1\ldots i_p},\mathcal F)$ son exactos en $\mathcal F$ \sref[I]{I.1.3.11}.
Si se tiene una sucesión exacta $0\to\mathcal F'\to\mathcal F\to\mathcal F''\to0$ de $\mathcal O_X$-módulos cuasicoherentes, la sucesión de complejos
\[
0\longrightarrow C^\bullet(\mathfrak U,\mathcal F')
\longrightarrow C^\bullet(\mathfrak U,\mathcal F)
\longrightarrow C^\bullet(\mathfrak U,\mathcal F'')
\longrightarrow0
\]
es exacta y da, por tanto, una sucesión exacta de cohomología
\[
\ldots\longrightarrow H^p(\mathfrak U,\mathcal F')
\longrightarrow H^p(\mathfrak U,\mathcal F)
\longrightarrow H^p(\mathfrak U,\mathcal F'')
\xrightarrow{\partial}H^{p+1}(\mathfrak U,\mathcal F')
\longrightarrow\ldots
\]
Por otra parte, si se pone $M'=\Gamma(X,\mathcal F')$ y $M''=\Gamma(X,\mathcal F'')$, la sucesión $0\to M'\to M\to M''\to0$ es exacta; como $C^\bullet((\mathbf f),M)$ es un funtor exacto en $M$, se tiene asimismo la sucesión exacta de cohomología
\[
\ldots\longrightarrow H^p((\mathbf f),M')
\longrightarrow H^p((\mathbf f),M)
\longrightarrow H^p((\mathbf f),M'')
\xrightarrow{\partial}H^{p+1}((\mathbf f),M')
\longrightarrow\ldots
\]

Así pues, como el diagrama
\[
\xymatrix{
0\ar[r]&C^\bullet(\mathfrak U,\mathcal F')\ar[r]\ar[d]
&C^\bullet(\mathfrak U,\mathcal F)\ar[r]\ar[d]
&C^\bullet(\mathfrak U,\mathcal F'')\ar[r]\ar[d]&0\\
0\ar[r]&C^\bullet((\mathbf f),M')\ar[r]
&C^\bullet((\mathbf f),M)\ar[r]
&C^\bullet((\mathbf f),M'')\ar[r]&0
}
\]
\oldpage[III]{88}
es conmutativo, se concluye que los diagramas
\begin{equation}
\label{III.1.2.5.1-fr}
\tag{1.2.5.1}
\xymatrix{
H^p(\mathfrak U,\mathcal F'')\ar[r]^{\partial}\ar[d]
&H^{p+1}(\mathfrak U,\mathcal F')\ar[d]\\
H^{p+1}((\mathbf f),M'')\ar[r]_{\partial}
&H^{p+2}((\mathbf f),M')
}
\end{equation}
son conmutativos para todo $p$~(G, I, 2.1.1).
\end{remark}

\subsection{Cohomología de un esquema afín}
\label{subsection:III.1.3}
\label{subsection:III.1.3-fr}

\begin{theorem}[1.3.1]
\label{III.1.3.1}
\label{III.1.3.1-fr}
Sea $X$ un esquema afín.
Para todo $\mathcal O_X$-módulo cuasicoherente $\mathcal F$, se tiene $H^p(X,\mathcal F)=0$ para todo $p>0$.
\end{theorem}

\begin{proof}
Sea $\mathfrak U$ un recubrimiento finito de $X$ por los abiertos afines $X_{f_i}=D(f_i)$ ($1\leq i\leq r$); se sabe entonces que el ideal de $A=\Gamma(X,\mathcal O_X)$ generado por los $f_i$ es igual a $A$.
Del Corolario~\sref{III.1.2.4} se concluye, pues, que $H^p(\mathfrak U,\mathcal F)=0$ para $p>0$.
Como existen recubrimientos finitos de $X$ por abiertos afines arbitrariamente finos \sref[I]{I.1.1.10}, la definición de la cohomología de \v Cech (G, II, 5.8) muestra que también $\check H^p(X,\mathcal F)=0$ para $p>0$.
Pero esto se aplica asimismo a todo preesquema $X_f$, para $f\in A$ \sref[I]{I.1.3.6}; por tanto, $\check H^p(X_f,\mathcal F)=0$ para $p>0$.
Como $X_f\cap X_g=X_{fg}$, se deduce también que $H^p(X,\mathcal F)=0$ para todo $p>0$, en virtud de (G, II, 5.9.2).
\end{proof}

\begin{corollary}[1.3.2]
\label{III.1.3.2}
\label{III.1.3.2-fr}
Sea $Y$ un preesquema y sea $f:X\to Y$ un morfismo afín \sref[II]{II.1.6.1}.
Para todo $\mathcal O_X$-módulo cuasicoherente $\mathcal F$, se tiene $R^qf_*(\mathcal F)=0$ para $q>0$.
\end{corollary}

\begin{proof}
Por definición, $R^qf_*(\mathcal F)$ es el $\mathcal O_Y$-módulo asociado al prehaz $U\leadsto H^q(f^{-1}(U),\mathcal F)$, donde $U$ recorre los abiertos de $Y$.
Pero los abiertos afines forman una base de $Y$, y, para un abierto tal $U$, $f^{-1}(U)$ es afín \sref[II]{II.1.3.2}; por tanto, $H^q(f^{-1}(U),\mathcal F)=0$ por el Teorema~\sref{III.1.3.1}, lo que prueba el corolario.
\end{proof}

\begin{corollary}[1.3.3]
\label{III.1.3.3}
\label{III.1.3.3-fr}
Sea $Y$ un preesquema y sea $f:X\to Y$ un morfismo afín.
Para todo $\mathcal O_X$-módulo cuasicoherente $\mathcal F$, el homomorfismo canónico $H^p(Y,f_*(\mathcal F))\to H^p(X,\mathcal F)$ \textup{(\textbf{0},~12.1.3.1)} es biyectivo para todo $p$.
\end{corollary}

\begin{proof}
Basta (por \sref[0]{0.12.1.7}) mostrar que los homomorfismos de borde $''E_2^{p0}=H^p(Y,f_*(\mathcal F))\to H^p(X,\mathcal F)$ de la segunda sucesión espectral del funtor compuesto $\Gamma f_*$ son biyectivos.
Pero el término $E_2$ de esta sucesión espectral viene dado por $''E_2^{pq}=H^p(Y,R^qf_*(\mathcal F))$ (G, II, 4.17.1); por el Corolario~\sref{III.1.3.2}, $''E_2^{pq}=0$ para $q>0$, y la sucesión espectral degenera; de ahí nuestra aserción \sref[0]{0.11.1.6}.
\end{proof}

\begin{corollary}[1.3.4]
\label{III.1.3.4}
\label{III.1.3.4-fr}
Sean $f:X\to Y$ un morfismo afín y $g:Y\to Z$ un morfismo.
\oldpage[III]{89}
Para todo $\mathcal O_X$-módulo cuasicoherente $\mathcal F$, el homomorfismo canónico $R^pg_*(f_*(\mathcal F))\to R^p(g\circ f)_*(\mathcal F)$ \textup{(\textbf{0},~12.2.5.1)} es biyectivo para todo $p$.
\end{corollary}

\begin{proof}
Basta observar que, según el Corolario~\sref{III.1.3.3}, para todo abierto afín $W$ de $Z$, el homomorfismo canónico $H^p(g^{-1}(W),f_*(\mathcal F))\to
H^p(f^{-1}(g^{-1}(W)),\mathcal F)$ es biyectivo; esto prueba que el homomorfismo de prehaces que define el homomorfismo canónico $R^pg_*(f_*(\mathcal F))\to R^p(g\circ f)_*(\mathcal F)$ es biyectivo \sref[0]{0.12.2.5}.
\end{proof}

\subsection{Aplicación a la cohomología de preesquemas arbitrarios}
\label{subsection:III.1.4}
\label{subsection:III.1.4-fr}

\begin{proposition}[1.4.1]
\label{III.1.4.1}
\label{III.1.4.1-fr}
Sea $X$ un esquema y sea $\mathfrak U=(U_\alpha)$ un recubrimiento de $X$ por abiertos afines.
Para todo $\mathcal O_X$-módulo cuasicoherente $\mathcal F$, los módulos de cohomología $H^\bullet(X,\mathcal F)$ y $H^\bullet(\mathfrak U,\mathcal F)$ (sobre $\Gamma(X,\mathcal O_X)$) son canónicamente isomorfos.
\end{proposition}

\begin{proof}
Como $X$ es un esquema, toda intersección finita $V$ de abiertos del recubrimiento $\mathfrak U$ es afín \sref[I]{I.5.5.6}; por tanto, $H^q(V,\mathcal F)=0$ para $q\geq1$ por el Teorema~\sref{III.1.3.1}.
La proposición se deduce entonces de un teorema de Leray (G, II, 5.4.1).
\end{proof}

\begin{remark}[1.4.2]
\label{III.1.4.2}
\label{III.1.4.2-fr}
Se observará que el resultado de la Proposición~\sref{III.1.4.1} sigue siendo válido cuando las intersecciones finitas de los conjuntos $U_\alpha$ son afines, incluso si no se supone necesariamente que $X$ sea un esquema.
\end{remark}

\begin{corollary}[1.4.3]
\label{III.1.4.3}
\label{III.1.4.3-fr}
Sea $X$ un esquema de espacio subyacente cuasicompacto, sea $A=\Gamma(X,\mathcal O_X)$ y sea $\mathbf f=(f_i)_{1\leq i\leq r}$ una sucesión finita de elementos de $A$ tal que los $X_{f_i}$ (notación de \sref{III.1.2.1}) sean afines.
Entonces (con las notaciones de \sref{III.1.2.1}), para todo $\mathcal O_X$-módulo cuasicoherente $\mathcal F$, se tiene un isomorfismo canónico funtorial en $\mathcal F$
\begin{equation}
\label{III.1.4.3.1-fr}
\tag{1.4.3.1}
H^q(U,\mathcal F)\xrightarrow{\sim}H^{q+1}((\mathbf f),M)
\qquad\text{para }q\geq1
\end{equation}
y una sucesión exacta funtorial en $\mathcal F$
\begin{equation}
\label{III.1.4.3.2-fr}
\tag{1.4.3.2}
0\longrightarrow H^0((\mathbf f),M)\longrightarrow M
\longrightarrow H^0(U,\mathcal F)
\longrightarrow H^1((\mathbf f),M)\longrightarrow0.
\end{equation}
\end{corollary}

\begin{proof}
Esto se deduce inmediatamente de las Proposiciones~\sref{III.1.4.1} y \sref{III.1.2.3}.
\end{proof}

\begin{env}[1.4.4]
\label{III.1.4.4}
\label{III.1.4.4-fr}
Si $X$ es un \emph{esquema afín}, de la Observación~\sref{III.1.2.5} y de la Proposición~\sref{III.1.4.1} se deduce que, para todo $q\geq0$, los diagramas
\begin{equation}
\label{III.1.4.4.1-fr}
\tag{1.4.4.1}
\xymatrix{
H^q(U,\mathcal F'')\ar[r]^{\partial}\ar[d]
&H^{q+1}(U,\mathcal F')\ar[d]\\
H^{q+1}((\mathbf f),M'')\ar[r]_{\partial}
&H^{q+2}((\mathbf f),M')
}
\end{equation}
correspondientes a una sucesión exacta $0\to\mathcal F'\to\mathcal F\to\mathcal F''\to0$ de $\mathcal O_X$-módulos cuasicoherentes (con las notaciones de la Observación~\sref{III.1.2.5}) son conmutativos.
\end{env}

\begin{proposition}[1.4.5]
\label{III.1.4.5}
\label{III.1.4.5-fr}
\oldpage[III]{90}
Sea $X$ un esquema cuasicompacto, sea $\mathcal L$ un $\mathcal O_X$-módulo invertible, y consideremos el anillo graduado $A_*=\Gamma_*(\mathcal L)$ \sref[0]{0.5.4.6}; entonces $H^\bullet(\mathcal F,\mathcal L)
=\bigoplus_{n\in\mathbb Z}H^\bullet(X,\mathcal F\otimes\mathcal L^{\otimes n})$ es un $A_*$-módulo graduado y, para todo $f\in A_n$, se tiene un isomorfismo canónico
\begin{equation}
\label{III.1.4.5.1-fr}
\tag{1.4.5.1}
H^\bullet(X_f,\mathcal F)
\xrightarrow{\sim}(H^\bullet(\mathcal F,\mathcal L))_{(f)}
\end{equation}
de $(A_*)_{(f)}$-módulos.
\end{proposition}

\begin{proof}
Como $X$ es un esquema cuasicompacto, puede calcularse la cohomología de todos los $\mathcal O_X$-módulos $\mathcal F\otimes\mathcal L^{\otimes n}$ utilizando un mismo recubrimiento finito $\mathfrak U=(U_i)$ por abiertos afines tales que la restricción
$\mathcal L\vert U_i$ es isomorfa a $\mathcal O_X\vert U_i$ para cada
$i$ \sref{III.1.4.1}.
Es entonces inmediato que los $U_i\cap X_f$ son abiertos afines \sref[I]{I.1.3.6}, y puede calcularse así la cohomología $H^\bullet(X_f,\mathcal F\otimes\mathcal L^{\otimes n})$ utilizando el recubrimiento $\mathfrak U\vert X_f=(U_i\cap X_f)$ \sref{III.1.4.1}.
Es inmediato que, para todo $f\in A_n$, la multiplicación por $f$ define un homomorfismo $C^\bullet(\mathfrak U,\mathcal F\otimes\mathcal L^{\otimes m})
\longrightarrow
C^\bullet(\mathfrak U,\mathcal F\otimes\mathcal L^{\otimes(m+n)})$ y, por tanto, un homomorfismo $H^\bullet(\mathfrak U,\mathcal F\otimes\mathcal L^{\otimes m})
\longrightarrow
H^\bullet(\mathfrak U,\mathcal F\otimes\mathcal L^{\otimes(m+n)})$, lo que establece la primera afirmación.
Por otra parte, para un $f\in A_n$ dado, de \sref[I]{I.9.3.2} se deduce que se tiene un isomorfismo de complejos de $(A_*)_{(f)}$-módulos
\[
C^\bullet(\mathfrak U\vert X_f,\mathcal F)
\xrightarrow{\sim}
\left(C^\bullet\left(\mathfrak U,
\bigoplus_{n\in\mathbb Z}\mathcal F\otimes\mathcal L^{\otimes n}
\right)\right)_{(f)}
\]
teniendo en cuenta \sref[I]{I.1.3.9}[ii].
Al pasar a la cohomología de estos complejos se obtiene el isomorfismo (\hyperref[III.1.4.5.1]{1.4.5.1}), recordando que el funtor $M\leadsto M_{(f)}$ es exacto sobre la categoría de $A_*$-módulos graduados.
\end{proof}

\begin{corollary}[1.4.6]
\label{III.1.4.6}
\label{III.1.4.6-fr}
Se supone que se satisfacen las hipótesis de la Proposición~\textup{\sref{III.1.4.5}} y se supone además que $\mathcal L=\mathcal O_X$.
Si se pone $A=\Gamma(X,\mathcal O_X)$, entonces, para todo $f\in A$, se tiene un isomorfismo canónico $H^\bullet(X_f,\mathcal F)
\xrightarrow{\sim}(H^\bullet(X,\mathcal F))_f$ de $A_f$-módulos.
\end{corollary}

\begin{corollary}[1.4.7]
\label{III.1.4.7}
\label{III.1.4.7-fr}
Sea $X$ un esquema cuasicompacto y sea $f$ un elemento de $\Gamma(X,\mathcal O_X)$.
\begin{enumerate}
  \item[{\rm(i)}] Supongamos que el abierto $X_f$ es afín.
    Entonces, para todo $\mathcal O_X$-módulo cuasicoherente $\mathcal F$, todo $i>0$ y todo $\xi\in H^i(X,\mathcal F)$, existe un entero $n>0$ tal que $f^n\xi=0$.
  \item[{\rm(ii)}] Recíprocamente, supongamos que $X_f$ es cuasicompacto y que, para todo haz cuasicoherente de ideales $\mathcal J$ de $\mathcal O_X$ y todo $\zeta\in H^1(X,\mathcal J)$, existe un $n>0$ tal que $f^n\zeta=0$.
    Entonces $X_f$ es afín.
\end{enumerate}
\end{corollary}

\begin{proof}
\medskip\noindent
\begin{enumerate}
  \item[(i)] Si $X_f$ es afín, se tiene $H^i(X_f,\mathcal F)=0$ para todo $i>0$ \sref{III.1.3.1}, de modo que la afirmación se deduce directamente del Corolario~\sref{III.1.4.6}.
  \item[(ii)] En virtud del criterio de Serre \sref[II]{II.5.2.1}, basta probar que, para todo haz cuasicoherente de ideales $\mathcal K$ de $\mathcal O_X\vert X_f$, se tiene $H^1(X_f,\mathcal K)=0$.
    Como $X_f$ es un abierto cuasicompacto de un esquema cuasicompacto $X$, existe un haz cuasicoherente de ideales $\mathcal J$ de $\mathcal O_X$ tal que $\mathcal K=\mathcal J\vert X_f$ \sref[I]{I.9.4.2}.
    Según el Corolario~\sref{III.1.4.6}, se tiene $H^1(X_f,\mathcal K)=(H^1(X,\mathcal J))_f$, y la hipótesis implica que el miembro de la derecha es nulo, lo que da la afirmación.
\end{enumerate}
\end{proof}

\begin{remark}[1.4.8]
\label{III.1.4.8}
\label{III.1.4.8-fr}
Se observará que el Corolario~\sref{III.1.4.7}[i] da una demostración más sencilla de la relación (\textbf{II},~4.5.13.2).
\end{remark}

\begin{lemma}[1.4.9]
\label{III.1.4.9}
\label{III.1.4.9-fr}
\oldpage[III]{91}
Sea $X$ un esquema cuasicompacto, sea $\mathfrak U=(U_i)_{1\leq i\leq n}$ un recubrimiento finito de $X$ por abiertos afines y sea $\mathcal F$ un $\mathcal O_X$-módulo cuasicoherente.
El complejo de haces $\mathcal C^\bullet(\mathfrak U,\mathcal F)$ definido por el recubrimiento $\mathfrak U$ \textup{(G, II, 5.2)} es entonces un $\mathcal O_X$-módulo cuasicoherente.
\end{lemma}

\begin{proof}
De las definiciones (G, II, 5.2) se deduce que $\mathcal C^p(\mathfrak U,\mathcal F)$ es la suma directa de los haces imagen directa de los $\mathcal F\vert U_{i_0\ldots i_p}$ por la inyección canónica $U_{i_0\ldots i_p}\to X$.
La hipótesis de que $X$ es un esquema implica que estas inyecciones son morfismos afines \sref[I]{I.5.5.6}; por tanto, los $\mathcal C^p(\mathfrak U,\mathcal F)$ son cuasicoherentes \sref[II]{II.1.2.6}.
\end{proof}

\begin{proposition}[1.4.10]
\label{III.1.4.10}
\label{III.1.4.10-fr}
Sea $u:X\to Y$ un morfismo separado y cuasicompacto.
Para todo $\mathcal O_X$-módulo cuasicoherente $\mathcal F$, los $R^qu_*(\mathcal F)$ son $\mathcal O_Y$-módulos cuasicoherentes.
\end{proposition}

\begin{proof}
La cuestión es local sobre $Y$, de modo que puede suponerse que $Y$ es afín.
Entonces $X$ es una unión finita de abiertos afines $U_i$ ($1\leq i\leq n$); sea $\mathfrak U$ el recubrimiento $(U_i)$.
Además, como $Y$ es un esquema, de \sref[I]{I.5.5.10} se deduce que, para todo abierto afín $V\subset Y$, la inyección canónica $u^{-1}(V)\to X$ es un morfismo afín; se concluye (Proposición~\sref{III.1.4.1} y (G, II, 5.2)) que se tiene un isomorfismo canónico
\begin{equation}
\label{III.1.4.10.1-fr}
\tag{1.4.10.1}
H^\bullet(u^{-1}(V),\mathcal F)
\xrightarrow{\sim}H^\bullet(\Gamma(V,\mathcal K^\bullet))
\end{equation}
donde se pone $\mathcal K^\bullet=u_*(\mathcal C^\bullet(\mathfrak U,\mathcal F))$.
Según el Lema~\sref{III.1.4.9} y \sref[I]{I.9.2.2}, $\mathcal K^\bullet$ es un $\mathcal O_Y$-módulo cuasicoherente; además, constituye un \emph{complejo de haces}, puesto que también lo es $\mathcal C^\bullet(\mathfrak U,\mathcal F)$.
De la definición de la cohomología $\mathcal H^\bullet(\mathcal K^\bullet)$ (G, II, 4.1) se deduce entonces que esta última está formada por $\mathcal O_Y$-módulos cuasicoherentes \sref[I]{I.4.1.1}.
Como (para $V$ afín en $Y$) el funtor $\Gamma(V,\mathcal G)$ es exacto en $\mathcal G$ sobre la categoría de $\mathcal O_Y$-módulos cuasicoherentes, se tiene (G, II, 4.1)
\begin{equation}
\label{III.1.4.10.2-fr}
\tag{1.4.10.2}
H^\bullet(\Gamma(V,\mathcal K^\bullet))
=\Gamma(V,\mathcal H^\bullet(\mathcal K^\bullet)).
\end{equation}
Observemos finalmente que resulta de la definición del homomorfismo canónico \[
H^\bullet(\mathfrak U,\mathcal F)\longrightarrow H^\bullet(X,\mathcal F)
\]
dado en (G, II, 5.2), se deduce que, si $V'\subset V$ es otro abierto afín de $Y$, el
diagrama
\[
\xymatrix{
H^\bullet(u^{-1}(V),\mathcal F)\ar[r]^{\sim}\ar[d]
&H^\bullet(\Gamma(V,\mathcal K^\bullet))\ar[d]\\
H^\bullet(u^{-1}(V'),\mathcal F)\ar[r]^{\sim}
&H^\bullet(\Gamma(V',\mathcal K^\bullet))
}
\]
es conmutativo.
De lo anterior se concluye que los isomorfismos (\hyperref[III.1.4.10.1]{1.4.10.1}) definen un isomorfismo de $\mathcal O_Y$-módulos
\begin{equation}
\label{III.1.4.10.3-fr}
\tag{1.4.10.3}
R^\bullet u_*(\mathcal F)
\xrightarrow{\sim}\mathcal H^\bullet(\mathcal K^\bullet)
\end{equation}
\oldpage[III]{92}
y, por consiguiente, $R^\bullet u_*(\mathcal F)$ es cuasicoherente.
\end{proof}

Además, de (\hyperref[III.1.4.10.3]{1.4.10.3}), (\hyperref[III.1.4.10.2]{1.4.10.2}) y (\hyperref[III.1.4.10.1]{1.4.10.1}) se deduce que:
\begin{corollary}[1.4.11]
\label{III.1.4.11}
\label{III.1.4.11-fr}
Bajo las hipótesis de la Proposición~\sref{III.1.4.10}, para todo abierto afín $V$ de $Y$, el homomorfismo canónico
\begin{equation}
\label{III.1.4.11.1-fr}
\tag{1.4.11.1}
H^q(u^{-1}(V),\mathcal F)\longrightarrow
\Gamma(V,R^qu_*(\mathcal F))
\end{equation}
es un isomorfismo para todo $q\geq0$.
\end{corollary}

\begin{corollary}[1.4.12]
\label{III.1.4.12-fr}
Supongamos verificadas las hipótesis de \textup{(\hyperref[III.1.4.10-fr]{1.4.10})} y supongamos además que $Y$ es cuasicompacto.
Entonces existe un entero $r>0$ tal que, para todo $\mathcal O_X$-módulo cuasicoherente $\mathcal F$ y todo entero $q>r$, se tiene $R^qu_*(\mathcal F)=0$.
Si $Y$ es afín, puede tomarse como $r$ un entero tal que exista un recubrimiento de $X$ formado por $r$ abiertos afines.
\end{corollary}

\begin{proof}
Como puede recubrirse $Y$ por un número finito de abiertos afines, en virtud del Corolario~\sref{III.1.4.11} puede reducirse la demostración a la segunda afirmación.
Si $\mathfrak U$ es un recubrimiento de $X$ por $r$ abiertos afines, se tiene $H^q(\mathfrak U,\mathcal F)=0$ para $q>r$, puesto que las cocadenas de $C^q(\mathfrak U,\mathcal F)$ son alternadas; la afirmación se deduce entonces de la Proposición~\sref{III.1.4.1}.
\end{proof}

\begin{corollary}[1.4.13]
\label{III.1.4.13-fr}
Supongamos verificadas las hipótesis de \textup{(\hyperref[III.1.4.10-fr]{1.4.10})} y supongamos además que $Y=\operatorname{Spec}(A)$ es afín.
Entonces, para todo $\mathcal O_X$-módulo cuasicoherente $\mathcal F$ y todo $f\in A$, se tiene
\[
\Gamma(Y_f,R^qu_*(\mathcal F))
=(\Gamma(Y,R^qu_*(\mathcal F)))_f
\]
salvo isomorfismo canónico.
\end{corollary}

\begin{proof}
Esto se deduce del hecho de que $R^qu_*(\mathcal F)$ es un $\mathcal O_Y$-módulo cuasicoherente \sref[I]{I.1.3.7}.
\end{proof}

\begin{proposition}[1.4.14]
\label{III.1.4.14}
\label{III.1.4.14-fr}
Sean $f:X\to Y$ un morfismo separado y cuasicompacto y $g:Y\to Z$ un morfismo afín.
Para todo $\mathcal O_X$-módulo cuasicoherente $\mathcal F$, el homomorfismo canónico $R^p(g\circ f)_*(\mathcal F)\to g_*(R^pf_*(\mathcal F))$ \textup{(\textbf{0},~12.2.5.2)} es biyectivo para todo $p$.
\end{proposition}

\begin{proof}
Para todo abierto afín $W$ de $Z$, $g^{-1}(W)$ es un abierto afín de $Y$.
El homomorfismo de prehaces que define el homomorfismo canónico
\[
R^p(g\circ f)_*(\mathcal F)\longrightarrow g_*(R^pf_*(\mathcal F))
\]
\sref[0]{0.12.2.5} es, por tanto, biyectivo en virtud del Corolario~\sref{III.1.4.11}.
\end{proof}

\begin{proposition}[1.4.15]
\label{III.1.4.15}
\label{III.1.4.15-fr}
Sean $u:X\to Y$ un morfismo separado de tipo finito y $v:Y'\to Y$ un morfismo plano de preesquemas \sref[0]{0.6.7.1}; sea $u'=u_{(Y')}$, de modo que se tiene el diagrama conmutativo
\begin{equation}
\label{III.1.4.15.1-fr}
\tag{1.4.15.1}
\xymatrix{
X\ar[d]_u & X'=X_{(Y')}\ar[l]_{v'}\ar[d]^{u'}\\
Y & Y'\ar[l]^v
}\quad.
\end{equation}
Entonces, para todo $\mathcal O_X$-módulo cuasicoherente $\mathcal F$, $R^qu_*'(\mathcal F')$, donde
$\mathcal F'=v'^*(\mathcal F)=\mathcal F\otimes_{\mathcal O_Y}\mathcal O_{Y'}$,
es canónicamente isomorfo a
$R^qu_*(\mathcal F)\otimes_{\mathcal O_Y}\mathcal O_{Y'}
=v^*(R^qu_*(\mathcal F))$ para todo $q\geq0$.
\end{proposition}

\oldpage[III]{93}
\begin{proof}
El homomorfismo canónico
$\rho:\mathcal F\to v_*'(v'^*(\mathcal F))$ (\textbf{0},~4.4.3.2) define por funtorialidad un homomorfismo
\begin{equation}
\label{III.1.4.15.2-fr}
\tag{1.4.15.2}
R^qu_*(\mathcal F)\longrightarrow R^qu_*(v_*'(\mathcal F')).
\end{equation}

Por otra parte, poniendo $w=u\circ v'=v\circ u'$, se tienen los homomorfismos canónicos (\textbf{0}~,12.2.5.1 y 12.2.5.2)
\begin{equation}
\label{III.1.4.15.3-fr}
\tag{1.4.15.3}
R^qu_*(v_*'(\mathcal F'))\longrightarrow R^qw_*(\mathcal F')
\longrightarrow v_*(R^qu_*'(\mathcal F')).
\end{equation}

Componiendo (\hyperref[III.1.4.15.3]{1.4.15.3}) y (\hyperref[III.1.4.15.2]{1.4.15.2}), se obtiene un homomorfismo
\[
\psi:R^qu_*(\mathcal F)\longrightarrow v_*(R^qu_*'(\mathcal F'))
\]
y se obtiene finalmente un homomorfismo canónico (cuya definición no impone \emph{hipótesis alguna} sobre $v$)
\begin{equation}
\label{III.1.4.15.4-fr}
\tag{1.4.15.4}
\psi^\sharp:v^*(R^qu_*(\mathcal F))\longrightarrow R^qu_*'(\mathcal F')
\end{equation}
y hay que probar que es un isomorfismo cuando $v$ es \emph{plano}.
Es claro que la cuestión es local sobre $Y$ y $Y'$; puede suponerse, por tanto, que $Y=\operatorname{Spec}(A)$ y $Y'=\operatorname{Spec}(B)$; utilizaremos además el
\begin{lemma}[1.4.15.5]
\label{III.1.4.15.5-fr}
Sean $\varphi:A\to B$ un homomorfismo de anillos, $Y=\operatorname{Spec}(A)$, $X=\operatorname{Spec}(B)$, $f:X\to Y$ el morfismo correspondiente a $\varphi$, y $M$ un $B$-módulo.
Para que el $\mathcal O_X$-módulo $\widetilde M$ sea $f$-plano \sref[0]{0.6.7.1}, es necesario y suficiente que $M$ sea un $A$-módulo plano.
En particular, para que el morfismo $f$ sea plano, es necesario y suficiente que $B$ sea un $A$-módulo plano.
\end{lemma}

Esto se deduce de la definición \sref[0]{0.6.7.1} y de \sref[0]{0.6.3.3}, teniendo en cuenta \sref[I]{I.1.3.4}.

Así pues, de (\hyperref[III.1.4.11.1]{1.4.11.1}) y de las definiciones de los homomorfismos (\hyperref[III.1.4.15.3]{1.4.15.3}) (cf.~\sref[0]{0.12.2.5}) se deduce que $\psi$ corresponde entonces al morfismo compuesto
\[
H^q(X,\mathcal F)\xrightarrow{\rho_q}
H^q(X,v_*'(v'^*(\mathcal F)))\xrightarrow{\theta_q}
H^q(X',v'^*(v_*'(v'^*(\mathcal F))))\xrightarrow{\sigma_q}
H^q(X',v'^*(\mathcal F))
\]
donde $\rho_q$ y $\sigma_q$ son los homomorfismos de cohomología correspondientes a los morfismos canónicos $\rho$ y $\sigma:v'^*(v_*'(\mathcal G'))\to\mathcal G'$, y $\theta_q$ es el $\varphi$-morfismo (\textbf{0},~12.1.3.1) relativo al $\mathcal O_X$-módulo $v_*'(v'^*(\mathcal F))$.
Pero, por la funtorialidad de $\theta_q$, se tiene el diagrama conmutativo
\[
\xymatrix{
H^q(X,\mathcal F)\ar[rr]^{\rho_q}\ar[dd]_{\theta_q}
&&H^q(X,v_*'(v'^*(\mathcal F)))\ar[dd]^{\theta_q}\\\\
H^q(X',v'^*(\mathcal F))\ar[rr]_{v'^*(\rho_q)}
&&H^q(X',v'^*(v_*'(v'^*(\mathcal F))))
}
\]
\oldpage[III]{94}
y, como por definición \sref[0]{0.4.4.3} $v'^*(\rho)$ es el inverso de $\sigma$, se ve que el morfismo compuesto considerado anteriormente no es finalmente otro que $\theta_q$; por consiguiente, $\psi^\sharp$ es el $B$-homomorfismo asociado $H^q(X,\mathcal F)\otimes_A B\to H^q(X',\mathcal F')$.
Como $u$ es de tipo finito, $X$ es una unión finita de abiertos afines $U_i$ ($1\leq i\leq r$); sea $\mathfrak U$ el recubrimiento $(U_i)$.
Como $v$ es un morfismo afín, también lo es $v'$ \sref[II]{II.1.6.2}[iii]; por consiguiente, los $U_i'=v'^{-1}(U_i)$ forman un recubrimiento abierto afín $\mathfrak U'$ de $X'$.
Se sabe entonces \sref[0]{0.12.1.4.2} que el diagrama
\[
\xymatrix{
H^q(\mathfrak U,\mathcal F)\ar[r]^{\theta_q}\ar[d]
&H^q(\mathfrak U',\mathcal F')\ar[d]\\
H^q(X,\mathcal F)\ar[r]_{\theta_q}&H^q(X',\mathcal F')
}
\]
es conmutativo, y las flechas verticales son isomorfismos puesto que $X$ y $X'$ son esquemas \sref{III.1.4.1}.
Por consiguiente, basta probar que el $\varphi$-morfismo canónico $\theta_q:H^q(\mathfrak U,\mathcal F)\to H^q(\mathfrak U',\mathcal F')$ es tal que el $B$-homomorfismo asociado
\[
H^q(\mathfrak U,\mathcal F)\otimes_A B
\longrightarrow H^q(\mathfrak U',\mathcal F')
\]
sea un isomorfismo.
Para toda sucesión $\mathbf s=(i_k)_{0\leq k\leq p}$ de $p+1$ índices de $[1,r]$, pongamos $U_{\mathbf s}=\bigcap\limits_{k=0}^pU_{i_k}$, $U_{\mathbf s}'=\bigcap\limits_{k=0}^{\boldsymbol{p}}U_{i_k}'=v'^{-1}(U_{\mathbf s})$, $M_{\mathbf s}=\Gamma(U_{\mathbf s},\mathcal F)$ y $M_{\mathbf s}'=\Gamma(U_{\mathbf s}',\mathcal F')$.
La aplicación canónica $M_{\mathbf s}\otimes_A B\to M_{\mathbf s}'$ es un isomorfismo \sref[I]{I.1.6.5}; por tanto, la aplicación canónica $C^p(\mathfrak U,\mathcal F)\otimes_A B\to
C^p(\mathfrak U',\mathcal F')$ es un isomorfismo, mediante el cual $d\otimes1$ se identifica con la aplicación coborde $C^p(\mathfrak U',\mathcal F')\to C^{p+1}(\mathfrak U',\mathcal F')$.
Como $B$ es un $A$-módulo \emph{plano}, de la definición de los módulos de cohomología se deduce que la aplicación canónica $H^q(\mathfrak U,\mathcal F)\otimes_A B\to
H^q(\mathfrak U',\mathcal F')$ es un isomorfismo \sref[0]{0.6.1.1}.
Este resultado se generalizará más adelante en \EGAUnavailableHyperref{section:III.6}{\textsection6}.
\end{proof}

\begin{corollary}[1.4.16]
\label{III.1.4.16}
\label{III.1.4.16-fr}
Sean $A$ un anillo, $X$ un $A$-esquema de tipo finito y $B$ una $A$-álgebra fielmente plana sobre $A$.
Para que $X$ sea afín, es necesario y suficiente que $X\otimes_A B$ lo sea.
\end{corollary}

\begin{proof}
La condición es evidentemente necesaria \sref[I]{I.3.2.2}; mostremos que es suficiente.
Como $X$ es separado sobre $A$ y el morfismo $\operatorname{Spec}(B)\to\operatorname{Spec}(A)$ es plano, de la Proposición~\sref{III.1.4.1} se deduce que
\begin{equation}
\label{III.1.4.16.1-fr}
\tag{1.4.16.1}
H^i(X\otimes_A B,\mathcal F\otimes_A B)
=H^i(X,\mathcal F)\otimes_A B
\end{equation}
para todo $i\geq0$ y todo $\mathcal O_X$-módulo cuasicoherente $\mathcal F$.
Si $X\otimes_A B$ es afín, el miembro izquierdo de (\hyperref[III.1.4.16.1]{1.4.16.1}) es nulo para $i=1$; por tanto, también lo es $H^1(X,\mathcal F)$, puesto que $B$ es un $A$-módulo fielmente plano.
Como $X$ es un esquema cuasicompacto, la demostración termina por el criterio de Serre \sref[II]{II.5.2.1}.
\end{proof}

\begin{proposition}[1.4.17]
\label{III.1.4.17}
\label{III.1.4.17-fr}
Sea $X$ un preesquema y sea $0\to\mathcal F\xrightarrow{u}\mathcal G\xrightarrow{v}\mathcal H\to0$ una sucesión exacta de $\mathcal O_X$-módulos.
Si $\mathcal F$ y $\mathcal H$ son cuasicoherentes, también lo es $\mathcal G$.
\end{proposition}

\oldpage[III]{95}
\begin{proof}
La cuestión es local sobre $X$, de modo que puede suponerse que $X=\operatorname{Spec}(A)$ es afín; basta entonces probar que $\mathcal G$ satisface las condiciones (\textit{d}1) y (\textit{d}2) de \sref[I]{I.1.4.1} (con $V=X$).
La verificación de (\textit{d}2) es inmediata, pues, si $t\in\Gamma(X,\mathcal G)$ se anula al restringirlo a $D(f)$, lo mismo ocurre con su imagen $v(t)\in\Gamma(X,\mathcal H)$; por tanto, existe un $m>0$ tal que $f^mv(t)=v(f^mt)=0$ \sref[I]{I.1.4.1}, y, como $\Gamma$ es exacto por la izquierda, $f^mt=u(s)$, donde $s\in\Gamma(X,\mathcal F)$; como $u$ es inyectivo, la restricción de $s$ a $D(f)$ es nula, luego por \sref[I]{I.1.4.1} existe un entero $n>0$ tal que $f^ns=0$; se deduce finalmente que $f^{m+n}t=u(f^ns)=0$.

Compruébese ahora (\textit{d}2); sea $t'\in\Gamma(D(f),\mathcal G)$; como $\mathcal H$ es cuasicoherente, existe un entero $m$ tal que $f^mv(t')=v(f^mt')$ se prolonga a una sección $z\in\Gamma(X,\mathcal H)$ \sref[I]{I.1.4.1}.
Pero, en virtud del Teorema~\sref{III.1.3.1} (o de (\textbf{I},~5.1.9.2)) aplicado al $\mathcal O_X$-módulo cuasicoherente $\mathcal F$, la sucesión $\Gamma(X,\mathcal G)\to\Gamma(X,\mathcal H)\to0$ es exacta; existe, pues, $t\in\Gamma(X,\mathcal G)$ tal que $z=v(t)$; se ve así que $v(f^mt'\setmin t'')=0$, denotando por $t''$ la restricción de $t$ a $D(f)$; por tanto, $f^mt'\setmin t''=u(s')$, donde $s'\in\Gamma(D(f),\mathcal F)$.
Pero, como $\mathcal F$ es cuasicoherente, existe un entero $n>0$ tal que $f^ns'$ se prolonga a una sección $s\in\Gamma(X,\mathcal F)$; como $f^{m+n}t'\setmin f^nt''=u(f^ns')$, se ve que $f^{m+n}t'$ es la restricción a $D(f)$ de una sección $f^nt+u(f^ns)\in\Gamma(X,\mathcal G)$, lo que termina la demostración.
\end{proof}

% END tree/repo/editions/III1_DIPLOMATIC/ega3/ega3-1.tex


% BEGIN tree/repo/editions/III1_DIPLOMATIC/ega3/ega3-2.tex
\section{Estudio cohomológico de los morfismos proyectivos}
\label{section:III.2}
\label{section:III.2-fr}


\subsection{Cálculos explícitos de ciertos grupos de cohomología}
\label{subsection:III.2.1}
\label{subsection:III.2.1-fr}

\begin{env}[2.1.1]
\label{III.2.1.1}
\label{III.2.1.1-fr}

Sea $X$ un preesquema y sea $\mathcal L$ un $\mathcal O_X$-módulo invertible; consideremos el anillo graduado \sref[0]{0.5.4.6}
\begin{equation}
\label{III.2.1.1.1-fr}
\tag{2.1.1.1}
S=\Gamma_*(X,\mathcal L)
=\bigoplus_{n\in\mathbb Z}\Gamma(X,\mathcal L^{\otimes n}).
\end{equation}

Sea $(f_i)_{1\leq i\leq r}$ una familia finita de elementos \emph{homogéneos} de $S$, sea $f_i\in S_{d_i}$; pongamos $U_i=X_{f_i}$, $U=\bigcup_iU_i$, y denotemos por $\mathfrak U$ el recubrimiento $(U_i)$ de $U$.
Para todo $\mathcal O_X$-módulo cuasicoherente $\mathcal F$, se pondrá
\begin{align}
\label{III.2.1.1.2-fr}
\tag{2.1.1.2}
H^\bullet(U,\mathcal F(*))
&=\bigoplus_{n\in\mathbb Z}
H^\bullet(U,\mathcal F\otimes\mathcal L^{\otimes n})\\
\label{III.2.1.1.3-fr}
\tag{2.1.1.3}
H^\bullet(\mathfrak U,\mathcal F(*))
&=\bigoplus_{n\in\mathbb Z}
H^\bullet(\mathfrak U,\mathcal F\otimes\mathcal L^{\otimes n}).
\end{align}

Los grupos abelianos (\hyperref[III.2.1.1.2-fr]{2.1.1.2}) y
(\hyperref[III.2.1.1.3-fr]{2.1.1.3}) son \emph{bigraduados}, tomando
\[
(H^\bullet(U,\mathcal F(*)))_{mn}
=H^m(U,\mathcal F\otimes\mathcal L^{\otimes n})
\]
y una definición análoga para (\hyperref[III.2.1.1.3]{2.1.1.3}).
Para la segunda graduación, está claro que estos grupos son $S$-módulos graduados, como se deduce, por ejemplo, del hecho de que $\mathcal F\leadsto H^m(U,\mathcal F)$ y $\mathcal F\leadsto H^m(\mathfrak U,\mathcal F)$ son funtores.
\end{env}

\begin{env}[2.1.2]
\label{III.2.1.2}
\label{III.2.1.2-fr}
Consideremos ahora el $S$-módulo graduado \sref[0]{0.5.4.6}
\begin{equation}
\label{III.2.1.2.1-fr}
\tag{2.1.2.1}
M=\Gamma_*(\mathcal L,\mathcal F)
=H^0(X,\mathcal F(*))
=\bigoplus_{n\in\mathbb Z}
\Gamma(X,\mathcal F\otimes\mathcal L^{\otimes n}).
\end{equation}

\oldpage[III]{96}
Si $X$ es un preesquema cuyo espacio subyacente es noetheriano, o un esquema cuasicompacto, se deduce de \sref[I]{I.9.3.1} que, poniendo como de costumbre $U_{i_0i_1\ldots i_p}=\bigcap_{k=0}^pU_{i_k}$, se tiene, salvo isomorfismo canónico,
\[
\Gamma(U_{i_0i_1\ldots i_p},\mathcal F(*))
=H^0(U_{i_0i_1\ldots i_p},\mathcal F(*))
=M_{f_{i_0}f_{i_1}\ldots f_{i_p}}.
\]

Además, con la notación de (\hyperref[III.1.2.2]{1.2.2}), se puede identificar $M_{f_{i_0}f_{i_1}\ldots f_{i_p}}$ con $\varinjlim_n M_{i_0i_1\ldots i_p}^{(n)}$.
Esta identificación es un isomorfismo de $S$-módulos \emph{graduados} si se define el grado de un elemento homogéneo $z\in\varinjlim_n M_{i_0i_1\ldots i_p}^{(n)}$ como sigue: $z$ es la imagen canónica de un elemento homogéneo $x\in M_{i_0i_1\ldots i_p}^{(n)}=M$ de grado $m$, y se toma como grado de $z$ el número $m-n(d_{i_0}+d_{i_1}+\ldots+d_{i_p})$.
Teniendo en cuenta la definición de los homomorfismos $\varphi_{kh}:M_{i_0i_1\ldots i_p}^{(h)}\to
M_{i_0i_1\ldots i_p}^{(k)}$ (\hyperref[III.1.2.2]{1.2.2}), se ve inmediatamente que esta definición no depende del ``representante'' $x$ de $z$ elegido.
Denotando, como en (\hyperref[III.1.2.2]{1.2.2}), por $C_n^p(M)$ el conjunto de aplicaciones alternadas de $[1,r]^{p+1}$ en $M$ (para todo $n$), se define de manera análoga una estructura de $S$-módulo \emph{graduado} sobre $\varinjlim_n C_n^p(M)$; como en (\hyperref[III.1.2.2]{1.2.2}), se tiene entonces
\begin{equation}
\label{III.2.1.2.2-fr}
\tag{2.1.2.2}
C^p(\mathfrak U,\mathcal F(*))=\varinjlim_n C_n^p(M)
\end{equation}
donde el isomorfismo entre ambos miembros \emph{respeta los grados}.
Se tiene también, como en (\hyperref[III.1.2.2]{1.2.2}),
\begin{equation}
\label{III.2.1.2.3-fr}
\tag{2.1.2.3}
C^p(\mathfrak U,\mathcal F(*))
=C^{p+1}((\mathbf f),M)
=\varinjlim_n K^{p+1}(\mathbf f^n,M)
\end{equation}
donde el isomorfismo \emph{conserva los grados}: el grado de un elemento de $\varinjlim_n K^{p+1}(\mathbf f^n,M)$ que sea la imagen canónica de una cocadena $g\in K^{p+1}(\mathbf f^n,M)$ cuyos valores $g(i_0,\ldots,i_p)$ pertenezcan a una \emph{misma} componente homogénea $M_m$ de $M$ es $m-n(d_{i_0}+\ldots+d_{i_p})$, independientemente de la elección de esta cocadena como representante del elemento considerado.

Como los isomorfismos anteriores son compatibles con los operadores de coborde, se concluye, como en (\hyperref[III.1.2.2]{1.2.2}), que:
\end{env}

\begin{proposition}[2.1.3]
\label{III.2.1.3}
\label{III.2.1.3-fr}
Sea $X$ un preesquema cuyo espacio subyacente es noetheriano, o un esquema cuasicompacto.
Existe un isomorfismo canónico, funtorial en $\mathcal F$,
\begin{equation}
\label{III.2.1.3.1-fr}
\tag{2.1.3.1}
H^p(\mathfrak U,\mathcal F(*))
\xrightarrow{\sim}H^{p+1}((\mathbf f),M)
\qquad\text{para todo }p\geq1.
\end{equation}
Además, existe una sucesión exacta, funtorial en $\mathcal F$,
\begin{equation}
\label{III.2.1.3.2-fr}
\tag{2.1.3.2}
0\longrightarrow H^0((\mathbf f),M)\longrightarrow M
\longrightarrow H^0(\mathfrak U,\mathcal F(*))
\longrightarrow H^1((\mathbf f),M)\longrightarrow0.
\end{equation}
Todos los homomorfismos introducidos son de grado $0$ para las estructuras de $S$-módulo graduado, donde $S$ es el anillo (\hyperref[III.2.1.1.1]{2.1.1.1}).
\end{proposition}

\begin{corollary}[2.1.4]
\label{III.2.1.4}
\label{III.2.1.4-fr}
\oldpage[III]{97}
Si $X$ es un esquema cuasicompacto y los $U_i=X_{f_i}$ son afines, existe un isomorfismo canónico, funtorial en $\mathcal F$ y de grado $0$,
\begin{equation}
\label{III.2.1.4.1-fr}
\tag{2.1.4.1}
H^p(U,\mathcal F(*))\xrightarrow{\sim}H^{p+1}((\mathbf f),M)
\qquad\text{para }p\geq1
\end{equation}
y una sucesión exacta, funtorial en $\mathcal F$,
\begin{equation}
\label{III.2.1.4.2-fr}
\tag{2.1.4.2}
0\longrightarrow H^0((\mathbf f),M)\longrightarrow M
\longrightarrow H^0(U,\mathcal F(*))
\longrightarrow H^1((\mathbf f),M)\longrightarrow0
\end{equation}
en la que todos los homomorfismos son de grado $0$.
\end{corollary}

\begin{proof}
Basta aplicar \sref{III.1.4.1} al resultado de \sref{III.2.1.3}.
\end{proof}

La proposición ``local'' análoga a \sref{III.2.1.3} es la siguiente.

\begin{proposition}[2.1.5]
\label{III.2.1.5}
\label{III.2.1.5-fr}
Sea $S$ un anillo graduado positivamente, sean los $f_i$ ($(1\leq i\leq r)$) elementos homogéneos de $S_+$ de grados $d_i$, y sea $M$ un $S$-módulo graduado.
Sea $X=\operatorname{Proj}(S)$ el espectro primo homogéneo de $S$, y pongamos $U_i=D_+(f_i)$, $U=\bigcup_iU_i$ y $H^\bullet(U,\widetilde M(*))
=\bigoplus_{n\in\mathbb Z}H^\bullet(U,(M(n))^\sim)$.
Existen isomorfismos canónicos, funtoriales en $M$ y de grado $0$ para las estructuras de $S$-módulo graduado,
\begin{equation}
\label{III.2.1.5.1-fr}
\tag{2.1.5.1}
H^p(U,\widetilde M(*))\xrightarrow{\sim}H^{p+1}((\mathbf f),M)
\qquad\text{para }p\geq1
\end{equation}
y una sucesión exacta, funtorial en $M$,
\begin{equation}
\label{III.2.1.5.2-fr}
\tag{2.1.5.2}
0\longrightarrow H^0((\mathbf f),M)\longrightarrow M
\longrightarrow H^0(U,\widetilde M(*))
\longrightarrow H^1((\mathbf f),M)\longrightarrow0
\end{equation}
en la que todos los homomorfismos son de grado $0$.
\end{proposition}

\begin{proof}
En efecto, por definición \sref[II]{II.2.5.2},
$\Gamma(U_{i_0i_1\ldots i_p},(M(n))^\sim)
=(M_{f_{i_0}\ldots f_{i_p}})_n$
y, por tanto, $\Gamma(U_{i_0i_1\ldots i_p},\widetilde M(*))
=M_{f_{i_0}\ldots f_{i_p}}$.
El resto del razonamiento es el mismo que en la demostración de \sref{III.2.1.4}, teniendo en cuenta que $X$ es un esquema.
\end{proof}

\begin{env}[2.1.6]
\label{III.2.1.6}
\label{III.2.1.6-fr}
\emph{Observaciones.} ---
(i) Bajo las hipótesis de \sref{III.2.1.5}, los funtores $\Gamma(U_{i_0i_1\ldots i_p},\widetilde M(*))$ son exactos en $M$, por \sref[0]{0.1.3.2}.
El mismo razonamiento que en (\hyperref[III.1.2.5]{1.2.5}) muestra, por tanto, que, si $0\to M'\to M\to M''\to0$ es una sucesión exacta de $S$-módulos graduados cuyos homomorfismos tienen grado $0$, entonces para todo $p\geq0$ existen diagramas conmutativos
\begin{equation}
\label{III.2.1.6.1-fr}
\tag{2.1.6.1}
\xymatrix{
H^p(U,\widetilde{M}''(*))\ar[r]^-\partial\ar[d] &
H^{p+1}(U,\widetilde{M}'(*))\ar[d]\\
H^{p+1}((\mathbf f),M'')\ar[r]_-\partial &
H^{p+2}((\mathbf f),M')
}
\end{equation}

(ii) La Proposición~\sref{III.2.1.5} es especialmente útil cuando $S$ es una $A$-álgebra engendrada por un número finito de elementos de grado $1$, con $A$ supuesto noetheriano; en
\oldpage[III]{98}
efecto, cuando es así, \emph{todo} $\mathcal O_X$-módulo cuasicoherente es de la forma $\widetilde M$ \sref[II]{II.2.7.7}.
\end{env}

\begin{env}[2.1.7]
\label{III.2.1.7}
\label{III.2.1.7-fr}

Vamos a aplicar \sref{III.2.1.5} al caso $S=A[T_0,\ldots,T_r]$, donde $A$ es un anillo arbitrario y las $T_i$ son indeterminadas, tomando $M=S$ y $f_i=T_i$. Se reduce pues esencialmente a calcular $H^\bullet((\mathbf T),S)$, donde $\mathbf T=(T_i)_{0\leq i\leq r}$.
\end{env}

\begin{lemma}[2.1.8]
\label{III.2.1.8}
\label{III.2.1.8-fr}
Si $S=A[T_0,\ldots,T_r]$, entonces, con $\mathbf T=(T_i)_{0\leq i\leq r}$
\begin{align}
\label{III.2.1.8.1-fr}
\tag{2.1.8.1}
H^i(\mathbf T^n,S)&=0 &&\text{si }i\neq r+1\\
\label{III.2.1.8.2-fr}
\tag{2.1.8.2}
H^{r+1}(\mathbf T^n,S)&=S/(\mathbf T^n).
\end{align}
El $A$-módulo $H^{r+1}(\mathbf T^n,S)$ tiene una $A$-base formada por las clases módulo $(\mathbf T^n)$ de los monomios
$\mathbf T_{\mathbf p}=T_0^{p_0}\ldots T_r^{p_r}$ con
$\mathbf p=(p_0,\ldots,p_r)$, $0\leq p_i<n$ para todo $i$.
\end{lemma}

\begin{proof}
Es una consecuencia inmediata de (1.1.3.5) y de la prop. (1.1.4), cuyas
hipótesis se verifican trivialmente.
\end{proof}

\begin{env}[2.1.9]
\label{III.2.1.9-fr}
Pasemos al límite inductivo sobre $n$, las relaciones (\hyperref[III.2.1.8.1]{2.1.8.1}) dan $H^i((\mathbf T),S)=0$ para $i\neq r+1$.
Para $i=r+1$, el sistema inductivo consta de los módulos $S/(\mathbf T^n)$, siendo el homomorfismo
$\varphi_{nm}:S/(\mathbf T^n)\longrightarrow S/(\mathbf T^m)$
para $0\leq n\leq m$
la multiplicación por $(T_0\ldots T_r)^{m-n}$.
Para $n\geq\sup(p_i)_{0\leq i\leq r}$, denotemos por $\xi_{\mathbf p}^{(n)}=\xi_{p_0\ldots p_r}^{(n)}$ la clase de
$T_0^{n-p_0}\ldots T_r^{n-p_r}\ \mathrm{mod.}(\mathbf T^n)$; se tiene entonces
$\varphi_{nm}(\xi_{\mathbf p}^{(n)})=\xi_{\mathbf p}^{(m)}$, de modo que estos elementos tienen la misma imagen canónica $\xi_{\mathbf p}=\xi_{p_0\ldots p_r}$ en el límite inductivo $H^{r+1}((\mathbf T),S)$.
Por la definición del grado en (\hyperref[III.2.1.2]{2.1.2}), el grado de $\xi_{\mathbf p}$ es $-|\mathbf p|=-(p_0+p_1+\ldots+p_r)$.
Está claro que los $\xi_{\mathbf p}^{(n)}$ con $0<p_i\leq n$ para $0\leq i\leq r$ forman una base de $S/(\mathbf T^n)$.
Se deduce, pues, inmediatamente de \sref{III.2.1.8}:
\end{env}

\begin{corollary}[2.1.10]
\label{III.2.1.10}
\label{III.2.1.10-fr}
Con la notación de \sref{III.2.1.8},
\begin{equation}
\label{III.2.1.10.1-fr}
\tag{2.1.10.1}
H^i((\mathbf T),S)=0\qquad\text{para }i\neq r+1
\end{equation}
y $H^{r+1}((\mathbf T),S)$ es un $A$-módulo libre con base los elementos $\xi_{p_0\ldots p_r}$ tales que $p_i>0$ para $0\leq i\leq r$.
\end{corollary}

\begin{env}[2.1.11]
\label{III.2.1.11}
\label{III.2.1.11-fr}
\emph{Observación.} ---
Sea $N$ un $A$-módulo cualquiera y sea $M=S\otimes_A N$; el razonamiento de
(\hyperref[III.2.1.8-fr]{2.1.8}) muestra que se tiene más generalmente
\begin{align}
\label{III.2.1.11.1-fr}
\tag{2.1.11.1}
H^i(\mathbf T^n,M)&=0 &&\text{si }i\neq r+1\\
\label{III.2.1.11.2-fr}
\tag{2.1.11.2}
H^{r+1}(\mathbf T^n,M)&=(S/(\mathbf T^n))\otimes_A N
\end{align}
pues la última fórmula se deduce directamente de (1.1.3.5), y por otra
parte está claro que
$M/(T_0^nM+\ldots+T_{i-1}^nM)$
se identifica con el producto tensorial
$(S/(T_0^nS+\ldots+T_{i-1}^nS))\otimes_A N$,
el ideal $T_0^nS+\ldots+T_{i-1}^nS$ es un sumando directo del $A$-módulo $S$.
Esto permite aplicar \sref{III.1.1.4} a $M$ y da (\hyperref[III.2.1.11.1]{2.1.11.1}).
\end{env}

Combinando \sref{III.2.1.10} y \sref{III.2.1.5}, se obtiene:

\begin{proposition}[2.1.12]
\label{III.2.1.12}
\label{III.2.1.12-fr}
Sea $A$ un anillo, sea $r$ un entero $>0$ y sea $X=\mathbf P_A^r$ \sref[II]{II.4.1.1}.
Entonces:
\begin{enumerate}
  \item[{\rm(i)}] $H^i(X,\mathcal O_X(*))=0$ para $i\neq0,r$.
  \item[{\rm(ii)}] El homomorfismo canónico $\alpha:S\to H^0(X,\mathcal O_X(*))$ \sref[II]{II.2.6.2} es biyectivo.
\oldpage[III]{99}
  \item[{\rm(iii)}] $H^r(X,\mathcal O_X(*))$ es un $A$-módulo libre con base los elementos $\xi_{p_0\ldots p_r}$, donde $p_i>0$ para $0\leq i\leq r$; el elemento $\xi_{p_0\ldots p_r}$ tiene grado $-|\mathbf p|=-(p_0+\ldots+p_r)$, y siendo el producto
$T_i\xi_{p_0\ldots p_r}$ igual a
$\xi_{p_0,\ldots,p_i-1,\ldots,p_r}$.
\end{enumerate}
\end{proposition}

\begin{proof}
Observemos en efecto que, en la sucesión exacta
(\hyperref[III.2.1.5.2-fr]{2.1.5.2}) aplicada a
\[
M=S=A[T_0,\ldots,T_r],
\],
se tiene
$H^0((\mathbf T),S)=0$ y $H^1((\mathbf T),S)=0$ por \sref{III.2.1.10.1}; además, la Proposición~\sref{III.2.1.5} se aplica con $U=X$, puesto que $X$ es la unión de los $D_+(T_i)$ \sref[II]{II.2.3.14}.
Queda identificar la aplicación $S\to H^0(X,\mathcal O_X(*))$ de (\hyperref[III.2.1.5.2]{2.1.5.1}) con la aplicación canónica $\alpha$.
Esto se deduce de la identificación canónica de $H^0(U,\mathcal O_X(*))$ y $H^0(\mathfrak U,\mathcal O_X(*))$.
\end{proof}

\begin{corollary}[2.1.13]
\label{III.2.1.13}
\label{III.2.1.13-fr}
Los únicos pares $(i,n)$ para los cuales puede ocurrir que $H^i(X,\mathcal O_X(n))\neq0$ son $i=0$, $n\geq0$, e $i=r$, $n\leq-(r+1)$.
\end{corollary}

Si $A\neq0$, entonces se tiene efectivamente $H^i(X,\mathcal O_X(n))\neq0$ para los pares enumerados en
(\hyperref[III.2.1.13-fr]{2.1.13}); esto resulta de
(\hyperref[III.2.1.12-fr]{2.1.12}), puesto que $S_n$ es entonces $\neq0$ para
todos los grados $n\geq0$.

En las aplicaciones de este capítulo, utilizaremos principalmente el siguiente resultado menos preciso.

\begin{corollary}[2.1.14]
\label{III.2.1.14}
\label{III.2.1.14-fr}
Los $A$-módulos $H^i(X,\mathcal O_X(n))$ son libres de tipo finito; si $i>0$, se anulan para $n>0$.
\end{corollary}

\begin{proposition}[2.1.15]
\label{III.2.1.15}
\label{III.2.1.15-fr}
Sea $Y$ un preesquema, sea $\mathcal E$ un $\mathcal O_Y$-módulo localmente libre de rango $r+1$, sea $X=\mathbf P(\mathcal E)$ el fibrado proyectivo definido por $E$, y sea $f:X\to Y$ el morfismo estructural.
Los únicos valores de $i$ y $n$ para los cuales puede ocurrir que $R^if_*(\mathcal O_X(n))\neq0$ son $i=0$, $n\geq0$, e $i=r$, $n\leq-(r+1)$.
Además, el homomorfismo canónico \sref[II]{II.3.3.2}
\[
\alpha:\mathbf S_{\mathcal O_Y}(\mathcal E)
\longrightarrow\bbGamma_*(\mathcal O_X)
=R^0f_*(\mathcal O_X(*))
=\bigoplus_{n\in\mathbb Z}f_*(\mathcal O_X(n))
\]
es un isomorfismo.
\end{proposition}

\begin{proof}
La cuestión es local sobre $Y$, por lo que se puede suponer que $Y$ es afín de anillo $A$ y que $\mathcal E=\widetilde E$, donde $E=A^{r+1}$; se reduce entonces inmediatamente a (\hyperref[III.2.1.12-fr]{2.1.12}), teniendo en cuenta (1.4.11).
\end{proof}

\begin{env}[2.1.16]
\label{III.2.1.16}
\label{III.2.1.16-fr}
\emph{Observación.} ---
Más adelante completaremos los resultados de \sref{III.2.1.15} demostrando las proposiciones siguientes.
Pongamos
$\boldsymbol\omega=f^*(\bigwedge^{r+1}\mathcal E)(-r-1)$,
que es un $\mathcal O_X$-Módulo invertible. Entonces:
\begin{enumerate}
\item[{\rm(i)}] Se tiene un isomorfismo canónico
\begin{equation}
\label{III.2.1.16.1-fr}
\tag{2.1.16.1}
\rho:R^rf_*(\boldsymbol\omega)\xleftarrow{\sim}\mathcal O_Y.
\end{equation}
\item[{\rm(ii)}] El apareamiento por producto cup (\textbf{0}, 12.2.2)
\begin{equation}
\label{III.2.1.16.2-fr}
\tag{2.1.16.2}
R^rf_*(\mathcal O_X(n))\times R^0f_*(\boldsymbol\omega(-n))
\longrightarrow R^rf_*(\boldsymbol\omega)
\end{equation}
  compuesto con $\rho^{-1}$, define un \emph{isomorfismo} de $R^rf_*(\mathcal O_X(n))$ sobre el \emph{dual} del $\mathcal O_Y$-módulo localmente libre
  \[
R^0f_*(\boldsymbol\omega(-n))
=\left(\bigwedge^{r+1}\mathcal E\right)
\otimes_{\mathcal O_Y}\bigl(\mathbf S_{\mathcal O_Y}(\mathcal E)\bigr)_{-n}.
\]
\end{enumerate}
\end{env}

\subsection{El teorema fundamental de los morfismos proyectivos}
\label{subsection:III.2.2}
\label{subsection:III.2.2-fr}
\oldpage[III]{100}

\begin{theorem}[2.2.1]
\label{III.2.2.1-fr}
(Serre). Sean $Y$ un preesquema noetheriano, $f:X\to Y$ un morfismo propio, $\mathcal L$ un $\mathcal O_X$-módulo invertible amplio para $f$. Para todo $\mathcal O_X$-módulo $\mathcal F$, pongamos
$\mathcal F(n)=\mathcal F\otimes_{\mathcal O_X}\mathcal L^{\otimes n}$
para todo $n\in\mathbb Z$. Entonces, para todo $\mathcal O_X$-módulo coherente $\mathcal F$:
\begin{enumerate}
\item[{\rm(i)}] Los $R^qf_*(\mathcal F)$ son $\mathcal O_Y$-módulos coherentes.
\item[{\rm(ii)}] Existe un entero $N$ tal que para $n\geq N$ se tenga $R^qf_*(\mathcal F(n))=0$ para todo $q>0$.
\item[{\rm(iii)}] Existe un entero $N$ tal que, para $n\geq N$, el homomorfismo canónico
$f^*(f_*(\mathcal F(n)))\longrightarrow\mathcal F(n)$
sea sobreyectivo.
\end{enumerate}
\end{theorem}

\begin{proof}
Observemos primero que si el teorema es verdadero cuando en él se reemplaza $\mathcal L$ por $\mathcal L^{\otimes d}$ $(d>0)$, es verdadero en su forma inicial. En efecto, se puede entonces escribir $\mathcal F(n)
=(\mathcal F\otimes\mathcal L^{\otimes r})
\otimes\mathcal L^{\otimes hd}$ con un $h>0$ y $0\leq r<d$, y por hipótesis para cada $r$ hay un entero $N_r$ tal que para $h\geq N_r$ las propiedades (ii) y (iii) se cumplan para el $\mathcal O_X$-módulo $\mathcal F\otimes\mathcal L^{\otimes r}$; tomando para $N$ el mayor de los $dN_r$, (ii) y (iii) se cumplirán para $n\geq N$. Se puede, pues, suponer $\mathcal L$ \emph{muy amplio} relativamente a $f$ (\textbf{II}, 4.6.11); existe, por consiguiente, una $Y$-inmersión abierta dominante $i:X\to P$, donde $P=\operatorname{Proj}(\mathcal S)$, donde $\mathcal S$ es una $\mathcal O_Y$-álgebra cuasicoherente graduada en grados positivos, en la que $\mathcal S_1$ es de tipo finito y engendra $\mathcal S$; además, $\mathcal L$ es isomorfo a $i^*(\mathcal O_P(1))$ (\textbf{II}, 4.4.7). Pero como $f$ es propio, lo mismo ocurre con $i$ (\textbf{II}, 5.4.4), por tanto $i$ es un \emph{isomorfismo} $X\xrightarrow{\sim}P$. Se puede, pues, limitarse al caso en que $X=\operatorname{Proj}(\mathcal S)$ y $\mathcal L=\mathcal O_X(1)$. El teorema (2.2.1) es entonces consecuencia de la
\end{proof}

\begin{proposition}[2.2.2]
\label{III.2.2.2-fr}
Sean $A$ un anillo noetheriano, $S$ una $A$-álgebra graduada en grados positivos, en la que $S_1$ es un $A$-módulo que tiene un sistema de $r+1$ generadores, y que engendra el álgebra $S$. Sea $X=\operatorname{Proj}(S)$. Para todo $\mathcal O_X$-módulo coherente $\mathcal F$:
\begin{enumerate}
\item[{\rm(i)}] Los $A$-módulos $H^q(X,\mathcal F)$ son de tipo finito.
\item[{\rm(ii)}] Se tiene $H^q(X,\mathcal F)=0$ para $q>r$.
\item[{\rm(iii)}] Existe un entero $N$ tal que para $n\geq N$ se tenga $H^q(X,\mathcal F(n))=0$ para todo $q>0$.
\item[{\rm(iv)}] Existe un entero $N$ tal que para $n\geq N$, $\mathcal F(n)$ esté engendrado por sus secciones sobre $X$.
\end{enumerate}
\end{proposition}

\begin{proof}
Mostremos primero cómo (2.2.2) implica (2.2.1): en (2.2.1) (reducido al caso particular $X=\operatorname{Proj}(\mathcal S)$ considerado más arriba), $Y$ es cuasicompacto, por tanto puede recubrirse por un número finito de abiertos afines, de anillos noetherianos, tales que la restricción de $\mathcal S_1$ a cada uno de estos abiertos $U_\alpha$ esté engendrada por un número finito de secciones de $\mathcal S_1$ sobre $U_\alpha$. Si se supone demostrado (2.2.2), bastará entonces tomar para $N$ en las partes (ii) y (iii) de (2.2.1) el mayor de los enteros análogos correspondientes a los $U_\alpha$ (teniendo en cuenta (1.4.11) y (\textbf{II}, 3.4.7)).

Para demostrar (2.2.2), observemos que $X$ se identifica con un subesquema cerrado de $P=\mathbf P_A^r$ (\textbf{II}, 3.6.2); además, si $j:X\to P$ es la inyección canónica, $j_*(\mathcal F)$ es un $\mathcal O_P$-módulo coherente y se tiene $j_*(\mathcal F(n))=(j_*(\mathcal F))(n)$ (\textbf{II}, 3.4.5 y 3.5.2). Habida cuenta de (G, \textbf{II}, cor. del teorema 4.9.1), se reduce, pues, a demostrar (2.2.2) en el caso particular en que $X=\mathbf P_A^r$ y $S=A[T_0,\ldots,T_r]$. Como $X$ está recubierto por los abiertos
\oldpage[III]{101}
afines $D_+(T_i)$ en número $r+1$, (ii) resulta de (1.4.12). Observemos por otra parte que (iv) ya ha sido demostrado (\textbf{II}, 2.7.9).

Vamos a demostrar simultáneamente (i) y (iii). Observemos que estas aserciones son verdaderas para $\mathcal F=\mathcal O_X(m)$ (\hyperref[III.2.1.13-fr]{2.1.13}); lo son, pues, también cuando $\mathcal F$ es suma directa de un número finito de $\mathcal O_X$-módulos de la forma $\mathcal O_X(m_j)$. Por otra parte, (i) y (iii) son trivialmente verdaderas para $q>r$ en virtud de (ii). Vamos a proceder por \emph{inducción descendente} sobre $q$. Se sabe que $\mathcal F$ es isomorfo a un cociente de una suma directa $\mathcal E$ de un número finito de haces $\mathcal O_X(m_j)$ (\textbf{II}, 2.7.10); es decir, se tiene una sucesión exacta
\[
0\longrightarrow\mathcal R\longrightarrow\mathcal E
\longrightarrow\mathcal F\longrightarrow0,
\]
donde $\mathcal R$ es coherente (\textbf{0}\textsubscript{I}, 5.3.3) y donde $\mathcal E$ satisface (i) y (iii). Como $\mathcal F(n)$ es un funtor exacto en $\mathcal F$, se tiene también la sucesión exacta
\[
0\longrightarrow\mathcal R(n)\longrightarrow\mathcal E(n)
\longrightarrow\mathcal F(n)\longrightarrow0
\]
para todo $n\in\mathbb Z$. Se deduce de ella la sucesión exacta de cohomología
\[
H^{q-1}(X,\mathcal E(n))\longrightarrow
H^{q-1}(X,\mathcal F(n))\longrightarrow
H^q(X,\mathcal R(n)).
\]

Como $\mathcal E(n)$ es suma directa de los $\mathcal O_X(n+m_j)$ (\textbf{II}, 2.5.14), $H^{q-1}(X,\mathcal E(n))$ es de tipo finito, y lo mismo ocurre con $H^q(X,\mathcal R(n))$ por la hipótesis de inducción; como $A$ es noetheriano, se concluye que $H^{q-1}(X,\mathcal F(n))$ es de tipo finito para todo $n\in\mathbb Z$, y en particular para $n=0$. Por otra parte, por la hipótesis de inducción, existe un entero $N$ tal que para $n\geq N$ se tenga $H^q(X,\mathcal R(n))=0$; además, se puede suponer también $N$ elegido de modo que $H^{q-1}(X,\mathcal E(n))=0$ para $n\geq N$, puesto que $\mathcal E$ satisface (iii); se concluye que $H^{q-1}(X,\mathcal F(n))=0$ para $n\geq N$, lo que termina la demostración.
\end{proof}

\begin{corollary}[2.2.3]
\label{III.2.2.3}
\label{III.2.2.3-fr}
Bajo las hipótesis de \textup{\sref{III.2.2.1}}, sea
$\mathcal F\to\mathcal G\to\mathcal H$
una sucesión exacta de $\mathcal O_X$-módulos coherentes.
Existe un entero $N$ tal que, para $n\geq N$, la sucesión
\[
f_*(\mathcal F(n))\longrightarrow f_*(\mathcal G(n))
\longrightarrow f_*(\mathcal H(n))
\]
es exacta.
\end{corollary}

\begin{proof}
Sean $\mathcal F'$, $\mathcal G'$ y $\mathcal G''$ el núcleo, la imagen y el conúcleo de $\mathcal F\to\mathcal G$.
Entonces $\mathcal G'$ es el núcleo y $\mathcal G''$ la imagen de $\mathcal G\to\mathcal H$; sea $\mathcal H''$ el conúcleo de este último homomorfismo.
Todos estos $\mathcal O_X$-Módulos son
coherentes (\textbf{0}\textsubscript{I}, 5.3.4). Como $\mathcal F(n)$ es un
funtor exacto en $\mathcal F$, basta probar que para $n$ suficientemente grande, cada una de las sucesiones
\begin{align*}
0&\longrightarrow f_*(\mathcal F'(n))
\longrightarrow f_*(\mathcal F(n))
\longrightarrow f_*(\mathcal G'(n))\longrightarrow0\\
0&\longrightarrow f_*(\mathcal G'(n))
\longrightarrow f_*(\mathcal G(n))
\longrightarrow f_*(\mathcal G''(n))\longrightarrow0\\
0&\longrightarrow f_*(\mathcal G''(n))
\longrightarrow f_*(\mathcal H(n))
\longrightarrow f_*(\mathcal H''(n))\longrightarrow0
\end{align*} es exacta; por consiguiente, se puede suponer que $0\to\mathcal F\to\mathcal G\to\mathcal H\to0$
es exacta.
La sucesión de cohomología es entonces
\[
0\longrightarrow f_*(\mathcal F(n))
\longrightarrow f_*(\mathcal G(n))
\longrightarrow f_*(\mathcal H(n))
\longrightarrow R^1f_*(\mathcal F(n))\longrightarrow\ldots
\]
y la conclusión se deduce de \sref{III.2.2.1}, (ii).
\end{proof}

\oldpage[III]{102}

\begin{corollary}[2.2.4]
\label{III.2.2.4}
\label{III.2.2.4-fr}
Sea $Y$ un preesquema noetheriano, sea $f:X\to Y$ un morfismo de tipo finito y sea $\mathcal L$ un $\mathcal O_X$-módulo invertible amplio respecto de $f$.
Para todo $\mathcal O_X$-módulo $\mathcal F$, se pone
$\mathcal F(n)=\mathcal F\otimes_{\mathcal O_X}\mathcal L^{\otimes n}$
(para $n\in\mathbb Z$). Sea
$\mathcal F\to\mathcal G\to\mathcal H$
una sucesión exacta de $\mathcal O_X$-módulos coherentes tal que los soportes de $\mathcal F$ y $\mathcal H$ sean propios sobre $Y$ \textup{\sref[II]{II.5.4.10}}.
Entonces existe un entero $N$ tal que, para $n\geq N$, la sucesión
\[
f_*(\mathcal F(n))\longrightarrow f_*(\mathcal G(n))
\longrightarrow f_*(\mathcal H(n))
\]
es exacta.
\end{corollary}

\begin{proof}
El razonamiento del comienzo de \sref{III.2.2.1} muestra que, si el corolario es cierto para $\mathcal L^{\otimes d}$, donde $(d>0)$, entonces es cierto para $\mathcal L$; se puede pues limitar al caso en que $\mathcal L$ es muy amplio respecto de $f$ (\textbf{II}, 4.6.11), y por consiguiente se puede identificar $X$ con un subconjunto abierto de un $Y$-esquema
$Z=\operatorname{Proj}(\mathcal S)$,
donde $\mathcal S$ es una $\mathcal O_Y$-álgebra cuasicoherente graduada en grados positivos, $\mathcal S_1$ es de tipo finito y engendra $\mathcal S$, y
$\mathcal L=i^*(\mathcal O_Z(1))$, donde $i$ es la inmersión
canónica $X\to Z$ (\textbf{II}, 4.4.7). Dicho esto, como
$\operatorname{Supp}(\mathcal G)$ es cerrado en $X$ y está contenido en
$\operatorname{Supp}(\mathcal F)\cap\operatorname{Supp}(\mathcal H)$,
por tanto, es propio sobre $Y$.
En consecuencia, los soportes de $\mathcal F$, $\mathcal G$ y $\mathcal H$ son cerrados en $Z$ (\textbf{II}, 5.4.10). Los haces
$\mathcal F'=i_*(\mathcal F)$, $\mathcal G'=i_*(\mathcal G)$,
$\mathcal H'=i_*(\mathcal H)$ son pues $\mathcal O_Z$-módulos coherentes, y la sucesión
$\mathcal F'\to\mathcal G'\to\mathcal H'$
es exacta.
Si $g:Z\to Y$ es el morfismo estructural, entonces $f=g\circ i$, y claramente
$\mathcal F'(n)=i_*(\mathcal F(n))$
con identidades análogas para $\mathcal G'$ y $\mathcal H'$.
La conclusión se deduce de \sref{III.2.2.3}, aplicado a $\mathcal F'$, $\mathcal G'$ y $\mathcal H'$.
\end{proof}

\begin{env}[2.2.5]
\label{III.2.2.5}
\label{III.2.2.5-fr}
\emph{Observaciones.}

(i) La afirmación (i) de \sref{III.2.2.1} sigue siendo cierta si solo se supone que $Y$ es \emph{localmente noetheriano}.
En efecto, la propiedad es manifiestamente local sobre $Y$.
Además, las hipótesis de \sref{III.2.2.1} implican que, para todo subconjunto abierto $U\subset Y$, la restricción de $f$ a $f^{-1}(U)$ es un morfismo proyectivo
$f^{-1}(U)\to U$
\sref[II]{II.5.5.5}[iii], y $\mathcal L|f^{-1}(U)$ es amplio para este morfismo \sref[II]{II.4.6.4}.

(ii) La afirmación (iii) de \sref{III.2.2.1} sigue siendo válida, como se ha visto, si solo se supone que $X$ es un esquema cuasicompacto, o un preesquema cuyo espacio subyacente es noetheriano, y que $f:X\to Y$ un morfismo cuasicompacto (\textbf{II}, 4.6.8). Pero hay que observar que incluso cuando se supone que $Y$ es el \emph{espectro de un cuerpo} $K$ y $f$ es \emph{cuasiproyectivo}, la afirmación (ii) de \sref{III.2.2.1} puede dejar de cumplirse.

Por ejemplo, sea
$X'=\operatorname{Spec}(K[T_0,\ldots,T_r])$
y sea $X$ la unión de los subconjuntos abiertos afines $D(T_i)$ de $X'$, para $(0\leq i\leq r)$.
Como la inmersión $X\to X'$ es cuasicompacta, el morfismo estructural $f:X\to Y$ es cuasiafín \sref[II]{II.5.1.10} y, por tanto, $\mathcal O_X$ es muy amplio respecto de $f$ (\textbf{II}, 5.1.6). Pero el anillo
$\Gamma(X,\mathcal O_X)$ se identifica con la intersección de los anillos de
fracciones $(K[T_0,\ldots,T_r])_{T_i}$ para $0\leq i\leq r$
(\textbf{I}, 8.2.1.1), es decir, con $K[T_0,\ldots,T_r]$. Por tanto, resulta
de las fórmulas (1.4.3.1) y (1.1.3.5) que se tiene
$H^r(X,\mathcal O_X^{\otimes n})=H^r(X,\mathcal O_X)=A\neq0$
para todo $n$.
\end{env}


\subsection{Aplicación a los haces graduados de álgebras y de módulos}
\label{subsection:III.2.3}
\label{subsection:III.2.3-fr}

\begin{theorem}[2.3.1]
\label{III.2.3.1}
\label{III.2.3.1-fr}
Sea $Y$ un preesquema noetheriano, sea $\mathcal S$ una $\mathcal O_Y$-álgebra cuasicoherente de tipo finito, graduada en grados positivos, sea
$X=\operatorname{Proj}(\mathcal S)$,
y sea $q:X\to Y$ el morfismo estructural.
Sea $\mathcal M$ un $\mathcal S$-módulo graduado cuasicoherente que satisface la condición \textup{\textbf{(TF)}}.
Existe entonces un entero $N$ tal que, para $n\geq N$, el homomorfismo canónico \textup{\sref[II]{II.8.14.5.1}}
\[
\alpha_n:\mathcal M_n\longrightarrow
q_*(\shProj_0(\mathcal M(n)))
=q_*((\shProj(\mathcal M))_n)
\]
\oldpage[III]{103}
es biyectivo.
En otras palabras, el homomorfismo canónico
\[
\alpha:\mathcal M\longrightarrow\bbGamma_*(\shProj(\mathcal M))
\]
es un \textup{\textbf{(TN)}}-isomorfismo.
\end{theorem}

\begin{proof}
Se puede restringir al caso en que $\mathcal M$ es un $\mathcal S$-módulo de tipo finito \sref[II]{II.3.4.2}.
Como $Y$ es cuasicompacto, existe un entero $d>0$ tal que $\mathcal S^{(d)}$ está engendrada por el $\mathcal O_Y$-módulo cuasicoherente $\mathcal S_d$, que es de tipo finito \sref[II]{II.3.1.10} y, por tanto, coherente porque $Y$ es noetheriano. Observemos ahora que $\mathcal M$ es la suma directa de los $\mathcal M^{(d,k)}$, para $0\leq k<d$, y cada $\mathcal M^{(d,k)}$ es un $\mathcal S^{(d)}$-módulo cuasicoherente de tipo finito, por \sref[II]{II.2.1.6}[iii].
Como la cuestión es local sobre $Y$, basta manifiestamente demostrar que cada homomorfismo canónico
\[
\alpha:\mathcal M^{(d,k)}
\longrightarrow\bbGamma_*((\shProj(\mathcal M))^{(d,k)})
\] es un \textbf{(TN)}-isomorfismo. Teniendo en cuenta (\textbf{II}, 8.14.13) (y en particular el diagrama (8.14.13.4)), se ve que se reduce a probar el teorema cuando $\mathcal S$ está engendrada por $\mathcal S_1$ y $\mathcal S_1$ es un $\mathcal O_Y$-módulo coherente.

Como $Y$ es noetheriano, el mismo razonamiento que al comienzo de
(\hyperref[III.2.2.2-fr]{2.2.2}) muestra que se puede limitar al caso en que
$Y=\operatorname{Spec}(A)$, $\mathcal S=\widetilde S$,
$\mathcal M=\widetilde M$, $A$ es un anillo noetheriano, $S_1$ es un $A$-módulo de tipo finito y $M$ es un $S$-módulo graduado de tipo finito.
Mostraremos que basta entonces demostrar el teorema para $M=S$.
En efecto, en el caso general existe una sucesión exacta
\[
L'\longrightarrow L\longrightarrow M\longrightarrow0,
\]
donde $L$ y $L'$ son sumas directas de módulos graduados de la forma $S(m)$.
Si el resultado se cumple para $M=S$, se cumple también para $M=S(m)$ y, por tanto, para $L$ y $L'$.
Consideremos el diagrama conmutativo
\[
\xymatrix{
\widetilde{L}'_n\ar[r]\ar[d]_{\alpha_n} &
\widetilde{L}_n\ar[r]\ar[d]_{\alpha_n} &
\widetilde{M}_n\ar[r]\ar[d]_{\alpha_n} & 0\\
q_*(\widetilde {L'}(n))\ar[r] &
q_*(\widetilde L(n))\ar[r] &
q_*(\widetilde M(n))\ar[r] & 0
}
\]
\oldpage[III]{104}
La segunda fila es exacta por \sref{III.2.2.3} cuando $n$ es suficientemente grande.
Lo mismo ocurre con la primera fila y, como las dos flechas verticales de la izquierda son isomorfismos, también lo es la tercera.

Queda, pues, demostrar el teorema para $M=S$.
Supongamos primero que
$S=A[T_0,\ldots,T_r]$ ($T_i$ indeterminadas); en este caso, nuestra aserción no es sino (\hyperref[III.2.1.11-fr]{2.1.11}, (ii)). En el caso general, $S$ se identifica con un cociente de un anillo
$S'=A[T_0,\ldots,T_r]$
por un ideal graduado, de modo que $X$ es un subesquema cerrado de
$X'=\mathbf P_A^r$ (\textbf{II}, 2.9.2). Si $j$ es la inclusión canónica
$X\to X'$, $j_*(\widetilde S(n))$ es el $\mathcal O_{X'}$-módulo
$(\shProj(\widetilde S))(n)$
donde $S$ se considera como $S'$-módulo graduado; esto se deduce inmediatamente de \sref[II]{II.2.8.7}.
Como $j_*(\widetilde S(n))$ es un $\mathcal O_{X'}$-módulo que satisface \textbf{(TF)}, el homomorfismo canónico
$\alpha_n:S_n\to\Gamma(X',j_*(\widetilde S(n)))$
es biyectivo para todo $n$ suficientemente grande, en virtud de lo que precede; esto termina la demostración, puesto que $\Gamma(X',j_*(\widetilde S(n)))=\Gamma(X,\widetilde S(n))$.
\end{proof}

\begin{corollary}[2.3.2]
\label{III.2.3.2}
\label{III.2.3.2-fr}
Bajo las hipótesis de \textup{\sref{III.2.3.1}}, sea
\[
\mathcal S_X=\bigoplus_{n\in\mathbb Z}\mathcal O_X(n),
\]
y sea $\mathcal F$ un $\mathcal S_X$-módulo graduado cuasicoherente de tipo finito.
Entonces $\bbGamma_*(\mathcal F)$ satisface la condición \textup{\textbf{(TF)}}.
\end{corollary}

\begin{proof}
En la demostración de \sref{III.2.3.1} se ha visto que $X$, que es isomorfo a $\operatorname{Proj}(\mathcal S^{(d)})$ \sref[II]{II.3.1.8}, es de tipo finito sobre $Y$ \sref[II]{II.3.4.1}.
Se deduce de \sref[II]{II.8.14.9} que $\mathcal F$ es isomorfo a un $\mathcal S_X$-módulo graduado de la forma $\shProj(\mathcal M)$, donde $\mathcal M$ es un $\mathcal S$-módulo graduado cuasicoherente de tipo finito.
Por \sref{III.2.3.1}, $\bbGamma_*(\mathcal F)$ es \textbf{(TN)}-isomorfo a $\mathcal M$ y, por consiguiente, satisface \textbf{(TF)}.
\end{proof}

\begin{env}[2.3.3]
\label{III.2.3.3}
\label{III.2.3.3-fr}
\emph{Escolio.}
Sea $Y$ un preesquema noetheriano, sea $\mathcal S$ una $\mathcal O_Y$-álgebra graduada que satisface las hipótesis de \sref{III.2.3.1}, y sea $X=\operatorname{Proj}(\mathcal S)$.
Sea $\boldsymbol{K}_{\mathcal S}$ la categoría abeliana de los $\mathcal S$-módulos graduados cuasicoherentes que satisfacen \textbf{(TF)}, y sea $\boldsymbol{K}'_{\mathcal S}$ la subcategoría de $\boldsymbol{K}_{\mathcal S}$ formada por los $\mathcal S$-módulos que satisfacen \textbf{(TN)}.
Por último, sea $\boldsymbol{K}_X$ la categoría de los $\mathcal S_X$-módulos graduados cuasicoherentes $\mathcal F$ de tipo finito; como $\mathcal S_X$ es periódica \sref[II]{II.8.14.4} y \sref[II]{II.8.14.12}, esto equivale a exigir que los $\mathcal F_i$ sean $\mathcal O_X$-módulos coherentes.

Por \sref[II]{II.8.14.8}, \sref[II]{II.8.14.10} y \sref{III.2.3.2}, los funtores
\[
\mathcal M\leadsto\shProj(\mathcal M)
\quad\text{en }\boldsymbol{K}_{\mathcal S}
\quad\text{y}\quad
\mathcal F\leadsto\bbGamma_*(\mathcal F)
\quad\text{en }\boldsymbol{K}_X
\] definen, en virtud de (\textbf{II}, 8.14.8 y 8.14.10) y (\hyperref[III.2.3.2-fr]{2.3.2}), \emph{una equivalencia} (T, I, 1.2) \emph{de la categoría cociente} $\boldsymbol{K}_{\mathcal S}/\boldsymbol{K}'_{\mathcal S}$
(T, I, 1.11) \emph{y la categoría} $\boldsymbol{K}_X$.
Cuando $\mathcal S$ está engendrada por $\mathcal S_1$, la categoría $\boldsymbol{K}_X$ puede sustituirse por la categoría de los $\mathcal O_X$-módulos coherentes \sref[II]{II.8.14.12}.
\end{env}

\begin{proposition}[2.3.4]
\label{III.2.3.4-fr}
Sea $Y$ un preesquema noetheriano.
\begin{enumerate}
\item[{\rm(i)}] Sea $\mathcal S$ una $\mathcal O_Y$-álgebra graduada en grados positivos, cuasicoherente de tipo finito. Sean $X=\operatorname{Proj}(\mathcal S)$, y
\[
\mathcal S_X=\shProj(\mathcal S)
=\bigoplus_{n\in\mathbb Z}\mathcal O_X(n).
\]
Entonces $\mathcal S_X$ es una $\mathcal O_X$-álgebra graduada periódica \textup{(\textbf{II}, 8.14.12)} cuyos componentes homogéneos $(\mathcal S_X)_n=\mathcal O_X(n)$ son $\mathcal O_X$-módulos coherentes, y si $d>0$ es un período de $\mathcal S_X$, $(\mathcal S_X)_d=\mathcal O_X(d)$ es un $\mathcal O_X$-módulo invertible $Y$-amplio. Además, el homomorfismo canónico
\[
\alpha:\mathcal S\longrightarrow\bbGamma_*(\mathcal S_X)
\]
es un \textup{\textbf{(TN)}}-isomorfismo.
\item[{\rm(ii)}] Recíprocamente, sea $q:X\to Y$ un morfismo proyectivo, y sea $\mathcal S'$ una $\mathcal O_X$-álgebra graduada, cuyos componentes homogéneos $\mathcal S'_n$ $(n\in\mathbb Z)$ son $\mathcal O_X$-módulos coherentes, y que admite un período $d>0$ tal que $\mathcal S'_d$ sea un $\mathcal O_X$-módulo invertible amplio para $q$. Entonces
\[
\mathcal S=\bigoplus_{n\geq0}q_*(\mathcal S'_n)
\]
es una $\mathcal O_Y$-álgebra graduada en grados positivos cuasicoherente y de tipo finito, y existe un $Y$-isomorfismo
\[
r:X\xrightarrow{\sim}\operatorname{Proj}(\mathcal S)
\]
tal que $r^*(\shProj(\mathcal S))$ sea isomorfo (como $\mathcal O_X$-álgebra graduada) a $\mathcal S'$.
\end{enumerate}
\end{proposition}

\begin{proof}
(i) Todas las aserciones ya han sido prácticamente demostradas, no siendo la última sino un caso particular de (\hyperref[III.2.3.2-fr]{2.3.2}). El hecho de que $\mathcal S_X$ es periódico se ha visto en (\textbf{II}, 8.14.14) y el hecho de que hay un período $d>0$ tal que $\mathcal O_X(d)$ sea invertible y $Y$-amplio no es sino (\textbf{II}, 4.6.18). Finalmente, para $0\leq k<d$, $(\mathcal S_X)^{(d,k)}$ es un $(\mathcal S_X)^{(d)}$-módulo de tipo finito (\textbf{II}, 8.14.14), por tanto cada uno de los $(\mathcal S_X)_n$ es un $\mathcal O_X$-módulo cuasicoherente de tipo finito en virtud de (\textbf{II}, 2.1.6, (ii)), siendo la cuestión local; como $\mathcal O_X$ es coherente, lo mismo ocurre con los $(\mathcal S_X)_n$.

(ii) Reemplazando, si es preciso, el período $d$ por uno de sus múltiplos, se puede suponer que $\mathcal L=\mathcal S'_d$ es un $\mathcal O_X$-módulo \emph{muy amplio} relativamente a $q$ (\textbf{II}, 4.6.11). Se tiene además
\[
\mathcal S'^{(d)}=\bigoplus_{n\in\mathbb Z}\mathcal L^{\otimes n}
\quad\text{por hipótesis, por tanto}\quad
\mathcal S^{(d)}=\bigoplus_{n\geq0}q_*(\mathcal L^{\otimes n})~;
\]
se sabe (\textbf{II}, 3.1.8 y 3.2.9) que hay un $Y$-isomorfismo $s$ de $X'=\operatorname{Proj}(\mathcal S)$ sobre $X''=\operatorname{Proj}(\mathcal S^{(d)})$ tal que
\oldpage[III]{105}
\[
s^*(\mathcal O_{X''}(n))=\mathcal O_{X'}(nd).
\]
Se establecerá, pues, la existencia de un $Y$-isomorfismo $X\xrightarrow{\sim}X'$ si se demuestra la

\begin{proposition}[2.3.4.1]
\label{III.2.3.4.1-fr}
Sean $Y$ un preesquema noetheriano, $q:X\to Y$ un morfismo proyectivo, $\mathcal L$ un $\mathcal O_X$-módulo invertible muy amplio para $q$. Entonces $\mathcal S=\bigoplus_{n\geq0}q_*(\mathcal L^{\otimes n})$ es una $\mathcal O_Y$-álgebra graduada cuasicoherente de tipo finito, tal que $\mathcal S_n=\mathcal S_1^n$ para $n$ suficientemente grande, y existe un $Y$-isomorfismo $r:X\xrightarrow{\sim}P=\operatorname{Proj}(\mathcal S)$ tal que $\mathcal L=r^*(\mathcal O_P(1))$.
\end{proposition}

Como $q$ es un morfismo proyectivo, resulta de (\textbf{II}, 5.4.4 y 4.4.7) que existe un $Y$-isomorfismo $r':X\xrightarrow{\sim}P'=\operatorname{Proj}(\mathcal T)$, donde $\mathcal T$ es una $\mathcal O_Y$-álgebra cuasicoherente tal que $\mathcal T_1$ sea un $\mathcal O_Y$-módulo de tipo finito y engendre $\mathcal T$, y se tiene $\mathcal L=r'^*(\mathcal O_{P'}(1))$. Se tiene entonces
\[
\mathcal S=\bigoplus_{n\geq0}q'_*(\mathcal O_{P'}(n)),
\]
donde $q':P'\to Y$ es el morfismo estructural, y resulta de (\hyperref[III.2.3.1-fr]{2.3.1}) que para $n$ suficientemente grande, el homomorfismo canónico $\alpha_n:\mathcal T_n\to\mathcal S_n=q_*(\mathcal L^{\otimes n})$ es \emph{biyectivo}; como $\mathcal T_n=\mathcal T_1^n$, se tiene \emph{a fortiori} $\mathcal S_n=\mathcal S_1^n$ en cuanto $n$ es suficientemente grande. Además, como el homomorfismo canónico $\alpha:\mathcal T\to\mathcal S$ de $\mathcal O_Y$-álgebras graduadas es un \textbf{(TN)}-isomorfismo,
\phantomsection\label{III.2.3.5.1-fr}
\[
\Phi=\operatorname{Proj}(\alpha):
\operatorname{Proj}(\mathcal S)\longrightarrow
\operatorname{Proj}(\mathcal T)
\]
es un isomorfismo (\textbf{II}, 3.6.1) y se tiene $\Phi_*(\widetilde{\mathcal S}(n))=(\mathcal S(n))\supertilde_{[\alpha]}$ (\textbf{II}, 3.5.2); pero como los $\mathcal T$-módulos graduados $(\mathcal S(n))_{[\alpha]}$ y $\mathcal T(n)$ son \textbf{(TN)}-isomorfos, se tiene \[
\Phi_*(\mathcal O_P(n))=\mathcal O_{P'}(n)
\] para todo $n$ (\textbf{II}, 3.4.2); para terminar de demostrar (\hyperref[III.2.3.4.1-fr]{2.3.4.1}), queda por mostrar que $\mathcal S$ es una $\mathcal O_Y$-álgebra de tipo finito; ahora bien, los $\mathcal S_n=q'_*(\mathcal O_{P'}(n))$ son $\mathcal O_Y$-módulos \emph{coherentes} en virtud de (\hyperref[III.2.2.1-fr]{2.2.1}) y como $\mathcal S_1^n=\mathcal S_n$ para $n\geq n_0$, $\mathcal S$ está engendrada por $\bigoplus_{i\leq n_0}\mathcal S_i$, que es coherente, de donde nuestra aserción (\textbf{I}, 9.6.2).

Volvamos a la demostración de (\hyperref[III.2.3.4-fr]{2.3.4}), cuyas notaciones retomamos. Hemos demostrado la existencia de un $Y$-isomorfismo $r'':X\xrightarrow{\sim}X''$ tal que $r''_*(\mathcal L^{\otimes n})=\mathcal O_{X''}(n)$ para todo $n\in\mathbb Z$; denotaremos por $q''$ el morfismo estructural $X''\to Y$. Observemos ahora que $\mathcal S'$ es suma directa de los $\mathcal S'^{(d)}$-módulos graduados $\mathcal S'^{(d,k)}$; cada uno de estos últimos es un $\mathcal S'^{(d)}$-módulo cuasicoherente de tipo finito, en virtud de la periodicidad de $\mathcal S'$ y de la hipótesis de que los $\mathcal S'_n$ son $\mathcal O_X$-módulos de tipo finito (\textbf{II}, 8.14.12). Pongamos $\mathcal F^{(k)}=r''_*(\mathcal S'^{(d,k)})$, de modo que los $\mathcal F^{(k)}$ son $\mathcal S_{X''}$-módulos graduados cuasicoherentes de tipo finito; por consiguiente (\textbf{II}, 8.14.8), el homomorfismo canónico
\[
\beta:\shProj(\bbGamma_*(\mathcal F^{(k)}))
\longrightarrow\mathcal F^{(k)}
\]
es un isomorfismo de $\mathcal S_{X''}$-módulos. Pero se tiene $q''_*((\mathcal F^{(k)})_n)=q_*((\mathcal S'^{(d,k)})_n)$ y para $n\geq0$, este último $\mathcal O_Y$-módulo es por definición igual a $(\mathcal S^{(d,k)})_n$. Es decir, la inyección canónica $\mathcal S^{(d,k)}\to\bbGamma_*(\mathcal F^{(k)})$ es un \textbf{(TN)}-isomorfismo, por tanto (\textbf{II}, 3.4.2) se tiene \[
\shProj(\mathcal S^{(d,k)})
=\shProj(\bbGamma_*(\mathcal F^{(k)})),
\] y por consiguiente $r''{}^*(\shProj(\mathcal S^{(d,k)}))=\mathcal S'^{(d,k)}$. Basta observar que \[
\shProj(\mathcal S^{(d,k)})
=s_*((\shProj(\mathcal S))^{(d,k)})
\] salvo un isomorfismo canónico (\textbf{II}, 8.14.13.1) para haber demostrado la isomorfía de $r^*(\shProj(\mathcal S))$ y de $\mathcal S'$. Finalmente, en virtud de (\hyperref[III.2.3.2-fr]{2.3.2}), cada uno de los $\bbGamma_*(\mathcal F^{(k)})$ satisface la condición \textbf{(TF)}, por tanto lo mismo ocurre con cada uno de los $\mathcal S^{(d,k)}$; además, como los $\mathcal S'_n$ son coherentes, lo mismo ocurre con los $\mathcal S_n=q_*(\mathcal S'_n)$ por (\hyperref[III.2.2.1-fr]{2.2.1}), y se concluye inmediatamente que los $\mathcal S^{(d,k)}$ son $\mathcal S^{(d)}$-módulos de tipo finito. Como se ha visto en (\hyperref[III.2.3.4.1-fr]{2.3.4.1}) que $\mathcal S^{(d)}$ es una $\mathcal O_Y$-álgebra de tipo finito, se concluye que $\mathcal S$ es también una $\mathcal O_Y$-álgebra de tipo finito.
\end{proof}

\oldpage[III]{106}

\begin{proposition}[2.3.5]
\label{III.2.3.5-fr}
Sean $Y$ un preesquema íntegro noetheriano, $X$ un preesquema íntegro, $f:X\to Y$ un morfismo \textup{proyectivo birracional}. Existe entonces un ideal fraccionario coherente $\mathcal J\subset\mathcal R(Y)$ \textup{(\textbf{II}, 8.1.2)} tal que $X$ sea $Y$-isomorfo al preesquema obtenido por explosión de $\mathcal J$ \textup{(\textbf{II}, 8.1.3)}. Además, existe un abierto $U$ de $Y$ tal que la restricción de $f$ a $f^{-1}(U)$ sea un isomorfismo de $f^{-1}(U)$ sobre $U$ \textup{(cf. \textbf{I}, 6.5.5)}, y que $\mathcal J|U$ sea invertible.
\end{proposition}

\begin{proof}
Como existe un $\mathcal O_X$-módulo invertible $\mathcal L$ muy amplio para $f$ (\textbf{II}, 4.4.2 y 5.3.2), se puede aplicar (\hyperref[III.2.3.4.1-fr]{2.3.4.1}), y se ve que $X$ se identifica con $\operatorname{Proj}(\mathcal S)$, donde $\mathcal S=\bigoplus_{n\geq0}f_*(\mathcal L^{\otimes n})$. Se sabe además que los $f_*(\mathcal L^{\otimes n})$ son $\mathcal O_Y$-módulos sin torsión (\textbf{I}, 7.4.5), por tanto lo mismo ocurre con el $\mathcal O_Y$-módulo $\mathcal S$, y por consiguiente $\mathcal S$ se identifica canónicamente con un sub-$\mathcal O_Y$-módulo de $\mathcal S\otimes_{\mathcal O_Y}\mathcal R(Y)$ (\textbf{I}, 7.4.1); este último es un haz \emph{simple} (\textbf{I}, 7.3.6) que queda determinado cuando se conoce su restricción a un abierto no vacío, por ejemplo a un abierto no vacío $U'\subset U$ tal que $\mathcal L|f^{-1}(U')$ sea isomorfo a $\mathcal O_X|f^{-1}(U')$. Como por hipótesis los $f_*(\mathcal L^{\otimes n})|U'$ son entonces isomorfos a $\mathcal O_Y|U'$, se ve que $\mathcal S\otimes_{\mathcal O_Y}\mathcal R(Y)$ es un $\mathcal R(Y)$-módulo isomorfo a $\mathcal R(Y)[T]$, donde $T$ es una indeterminada, y $\mathcal S$ es \textbf{(TN)}-isomorfo a la sub-$\mathcal O_Y$-álgebra engendrada por la imagen canónica de $f_*(\mathcal L)$ en $\mathcal S\otimes_{\mathcal O_Y}\mathcal R(Y)$ (\hyperref[III.2.3.4.1-fr]{2.3.4.1}); pero si se identifica $\mathcal S\otimes_{\mathcal O_Y}\mathcal R(Y)$ con $\mathcal R(Y)[T]$, la imagen de $f_*(\mathcal L)$ se identifica con $\mathcal J.T$, donde $\mathcal J$ es un sub-$\mathcal O_Y$-módulo coherente (\hyperref[III.2.2.1-fr]{2.2.1}) de $\mathcal R(Y)$, cuya restricción a $U'$ es isomorfa a $\mathcal O_Y|U'$, y que por consiguiente es tal que $\mathcal J|U$ sea invertible. Se ve entonces que $\mathcal S$ es \textbf{(TN)}-isomorfo a $\bigoplus_{n\geq0}\mathcal J^n$, lo que termina la demostración.
\end{proof}

\begin{corollary}[2.3.6]
\label{III.2.3.6}
\label{III.2.3.6-fr}
Bajo las hipótesis de \sref{III.2.3.5}, supongamos además que para todo sub-$\mathcal O_Y$-Módulo coherente
$\mathcal J\neq0$ de $\mathcal R(Y)$, existe un
$\mathcal O_Y$-módulo invertible $\mathcal L$ tal que
\[
\Gamma(Y,\mathcal L\otimes_{\mathcal O_Y}
\mathcal{H}om(\mathcal J,\mathcal O_Y))\neq0~;
\]
Entonces, en el enunciado de \sref{III.2.3.5}, puede tomarse $\mathcal J$ como un ideal de $\mathcal O_Y$.
Esta condición adicional se satisface siempre si existe un $\mathcal O_Y$-módulo amplio.
\end{corollary}

\begin{proof}
En efecto, por \sref[0]{0.5.4.2},
\[
\mathcal L\otimes\mathcal{H}om(\mathcal J,\mathcal O_Y)
=\mathcal{H}om(\mathcal L^{-1},\mathcal{H}om(\mathcal J,\mathcal O_Y))
=\mathcal{H}om(\mathcal J\otimes\mathcal L^{-1},\mathcal O_Y)~;
\]
la hipótesis significa pues que hay un homomorfismo no nulo $u$ de
$\mathcal J\otimes\mathcal L^{-1}$ en $\mathcal O_Y$. Como, para todo
$y\in Y$, el módulo $(\mathcal J\otimes\mathcal L^{-1})_y$ se identifica con un sub-$\mathcal O_y$-módulo del cuerpo de fracciones $(\mathcal R(Y))_y$ de $\mathcal O_y$ \sref[I]{I.7.1.5}.
Así, $u_y$ es necesariamente inyectivo, y $u$ identifica $\mathcal J\otimes\mathcal L^{-1}$ con un ideal $\mathcal J'$ de $\mathcal O_Y$.
Pero
\[
\operatorname{Proj}\left(\bigoplus_{n\geq0}\mathcal J^n\right)
\quad\text{y}\quad
\operatorname{Proj}\left(
\bigoplus_{n\geq0}(\mathcal J\otimes\mathcal L^{-1})^n\right)
\]
son $Y$-isomorfos \sref[II]{II.3.1.8}; esto demuestra la primera aserción.

Para la segunda, observemos que
$\mathcal F=\mathcal{H}om_{\mathcal O_Y}(\mathcal J,\mathcal O_Y)$ es
coherente y $\neq0$, puesto que existe un abierto $U$ de $Y$ sobre el cual $\mathcal J|U$ es invertible.
Si $\mathcal L$ es un $\mathcal O_Y$-módulo amplio, existe un entero $n$ tal que
$\mathcal F(n)=\mathcal F\otimes\mathcal L^{\otimes n}$
está engendrado por sus secciones sobre $Y$ \sref[II]{II.4.5.5}.
\emph{A fortiori},
$\Gamma(Y,\mathcal F(n))\neq0$,
lo que demuestra la aserción.
\end{proof}

\begin{corollary}[2.3.7]
\label{III.2.3.7}
\label{III.2.3.7-fr}
Sean $X$ e $Y$ esquemas íntegros proyectivos sobre un cuerpo $k$, y sea
$f:X\to Y$
un $k$-morfismo birracional.
Entonces $X$ es $k$-isomorfo a un $Y$-esquema obtenido por explosión de un subesquema cerrado $Y'$ de $Y$, no necesariamente reducido.
\end{corollary}

\begin{proof}
\oldpage[III]{107}
En efecto, $f$ es proyectivo \sref[II]{II.5.5.5}[v], y como $Y$ es proyectivo sobre $k$, se satisface la condición adicional de \sref{III.2.3.6}.
Basta, por tanto, tomar el subesquema cerrado $Y'$ de $Y$ definido por el ideal coherente $\mathcal J$ de \sref{III.2.3.6}.
\end{proof}

\begin{env}[2.3.8]
\label{III.2.3.8}
\label{III.2.3.8-fr}
\emph{Observación.}
En el Capítulo~IV, al estudiar la noción de divisor, veremos que si, en el enunciado de \sref{III.2.3.5}, se supone que los anillos locales $\mathcal O_y$, para $(y\in Y)$, son factoriales---como ocurre, por ejemplo, cuando $Y$ es no singular---, entonces $X$ puede obtenerse de $Y$ por explosión de un subpreesquema cerrado $Y'$ de $Y$ cuyo espacio subyacente está contenido en $Y\setmin U$.
\end{env}


\subsection{Una generalización del teorema fundamental}
\label{subsection:III.2.4}
\label{subsection:III.2.4-fr}

\begin{theorem}[2.4.1]
\label{III.2.4.1-fr}
Sean $Y$ un preesquema noetheriano, $\mathcal S$ una
$\mathcal O_Y$-Álgebra cuasicoherente de tipo finito. Sean
$f:X\to Y$ un morfismo proyectivo, $\mathcal S'=f^*(\mathcal S)$,
$\mathcal M$ un $\mathcal S'$-Módulo cuasicoherente de tipo finito.
Entonces:
\begin{enumerate}
\item[\textup{(i)}] Para todo $p\in\mathbb Z$, $R^pf_*(\mathcal M)$ es un
$\mathcal S$-Módulo de tipo finito.
\item[\textup{(ii)}] Sea además $\mathcal L$ un $\mathcal O_X$-Módulo invertible
amplio para $f$, y pongamos
$\mathcal M(n)=\mathcal M\otimes\mathcal L^{\otimes n}$ para todo
$n\in\mathbb Z$. Existe un entero $N$ tal que, para $n\geq N$, se tiene
\[
R^pf_*(\mathcal M(n))=0
\tag{2.4.1.1}\label{III.2.4.1.1-fr}
\]
para todo $p>0$, y el homomorfismo canónico
$f^*(f_*(\mathcal M(n))\longrightarrow\mathcal M(n)$
\textup{(\textbf{0}\textsubscript{I}, 4.4.3)} es sobreyectivo.
\end{enumerate}
\end{theorem}

\begin{proof}
Pongamos $Y'=\operatorname{Spec}(\mathcal S)$,
$X'=\operatorname{Spec}(\mathcal S')$ de modo que
$X'=X\times_Y Y'$ (\textbf{II}, 1.5.5); sean $g:Y'\to Y$,
$g':X'\to X$ los morfismos estructurales, que son afines por
definición, y $f'=f_{(Y')}:X'\to Y'$; se tiene pues un diagrama
conmutativo
\[
\xymatrix{
X\ar[d]_f & X'\ar[l]_{g'}\ar[d]^{f'}\ar[dl]^{h}\\
Y & Y'\ar[l]^{g}
}
\]
y el morfismo $f'$ es proyectivo (\textbf{II}, 5.5.5, (iii)); pongamos
$h=f\circ g'=g\circ f'$.

(i) Sea $\widetilde{\mathcal M}$ el $\mathcal O_{X'}$-Módulo asociado al
$\mathcal S'$-Módulo cuasicoherente $\mathcal M$, cuando $X'$ se
considera como un $X$-esquema afín (\textbf{II}, 1.4.3), de modo que se tiene
$\mathcal M=g'_*(\widetilde{\mathcal M})$; como $\mathcal M$ es un
$\mathcal S'$-Módulo de tipo finito, $\widetilde{\mathcal M}$ es un
$\mathcal O_{X'}$-Módulo de tipo finito (\textbf{II}, 1.4.5); como $h$ es de tipo
finito, puesto que $g$ y $f'$ lo son (\textbf{II}, 1.3.7 y \textbf{I}, 6.3.4, (ii)), $X'$ es
noetheriano (\textbf{I}, 6.3.7) y $\widetilde{\mathcal M}$ es por tanto coherente.
Dicho esto, como $g'$ es afín, el homomorfismo canónico
$R^pf_*(\mathcal M)\to R^ph_*(\widetilde{\mathcal M})$ es biyectivo
(\hyperref[III.1.3.4-fr]{1.3.4}). Además, este homomorfismo es un
homomorfismo de $\mathcal S$-\emph{Módulos}; en efecto, del homomorfismo
canónico
\[
g'_*(\mathcal O_{X'})\otimes_{\mathcal O_X}
g'_*(\widetilde{\mathcal M})\longrightarrow
g'_*(\widetilde{\mathcal M})
\tag{2.4.1.2}\label{III.2.4.1.2-fr}
\]
que define la estructura de $\mathcal S'$-Módulo de $\mathcal M$ (recordando
que $\mathcal S'=g'_*(\mathcal O_{X'})$), se deduce
canónicamente un homomorfismo
\[
f_*(g'_*(\mathcal O_{X'}))\otimes
R^pf_*(g'_*(\widetilde{\mathcal M}))\longrightarrow
R^pf_*(g'_*(\widetilde{\mathcal M}))
\]

\oldpage[III]{108}
(\textbf{0}, 12.2.2), y como (\hyperref[III.2.4.1.2-fr]{2.4.1.2}) procede
a su vez (por aplicación de (\textbf{0}\textsubscript{I}, 4.2.2.1)) del
homomorfismo
$\mathcal O_{X'}\otimes\widetilde{\mathcal M}
\longrightarrow\widetilde{\mathcal M}$
que define la estructura de $\mathcal O_{X'}$-Módulo de
$\widetilde{\mathcal M}$, el diagrama
\[
\xymatrix{
f_*(g'_*(\mathcal O_{X'}))\otimes
R^pf_*(g'_*(\widetilde{\mathcal M}))
\ar[r]\ar[d] &
R^pf_*(g'_*(\widetilde{\mathcal M}))\ar[d]\\
h_*(\mathcal O_{X'})\otimes R^ph_*(\widetilde{\mathcal M})
\ar[r] & R^ph_*(\widetilde{\mathcal M})
}
\]
es conmutativo (\textbf{0}, 12.2.6); componiendo las flechas horizontales con
el homomorfismo procedente del homomorfismo canónico
$\mathcal S\longrightarrow f_*(f^*(\mathcal S))=f_*(\mathcal S')
=f_*(g'_*(\mathcal O_{X'}))=h_*(\mathcal O_{X'})$,
se obtiene nuestra afirmación. Por otra parte, puesto que $g$ es afín y
$f'$ es separado y cuasicompacto, el homomorfismo canónico
$R^ph_*(\widetilde{\mathcal M})\to
g_*(R^pf'_*(\widetilde{\mathcal M}))$ es biyectivo
(\hyperref[III.1.4.14-fr]{1.4.14}), y se demuestra como antes que
es un isomorfismo de $\mathcal S$-\emph{Módulos} (utilizando esta vez la
conmutatividad de (\textbf{0}, 12.2.6.2)). Ahora bien, al ser $f'$ proyectivo y
$\widetilde{\mathcal M}$ coherente,
$R^pf'_*(\widetilde{\mathcal M})$ es un
$\mathcal O_{Y'}$-Módulo coherente en virtud de
(\hyperref[III.2.2.1-fr]{2.2.1}); se concluye que
$g_*(R^pf'_*(\widetilde{\mathcal M}))$ es un $\mathcal S$-Módulo \emph{de tipo
finito} (\textbf{II}, 1.4.5).

(ii) Sea $\mathcal L'=g'^*(\mathcal L)$, que es un
$\mathcal O_{X'}$-Módulo invertible; para todo $n\in\mathbb Z$, se tiene
$g'_*(\widetilde{\mathcal M}\otimes\mathcal L'^{\otimes n})
=g'_*(\widetilde{\mathcal M})\otimes\mathcal L^{\otimes n}
=\mathcal M(n)$
(\textbf{0}\textsubscript{I}, 5.4.10) salvo isomorfismo; se puede aplicar
a $\widetilde{\mathcal M}\otimes\mathcal L'^{\otimes n}$ el
razonamiento hecho en (i) para $\widetilde{\mathcal M}$, que prueba que
$R^pf_*(g'_*(\widetilde{\mathcal M}\otimes\mathcal L'^{\otimes n}))$
es isomorfo a
$g_*(R^pf'_*(\widetilde{\mathcal M}\otimes\mathcal L'^{\otimes n}))$.
Ahora bien, $\mathcal L'$ es amplio para $f'$ (\textbf{II}, 4.6.13, (iii)), luego resulta
de (\hyperref[III.2.2.1-fr]{2.2.1}) que existe un entero $N$ tal que
$R^pf'_*(\widetilde{\mathcal M}\otimes\mathcal L'^{\otimes n})=0$
para todo $p$ y todo $n\geq N$, lo que prueba
(\hyperref[III.2.4.1.1-fr]{2.4.1.1}). Finalmente, resulta también de
(\hyperref[III.2.2.1-fr]{2.2.1}) que se puede suponer $N$ tal que para
$n\geq N$ el homomorfismo canónico
$f'^*(f'_*(\widetilde{\mathcal M}\otimes\mathcal L'^{\otimes n}))
\longrightarrow\widetilde{\mathcal M}\otimes\mathcal L'^{\otimes n}$
sea sobreyectivo; como $g'_*$ es un funtor exacto (\textbf{II}, 1.4.4),
el homomorfismo correspondiente
\[
g'_*(f'^*(f'_*(\widetilde{\mathcal M}\otimes
\mathcal L'^{\otimes n})))
\longrightarrow
g'_*(\widetilde{\mathcal M}\otimes\mathcal L'^{\otimes n})
=\mathcal M(n)
\]
es sobreyectivo. Ahora bien, se tiene
$g'_*(f'^*(f'_*(\widetilde{\mathcal M}\otimes
\mathcal L'^{\otimes n})))
=f^*(g_*(f'_*(\widetilde{\mathcal M}\otimes
\mathcal L'^{\otimes n}))$
(\textbf{II}, 1.5.2) y como $g_*\circ f'_* = f_*\circ g'_*$, se ve
finalmente que se \textbf{tiene}
\[
g'_*(f'^*(f'_*(\widetilde{\mathcal M}\otimes\mathcal L'^{\otimes n})))
=f^*(f_*(g'_*(\widetilde{\mathcal M}\otimes
\mathcal L'^{\otimes n})))=f^*(f_*(\mathcal M(n))),
\]
lo que termina la demostración.
\end{proof}

\begin{env}[2.4.2]
\label{III.2.4.2}
\label{III.2.4.2-fr}
Aplicaremos \sref{III.2.4.1}, en particular, cuando $\mathcal S$ sea una $\mathcal O_Y$-álgebra \emph{graduada} en grados positivos y
$\mathcal M=\bigoplus_{\boldsymbol{k}\in\mathbb Z}\mathcal M_{\boldsymbol{k}}$
sea un $\mathcal S'$-módulo \emph{graduado}.
\oldpage[III]{109}
Bajo las mismas hipótesis de finitud sobre $\mathcal S$ y $\mathcal M$, se deduce de \sref{III.2.4.1} que
$\bigoplus_{k\in\mathbb Z}R^pf_*(\mathcal M_k)$
es un $\mathcal S$-módulo \emph{de tipo finito} para todo $p$.
Bajo las hipótesis adicionales de \sref{III.2.4.1}[ii], existe además un entero $N$ tal que, para $n\geq N$,
$R^pf_*(\mathcal M_k(n))=0$
para todo $p>0$ y \emph{todo} $k\in\mathbb Z$, y el homomorfismo canónico
$f^*(f_*(\mathcal M_k(n)))\to\mathcal M_{\boldsymbol{k}}(n)$
es sobreyectivo para \emph{todo} $k\in\mathbb Z$.
\end{env}


\subsection{Característica de Euler--Poincar\'e y polinomio de Hilbert}
\label{subsection:III.2.5}
\label{subsection:III.2.5-fr}

\begin{env}[2.5.1]
\label{III.2.5.1}
\label{III.2.5.1-fr}
Sea $A$ un anillo \emph{artiniano} y sea $X$ un $A$-esquema \emph{proyectivo} sobre
$Y=\operatorname{Spec}(A)$.
Para todo $\mathcal O_X$-módulo coherente $\mathcal F$, los $H^i(X,\mathcal F)$ $(i\geq0)$ son $A$-módulos de
tipo finito (\hyperref[III.2.2.1-fr]{2.2.1}), luego aquí \emph{de longitud finita}
puesto que $A$ es artiniano.
Además, \sref{III.2.2.1} muestra que $H^i(X,\mathcal F)=0$ salvo para un número finito de $i\geq0$.
Así, el entero
\[
\chi_A(\mathcal F)=\sum_{i=0}^{\infty}(-1)^i
\operatorname{long}(H^i(X,\mathcal F))
\tag{2.5.1.1}\label{III.2.5.1.1-fr}
\]
está definido para todo $\mathcal O_X$-módulo coherente $\mathcal F$.
Cuando $A$ es un \emph{anillo local artiniano}, $\chi_A(\mathcal F)$ se llama la \emph{característica de Euler--Poincar\'e de} $\mathcal F$ \emph{con respecto a} $A$.
Para $\mathcal F=\mathcal O_X$, el entero $\chi_A(\mathcal O_X)$ se llama el \emph{género aritmético de} $X$ con respecto a $A$.
\end{env}

\begin{proposition}[2.5.2]
\label{III.2.5.2}
\label{III.2.5.2-fr}
Sea
$0\longrightarrow\mathcal F'\longrightarrow\mathcal F
\longrightarrow\mathcal F''\longrightarrow0$
una sucesión exacta de $\mathcal O_X$-módulos coherentes.
Entonces
\[
\chi_A(\mathcal F)=\chi_A(\mathcal F')+\chi_A(\mathcal F'').
\tag{2.5.2.1}\label{III.2.5.2.1-fr}
\]
\end{proposition}

\begin{proof}
Los módulos de cohomología de $\mathcal F$, $\mathcal F'$ y $\mathcal F''$ se anulan en todos los grados suficientemente grandes.
Existe, por tanto, un entero $r>0$ para el cual la sucesión exacta larga de cohomología tiene la forma
\[
0\to H^0(X,\mathcal F')\to H^0(X,\mathcal F)
\to H^0(X,\mathcal F'')\to H^1(X,\mathcal F')\to\ldots
\]
\[
\ldots\to H^r(X,\mathcal F')\to H^r(X,\mathcal F)
\to H^r(X,\mathcal F'')\to0.
\]
En una sucesión exacta de $A$-módulos de longitud finita con términos nulos en ambos extremos, la suma alternada de las longitudes es nula \sref[0]{0.11.10.1}.
La aplicación de este hecho da inmediatamente \eqref{III.2.5.2.1}.
\end{proof}

El resultado de \sref{III.2.5.2} se aplica siempre que $X$ sea un $A$-esquema cuasicompacto y que los $A$-módulos $H^i(X,\mathcal F)$ sean \emph{de tipo finito} para todo $\mathcal O_X$-módulo coherente $\mathcal F$ \sref{III.1.4.12}.

\begin{theorem}[2.5.3]
\label{III.2.5.3-fr}
Sean $A$ un anillo local artiniano, $X$ un esquema proyectivo sobre
$Y=\operatorname{Spec}(A)$, $\mathcal L$ un $\mathcal O_X$-Módulo
invertible muy amplio con respecto a $Y$, $\mathcal F$ un
$\mathcal O_X$-Módulo coherente; se pone
$\mathcal F(n)=\mathcal F\otimes_{\mathcal O_X}\mathcal L^{\otimes n}$
para todo $n\in\mathbb Z$.
\begin{enumerate}
\item[\textup{(i)}] Existe un polinomio único $P\in\mathbb Q[T]$ tal que
$\chi_A(\mathcal F(n))=P(n)$ para todo $n\in\mathbb Z$ \textup{(se dice que $P$
es el \emph{polinomio de Hilbert} de $\mathcal F$ con respecto a $A$)}.
\item[\textup{(ii)}] Para $n$ suficientemente grande, se tiene
$\chi_A(\mathcal F(n))=\operatorname{long}_A
\Gamma(X,\mathcal F(n))$.
\item[\textup{(iii)}] El término de mayor grado de $\chi_A(\mathcal F(n))$ tiene un
coeficiente $\geq0$.
\end{enumerate}
\end{theorem}

Añadamos que en el cap. IV, en el parágrafo dedicado a la noción de
dimensión, demostraremos además que \emph{el grado de
$\chi_A(\mathcal F(n))$ es igual a la dimensión del soporte de}
$\mathcal F$.

\begin{proof}
Como se tiene $H^i(X,\mathcal F(n))=0$ para todo $i>0$ en cuanto $n$ es
suficientemente grande (\hyperref[III.2.2.1-fr]{2.2.1}),

\oldpage[III]{110}
se tiene
$\chi_A(\mathcal F(n))=\operatorname{long}H^0(X,\mathcal F(n))
=\operatorname{long}\Gamma(X,\mathcal F(n))$
para $n$ suficientemente grande, de donde (ii); esto implica
$\chi_A(\mathcal F(n))\geq0$ para $n$ suficientemente grande, y (iii) resulta pues
de (i); como, por lo demás, la afirmación de unicidad de (i) es inmediata,
queda por probar la existencia del polinomio $P$.

Mostremos primero que se puede suponer $\mathfrak m\mathcal F=0$, donde
$\mathfrak m$ es el ideal maximal de $A$. En efecto, existe un entero
$s>0$ tal que $\mathfrak m^s=0$, y $\mathcal F(n)$ admite por tanto una
filtración finita
\[
\mathcal F(n)\supset\mathfrak m\mathcal F(n)\supset\ldots
\supset\mathfrak m^{s-1}\mathcal F(n)\supset0.
\]
Por inducción, se deduce de
(\hyperref[III.2.5.2.1-fr]{2.5.2.1}) que
\[
\chi_A(\mathcal F(n))=\sum_{k=1}^{s}
\chi_A(\mathfrak m^{k-1}\mathcal F(n)/\mathfrak m^k\mathcal F(n))~;
\]
como
$\mathfrak m^{k-1}\mathcal F(n)/\mathfrak m^k\mathcal F(n)
=\mathcal F'_k(n)$, donde
$\mathcal F'_k=\mathfrak m^{k-1}\mathcal F/\mathfrak m^k\mathcal F$,
esto prueba nuestra afirmación.

Supongamos pues $\mathfrak m\mathcal F=0$; si $X'$ es el subesquema
cerrado de $X$, imagen inversa por el morfismo estructural
$X\to\operatorname{Spec}(A)$ del único punto cerrado de
$\operatorname{Spec}(A)$, y $j:X'\to X$ la inclusión canónica, se tiene
$\mathcal F=j_*(\mathcal F')$, donde $\mathcal F'$ es un
$\mathcal O_{X'}$-Módulo coherente; $X'$ es un esquema proyectivo sobre
$\operatorname{Spec}(K)$, donde $K=A/\mathfrak m$. Si
$\mathcal L'=j^*(\mathcal L)$, $\mathcal L'$ es muy amplio con respecto
a $\operatorname{Spec}(K)$ (\textbf{II}, 4.4.10), y se tiene
$\mathcal F(n)=j_*(\mathcal F'(n))$, donde
$\mathcal F'(n)=\mathcal F'\otimes_{\mathcal O_{X'}}
\mathcal L'^{\otimes n}$ (\textbf{0}\textsubscript{I}, 5.4.10). Se concluye que
$\chi_A(\mathcal F(n))=\chi_K(\mathcal F'(n))$ (G, II, 4.9.1), y se está
pues reducido al caso en que $A$ es un \emph{cuerpo}.

Observemos ahora que $X$ puede considerarse como un subesquema
cerrado de $P=\mathbf P_A^r$ para un $r$ apropiado (\textbf{II}, 5.5.4, (ii));
si $i:X\to P$ es la inclusión canónica, se ve como antes que se
tiene $\chi_A(\mathcal F(n))=\chi_A(i_*(\mathcal F)(n))$, de modo que se
puede limitar al caso en que $X=\mathbf P_A^r=\operatorname{Proj}(S)$ con
$S=A[T_0,\ldots,T_r]$, siendo $A$ un cuerpo.

Dicho esto, se tiene $\mathcal F=\widetilde M$, donde $M$ es un $S$-módulo
graduado de tipo finito (\textbf{II}, 2.7.8); existe por tanto una resolución finita
de $M$ por $S$-módulos graduados \emph{libres} de tipo finito
\[
0\to L_q\to L_{q-1}\to\ldots\to L_1\to M\to0
\]
en virtud del teorema de las sicigias de Hilbert (M, VIII, 6.5); como
$\widetilde M$ es un funtor exacto en $M$ (\textbf{II}, 2.5.4), se tiene
también una sucesión exacta
\[
0\to\widetilde L_q\to\widetilde L_{q-1}\to\ldots
\to\widetilde L_1\to\widetilde M\to0
\]
y por tanto, para todo $n\in\mathbb Z$, la sucesión
\[
0\to\widetilde L_q(n)\to\widetilde L_{q-1}(n)\to\ldots
\to\widetilde L_1(n)\to\widetilde M(n)\to0
\]
es exacta; aplicando por inducción sobre $q$ la prop. (2.5.1), resulta
\[
\chi_A(\widetilde M(n))=\sum_{j=1}^{q}(-1)^{j+1}
\chi_A(\widetilde L_j(n))
\]
y para probar (i), se está pues reducido al caso en que $M$ es \emph{libre} y graduado
de tipo finito, luego al caso en que $M=S(h)$ para un $h\in\mathbb Z$. Como
entonces $\widetilde{\mathcal M}(n)=(M(n))^\sim=(S(n+h))^\sim$ (\textbf{II}, 2.5.15), se ve
finalmente que el teorema resultará del
\end{proof}

\oldpage[III]{111}
\begin{lemma}[2.5.3.1]
\label{III.2.5.3.1-fr}
Sean $A$ un cuerpo, $r$ un entero $>0$, y $X=\mathbf P_A^r$; se tiene
entonces
$\chi_A(\mathcal O_X(n))=\binom{n+r}{r}$
para todo $n\in\mathbb Z$.
\end{lemma}

\begin{proof}
En efecto, para $n>0$, se tiene
$\chi_A(\mathcal O_X(n))=\operatorname{long}H^0(X,\mathcal O_X(n))$, que
es el número de monomios en los $T_i$ de grado total $n$, es decir,
$\binom{n+r}{r}$ (\hyperref[III.2.1.12-fr]{2.1.12}). Para
$n\leq-r-1$, se tiene análogamente
$\chi_A(\mathcal O_X(n))=(-1)^r
\operatorname{long}H^r(X,\mathcal O_X(n))$;
si $n=-r-h$, la dimensión de $H^r(X,\mathcal O_X(n))$ sobre $A$ es el
número de sucesiones $(p_i)_{0\leq i\leq r}$ de enteros $p_i>0$ tales que
$\sum_{i=0}^{r}p_i=r+h$ (\hyperref[III.2.1.12-fr]{2.1.12}), o también el
número de sucesiones de enteros $q_i\geq0$ $(0\leq i\leq r)$ tales que
$\sum_{i=0}^{r}q_i=h-1$; es pues el número
$\binom{h+r-1}{r}=(-1)^r\binom{n+r}{r}$.
Finalmente, para $-r\leq n\leq0$, se tiene $\binom{n+r}{r}=0$ y por otra parte
$H^i(X,\mathcal O_X(n))=0$ para todo $i\geq0$
(\hyperref[III.2.1.12-fr]{2.1.12}), lo que prueba el lema.
\end{proof}

\begin{corollary}[2.5.4]
\label{III.2.5.4}
\label{III.2.5.4-fr}
Sea $A$ un anillo local artiniano, sea $S$ una $A$-álgebra graduada de tipo finito engendrada por $S_1$, sea $M$ un $S$-módulo graduado de tipo finito y sea
$X=\operatorname{Proj}(S)$. Se tiene entonces
$\chi_A(\widetilde M(n))=\operatorname{long}M_n$
para $n$ suficientemente grande.
\end{corollary}

\begin{proof}
Esto resulta de que $M_n$ y $\Gamma(X,\widetilde M(n))$ son
isomorfos para $n$ suficientemente grande
(\hyperref[III.2.3.1-fr]{2.3.1}).
\end{proof}

\subsection{Aplicación: criterios de amplitud}
\label{subsection:III.2.6-fr}

\begin{proposition}[2.6.1]
\label{III.2.6.1-fr}
Sean $Y$ un preesquema noetheriano, $f:X\to Y$ un morfismo propio,
$\mathcal L$ un $\mathcal O_X$-Módulo invertible. Las condiciones
siguientes son equivalentes:
\begin{enumerate}
\item[\textup{a)}] $\mathcal L$ es amplio para $f$.
\item[\textup{b)}] Para todo $\mathcal O_X$-Módulo coherente $\mathcal F$,
existe un entero $N$ tal que para $n\geq N$ se tiene
$R^qf_*(\mathcal F\otimes\mathcal L^{\otimes n})=0$
para todo $q>0$.
\item[\textup{c)}] Para todo Ideal coherente $\mathcal J$ de $\mathcal O_X$,
existe un entero $N$ tal que para $n\geq N$ se tiene
$R^1f_*(\mathcal J\otimes\mathcal L^{\otimes n})=0$.
\end{enumerate}
\end{proposition}

\begin{proof}
Se ha visto que \emph{a)} implica \emph{b)} (\hyperref[III.2.2.1-fr]{2.2.1}, (ii)). Es
trivial que \emph{b)} implica \emph{c)}, y queda por probar que \emph{c)} implica
\emph{a)}. Se puede limitar al caso en que $Y$ es afín (\textbf{II}, 4.6.4), y probar
en este caso que $\mathcal L$ es amplio; bastará mostrar que cuando
$h$ recorre el conjunto de secciones de los $\mathcal L^{\otimes n}$
$(n>0)$ sobre $X$, aquellos $X_h$ que son afines forman un
recubrimiento de $X$ (\textbf{II}, 4.5.2). Para ello, mostremos que para todo punto
\emph{cerrado} $x$ de $X$ y toda vecindad abierta afín $U$ de $x$, existe un
$n$ y un $h\in\Gamma(X,\mathcal L^{\otimes n})$ tales que
$x\in X_h\subset U$; $X_h$ será necesariamente afín (\textbf{I}, 1.3.6) y la
unión de estos $X_h$ será un abierto de $X$ que contiene todos los puntos
cerrados de $X$, y por tanto el propio $X$ puesto que $X$ es noetheriano
(\textbf{I}, 6.3.7 y \textbf{0}\textsubscript{I}, 2.1.3). Sea $\mathcal J$ (resp.
$\mathcal J'$) el haz cuasicoherente de ideales de $\mathcal O_X$
que define el subpreesquema cerrado reducido de $X$ cuyo espacio
subyacente es $X\setmin U$ (resp. $(X\setmin U)\cup\{x\}$) (\textbf{I}, 5.2.1); está claro que
$\mathcal J$ y $\mathcal J'$ son coherentes (\textbf{I}, 6.1.1), que
$\mathcal J'\subset\mathcal J$ y que
$\mathcal J''=\mathcal J/\mathcal J'$ es un $\mathcal O_X$-Módulo
coherente (\textbf{0}\textsubscript{I}, 5.3.3) cuyo soporte es $\{x\}$ y tal
que $\mathcal J''_x=\boldsymbol{k}(x)$. Como $\mathcal L$ es localmente libre, la
sucesión
$0\to\mathcal J'\otimes\mathcal L^{\otimes n}
\to\mathcal J\otimes\mathcal L^{\otimes n}
\to\mathcal J''\otimes\mathcal L^{\otimes n}\to0$
es exacta para todo $n$, y por hipótesis existe $n$ suficientemente grande tal
que $H^1(X,\mathcal J'\otimes\mathcal L^{\otimes n})=0$; la sucesión
exacta de cohomología
\oldpage[III]{112}
prueba pues que el homomorfismo
$\Gamma(X,\mathcal J\otimes\mathcal L^{\otimes n})
\to
\Gamma(X,\mathcal J''\otimes\mathcal L^{\otimes n})$
es \emph{sobreyectivo}. Una sección $g$ de
$\mathcal J''\otimes\mathcal L^{\otimes n}$ sobre $X$ tal que
$g(x)\neq0$ es pues la imagen de una sección
$h\in\Gamma(X,\mathcal J\otimes\mathcal L^{\otimes n})
\subset\Gamma(X,\mathcal L^{\otimes n})$ (pues en virtud de
(\textbf{0}\textsubscript{I}, 5.4.1),
$\mathcal J\otimes\mathcal L^{\otimes n}$ es un sub-$\mathcal O_X$-Módulo
de $\mathcal L^{\otimes n}$); se tiene por definición $h(x)\neq0$ y
$h(z)=0$ para $z\notin U$, lo que termina la demostración.
\end{proof}

\begin{proposition}[2.6.2]
\label{III.2.6.2-fr}
Sean $Y$ un preesquema noetheriano, $f:X\to Y$ un morfismo de tipo
finito, $g:X'\to X$ un morfismo finito sobreyectivo, $\mathcal L$ un
$\mathcal O_X$-Módulo invertible y $\mathcal L'=g^*(\mathcal L)$. Se
supone verificada la condición siguiente: Existe una parte $Z$ de
$X$, propia sobre $Y$ (\textbf{II}, 5.4.10) tal que para todo $x\in X\setmin Z$, o bien
$X$ es normal en el punto $x$, o bien $(g_*(\mathcal O_{X'}))_x$ es un
$\mathcal O_x$-módulo libre. En estas condiciones, para que $\mathcal L$
sea amplio para $f$, es necesario y suficiente que $\mathcal L'$ sea amplio
para $f\circ g$.
\end{proposition}

\begin{proof}
\begin{env}[2.6.2.1]
\label{III.2.6.2.1-fr}
Puesto que $g$ es afín, la condición es necesaria (\textbf{II}, 5.1.12). Para
ver que es suficiente, se puede suponer $Y$ afín (\textbf{II}, 4.6.4).
Mostremos además que se puede limitar al caso en que $X$ es \emph{reducido}. En
efecto, sea $j:X_{\mathrm{red}}\to X$ la inclusión canónica, y pongamos
$X_1=X_{\mathrm{red}}$, $X'_1=X'\times_X X_1$, de modo que se tiene el
diagrama conmutativo
\[
\xymatrix{
X'\ar[d]_g & X'_1\ar[l]_{j'}\ar[d]^{g_1\, .}\\
X & X_1\ar[l]^{j}
}
\tag{2.6.2.2}\label{III.2.6.2.2-fr}
\]

El morfismo $f\circ j$ es entonces de tipo finito (\textbf{I}, 6.3.4) y $g_1$ es
un morfismo finito (\textbf{II}, 6.1.5, (iii)); si $\mathcal L'$ es amplio para
$f\circ g$, $j'^*(\mathcal L')$ es amplio para $f\circ g\circ j'$ puesto que
$j'$ es una inmersión cerrada (\textbf{II}, 5.1.12 y \textbf{I}, 4.3.2). Si se pone
$Z_1=j^{-1}(Z)$, $Z_1$ es propio sobre $Y$ (\textbf{II}, 5.4.10); por otra parte,
si $X$ es normal en un punto $x$, lo mismo ocurre evidentemente con
$X_{\mathrm{red}}$; finalmente, si $(g_*(\mathcal O_{X'}))_x$ es un
$\mathcal O_x$-módulo libre, resulta inmediatamente de (\textbf{II}, 1.5.2) que
$((g_1)_*(\mathcal O_{X'_1}))_x$ es un $\mathcal O_{X_1,x}$-módulo
libre. Finalmente, como $X$ es noetheriano (\textbf{I}, 6.3.7), si
$j^*(\mathcal L)$ es amplio, $\mathcal L$ es amplio (\textbf{II}, 4.5.14), y
como $j'^*(\mathcal L')=g_1^*(j^*(\mathcal L))$, esto termina la
reducción anunciada. Suponemos pues en adelante $Y$ \emph{afín} y $X$
\emph{reducido}.
\end{env}

Al verificarse entonces las hipótesis de (\textbf{II}, 6.6.11), existe un
$Y$-preesquema reducido $X_2$, y un $Y$-morfismo $h:X_2\to X$ \emph{finito} y
\emph{birracional} tal que la restricción de $h$ a $h^{-1}(X\setmin Z)$ es un
\emph{isomorfismo} sobre $X\setmin Z$ y $h^*(\mathcal L)$ es amplio. Reemplazando
$X'$ por $X_2$, se ve que se está reducido a demostrar la proposición
suponiendo además que \emph{$g$ posee las propiedades que se acaban de
enumerar para $h$}. Designaremos también por $Z$ un subpreesquema de
$X$ cuyo espacio subyacente es $Z$, que es propio sobre $Y$
(\textbf{II}, 5.4.10).

\begin{env}[2.6.2.3]
\label{III.2.6.2.3-fr}
Sean ahora $X_1$ un \emph{subpreesquema cerrado} de $X$,
$j:X_1\to X$ la inclusión canónica,
$X'_1=g^{-1}(X_1)=X'\times_X X_1$ su imagen inversa,
$j':X'_1\to X'$ la inclusión canónica, de modo que se tiene el diagrama
conmutativo (\hyperref[III.2.6.2.2-fr]{2.6.2.2}); pongamos
$\mathcal L_1=j^*(\mathcal L)$,
$\mathcal L'_1=j'^*(\mathcal L')=g_1^*(\mathcal L_1)$, de modo que
$\mathcal L'_1$ es amplio para $f\circ g\circ j'$ (\textbf{II}, 5.1.12). Si se
pone $Z_1=j^{-1}(Z)$, el subpreesquema cerrado $Z_1$ de $X_1$ es propio
sobre $Y$ (\textbf{II}, 5.4.2, (ii)). En otros términos, \emph{las hipótesis de}
(\hyperref[III.2.6.2-fr]{2.6.2}) \emph{se verifican para} $X_1$,
$\mathcal L_1$, $g_1$ y $Z_1$.
\end{env}

\oldpage[III]{113}
Esto nos permitirá demostrar
(\hyperref[III.2.6.2-fr]{2.6.2}) por \emph{inducción noetheriana}
(\textbf{0}\textsubscript{I}, 2.2.2) en el caso en que la restricción de $g$ a
$g^{-1}(X\setmin Z)$ es un \emph{isomorfismo} sobre $X\setmin Z$ (lo que es suficiente para
nuestro propósito, como se ha visto en
(\hyperref[III.2.6.2.2-fr]{2.6.2.2})): bastará establecer que si, para
todo subpreesquema cerrado $X_1$ de $X$ cuyo espacio subyacente es
$\neq X$, la conclusión de (\hyperref[III.2.6.2-fr]{2.6.2}) es verdadera
para el haz $\mathcal L_1$, entonces también es verdadera para el
haz $\mathcal L$.

\begin{env}[2.6.2.4]
\label{III.2.6.2.4-fr}
Sean entonces $\mathcal A=\mathcal O_X$, $\mathcal B=g_*(\mathcal O_X)$,
de modo que $\mathcal B$ es una sub-$\mathcal A$-Álgebra de
$\mathcal R(X)$, que es un $\mathcal A$-Módulo coherente; además, la
restricción $\mathcal B|(X\setmin Z)$ es igual a $\mathcal A|(X\setmin Z)$. Sea
$\mathcal K$ el \emph{conductor} de $\mathcal B$ sobre $\mathcal A$, es decir,
el mayor sub-$\mathcal A$-Módulo cuasicoherente de $\mathcal A$ tal
que $\mathcal B.\mathcal K\subset\mathcal A$ (o también el \emph{anulador} del
$\mathcal A$-Módulo $\mathcal B/\mathcal A$
(\textbf{0}\textsubscript{I}, 5.3.7)), lo que implica
$\mathcal B.\mathcal K=\mathcal K$. Está claro que
$\mathcal K_x=\mathcal A_x$ en todos los puntos que admiten una vecindad
$W_x$ tal que $g$ es un isomorfismo de $g^{-1}(W_x)$ sobre $W_x$, y en
particular en todos los puntos de $X\setmin Z$ y en una vecindad de todo
punto genérico de una componente irreducible de $X$. Consideremos entonces
el subpreesquema cerrado $Z_1=\operatorname{Spec}(\mathcal A/\mathcal K)$
de $X$ definido por $\mathcal K$; es también \emph{propio} sobre $Y$ pues el
subespacio $Z_1$ es cerrado en $Z$ (\textbf{II}, 5.4.10). Además, la definición
de $\mathcal K$ muestra que
$\mathcal B|(X\setmin Z_1)=\mathcal A|(X\setmin Z_1)$; se ve pues que siempre es posible
reducirse al caso en que
$Z=\operatorname{Spec}(\mathcal A/\mathcal K)$, y como se ha visto que
$X\setmin Z_1$ es un abierto no vacío de $X$, siempre se puede suponer que
el espacio $Z$ es \emph{distinto} de $X$.
\end{env}

\begin{env}[2.6.2.5]
\label{III.2.6.2.5-fr}
Consideremos $X'$ como igual a $\operatorname{Spec}(\mathcal B)$ (puesto que
$g$ es afín) y sea $\mathcal K'=\widetilde{\mathcal K}$, Ideal
coherente de $\mathcal O_{X'}$ tal que $g_*(\mathcal K')=\mathcal K$
(\textbf{II}, 1.4.1); el subpreesquema cerrado
$Z'=g^{-1}(Z)=Z\times_X X'$ de $X'$ está definido por $\mathcal K'$ y es igual
a $\operatorname{Spec}(\mathcal B/\mathcal K)$ (\textbf{II}, 1.4.10); como
$h:Z'\to Z$ es un morfismo finito (\textbf{II}, 6.1.5, (iii)), $Z'$ es \emph{propio}
sobre $Y$ (\textbf{II}, 6.1.11 y 5.4.2, (ii)).

Dicho esto, debemos probar que para todo $x\in X$ y toda vecindad
abierta $U$ de $x$, existe una sección $s$ de un
$\mathcal L^{\otimes n}$ $(n>0)$ sobre $X$ tal que
$x\in X_s\subset U$ (\textbf{II}, 4.5.2); distinguiremos dos casos:

$1^\circ$ Se tiene $x\in X\setmin Z$; se puede evidentemente suponer entonces que se tiene
también $U\subset X\setmin Z$, luego el abierto $U'=g^{-1}(U)$ no interseca
$Z'$. Como $\mathcal L'$ es amplio por hipótesis, existe un $n>0$ y
una sección $s'$ de $\mathcal L'^{\otimes n}$ sobre $X'$ tal
que $x'=g^{-1}(x)\in X'_{s'}\subset g^{-1}(U)$ (\textbf{II}, 4.5.2). Además, se
puede suponer que $\mathcal K'\otimes\mathcal L'^{\otimes n}$ esté
generado por sus secciones sobre $X'$ (\textbf{II}, 4.5.5), luego, como
$\mathcal K'_{x'}=\mathcal O_{x'}$, hay una de estas secciones $s''$
tal que $s''(x')\neq0$; multiplicándola por $s'$ (lo que equivale a
reemplazar $n$ por $2n$), se ve que se puede suponer también que
$x'\in X'_{s''}\subset g^{-1}(U)$. Dicho esto, resulta de
(\textbf{0}\textsubscript{I}, 5.4.10) que se tiene un isomorfismo canónico
\[
\Gamma(X,\mathcal K\otimes\mathcal L^{\otimes n})
\xrightarrow{\sim}
\Gamma(X',\mathcal K'\otimes\mathcal L'^{\otimes n}).
\]

La sección $s$ de $\mathcal K\otimes\mathcal L^{\otimes n}$ que
corresponde a $s''$ por este isomorfismo tiene evidentemente las propiedades
requeridas.

$2^\circ$ Se tiene $x\in Z$. Sea $\mathcal J$ el Ideal coherente de
$\mathcal O_X$ que define el subpreesquema cerrado reducido de $X$ cuyo
espacio subyacente es $X\setmin U$, y consideremos \emph{en} $\mathcal B$ los Ideales
\oldpage[III]{114}
coherentes $\mathcal J\mathcal B$ y
$\mathcal J\mathcal K=\mathcal J(\mathcal K\mathcal B)
=\mathcal K(\mathcal J\mathcal B)$, de modo que se tiene el diagrama
de inclusiones
\[
\xymatrix@R=1.5em{
\mathcal J\mathcal B\ar[r] & \mathcal B\\
\mathcal J\ar[r]\ar[u] &
\mathcal A\ar[u]\\
\mathcal J\mathcal K\mathcal B=\mathcal J\mathcal K
\ar[r]\ar[u] &
\mathcal K\ar[u]
}
\tag{2.6.2.6}\label{III.2.6.2.6-fr}
\]

Sea $\mathcal J'$ el Ideal coherente $(\mathcal J\mathcal B)^\sim$ de
$\mathcal O_{X'}$ de modo que $\mathcal J\mathcal B=g_*(\mathcal J')$,
$\mathcal J'\mathcal K'=(\mathcal J\mathcal K\mathcal B)^\sim$, y por
tanto
$\mathcal J'/\mathcal J'\mathcal K'
=(\mathcal J\mathcal B/\mathcal J\mathcal K\mathcal B)^\sim$
(\textbf{II}, 1.4.4). Como $\mathcal J|V=\mathcal J\mathcal K|V$ para todo
abierto $V$ que no interseca $Z$, se ve que el soporte de
$\mathcal J'/\mathcal J'\mathcal K'$ está \emph{contenido en} $Z'$. Como $Z'$ es
propio sobre $Y$, se puede aplicar (\hyperref[III.2.2.4-fr]{2.2.4}) y se
ve que para $n$ suficientemente grande, la aplicación canónica
\[
\Gamma(X',\mathcal J'\otimes\mathcal L'^{\otimes n})
\longrightarrow
\Gamma(X',(\mathcal J'/\mathcal J'\mathcal K')
\otimes\mathcal L'^{\otimes n})
\]
es \emph{sobreyectiva}.

Pero en virtud de (\textbf{0}\textsubscript{I}, 5.4.10), se
concluye que la aplicación canónica
\[
\Gamma(X,\mathcal J\mathcal B\otimes\mathcal L^{\otimes n})
\longrightarrow
\Gamma(X,(\mathcal J\mathcal B/\mathcal J\mathcal K\mathcal B)
\otimes\mathcal L^{\otimes n})
\]
es \emph{sobreyectiva}.

Dicho esto, sean $i:Z\to X$ la inclusión canónica,
$i':Z'\to X'$ la inclusión canónica, de modo que se tiene el diagrama
conmutativo
\[
\xymatrix{
X'\ar[d]_g & Z'\ar[l]_{i'}\ar[d]^{h}\\
X & Z\ar[l]^{i}
}
\]

Sean $\mathcal M=i^*(\mathcal L)$, $\mathcal M'=i'^*(\mathcal L')$;
como $\mathcal L'$ es amplio, $\mathcal M'$ es amplio (\textbf{II}, 5.1.12), y
por otra parte $\mathcal M'=h^*(\mathcal M)$; se concluye pues de
la hipótesis de inducción noetheriana (puesto que $Z\neq X$) que
$\mathcal M$ es \emph{amplio}. Por tanto
$i^*(\mathcal J/\mathcal J\mathcal K)\otimes\mathcal M^{\otimes n}$ está
generado por sus secciones sobre $Z$ para $n$ suficientemente grande
(\textbf{II}, 4.5.5). Como
$\mathcal J/\mathcal J\mathcal K=i_*(i^*(\mathcal J/\mathcal J\mathcal K))$,
se deduce también de (\textbf{0}\textsubscript{I}, 5.4.10) que existe una sección
$s$ de $(\mathcal J/\mathcal J\mathcal K)\otimes
\mathcal L^{\otimes n}$ sobre $X$ (para un $n$ suficientemente grande) tal
que $s(x)\neq0$, puesto que se tiene $\mathcal J_x=\mathcal A_x$ por
definición de $\mathcal J$ y $\mathcal K_x\neq\mathcal A_x$ por
hipótesis. El diagrama (\hyperref[III.2.6.2.6-fr]{2.6.2.6}) muestra que
$s$ es también una sección de
$(\mathcal J\mathcal B/\mathcal J\mathcal K\mathcal B)
\otimes\mathcal L^{\otimes n}$ sobre $X$, luego $s$ es la \emph{imagen
canónica de una sección $t$ de}
$(\mathcal J\mathcal B)\otimes\mathcal L^{\otimes n}$ sobre $X$.
Pero por definición, la imagen canónica $s$ de $t$ mód.
$(\mathcal J\mathcal K\mathcal B)\otimes\mathcal L^{\otimes n}$ está
en
$(\mathcal J\otimes\mathcal L^{\otimes n})/
(\mathcal J\mathcal K\mathcal B\otimes\mathcal L^{\otimes n})$
luego, en virtud de (\hyperref[III.2.6.2.6-fr]{2.6.2.6}), esto implica que
$t$ es una \emph{sección de} $\mathcal J\otimes\mathcal L^{\otimes n}$
sobre $X$, y \emph{a fortiori} una sección de
$\mathcal L^{\otimes n}$. Se ha visto antes que $t(x)\neq0$, luego
$x\in X_t$, y por definición de $\mathcal J$, $t(y)=0$ en $X\setmin U$ que
es el soporte de $\mathcal O_X/\mathcal J$; luego $X_t\subset U$, lo que
termina la demostración.
\end{env}
\end{proof}

\oldpage[III]{115}

\begin{remark}[2.6.3]
\label{III.2.6.3}
\label{III.2.6.3-fr}
Cuando $X$ es \emph{propio} sobre $Y$, la Proposición~\sref{III.2.6.2} puede demostrarse
más sencillamente razonando como en el teorema de Chevalley
\sref[II]{II.6.7.1}, con ayuda de la Proposición~\sref{III.2.6.1} y
del lema \sref[II]{II.6.7.1.1}.
\end{remark}

% END tree/repo/editions/III1_DIPLOMATIC/ega3/ega3-2.tex


% BEGIN tree/repo/editions/III1_DIPLOMATIC/ega3/ega3-3.tex
\section{Teorema de finitud para los morfismos propios}
\label{section:III.3}
\label{section:III.3-fr}

\subsection{El lema de d\'evissage}
\label{subsection:III.3.1}
\label{subsection:III.3.1-fr}

\begin{definition}[3.1.1]
\label{III.3.1.1-fr}
Sea $\boldsymbol{K}$ una categoría abeliana. Se dice que un subconjunto $\boldsymbol{K}'$ del
conjunto de objetos de $\boldsymbol{K}$ es exacto si $0\in\boldsymbol{K}'$ y si, para toda
sucesión exacta
$0\to A'\to A\to A''\to0$
en $\boldsymbol{K}$ tal que dos de los objetos $A$, $A'$, $A''$ están en $\boldsymbol{K}'$, entonces el
tercero está también en $\boldsymbol{K}'$.
\end{definition}

\begin{theorem}[3.1.2]
\label{III.3.1.2-fr}
Sea $X$ un preesquema noetheriano; se designa por $\boldsymbol{K}$ la categoría abeliana
de los $\mathcal O_X$-Módulos coherentes. Sean $\boldsymbol{K}'$ un subconjunto exacto de
$\boldsymbol{K}$, $X'$ una parte cerrada del espacio subyacente a $X$. Se supone que
para toda parte cerrada irreducible $Y$ de $X'$, de punto genérico $y$,
existe un $\mathcal O_X$-Módulo $\mathcal G\in\boldsymbol{K}'$ tal que $\mathcal G_y$ es
un $\boldsymbol{k}(y)$-espacio vectorial de dimensión $1$. Entonces todo $\mathcal O_X$-Módulo
coherente de soporte contenido en $X'$ pertenece a $\boldsymbol{K}'$ (y en particular,
si $X'=X$, se tiene $\boldsymbol{K}'=\boldsymbol{K}$).
\end{theorem}

\begin{proof}
Consideremos la propiedad siguiente $\boldsymbol{P}(Y)$ de una parte cerrada $Y$ de $X'$:
todo $\mathcal O_X$-Módulo coherente de soporte contenido en $Y$ pertenece a
$\boldsymbol{K}'$. En virtud del principio de inducción noetheriana (\textbf{0}\textsubscript{I},
2.2.2), se ve que se está reducido a demostrar que \emph{si $Y$ es una parte
cerrada de $X'$ tal que la propiedad $\boldsymbol{P}(Y')$ es verdadera para toda parte
cerrada $Y'$ de $Y$, distinta de $Y$, entonces $\boldsymbol{P}(Y)$ es verdadera}.

Sea pues $\mathcal F\in\boldsymbol{K}$ de soporte contenido en $Y$ y probemos que
$\mathcal F\in\boldsymbol{K}'$. Designemos también por $Y$ el subpreesquema cerrado reducido de
$X$ cuyo espacio subyacente es $Y$ por espacio subyacente (\textbf{I}, 5.2.1); está definido por un Ideal coherente $\mathcal J$ de $\mathcal O_X$. Se sabe (\textbf{I}, 9.3.4) que existe un
entero $n>0$ tal que $\mathcal J^n\mathcal F=0$; para $1\leq k\leq n$ se tiene
pues una sucesión exacta
\[
0\longrightarrow
\mathcal J^{k-1}\mathcal F/\mathcal J^k\mathcal F
\longrightarrow \mathcal F/\mathcal J^k\mathcal F
\longrightarrow \mathcal F/\mathcal J^{k-1}\mathcal F
\longrightarrow0
\]
de $\mathcal O_X$-Módulos coherentes (\textbf{0}\textsubscript{I}, 5.3.6 y 5.3.3);
como $\boldsymbol{K}'$ es exacto, se ve, por inducción sobre $k$, que basta mostrar
que cada uno de los
$\mathcal F_k=\mathcal J^{k-1}\mathcal F/\mathcal J^k\mathcal F$ está en
$\boldsymbol{K}'$. Se está pues reducido a probar que $\mathcal F\in\boldsymbol{K}'$ bajo la hipótesis
suplementaria de que $\mathcal J\mathcal F=0$; equivale a decir que
$\mathcal F=j_*(j^*(\mathcal F))$, donde $j$ es la inclusión canónica
$Y\to X$. Distingamos ahora dos casos:
\begin{enumerate}
\item[\textit{a)}] $Y$ es \emph{reducible}. Sean $Y=Y'\cup Y''$, siendo $Y'$ y $Y''$
partes cerradas de $Y$, distintas de $Y$; sean también $Y'$,
$Y''$ los subpreesquemas cerrados reducidos de $X$ cuyos espacios
subyacentes son $Y'$, $Y''$ respectivamente, que están definidos respectivamente por
los Ideales coherentes $\mathcal J'$, $\mathcal J''$ de $\mathcal O_X$. Pongamos
$\mathcal F'=\mathcal F\otimes_{\mathcal O_X}
(\mathcal O_X/\mathcal J')$,
$\mathcal F''=\mathcal F\otimes_{\mathcal O_X}
(\mathcal O_X/\mathcal J'')$.
Los homomorfismos canónicos $\mathcal F\to\mathcal F'$,
$\mathcal F\to\mathcal F''$ definen pues un homomorfismo
$u:\mathcal F\to\mathcal F'\oplus\mathcal F''$. Mostremos que para todo
$z\notin Y'\cap Y''$ el homomorfismo
$u_z:\mathcal F_z\to\mathcal F'_z\oplus\mathcal F''_z$ es \emph{biyectivo}.
En efecto, se tiene $\mathcal J'\cap\mathcal J''=\mathcal J$, pues la cuestión es
local y
\oldpage[III]{116}
la igualdad anterior resulta de (\textbf{I}, 5.2.1 y 1.1.5); si $z\notin Y''$, se tiene
pues $\mathcal J'_z=\mathcal J_z$, de donde $\mathcal F'_z=\mathcal F_z$ y
$\mathcal F''_z=0$, lo que establece nuestra afirmación en este caso; se razona
del mismo modo si $z\notin Y'$. Por tanto, el núcleo y el conúcleo de $u$, que están
en $\boldsymbol{K}$ (\textbf{0}\textsubscript{I}, 5.3.4), tienen sus soportes en $Y'\cap Y''$,
y están pues en $\boldsymbol{K}'$ por hipótesis; por la misma razón, $\mathcal F'$ y
$\mathcal F''$ están en $\boldsymbol{K}'$, luego también $\mathcal F'\oplus\mathcal F''$,
puesto que $\boldsymbol{K}'$ es exacto. La conclusión resulta entonces de considerar las
dos sucesiones exactas
\[
0\longrightarrow\operatorname{Im}u\longrightarrow
\mathcal F'\oplus\mathcal F''\longrightarrow
\operatorname{Coker}u\longrightarrow0
\]
\[
0\longrightarrow\operatorname{Ker}u\longrightarrow\mathcal F
\longrightarrow\operatorname{Im}u\longrightarrow0
\]
y de la hipótesis de que $\boldsymbol{K}'$ es exacto.

\item[\textit{b)}] $Y$ es irreducible, y por tanto el subpreesquema $Y$ de $X$ es
\emph{íntegro}. Si $y$ es su punto genérico, se tiene
$(\mathcal O_Y)_y=\boldsymbol{k}(y)$, y como $j^*(\mathcal F)$ es un
$\mathcal O_Y$-Módulo coherente, $\mathcal F_y=(j^*(\mathcal F))_y$ es un
$\boldsymbol{k}(y)$-espacio vectorial de dimensión finita $m$. Por hipótesis, hay un
$\mathcal O_X$-Módulo coherente $\mathcal G$ (necesariamente de soporte $Y$)
tal que $\mathcal G_y$ es un $\boldsymbol{k}(y)$-espacio vectorial de dimensión $1$. Por
tanto, hay un $\boldsymbol{k}(y)$-isomorfismo
$(\mathcal G_y)^m\xrightarrow{\sim}\mathcal F_y$, que es también un
$\mathcal O_Y$-isomorfismo, y como $\mathcal G^m$ y $\mathcal F$ son
coherentes, existe una vecindad abierta $W$ de $y$ en $X$ y un
isomorfismo $\mathcal G^m|W\xrightarrow{\sim}\mathcal F|W$
(\textbf{0}\textsubscript{I}, 5.2.7). Sea $\mathcal H$ el grafo de este isomorfismo,
que es un sub-$(\mathcal O_X|W)$-Módulo coherente de
$(\mathcal G^m\oplus\mathcal F)|W$, canónicamente isomorfo a
$\mathcal G^m|W$ y a $\mathcal F|W$; existe pues un sub-$\mathcal O_X$-%
Módulo coherente $\mathcal H_0$ de $\mathcal G^m\oplus\mathcal F$, que induce
$\mathcal H$ sobre $W$ y $0$ sobre $X\setmin Y$ puesto que $\mathcal G^m$ y
$\mathcal F$ tienen por soporte $Y$ (\textbf{I}, 9.4.7). Las restricciones
$v:\mathcal H_0\to\mathcal G^m$ y $w:\mathcal H_0\to\mathcal F$ de las
proyecciones canónicas de $\mathcal G^m\oplus\mathcal F$ son entonces
homomorfismos de $\mathcal O_X$-Módulos coherentes, que en $W$ y en
$X\setmin Y$ se reducen a isomorfismos; en otros términos, los núcleos y
conúcleos de $v$ y $w$ tienen su soporte en el conjunto cerrado
$Y\setmin(Y\cap W)$, distinto de $Y$. Pertenecen pues a $\boldsymbol{K}'$; por otra
parte, se tiene $\mathcal G^m\in\boldsymbol{K}'$ puesto que $\mathcal G\in\boldsymbol{K}'$ y $\boldsymbol{K}'$ es
exacto. Se concluye sucesivamente, por la exactitud de $\boldsymbol{K}'$, que
$\mathcal H_0\in\boldsymbol{K}'$ y después que $\mathcal F\in\boldsymbol{K}'$. C.Q.D.
\end{enumerate}
\end{proof}

\begin{corollary}[3.1.3]
\label{III.3.1.3-fr}
Supongamos que el subconjunto exacto $\boldsymbol{K}'$ de $\boldsymbol{K}$ tiene además la propiedad
de que todo factor directo coherente de un $\mathcal O_X$-Módulo coherente
$\mathcal M\in\boldsymbol{K}'$ pertenece también a $\boldsymbol{K}'$. En estas condiciones, la
conclusión de (\hyperref[III.3.1.2-fr]{3.1.2}) sigue siendo válida cuando la
condición «$\mathcal G_y$ es un espacio vectorial sobre $\boldsymbol{k}(y)$ de dimensión~$1$» se sustituye por $\mathcal G_y\neq0$ (lo que equivale a $\operatorname{Supp}(\mathcal G)=Y$).
\end{corollary}

\begin{proof}
El razonamiento del Teorema~\sref{III.3.1.2} solo debe modificarse en el caso~\textit{(b)}; ahora $\mathcal G_y$ es un espacio vectorial sobre $\boldsymbol{k}(y)$ de dimensión $q>0$ y, por tanto, se tiene un $\mathcal O_Y$-isomorfismo $(\mathcal G_y)^m\xrightarrow{\sim}(\mathcal F_y)^q$; el final
del razonamiento de (\hyperref[III.3.1.2-fr]{3.1.2}) prueba entonces que
$\mathcal F^q\in\boldsymbol{K}'$, y la hipótesis adicional sobre $\boldsymbol{K}'$ implica que
$\mathcal F\in\boldsymbol{K}'$.
\end{proof}

\subsection{El teorema de finitud: caso de los esquemas usuales}
\label{subsection:III.3.2-fr}

\begin{theorem}[3.2.1]
\label{III.3.2.1-fr}
Sean $Y$ un preesquema localmente noetheriano, $f:X\to Y$ un morfismo
propio. Para todo $\mathcal O_X$-Módulo coherente $\mathcal F$, los
$\mathcal O_Y$-Módulos $R^qf_*(\mathcal F)$ son coherentes para $q\geq0$.
\end{theorem}

\begin{proof}
Al ser la cuestión local sobre $Y$, se puede suponer $Y$ noetheriano, luego $X$
noetheriano (\textbf{I}, 6.3.7). Los $\mathcal O_X$-Módulos coherentes $\mathcal F$ para
los cuales la conclusión del teorema (\hyperref[III.3.2.1-fr]{3.2.1}) es verdadera
forman un subconjunto \emph{exacto} $\boldsymbol{K}'$ de la categoría $\boldsymbol{K}$ de los
$\mathcal O_X$-Módulos coherentes.
\oldpage[III]{117}
En efecto, sea
$0\to\mathcal F'\to\mathcal F\to\mathcal F''\to0$ una sucesión exacta de
$\mathcal O_X$-Módulos coherentes; supongamos por ejemplo que $\mathcal F'$
y $\mathcal F''$ pertenecen a $\boldsymbol{K}'$; se tiene la sucesión exacta de cohomología
\[
R^{q-1}f_*(\mathcal F'')\xrightarrow{\partial}R^qf_*(\mathcal F')
\longrightarrow R^qf_*(\mathcal F)\longrightarrow R^qf_*(\mathcal F'')
\xrightarrow{\partial}R^{q+1}f_*(\mathcal F')
\]
en la que, por hipótesis, los cuatro términos extremos son coherentes;
lo mismo ocurre pues con el término central $R^qf_*(\mathcal F)$ por
(\textbf{0}\textsubscript{I}, 5.3.4 y 5.3.3). Se muestra de la misma manera que
cuando $\mathcal F$ y $\mathcal F'$ (resp. $\mathcal F$ y
$\mathcal F''$) están en $\boldsymbol{K}'$, lo mismo ocurre con $\mathcal F''$ (resp.
$\mathcal F'$). Además, todo \emph{factor directo} coherente $\mathcal F'$
de un $\mathcal O_X$-Módulo $\mathcal F\in\boldsymbol{K}'$ pertenece también a $\boldsymbol{K}'$: en
efecto, $R^qf_*(\mathcal F')$ es entonces factor directo de
$R^qf_*(\mathcal F)$ (G, II, 4.4.4), luego es de tipo finito, y como es
cuasicoherente (\hyperref[III.1.4.10-fr]{1.4.10}), es coherente, al ser $Y$
noetheriano. En virtud de (\hyperref[III.3.1.3-fr]{3.1.3}), se está
reducido a probar que cuando $X$ es \emph{irreducible} de punto genérico
$x$, existe \emph{un} $\mathcal O_X$-Módulo coherente $\mathcal F$ que pertenece a
$\boldsymbol{K}'$, tal que $\mathcal F_x\neq0$: en efecto, si se establece este punto, se
podrá aplicar a todo subpreesquema cerrado irreducible $Y$ de $X$, pues
si $j:Y\to X$ es la inclusión canónica, $f\circ j$ es propio (\textbf{II}, 5.4.2),
y si $\mathcal G$ es un $\mathcal O_Y$-Módulo coherente de soporte $Y$,
$j_*(\mathcal G)$ es un $\mathcal O_X$-Módulo coherente tal que
\[
R^q(f\circ j)_*(\mathcal G)=R^qf_*(j_*(\mathcal G))
\]
(G, II, 4.9.1), luego se está efectivamente en las condiciones de aplicación de
(\hyperref[III.3.1.3-fr]{3.1.3}).

Ahora bien, en virtud del lema de Chow (\textbf{II}, 5.6.2), existe un preesquema
irreducible $X'$ y un morfismo \emph{proyectivo} y sobreyectivo
$g:X'\to X$ tal que $f\circ g:X'\to Y$ es \emph{proyectivo}. Existe un
$\mathcal O_{X'}$-Módulo $\mathcal L$ amplio para $g$ (\textbf{II}, 5.3.1);
apliquemos el teorema fundamental de los morfismos proyectivos
(\hyperref[III.2.2.1-fr]{2.2.1}) a $g:X'\to X$ y a $\mathcal L$:
existe pues un entero $n$ tal que
$\mathcal F=g_*(\mathcal O_{X'}(n))$ es un $\mathcal O_X$-Módulo coherente
y $R^qg_*(\mathcal O_{X'}(n))=0$ para todo $q>0$; además, como
\[
g^*(g_*(\mathcal O_{X'}(n)))\longrightarrow\mathcal O_{X'}(n)
\]
es sobreyectivo para $n$ suficientemente grande (\hyperref[III.2.2.1-fr]{2.2.1}), se
ve que se puede suponer que, en el punto genérico $x$ de $X$, se tiene
$\mathcal F_x\neq0$ (\textbf{II}, 3.4.7). Por otra parte, como $f\circ g$ es
proyectivo y $Y$ noetheriano, los
$R^p(f\circ g)_*(\mathcal O_{X'}(n))$ son \emph{coherentes}
(\hyperref[III.2.2.1-fr]{2.2.1}). Dicho esto,
$R^*(f\circ g)_*(\mathcal O_{X'}(n))$ es el resultado de una sucesión
espectral de Leray, cuyo término $E_2$ viene dado por
\[
E_2^{pq}=R^pf_*(R^qg_*(\mathcal O_{X'}(n))~;
\]
lo anterior muestra que esta sucesión espectral es degenerada, y se sabe
entonces (\textbf{0}, 11.1.6) que $E_2^{p0}=R^pf_*(\mathcal F)$ es isomorfo a
$R^p(f\circ g)_*(\mathcal O_{X'}(n)))$, lo que termina la demostración.
\end{proof}

\begin{corollary}[3.2.2]
\label{III.3.2.2}
\label{III.3.2.2-fr}
Sea $Y$ un preesquema localmente noetheriano.
Para todo morfismo propio $f:X\to Y$, la imagen directa por $f$ de cualquier $\mathcal O_X$-módulo coherente es un $\mathcal O_Y$-módulo coherente.
\end{corollary}

\begin{corollary}[3.2.3]
\label{III.3.2.3}
\label{III.3.2.3-fr}
Sea $A$ un anillo noetheriano y sea $X$ un esquema propio sobre $A$; para todo $\mathcal O_X$-módulo coherente $\mathcal F$, los $H^p(X,\mathcal F)$ son $A$-módulos de tipo finito, y existe un entero $r>0$ tal que, para todo $\mathcal O_X$-módulo coherente $\mathcal F$ y todo $p>r$, $H^p(X,\mathcal F)=0$.
\end{corollary}

\begin{proof}
La segunda afirmación ya se demostró en~\sref{III.1.4.12}; la primera se deduce del teorema de finitud~\sref{III.3.2.1}, teniendo en cuenta el Corolario~\sref{III.1.4.11}.
\end{proof}

En particular, si $X$ es un \emph{esquema algebraico propio} sobre un cuerpo $k$, entonces, para todo $\mathcal O_X$-módulo coherente $\mathcal F$, los $H^p(X,\mathcal F)$ son espacios vectoriales sobre $k$ de \emph{dimensión finita}.

\begin{corollary}[3.2.4]
\label{III.3.2.4}
\label{III.3.2.4-fr}
Sea $Y$ un preesquema localmente noetheriano y sea $f:X\to Y$ un morfismo de tipo finito.
Para todo $\mathcal O_X$-módulo coherente $\mathcal F$ cuyo soporte sea propio sobre $Y$~\sref[II]{II.5.4.10}, los $\mathcal O_Y$-módulos $R^qf_*(\mathcal F)$ son coherentes.
\end{corollary}

\oldpage[III]{118}
\begin{proof}
Como la cuestión es local sobre $Y$, se puede suponer que $Y$ es noetheriano, y lo mismo ocurre entonces con $X$~\sref[I]{I.6.3.7}.
Por hipótesis, todo subpreesquema cerrado $Z$ de $X$ cuyo espacio subyacente sea $\operatorname{Supp}(\mathcal F)$ es propio sobre $Y$; en otras palabras, si $j:Y\to X$ es la inmersión canónica, entonces $f\circ j:Z\to Y$ es propio.
Se puede suponer que $Z$ es tal que $\mathcal F=j_*(\mathcal G)$, donde $\mathcal G=j^*(\mathcal F)$ es un $\mathcal O_Z$-módulo coherente~\sref[I]{I.9.3.5}; como $R^qf_*(\mathcal F)=R^q(f\circ j)_*(\mathcal G)$ por el Corolario~\sref{III.1.3.4}, la conclusión se deduce inmediatamente del Teorema~\sref{III.3.2.1}.
\end{proof}

\subsection{Generalización del teorema de finitud (esquemas ordinarios)}
\label{subsection:III.3.3}
\label{subsection:III.3.3-fr}

\begin{proposition}[3.3.1]
\label{III.3.3.1}
\label{III.3.3.1-fr}
Sea $Y$ un preesquema noetheriano, sea $\mathcal S$ una $\mathcal O_Y$-álgebra cuasicoherente de tipo finito, graduada en grados positivos, sea $Y'=\operatorname{Proj}(\mathcal S)$, y sea $g:Y'\to Y$ el morfismo estructural.
Sea $f:X\to Y$ un morfismo propio, sea $\mathcal S'=f^*(\mathcal S)$, y sea $\mathcal M=\bigoplus_{k\in\mathbb Z}\mathcal M_k$ un $\mathcal S'$-módulo graduado cuasicoherente de tipo finito.
Entonces los \[
R^pf_*(\mathcal M)=\bigoplus_{k\in\mathbb Z}R^pf_*(\mathcal M_k)
\] son $\mathcal S$-módulos graduados de tipo finito para todo $p$.
Supongamos además que $\mathcal S$ está generada por $\mathcal S_1$; entonces, para todo $p\in\mathbb Z$, existe un entero $k_p$ tal que, para todo $k\geq k_p$ y todo $r\geq0$, se tiene
\[
R^pf_*(\mathcal M_{k+r})=\mathcal S_rR^pf_*(\mathcal M_k).
\tag{3.3.1.1}\label{III.3.3.1.1-fr}
\] 
\end{proposition}

\begin{proof}
La primera afirmación es idéntica al enunciado del Teorema~\sref{III.2.4.1}[i], en el que simplemente se sustituye «morfismo proyectivo» por «morfismo propio».
En la demostración del Teorema~\sref{III.2.4.1}[i], la hipótesis sobre $f$ solo se utilizó para demostrar (con la notación de aquella demostración) que $R^pf'_*(\widetilde{\mathcal M})$ es un $\mathcal O_{Y'}$-módulo coherente.
Bajo las hipótesis de la Proposición~\sref{III.3.3.1}, $f'$ es propio~\sref[II]{II.5.4.2}[iii], de modo que se puede retomar sin modificación la demostración del Teorema~\sref{III.2.4.1}[i], gracias al teorema de finitud~\sref{III.3.2.1}.

Para la segunda afirmación, basta observar que existe un recubrimiento abierto afín finito $(U_i)$ de $Y$ tal que las restricciones a los $U_i$ de los dos miembros de~(\hyperref[III.3.3.1.1]{3.3.1.1}) son iguales para todo $k\geq k_{p,i}$~\sref[II]{II.2.1.6}[ii]; basta tomar como $k_p$ el mayor de los $k_{p,i}$.
\end{proof}

\begin{corollary}[3.3.2]
\label{III.3.3.2}
\label{III.3.3.2-fr}
Sea $A$ un anillo noetheriano, sea $\mathfrak m$ un ideal de $A$, sea $X$ un $A$-esquema propio y sea $\mathcal F$ un $\mathcal O_X$-módulo coherente.
Entonces, para todo $p\geq0$, la suma directa $\bigoplus_{k\geq0}H^p(X,\mathfrak m^k\mathcal F)$ es un módulo de tipo finito sobre el anillo $S=\bigoplus_{k\geq0}\mathfrak m^k$; en particular, existe un entero $k_p\geq0$ tal que, para todo $k\geq k_p$ y todo $r\geq0$, se tiene
\[
H^p(X,\mathfrak m^{k+r}\mathcal F)
=\mathfrak m^rH^p(X,\mathfrak m^k\mathcal F).
\tag{3.3.3.1}\label{III.3.3.2.1-fr}
\]
\end{corollary}

\begin{proof}
Basta aplicar el Corolario~\sref{III.3.3.2} con $Y=\operatorname{Spec}(A)$, $\mathcal S=\widetilde{\mathcal S}$ y $\mathcal M_k=\mathfrak m^k\mathcal F$, teniendo en cuenta el Corolario~\sref{III.1.4.11}.
\end{proof}

Conviene recordar que la estructura de $S$-módulo sobre $\bigoplus_{k\geq0}H^p(X,\mathfrak m^k\mathcal F)$ se obtiene considerando, para todo $a\in\mathfrak m^r$, la aplicación \[
H^p(X,\mathfrak m^k\mathcal F)
\longrightarrow H^p(X,\mathfrak m^{k+r}\mathcal F)
\] que procede del paso a cohomología de la aplicación de multiplicación $\mathfrak m^k\mathcal F\to\mathfrak m^{k+r}\mathcal F$ definida por $a$~\sref{III.2.4.1}.

\subsection{Teorema de finitud: caso de los esquemas formales}
\label{subsection:III.3.4}
\label{subsection:III.3.4-fr}

\oldpage[III]{119}
Los resultados de esta sección (salvo la definición~\sref{III.3.4.1}) no se utilizarán en el resto de este capítulo.

\begin{env}[3.4.1]
\label{III.3.4.1}
\label{III.3.4.1-fr}
Sean $\mathfrak X$ y $\mathfrak S$ dos preesquemas formales localmente noetherianos~\sref[I]{I.10.4.2}, y sea $f:\mathfrak X\to\mathfrak S$ un morfismo de preesquemas formales.
Diremos que $f$ es un morfismo \emph{propio} si
verifica las condiciones siguientes:
\begin{enumerate}
\item[$1^\circ$] \emph{$f$ es un morfismo de tipo finito} (\textbf{I}, 10.13.3).
\item[$2^\circ$] \emph{Si $\mathcal K$ es un haz de ideales de definición de $\mathfrak S$ y se pone $\mathcal J=f^*(\mathcal K)\mathcal O_{\mathfrak X}$, $X_0=(\mathfrak X,\mathcal O_{\mathfrak X}/\mathcal J)$, $S_0=(\mathfrak S,\mathcal O_{\mathfrak S}/\mathcal K)$, entonces el morfismo $f_0:X_0\to Y_0$ inducido por $f$~\sref[I]{I.10.5.6} es propio}.
\end{enumerate}
Es inmediato que esta definición no depende del haz de ideales de definición $\mathcal K$ de $\mathfrak S$ considerado; en efecto, si $\mathcal K'$ es un segundo haz de ideales de definición tal que $\mathcal K'\subset\mathcal K$, y si se pone $\mathcal J'=f^*(\mathcal K')\mathcal O_{\mathfrak X}$, $X'_0=(\mathfrak X,\mathcal O_{\mathfrak X}/\mathcal J')$, $S'_0=(\mathfrak S,\mathcal O_{\mathfrak S}/\mathcal K')$, entonces el morfismo $f'_0:X'_0\to S'_0$ inducido por $f$ es tal que el diagrama
\[
\xymatrix{
X_0\ar[r]^{f_0}\ar[d]_i & S_0\ar[d]^j\\
X'_0\ar[r]_{f'_0} & S'_0
}
\]
es conmutativo, siendo $i$ y $j$ inmersiones sobreyectivas; en virtud de~\sref[II]{II.5.4.5}, es equivalente decir que $f_0$ o $f'_0$ es propio, en virtud de (\textbf{II}, 5.4.5). Se observará que, para todo $n\geq0$, si se pone $X_n=(\mathfrak X,\mathcal O_{\mathfrak X}/\mathcal J^{n+1})$, $S_n=(\mathfrak S,\mathcal O_{\mathfrak S}/\mathcal K^{n+1})$, entonces el morfismo $f_n:X_n\to Y_n$ inducido por $f$~\sref[I]{I.10.5.6} es propio para todo $n$ en cuanto lo es para $n=0$~\sref[II]{II.5.4.6}.

Si $g:Y\to Z$ es un morfismo propio de preesquemas ordinarios localmente noetherianos, $Z'$ es un subconjunto cerrado de $Z$, e $Y'$ es un subconjunto cerrado de $Y$ tal que $g(Y')\subset Z'$, entonces la extensión $\widehat g:Y_{/Y'}\to Z_{/Z'}$ de $g$ a las compleciones~\sref[I]{I.10.9.1} es un morfismo propio de preesquemas formales, como se deduce de la definición y de~\sref[II]{II.5.4.5}.

Sean $\mathfrak X$ y $\mathfrak S$ dos preesquemas formales localmente noetherianos y sea $f:\mathfrak X\to\mathfrak S$ un morfismo \emph{de tipo finito} (\textbf{I}, 10.13.3); siendo las notaciones las mismas que arriba, se dice que una parte $Z$ del espacio subyacente a $\mathfrak X$ es \emph{propio} sobre $\mathfrak S$ (o propio para $f$) si, considerado como subconjunto de $X_0$, $Z$ es \emph{propio sobre $S_0$}~\sref[II]{II.5.4.10}.
Todas las propiedades de los subconjuntos propios de preesquemas ordinarios enunciadas en~\sref[II]{II.5.4.10} siguen siendo válidas para los subconjuntos propios de preesquemas formales, como se deduce inmediatamente de las definiciones.
\end{env}

\begin{theorem}[3.4.2]
\label{III.3.4.2}
\label{III.3.4.2-fr}
Sean $\mathfrak X$ y $\mathfrak Y$ preesquemas formales localmente noetherianos, y sea $f:\mathfrak X\to\mathfrak Y$ un morfismo propio.
Para todo $\mathcal O_{\mathfrak X}$-módulo coherente $\mathcal F$, los $\mathcal O_{\mathfrak Y}$-módulos $R^qf_*(\mathcal F)$ son coherentes para todo $q\geq0$.
\end{theorem}

Sea $\mathcal J$ un haz de ideales de definición de $\mathfrak Y$, sea $\mathcal K=f^*(\mathcal J)\mathcal O_{\mathfrak X}$, y consideremos los $\mathcal O_{\mathfrak X}$-módulos
\[
\mathcal F_k=\mathcal F\otimes_{\mathcal O_{\mathfrak Y}}
(\mathcal O_{\mathfrak Y}/\mathcal J^{k+1})
=\mathcal F/\mathcal K^{k+1}\mathcal F\qquad(k\geq0)
\tag{3.4.2.1}\label{III.3.4.2.1-fr}
\]
que forman evidentemente un \emph{sistema proyectivo} de $\mathcal O_{\mathfrak X}$-módulos topológicos, tal que
\oldpage[III]{120}
\[
\mathcal F=\varprojlim_k\mathcal F_k
\]
por \sref[I]{I.10.11.3}.
Por otra parte, del Teorema~\sref{III.3.4.2} se deduce que cada $R^qf_*(\mathcal F)$, por ser coherente, está dotado naturalmente de una estructura de $\mathcal O_{\mathfrak Y}$-módulo topológico~\sref[I]{I.10.11.6}, y lo mismo ocurre con los $R^qf_*(\mathcal F_k)$.
A los homomorfismos canónicos $\mathcal F\to\mathcal F_k=\mathcal F/\mathcal K^{k+1}\mathcal F$ corresponden canónicamente homomorfismos
\[
R^qf_*(\mathcal F)\longrightarrow R^qf_*(\mathcal F_k)
\]
que son necesariamente continuos para las estructuras de $\mathcal O_{\mathfrak Y}$-módulo topológico anteriores~\sref[I]{I.10.11.6} y forman un sistema proyectivo, dando en el límite un homomorfismo funtorial canónico
\[
R^qf_*(\mathcal F)\longrightarrow\varprojlim_kR^qf_*(\mathcal F_k)
\tag{3.4.2.2}\label{III.3.4.2.2-fr}
\]
que es un homomorfismo continuo de $\mathcal O_{\mathfrak Y}$-módulos topológicos.
Junto con el Teorema~\sref{III.3.4.2}, demostraremos el
\begin{corollary}[3.4.3]
\label{III.3.4.3}
\label{III.3.4.3-fr}
Cada uno de los homomorfismos~(\hyperref[III.3.4.2.2]{3.4.2.2}) es un isomorfismo topológico.
Además, si $\mathfrak Y$ es noetheriano, el sistema proyectivo $(R^qf_*(\mathcal F/\mathcal K^{k+1}\mathcal F))_{k\geq0}$ satisface la condición \emph{(ML)}~\sref[0]{0.13.1.1}.
\end{corollary}
Comenzaremos estableciendo el Teorema~\sref{III.3.4.2} y el Corolario~\sref{III.3.4.3} cuando $Y$ es un esquema formal afín noetheriano~\sref[I]{I.10.4.1}:
\begin{corollary}[3.4.4]
\label{III.3.4.4}
\label{III.3.4.4-fr}
Bajo las hipótesis del Teorema~\sref{III.3.4.2}, supongamos además que $\mathfrak Y=\operatorname{Spf}(A)$, donde $A$ es un anillo ádico noetheriano.
Sea $\mathfrak J$ un ideal de definición de $A$, y pongamos $\mathcal F_k=\mathcal F/\mathfrak J^{k+1}\mathcal F$ para $k\geq0$.
Entonces los $H^n(\mathfrak X,\mathcal F)$ son $A$-módulos de tipo finito; el sistema proyectivo $(H^n(\mathfrak X,\mathcal F_k))_{k\geq0}$ satisface la condición \emph{(ML)} para todo $n$; si se pone
\[
N_{n,k}=\operatorname{Ker}\bigl(H^n(\mathfrak X,\mathcal F)
\longrightarrow H^n(\mathfrak X,\mathcal F_k)\bigr)
\tag{3.4.4.1}\label{III.3.4.4.1-fr}
\]
(igual también a $\operatorname{Im}(H^n(\mathfrak X,\mathfrak J^{k+1}\mathcal F)
\to H^n(\mathfrak X,\mathcal F))$ por la sucesión exacta de cohomología), entonces los $N_{k,n}$ definen sobre $H^n(\mathfrak X,\mathcal F)$ una filtración $\mathfrak J$-buena~\sref[0]{0.13.7.7}; finalmente, el homomorfismo canónico
\[
H^n(\mathfrak X,\mathcal F)
\longrightarrow\varprojlim_kH^n(\mathfrak X,\mathcal F_k)
\tag{3.4.4.2}\label{III.3.4.4.2-fr}
\]
es un isomorfismo topológico para todo $n$ (el miembro izquierdo está dotado de la topología $\mathfrak J$-ádica y los $H^n(\mathfrak X,\mathcal F_k)$ de la topología discreta).
\end{corollary}

Pongamos
\[
S=\operatorname{gr}(A)=\bigoplus_{k\geq0}
\mathfrak J^k/\mathfrak J^{k+1},\qquad
\mathcal M=\operatorname{gr}(\mathcal F)=\bigoplus_{k\geq0}
\mathfrak J^k\mathcal F/\mathfrak J^{k+1}\mathcal F.
\tag{3.4.4.3}\label{III.3.4.4.3-fr}
\]
Se sabe que $\mathfrak J^\Delta$ es un haz de ideales de definición de $\mathfrak Y$~\sref[I]{I.10.3.1}; sean $\mathcal K=f^*(\mathfrak J^\Delta)\mathcal O_{\mathfrak X}$, $X_0=(\mathfrak X,\mathcal O_{\mathfrak X}/\mathcal K)$, $Y_0=(\mathfrak Y,\mathcal O_{\mathfrak Y}/\mathfrak J^\Delta)
=\operatorname{Spec}(A_0)$, con $A_0=A/\mathfrak J$.
Está claro que los $\mathcal M_k=\mathfrak J^k\mathcal F/\mathfrak J^{k+1}\mathcal F$ son $\mathcal O_{X_0}$-módulos coherentes~\sref[I]{I.10.11.3}.
Consideremos, por otra parte, la $\mathcal O_{X_0}$-álgebra graduada cuasicoherente
\[
\mathcal S=\mathcal O_{X_0}\otimes_{A_0}S
=\operatorname{gr}(\mathcal O_{\mathfrak X})
=\bigoplus_{k\geq0}\mathcal K^k/\mathcal K^{k+1}.
\tag{3.4.4.4}\label{III.3.4.4.4-fr}
\]

La hipótesis de que $\mathcal F$ es un $\mathcal O_{\mathfrak X}$-módulo de tipo finito implica primero que $\mathcal M$ es
\oldpage[III]{121}
un $\mathcal S$-módulo graduado \emph{de tipo finito}.
En efecto, la cuestión es local sobre $\mathfrak X$, y se puede pues suponer, para tratarla, que $\mathfrak X=\operatorname{Spf}(B)$, donde $B$ es un anillo ádico noetheriano, y que $\mathcal F=N^\Delta$, donde $N$ es un $B$-módulo de tipo finito~\sref[I]{I.10.10.5}; se tiene además $X_0=\operatorname{Spec}(B_0)$, donde $B_0=B/\mathfrak J B$, y los $\mathcal O_{X_0}$-módulos cuasicoherentes $\mathcal S$ y $\mathcal M$ son respectivamente iguales a $\widetilde{S}'$ y $\widetilde{M}'$, donde
\[
S'=\bigoplus_{k\geq0}
((\mathfrak J^k/\mathfrak J^{k+1})\otimes_{A_0}B_0)\quad\text{y}\quad
M'=\bigoplus_{k\geq0}
((\mathfrak J^k/\mathfrak J^{k+1})\otimes_{A_0}N_0),
\]
con $N_0=N/\mathfrak JN$; se tiene evidentemente $M'=S'\otimes_{B_0}N_0$, y como $N_0$ es un $B_0$-módulo de tipo finito, $M'$ es un $S'$-módulo de tipo finito, lo que prueba nuestra afirmación~\sref[I]{I.1.3.13}.

Como el morfismo $f_0:X_0\to Y_0$ es \emph{propio} por hipótesis, se puede aplicar el Corolario~\sref{III.3.3.2} a $\mathcal S$, $\mathcal M$ y al morfismo $f_0$: teniendo en cuenta el Corolario~\sref{III.1.4.11}, se concluye que, para \emph{todo $n\geq0$}, $\bigoplus_{k\geq0}H^n(X_0,\mathcal M_k)$ es un $S$-módulo graduado \emph{de
tipo finito}. Esto prueba que la condición $(\mathrm F_n)$ de (\textbf{0}, 13.7.7) se
verifica para \emph{todo $n\geq0$}, cuando se considera el sistema proyectivo estricto $(\mathcal F/\mathfrak J^k\mathcal F)_{k\geq0}$ de haces de grupos abelianos sobre $X_0$, cada uno dotado de su estructura natural de «$A$-módulo filtrado». Se puede pues aplicar (\textbf{0}, 13.7.7), que prueba que:
\begin{enumerate}
\item[$1^\circ$] El sistema proyectivo
$(H^n(\mathfrak X,\mathcal F_k))_{k\geq0}$ verifica la condición (ML).
\item[$2^\circ$] Si
$H'^{n}=\varprojlim_kH^n(\mathfrak X,\mathcal F_k)$, entonces $H'^{n}$ es un $A$-módulo de tipo finito.
\item[$3^\circ$] La filtración definida sobre $H'^{n}$ por los núcleos de los homomorfismos canónicos $H'^{n}\to H^n(\mathfrak X,\mathcal F_k)$ es $\mathfrak J$-buena.
\end{enumerate}

Observemos, por otra parte, que si se pone $X_k=(\mathfrak X,\mathcal O_{\mathfrak X}/\mathcal K^{k+1})$, entonces $\mathcal F_k$ es un $\mathcal O_{X_k}$-módulo coherente~\sref[I]{I.10.11.3}; y si $U$ es un abierto afín de $X_0$, entonces $U$ es también un abierto afín de cada $X_k$~\sref[I]{I.5.1.9}, de modo que $H^n(U,\mathcal F_k)=0$ para todo $n>0$ y todo $k$~\sref{III.1.3.1}, y $H^0(U,\mathcal F_k)\to H^0(U,\mathcal F_h)$ es sobreyectivo para $h\leq k$~\sref[I]{I.1.3.9}.
Se está, pues, en las condiciones de~\sref[0]{0.13.3.2}, y la aplicación de~\sref[0]{0.13.3.1} demuestra que $H'^{n}$ se identifica canónicamente con $H^n(\mathfrak X,\varprojlim_k\mathcal F_k)=H^n(\mathfrak X,\mathcal F)$; esto termina la demostración del Corolario~\sref{III.3.4.4}.

\begin{env}[3.4.5]
\label{III.3.4.5}
\label{III.3.4.5-fr}
Volvamos a la demostración de~\sref{III.3.4.2} y~\sref{III.3.4.3}.
Demostremos primero las proposiciones en el caso $\mathfrak Y=\operatorname{Spf}(A)$ considerado en~\sref{III.3.4.4}; para ello, para todo $g\in A$, apliquemos~\sref{III.3.4.4} al esquema formal afín noetheriano inducido sobre el abierto $\mathfrak Y_g=\mathfrak D(g)$ de $\mathfrak Y$, que es igual a $\operatorname{Spf}(A_{\{g\}})$, y al preesquema formal inducido por $\mathfrak X$ sobre $f^{-1}(\mathfrak Y_g)$; observemos que $\mathfrak Y_g$ es también un abierto afín del preesquema $Y_k=(\mathfrak Y,\mathcal O_{\mathfrak Y}/(\mathfrak J^\Delta)^{k+1})$, y como $\mathcal F_k$ es un $\mathcal O_{X_k}$-módulo coherente, se tiene
\[
H^n(f^{-1}(\mathfrak Y_g),\mathcal F_k)
=\Gamma(\mathfrak Y_g,R^nf_*(\mathcal F_k))
\]
para todo $k\geq0$, en virtud del Corolario~\sref{III.1.4.11}.
El homomorfismo canónico
\[
H^n(f^{-1}(\mathfrak Y_g),\mathcal F)
\longrightarrow\varprojlim_k\Gamma(\mathfrak Y_g,R^nf_*(\mathcal F_k))
\]
es un isomorfismo; pero se tiene~\sref[0]{0.3.2.6}
\[
\varprojlim_k\Gamma(\mathfrak Y_g,R^nf_*(\mathcal F_k))
=\Gamma(\mathfrak Y_g,\varprojlim_kR^nf_*(\mathcal F_k))
\]
\oldpage[III]{122}
y como el haz $R^nf_*(\mathcal F)$ es el haz asociado al prehaz $\mathfrak Y_g\leadsto H^n(f^{-1}(\mathfrak Y_g),\mathcal F)$ sobre los $\mathfrak Y_g$~\sref[0]{0.3.2.1}, se ha demostrado que el homomorfismo~(\hyperref[III.3.4.2.2]{3.4.2.2}) es \emph{biyectivo}.
Probemos después que $R^nf_*(\mathcal F)$ es un $\mathcal O_{\mathfrak Y}$-módulo coherente y, más precisamente, que se tiene
\[
R^nf_*(\mathcal F)=(H^n(\mathfrak X,\mathcal F))^\Delta.
\tag{3.4.5.1}\label{III.3.4.5.1-fr}
\]

Con la notación anterior, como $\mathcal F_k$ es un $\mathcal O_{X_k}$-módulo coherente~\sref{III.1.4.13}, se tiene
\[
\Gamma(\mathfrak Y_g,R^nf_*(\mathcal F_k))
=(\Gamma(\mathfrak Y,R^nf_*(\mathcal F_k)))_g
=(H^n(\mathfrak X,\mathcal F_k))_g.
\]

Ahora bien, los $H^n(\mathfrak X,\mathcal F_k)$ forman un sistema proyectivo que satisface (ML), y su límite proyectivo $H^n(\mathfrak X,\mathcal F)$ es un $A$-módulo de tipo finito.
Se concluye~\sref[0]{0.13.7.8} que se tiene
\[
\varprojlim_k((H^n(\mathfrak X,\mathcal F_k))_g)
=H^n(\mathfrak X,\mathcal F)\otimes_AA_{\{g\}}
=\Gamma(\mathfrak Y_g,(H^n(\mathfrak X,\mathcal F))^\Delta)
\]
teniendo en cuenta~\sref[I]{I.10.10.8} aplicado a $A$ y $A_{\{g\}}$; esto demuestra~(\hyperref[III.3.4.5.1]{3.4.5.1}), pues $\Gamma(\mathfrak Y_g,R^nf_*(\mathcal F))
=\varprojlim_k\Gamma(\mathfrak Y_g,R^nf_*(\mathcal F_k))$.

Como (\hyperref[III.3.4.2.2]{3.4.2.2}) es entonces un isomorfismo de $\mathcal O_{\mathfrak Y}$-módulos coherentes, es necesariamente un isomorfismo \emph{topológico}~\sref[I]{I.10.11.6}.
Finalmente, de las relaciones $R^nf_*(\mathcal F_k)=(H^n(X,\mathcal F_k))^\Delta$ se deduce que el sistema proyectivo $(R^nf_*(\mathcal F_k))_{k\geq0}$ satisface (ML)~\sref[I]{I.10.10.2}.

Una vez demostrados \sref{III.3.4.2} y \sref{III.3.4.3} en el caso en que el preesquema formal $\mathfrak Y$ es afín noetheriano, el paso al caso general es inmediato para~\sref{III.3.4.2} y para la primera afirmación de~\sref{III.3.4.3}, que son locales sobre $\mathfrak Y$.
En cuanto a la segunda afirmación de~\sref{III.3.4.3}, como $\mathfrak Y$ es noetheriano, basta recubrirlo por un número finito de abiertos afines noetherianos $U_i$ y observar que las restricciones del sistema proyectivo $(R^qf_*(\mathcal F_k))$ a cada uno de los $U_i$ satisfacen (ML).
\end{env}

En el curso de la demostración hemos probado además:
\begin{corollary}[3.4.6]
\label{III.3.4.6}
\label{III.3.4.6-fr}
Bajo las hipótesis del Corolario~\sref{III.3.4.4}, el homomorfismo canónico
\[
H^q(\mathfrak X,\mathcal F)
\longrightarrow\Gamma(\mathfrak Y,R^qf_*(\mathcal F))
\tag{3.4.6.1}\label{III.3.4.6.1-fr}
\]
es biyectivo.
\end{corollary}

% END tree/repo/editions/III1_DIPLOMATIC/ega3/ega3-3.tex


% BEGIN tree/repo/editions/III1_DIPLOMATIC/ega3/ega3-4.tex
\section{El teorema fundamental de los morfismos propios. Aplicaciones}
\label{section:III.4}
\label{section:III.4-fr}


\subsection{El teorema fundamental}
\label{subsection:III.4.1}
\label{subsection:III.4.1-fr}

\begin{env}[4.1.1]
\label{III.4.1.1}
\label{III.4.1.1-fr}
Sean $X$ e $Y$ preesquemas ordinarios noetherianos, sea $f:X\to Y$ un
morfismo \emph{propio}, sea $Y'$ un subconjunto cerrado de $Y$, $X'$ su imagen inversa
$f^{-1}(Y')$. Designaremos por $\widehat X$ y $\widehat Y$ los preesquemas formales
$X_{/X'}$ y $Y_{/Y'}$ \emph{obtenidos completando} $X$ e $Y$ a lo largo de
$X'$ e $Y'$ respectivamente (\textbf{I}, 10.8.5), por $\widehat f$ la \emph{prolongación}
de $f$ a estos completados (\textbf{I}, 10.9.1), que es un morfismo
$\widehat X\to\widehat Y$ de preesquemas formales. Para todo
$\mathcal O_X$-módulo \emph{coherente} $\mathcal F$,
\oldpage[III]{123}
designaremos por $\widehat{\mathcal F}$ su \emph{completado} $\mathcal F_{/X'}$ a lo
largo de $X'$ \sref[I]{I.10.8.4}; este es un
$\mathcal O_{\widehat X}$-módulo coherente \sref[I]{I.10.8.8}.
\end{env}

\begin{env}[4.1.2]
\label{III.4.1.2}
\label{III.4.1.2-fr}
Sea $\mathcal J$ un ideal coherente de $\mathcal O_Y$ tal que
$\operatorname{Supp}(\mathcal O_Y/\mathcal J)=Y'$ (\textbf{I}, 5.2.1); se sabe (\textbf{I}, 4.4.5) que $\mathcal K=f^*(\mathcal J)\mathcal O_X$
es un ideal coherente de $\mathcal O_X$ que satisface
\[
\operatorname{Supp}(\mathcal O_X/\mathcal K)=X'.
\]
Consideraremos para todo $k\geq0$ los $\mathcal O_X$-módulo coherente
\[
\mathcal F_k
=\mathcal F\otimes_{\mathcal O_Y}(\mathcal O_Y/\mathcal J^{k+1})
=\mathcal F/\mathcal K^{k+1}\mathcal F.
\]
Los $\mathcal O_Y$-Módulos $R^nf_*(\mathcal F)$ y
$R^nf_*(\mathcal F_k)$ son \emph{coherentes} para todo $n$ \sref{III.3.2.1}.
Para todo $k\geq0$ y todo $n$, el homomorfismo canónico
$\mathcal F\to\mathcal F_k$ induce, por funtorialidad, un homomorfismo
\[
R^nf_*(\mathcal F)\longrightarrow R^nf_*(\mathcal F_k).
\tag{4.1.2.1}\label{III.4.1.2.1-fr}
\]
Como $\mathcal F_k$ es un $\mathcal O_X/\mathcal K^{k+1}$-módulo,
$R^nf_*(\mathcal F_k)$ es un
$\mathcal O_Y/\mathcal J^{k+1}$-Módulo (\textbf{0}, 12.2.1) y se deduce por tanto de (\hyperref[III.4.1.2.1-fr]{4.1.2.1}) un homomorfismo \[
R^nf_*(\mathcal F)\otimes_{\mathcal O_Y}
(\mathcal O_Y/\mathcal J^{k+1})
\longrightarrow R^nf_*(\mathcal F_k).
\tag{4.1.2.2}\label{III.4.1.2.2-fr}
\]
Los dos miembros de \eqref{III.4.1.2.2} forman sistemas proyectivos.
El límite proyectivo del miembro izquierdo es la \emph{completación}
$(R^nf_*(\mathcal F))_{/Y'}$, que denotaremos por
$(R^nf_*(\mathcal F))^{\widehat{\,}}$. Además, es inmediato que los
homomorfismos (\hyperref[III.4.1.2.2-fr]{4.1.2.2}) forman un sistema
proyectivo, de donde, por paso al límite, un \emph{homomorfismo canónico}
\[
\varphi_n:(R^nf_*(\mathcal F))^{\widehat{\,}}
\longrightarrow\varprojlim_kR^nf_*(\mathcal F_k).
\tag{4.1.2.3}\label{III.4.1.2.3-fr}
\]

Por lo demás, (\hyperref[III.4.1.2.2-fr]{4.1.2.2}) es un homomorfismo de
$(\mathcal O_Y/\mathcal J^{k+1})$-módulos.
Por \sref[I]{I.10.8.3}, puede considerarse, por tanto, como un
homomorfismo continuo de $\mathcal O_{\widehat Y}$-módulos topológicos
provistos de sus topologías seudodiscretas \sref[0]{0.3.8.1}.
En consecuencia, $\varphi_n$ es un homomorfismo \emph{continuo} de
$\mathcal O_Y$-módulos topológicos.
\end{env}

\begin{env}[4.1.3]
\label{III.4.1.3}
\label{III.4.1.3-fr}
Sea $i:\widehat X\to X$ el morfismo canónico de espacios anillados definido en (\textbf{I}, 10.8.7), de modo que se tiene el diagrama conmutativo \[
\xymatrix{
X_k\ar[r]^{h_k}\ar[dr]_{i_k}&\widehat X\ar[d]^i\\
&X
}
\tag{4.1.3.1}\label{III.4.1.3.1-fr}
\]
donde $X_k$ es el subpreesquema cerrado de $X$ definido por
$\mathcal K^{k+1}$, $i_k$ es la inmersión canónica, y $h_k$ es el
morfismo de espacios anillados que corresponde a la identidad sobre
los espacios subyacentes y al homomorfismo canónico
$\mathcal O_{\widehat X}\to\mathcal O_X/\mathcal K^{k+1}$
\sref[I]{I.10.5.2}.
Además,
$\widehat{\mathcal F}=i^*(\mathcal F)$ (\textbf{I}, 10.8.8) salvo isomorfismo canónico. Se sabe que se tiene \[
H^n(X_k,i_k^*(\mathcal F_k))=H^n(X,\mathcal F_k)
\tag{4.1.3.2}\label{III.4.1.3.2-fr}
\]
salvo isomorfismo canónico.

\oldpage[III]{124}

En efecto,
$\mathcal F_k=(i_k)_*(i_k^*(\mathcal F_k))$
(G,~II,~4.9.1).
El homomorfismo canónico
$H^n(\widehat X,\widehat{\mathcal F})\to
H^n(X_k,h_k^*(\widehat{\mathcal F}))$
\sref[0]{0.12.1.3.5} puede escribirse también
\[
H^n(\widehat X,\widehat{\mathcal F})\longrightarrow H^n(X,\mathcal F_k)
\tag{4.1.3.3}\label{III.4.1.3.3-fr}
\]
Estos homomorfismos forman un sistema proyectivo y definen, por tanto,
un homomorfismo canónico
\[
\psi_{n,X}:H^n(\widehat X,\widehat{\mathcal F})
\longrightarrow\varprojlim_kH^n(X,\mathcal F_k).
\tag{4.1.3.4}\label{III.4.1.3.4-fr}
\]
Reemplazando $X$ por un subconjunto abierto de la forma $f^{-1}(V)$,
donde $V$ es un abierto afín de $Y$, se tienen, teniendo en cuenta (\hyperref[III.1.4.11-fr]{1.4.11}), homomorfismos \[
\psi_{n,V}:H^n(\widehat X\cap f^{-1}(V),\widehat{\mathcal F})
\longrightarrow\varprojlim_k\Gamma(V,R^nf_*(\mathcal F_k))~;
\tag{4.1.3.5}\label{III.4.1.3.5-fr}
\]
Estos homomorfismos conmutan manifiestamente con la restricción de $V$
a un subconjunto abierto afín más pequeño.
Definen, por tanto, un \emph{homomorfismo canónico} de haces
\[
\psi_n:R^n\widehat f_*(\widehat{\mathcal F})
\longrightarrow\varprojlim_kR^nf_*(\mathcal F_k).
\tag{4.1.3.6}\label{III.4.1.3.6-fr}
\]
\end{env}

\begin{env}[4.1.4]
\label{III.4.1.4}
\label{III.4.1.4-fr}
Por último, sea $j:\widehat Y\to Y$ el morfismo canónico de espacios
anillados \sref[I]{I.10.8.7}.
Como $R^nf_*(\mathcal F)$ es un $\mathcal O_Y$-Módulo coherente (\hyperref[III.3.2.1-fr]{3.2.1}), se tiene $j^*(R^nf_*(\mathcal F))=R^nf_*(\mathcal F))^{\widehat{\,}}$ salvo isomorfismo canónico (\textbf{I}, 10.8.8) y se tiene por tanto un homomorfismo canónico \[
\rho_n:(R^nf_*(\mathcal F))^{\widehat{\,}}=j^*(R^nf_*(\mathcal F))
\longrightarrow R^n\widehat f_*(i^*(\mathcal F))
=R^n\widehat f_*(\widehat{\mathcal F})
\tag{4.1.4.1}\label{III.4.1.4.1-fr}
\]
definido de manera general para los espacios anillados (véase la demostración
de (\hyperref[III.1.4.15-fr]{1.4.15})). Mostremos que el diagrama
\[
\xymatrix{
(R^nf_*(\mathcal F))^{\widehat{\,}}\ar[rr]^{\rho_n}\ar[dr]_{\varphi_n}
&&R^n\widehat f_*(\widehat{\mathcal F})\ar[dl]^{\psi_{\boldsymbol{n}}}\\
&\displaystyle\varprojlim_kR^nf_*(\mathcal F_k)&
}
\tag{4.1.4.2}\label{III.4.1.4.2-fr}
\] es \emph{conmutativo}. Basta evidentemente probar la conmutatividad del diagrama correspondiente de homomorfismos de prehaces, luego se puede limitar al caso en que $Y$ es afín, y todo se reduce a probar que el diagrama \[
\xymatrix{
(H^n(X,\mathcal F))^{\widehat{\,}}\ar[rr]^{\rho_n}\ar[dr]_{\varphi_n}
&&H^n(\widehat X,\widehat{\mathcal F})\ar[dl]^{\psi_{n,X}}\\
&\displaystyle\varprojlim_kH^n(X,\mathcal F_k)&
}
\tag{4.1.4.3}\label{III.4.1.4.3-fr}
\]
es conmutativo. Pero la conmutatividad de
(\hyperref[III.4.1.3.1-fr]{4.1.3.1}) y las relaciones vistas en
(\hyperref[III.4.1.3-fr]{4.1.3}) entre los grupos de cohomología dan
inmediatamente el diagrama conmutativo
\oldpage[III]{125}
\[
\xymatrix{
H^n(X,\mathcal F)\ar[rr]\ar[dr]&&
H^n(\widehat X,\widehat{\mathcal F})
=H^n(\widehat X,i^*(\mathcal F))\ar[dl]\\
&H^n(X,\mathcal F_k)=H^n(X_k,i_k^*(\mathcal F))&
}
\]
de donde se deduce inmediatamente la conmutatividad de
(\hyperref[III.4.1.4.3-fr]{4.1.4.3}).
\end{env}

\begin{theorem}[4.1.5]
\label{III.4.1.5}
\label{III.4.1.5-fr}
Sea $f:X\to Y$ un morfismo propio de preesquemas noetherianos, sea
$Y'$ un subconjunto cerrado de $Y$, $X'=f^{-1}(Y')$.
Entonces, para todo $\mathcal O_X$-módulo coherente $\mathcal F$,
$R^n\widehat f_*(\widehat{\mathcal F})$
es un $\mathcal O_{\widehat X}$-módulo coherente, y los homomorfismos
$\varphi_n$, $\psi_n$ y $\rho_n$ del diagrama
\eqref{III.4.1.4.2} son isomorfismos topológicos.
\end{theorem}

\begin{proof}
Basta demostrar que $\varphi_n$ y $\psi_n$ son isomorfismos.
En efecto, $R^nf_*(\mathcal F)$ es coherente \sref{III.3.2.1}, luego
su completación $(R^nf_*(\mathcal F))^{\widehat{\,}}$ es coherente
\sref[I]{I.10.8.8}; la bicontinuidad de $\varphi_n$, $\psi_n$ y
$\rho_n$ es entonces automática \sref[I]{I.10.11.6}.
\end{proof}

\begin{remarks}[4.1.6]
\label{III.4.1.6-fr}
\upshape
\begin{enumerate}
\item[(i)] Si se pone $\mathcal F_k=
\widehat{\mathcal F}/\widehat{\mathcal K}^{k+1}
\widehat{\mathcal F}$,
  Entonces, inmediatamente,
  $\mathcal F_k=i_*(\widehat{\mathcal F}_k)$,
  y el homomorfismo canónico \eqref{III.4.1.3.6} es precisamente el
  homomorfismo ya definido en (\hyperref[III.3.4.2.2]{3.4.2.2}):
  \[
R^n\widehat f_*(\widehat{\mathcal F})
\longrightarrow\varprojlim_kR^n\widehat f_*(\widehat{\mathcal F}_k)~;
\tag{4.1.6.1}\label{III.4.1.6.1-fr}
\]
  Así, el hecho de que $\psi_n$ sea un isomorfismo es un caso particular
  de \sref{III.3.4.3}.
  Daremos, no obstante, más adelante una demostración directa, evitando
  las delicadas consideraciones relativas a los límites proyectivos de
  sucesiones espectrales \sref[0]{0.13.7} en las que se apoya el teorema
  general \sref{III.3.4.3}.

  \item[(ii)] Como los $\psi_n$ son isomorfismos, es equivalente decir que los $\varphi_n$ son isomorfismos equivale a decir que
  $\rho_n=\psi_n^{-1}\circ\varphi_n$
  son isomorfismos.
  El Teorema~\sref{III.4.1.5} afirma, en particular, que \emph{la formación de}
  $R^nf_*$ \emph{conmuta con la completación}; puede llamarse, por tanto, el
  \emph{primer teorema de comparación} entre las teorías ``algebraica'' y
  ``formal''.
\end{enumerate}
Comenzaremos demostrando la forma afín de \sref{III.4.1.5}.
\end{remarks}

\begin{corollary}[4.1.7]
\label{III.4.1.7-fr}
Siendo las hipótesis las de \textup{(\hyperref[III.4.1.5-fr]{4.1.5})}, supongamos
además $Y=\operatorname{Spec}(A)$, donde $A$ es noetheriano, y
$\mathcal J=\widetilde{\mathfrak J}$, donde $\mathfrak J$ es un ideal de
$A$, de modo que
$\mathcal F_k=\mathcal F/\mathfrak J^{k+1}\mathcal F$.
El homomorfismo canónico
\[
\varphi_n:(H^n(X,\mathcal F))^{\widehat{\,}}
\longrightarrow\varprojlim_kH^n(X,\mathcal F_k)
\tag{4.1.7.1}\label{III.4.1.7.1-fr}
\]
(donde el primer miembro es el separado completado de $H^n(X,\mathcal F)$ para
la topología $\mathfrak J$-preádica) es un isomorfismo. El sistema
proyectivo $(H^n(X,\mathcal F_k))_{k\geq0}$ verifica para todo $n$ la
condición \textup{(ML)}, y el homomorfismo canónico
\[
\psi_n:H^n(\widehat X,\widehat{\mathcal F})
\longrightarrow\varprojlim_{\boldsymbol{k}}H^n(X,\mathcal F_k)
\tag{4.1.7.2}\label{III.4.1.7.2-fr}
\]
es un isomorfismo. Finalmente, la filtración sobre $H^n(X,\mathcal F)$, definida
por los núcleos de los homomor%
\oldpage[III]{126}
fismos canónicos
$H^n(X,\mathcal F)\to H^n(X,\mathcal F_k)$, es $\mathfrak J$-buena
\textup{(\textbf{0}, 13.7.7)} y $\varphi_n$ es un isomorfismo topológico\footnote{La demostración que sigue, más sencilla que la demostración
inicial, y el complemento relativo a la filtración de
$H^n(X,\mathcal F)$, nos fueron comunicados por J.-P. Serre.}.
\end{corollary}

\begin{proof}
Al estar fijado el entero $n\geq0$ en esta demostración, pondremos para
simplificar
\[
\mathrm H=H^n(X,\mathcal F),\qquad
\mathrm H_k=H^n(X,\mathcal F_k)
\tag{4.1.7.3}\label{III.4.1.7.3-fr}
\]
\[
\mathrm R_k=\operatorname{Ker}(\mathrm H\longrightarrow\mathrm H_k),
\qquad\text{sub-$A$-módulo de $\mathrm H$.}
\tag{4.1.7.4}\label{III.4.1.7.4-fr}
\]

La sucesión exacta de cohomología
\[
H^n(X,\mathfrak J^{k+1}\mathcal F)\longrightarrow H^n(X,\mathcal F)
\longrightarrow H^n(X,\mathcal F_k)\xrightarrow{\partial}
H^{n+1}(X,\mathfrak J^{k+1}\mathcal F)
\longrightarrow H^{n+1}(X,\mathcal F)
\]
muestra que se tiene también
$\mathrm R_k=\operatorname{Im}(
H^n(X,\mathfrak J^{k+1}\mathcal F)\longrightarrow H^n(X,\mathcal F)
)$;
pondremos
\[
\begin{aligned}
\mathrm Q_k
&=\operatorname{Ker}\bigl(H^{n+1}(X,\mathfrak J^{k+1}\mathcal F)
\longrightarrow H^{n+1}(X,\mathcal F)\bigr)\\
&=\operatorname{Im}\bigl(H^n(X,\mathcal F_k)
\longrightarrow H^{n+1}(X,\mathfrak J^{k+1}\mathcal F)\bigr).
\end{aligned}
\tag{4.1.7.5}\label{III.4.1.7.5-fr}
\]

Se tiene pues la sucesión exacta
\[
0\longrightarrow\mathrm R_k\longrightarrow\mathrm H
\longrightarrow\mathrm H_k\longrightarrow\mathrm Q_k\longrightarrow0.
\tag{4.1.7.6}\label{III.4.1.7.6-fr}
\]

\begin{env}[4.1.7.7]
\label{III.4.1.7.7-fr}
Sea $x$ un elemento de $\mathfrak J^m$ ($m\geq0$); la multiplicación por
$\boldsymbol{x}$ en $\mathfrak J^k\mathcal F$ es un homomorfismo
$\mathfrak J^k\mathcal F\to\mathfrak J^{k+m}\mathcal F$, y da lugar
por tanto a un homomorfismo
\[
\mu_{x,m}:H^n(X,\mathfrak J^k\mathcal F)
\longrightarrow H^n(X,\mathfrak J^{k+m}\mathcal F).
\tag{4.1.7.8}\label{III.4.1.7.8-fr}
\]

Si se designa por $S$ la $A$-álgebra graduada
$\bigoplus_{k\geq0}\mathfrak J^k$, se sabe que las multiplicaciones
$\mu_{x,m}$ definen sobre
$\mathrm E=\bigoplus_{k\geq0}H^n(X,\mathfrak J^k\mathcal F)$ una
estructura de \emph{módulo graduado de tipo finito} sobre el anillo graduado $S$
(\hyperref[III.3.3.2-fr]{3.3.2}), que es \emph{noetheriano} (\textbf{II}, 2.1.5).
\end{env}

\begin{lemma}[4.1.7.9]
\label{III.4.1.7.9-fr}
Los submódulos $\mathrm R_k$ de $\mathrm H$ definen sobre $\mathrm H$
una filtración $\mathfrak J$-buena.
\end{lemma}

\begin{proof}
En primer lugar, mostremos que se tiene
\[
\mathfrak J^m\mathrm R_k\subset\mathrm R_{\boldsymbol{k}+m},
\tag{4.1.7.10}\label{III.4.1.7.10-fr}
\]
siendo pues la multiplicación en $\mathrm H=H^n(X,\mathcal F)$ por un elemento
$x\in\mathfrak J^m$ la aplicación $\mu_{x,0}$.

Para todo
$x\in\mathfrak J^m$, el diagrama
\[
\xymatrix{
\mathfrak J^{k+1}\mathcal F\ar[r]^x\ar[d]&
\mathfrak J^{k+m+1}\mathcal F\ar[d]\\
\mathcal F\ar[r]_x&\mathcal F
}
\]
\oldpage[III]{127}
(donde las flechas horizontales son la multiplicación por $x$, y las
flechas verticales las inclusiones canónicas) es conmutativo; luego el
diagrama correspondiente
\[
\xymatrix{
H^n(X,\mathfrak J^{k+1}\mathcal F)\ar[r]^{\mu_{x,m}}\ar[d]&
H^n(X,\mathfrak J^{k+m+1}\mathcal F)\ar[d]\\
H^n(X,\mathcal F)\ar[r]_{\mu_{x,0}}&H^n(X,\mathcal F)
}
\tag{4.1.7.11}\label{III.4.1.7.11-fr}
\]
es conmutativo, lo que, teniendo en cuenta la interpretación de $\mathrm R_k$
como imagen de
$H^n(X,\mathfrak J^{k+1}\mathcal F)\to H^n(X,\mathcal F)$, prueba
(\hyperref[III.4.1.7.10-fr]{4.1.7.10}) y muestra además que el $S$-módulo
graduado $\mathrm R=\bigoplus_{k\geq0}\mathrm R_k$ es un \emph{cociente} del
sub-$S$-módulo
$\mathrm M=\bigoplus_{k\geq0}H^n(X,\mathfrak J^{k+1}\mathcal F)$ de
$\mathrm E$; la observación hecha más arriba muestra pues que $\mathrm R$ es
un $S$-módulo \emph{de tipo finito}, lo que equivale a la afirmación de
(\hyperref[III.4.1.7.9-fr]{4.1.7.9}) (Bourbaki, \emph{Alg. comm.},
cap. III, §~3, n\textsuperscript{o}~1, teorema~1).
\end{proof}

\begin{env}[4.1.7.12]
\label{III.4.1.7.12-fr}
Consideremos ahora el $S$-módulo graduado
$\mathrm N=\bigoplus_{k\geq0}H^{n+1}(X,\mathfrak J^{k+1}\mathcal F)$
definido como en (4.7.1.8); es también un $S$-módulo \emph{de tipo finito} en
virtud de (\hyperref[III.3.3.2-fr]{3.3.2}); se tiene
$\mathrm Q_k\subset\mathrm N_k$ para todo $k$ por
(\hyperref[III.4.1.7.5-fr]{4.1.7.5}) y el diagrama
(\hyperref[III.4.1.7.11-fr]{4.1.7.11}), donde se reemplaza $n$ por $n+1$,
muestra que
$S_m\mathrm Q_k=\mathfrak J^m\mathrm Q_k\subset\mathrm Q_{k+m}$.
En otros términos, $\mathrm Q$ es un \emph{sub-$S$-módulo graduado} de $\mathrm N$,
y por tanto es \emph{de tipo finito}.
\end{env}

\begin{env}[4.1.7.13]
\label{III.4.1.7.13-fr}
Designemos por $\alpha_m$ la inclusión canónica
$\mathfrak J^m\to A$, que se puede escribir $S_m\to S_0$. Como
$\mathfrak J^{k+1}\mathcal F_k=0$, el $A$-módulo
$H^n(X,\mathcal F_k)$ está anulado por $\mathfrak J^{k+1}$; como
$\mathrm Q_k$ es la imagen del $A$-homomorfismo
$H^n(X,\mathcal F_k)\to H^{n+1}(X,\mathfrak J^{k+1}\mathcal F)$,
$\mathrm Q_k$, como $A$-módulo, está también anulado por
$\mathfrak J^{k+1}$; esto significa también que, en el $S$-\emph{módulo}
$\mathrm Q$, se tiene
\[
\alpha_{k+1}(S_{k+1})\mathrm Q_k=0.
\tag{4.1.7.14}\label{III.4.1.7.14-fr}
\]

Como $\mathrm Q$ es un $S$-módulo de tipo finito, existe un entero
$k_0$ y un entero $h$ tales que
$\mathrm Q_{k+h}=S_h\mathrm Q_k$ para $k\geq k_0$ (\textbf{II}, 2.1.6, (ii));
se deduce de esta relación y de
(\hyperref[III.4.1.7.14-fr]{4.1.7.14}) que existe un entero $r>0$ tal
que
\[
\alpha_r(S_r)\mathrm Q=0.
\tag{4.1.7.15}\label{III.4.1.7.15-fr}
\]
\end{env}

\begin{env}[4.1.7.16]
\label{III.4.1.7.16-fr}
Observemos ahora que la inclusión canónica
$\mathfrak J^{k+m}\mathcal F\to\mathfrak J^k\mathcal F$ da, pasando
a la cohomología, un $A$-homomorfismo
\[
\nu_m:H^{n+1}(X,\mathfrak J^{k+m}\mathcal F)
\longrightarrow H^{n+1}(X,\mathfrak J^k\mathcal F)
\tag{4.1.7.17}\label{III.4.1.7.17-fr}
\]
y, para todo $x\in\mathfrak J^m$, se tiene evidentemente la factorización
\[
\xymatrix@C=4.5em{
\mu_{x,0}:H^{n+1}(X,\mathfrak J^k\mathcal F)\ar[r]^{\mu_{x,m}}&
H^{n+1}(X,\mathfrak J^{k+m}\mathcal F)\ar[r]^{\nu_m}&
H^{n+1}(X,\mathfrak J^k\mathcal F)
}
\tag{4.1.7.18}\label{III.4.1.7.18-fr}
\]
\oldpage[III]{128}
de donde se concluye que, para todo sub-$A$-módulo $P$ de
$H^{n+1}(X,\mathfrak J^k\mathcal F)$, se tiene, en el $S$-módulo
$\mathrm N$,
\[
\nu_m(S_mP)=\alpha_m(S_m)P.
\tag{4.1.7.19}\label{III.4.1.7.19-fr}
\]
\end{env}

\begin{lemma}[4.1.7.20]
\label{III.4.1.7.20-fr}
Existe un entero $m>0$ tal que $\nu_m(\mathrm Q_{k+m})=0$ para todo
$k\geq k_0$.
\end{lemma}

\begin{proof}
Tomemos en efecto como $m$ un múltiplo de $h$ que sea $\geq r$; como
$\mathrm Q_{k+m}=S_m\mathrm Q_k$ para $k\geq k_0$, se tiene en virtud de
(\hyperref[III.4.1.7.19-fr]{4.1.7.19}) y
(\hyperref[III.4.1.7.15-fr]{4.1.7.15})
$\nu_m(\mathrm Q_{k+m})=\alpha_m(S_m)\mathrm Q_k
\subset\alpha_r(S_r)\mathrm Q_k=0$.
\end{proof}

\begin{env}[4.1.7.21]
\label{III.4.1.7.21-fr}
Observemos que del diagrama conmutativo
\[
\xymatrix@C=2.5em{
H^n(X,\mathcal F)\ar[r]&H^n(X,\mathcal F_k)\ar[r]^-\partial&
H^{n+1}(X,\mathfrak J^{k+1}\mathcal F)\ar[r]&H^{n+1}(X,\mathcal F)\\
H^n(X,\mathcal F)\ar[r]\ar[u]&H^n(X,\mathcal F_{k+m})\ar[r]^-\partial\ar[u]&
H^{n+1}(X,\mathfrak J^{k+m+1}\mathcal F)\ar[r]\ar[u]&
H^{n+1}(X,\mathcal F)\ar[u]
}
\]
procedente a su vez del diagrama conmutativo
\[
\xymatrix{
0\ar[r]&\mathfrak J^{k+1}\mathcal F\ar[r]&\mathcal F\ar[r]&
\mathcal F_k\ar[r]&0\\
0\ar[r]&\mathfrak J^{k+m+1}\mathcal F\ar[r]\ar[u]&
\mathcal F\ar[r]\ar[u]&\mathcal F_{\boldsymbol{k}+m}\ar[r]\ar[u]&0
}
\]
donde las flechas verticales son las aplicaciones canónicas, se deduce un
diagrama conmutativo
\[
\xymatrix@C=2.2em{
0\ar[r]&\mathrm R_k\ar[r]&\mathrm H\ar[r]&\mathrm H_k\ar[r]&
\mathrm Q_k\ar[r]&0\\
0\ar[r]&\mathrm R_{k+m}\ar[r]\ar[u]&\mathrm H\ar[r]\ar[u]^{\textit{id.}}&
\mathrm H_{k+m}\ar[r]\ar[u]&\mathrm Q_{k+m}\ar[r]\ar[u]_{\nu_m}&0
}
\]
donde las filas son exactas. Como la última flecha vertical es nula
para $k\geq k_0$ (\hyperref[III.4.1.7.20-fr]{4.1.7.20}), la imagen de
$\mathrm H_{k+m}$ en $\mathrm H_k$ está contenida en
$\operatorname{Ker}(\mathrm H_k\to\mathrm Q_k)
=\operatorname{Im}(\mathrm H\to\mathrm H_k)$, pero por otra parte
contiene $\operatorname{Im}(\mathrm H\to\mathrm H_k)$ por la conmutatividad
del diagrama, luego es igual a él; lo mismo ocurre pues con las imágenes
en $\mathrm H_k$ de los $\mathrm H_{k'}$ para $k'\geq k+m$, lo que demuestra
la condición (ML) para el sistema proyectivo $(\mathrm H_k)_{k\geq0}$. Además,
para todo abierto afín $U$ de $X$ se tiene
$H^i(U,\mathcal F_k)=0$ para $i>0$
(\hyperref[III.1.3.1-fr]{1.3.1}), y para $m>0$ la aplicación
$H^0(U,\mathcal F_{k+m})\to H^0(U,\mathcal F_k)$ es sobreyectiva
(\textbf{I}, 1.3.9). Se puede pues apli%
\oldpage[III]{129}
car (\textbf{0}, 13.3.1), y el homomorfismo canónico
$H^n(\widehat X,\widehat{\mathcal F})\to
\varprojlim_kH^n(X,\mathcal F_k)$ es \emph{biyectivo} para todo $n\geq0$.

Como el sistema proyectivo $(\mathrm H/\mathrm R_k)_{k\geq0}$ es estricto,
se puede pasar al límite proyectivo en las sucesiones exactas
\[
0\longrightarrow\mathrm H/\mathrm R_k\longrightarrow\mathrm H_k
\longrightarrow\mathrm Q_k\longrightarrow0
\tag{4.1.7.22}\label{III.4.1.7.22-fr}
\]
(\textbf{0}, 13.2.2); como $\nu_m(\mathrm Q_{k+m})=0$, se tiene
$\varprojlim_k\mathrm Q_k=0$, de donde un isomorfismo topológico
$\varprojlim_k(\mathrm H/\mathrm R_k)\xrightarrow{\sim}
\varprojlim_k\mathrm H_k$.
Pero como la filtración $(\mathrm R_k)$ de $\mathrm H$ es
$\mathfrak J$-buena, define sobre $\mathrm H$ la topología
$\mathfrak J$-preádica; luego $\varprojlim_k(\mathrm H/\mathrm R_k)$ es
el separado completado de $\mathrm H$ para la topología
$\mathfrak J$-preádica, lo que termina de demostrar
(\hyperref[III.4.1.7-fr]{4.1.7}).
\end{env}
\end{proof}

\begin{env}[4.1.8]
\label{III.4.1.8-fr}
Pasemos finalmente a la demostración de (\hyperref[III.4.1.5-fr]{4.1.5}): para todo abierto afín $V$ de $Y$,
$\Gamma(V,(R^nf_*(\mathcal F))^{\widehat{\,}})$
es la completación separada de
$\Gamma(V,R^nf_*(\mathcal F))$ para la topología
$\mathfrak J$-preádica, si
$\mathcal J|V=\widetilde{\mathfrak J}$, porque
$R^nf_*(\mathcal F)$ es un $\mathcal O_Y$-Módulo coherente (\textbf{I}, 10.8.4),
y $\Gamma(V,\varprojlim_kR^nf_*(\mathcal F_k))$ es igual a
$\varprojlim_k\Gamma(V,R^nf_*(\mathcal F_k))$
(\textbf{0}\textsubscript{I}, 3.2.6); el hecho de que $\varphi_n$ es un isomorfismo topológico resulta entonces de (\hyperref[III.4.1.7-fr]{4.1.7}) y de (\hyperref[III.1.4.11-fr]{1.4.11}). Por otra parte (siempre en virtud de (\hyperref[III.1.4.11-fr]{1.4.11})), resulta de (\hyperref[III.4.1.7-fr]{4.1.7}) que el homomorfismo $\psi_{n,V}$ de
(\hyperref[III.4.1.3.3-fr]{4.1.3.3}) es un \emph{isomorfismo}, luego $\psi_n$
es un isomorfismo por definición de
$R^n\widehat f_*(\widehat{\mathcal F})$.
\end{env}

\begin{corollary}[4.1.9]
\label{III.4.1.9}
\label{III.4.1.9-fr}
Bajo las hipótesis de \sref{III.4.1.4}, para todo subconjunto abierto
afín $V$ de $Y$, el homomorfismo canónico
\[
H^n(\widehat X\cap f^{-1}(V),\widehat{\mathcal F})
\longrightarrow\Gamma(\widehat Y\cap V,
R^n\widehat f_*(\widehat{\mathcal F}))
\]
es biyectivo.
\end{corollary}

\begin{remark}[4.1.10]
\label{III.4.1.10}
\label{III.4.1.10-fr}
Sea $f:X\to Y$ un morfismo \emph{de tipo finito} entre preesquemas ordinarios
noetherianos, y sea $\mathcal F$ un $\mathcal O_X$-módulo coherente cuyo
soporte es \emph{propio sobre} $Y$ \sref[II]{II.5.4.10}.
Se sabe entonces \sref{III.3.2.4} que $R^nf_*(\mathcal F)$ es un
$\mathcal O_Y$-módulo coherente para todo $n\geq0$. Además, se puede
siempre suponer que $\mathcal F=u_*(\mathcal G)$, donde
$\mathcal G=u^*(\mathcal F)$ es un $\mathcal O_Z$-módulo coherente, $Z$ es un
subpreesquema cerrado adecuado de $X$ cuyo espacio subyacente es
$\operatorname{Supp}(\mathcal F)$, y $u:Z\to X$ es la inyección canónica
\sref[I]{I.9.3.5}.
Si
$\mathcal G_k=\mathcal G\otimes_{\mathcal O_Y}
(\mathcal O_Y/\mathcal J^{k+1})$,
se tiene $\mathcal G_k=u^*(\mathcal F_k)$,
$R^nf_*(\mathcal F_k)=R^n(f\circ u)_*(\mathcal G_k)$ y
$R^nf_*(\mathcal F)=R^n(f\circ u)_*(\mathcal G)$
por \sref{III.1.3.4}; y, teniendo en cuenta \sref[I]{I.10.9.5},
\[
R^n\widehat f_*(\widehat{\mathcal F})
=R^n(f\circ u)^{\widehat{\,}}_*(\mathcal G).
\]
Se puede, por tanto, aplicar \sref{III.4.1.5} a $\mathcal G$ y al morfismo
propio $f\circ u$.
Se sigue que, bajo estas hipótesis, \emph{los resultados de}
\sref{III.4.1.5} \emph{siguen siendo válidos para} $\mathcal F$ \emph{y} $f$.
\end{remark}


\subsection{Casos particulares y variantes}
\label{subsection:III.4.2}
\label{subsection:III.4.2-fr}

La forma más útil del teorema de comparación \sref{III.4.1.5} es la
siguiente.

\begin{proposition}[4.2.1]
\label{III.4.2.1-fr}
Sean $Y$ un preesquema localmente noetheriano, $f:X\to Y$ un morfismo
propio, $\mathcal F$ un $\mathcal O_X$-Módulo coherente. Entonces, para todo
$y\in Y$ y todo $p$, $(R^pf_*(\mathcal F))_y$ es
\oldpage[III]{130}
un $\mathcal O_y$-módulo de tipo finito, luego separado para la topología
$\mathfrak m_y$-preádica, y se tiene un isomorfismo topológico canónico
\[
((R^pf_*(\mathcal F))_y)^{\widehat{\,}}
\xrightarrow{\sim}
\varprojlim_kH^p\bigl(f^{-1}(y),
\mathcal F\otimes_{\mathcal O_Y}(\mathcal O_y/\mathfrak m_y^k)\bigr)
\tag{4.2.1.1}\label{III.4.2.1.1-fr}
\]
donde el primer miembro es el completado de $(R^pf_*(\mathcal F))_y$ para la
topología $\mathfrak m_y$-preádica, y en el segundo miembro $f^{-1}(y)$ se
considera, para todo $k\geq0$, como espacio subyacente al preesquema
$X\times_Y\operatorname{Spec}(\mathcal O_y/\mathfrak m_y^k)$
\textup{(\textbf{I}, 3.6.1)}.
\end{proposition}

\begin{proof}
Como $\mathcal O_y$ es un anillo local noetheriano y
$(R^pf_*(\mathcal F))_y$ un $\mathcal O_y$-módulo de tipo finito
(\hyperref[III.3.2.1-fr]{3.2.1}), la topología $\mathfrak m_y$-preádica
sobre $(R^pf_*(\mathcal F))_y$ es separada (\textbf{0}\textsubscript{I}, 7.3.5). Las
demás afirmaciones son consecuencias de (\hyperref[III.4.1.7-fr]{4.1.7})
cuando $Y$ es noetheriano y el punto $y$ cerrado, reemplazando $Y$ por una
vecindad afín de $y$ y tomando $Y'=\{y\}$, teniendo en cuenta
(G, II, 4 9.1). En el caso general, pongamos
$Y_1=\operatorname{Spec}(\mathcal O_y)$, $X_1=X\times_YY_1$,
$\mathcal F_1=\mathcal F\otimes_{\mathcal O_Y}\mathcal O_{Y_1}$
y sea $f_1=f\times1_{Y_1}:X_1\to Y_1$; $Y_1$ es noetheriano y $f_1$
propio (\textbf{II}, 5.4.2, (iii)) y $\mathcal F_1$ es coherente
(\textbf{0}\textsubscript{I}, 5.3.11). Sea $y_1$ el único punto cerrado de $Y_1$;
la proposición es válida para $f_1$, $\mathcal F_1$ y $y_1$; se tiene
$\mathcal O_{y_1}=\mathcal O_y$, $f_1^{-1}(y_1)=f^{-1}(y)$ (\textbf{I}, 3.6.5),
estando canónicamente identificados los preesquemas
$X\times_Y\operatorname{Spec}(\mathcal O_y/\mathfrak m_y^k)$ y
$X_1\times_{Y_1}\operatorname{Spec}(\mathcal O_y/\mathfrak m_y^k)$
(\textbf{I}, 3.3.9); además,
$\mathcal F_1\otimes_{\mathcal O_{Y_1}}
(\mathcal O_y/\mathfrak m_y^k)$ se identifica con
$\mathcal F\otimes_{\mathcal O_Y}(\mathcal O_y/\mathfrak m_y^k)$
(\textbf{I}, 9.1.6). Queda por ver que $R^pf_{1*}(\mathcal F_1)$ es
canónicamente isomorfo a
$R^pf_*(\mathcal F)\otimes_{\mathcal O_Y}\mathcal O_{Y_1}$, lo que
resulta de (\hyperref[III.1.4.15-fr]{1.4.15}), siendo plano el morfismo local
$\operatorname{Spec}(\mathcal O_y)\to Y$
(\textbf{0}\textsubscript{I}, 6.7.1 y \textbf{I}, 2.4.2).
\end{proof}

El corolario siguiente utiliza la terminología de la teoría de la dimensión
(cap. IV) y no se aplicará antes del cap. IV.

\begin{corollary}[4.2.2]
\label{III.4.2.2}
\label{III.4.2.2-fr}
Sea $Y$ un preesquema localmente noetheriano, sea $f:X\to Y$ un morfismo
propio, $y$ un punto de $Y$, $r$ la dimensión de
$f^{-1}(y)$.
Entonces, para todo $\mathcal O_X$-módulo coherente $\mathcal F$, los haces
$R^pf_*(\mathcal F)$ se anulan en una vecindad de $y$ para todo
$p>r$.
\end{corollary}

\begin{proof}
En efecto, se tiene entonces
$H^p(f^{-1}(y),
\mathcal F\otimes(\mathcal O_y/\mathfrak m_y^k))=0$
(G, II, 4.15.2) para todo $k$, luego
(\hyperref[III.4.2.1-fr]{4.2.1}) el separado completado de
$(R^pf_*(\mathcal F))_y$ para la topología
$\mathfrak m_y$-preádica es nulo, y como esta topología es separada, se tiene también $(R^pf_*(\mathcal F))_y=0$
La conclusión se sigue porque $R^pf_*(\mathcal F)$ es coherente
\sref[0]{0.5.2.2}.
\end{proof}

\begin{env}[4.2.3]
\label{III.4.2.3}
\label{III.4.2.3-fr}
El resultado \sref{III.4.2.1} se utiliza especialmente cuando $p=0$; se obtiene pues el corolario siguiente:
\end{env}

\begin{corollary}[4.2.4]
\label{III.4.2.4}
\label{III.4.2.4-fr}
Bajo las hipótesis de \sref{III.4.2.1}, existe un isomorfismo
topológico canónico
\[
((f_*(\mathcal F))_y)^{\widehat{\,}}
\xrightarrow{\sim}
\varprojlim_k\Gamma\bigl(f^{-1}(y),
\mathcal F_y/\mathfrak m_y^k\mathcal F_y\bigr).
\]
\end{corollary}


\subsection{Teorema de conexidad de Zariski}
\label{subsection:III.4.3}
\label{subsection:III.4.3-fr}
\label{III.4.3}
\label{III.4.3-fr}

Los resultados de esta sección y de la siguiente generalizan teoremas
bien conocidos de Zariski, y todos pueden deducirse de
\sref{III.4.2.4}. Son consecuencias del teorema siguiente.

\begin{theorem}[4.3.1]
\label{III.4.3.1-fr}
\textnormal{(teorema de conexión).}
Sean $Y$ un preesquema localmente noetheriano, $f:X\to Y$ un morfismo
propio. Entonces $\mathcal A(X)=f_*(\mathcal O_X)$ es una
$\mathcal O_Y$-Álgebra coherente.
\oldpage[III]{131}
Sea $Y'$ el $Y$-esquema finito sobre $Y$ tal que
$\mathcal A(Y')=\mathcal A(X)$, que está determinado salvo $Y$-isomorfismo
\textup{(\textbf{II}, 1.3.1 y 6.1.3)}; si $f'=\mathcal A(e)$ es el $Y$-morfismo
$X\to Y'$ deducido del isomorfismo identidad
$e:\mathcal A(Y')\to\mathcal A(X)$ \textup{(\textbf{II}, 1.2.7)}, $f'$ es propio,
$f'_*(\mathcal O_X)$ es isomorfo a $\mathcal O_{Y'}$, y las fibras
$f'^{-1}(y')$ del morfismo $f'$ son conexas y no vacías para todo
$y'\in Y'$.
\end{theorem}

\begin{proof}
Sea $g:Y'\to Y$ el morfismo estructural. Para probar que el homomorfismo
$\theta:\mathcal O_{Y'}\to f'_*(\mathcal O_X)$ que interviene en la definición
del morfismo $f'$ es biyectivo, basta, puesto que $Y'$ es afín sobre $Y$,
probar que
$g_*(\theta):g_*(\mathcal O_{Y'})\longrightarrow
g_*(f'_*(\mathcal O_X))=f_*(\mathcal O_X)$
es la identidad (\textbf{II}, 1.4.2); pero esto resulta de las definiciones puesto que
$g_*(\mathcal O_{Y'})=\mathcal A(Y')$ y
$f_*(\mathcal O_X)=\mathcal A(X)$. El hecho de que $\mathcal A(X)$ es
coherente es un caso particular del teorema de finitud
(\hyperref[III.3.2.1-fr]{3.2.1}). Como $f$ es propio y $g$ separado,
$f'$ es propio (\textbf{II}, 5.4.3, (i)); para terminar la demostración de
(\hyperref[III.4.3.1-fr]{4.3.1}), basta pues probar el
\end{proof}

\begin{corollary}[4.3.2]
\label{III.4.3.2-fr}
Bajo las hipótesis de \textup{(\hyperref[III.4.3.1-fr]{4.3.1})}, supongamos además
que $f_*(\mathcal O_X)$ sea isomorfo a $\mathcal O_Y$. Entonces las fibras
$f^{-1}(y)$ de $f$ son conexas y no vacías para todo $y\in Y$.
\end{corollary}

\begin{proof}
La hipótesis de que $f_*(\mathcal O_X)$ es isomorfo a $\mathcal O_Y$
implica ya que $f$ es dominante, luego sobreyectivo puesto que $f$ es una
aplicación cerrada. Se puede reducir, como en
(\hyperref[III.4.2.1-fr]{4.2.1}), al caso en que $y$ es \emph{cerrado} en $Y$;
$f^{-1}(y)$, siendo un espacio noetheriano, tiene un número finito de componentes
conexas, y es el espacio subyacente del completado $\widehat X$ a lo largo de
$f^{-1}(y)$. Si $Z_i$ ($1\leq i\leq n$) son estas componentes conexas, está
claro que $\Gamma(\widehat X,\mathcal O_{\widehat X})$ es producto
directo de los anillos $\Gamma(Z_i,\mathcal O_{\widehat X})$, y ninguno de estos
últimos se reduce a $0$, puesto que la sección unidad es distinta de
$0$ en cada punto de $\widehat X$. Ahora bien, si se aplica
(\hyperref[III.4.1.5-fr]{4.1.5}) a $\mathcal F=\mathcal O_X$, cuyo
completado a lo largo de $f^{-1}(y)$ es $\mathcal O_{\widehat X}$, se ve que
$\Gamma(\widehat X,\mathcal O_{\widehat X})$ es isomorfo al separado
completado $\mathfrak m_y$-ádico $\widehat{\mathcal O}_y$ del anillo local
$\mathcal O_y$; es por tanto un \emph{anillo local} que no puede ser producto
directo de varios anillos no reducidos a $0$ (pues tendría
varios ideales maximales distintos). Se tiene pues $n=1$, lo que demuestra el
corolario.
\end{proof}

\begin{corollary}[4.3.3]
\label{III.4.3.3}
\label{III.4.3.3-fr}
Bajo las hipótesis de \sref{III.4.3.1}, para todo $y\in Y$ el conjunto
de componentes conexas de la fibra $f^{-1}(y)$ está en correspondencia
biyectiva con el conjunto finito de puntos de la fibra
$g^{-1}(y)$, donde $g:Y'\to Y$ es el morfismo estructural; en otros
términos, corresponde al conjunto de \emph{ideales maximales} de
$(f_*(\mathcal O_X))_y$.
\end{corollary}

\begin{proof}
Como $Y'$ es finito sobre $Y$, el espacio $g^{-1}(y)$ es finito y
discreto \sref[II]{II.6.1.7}. Como
$f^{-1}(y)=f'^{-1}(g^{-1}(y))$,
la afirmación se sigue de esta observación y de
\sref{III.4.3.1}.
\end{proof}

Esto da una interpretación notable del $Y$-preesquema $Y'$ definido en
\sref{III.4.3.1}. La factorización
$f=g\circ f'$ del morfismo propio $f$ es análoga a la factorización
obtenida por K.~Stein para aplicaciones holomorfas de espacios
analíticos, y lo llamaremos en adelante la \emph{factorización de Stein}
de $f$.

\begin{remark}[4.3.4]
\label{III.4.3.4}
\label{III.4.3.4-fr}
Sea $k$ una extensión del cuerpo $\boldsymbol{k}(y)$: si el preesquema $f^{-1}(y)\otimes_{\boldsymbol{k}(y)}k=X\times_Y\operatorname{Spec}(k)$
es conexo, entonces $f^{-1}(y)$, que es su imagen por un morfismo de proyección (\textbf{I}, 3.4.7). Diremos que, para un morfismo $f:X\to Y$ de preesquemas y un punto
$y\in Y$, la fibra $f^{-1}(y)$ es
\emph{geométricamente conexa} si, para \emph{toda} extensión $k$ de
$\boldsymbol{k}(y)$, el preesquema
$f^{-1}(y)\otimes_{\boldsymbol{k}(y)}k=X\times_Y\operatorname{Spec}(k)$ es conexo.

\oldpage[III]{132}

Bajo las hipótesis de \sref{III.4.3.2}, la conclusión puede
reforzarse: las fibras $f^{-1}(y)$ son en efecto \emph{geométricamente conexas}. Para verlo, observemos que para toda extensión $k$ de $\boldsymbol{k}(y)$ existen un anillo
local noetheriano $A$ y un homomorfismo local
$\varphi:\mathcal O_y\to A$
tal que $A$ es un $\mathcal O_y$-módulo plano y el cuerpo residual de
$A$ es $\boldsymbol{k}(y)$-isomorfo a $k$ \sref[0]{0.10.3.1}.
Sea $Y_1=\operatorname{Spec}(A)$, y sea $h:Y_1\to Y$ el morfismo local
correspondiente a $\varphi$, que lleva el único punto cerrado $y_1$ de
$Y_1$ a $y$ (\textbf{I}, 2.4.1); pongamos
$X_1=X\times_YY_1$, y $f_1=f\times1_{Y_1}$; $f_1$ es propio
\sref[II]{II.5.4.2}, y $f_1^{-1}(y_1)$ es un
$\boldsymbol{k}(y_1)$-preesquema isomorfo a $X\times_Y\operatorname{Spec}(k)$. Basta, por
tanto, demostrar que
$f_{1*}(\mathcal O_{X_1})=\mathcal O_{Y_1}$
de modo que \sref{III.4.3.2} pueda aplicarse a $f_1$. Ahora $g$ es un morfismo \emph{plano},
como resulta de (\textbf{I}, 2.4.2) y de
(\hyperref[III.1.4.15-fr]{1.4.15.5}); se tiene pues
$f_{1*}(\mathcal O_{X_1})
=h^*(f_*(\mathcal O_X))
=h^*(\mathcal O_Y)
=\mathcal O_{Y_1}$
en virtud de (\hyperref[III.1.4.15-fr]{1.4.15}) aplicado para $q=0$.

En el caso general (\hyperref[III.4.3.1-fr]{4.3.1}), el mismo razonamiento
muestra que se tiene (con las notaciones de
(\hyperref[III.4.3.1-fr]{4.3.1}))
$f_{1*}(\mathcal O_{X_1})=h^*(g_*(\mathcal O_{Y'}))$,
y que la factorización de Stein $f_1=g_1\circ f'_1$ de $f_1$
satisface
$g_1=g\times1_{Y_1}$
\sref[II]{II.1.5.2}, siendo el $Y_1$-esquema finito correspondiente
$Y'_1=Y'\times_YY_1$. Teniendo en cuenta la transitividad de las
fibras \sref[I]{I.3.6.4}, se sigue de
\sref{III.4.3.3} que el número de componentes conexas de
$f_1^{-1}(y_1)$ es, en virtud de (\hyperref[III.4.3.3-fr]{4.3.3}), igual al número de elementos de $g_1^{-1}(y_1)=g^{-1}(y)\otimes_{\boldsymbol{k}(y)}k$.
Si $k$ una extensión \emph{algebraicamente cerrada} de $\boldsymbol{k}(y)$, este número
es independiente de la extensión algebraicamente cerrada considerada e igual al
\emph{número geométrico de puntos de} $g^{-1}(y)$
\sref[I]{I.6.4.7}, o equivalentemente a la \emph{suma de los grados separables}
$[\boldsymbol{k}(y'_i):\boldsymbol{k}(y)]_s$
donde los $y'_i$ recorren el conjunto finito $g^{-1}(y)$. Este número
se llama \emph{número geométrico de componentes conexas} de
$f^{-1}(y)$. Los cuerpos $\boldsymbol{k}(y'_i)$ son precisamente los cuerpos
residuales del anillo semilocal $(f_*(\mathcal O_X))_y$.
\end{remark}

\begin{proposition}[4.3.5]
\label{III.4.3.5}
\label{III.4.3.5-fr}
Sean $X$ e $Y$ preesquemas íntegros localmente noetherianos, y sea
$f:X\to Y$ un morfismo propio dominante. Para todo $y\in Y$, el
número de componentes conexas de $f^{-1}(y)$ es a lo sumo el número
de ideales maximales de la clausura entera $\mathcal O'_y$ de
$\mathcal O_y$ en el cuerpo de funciones racionales $\mathrm R(X)$.
\end{proposition}

\begin{proof}
Para todo abierto $U$ de $Y$,
$\Gamma(U,f_*(\mathcal O_X))=\Gamma(f^{-1}(U),\mathcal O_X)$
es la intersección de los anillos locales $\mathcal O_x$ para
$x\in f^{-1}(U)$ \sref[I]{I.8.2.1.1}. Se sigue inmediatamente que
$(f_*(\mathcal O_X))_y$ es un subanillo de $\mathrm R(X)$ que
contiene a $\mathcal O_y$. Como $f_*(\mathcal O_X)$ es un
$\mathcal O_Y$-módulo coherente, $(f_*(\mathcal O_X))_y$ es un
$\mathcal O_y$-módulo finito y, por tanto, está contenido en $\mathcal O'_y$; se
sabe ([13], vol. I, pp. 257 y 259) que todo ideal maximal de tal anillo
$A$ es la intersección de $A$ y de un ideal maximal de $\mathcal O'_y$
\cite[vol.~I, pp.~257 and 259]{I-13}, lo que demuestra la proposición.
\end{proof}

\begin{definition}[4.3.6]
\label{III.4.3.6}
\label{III.4.3.6-fr}
Un anillo local íntegro se dice unirrama si su clausura entera
es de nuevo un anillo local. Un punto $y$ de un preesquema íntegro
$Y$ se dice unirrama si el anillo local $\mathcal O_y$ es
unirrama; este es, en particular, el caso cuando $Y$ es \emph{normal} en
$y$.
\end{definition}

Sea $A$ un anillo local íntegro y sea $K$ su cuerpo de fracciones. El
anillo $A$ es unirrama si y solo si todo subanillo $A_1$ de $K$ que
contenga a $A$ y sea una $A$-álgebra finita es local. En efecto, sea
$A'$ la clausura entera de $A$. Por el primer teorema de
Cohen--Seidenberg (Bourbaki, \emph{Alg\`ebre commutative}, capítulo~V,
\S2, n.º~1, teorema~1), todo ideal maximal de $A_1$ es la traza
de un ideal maximal de $A'$; así, si $A'$ es local, también
lo es $A_1$.

\oldpage[III]{133}

Recíprocamente, $A'$ es el límite inductivo de la familia dirigida de
$A$-subálgebras finitas $A_\alpha$ de $A'$. Si cada $A_\alpha$ es un anillo local, el ideal maximal
de $A_\alpha$ es la traza sobre $A_\alpha$ es la contracción del ideal maximal de $A_\beta$ para
$A_\alpha\subset A_\beta$ por el mismo razonamiento que antes, luego $A'$ es un anillo local (\textbf{0}, 10.3.1.3). Se observará que si el completado de un anillo local \emph{noetheriano} $A$ es íntegra ---propiedad que se expresa diciendo que
$A$ es \emph{analíticamente íntegro}---, entonces $A$ es \emph{unirrama}.
Sea $\mathfrak m$ el ideal maximal de $A$, sea $K$ su cuerpo de
fracciones, y sea $K'$ el cuerpo de fracciones de $\widehat A$;
entonces
$K'=K\otimes_A\widehat A$.
Sea $A'_F$ una $A$-subálgebra finita de $K$. El subanillo $B_F$ de
$K'$ generado por $\widehat A$ y $A'_F$ es isomorfo a
$A'_F\otimes_A\widehat A$
es un $\widehat A$-módulo finito y es la \emph{completación} de $A'_F$
para la topología $\mathfrak m$-ádica
\sref[0]{0.7.3.3}, \sref[0]{0.7.3.6}. Como $A'_F$ es semilocal
(Bourbaki, \emph{Alg\`ebre commutative}, capítulo~IV, \S2, n.º~5,
corolario~3 de la proposición~9), mientras que su completación es
íntegra, $A'_F$ solo puede tener un ideal maximal $\mathfrak m'_F$, y
$\mathfrak m'_F\cap A=\mathfrak m$; de donde nuestra aserción.
\begin{corollary}[4.3.7]
\label{III.4.3.7}
\label{III.4.3.7-fr}
Bajo las hipótesis de \sref{III.4.3.5}, supongamos que la clausura
algebraica de $\mathrm R(Y)$ en $\mathrm R(X)$ tiene grado separable
$n$ y que $y\in Y$ es unirrama. Entonces la fibra $f^{-1}(y)$ tiene a
lo sumo $n$ componentes conexas. En particular, si la clausura
algebraica de $\mathrm R(Y)$ en $\mathrm R(X)$ es radicial sobre
$\mathrm R(X)$, entonces $f^{-1}(y)$ es conexa.
\end{corollary}

\begin{proof}
Sea $\mathcal O''_y$ la clausura entera de $\mathcal O_y$, y sea
$\mathcal O'_y$ la clausura entera de $\mathcal O_y$ en $\mathrm R(X)$.
Esta última es también la clausura entera de $\mathcal O''_y$; pero se sabe que si $\mathcal O''_y$ es local, entonces $\mathcal O'_y$ es
un anillo semilocal que tiene a lo sumo $n$ ideales maximales
\cite[vol.~I, p.~289, Theorem~22]{I-13}.

Este corolario es esencialmente la forma en la que Zariski enuncia su
``teorema de conexidad'' para esquemas algebraicos.
\end{proof}

\begin{remark}[4.3.8]
\label{III.4.3.8}
\label{III.4.3.8-fr}
Si, además de las hipótesis de \sref{III.4.3.7}, $Y$ es \emph{normal} en
$y$, entonces la fibra $f^{-1}(y)$ es geométricamente conexa: con la
notación de \sref{III.4.3.4}, $g^{-1}(y)$ es un solo punto $y'$, y
$\boldsymbol{k}(y')$ es radicial sobre $\boldsymbol{k}(y)$.
\end{remark}

\begin{definition}[4.3.9]
\label{III.4.3.9-fr}
Dado un preesquema localmente noetheriano $Y$, se dice que un morfismo \textup{de
tipo finito} $f:X\to Y$ es universalmente abierto si, para todo
preesquema irreducible localmente noetheriano $Y'$ y todo morfismo dominante
$g:Y'\to Y$, toda componente irreducible de $X'=X\times_YY'$ domina a $Y'$.
\end{definition}

Si $Y$ es irreducible, esto equivale a decir que si $\eta$, $\eta'$ son los
puntos genéricos de $Y$ y $Y'$ respectivamente (de modo que
$g(\eta')=\eta$), y si se pone $f'=f_{(Y')}$, toda componente irreducible
de $X'$ corta a $f'^{-1}(\eta')$ (\textbf{0}\textsubscript{I}, 2.1.8); esto
implica pues que para todo abierto $U$ de $Y$ el morfismo
$f^{-1}(U)\to U$, restricción de $f$, es universalmente abierto,

\begin{corollary}[4.3.10]
\label{III.4.3.10}
\label{III.4.3.10-fr}
Sean $X$ e $Y$ preesquemas íntegros localmente noetherianos, y sea
$f:X\to Y$ propio, dominante y universalmente abierto. Si la clausura
algebraica de $\mathrm R(Y)$ en $\mathrm R(X)$ es radicial sobre
$\mathrm R(Y)$, entonces toda fibra $f^{-1}(y)$ ($y\in Y$) es geométricamente conexa.
\end{corollary}

\begin{proof}
Se puede limitar al caso en que $Y=\operatorname{Spec}(B)$, donde $B$ es un anillo
íntegro noetheriano. Por \sref[II]{II.7.1.7}, existe un anillo local
\emph{noetheriano integralmente cerrado} $A$ que domina a $\mathcal O_y$ y tiene a
$\mathrm R(Y)$ por cuerpo de fracciones. Sea $Y'=\operatorname{Spec}(A)$, y
sea $h:Y'\to Y$ el morfismo correspondiente a la inyección canónica
$B\to A$. Es birracional, luego dominante; además, si $y'$ es el
único punto cerrado de $Y'$, entonces $h(y')=y$. Sean
\oldpage[III]{134}
$X'=X\times_YY'$, $f'=f\times1_{Y'}$; designemos por $\eta$, $\eta'$ y $\xi$ los puntos genéricos de $Y$, $Y'$ y
$X$, respectivamente. Entonces $f(\xi)=\eta$ y
$h^{-1}(\eta)=\{\eta'\}$; además,
$\boldsymbol{k}(\eta)=\boldsymbol{k}(\eta')=\mathrm R(Y)$. Así, $f'^{-1}(\eta')$ es
isomorfo a $f^{-1}(\eta)$ \sref[I]{I.3.6.4}, y tiene un único punto
genérico, pues $\xi$ es el punto genérico de
$f^{-1}(\eta)$ (\textbf{0}\textsubscript{I}, 2.1.8),
$f'^{-1}(\eta')$ tiene un solo punto genérico. Pero por hipótesis, toda
componente irreducible de $X'$ tiene su punto genérico en
$f'^{-1}(\eta')$. Por tanto, $X'$ es irreducible; su punto genérico
$\xi'$ es el de $f'^{-1}(\eta')$, y
$\boldsymbol{k}(\xi')=\boldsymbol{k}(\xi)$.



Pongamos $X''=X'_{\mathrm{red}}$ y $f''=f'_{\mathrm{red}}$. Entonces
$X''$ es íntegro y noetheriano, $f''$ es propio
\sref[II]{II.5.4.6}, y los espacios subyacentes de las fibras
$f'^{-1}(y')$ y $f''^{-1}(y')$ son iguales. Además,
$\mathrm R(X'')=\boldsymbol{k}(\xi')=\mathrm R(X)$,
Así, $f''$ satisface las hipótesis de \sref{III.4.3.8}, y
$f''^{-1}(y')$ es geométricamente conexa.

Sea ahora $k$ una extensión cualquiera de $\boldsymbol{k}(y)$. Existe una
extensión $k_1$ de $\boldsymbol{k}(y)$ en la que tanto $\boldsymbol{k}(y')$ como $k$ pueden considerarse como subextensiones (Bourbaki, \emph{Alg.}, cap. V, \S~4, prop. 2). Por hipótesis, $f''^{-1}(y')\times_{Y'}\operatorname{Spec}(k_1)$ es conexo, y tiene el mismo preesquema reducido que $f'^{-1}(y')\times_{Y'}\operatorname{Spec}(k_1)$ (\textbf{I}, 5.1.8), luego este último es conexo, y como es isomorfo a $f^{-1}(y)\times_Y\operatorname{Spec}(k_1)$ (\textbf{I}, 3.6.4), se concluye que este último es conexo; \emph{a fortiori}, lo mismo ocurre con $f^{-1}(y)\times_Y\operatorname{Spec}(k)$ según la observación del comienzo de (\hyperref[III.4.3.4-fr]{4.3.4}), lo que termina la demostración.
\end{proof}

\begin{env}[4.3.11]
\label{III.4.3.11-fr}
\emph{Observaciones} (4.3.11). ---
\begin{enumerate}
\item[(i)] El razonamiento anterior se debe en sustancia a Zariski [20],
salvo que puede tomar como $A$ el cierre integral de
$\mathcal O_y$, que es un anillo noetheriano para los anillos locales
de la geometría algebraica clásica. Por otra parte, Zariski prueba que si
$Y$ es la variedad de Chow de un espacio proyectivo $\mathbf P_k^r$ sobre un
cuerpo $k$, y si $X$ es la parte cerrada de $\mathbf P_k^r\times_kY$ que
define la correspondencia de Chow entre $\mathbf P_k^r$ y $Y$, entonces la
proyección $X\to Y$ es un morfismo universalmente abierto
(\emph{loc. cit.}, lema de la p. 82). Parece que esta es la única
propiedad formal de las «coordenadas de Chow» que se ha utilizado en
ciertas aplicaciones; interesa por tanto, en tal situación,
sustituir el lenguaje: fibras de un morfismo propio (eventualmente
supuesto universalmente abierto o sometido a otras restricciones análogas
de regularidad local) al lenguaje: especialización de ciclos en el espacio
proyectivo.

\item[(ii)] En el cap. IV veremos que un morfismo universalmente abierto
$f:X\to Y$ puede definirse también de la manera siguiente (que justifica la
terminología): para \emph{todo} morfismo $Z\to Y$, el morfismo
$f_{(Z)}:X_{(Z)}\to Z$ es \emph{abierto}. Se puede mostrar además que si $f$
verifica las hipótesis de (\hyperref[III.4.3.10-fr]{4.3.10}), entonces, si
$y,y'$ son dos puntos de $Y$ tales que $y$ es una \emph{especialización} de $y'$,
el número geométrico de componentes conexas de $f^{-1}(y)$ es a lo sumo
igual al de las componentes conexas de $f^{-1}(y')$.
\end{enumerate}
\end{env}

\begin{corollary}[4.3.12]
\label{III.4.3.12}
\label{III.4.3.12-fr}
Bajo las hipótesis de \sref{III.4.3.5}, supongamos además que
$\mathrm R(Y)$ es algebraicamente cerrado en $\mathrm R(X)$, y sea
$y$ un punto normal de $Y$. Entonces $f^{-1}(y)$ es geométricamente
conexa, y existe una vecindad abierta $U$ de $y$ en $Y$ tal que
$f_*(\mathcal O_X|f^{-1}(U))$ es isomorfo a $\mathcal O_Y|U$. Si, más
particularmente, se supone $Y$ normal (y $\mathrm R(Y)$ es algebraicamente
cerrado en $\mathrm R(X)$) $f_*(\mathcal O_X)$ es isomorfo a
$\mathcal O_Y$.
\end{corollary}

\begin{proof}
La primera afirmación relativa a $f^{-1}(y)$ es un caso particular de (\hyperref[III.4.3.8-fr]{4.3.8}). Se
\oldpage[III]{135}
deduce que si $f:X\xrightarrow{f'}Y'\xrightarrow{g}Y$
es la factorización de Stein de $f$ \sref{III.4.3.3}, se sigue que
$g^{-1}(y)$ consta de un solo punto $y'$. Además,
$\mathcal O_y\subset\mathcal O_{y'}=(f_*(\mathcal O_X))_y
\subset\mathrm R(X)$,
Como $\mathcal O_{y'}$ es finito sobre $\mathcal O_y$ (y \emph{a fortiori} sobre $\mathrm R(Y)$) está contenido en $\mathrm R(Y)$ en virtud de la hipótesis. Como $y$ es normal, necesariamente
$\mathcal O_{y'}=\mathcal O_y$. Así, $g$ es un isomorfismo local en $y'$
\sref[I]{I.6.5.4}, lo que demuestra la primera parte. La segunda
afirmación se sigue de la primera: la hipótesis adicional hace que $g$
sea biyectivo y un isomorfismo local en una vecindad de cada punto de
$Y'$, luego un isomorfismo.
\end{proof}

El hecho de que \sref{III.4.3.7} esté disponible para esquemas da, por
ejemplo, la aplicación siguiente.

\begin{proposition}[4.3.13]
\label{III.4.3.13-fr}
Sean $A$ un anillo local noetheriano unirramificado, $\mathfrak a$ un ideal de
definición de $A$, $A_0=A/\mathfrak a$,
$S=\operatorname{gr}_{\mathfrak a}(A)$ el anillo graduado asociado a $A$ para
la filtración $\mathfrak a$-preádica; $S$ es una $A_0$-álgebra graduada
generada por $S_1$, siendo $S_1$ un $A_0$-módulo de tipo finito. Entonces
$\operatorname{Proj}(S)$ es un $A_0$-esquema conexo.
\end{proposition}

\begin{proof}
Sea $\mathfrak m$ el ideal maximal de $A$; $Y=\operatorname{Spec}(A)$ es
un esquema íntegro cuyo punto $y$ correspondiente a $\mathfrak m$ es
el único punto cerrado. Por hipótesis se tiene
$\mathfrak m^p\subset\mathfrak a\subset\mathfrak m$ para un entero $p$,
luego $V(\mathfrak a)=\{\mathfrak m\}$. Sea
$S'=\bigoplus_{n\geq0}\mathfrak a^n$, y sea
$X=\operatorname{Proj}(S')$, que es el $Y$-esquema obtenido haciendo explotar
el ideal $\mathfrak a$; $X$ es íntegro y el morfismo estructural
$f:X\to Y$ es birracional (\textbf{II}, 8.1.4) y evidentemente proyectivo. Por tanto,
(\hyperref[III.4.3.7-fr]{4.3.7}) es aplicable y muestra que $f^{-1}(y)$
es conexo; pero el espacio $f^{-1}(y)$ es subyacente a
$\operatorname{Proj}(S'\otimes_AA_0)$ (\textbf{I}, 3.6.1 y \textbf{II}, 2.8.10); como
$S'\otimes_AA_0=S$ por definición, la proposición queda demostrada.
\end{proof}

\subsection{El «main theorem» de Zariski.}
\label{subsection:III.4.4-fr}

\begin{proposition}[4.4.1]
\label{III.4.4.1}
\label{III.4.4.1-fr}
Sea $Y$ un preesquema localmente noetheriano y sea $f:X\to Y$ un
morfismo propio. Sea $X'$ el conjunto de puntos $x\in X$ que están
aislados en su fibra $f^{-1}(f(x))$. Entonces $X'$ es abierto en $X$.
Si
$f=g\circ f'$
es la factorización de Stein de $f$ \sref{III.4.3.3}, la restricción
de $f'$ a $X'$ es un isomorfismo de $X'$ sobre el subpreesquema inducido
en un subconjunto abierto $U$ de $Y'$, y
$X'=f'^{-1}(U)$.
\end{proposition}

\begin{proof}
Como $g^{-1}(f(x))$ es finito y discreto
\sref{III.4.3.3}, \sref[II]{II.6.1.7}, el punto $x$ está aislado en
$f^{-1}(f(x))$ si y solo si está aislado en
$f'^{-1}(f'(x))$; se puede pues limitar al caso en que $f'=f$, de
modo que $f_*(\mathcal O_X)=\mathcal O_Y$. Si $x\in X'$, la fibra
$f^{-1}(f(x))$, que es conexa por \sref{III.4.3.2}, debe constar
entonces del único punto $x$.

Como $f$ es cerrado, para toda vecindad abierta $V$ de $x$ en $X$, el
conjunto cerrado $f(X\setmin V)$ es cerrado en $Y$ y no contiene
$y=f(x)$, puesto que $f^{-1}(y)=\{x\}$; si $U$ es el complementario de
$f(X\setmin V)$ en $Y$, entonces $f^{-1}(U)\subset V$. Así, las
imágenes inversas por $f$ de un sistema fundamental de vecindades
abiertas de $y$ forman un sistema fundamental de vecindades abiertas de
$x$. La igualdad $f_*(\mathcal O_X)=\mathcal O_Y$ y la definición de la
imagen directa de un haz \sref[0]{0.3.4.1}, \sref{III.4.2.1} implican
ahora que, si $f=(\psi,\theta)$, el homomorfismo
$\theta_y^\sharp:\mathcal O_y\to\mathcal O_x$
es un \emph{isomorfismo}. Por tanto, existen vecindades abiertas $V$ de $x$ y
$U$ de $y$ tales que la restricción de $f$ a $V$ es un isomorfismo de $V$ sobre
$U$ (\textbf{I}, 6.5.4); además, por lo que se acaba de ver,
\oldpage[III]{136}
se puede suponer que $f^{-1}(U)=V$. Se sigue inmediatamente que $V\subset X'$, lo que
termina la demostración.
\end{proof}

La proposición siguiente fue demostrada por Chevalley para esquemas
algebraicos.

\begin{proposition}[4.4.2]
\label{III.4.4.2}
\label{III.4.4.2-fr}
Sea $Y$ un preesquema localmente noetheriano y sea $f:X\to Y$ un
morfismo. Las condiciones siguientes son equivalentes:
\begin{enumerate}
  \item[{\rm a)}] $f$ es finito;
  \item[{\rm b)}] $f$ es afín y propio;
  \item[{\rm c)}] $f$ es propio y, para todo $y\in Y$, el conjunto
    $f^{-1}(y)$ es finito.
\end{enumerate}
\end{proposition}

\begin{proof}
La condición~\emph{a)} implica \emph{b)} por \sref[II]{II.6.1.2} y
\sref[II]{II.6.1.11}. Si $f$ es propio y afín, también lo es
$f^{-1}(y)\to\operatorname{Spec}(\boldsymbol{k}(y))$ (\textbf{II}, 1.6.2, (iii) y 5.4.2, (iii)), y el teorema de finitud (\hyperref[III.3.2.1-fr]{3.2.1}) aplicado al haz estructural de $f^{-1}(y)$, muestra
que $f^{-1}(y)=\operatorname{Spec}(A)$, donde
$A$ es una $\boldsymbol{k}(y)$-álgebra \emph{finita}; luego $f^{-1}(y)$ es un conjunto finito (\textbf{II}, 6.1.7), y se ve que \emph{b)} implica \emph{c)}. Finalmente, como $f^{-1}(y)$ es un preesquema
algebraico sobre $\boldsymbol{k}(y)$, la hipótesis de que el conjunto $f^{-1}(y)$ es
finito implica que el espacio $f^{-1}(y)$ es discreto (\textbf{I}, 6.4.4). Con las
notaciones de (\hyperref[III.4.4.1-fr]{4.4.1}), se tiene pues $X'=X$, luego
$f':X\to Y'$ es un isomorfismo. Como $g$ es finito, \emph{c)} implica \emph{a)}.
\end{proof}

\begin{theorem}[4.4.3]
\label{III.4.4.3-fr}
\textnormal{(«Teorema principal» de Zariski).}
Sean $Y$ un preesquema noetheriano, $f:X\to Y$ un morfismo cuasiproyectivo,
$X'$ el conjunto de los puntos $x\in X$ que están aislados en su fibra
$f^{-1}(f(x))$. Entonces $X'$ es una parte abierta de $X$, y el
subpreesquema inducido $X'$ es isomorfo a un preesquema inducido sobre una parte
abierta de un $Y$-preesquema $Y'$ finito sobre $Y$.
\end{theorem}

\begin{proof}
La hipótesis implica que existe un $Y$-preesquema proyectivo $Z$ tal que $X$
sea $Y$-isomorfo a un subpreesquema inducido sobre un abierto de $Z$
(\textbf{II}, 5.3.2 y 5.5.1). Se reduce pues a demostrar el teorema cuando $f$
es un morfismo proyectivo, luego propio (\textbf{II}, 5.5.3), y resulta entonces
inmediatamente de (\hyperref[III.4.4.1-fr]{4.4.1}).
\end{proof}

\begin{remark}[4.4.4]
\label{III.4.4.4-fr}
Si $X$ es reducido (resp. irreducible y $X'$ no
es vacío), puede exigirse en \sref{III.4.4.3} que $Y'$ es reducido (resp. irreducible). En efecto, se puede siempre reemplazar $Y'$ por el subpreesquema \emph{adherencia} $\overline{X'}$ de $X'$ en $Y'$ (\textbf{I}, 9.5.11 y \textbf{II}, 6.1.5, (i) y (ii)), y se sabe que si $X'$ es reducido, también
lo es $\overline{X}'$ \sref[I]{I.9.5.9}; si $X'$ no es vacío, es
irreducible si $X$ lo es, y $\overline{X'}$ también es irreducible.
\end{remark}

\begin{corollary}[4.4.5]
\label{III.4.4.5}
\label{III.4.4.5-fr}
Sea $Y$ un esquema localmente noetheriano, sea $f:X\to Y$ un
morfismo de tipo finito, y sea $x$ un punto de $X$ aislado en su fibra
$f^{-1}(f(x))$. Entonces existe una vecindad abierta de $x$ en $X$
isomorfa a un subconjunto abierto de un $Y$-preesquema finito sobre $Y$.
\end{corollary}

\begin{proof}
Sean en efecto $y=f(x)$, sea $U$ una vecindad abierta afín de $y$ en $Y$, y
sea $V$ una vecindad abierta afín de $x$ en $X$, contenida en
$f^{-1}(U)$. Como $Y$ es separado, la inclusión $U\to Y$ es afín
\sref[II]{II.1.6.3}; y como $V$ es afín sobre $U$ (\emph{ibid.}), la
restricción de $f$ a $V$ es un morfismo afín $V\to Y$
\sref[II]{II.1.6.2}. A fortiori, esta restricción es
cuasiproyectiva, pues es de tipo finito
\sref[I]{I.6.3.5}, \sref[II]{II.5.3.4}. Solo queda aplicar
\sref{III.4.4.3} a esta restricción.
\end{proof}

El Corolario \sref{III.4.4.5} puede enunciarse en el lenguaje del
álgebra conmutativa como sigue.

\oldpage[III]{137}

\begin{corollary}[4.4.6]
\label{III.4.4.6}
\label{III.4.4.6-fr}
Sea $A$ un anillo noetheriano, sea $B$ una $A$-álgebra de tipo
finito, sea $\mathfrak q$ un ideal primo de $B$, y sea $\mathfrak p$
su imagen inversa en $A$. Se supone que $\mathfrak q$ es a la vez
maximal y minimal en el conjunto de ideales primos de $B$ cuya imagen
inversa es $\mathfrak p$. Entonces existen un elemento
$g\in B\setmin\mathfrak q$, una $A$-álgebra finita $A'$, y un elemento
$f'\in A'$ tales que las $A$-álgebras $B_g$ y $A'_{f'}$ sean isomorfos.
\end{corollary}

\begin{proof}
Basta en efecto aplicar (\hyperref[III.4.4.5-fr]{4.4.5}) a $Y=\operatorname{Spec}(A)$ y $X=\operatorname{Spec}(B)$. La
hipótesis sobre $\mathfrak q$ significa exactamente que está aislado en su fibra
(\textbf{I}, 1.1.7).
\end{proof}

Se deduce el resultado siguiente, menos general en apariencia:

\begin{corollary}[4.4.7]
\label{III.4.4.7}
\label{III.4.4.7-fr}
Sea $A$ un anillo local noetheriano, sea $B$ una $A$-álgebra de tipo
finito, y sea $\mathfrak n$ un ideal primo de $B$ cuya imagen inversa
en $A$ es el ideal maximal $\mathfrak m$. Se supone que $\mathfrak n$
es maximal en $B$ y minimal en el conjunto de ideales primos de $B$
cuya imagen inversa es $\mathfrak m$ (lo que significa también que
$B\mathfrak m$ es primario para $\mathfrak n$-primario). Entonces existen una
$A$-álgebra finita $A'$ y un ideal maximal $\mathfrak m'$ de $A'$ (del cual
$\mathfrak m$ es la imagen inversa en $A$) tales que $B_{\mathfrak n}$ es isomorfo, como $A$-álgebra, a $A'_{\mathfrak m'}$.
\end{corollary}

El caso particular siguiente de \sref{III.4.4.7} se llama también a
veces el ``teorema principal''.

\begin{corollary}[4.4.8]
\label{III.4.4.8}
\label{III.4.4.8-fr}
Bajo las hipótesis de \sref{III.4.4.7}, supongamos además que $A$ y
$B$ son dominios íntegros con el mismo cuerpo de fracciones $K$.
Si $A$ es
integralmente cerrado, se tiene $B=A$.
\end{corollary}

\begin{proof}
En efecto, la \emph{Observación} (\hyperref[III.4.4.4-fr]{4.4.4}) muestra que se puede
suponer, en la aplicación de (\hyperref[III.4.4.7-fr]{4.4.7}), que $A'$ es un dominio íntegro con
cuerpo de fracciones $K$. La hipótesis sobre $A$ da entonces
$A'=A$, luego $B_{\mathfrak n}=A$. Como
$A\subset B\subset B_{\mathfrak n}$, se sigue que $A=B$.
\end{proof}

El enunciado \sref{III.4.4.8} es la forma en la que Zariski dio su
``teorema principal'' (extendido aquí a anillos locales íntegros
noetherianos arbitrarios).

Los corolarios anteriores eran variantes locales de
\sref{III.4.4.3}, que es un resultado global. He aquí otra
consecuencia de naturaleza global.

\begin{corollary}[4.4.9]
\label{III.4.4.9}
\label{III.4.4.9-fr}
Sea $Y$ un preesquema íntegro localmente noetheriano, y sea
$f:X\to Y$ un morfismo separado, de tipo finito y birracional.
Supongamos además que $Y$ es normal y que toda fibra $f^{-1}(y)$,
para $y\in Y$, es finita. Entonces $f$ es una inmersión abierta. Si,
además, $f$ es cerrado (en particular, si $f$ es propio), entonces
$f$ es un isomorfismo.
\end{corollary}

\begin{proof}
Sea $x\in X$ y pongamos $y=f(x)$. Como $f^{-1}(y)$ es un esquema
algebraico sobre $\boldsymbol{k}(y)$, la hipótesis de que es finito implica que
es discreto \sref[I]{I.6.4.4}. Además, $\mathcal O_y$ es integralmente
cerrado, y $\mathcal O_x$ y $\mathcal O_y$ tienen el mismo cuerpo de fracciones (\textbf{I}, 7.1.5). Se puede pues aplicar (\hyperref[III.4.4.8-fr]{4.4.8}), y si $f=(\psi,\theta)$, el homomorfismo
$\theta_y^\sharp:\mathcal O_y\to\mathcal O_x$
es biyectivo. De \sref[I]{I.6.5.4} se sigue que $f$ es un isomorfismo
local. Como $f$ es separado y $X$ es íntegro, $f$ es una inmersión
abierta \sref[I]{I.8.2.8}. La última afirmación se sigue del hecho de
que $f$ es dominante.
\end{proof}

\begin{proposition}[4.4.10]
\label{III.4.4.10}
\label{III.4.4.10-fr}
Sea $Y$ un preesquema localmente noetheriano, y sea $f:X\to Y$ un
morfismo localmente de tipo finito. El conjunto $X'$ de puntos
$x\in X$ que están aislados en su fibra $f^{-1}(f(x))$ es abierto en
$X$.
\end{proposition}

\begin{proof}
La cuestión es local sobre $X$ e $Y$, por lo que se puede suponer que
$X$ e $Y$ son afines noetherianos y que $f$ es de tipo finito. El
morfismo $f$ es entonces afín y de tipo finito, luego
cuasiproyectivo \sref[II]{II.5.3.4}; basta aplicar
\sref{III.4.4.3}.
\end{proof}

\oldpage[III]{138}

\begin{corollary}[4.4.11]
\label{III.4.4.11}
\label{III.4.4.11-fr}
Sea $Y$ un preesquema localmente noetheriano, y sea $f:X\to Y$ un
morfismo propio. El conjunto $U$ de puntos $y\in Y$ para los que
$f^{-1}(y)$ es discreto es abierto en $Y$, y el morfismo restringido
$f^{-1}(U)\to U$, restricción de $f$, es
finito. En particular, un morfismo propio y cuasifinito $X\to Y$ es finito.
\end{corollary}

\begin{proof}
El complementario de $U$ en $Y$ es la imagen por $f$ de $X\setmin X'$, que
es cerrado en $X$ por \sref{III.4.4.10}; como $f$ es una aplicación
cerrada, $U$ es abierto. Además, \sref[II]{II.6.2.2} muestra que
$f^{-1}(y)$ es finita para todo $y\in U$. Como el morfismo restringido
$f^{-1}(U)\to U$, restricción de $f$, es propio (\textbf{II}, 5.4.1), es finito en
virtud de (\hyperref[III.4.4.2-fr]{4.4.2}).
\end{proof}

\begin{env}[4.4.12]
\label{III.4.4.12-fr}
\emph{Observaciones} (4.4.12). ---
\begin{enumerate}
\item[(i)] Como se ha anunciado en (\textbf{II}, 6.2.7), mostraremos en el cap. V
que si $Y$ es localmente noetheriano, todo morfismo cuasifinito y separado
$f:X\to Y$ es cuasiafín, luego cuasiproyectivo. Resultará que, en el
Main Theorem (\hyperref[III.4.4.3-fr]{4.4.3}), la conclusión sigue siendo válida
cuando se supone solamente $f$ \emph{separado y de tipo finito}. En efecto, resulta
de (\hyperref[III.4.4.10-fr]{4.4.10}) que $X'$ es abierto en $X$, y como
$X$ es localmente noetheriano, la restricción de $f$ a $X'$ es también de
tipo finito (\textbf{I}, 6.3.5), luego cuasifinita por definición de $X'$, y evidentemente
separada; se puede pues aplicar (\hyperref[III.4.4.3-fr]{4.4.3}) a esta
restricción, de donde la conclusión.

\item[(ii)] Daremos en el cap. IV una demostración más elemental de
(\hyperref[III.4.4.10-fr]{4.4.10}), utilizando la teoría de la dimensión.
\end{enumerate}
\end{env}

\subsection{Completados de módulos de homomorfismos}
\label{subsection:III.4.5-fr}

\begin{proposition}[4.5.1]
\label{III.4.5.1-fr}
Sean $A$ un anillo noetheriano, $\mathfrak J$ un ideal de $A$, $X$ un
$A$-preesquema de tipo finito, $\mathcal F$, $\mathcal G$ dos
$\mathcal O_X$-Módulos coherentes cuyos soportes tienen una intersección
propia sobre $Y=\operatorname{Spec}(A)$ \textup{(\textbf{II}, 5.4.10)}. Entonces, para todo entero
$n\geq0$, $\operatorname{Ext}^n_{\mathcal O_X}(X;\mathcal F,\mathcal G)$ es un $A$-módulo
de tipo finito, y su separado completado para la topología $\mathfrak J$-preádica
se identifica canónicamente (con las notaciones de
\textup{(\hyperref[III.4.1.7-fr]{4.1.7})}) con
$\operatorname{Ext}^n_{\mathcal O_{\widehat X}}(\widehat X;\widehat{\mathcal F},
\widehat{\mathcal G})$.
\end{proposition}

\begin{proof}
Se sabe (T, 4.2) que existe una sucesión espectral birregular
$\mathrm E(\mathcal F,\mathcal G)$ cuyo resultado es
$\operatorname{Ext}_{\mathcal O_X}(X;\mathcal F,\mathcal G)$ y cuyos términos
$\mathrm E_2$ vienen dados por
$\mathrm E_2^{pq}=\mathrm H^p\bigl(X,
\mathcal{E}\!xt^q_{\mathcal O_X}(\mathcal F,\mathcal G)\bigr)$.
Se sabe que $\mathcal{E}\!xt^q_{\mathcal O_X}(\mathcal F,\mathcal G)$ es un
$\mathcal O_X$-Módulo coherente (\textbf{0}, 12.3.3) cuyo soporte está contenido en
la intersección de los de $\mathcal F$ y $\mathcal G$ (T, 4.2.2), y es
por tanto propio sobre $Y$ (\textbf{II}, 5.4.10). Se concluye de
(\hyperref[III.3.2.4-fr]{3.2.4}) que los $\mathrm E_2^{pq}$ son
$A$-módulos de tipo finito, y por tanto (\textbf{0}, 11.1.8) lo mismo ocurre con todos
los términos $\mathrm E_r^{pq}$ de la sucesión espectral y con su resultado.
Por otra parte, si $i:\widehat X\to X$ es el morfismo canónico,
$\widehat{\mathcal F}$ y $\widehat{\mathcal G}$ se identifican canónicamente
con $i^*(\mathcal F)$ y $i^*(\mathcal G)$, y $i$ es \emph{plano} (\textbf{I}, 10.8.8 y
10.8.9). Se sabe entonces (\textbf{0}, 12.3.4) que existe, para todo $q\geq0$, un
$i$-morfismo canónico
$u_q:\mathcal{E}\!xt^q_{\mathcal O_X}(\mathcal F,\mathcal G)
\longrightarrow
\mathcal{E}\!xt^q_{\mathcal O_{\widehat X}}
(\widehat{\mathcal F},\widehat{\mathcal G})$,
y que el $\mathcal O_{\widehat X}$-homomorfismo correspondiente
$u^\sharp:i^*\bigl(\mathcal{E}\!xt^q_{\mathcal O_X}(\mathcal F,\mathcal G)\bigr)
\longrightarrow
\mathcal{E}\!xt^q_{\mathcal O_{\widehat X}}
(\widehat{\mathcal F},\widehat{\mathcal G})$
es un \emph{isomorfismo} (\textbf{0}, 12.3.5); en otros términos (\textbf{I}, 10.8.8),
$\mathcal{E}\!xt^q_{\mathcal O_{\widehat X}}
(\widehat{\mathcal F},\widehat{\mathcal G})$ se identifica canónicamente con el
completado $(\mathcal{E}\!xt^q_{\mathcal O_X}(\mathcal F,\mathcal G))^{\widehat{\,}}$ (con las
notaciones de (\hyperref[III.4.1.7-fr]{4.1.7})). Se concluye entonces del teorema de
comparación (\hyperref[III.4.1.10-fr]{4.1.10}) que para todo $p\geq0$,
$\mathrm H^p\bigl(\widehat X,
\mathcal{E}\!xt^q_{\mathcal O_{\widehat X}}
(\widehat{\mathcal F},\widehat{\mathcal G})\bigr)$
se identifica canónicamente con el separado completado
\oldpage[III]{139}
de
$\mathrm H^p\bigl(X,\mathcal{E}\!xt^q_{\mathcal O_X}
(\mathcal F,\mathcal G)\bigr)$ para la topología $\mathfrak J$-preádica.
Si se designa por $\mathrm E(\widehat{\mathcal F},
\widehat{\mathcal G})$ la sucesión espectral birregular definida en (T, 4.2)
relativa a $\widehat{\mathcal F}$ y $\widehat{\mathcal G}$, se ve pues que
si $\widehat A$ designa el separado completado de $A$ para la topología
$\mathfrak J$-preádica, se tiene, salvo isomorfismo canónico,
$\mathrm E_2^{pq}(\widehat{\mathcal F},\widehat{\mathcal G})
=\mathrm E_2^{pq}(\mathcal F,\mathcal G)\otimes_A\widehat A$
(\textbf{0}\textsubscript{I}, 7.3.3).

Dicho esto, se sabe que el dato del morfismo plano $i$ define un
homomorfismo canónico de sucesiones espectrales
\[
\varphi:\mathrm E(\mathcal F,\mathcal G)\longrightarrow
\mathrm E(\widehat{\mathcal F},\widehat{\mathcal G})
=\mathrm E(i^*(\mathcal F),i^*(\mathcal G))
\]
que, para los términos $\mathrm E_2$ (resp. el resultado), se reduce al
homomorfismo
\[
\varphi_2^{pq}:\mathrm H^p\bigl(X,
\mathcal{E}\!xt^q_{\mathcal O_X}(\mathcal F,\mathcal G)\bigr)
\longrightarrow
\mathrm H^p\bigl(\widehat X,
\mathcal{E}\!xt^q_{\mathcal O_{\widehat X}}
(\widehat{\mathcal F},\widehat{\mathcal G})\bigr)
\]
\[
\text{(resp. }\varphi^n:\operatorname{Ext}^n_{\mathcal O_X}
(X;\mathcal F,\mathcal G)\longrightarrow
\operatorname{Ext}^n_{\mathcal O_{\widehat X}}
(\widehat X;\widehat{\mathcal F},\widehat{\mathcal G})\text{)}
\]
deducido de $u_q$ (resp. $u_0$) por funtorialidad (\textbf{0}, 12.3.4). Por
tensorización con $\widehat A$, los $\varphi_r^{pq}$ y $\varphi^n$ dan
homomorfismos de $\widehat A$-módulos
\[
\psi_r^{pq}:\mathrm E_r^{pq}(\mathcal F,\mathcal G)
\otimes_A\widehat A\longrightarrow
\mathrm E_r^{pq}(\widehat{\mathcal F},\widehat{\mathcal G})
\]
\[
\psi^n:\operatorname{Ext}^n_{\mathcal O_X}(X;\mathcal F,\mathcal G)
\otimes_A\widehat A\longrightarrow
\operatorname{Ext}^n_{\mathcal O_{\widehat X}}
(\widehat X;\widehat{\mathcal F},\widehat{\mathcal G}).
\]

Como $\widehat A$ es un $A$-módulo \emph{plano} (\textbf{0}\textsubscript{I}, 7.3.3), los
$\widehat A$-módulos
$\mathrm E_r^{pq}(\mathcal F,\mathcal G)\otimes_A\widehat A$ forman una
sucesión espectral birregular cuyo resultado son los
$\operatorname{Ext}^n_{\mathcal O_X}(X;\mathcal F,\mathcal G)\otimes_A\widehat A$, y los
$\psi_r^{pq}$ y $\psi_n$ un morfismo de sucesiones espectrales. Como los
$\psi_2^{pq}$ son \emph{isomorfismos}, lo mismo ocurre con los $\psi^n$
(\textbf{0}, 11.1.5).
\end{proof}

\begin{corollary}[4.5.2]
\label{III.4.5.2}
\label{III.4.5.2-fr}
Bajo las hipótesis de \sref{III.4.5.1}, supongamos además que $A$ es
un anillo noetheriano $\mathfrak J$-ádico. Entonces, para todo entero
$n\geq0$, $\operatorname{Ext}^n_{\mathcal O_X}(X;\mathcal F,\mathcal G)$
se identifica canónicamente con
$\operatorname{Ext}^n_{\mathcal O_{\widehat X}}
(\widehat X;\widehat{\mathcal F},\widehat{\mathcal G})$.
\end{corollary}

\begin{proof}
Basta observar que $\operatorname{Ext}^n_{\mathcal O_X}
(X;\mathcal F,\mathcal G)$, al ser un $A$-módulo de tipo finito, es separado y
completo para la topología $\mathfrak J$-preádica \sref[0]{0.7.3.6}.
\end{proof}

El caso particular $n=0$ de \sref{III.4.5.1} puede enunciarse como
sigue.

\begin{corollary}[4.5.3]
\label{III.4.5.3}
\label{III.4.5.3-fr}
Bajo las hipótesis de \sref{III.4.5.1}, para todo homomorfismo
$u:\mathcal F\to\mathcal G$, designemos por
$\widehat u$ el homomorfismo
completado $\widehat{\mathcal F}\to\widehat{\mathcal G}$ \textup{(\textbf{I}, 10.8.4)}. Entonces se tiene un isomorfismo canónico
\[
\label{III.4.5.3.1-fr}
\left(\operatorname{Hom}_{\mathcal O_X}(\mathcal F,\mathcal G)\right)^{\widehat{\,}}
\xrightarrow{\sim}
\operatorname{Hom}_{\mathcal O_{\widehat X}}
(\widehat{\mathcal F},\widehat{\mathcal G})
\tag{4.5.3.1}
\]
donde el miembro izquierdo es la completación separada, para la
topología $\mathfrak J$-preádica, del $A$-módulo
$\operatorname{Hom}_{\mathcal O_X}(\mathcal F,\mathcal G)$. Este isomorfismo se obtiene pasando
a las completaciones separadas en el homomorfismo
$u\to\widehat u$.
\end{corollary}


\subsection{Relaciones entre morfismos formales y morfismos ordinarios}
\label{subsection:III.4.6}

\begin{proposition}[4.6.1]
\label{III.4.6.1-fr}
Sean $Y$ un preesquema localmente noetheriano, $f:X\to Y$ un morfismo
propio, $\mathcal F$ un $\mathcal O_X$-Módulo coherente y $f$-plano, $y$ un
punto de $Y$. Supongamos que para un
\oldpage[III]{140}
entero $n$ se tiene
$H^n\bigl(f^{-1}(y),\mathcal F\otimes_{\mathcal O_Y}\boldsymbol{k}(y)\bigr)=0$.
Entonces existe una vecindad $U$ de $y$ en $Y$ tal que
$R^nf_*(\mathcal F)|U=0$, y para todo entero $p\geq0$ el homomorfismo
canónico
\[
\bigl(R^{n-1}f_*(\mathcal F)\bigr)_y
\longrightarrow
H^{n-1}\bigl(f^{-1}(y),
\mathcal F\otimes_{\mathcal O_Y}
(\mathcal O_y/\mathfrak m_y^{p+1})\bigr)
\]
\textup{(\hyperref[III.4.2.1.1-fr]{4.2.1.1})} es sobreyectivo.
\end{proposition}

\begin{proof}
Como $R^nf_*(\mathcal F)$ es un $\mathcal O_Y$-Módulo coherente
(\hyperref[III.3.2.1-fr]{3.2.1}), la primera afirmación de la proposición
quedará establecida si se prueba que se tiene
$(R^nf_*(\mathcal F))_y=0$ (\textbf{0}\textsubscript{I}, 5.2.2); en virtud de
(\hyperref[III.4.2.1-fr]{4.2.1}), bastará probar que
$H^n\bigl(f^{-1}(y),\mathcal F\otimes_{\mathcal O_Y}
(\mathcal O_y/\mathfrak m_y^{p+1})\bigr)=0$
para todo $p$. Esto es verdadero por hipótesis para $p=0$; vamos a demostrarlo
por inducción sobre $p$. Pongamos
$X_p=X\times_Y\operatorname{Spec}
(\mathcal O_y/\mathfrak m_y^{p+1})$,
de modo que $X_{p-1}$ es un subpreesquema cerrado de $X_p$, con el mismo
espacio subyacente (\textbf{I}, 3.6.1); la hipótesis de inducción
$H^n\bigl(X_{p-1},\mathcal F\otimes_{\mathcal O_Y}
(\mathcal O_y/\mathfrak m_y^p)\bigr)=0$
implica pues
$H^n\bigl(X_p,\mathcal F\otimes_{\mathcal O_Y}
(\mathcal O_y/\mathfrak m_y^p)\bigr)=0$;
por otra parte, la sucesión exacta de cohomología da, a partir de la sucesión
exacta
\[
0\longrightarrow
\mathfrak m_y^p\mathcal F/\mathfrak m_y^{p+1}\mathcal F
\longrightarrow
\mathcal F/\mathfrak m_y^{p+1}\mathcal F
\longrightarrow
\mathcal F/\mathfrak m_y^p\mathcal F
\longrightarrow0
\]
de $\mathcal O_{X_p}$-Módulos, la sucesión exacta
\[
H^n\bigl(X_p,
\mathfrak m_y^p\mathcal F/\mathfrak m_y^{p+1}\mathcal F\bigr)
\longrightarrow
H^n\bigl(X_p,\mathcal F/\mathfrak m_y^{p+1}\mathcal F\bigr)
\longrightarrow
H^n\bigl(X_p,\mathcal F/\mathfrak m_y^p\mathcal F\bigr)
\]
y bastará mostrar que se tiene
\[
\label{III.4.6.1.1-fr}
H^n\bigl(X_p,
\mathfrak m_y^p\mathcal F/\mathfrak m_y^{p+1}\mathcal F\bigr)=0
\tag{4.6.1.1}
\]
pues entonces $H^n(X_p,\mathcal F/\mathfrak m_y^{p+1}\mathcal F)$ será un
submódulo de $H^n(X_p,\mathcal F/\mathfrak m_y^p\mathcal F)$, luego $0$ en
virtud de la hipótesis de inducción.

Observemos ahora que la fibra
$Z=f^{-1}(y)=X\times_Y\operatorname{Spec}(\boldsymbol{k}(y))$ es un subpreesquema cerrado
de $X_p$, y que
$\mathfrak m_y^p\mathcal F/\mathfrak m_y^{p+1}\mathcal F$ está anulado por
$\mathfrak m_y$, luego puede considerarse como un $\mathcal O_Z$-Módulo, de
modo que
$H^n\bigl(Z,\mathfrak m_y^p\mathcal F/
\mathfrak m_y^{p+1}\mathcal F\bigr)
=H^n\bigl(X_p,\mathfrak m_y^p\mathcal F/
\mathfrak m_y^{p+1}\mathcal F\bigr)$.
Dicho esto, vamos a mostrar que el $\mathcal O_Z$-homomorfismo canónico
\[
\label{III.4.6.1.2-fr}
(\mathcal F/\mathfrak m_y\mathcal F)
\otimes_{\mathbf{k}(y)}(\mathfrak m_y^p/\mathfrak m_y^{p+1})
\longrightarrow
\mathfrak m_y^p\mathcal F/\mathfrak m_y^{p+1}\mathcal F
\tag{4.6.1.2}
\]
es \emph{biyectivo}; establecido esto, resultará, puesto que
$\mathfrak m_y^p/\mathfrak m_y^{p+1}$ es un $\boldsymbol{k}(y)$-módulo libre, que se tiene
\[
\begin{split}
H^n\bigl(Z,\mathfrak m_y^p\mathcal F/
\mathfrak m_y^{p+1}\mathcal F\bigr)
&=H^n\bigl(Z,\mathcal F/\mathfrak m_y\mathcal F\bigr)
\otimes_{\mathbf{k}(y)}(\mathfrak m_y^p/\mathfrak m_y^{p+1})\\
&=0
\end{split}
\]
(\textbf{0}, 12.2.3), puesto que
$H^n(Z,\mathcal F/\mathfrak m_y\mathcal F)=0$ por hipótesis, de donde
(\hyperref[III.4.6.1.1-fr]{4.6.1.1}). Para establecer la primera afirmación,
queda pues por probar que (\hyperref[III.4.6.1.2-fr]{4.6.1.2}) es biyectivo;
como la cuestión es puntual sobre $X$ y $\mathcal F_x$ es un
$\mathcal O_y$-módulo \emph{plano} por hipótesis para todo $x\in f^{-1}(y)$,
basta aplicar (\textbf{0}, 10.2.1, \emph{c)}), pues
$\mathcal F_x/\mathfrak m_y\mathcal F_x$ es un módulo plano sobre el cuerpo
$\boldsymbol{k}(y)=\mathcal O_y/\mathfrak m_y$.

Para demostrar la segunda afirmación de (\hyperref[III.4.6.1-fr]{4.6.1}), se
reduce inmediatamente, como en (\hyperref[III.4.2.1-fr]{4.2.1}), al caso en que
$Y$ es afín y $y$ cerrado. Observemos que
(\hyperref[III.4.6.1.1-fr]{4.6.1.1}) da, por un razonamiento análogo,
para todo $k>0$, la relación
\[
\label{III.4.6.1.3-fr}
H^n\bigl(X_{k+1},
\mathfrak m_y^k\mathcal F/
\mathfrak m_y^{k+p+1}\mathcal F\bigr)=0
\tag{4.6.1.3}
\]
\oldpage[III]{141}
de donde se deduce, por (\hyperref[III.4.2.1-fr]{4.2.1}), que se tiene también
\[
\label{III.4.6.1.4-fr}
\bigl(R^nf_*(\mathfrak m_y^k\mathcal F)\bigr)_y=0.
\tag{4.6.1.4}
\]

Dicho esto, se obtiene de la sucesión exacta de cohomología la exactitud de la
sucesión
\[
\bigl(R^{n-1}f_*(\mathcal F)\bigr)_y
\longrightarrow
\bigl(R^{n-1}f_*(\mathcal F/\mathfrak m_y^p\mathcal F)\bigr)_y
\longrightarrow
\bigl(R^nf_*(\mathfrak m_y^p\mathcal F)\bigr)_y=0
\]
y como $y$ es cerrado y $Y$ afín, se tiene (\hyperref[III.1.4.11-fr]{1.4.11})
\[
R^{n-1}f_*(\mathcal F/\mathfrak m_y^p\mathcal F)
=\bigl(H^{n-1}(X,\mathcal F/\mathfrak m_y^p\mathcal F)\bigr)^\sim
=\bigl(H^{n-1}(f^{-1}(y),
\mathcal F/\mathfrak m_y^p\mathcal F)\bigr)^\sim
\]
(G, II, 4.9.1); ahora bien,
$H^{n-1}(f^{-1}(y),\mathcal F/\mathfrak m_y^p\mathcal F)$ es un
$(\mathcal O_y/\mathfrak m_y^p)$-módulo, de donde
\[
\bigl(R^{n-1}f_*(\mathcal F/\mathfrak m_y^p\mathcal F)\bigr)_y
=H^{n-1}(f^{-1}(y),\mathcal F/\mathfrak m_y^p\mathcal F)
\]
y esto termina de probar (\hyperref[III.4.6.1-fr]{4.6.1}).
\end{proof}

\begin{corollary}[4.6.2]
\label{III.4.6.2-fr}
Sean $Y$ un preesquema localmente noetheriano, $f:X\to Y$ un morfismo propio
y plano, $\mathcal F$, $\mathcal G$ dos $\mathcal O_X$-Módulos localmente
libres, $y$ un punto de $Y$. Pongamos
$X_y=f^{-1}(y)=X\otimes_Y\boldsymbol{k}(y)$,
$\mathcal F_y=\mathcal F\otimes_{\mathcal O_Y}\boldsymbol{k}(y)$,
$\mathcal G_y=\mathcal G\otimes_{\mathcal O_Y}\boldsymbol{k}(y)$,
y supongamos que se tiene
\[
\label{III.4.6.2.1-fr}
H^1\bigl(X_y,
\mathcal{H}\!om_{\mathcal O_{X_y}}(\mathcal F_y,\mathcal G_y)\bigr)=0.
\tag{4.6.2.1}
\]

Entonces, para todo homomorfismo $u_0:\mathcal F_y\to\mathcal G_y$, existe
una vecindad abierta $U$ de $Y$ y un homomorfismo
$u:F|f^{-1}(U)\longrightarrow\mathcal G|f^{-1}(U)$
tales que $u_0$ es igual al homomorfismo $u\otimes1$.
\end{corollary}

\begin{proof}
En efecto, la hipótesis permite aplicar (\hyperref[III.4.6.1-fr]{4.6.1}) al
$\mathcal O_X$-Módulo coherente
$\mathcal H=\mathcal{H}\!om_{\mathcal O_X}(\mathcal F,\mathcal G)$ para
$n=1$ y $p=0$, pues $\mathcal H$ es localmente libre y \emph{a fortiori}
$f$-plano, y el $\mathcal O_{X_y}$-Módulo
$\mathcal H\otimes_{\mathcal O_Y}\boldsymbol{k}(y)$ se identifica entonces con
$\mathcal{H}\!om_{\mathcal O_{X_y}}(\mathcal F_y,\mathcal G_y)$
(\textbf{0}\textsubscript{I}, 6.2.2). Se puede suponer $Y=\operatorname{Spec}(A)$
afín, y entonces (\hyperref[III.1.4.11-fr]{1.4.11})
$R^0f_*(\mathcal H)=
\bigl(\operatorname{Hom}_{\mathcal O_X}(\mathcal F,\mathcal G)\bigr)^\sim$,
luego
$\bigl(R^0f_*(\mathcal H)\bigr)_y
=\operatorname{Hom}_{\mathcal O_X}(\mathcal F,\mathcal G)
\otimes_A\mathcal O_y$;
al ser el homomorfismo canónico
\[
\operatorname{Hom}_{\mathcal O_X}(\mathcal F,\mathcal G)
\otimes_A\mathcal O_y
\longrightarrow
\operatorname{Hom}_{\mathcal O_{X_y}}(\mathcal F_y,\mathcal G_y)
\]
sobreyectivo por (\hyperref[III.4.6.1-fr]{4.6.1}), esto establece el
corolario, pues todo elemento de
$\operatorname{Hom}_{\mathcal O_X}(\mathcal F,\mathcal G)
\otimes_A\mathcal O_y$ puede siempre ponerse en la forma
$u\otimes(1/s)$, donde $s\notin\mathfrak m_y$ es un elemento de $A$.
\end{proof}

Este corolario puede completarse con el siguiente:

\begin{corollary}[4.6.3]
\label{III.4.6.3}
\label{III.4.6.3-fr}
Bajo las hipótesis de \sref{III.4.6.2}, si $u_0$ es inyectivo \textup{(resp.} sobreyectivo, biyectivo), se puede suponer que lo mismo ocurre con $u$.
\end{corollary}

\begin{proof}
Se puede limitar al caso en que $U=Y$. Basta probar
que, si $u_0$ es
inyectivo (resp. sobreyectivo), $\operatorname{Ker}u_x=0$ (resp.
$\operatorname{Coker}u_x=0$) para todo $x\in f^{-1}(y)$. En efecto, $\operatorname{Ker}u$ y $\operatorname{Coker}u$ son
$\mathcal O_X$-módulos coherentes \sref[0]{0.5.3.4}; por tanto, existe una
vecindad $V$ de $f^{-1}(y)$ en $X$ sobre la cual la restricción de
$\operatorname{Ker}u$ (resp. $\operatorname{Coker}u$) a $V$ es $0$
(\textbf{0}\textsubscript{I}, 5.2.2); como $f$ es cerrado, existe una vecindad abierta
$U'\subset U$ de $y$ tal que $f^{-1}(U')\subset V$, y
(\hyperref[III.4.6.3-fr]{4.6.3}) quedará demostrado. Por hipótesis,
$u_x\otimes1:
\mathcal F_x\otimes_{\mathcal O_y}\boldsymbol{k}(y)
\longrightarrow
\mathcal G_x\otimes_{\mathcal O_y}\boldsymbol{k}(y)$
es inyectivo (resp. sobreyectivo),
\oldpage[III]{142}
$\mathcal F_x$ y $\mathcal G_x$
son $\mathcal O_x$-módulos libres finitos, y $\mathcal O_x$ es plano sobre
$\mathcal O_y$. Si $u_x\otimes1$ es inyectivo, la inyectividad de $u_x$
se sigue de \sref[0]{0.10.2.4}. Si $u_x\otimes1$ es sobreyectivo,
entonces, con mayor razón, el homomorfismo inducido
$\mathcal F_x/\mathfrak m_x\mathcal F_x
\longrightarrow
\mathcal G_x/\mathfrak m_x\mathcal G_x$, que se deduce por paso a los cocientes, es sobreyectivo; como $\mathcal G_x$ es un $\mathcal O_x$-módulo finito y
$\mathcal O_x$ es un anillo local de ideal maximal $\mathfrak m_x$, la
conclusión se sigue del lema de Nakayama (Bourbaki,
\emph{Alg\`ebre}, cap.~VIII, \S\,6, n.º~3, cor.~4 de la prop.~6).
\end{proof}

\begin{corollary}[4.6.4]
\label{III.4.6.4}
\label{III.4.6.4-fr}
Sean $Y$ un preesquema localmente noetheriano, $f:X\to Y$ un
morfismo propio y plano, e $y$ un punto de $Y$, $X_y=X\otimes_Y\boldsymbol{k}(y)$. Sea $\mathcal E_0$ un
$\mathcal O_{X_y}$-módulo localmente libre tal que
\[
\label{III.4.6.4.1-fr}
H^1\bigl(X_y,
\mathcal{H}\!om_{\mathcal O_{X_y}}(\mathcal E_0,\mathcal E_0)\bigr)=0.
\tag{4.6.4.1}
\]
Sean $\mathcal F$, $\mathcal G$ dos $\mathcal O_X$-módulos localmente libres
tales que $\mathcal F_y$ y $\mathcal G_y$, con la notación de
\sref{III.4.6.2}, son isomorfos a $\mathcal E_0$. Entonces existe una
vecindad abierta $U$ de $y$ tal que $\mathcal F|f^{-1}(U)$ y
$\mathcal G|f^{-1}(U)$ son isomorfos.
\end{corollary}

Más particularmente:

\begin{corollary}[4.6.5]
\label{III.4.6.5}
\label{III.4.6.5-fr}
Bajo las hipótesis de \sref{III.4.6.4} sobre $f$, $X$ e $Y$,
supongamos que
$H^1(X_y,\mathcal O_{X_{\boldsymbol{y}}})=0$.
Si $\mathcal F$ y $\mathcal G$ son dos $\mathcal O_X$-módulos invertibles tales
que $\mathcal F_y$ y $\mathcal G_y$ son isomorfos, entonces existe una vecindad
abierta $U$ de $y$
tal que $\mathcal F|f^{-1}(U)$ y $\mathcal G|f^{-1}(U)$ son isomorfos.
\end{corollary}

\begin{proof}
Basta aplicar (\hyperref[III.4.6.4-fr]{4.6.4}) a los módulos
$\mathcal F^{-1}\otimes\mathcal G$ y $\mathcal O_X$.
\end{proof}

\begin{remarks}[4.6.6]
\label{III.4.6.6}
\label{III.4.6.6-fr}
\begin{enumerate}
\item[(i)] Utilizando \sref{III.4.6.5}, estableceremos en el capítulo~V
  la clasificación de los haces invertibles sobre un fibrado proyectivo
  anunciada en \sref[II]{II.4.2.7}.
\item[(ii)] El resultado de \sref{III.4.6.1} aparecerá en \S\,7 como
  consecuencia de proposiciones más generales.
\end{enumerate}
\end{remarks}

\begin{proposition}[4.6.7]
\label{III.4.6.7-fr}
Sean $Z$ un preesquema localmente noetheriano, $X$, $Y$ dos $Z$-preesquemas
tales que los morfismos estructurales $g:X\to Z$, $h:Y\to Z$ son propios.
Sean $f:X\to Y$ un $Z$-morfismo, $z$ un punto de $Z$, y sea
$f_z=f\times_Z1:
X\otimes_Z\boldsymbol{k}(z)\longrightarrow Y\otimes_Z\boldsymbol{k}(z)$.
\begin{enumerate}
\item[\textup{(i)}] Si $f_z$ es un morfismo finito \textup{(resp.} una inmersión cerrada),
existe una vecindad abierta $U$ de $z$ tal que el morfismo
$g^{-1}(U)\to h^{-1}(U)$, restricción de $f$, es un morfismo finito \textup{(resp.}
una inmersión cerrada).
\item[\textup{(ii)}] Se supone además que $g$ es un morfismo plano. Entonces, si
$f_z$ es un isomorfismo, existe una vecindad abierta $U$ de $z$ tal que
el morfismo $g^{-1}(U)\to h^{-1}(U)$, restricción de $f$, es un
isomorfismo.
\end{enumerate}
\end{proposition}

\begin{proof}
En ambos casos, bastará probar que para todo $y\in h^{-1}(z)$
existe una vecindad $V_{\boldsymbol{y}}$ de $y$ tal que la restricción
$f^{-1}(V_y)\to V_y$ de $f$ es un morfismo finito (resp. una inmersión
cerrada, un isomorfismo); resultará entonces en efecto que si $V$ es la
unión de los $V_{\boldsymbol{y}}$, la restricción $f^{-1}(V)\to V$ de $f$ es un morfismo
finito (resp. una inmersión cerrada, un isomorfismo) (\textbf{II}, 6.1.1 y I, 4.2.4).
Como $h$ es un morfismo cerrado, existirá una vecindad abierta $U$ de $z$
tal que $h^{-1}(U)\subset V$, y la proposición quedará demostrada.

\noindent (i) Observemos ante todo que $f$ es un morfismo propio
(II, 5.4.3); si se
\oldpage[III]{143}
supone $f_z$ finito, la existencia para todo
$y\in h^{-1}(z)$ de una vecindad $V_y$ tal que $f^{-1}(V_y)\to V_y$ es finito
resulta de (\hyperref[III.4.4.11-fr]{4.4.11}).
Para tratar el caso en que $f_z$ es una inmersión cerrada, se puede pues ya
suponer que el morfismo $f$ es \emph{finito}, luego que
$X=\operatorname{Spec}(\mathcal B)$, donde $\mathcal B$ es una
$\mathcal O_Y$-Álgebra coherente, correspondiendo el morfismo $f$ (\textbf{II}, 1.2.7)
al homomorfismo canónico $u:\mathcal O_Y\to\mathcal B$. Si se prueba
que para todo $y\in h^{-1}(z)$ el homomorfismo
$u_y:\mathcal O_y\to\mathcal B_y$ es sobreyectivo, resultará que para una
vecindad $V_y$ de $y$, $u|V_y$ será sobreyectivo, al ser el haz
$\operatorname{Coker}u$ coherente (\textbf{0}\textsubscript{I}, 5.3.4 y 5.2.2).
Dicho esto, el morfismo finito $f_z$ corresponde al homomorfismo
$v=u\otimes1:
\mathcal O_Y\otimes_{\mathcal O_Z}\boldsymbol{k}(z)
\longrightarrow
\mathcal B\otimes_{\mathcal O_Z}\boldsymbol{k}(z)$,
y la hipótesis de que $f_z$ es una inmersión cerrada implica que
el homomorfismo
$u_y\otimes1:
\mathcal O_y\otimes_{\mathcal O_z}\boldsymbol{k}(z)
\longrightarrow
\mathcal B_y\otimes_{\mathcal O_z}\boldsymbol{k}(z)$
es sobreyectivo. Como $\mathcal B_y$ es un $\mathcal O_y$-módulo de tipo
finito y $\mathcal O_y$ un anillo local noetheriano, la conclusión resulta,
como en (\hyperref[III.4.6.3-fr]{4.6.3}), del lema de Nakayama.

\noindent (ii) El mismo razonamiento que antes muestra que
esta vez basta probar que $u_y$ es biyectivo, sabiendo que
$u_y\otimes1$ es biyectivo.

Esto resultará del lema siguiente:
\end{proof}

\begin{lemma}[4.6.7.1]
\label{III.4.6.7.1-fr}
Sean $A$, $B$ dos anillos locales noetherianos,
$\rho:A\to B$ un homomorfismo local, $u:N\to M$ un homomorfismo de
$B$-módulos. Se supone que $M$ es un $A$-módulo plano, $N$ un $B$-módulo
de tipo finito y que
$u\otimes1:N\otimes_Ak\longrightarrow M\otimes_Ak$
(donde $k$ es el cuerpo residual de $A$) es inyectivo. Entonces $N$ es un
$A$-módulo plano y $u$ es inyectivo.
\end{lemma}

\begin{proof}
Para establecer la primera afirmación, hay que mostrar que para todo par de
$A$-módulos \emph{de tipo finito} $P$, $Q$ y todo $A$-homomorfismo inyectivo
$v:P\to Q$,
$1_N\otimes v:N\otimes_AP\to N\otimes_AQ$ es inyectivo. Ahora bien, se tiene el
diagrama conmutativo
\[
\xymatrix{
N\otimes_AP\ar[r]^{1_N\otimes v}\ar[d]_{u\otimes1_P}
  &N\otimes_AQ\ar[d]^{u\otimes1_Q}\\
M\otimes_AP\ar[r]_{1_M\otimes v}
  &M\otimes_AQ
}
\]
y como $1_M\otimes v$ es inyectivo por hipótesis, bastará probar
que lo mismo ocurre con $u\otimes1_P$. Sea $\mathfrak m$ el ideal maximal
de $A$; la filtración $\mathfrak m$-ádica sobre el $A$-módulo
$N\otimes_AP$ es también su filtración $\mathfrak mB$-ádica como
$B$-módulo; la topología definida por esta filtración es pues \emph{separada},
puesto que $B$ es noetheriano, $\mathfrak mB$ está contenido en el radical de
$B$, y $N\otimes_AP$ es un $B$-módulo de tipo finito, al ser $N$ un
$B$-módulo de tipo finito y $P$ un $A$-módulo de tipo finito
(\textbf{0}\textsubscript{I}, 7.3.5). Basta pues probar que el homomorfismo
$\operatorname{gr}(u\otimes1_P):
\operatorname{gr}_\bullet(N\otimes_AP)
\longrightarrow
\operatorname{gr}_\bullet(M\otimes_AP)$
(donde los módulos graduados son relativos a las filtraciones $\mathfrak m$-ádicas)
es inyectivo (Bourbaki, \emph{Alg. comm.}, cap. III, \S\,2, n\textsuperscript{o}~8, cor. 1 del
teorema 1). Observemos ahora que, puesto que $M$ es un $A$-módulo plano, los
homomorfismos
$M\otimes_A(\mathfrak m^nP)
\longrightarrow
\mathfrak m^n(M\otimes_AP)$
son biyectivos; lo mismo ocurre pues con el homomorfismo canónico
\[
\varphi_M:
\operatorname{gr}_0(M)\otimes_A\operatorname{gr}_\bullet(P)
\longrightarrow
\operatorname{gr}_\bullet(M\otimes_AP).
\]

\oldpage[III]{144}

Ahora bien, se tiene un diagrama conmutativo
\[
\xymatrix{
\operatorname{gr}_0(N)\otimes_A\operatorname{gr}_\bullet(P)
  \ar[r]^{\operatorname{gr}_0(u)\otimes1}\ar[d]_{\varphi_N}
  &\operatorname{gr}_0(M)\otimes_A\operatorname{gr}_\bullet(P)
  \ar[d]^{\varphi_M}\\
\operatorname{gr}_\bullet(N\otimes_AP)
  \ar[r]_{\operatorname{gr}(u\otimes1)}
  &\operatorname{gr}_\bullet(M\otimes_AP)
}
\]
en el cual $\varphi_M$ es biyectivo, $\varphi_N$ sobreyectivo; además,
$\operatorname{gr}_0(u)$ es inyectivo por hipótesis, y como
$\operatorname{gr}_0(N)\otimes_A\operatorname{gr}_\bullet(P)
=\operatorname{gr}_0(N)\otimes_k\operatorname{gr}_\bullet(P)$,
$\operatorname{gr}_0(M)\otimes_A\operatorname{gr}_\bullet(P)
=\operatorname{gr}_0(M)\otimes_k\operatorname{gr}_\bullet(P)$,
$\operatorname{gr}_0(u)\otimes1$ es también inyectivo. Se concluye que
$\operatorname{gr}(u\otimes1)$ es inyectivo, lo que termina de demostrar la
primera afirmación. La segunda se deduce del razonamiento anterior haciendo
$P=A$.
\end{proof}

\begin{proposition}[4.6.8]
\label{III.4.6.8-fr}
Sean $Z$ un preesquema localmente noetheriano, $X$, $Y$, dos $Z$-preesquemas
tales que los morfismos estructurales $g:X\to Z$, $h:Y\to Z$ son propios,
$Z'$ una parte cerrada de $Z$, $X'=g^{-1}(Z')$, $Y'=h^{-1}(Z')$ sus imágenes
inversas,
$\widehat X=X_{/X'}$, $\widehat Y=Y_{/Y'}$, $\widehat Z=Z_{/Z'}$
los completados formales de $X$, $Y$, $Z$ a lo largo de estas partes cerradas,
$f:X\to Y$ un $Z$-morfismo, $\widehat f:\widehat X\to\widehat Y$ su
prolongación a los completados. Para que $\widehat f$ sea un isomorfismo (\textup{resp.}
una inmersión cerrada), es necesario y suficiente que exista una vecindad abierta
$U$ de $Z'$ tal que el morfismo $g^{-1}(U)\to h^{-1}(U)$, restricción de
$f$, sea un isomorfismo (\textup{resp.} una inmersión cerrada).
\end{proposition}

\begin{proof}
La suficiencia de la condición es inmediata (\textbf{I}, 10.14.7). Para mostrar su
necesidad, basta también probar que para todo $y\in Y'$ existe una
vecindad abierta $V_y$ de $y$ tal que la restricción
$f^{-1}(V_y)\to V_y$ de $f$ es un isomorfismo (resp. una inmersión
cerrada), por el mismo razonamiento que en
(\hyperref[III.4.6.7-fr]{4.6.7}). Se está así reducido al caso en que $Y=Z$,
siendo $Y=\operatorname{Spec}(A)$ afín noetheriano. Por hipótesis
(I, 10.9.1 y 10.14.2), la fibra $f^{-1}(y)$ se reduce a un punto para
$y\in Y'$, luego, como $f$ es propio (\textbf{II}, 5.4.3), existe una vecindad
abierta $U$ de $y$ tal que la restricción $f^{-1}(U)\to U$ de $f$ es un
morfismo \emph{finito} (\hyperref[III.4.4.11-fr]{4.4.11}). Se puede pues ya suponer
que $f$ es un morfismo finito, luego $X=\operatorname{Spec}(B)$, donde $B$ es
una $A$-álgebra \emph{finita} sobre $A$. Si $Y'=V(\mathfrak J)$, se tiene entonces
$\widehat Y=\operatorname{Spf}(\widehat A)$, $\widehat X=\operatorname{Spf}(\widehat B)$,
siendo $\widehat A$ el separado completado de $A$ para la topología
$\mathfrak J$-preádica, $\widehat B$ el separado completado de $B$ para la
topología $\mathfrak JB$-preádica, o (lo que equivale a lo mismo), el separado
completado del $A$-módulo $B$ para la topología $\mathfrak J$-preádica;
además, $\widehat f$ es el morfismo de esquemas formales afines correspondiente
a la prolongación continua $\widehat\varphi:\widehat A\to\widehat B$ del
homomorfismo canónico de anillos $\varphi:A\to B$, y la hipótesis es que
$\widehat\varphi$ es sobreyectivo (resp. biyectivo) (\textbf{I}, 10.14.2). Ahora bien,
$\widehat\varphi$ es también la prolongación continua de $\varphi$ considerado
como homomorfismo de $A$-módulos; se sabe entonces (\textbf{I}, 10.8.14) que existe
una vecindad abierta $U$ de $Y'$ tal que la restricción a $U$ del
homomorfismo $\widetilde\varphi:\widetilde A\to\widetilde B$ de
$\mathcal O_Y$-Módulos es sobreyectiva (resp. biyectiva), lo que termina la
demostración.
\end{proof}

\oldpage[III]{145}

\subsection{Un criterio de amplitud}

\begin{theorem}[4.7.1]
\label{III.4.7.1-fr}
Sean $Y$ un preesquema localmente noetheriano, $f:X\to Y$ un morfismo
propio, $\mathcal L$ un $\mathcal O_X$-Módulo invertible, $y$ un punto de
$Y$, $X_y=X\otimes_Y\boldsymbol{k}(y)=f^{-1}(y)$, $g$ la proyección de $X_y$ en $X$.
Si $\mathcal L_y=g^*(\mathcal L)=\mathcal L\otimes_{\mathcal O_Y}\boldsymbol{k}(y)$ es
amplio sobre $X_y$, existe una vecindad abierta $U$ de $y$ en $Y$ tal que
$\mathcal L|f^{-1}(U)$ es amplio para la restricción de $f$ a $f^{-1}(U)$.
\end{theorem}

\begin{proof}
\noindent I) Pongamos
$Y'=\operatorname{Spec}(\mathcal O_y)$, $X'=X\times_YY'$, y sea
$\mathcal L'=\mathcal L\otimes_{\mathcal O_Y}\mathcal O_y$;
vamos primero a probar que $\mathcal L'$ es \emph{amplio para}
$f'=f_{(Y')}$. Se tiene el diagrama conmutativo
\[
\xymatrix{
X\ar[d]_f & X'\ar[l]\ar[d]_{f'} & X_y\ar[l]\ar[d]_{f_y}\\
Y & Y'\ar[l] & \operatorname{Spec}(\boldsymbol{k}(y))\ar[l]
}
\]
Como $f'$ es propio (\textbf{II}, 5.4.2, (iii)) y $\mathcal O_y$ noetheriano, se
ve que se puede limitar al caso en que
$Y=Y'=\operatorname{Spec}(\mathcal O_y)$, luego $X=X'$, suponer que
$\mathcal L_y$ es amplio para $f_y$ y probar que $\mathcal L$ es amplio
para $f$ (\textbf{II}, 4.6.6). Vamos a aplicar el criterio
(\hyperref[III.2.6.1-fr]{2.6.1, \emph{c)}}) y a mostrar de hecho que para todo
$\mathcal O_X$-Módulo coherente $\mathcal F$ existe un entero $N$ tal que
$H^1(X,\mathcal F(n))=0$ para todo $n\geq N$, con
$\mathcal F(n)=\mathcal F\otimes\mathcal L^{\otimes n}$.
Observemos que $y$ es un punto cerrado de $Y$ correspondiente al ideal maximal
$\mathfrak m$ de $\mathcal O_y$; $X_y$ es pues un subpreesquema cerrado de
$X$ definido \textbf{por} el Ideal coherente
$\mathcal J=f^*(\widetilde{\mathfrak m})\mathcal O_X=
\mathfrak m\mathcal O_X$ de $\mathcal O_X$ (\textbf{I}, 4.4.5), y $g$ la inclusión
canónica. Consideremos entonces la $\boldsymbol{k}(y)$-álgebra graduada
$S=\operatorname{gr}(\mathcal O_y)
=\bigoplus_{j\geq0}\mathfrak m^j/\mathfrak m^{j+1}$,
que es de tipo finito puesto que $\mathcal O_y$ es noetheriano; la
$\mathcal O_X$-Álgebra $\mathcal S=f^*(\widetilde S)$ es pues
cuasicoherente y de tipo finito, y está evidentemente anulada por
$\mathcal J$, luego, si se pone $\mathcal S_y=g^*(\mathcal S)$,
$\mathcal S_y$ es una $\mathcal O_{X_y}$-Álgebra cuasicoherente de tipo
finito, y $\mathcal S=g_*(\mathcal S_y)$. Pongamos por otra parte
$\mathcal M_j=\mathfrak m^j\mathcal F/\mathfrak m^{j+1}\mathcal F$ y
$\mathcal M=\bigoplus_{j\geq0}\mathcal M_j
=\operatorname{gr}(\mathcal F)$;
como $\mathcal F$ es coherente, $\mathcal M$ es un $\mathcal S$-Módulo
cuasicoherente de tipo finito (\textbf{0}, 10.1.1) que está también anulado por
$\mathcal J$, de modo que si se pone
$\mathcal M'_j=g^*(\mathcal M_j)$, $\mathcal M'=g^*(\mathcal M)=\bigoplus_{j\geq0}\mathcal M'_j$
es un $\mathcal S_y$-Módulo graduado cuasicoherente de tipo
finito tal que $\mathcal M=g_*(\mathcal M')$. Además, si se pone
$\mathcal M'_j(n)=\mathcal M'_j\otimes\mathcal L_y^{\otimes n}$, se tiene
$\mathcal M'_j(n)=g^*((\mathcal M_j(n)))$. Dicho esto, $f_y$ es propio
(\textbf{II}, 5.4.2, (iii)) y $\mathcal L_y$ es amplio, luego $f_y$ es \emph{proyectivo}
(\textbf{II}, 5.5.4 y 4.6.11), y se puede aplicar a
$\operatorname{Spec}(\boldsymbol{k}(y))$, $f_y$, $\mathcal S_y$, $\mathcal L_y$ y
$\mathcal M'$ el teorema (\hyperref[III.2.4.1-fr]{2.4.1, (ii)}):
existe un entero $N$ tal que para $n\geq N$ se tiene
$H^q(X_y,\mathcal M'_j(n))=0$ para todo $q>0$ \emph{y todo} $j$; por tanto, se tiene
también $H^q(X,\mathcal M_j(n))=0$ para todo $q>0$ y todo $j$
(G, II, 4.9.1). Pongamos entonces
$\mathcal F(n)_j=\mathcal F(n)/\mathfrak m^{j+1}\mathcal F(n)$,
de modo que $\mathcal F(n)_{j-1}=\mathcal F(n)_j/\mathcal M_j(n)$ para
$j\geq1$ y $\mathcal F(n)_0=\mathcal M_0(n)$. Se tiene
$H^1(X,\mathcal F(n)_0)=0$, y, por la sucesión exacta de cohomología,
$H^1(X,\mathcal F(n)_j)=H^1(X,\mathcal F(n)_{j-1})$ para todo $j\geq1$,
luego $H^1(X,\mathcal F(n)_j)=0$ para todo $j\geq0$. Se concluye pues de
(\hyperref[III.4.2.1-fr]{4.2.1}) que $H^1(X,\mathcal F(n))=0$, lo que
termina de probar nuestra afirmación.

\noindent II) Volvamos a las notaciones del comienzo de la demostración,
y observemos que
\oldpage[III]{146}
siempre se puede suponer $Y=\operatorname{Spec}(A)$
afín; como $f'$ es de tipo finito y $\mathcal L'$ amplio para $f'$,
existe un entero $m>0$ tal que $\mathcal L'^{\otimes m}$ es muy amplio para
$f'$ (\textbf{I}I, 4.6.11); reemplazando si es necesario $\mathcal L$ por
$\mathcal L^{\otimes m}$, se puede limitar a considerar el caso en que
$\mathcal L'$ es \emph{muy amplio} para $f'$, y a probar que
$\mathcal L|f^{-1}(U)$ es entonces muy amplio para $f$. Como $f'$ es propio,
existe entonces una $Y'$-inmersión \emph{cerrada}
$j:X'\to P=\mathbf P_{Y'}^r$ para un entero $r>0$ apropiado, tal que
$\mathcal L'$ es isomorfo a $j^*(\mathcal O_P(1))$ (\textbf{II}, 5.5.4, (ii));
esta inmersión corresponde canónicamente a un $\mathcal O_{X'}$-homomorfismo
\emph{sobreyectivo}
$u:\mathcal O_{X'}^{r+1}\longrightarrow\mathcal L'$
(\textbf{II}, 4.2.3). Este último corresponde (\textbf{0}\textsubscript{I}, 5.1.1) al dato
de $r+1$ secciones $s'_i$ ($0\leq i\leq r$) de $\mathcal L'$ sobre
$X'$ que generan este $\mathcal O_{X'}$-Módulo. Estas secciones son también por
definición secciones de $f'_*(\mathcal L')$ sobre $Y'$; se tiene
$f'_*(\mathcal L')=f_*(\mathcal L)\otimes_{\mathcal O_Y}\mathcal O_{Y'}$
(\textbf{0}\textsubscript{I}, 5.4.10), $Y$ es afín y $\mathcal O_y$ es el anillo
local en el ideal primo $\mathfrak j_y$ de $A$, luego se tiene
$s'_i=s''_i/t_i$, donde los $s''_i$ son secciones de $f_*(\mathcal L)$
sobre $Y$ y los $t_i$ elementos de $A$ que no pertenecen a
$\mathfrak j_y$; se concluye que existe una vecindad abierta afín $V$
de $y$ en $Y$ y secciones $s_i$ de $f_*(\mathcal L)|V$ tales que
$s'_i=s_i/1$ (se recuerda que el espacio $Y'$ está contenido en $V$, cf. \textbf{I},
2.4.2). Los $s_i$ son entonces secciones de $\mathcal L$ sobre
$f^{-1}(V)$, que definen pues un homomorfismo
$v:\bigl(\mathcal O_X|f^{-1}(V)\bigr)^{r+1}
\longrightarrow\mathcal L|f^{-1}(V)$
que, por hipótesis, es sobreyectivo en todos los puntos de $f^{-1}(y)$; como
$\operatorname{Coker}(v)$ es coherente (\textbf{0}\textsubscript{I}, 5.3.4), su
soporte es cerrado (\textbf{0}\textsubscript{I}, 5.2.2) y por tanto existe una
vecindad abierta $W\subset f^{-1}(V)$ de $f^{-1}(y)$ tal que la restricción
de $v$ a $W$ es un homomorfismo \emph{sobreyectivo}. Puesto que el morfismo $f$ es
cerrado, se puede suponer que $W$ es de la forma $f^{-1}(U)$, donde $U$ es una
vecindad abierta de $y$, y la conclusión resulta entonces de (\textbf{II}, 4.2.3).
\end{proof}

\subsection{Morfismos finitos de preesquemas formales}

\begin{proposition}[4.8.1]
\label{III.4.8.1}
\label{III.4.8.1-fr}
Sea $\mathfrak Y$ un preesquema formal localmente noetheriano,
$\mathcal K$ un ideal de definición de $\mathfrak Y$, y
$f:\mathfrak X\to\mathfrak Y$ un morfismo de preesquemas formales.
Las condiciones siguientes son equivalentes:
\begin{enumerate}
  \item[\rm{a)}] $\mathfrak X$ es localmente noetheriano, $f$ es un morfismo ádico \textup{(\textbf{I}, 10.12.1)} y si se pone $\mathcal J=f^*(\mathcal K)\mathcal O_{\mathfrak X}$, el morfismo
    $f_0:(\mathfrak X,\mathcal O_{\mathfrak X}/\mathcal J)
  \longrightarrow
  (\mathfrak Y,\mathcal O_{\mathfrak Y}/\mathcal K)$
    inducido por $f$ es finito.
  \item[\rm{b)}] $\mathfrak X$ es localmente noetheriano y es el
    límite inductivo de un sistema inductivo ádico de $(Y_n)$, $(X_n)$,
    tal que el morfismo $X_0\to Y_0$ es finito.
  \item[\rm{c)}] Todo punto de $\mathfrak Y$ posee una vecindad abierta
    formal afín noetheriana $V$ tal que $f^{-1}(V)$ es
    un abierto formal afín y
    $\Gamma(f^{-1}(V),\mathcal O_{\mathfrak X})$ es un
    $\Gamma(V,\mathcal O_{\mathfrak Y})$-módulo de tipo finito.
\end{enumerate}
\end{proposition}

\begin{proof}
De \sref[I]{I.10.12.3} se sigue inmediatamente que \emph{a)} implica
\emph{b)}. Para ver que \emph{b)} implica \emph{c)}, se puede suponer
que
$\mathfrak Y=\operatorname{Spf}(B)$,
donde $B$ es ádico noetheriano y
$\mathcal K=\mathfrak R^\Delta$, donde $\mathfrak R$ es un ideal
de definición de $B$. Por hipótesis, $X_0$ es un esquema afín
cuyo anillo $A_0$ es un $B/\mathfrak R$-módulo de tipo finito
\sref[II]{II.6.1.3}. Por \sref[I]{I.5.1.9}, cada $X_n$ es un esquema
afín; si $A_n$ denota su anillo, la condición \emph{b)} implica que,
para $m\leq n$, $A_m$ es isomorfo a
$A_n/\mathfrak R^{m+1}A_n$. Se deduce que $\mathfrak X$ es isomorfo a
$\operatorname{Spf}(A)$, donde $A=\varprojlim_n A_n$; se concluye en virtud de
(\textbf{0}\textsubscript{I}, 7.2.9). Finalmente, para probar que \emph{c)}
\oldpage[III]{147}
implica \emph{a)}, se puede limitar también al caso
$\mathfrak Y=\operatorname{Spf}(B)$, $\mathfrak X=\operatorname{Spf}(A)$,
$A$ es una $B$-álgebra finita. Como
$A/\mathfrak RA$ es entonces una
$B/\mathfrak R$-álgebra finita, \sref[I]{I.10.10.9} muestra que se
satisfacen las condiciones de \emph{a)}.
\end{proof}

\begin{definition}[4.8.2]
\label{III.4.8.2}
\label{III.4.8.2-fr}
Cuando se satisfacen las condiciones equivalentes \textup{a)}, \textup{b)} y \textup{c)}
de \sref{III.4.8.1}, se dice que el morfismo $f$ es
finito; también se dice que $\mathfrak X$ es un
$\mathfrak Y$-preesquema formal finito, o un preesquema
formal finito sobre $\mathfrak Y$.
\end{definition}

\begin{proposition}[4.8.3]
\label{III.4.8.3}
\label{III.4.8.3-fr}
\begin{enumerate}
  \item[\rm{(i)}] Una inmersión cerrada de preesquemas formales localmente
    noetherianos es un morfismo finito.
  \item[\rm{(ii)}] El compuesto de dos morfismos finitos de preesquemas formales
    localmente noetherianos es un morfismo finito.
  \item[\rm{(iii)}] Sean $\mathfrak X$, $\mathfrak Y$ y
    $\mathfrak S$ preesquemas formales localmente noetherianos,
    $f:\mathfrak X\to\mathfrak S$ un morfismo finito, y
    $g:\mathfrak Y\to\mathfrak S$ un morfismo. Entonces la proyección
    $\mathfrak X\times_{\mathfrak S}\mathfrak Y\to\mathfrak Y$
    es finita.
  \item[\rm{(iv)}] Sea $\mathfrak S$ un preesquema formal localmente
    noetheriano, y sean $\mathfrak X'$ y $\mathfrak Y'$
    preesquemas formales localmente noetherianos tales que
    $\mathfrak X'\times_{\mathfrak S}\mathfrak Y'$ es localmente
    noetheriano. Si $\mathfrak X$ y $\mathfrak Y$ son
    $\mathfrak S$-preesquemas formales localmente noetherianos,
  $f:\mathfrak X\to\mathfrak X'$, $g:\mathfrak Y\to\mathfrak Y'$ dos
  $S$-morfismos finitos, entonces
    $f\times_{\mathfrak S}g$ es un
  morfismo finito.
\item[\textup{(v)}] Sean $f:\mathfrak X\to\mathfrak Y$,
  $g:\mathfrak Y\to\mathfrak Z$ dos morfismos de preesquemas formales
  localmente noetherianos tales que $g$ es de tipo finito y separado. Si $g\circ f$ es finito,
    entonces $f$ es finito.
\end{enumerate}
\end{proposition}

\begin{proof}
La afirmación \emph{(i)} es inmediata. Utilizando el criterio \emph{a)} de
\sref{III.4.8.1}, las demás afirmaciones se reducen de inmediato a las
proposiciones correspondientes para morfismos de preesquemas ordinarios
\sref[II]{II.6.1.5}. Dejamos los detalles al lector, siguiendo
el modelo de \sref[I]{I.10.13.5}.
\end{proof}

\begin{corollary}[4.8.4]
\label{III.4.8.4}
\label{III.4.8.4-fr}
Bajo las hipótesis de \sref[I]{I.10.9.9}, si $f$ es un morfismo
finito, también lo es su extensión $\widehat f$ a las
completaciones.
\end{corollary}

\begin{corollary}[4.8.5]
\label{III.4.8.5}
\label{III.4.8.5-fr}
Sea $\mathfrak X$ un preesquema formal finito sobre
$\mathfrak Y$, y sea
$f:\mathfrak X\to\mathfrak Y$ el morfismo estructural. Para
todo subconjunto abierto $U\subset\mathfrak Y$, $f^{-1}(U)$ es finito sobre $U$.
\end{corollary}

\begin{proposition}[4.8.6]
\label{III.4.8.6}
\label{III.4.8.6-fr}
Si $f:\mathfrak X\to\mathfrak Y$ es un morfismo finito de preesquemas
formales localmente noetherianos, entonces
$f_*(\mathcal O_{\mathfrak X})$ es una
$\mathcal O_{\mathfrak Y}$-álgebra coherente.
\end{proposition}

\begin{proof}
Se puede considerar $f$ como el límite inductivo de un sistema inductivo
$(f_n)$ de morfismos $f_{\boldsymbol{n}}:X_n\to Y_{\boldsymbol{n}}$. Demostremos que los
$f_n$ son finitos y $f_*(\mathcal O_{\mathfrak X})$ es
isomorfo al límite proyectivo de los
$(f_n)_*(\mathcal O_{X_n})$, lo que establecerá nuestra afirmación (\textbf{I}, 10.10.5).
Basta limitarse al caso en que $\mathfrak Y=\operatorname{Spf}(B)$,
$\mathfrak X=\operatorname{Spf}(A)$, y observar que si $\mathfrak R$ es un ideal de definición
de $B$ y $A$ es un $B$-módulo de tipo finito, entonces
$A/\mathfrak R^{n+1}A$ es un módulo de tipo finito sobre
$B/\mathfrak R^{n+1}B$, y $A$ es el límite proyectivo de los
$A/\mathfrak R^{n+1}A$.
\end{proof}

Recíprocamente:

\begin{proposition}[4.8.7]
\label{III.4.8.7}
\label{III.4.8.7-fr}
Sea $\mathfrak Y$ un preesquema formal localmente noetheriano y
$\mathcal A$ una $\mathcal O_{\mathfrak Y}$-álgebra coherente. Existe
un preesquema formal $\mathfrak X$ finito sobre $\mathfrak Y$,
único salvo $\mathfrak Y$-isomorfismo único, tal que
$f_*(\mathcal O_{\mathfrak X})=\mathcal A$,
donde $f:\mathfrak X\to\mathfrak Y$ es el morfismo estructural.
\end{proposition}

\begin{proof}
\oldpage[III]{148}
Sea $\mathcal K$ un ideal de definición de $\mathfrak Y$, y pongamos
$Y_n=(\mathfrak Y,\mathcal O_{\mathfrak Y}/\mathcal K^{n+1})$ y
$\mathcal A_n=\mathcal A/\mathcal K^{n+1}\mathcal A$;
está claro que $\mathcal A_n$ es una $\mathcal O_{Y_n}$-álgebra finita y, por tanto,
define un $Y_n$-preesquema finito
$X_n=\operatorname{Spec}(\mathcal A_n)$
\sref[II]{II.6.1.3}. Para $m\leq n$, el homomorfismo canónico
sobreyectivo
$h_{mn}:\mathcal A_n\to\mathcal A_m$
define un morfismo $u_{mn}:X_m\to X_n$ para el cual el diagrama
\[
\xymatrix{
X_n\ar[d]_{f_n} & X_m\ar[l]_{u_{mn}}\ar[d]^{f_m}\\
Y_n & Y_m\ar[l]
}
\]
es conmutativo, donde $f_n$ denota el morfismo estructural; además,
identifica $X_m$ con
$X_n\times_{Y_n}Y_m$, como se sigue de inmediato de
\sref[II]{II.1.4.6}. El preesquema formal $X$ obtenido
como límite inductivo del sistema inductivo $(X_n)$ es, por tanto,
localmente noetheriano, y su morfismo estructural
$f:X\to Y$,
el límite inductivo del sistema $(f_n)$, sea finito (\hyperref[III.4.8.1-fr]{4.8.1} y \textbf{II}, 10.12.3.1); se ha visto además en la demostración de (\hyperref[III.4.8.6-fr]{4.8.6}) que $f_*(\mathcal O_X)$ es el límite proyectivo de las
$\mathcal A_n$, y es, por tanto, igual a $\mathcal A$
\sref[I]{I.10.10.6}. La afirmación de unicidad se sigue del
resultado más general siguiente.
\end{proof}

\begin{proposition}[4.8.8]
\label{III.4.8.8}
\label{III.4.8.8-fr}
Sea $\mathfrak Y$ un preesquema formal localmente noetheriano, y
sean $\mathfrak X$ y $\mathfrak X'$
$\mathfrak Y$-preesquemas formales finitos sobre $\mathfrak Y$, $f:\mathfrak X\to\mathfrak Y$,
$f':\mathfrak X'\to\mathfrak Y$ los morfismos estructurales. Existe una
biyección canónica de
$\operatorname{Hom}_{\mathfrak Y}(\mathfrak X,\mathfrak X')$
sobre
$\operatorname{Hom}_{\mathcal O_{\mathfrak Y}}
\bigl(f'_*(\mathcal O_{\mathfrak X'}),f_*(\mathcal O_{\mathfrak X})\bigr)$\footnotemark.
\footnotetext{La última expresión designa el conjunto de homomorfismos de
$\mathcal O_{\mathfrak Y}$-\emph{álgebras}
$f'_*(\mathcal O_{\mathfrak X'})\to f_*(\mathcal O_{\mathfrak X})$.}
\end{proposition}

\begin{proof}
La aplicación $h\to\mathcal A(h)$ es la misma que en (\textbf{II}, 1.1.2), y para ver que es biyectiva, se reduce inmediatamente al caso en que $\mathfrak Y=\operatorname{Spf}(B)$ es un esquema formal afín noetheriano.
Pero entonces $\mathfrak X=\operatorname{Spf}(A)$,
$\mathfrak X'=\operatorname{Spf}(A')$, donde $A$ y $A'$ son $B$-álgebras
finitas y $f_*(\mathcal O_{\mathfrak X})=A^\Delta$,
$f'_*(\mathcal O_{\mathfrak X'})=A'^\Delta$. La conclusión resulta entonces de
la correspondencia biunívoca, por una parte entre los $\mathfrak Y$-morfismos
$\mathfrak X\to\mathfrak X'$ y homomorfismos continuos de
$B$-homomorfismos (necesariamente continuos) $A'\to A$ que son homomorfismos de álgebras (\textbf{I}, 10.2.2), y por otra parte entre los homomorfismos de $B$-módulos
$A'\to A$ y homomorfismos de
$\mathcal O_{\mathfrak Y}$-módulos
$A'^\Delta\to A^\Delta$ \sref[I]{I.10.10.2.3}.
\end{proof}

\begin{corollary}[4.8.9]
\label{III.4.8.9}
\label{III.4.8.9-fr}
Bajo la biyección canónica de \sref{III.4.8.8}, las inmersiones
cerradas $\mathfrak X\to\mathfrak X'$ corresponden a
homomorfismos sobreyectivos de $\mathcal O_{\mathfrak Y}$-álgebras
$f'_*(\mathcal O_{\mathfrak X'})\to f_*(\mathcal O_{\mathfrak X})$.
\end{corollary}

\begin{proof}
La cuestión es local sobre $\mathfrak Y$ y, por tanto, se reduce a la
definición de las inmersiones cerradas de preesquemas formales localmente
noetherianos \sref[I]{I.10.14.2}.
\end{proof}

\begin{corollary}[4.8.10]
\label{III.4.8.10}
\label{III.4.8.10-fr}
Con la notación y las hipótesis de \sref{III.4.8.1}, un morfismo
ádico $f$ es una inmersión cerrada si y solo si $f_0$ es una
inmersión cerrada de preesquemas ordinarios.
\end{corollary}

\begin{proof}
Esto se sigue de inmediato de \sref{III.4.8.9} y del criterio de
sobreyectividad para un homomorfismo de
$\mathcal O_{\mathfrak Y}$-módulos coherentes \sref[I]{I.10.11.5}.
\end{proof}

\begin{proposition}[4.8.11]
\label{III.4.8.11}
\label{III.4.8.11-fr}
Un morfismo $f:\mathfrak X\to\mathfrak Y$ de preesquemas formales
localmente noetherianos
\oldpage[III]{149}
es finito si y solo si es propio y todas sus fibras
$f^{-1}(y)$, para $y\in\mathfrak Y$, son finitas.
\end{proposition}

\begin{proof}
Por las definiciones \sref{III.3.4.1} y \sref{III.4.8.2}, la
afirmación se reduce inmediatamente a la misma afirmación para $f_0$
(con la notación de \sref{III.4.8.1}), que es precisamente
\sref{III.4.4.2}.
\end{proof}

% END tree/repo/editions/III1_DIPLOMATIC/ega3/ega3-4.tex


% BEGIN tree/repo/editions/III1_DIPLOMATIC/ega3/ega3-5.tex
\section{Un teorema de existencia para los haces algebraicos coherentes}
\label{section:III.5}
\label{section:III.5-fr}


\subsection{Enunciado del teorema}
\label{subsection:III.5.1}
\label{subsection:III.5.1-fr}

\begin{env}[5.1.1]
\label{III.5.1.1}
\label{III.5.1.1-fr}
Sea $A$ un anillo \emph{noetheriano ádico} y sea $\mathfrak J$ un ideal de
definición de $A$, de modo que $A$ sea separado y completo para la
topología $\mathfrak J$-ádica. Si $Y=\operatorname{Spec}(A)$, el esquema formal
afín $\operatorname{Spf}(A)$ se identifica con la \emph{compleción}
$\widehat Y$ de $Y$ a lo largo del subconjunto cerrado
$Y'=\mathrm V(\mathfrak J)$ \sref[I]{I.10.10.1}. Sea $X$ un
$Y$-preesquema ordinario de tipo finito, de morfismo estructural
$f:X\to Y$. Denotemos por $\widehat X$ la \emph{compleción} de $X$ a lo largo
del subconjunto cerrado $X'=f^{-1}(Y')$, o, equivalentemente, el
$\widehat Y$-preesquema formal $X\times_Y\widehat Y$, y por
$\widehat f:\widehat X\to\widehat Y$ la extensión de $f$ a las
compleciones. Finalmente, para todo $\mathcal O_X$-módulo coherente
$\mathcal F$, denotemos por $\widehat{\mathcal F}$ su compleción
$\mathcal F_{/X'}$, que es un
$\mathcal O_{\widehat X}$-módulo coherente.
\end{env}

\begin{proposition}[5.1.2]
\label{III.5.1.2}
\label{III.5.1.2-fr}
Con las hipótesis y la notación de \sref{III.5.1.1}, sea $\mathcal F$
un $\mathcal O_X$-módulo coherente cuyo soporte es propio sobre $Y$
\sref[II]{II.5.4.10}. Los homomorfismos canónicos
\sref{III.4.1.4}
\[
\rho_i:\mathrm H^i(X,\mathcal F)\longrightarrow
\mathrm H^i(\widehat X,\widehat{\mathcal F})
\]
son entonces isomorfismos.
\end{proposition}

\begin{proof}
Como $\mathrm H^i(X,\mathcal F)$ es un $A$-módulo de tipo finito (3.2.4), luego
idéntico a su separado completado para la topología $\mathfrak J$-preádica
($\mathbf{0}_{\mathrm I}$, 7.3.6), la proposición no es sino un caso particular de (4.1.10).
\end{proof}

Recordemos que los isomorfismos canónicos $\rho_i$ conmutan con los
homomorfismos de conexión asociados a toda sucesión exacta
$0\to\mathcal F'\to\mathcal F\to\mathcal F''\to0$
de $\mathcal O_X$-módulos coherentes
\sref[0]{0.12.1.6} \sref[I]{I.10.8.9}.

\begin{corollary}[5.1.3]
\label{III.5.1.3-fr}
Sean $\mathcal F$, $\mathcal G$ dos $\mathcal O_X$-Módulos coherentes tales que
la intersección de sus soportes sea propia sobre $Y$. Entonces el homomorfismo canónico
\begin{equation}
\operatorname{Hom}_{\mathcal O_X}(\mathcal F,\mathcal G)
\longrightarrow
\operatorname{Hom}_{\mathcal O_{\widehat X}}
(\widehat{\mathcal F},\widehat{\mathcal G})
\tag{5.1.3.1}
\label{III.5.1.3.1.eq-fr}
\end{equation}
que, a todo homomorfismo $u:\mathcal F\to\mathcal G$, hace corresponder su
completado $\widehat u:\mathcal F\to\mathcal G$, es un
isomorfismo. Además, cuando el morfismo $f$ es cerrado, para que $\widehat u$ sea
inyectivo (\textup{resp.} sobreyectivo), es necesario y suficiente que $u$ lo sea.
\end{corollary}

\begin{proof}
La primera aserción es un caso particular de (4.5.3), debido también a que el
primer miembro de (\ref{III.5.1.3.1.eq-fr}) es un $A$-módulo de tipo finito, luego
idéntico a su separado completado. Para demostrar la segunda, observemos en virtud de
(\textbf{I}, 10.8.14) que $\widehat u$ es inyectivo (resp. sobreyectivo) si y solo si existe
una vecindad de $X'$ en la cual $u$ sea inyectivo (resp. sobreyectivo).

\oldpage[III]{150}
La conclusión
resulta pues del lema siguiente:
\end{proof}

\begin{lemma}[5.1.3.1]
\label{III.5.1.3.1-fr}
Bajo las hipótesis de \textup{(\hyperref[III.5.1.1-fr]{5.1.1})}, si se supone además el
morfismo $f$ cerrado, toda vecindad de $X'$ en $X$ es idéntica a $X$.
\end{lemma}

\begin{proof}
En primer lugar, se puede reducir al caso en que $f(X)=Y$. En efecto, por hipótesis,
$f(X)$ es una parte cerrada $Z$ de $Y$; se puede además reemplazar $f$ por
$f_{\mathrm{red}}$ (\textbf{I}, 6.3.4), y suponer por consiguiente $X$ e $Y$ reducidos; se puede entonces
reemplazar $Y$ por el subpreesquema cerrado reducido de $Y$ que tiene $Z$ por espacio
subyacente (\textbf{I}, 5.2.2), pues todo ideal de $A$ es cerrado, y todo anillo cociente de
$A$ es por tanto ádico y noetheriano. Se tiene entonces $f(X')=Y'$; si $V$ es una vecindad
abierta de $X'$ en $X$, $f(X\setmin V)$ es cerrado en $Y$ por hipótesis, y no corta
a $Y'$; pero esto es imposible a menos que $X\setmin V$ sea vacío, puesto que
$\mathfrak J$ está contenido en el radical de $A$ (\textbf{I}, 1.1.15 y $\mathbf{0}_{\mathrm I}$,
7.1.10), de donde la conclusión.
\end{proof}

Cuando se limita a los $\mathcal O_X$-Módulos coherentes cuyo soporte es propio
sobre $Y$, (\hyperref[III.5.1.3-fr]{5.1.3}) puede enunciarse, en el lenguaje de las
categorías, diciendo que el funtor
$\mathcal F\leadsto\widehat{\mathcal F}$ es \emph{plenamente fiel} de la categoría de los
$\mathcal O_X$-Módulos del tipo precedente en la categoría de los
$\mathcal O_{\widehat X}$-Módulos coherentes, y establece por consiguiente una \emph{equivalencia} de
la primera de estas categorías con una subcategoría \emph{plena} de la segunda
(\textbf{0}, 8.1.6). El teorema de existencia probará que, cuando $X$ es propio sobre $Y$,
esta subcategoría es de hecho la categoría de \emph{todos} los
$\mathcal O_{\widehat X}$-Módulos coherentes. De manera precisa:

\begin{theorem}[5.1.4]
\label{III.5.1.4}
\label{III.5.1.4-fr}
Sea $A$ un anillo noetheriano ádico, sea $Y=\operatorname{Spec}(A)$, sea
$\mathfrak J$ un ideal de definición de $A$, $Y'=\mathrm V(\mathfrak J)$. Sea $f:X\to Y$ un morfismo separado de tipo finito, $X'=f^{-1}(Y')$. Sean
$\widehat Y=Y_{/Y'}=\operatorname{Spf}(A)$, $\widehat X=X_{/X'}$ los completados de
$Y$ y $X$ a lo largo de $Y'$ y $X'$, $\widehat f:\widehat X\to\widehat Y$ la extensión de $f$ a
las compleciones. Entonces el funtor
$\mathcal F\leadsto\mathcal F_{/X'}=\widehat{\mathcal F}$
es una equivalencia de la categoría de los
$\mathcal O_X$-módulos coherentes cuyo soporte es propio sobre $\operatorname{Spec}(A)$ en la
categoría de los $\mathcal O_{\widehat X}$-módulos coherentes cuyo soporte
es propio sobre $\operatorname{Spf}(A)$.
\end{theorem}

En otras palabras, teniendo en cuenta \sref{III.5.1.3}:

\begin{corollary}[5.1.5]
\label{III.5.1.5}
\label{III.5.1.5-fr}
Un $\mathcal O_{\widehat X}$-módulo es isomorfo a la compleción de
un $\mathcal O_X$-módulo coherente cuyo soporte es propio sobre
$\operatorname{Spec}(A)$ si y solo si es coherente y su soporte es propio
sobre $\operatorname{Spf}(A)$.
\end{corollary}

El caso más importante es el siguiente:

\begin{corollary}[5.1.6]
\label{III.5.1.6}
\label{III.5.1.6-fr}
Supongamos que $X$ es propio sobre $Y=\operatorname{Spec}(A)$. Entonces el funtor
$\mathcal F\leadsto\widehat{\mathcal F}$ es una equivalencia entre la
categoría de los $\mathcal O_X$-módulos coherentes y la categoría de los
$\mathcal O_{\widehat X}$-módulos coherentes.
\end{corollary}

\begin{scholium}[5.1.7]
\label{III.5.1.7}
\label{III.5.1.7-fr}
Teniendo en cuenta la caracterización de los haces coherentes sobre
preesquemas formales \sref[I]{I.10.11.3}, se ve que, bajo las
condiciones de \sref{III.5.1.1}, dar un
$\mathcal O_X$-módulo coherente cuyo soporte es propio sobre $\operatorname{Spec}(A)$ \emph{equivale} (poniendo
$Y_n=\operatorname{Spec}(A/\mathfrak J^{n+1})$ y $X_n=X\times_Y Y_n$) al dato
de un sistema proyectivo de $\mathcal O_{X_n}$-Módulos coherentes $(\mathcal F_n)$ tal
que para $m\leq n$ se tiene
$\mathcal F_m=\mathcal F_n\otimes_{Y_n}\mathcal O_{Y_m}$
(o también $\mathcal F_m=\mathcal F_n/\mathfrak J^{m+1}\mathcal F_n$)
y el soporte de $\mathcal F_0$ sea un subconjunto de $X_0$
propio sobre $Y_0$. Por medio de (\textbf{I}, 10.11.4), se interpretan análogamente los homomorfismos de $\mathcal O_X$-Módulos coherentes como homomorfismos de sistemas proyectivos de $\mathcal O_{X_n}$-módulos coherentes.

\oldpage[III]{151}

En todas las aplicaciones conocidas, $A$ es de hecho un anillo
noetheriano ádico \emph{local}. Así, los $Y_n$ son espectros de anillos
\emph{artinianos locales}, y los resultados de esta subsección y de las
anteriores reducen, en gran medida, la geometría algebraica sobre un
anillo noetheriano ádico local a la geometría algebraica sobre anillos
artinianos locales.
\end{scholium}

\begin{corollary}[5.1.8]
\label{III.5.1.8}
\label{III.5.1.8-fr}
Bajo las hipótesis de \sref{III.5.1.4}, la aplicación
$Z\leadsto\widehat Z=Z_{/(Z\cap X')}$
es una biyección del conjunto de los subpreesquemas cerrados $Z$ de $X$
que son propios sobre $Y$ sobre el conjunto de los subpreesquemas
formales cerrados de $\widehat X$ que son propios sobre
$\widehat Y$.
\end{corollary}

\begin{proof}
Un subpreesquema formal cerrado de $\widehat X$ es de la forma
$(T,(\mathcal O_{\widehat X}/\mathcal A)|T)$,
donde $\mathcal A$ es un ideal coherente de
$\mathcal O_X$ \sref[I]{I.10.14.2}. Si $T$ es propio sobre
$\widehat Y$, \sref{III.5.1.4} muestra que
$\mathcal O_{\widehat X}/\mathcal A$ es isomorfo a un
$\mathcal O_{\widehat X}$-módulo de la forma
$\widehat{\mathcal F}$, donde $\mathcal F$ es un
$\mathcal O_X$-módulo coherente cuyo soporte es propio sobre $Y$. Además,
\sref{III.5.1.3} muestra que el homomorfismo canónico
$\mathcal O_{\widehat X}\to\mathcal O_{\widehat X}/\mathcal A$
es de la forma $\widehat u$, donde
$u:\mathcal O_X\to\mathcal F$ es un homomorfismo sobreyectivo de
$\mathcal O_X$-Módulos. Luego $\mathcal F$ es de la forma
$\mathcal O_X/\mathcal N$, donde $\mathcal N$ de
$\mathcal O_X$, y
$\mathcal A=\widehat{\mathcal N}$ \sref[I]{I.10.8.8}. La conclusión se
deduce ahora de \sref[I]{I.10.14.7}.
\end{proof}


\subsection{Demostración del teorema de existencia: caso proyectivo y cuasiproyectivo}
\label{subsection:III.5.2}
\label{subsection:III.5.2-fr}

\begin{env}[5.2.1]
\label{III.5.2.1}
\label{III.5.2.1-fr}
Bajo las condiciones de \sref{III.5.1.4}, diremos
provisionalmente que un $\mathcal O_{\widehat X}$-módulo coherente es
\emph{algebraizable} si es isomorfo a la compleción
$\widehat{\mathcal F}$ de un $\mathcal O_X$-módulo coherente $\mathcal F$ cuyo
soporte es propio sobre $Y$.
\end{env}

\begin{lemma}[5.2.2]
\label{III.5.2.2}
\label{III.5.2.2-fr}
Sean $\mathcal F'$ y $\mathcal G'$ $\mathcal O_{\widehat X}$-módulos
algebraizables. Para todo homomorfismo
$u:\mathcal F'\to\mathcal G'$, los módulos
$\operatorname{Ker}(u)$, $\operatorname{Im}(u)$ y $\operatorname{Coker}(u)$ son
algebraizables.
\end{lemma}

\begin{proof}
En efecto, si $\mathcal F'=\widehat{\mathcal F}$,
$\mathcal G'=\widehat{\mathcal G}$, donde $\mathcal F$ y $\mathcal G$ son $\mathcal O_X$-módulos coherentes cuyos
soportes son propios sobre $Y$. Por \sref{III.5.1.3},
$u=\widehat v$ para un homomorfismo
$v:\mathcal F\to\mathcal G$ es un homomorfismo
(\hyperref[III.5.1.3-fr]{5.1.3}). En virtud de la exactitud del funtor
$\mathcal F\leadsto\widehat{\mathcal F}$,
$\operatorname{Ker}(\widehat v)$ es isomorfo a
$(\operatorname{Ker}(v))^{\wedge}$ y como el soporte de
$\operatorname{Ker}(v)$ está contenido en el de $\mathcal F$, se deduce que
$\operatorname{Ker}(u)$ es algebraizable. Los argumentos para $\operatorname{Im}(u)$ y
$\operatorname{Coker}(u)$ son los mismos.
\end{proof}

\begin{proposition}[5.2.3]
\label{III.5.2.3-fr}
Sean $A$ un anillo ádico noetheriano, $\mathfrak J$ un ideal de definición de
$A$, $\mathfrak Y=\operatorname{Spf}(A)$,
$f:\mathfrak X\to\mathfrak Y$ un morfismo propio de preesquemas formales. Se pone
$Y_k=\operatorname{Spec}(A/\mathfrak J^{k+1})$,
$X_k=\mathfrak X\times_{\mathfrak Y}Y_k$, y para todo
$\mathcal O_{\mathfrak X}$-Módulo $\mathcal F$,
$\mathcal F_k=\mathcal F\otimes_{\mathcal O_{\mathfrak X}}\mathcal O_{X_n}
=\mathcal F/\mathfrak J^{k+1}\mathcal F$.
Sea $\mathcal L$ un $\mathcal O_{\mathfrak X}$-Módulo invertible, y supongamos que
$\mathcal L_0=\mathcal L/\mathfrak J\mathcal L$ es un
$\mathcal O_{X_0}$-Módulo amplio; para todo $\mathcal O_{\mathfrak X}$-Módulo
$\mathcal F$ y todo entero $n$, pongamos
$\mathcal F(n)=\mathcal F\otimes\mathcal L^{\otimes n}$. Entonces, para todo
$\mathcal O_{\mathfrak X}$-Módulo coherente $\mathcal F$, existe un entero $n_0$
tal que, para todo $n\geq n_0$, se verifican las propiedades siguientes:
\begin{enumerate}
\item[\textup{(i)}] El homomorfismo canónico
$\mathrm H^0(\mathfrak X,\mathcal F(n))\to
\mathrm H^0(\mathfrak X,\mathcal F_k(n))$ es sobreyectivo para todo $k\geq0$.
\item[\textup{(ii)}] Se tiene $\mathrm H^q(\mathfrak X,\mathcal F(n))=0$ para todo $q>0$.
\end{enumerate}
\end{proposition}

\begin{proof}
Se sabe que los espacios subyacentes a $\mathfrak X$ y a $X_0$ son los mismos;
los haces
$\mathcal M_k=\mathfrak J^k\mathcal F/\mathfrak J^{k+1}\mathcal F$,
al estar anulados por $\mathfrak J$, pueden considerarse como
$\mathcal O_{X_0}$-Módulos coherentes ($\mathbf{0}_{\mathrm I}$, 5.3.10); además, si se
pone $\mathcal M_k(n)=\mathcal M_k\otimes_{\mathcal O_{X_0}}
\mathcal L_0^{\otimes n}$, se ve inmediatamente que
$\mathcal M_k(n)=\mathfrak J^k\mathcal F(n)/
\mathfrak J^{k+1}\mathcal F(n)$.
Observemos que, puesto que $\mathcal L_0$ es amplio para $f_0:X_0\to Y_0$, y
\oldpage[III]{152}
$f_0$ es propio, se concluye que $f_0$ es \emph{proyectivo} (\textbf{II}, 5.5.4). Sea $S$ la $A$-álgebra
graduada $\bigoplus_{k\geq0}\mathfrak J^k/\mathfrak J^{k+1}$ asociada a la filtración $\mathfrak J$-ádica de $A$, que es
de tipo finito puesto que $A$ es noetheriano; si se pone
$\mathcal S'=f_0^*(\widetilde{\mathbf S})$,
$\mathcal S'$ es una
$\mathcal O_{X_0}$-Álgebra cuasicoherente de tipo finito, y
$\mathcal M=\bigoplus_{k\geq0}\mathcal M_k$ un
$\mathcal S'$-Módulo graduado cuasicoherente y de tipo finito (puesto que $\mathcal F_0$
es coherente y genera el $\mathcal S'$-Módulo $\mathcal M$). Se está pues en las
condiciones de aplicación del teorema (2.4.1, (ii)), y se concluye que existe
$n_0$ tal que, para $n\geq n_0$ y para \emph{todo} $k$, se tiene
\begin{equation}
\mathrm H^q(X_0,\mathcal M_k(n))=0\qquad\text{\emph{para todo} }q>0.
\tag{5.2.3.1}
\label{III.5.2.3.1-fr}
\end{equation}
Se tiene pues también $\mathrm H^q(\mathfrak X,\mathcal M_k(n))=0$ para $q>0$ y
$n\geq n_0$, considerando esta vez $\mathcal M_k(n)$ como
$\mathcal O_{\mathfrak X}$-Módulo. Aplicando la sucesión exacta de cohomología a
\[
0\to
\mathfrak J^k\mathcal F(n)/\mathfrak J^{k+1}\mathcal F(n)
\to
\mathfrak J^h\mathcal F(n)/\mathfrak J^{k+1}\mathcal F(n)
\to
\mathfrak J^h\mathcal F(n)/\mathfrak J^k\mathcal F(n)
\to0
\]
se deduce ante todo que para $0\leq h<k$, $n\geq n_0$ y $q>0$ se tiene, por
inducción sobre $h-k$,
\begin{equation}
\mathrm H^q(\mathfrak X,
\mathfrak J^h\mathcal F(n)/\mathfrak J^k\mathcal F(n))=0
\tag{5.2.3.2}
\label{III.5.2.3.2-fr}
\end{equation}
y en particular para $h=0$
\begin{equation}
\mathrm H^q(\mathfrak X,\mathcal F_k(n))=0
\qquad\text{para }n\geq n_0,\ k\geq0\text{ y }q>0.
\tag{5.2.3.3}
\label{III.5.2.3.3-fr}
\end{equation}
Otra porción de la sucesión exacta de cohomología, para $h=0$, da la sucesión
exacta
\begin{equation}
\mathrm H^0(\mathfrak X,\mathcal F_{k+1}(n))\to
\mathrm H^0(\mathfrak X,\mathcal F_k(n))\to
\mathrm H^1(\mathfrak X,
\mathfrak J^k\mathcal F(n)/\mathfrak J^{k+1}\mathcal F(n))=0
\tag{5.2.3.4}
\label{III.5.2.3.4-fr}
\end{equation}
de donde se deduce que para $h\leq k$ la aplicación canónica
\begin{equation}
\mathrm H^0(\mathfrak X,\mathcal F_k(n))\longrightarrow
\mathrm H^0(\mathfrak X,\mathcal F_h(n))
\tag{5.2.3.5}
\label{III.5.2.3.5-fr}
\end{equation}
es \emph{sobreyectiva}. Para \emph{todo} $q$, el sistema proyectivo
$(\mathrm H^q(\mathfrak X,\mathcal F_k(n))\quad(k\geq0)$ verifica pues la condición
(ML) para $n\geq n_0$. Por otra parte, todo abierto formal afín $U$ de $\mathfrak X$
es también un abierto afín en cada uno de los $X_k$ (\textbf{I}, 10.5.2), luego se tiene
$\mathrm H^q(U,\mathcal F_k(n))=0$ para todo $q>0$ (1.3.1), y
$\mathrm H^0(U,\mathcal F_k(n))\to\mathrm H^0(U,\mathcal F_h(n))$ es sobreyectiva
para $h\leq k$ (\textbf{I}, 1.3.9). Las condiciones de aplicación de (\textbf{0}, 13.3.1) se
cumplen por tanto, y se concluye que, para $n\geq n_0$:
\begin{enumerate}
\item[$1^\circ$] Para todo $q>0$,
$\mathrm H^q(\mathfrak X,\mathcal F(n))\longrightarrow
\varprojlim_k\mathrm H^q(\mathfrak X,\mathcal F_k(n))$
es biyectivo, luego, en virtud de (\ref{III.5.2.3.3-fr}),
$\mathrm H^q(\mathfrak X,\mathcal F(n))=0$.
\item[$2^\circ$] El homomorfismo
$\mathrm H^0(\mathfrak X,\mathcal F(n))\longrightarrow
\varprojlim_k\mathrm H^0(\mathfrak X,\mathcal F_k(n))$
es biyectivo; por otra parte, como los homomorfismos
(\ref{III.5.2.3.5-fr}) son sobreyectivos, lo mismo ocurre con cada uno de los
homomorfismos
\[
\varprojlim_k\mathrm H^0(\mathfrak X,\mathcal F_k(n))
\longrightarrow\mathrm H^0(\mathfrak X,\mathcal F_h(n))
\]
\end{enumerate}
lo que termina la demostración.
\end{proof}

\begin{corollary}[5.2.4]
\label{III.5.2.4-fr}
Siendo las hipótesis las de \textup{(\hyperref[III.5.2.3-fr]{5.2.3})}, para todo
$\mathcal O_{\mathfrak X}$-Módulo coherente $\mathcal F$ existe un entero $N$
tal que para $n\geq N$, $\mathcal F(n)$ está generado por sus secciones sobre
\oldpage[III]{153}
$\mathfrak X$; en otros términos, $\mathcal F$ es isomorfo al cociente de un
$\mathcal O_{\mathfrak X}$-Módulo de la forma
$(\mathcal O_{\mathfrak X}(\setmin n))^k$.
\end{corollary}

\begin{proof}
Como $X_0$ es noetheriano, resulta de la hipótesis sobre $\mathcal L_0$ y de
(\textbf{II}, 4.5.5) que existe $n_0$ tal que, para $n\geq n_0$, $\mathcal F_0(n)$ está
generado por sus secciones sobre $\mathfrak X$; por otra parte, se puede suponer
$n_0$ suficientemente grande para que el homomorfismo
$\Gamma(\mathfrak X,\mathcal F(n))\longrightarrow
\Gamma(\mathfrak X,\mathcal F_0(n))$
sea sobreyectivo para $n\geq n_0$ (\hyperref[III.5.2.3-fr]{5.2.3}). Existe pues
un número finito de secciones $s_i\in\Gamma(\mathfrak X,\mathcal F(n))$ cuyas
imágenes en $\Gamma(\mathfrak X,\mathcal F_0(n))$ generan $\mathcal F_0(n)$
($\mathbf{0}_{\mathrm I}$, 5.2.3). Como $\mathfrak J$ está contenido en el ideal maximal del
anillo local en todo punto de $\mathfrak X$, resulta del lema de Nakayama,
aplicado a estos anillos locales, que los $s_i$ generan $\mathcal F(n)$
($\mathbf{0}_{\mathrm I}$, 5.1.1).
\end{proof}

\begin{env}[5.2.5]
\label{III.5.2.5}
\label{III.5.2.5-fr}
\textbf{Demostración del teorema de existencia: caso proyectivo.}
Con la notación de \sref{III.5.1.4}, supongamos que $f$ es
\emph{proyectivo}, de modo que existe un $\mathcal O_X$-módulo amplio $\mathcal L$
\sref[I]{I.5.5.4}. Por definición,
$X_n=\widehat X\times_{\widehat Y}Y_n$
es el subpreesquema cerrado
$X_n\times_Y Y_n=f^{-1}(Y_n)$ de $X$; si
$\mathcal L'$ es el
completado $\mathcal L\otimes_{\mathcal O_X}\mathcal O_{\widehat X}$ de $\mathcal L$, entonces
$\mathcal L'_0=\mathcal L/\mathfrak J\mathcal L$, considerado como
$\mathcal O_{X_0}$-Módulo; se sabe que $\mathcal L'_0$ es \emph{amplio}
(\textbf{II}, 4.6.13, (i \emph{bis})). Se puede pues aplicar a $\mathcal L'$ y a todo
$\mathcal O_{\widehat X}$-módulo coherente $\mathcal F$. Se deduce que
$\mathcal F$ es isomorfo a un cociente de
$\mathcal G=(\mathcal L'^{\otimes(-n)})^k$
para enteros adecuados $n>0$ y $k>0$. El módulo $\mathcal G$ es la
compleción de $(\mathcal L^{\otimes(-n)})^k$ (\textbf{I}, 10.8.10), luego es algebraizable. Consideremos después el homomorfismo canónico sobreyectivo $u:\mathcal G\to\mathcal F$, y sea $\mathcal H=\operatorname{Ker}(u)$; este es un
$\mathcal O_{\widehat X}$-Módulo coherente ($\mathbf{0}_{\mathrm I}$, 5.3.4). Se ve del
mismo modo que existe un $\mathcal O_{\widehat X}$-módulo algebraizable $\mathcal K$ y un homomorfismo
$v:\mathcal K\to\mathcal G$ tal que $\mathcal H=\operatorname{Im}(v)$. Así,
$\mathcal F=\operatorname{Coker}(v)$, y $\mathcal F$ es algebraizable por
\sref{III.5.2.2}.
\end{env}

\begin{env}[5.2.6]
\label{III.5.2.6}
\label{III.5.2.6-fr}
\textbf{Demostración del teorema de existencia: caso cuasiproyectivo.}
Conservemos la notación de \sref{III.5.1.4} y supongamos ahora que $f$
es \emph{cuasiproyectivo}. Existe un morfismo proyectivo
$g:Z\to Y$ tal que $X$ se identifica con el $Y$-preesquema
inducido sobre un subconjunto abierto de $Z$ \sref[II]{II.5.3.2}. Si
$Z'=g^{\phantom{-}1}(Y')$, entonces $X'=X\cap Z'$. Por consiguiente, la compleción
$\widehat X=X_{/X'}$ se identifica con el preesquema formal
inducido por $\widehat Z=Z_{/Z'}$ sobre el subconjunto abierto $X\cap Z'$ de
$\widehat Z$ \sref[I]{I.10.8.5}.

Sea $\mathcal F'$ un $\mathcal O_{\widehat X}$-módulo coherente cuyo
soporte $T'$ sea \emph{propio sobre} $\widehat Y$. Por definición, existe
un subpreesquema cerrado de $X'$ de espacio subyacente
$T'\subset X'$ tal que la restricción $T'\to Y'$ de $f$ es
propia. Por tanto, $T'$ es propio \emph{sobre} $Y$ y, en consecuencia, \emph{cerrado en}
$Z'$ \sref[II]{II.5.4.10}. Se deduce que $\mathcal F'$ es inducido
sobre $\widehat X$ por el $\mathcal O_Z$-módulo $\mathcal G'$
obtenido pegando $\mathcal F'$ (definido sobre el abierto $\widehat X$ de $\widehat Z$) y
del haz $0$ sobre el abierto $\widehat Z\setmin T'$ de $\widehat Z$, coincidiendo estos dos haces
en el abierto intersección $\widehat X\setmin T'$.

El módulo $\mathcal G'$ es coherente. Por \sref{III.5.2.5}, existe un
$\mathcal O_Z$-módulo coherente $\mathcal G$ tal que
$\mathcal G'=\widehat{\mathcal G}$. Sea $T$ el soporte de $\mathcal G$;
entonces $T'=T\cap Z'$ \sref[I]{I.10.8.12}. Si $h$ es la restricción
de $g$ al subpreesquema cerrado reducido de $Z$ cuyo espacio subyacente
es $T$, entonces
$T'=h^{-1}(Y')=T\cap g^{-1}(Y')$,
Así, $X\cap T$ es un subconjunto abierto de $T$ que contiene a $T'$.

\oldpage[III]{154}
Como $h$ es propio \sref[II]{II.5.4.2}, y por tanto cerrado,
\sref{III.5.1.3.1} da $X\cap T=T$. Así, $T\subset X$; y
como $T$ es cerrado en $Z$, $T$ es propio sobre $Y$. Si $\mathcal F$ es
el haz inducido sobre $X$ por $\mathcal G$, su compleción
$\widehat{\mathcal F}$ es inducida sobre $\widehat X$ por
$\widehat{\mathcal G}$ \sref[I]{I.10.8.4} y, por tanto, es igual a
$\mathcal F'$. Esto termina la demostración.
\end{env}


\subsection{Demostración del teorema de existencia: caso general}
\label{subsection:III.5.3}
\label{subsection:III.5.3-fr}

\begin{lemma}[5.3.1]
\label{III.5.3.1}
\label{III.5.3.1-fr}
Bajo las condiciones de \sref{III.5.1.4}, sea
$0\to\mathcal H\to\mathcal F\to\mathcal G\to0$
una sucesión exacta de $\mathcal O_{\widehat X}$-módulos coherentes.
Si $\mathcal G$ y $\mathcal H$ son algebraizables, entonces $\mathcal F$ es
algebraizable.
\end{lemma}

\begin{proof}
En efecto, supongamos que $\mathcal G=\widehat{\mathcal B}$,
$\mathcal H=\widehat{\mathcal C}$, $\mathcal B$ y $\mathcal C$ son $\mathcal O_X$-módulos coherentes cuyos
soportes son propios sobre $Y$. El módulo $\mathcal F$ define canónicamente
un elemento del $A$-módulo
$\operatorname{Ext}^1_{\mathcal O_{\widehat X}}
(\widehat X;\widehat{\mathcal B},\widehat{\mathcal C})$
\sref[0]{0.12.3.2}. Las hipótesis implican que este $A$-módulo es
canónicamente isomorfo a
$\operatorname{Ext}^1_{\mathcal O_X}(X;\mathcal B,\mathcal C)$
\sref{III.4.5.2}. Existe, por tanto, una sucesión exacta
$0\to\mathcal C\to\mathcal A\to\mathcal B\to0$
de $\mathcal O_X$-Módulos
coherentes tal que la imagen canónica del elemento de
$\operatorname{Ext}^1_{\mathcal O_X}(X;\mathcal B,\mathcal C)$ correspondiente a
$\mathcal A$ es el elemento de
$\operatorname{Ext}^1_{\mathcal O_{\widehat X}}
(\widehat X;\widehat{\mathcal B},\widehat{\mathcal C})$ correspondiente a
$\mathcal F$. Por definición, teniendo en cuenta
\sref[I]{I.10.8.8}, (ii), esto significa que $\mathcal F$ es isomorfo a
$\widehat{\mathcal A}$. El lema se deduce, pues
$\operatorname{Supp}(\mathcal A)$ está contenido en la unión de
$\operatorname{Supp}(\mathcal B)$ y $\operatorname{Supp}(\mathcal C)$, luego es
propio sobre $Y$.
\end{proof}

\begin{corollary}[5.3.2]
\label{III.5.3.2}
\label{III.5.3.2-fr}
Bajo las condiciones de \sref{III.5.1.1}, sea
$u:\mathcal F\to\mathcal G$ un homomorfismo de
$\mathcal O_{\widehat X}$-módulos coherentes. Si $\mathcal G$, $\operatorname{Ker}(u)$ y
$\operatorname{Coker}(u)$ son algebraizables, entonces $\mathcal F$ es algebraizable.
\end{corollary}

\begin{proof}
El Lema \sref{III.5.2.2}, aplicado a
$\mathcal G\to\operatorname{Coker}(u)$, muestra que $\operatorname{Im}(u)$ es algebraizable. Basta entonces
aplicar el Lema \sref{III.5.3.1} a la sucesión exacta
$0\to\operatorname{Ker}(u)\to\mathcal F\to\operatorname{Im}(u)\to0$.
\end{proof}

\begin{lemma}[5.3.3]
\label{III.5.3.3}
\label{III.5.3.3-fr}
Bajo las condiciones de \sref{III.5.1.1}, sea $h:Z\to Y$ un
morfismo de tipo finito, sea $\widehat Z$ la compleción de $Z$
a lo largo de $Z'=h^{-1}(Y')$, sea $g:Z\to X$ un $Y$-morfismo propio,
y sea $\widehat g:\widehat Z\to\widehat X$ su extensión a
las compleciones. Para todo $\mathcal O_{\widehat Z}$-módulo
algebraizable $\mathcal F'$, el módulo
$\widehat g_*(\mathcal F')$ es un
$\mathcal O_{\widehat X}$-módulo algebraizable.
\end{lemma}

\begin{proof}
Si $\mathcal F$ es un $\mathcal O_Z$-módulo coherente tal que
$\mathcal F'=\widehat{\mathcal F}$, el primer teorema de comparación
\sref{III.4.1.5} identifica
$\widehat g_*(\widehat{\mathcal F})$ con la compleción de
$g_*(\mathcal F)$.
\end{proof}

\begin{lemma}[5.3.4]
\label{III.5.3.4-fr}
Sean $X$ un esquema (usual) noetheriano, $X'$ una parte cerrada de $X$,
$f:Z\to X$ un morfismo propio, $Z'=f^{-1}(X')$,
$\widehat X=X_{/X'}$, $\widehat Z=Z_{/Z'}$,
$\widehat f:\widehat Z\to\widehat X$ la prolongación de $f$ a los completados. Sea
$\mathcal M$ un Ideal coherente de $\mathcal O_X$ tal que, si
$U=X\setmin\operatorname{Supp}(\mathcal O_X/\mathcal M)$, la restricción
$f^{-1}(U)\to U$ de $f$ es un isomorfismo. Entonces, para todo
$\mathcal O_{\widehat X}$-Módulo coherente $\mathcal F$ existe un entero $n>0$
tal que el núcleo y el conúcleo del homomorfismo canónico
$\rho_{\mathcal F}:\mathcal F\longrightarrow
\widehat f_*\bigl(\widehat f^{*}(\mathcal F)\bigr)$
\textup{($\mathbf{0}_{\mathrm I}$, 4.4.3)} están anulados por $\widehat{\mathcal M}^{n}$.
\end{lemma}

\begin{proof}
Se puede limitar al caso en que $X=\operatorname{Spec}(B)$, donde $B$ es un anillo
noetheriano, luego $X'=\mathrm V(\mathfrak K)$, donde $\mathfrak K$ es un ideal de
$B$. Vamos a ver que se puede reducir al caso en que $B$ es un anillo \emph{ádico}
noetheriano y $\mathfrak K$ un ideal de definición de $B$. Sea en efecto $B_1$ el
separado completado de $B$ para la topología $\mathfrak K$-preádica; si
$\mathfrak K_1=\mathfrak K B_1$, $B_1$ es pues un
\oldpage[III]{155}
anillo ádico noetheriano del cual $\mathfrak K_1$ es un ideal de definición. Pongamos
$X_1=\operatorname{Spec}(B_1)$ y sea $h:X_1\to X$ el morfismo correspondiente al
homomorfismo canónico $B\to B_1$; si $X'_1=h^{-1}(X')$, se tiene pues
$X'_1=\mathrm V(\mathfrak K_1)$. Pongamos finalmente
$Z_1=Z\times_X X_1=Z_{(X_1)}$, $f_1=f_{(X_1)}:Z_1\to X_1$, que es un morfismo
propio (\textbf{II}, 5.4.2), y designemos por $\widehat X_1$ el completado de $X_1$ a lo largo
de $X'_1$, por
$\widehat Z_1=Z_1\times_{X_1}\widehat X_1$ el completado de $Z_1$ a lo largo de
$Z'_1=f_1^{-1}(X'_1)$, por $\widehat f_1$ la prolongación de $f_1$ a los completados.
Es inmediato que la prolongación $\widehat h:\widehat X_1\to\widehat X$ de $h$
a los completados es un isomorfismo, correspondiente a la aplicación identidad de $B_1$
(\textbf{I}, 10.9.1); se concluye que el homomorfismo correspondiente
$\widehat Z_1\to\widehat Z$ es también un isomorfismo, identificando estos isomorfismos
$\widehat f_1$ y $\widehat f$. Finalmente,
$\mathcal M_1=h^*(\mathcal M)$ es un Ideal coherente de $\mathcal O_{X_1}$ y
$\operatorname{Supp}(\mathcal O_{X_1}/\mathcal M_1)
=h^{-1}(\operatorname{Supp}(\mathcal O_X/\mathcal M))$
(\textbf{I}, 9.1.13), luego, si
$U_1=X_1\setmin\operatorname{Supp}(\mathcal O_{X_1}/\mathcal M_1)$, se tiene
$U_1=h^{-1}(U)$, de donde resulta inmediatamente que la restricción
$f_1^{-1}(U_1)\to U_1$ de $f_1$ es un isomorfismo (\textbf{I}, 3.2.7); además, los
completados $\widehat{\mathcal M}$ y $\widehat{\mathcal M}_1$ se identifican por
$\widehat h$ (\textbf{I}, 10.9.5). Todas las hipótesis de
(\hyperref[III.5.3.4-fr]{5.3.4}) se cumplen pues para $X_1$, $X'_1$, $f_1$ y
$\mathcal M_1$, y se puede por tanto suponer en adelante $B$ ádico noetheriano y
$\mathfrak K$ un ideal de definición de $B$. Se tiene entonces
$\widehat X=\operatorname{Spf}(B)$, y $\mathcal F=N^\Delta$, donde $N$ es un
$B$-módulo de tipo finito, de donde $\mathcal F=\widehat{\mathcal G}$, donde $\mathcal G$
es el $\mathcal O_X$-Módulo coherente $\widetilde N$ (\textbf{I}, 10.10.5), y por tanto
$\widehat f^{*}(\mathcal F)=(f^*(\mathcal G))^{\widehat{\,}}$
(\textbf{I}, 10.9.5). Además, en virtud del primer teorema de comparación (4.1.5),
$\widehat f_*((f^*(\mathcal G))^{\widehat{\,}})$ se identifica canónicamente con
$(f_*(f^*(\mathcal G)))^{\widehat{\,}}$,
y el homomorfismo canónico $\rho_{\mathcal F}$ no es otro que
$\widehat{\rho}_{\mathcal G}$ en virtud de
(\hyperref[III.5.1.3-fr]{5.1.3}). Ahora bien, el núcleo $\mathcal P$ y el conúcleo
$\mathcal R$ de
$\rho_{\mathcal G}:\mathcal G\longrightarrow f_*\bigl(f^*(\mathcal G)\bigr)$
son $\mathcal O_X$-Módulos coherentes, y por hipótesis sus restricciones a
$U$ son evidentemente nulas. Existe por tanto un entero $n>0$ tal que
$\mathcal M^n\mathcal P=\mathcal M^n\mathcal R=0$ (\textbf{I}, 9.3.4); se concluye que
$\widehat{\mathcal M}^{n}\widehat{\mathcal P}=
\widehat{\mathcal M}^{n}\widehat{\mathcal R}=0$
(\textbf{I}, 10.8.8 y 10.8.10).
\end{proof}

\par\medskip\noindent\textbf{5.3.5.}\ \label{III.5.3.5-fr}
\textbf{Fin de la demostración del teorema de existencia.}

Siendo las hipótesis las de (\hyperref[III.5.1.4-fr]{5.1.4}), vamos a
utilizar el principio de inducción noetheriana ($\mathbf{0}_{\mathrm I}$, 2.2.2), suponiendo
pues el teorema verdadero para todo subpreesquema cerrado $T$ de $X$ cuyo espacio subyacente es distinto
de $X$; aquí $\widehat T$ denota la compleción de $T$ a lo largo de
$T\cap X'$. Se puede suponer que $X$ no es vacío.

Como $f$ es separado y de tipo finito, el lema de Chow
\sref[II]{II.5.6.1} proporciona un $Y$-esquema $Z$ y un $Y$-morfismo
$g:Z\to X$ tales que el morfismo estructural $h:Z\to Y$ es
cuasiproyectivo, $g$ es proyectivo y sobreyectivo, y existe un
subconjunto abierto no vacío $U$ de $X$ para el cual
$g^{-1}(U)\to U$ es un isomorfismo. Sea $\mathcal M$ un ideal coherente
de $\mathcal O_X$ que define un subpreesquema cerrado de espacio subyacente
$X\setmin U$ \sref[I]{I.5.2.2}, y sea $\mathcal F$ un
$\mathcal O_{\widehat X}$-módulo coherente cuyo soporte $E$ es propio sobre $Y$.
Denotemos por $\widehat Z$ la compleción de $Z$ a lo largo de
$h^{-1}(Y')$, y por
$\widehat g:\widehat Z\to\widehat X$ la extensión de $g$ a
las compleciones.

El módulo $\widehat g^{*}(\mathcal F)$ es un
$\mathcal O_{\widehat Z}$-módulo coherente cuyo soporte está contenido en
$g^{-1}(E)$ y es, por tanto, propio sobre $Y$, pues $g$ es
proyectivo y, por tanto, propio \sref[II]{II.5.4.6}. Como $h$ es
cuasiproyectivo, $\widehat g^{*}(\mathcal F)$ es
algebraizable en virtud de (\hyperref[III.5.2.6-fr]{5.2.6}). Se concluye
\oldpage[III]{156}
que $\widehat g_*\bigl(\widehat g^{*}(\mathcal F)\bigr)$
es un $\mathcal O_{\widehat X}$-Módulo algebraizable
(\hyperref[III.5.3.3-fr]{5.3.3}) puesto que $g$ es propio. Se puede ahora
aplicar a $\mathcal F$ y $g$. El núcleo $\mathcal P$
y el conúcleo $\mathcal R$ de
$\rho_{\mathcal F}:\mathcal F\to
\widehat g_*\bigl(\widehat g^{*}(\mathcal F)\bigr)$
son anulados por una potencia $\widehat{\mathcal M}^{n}$. Sea $T$
el subpreesquema cerrado de $X$ definido por $\mathcal M^n$, cuyo
espacio subyacente es $X\setmin U$, y sea $j:T\to X$ la inyección canónica, de modo que la prolongación a los completados $\widehat j:\widehat T\to\widehat X$ es la inclusión canónica (\textbf{I}, 10.14.7). Se
puede pues escribir
$\mathcal P=\widehat j_*\bigl(\widehat j^{*}(\mathcal P)\bigr))$ y
$\mathcal R=\widehat j_*\bigl(\widehat j^{*}(\mathcal R)\bigr)$ y como $U$ no es vacío, la hipótesis de inducción muestra que
$\widehat j^{*}(\mathcal P)$ y
$\widehat j^{*}(\mathcal R)$ son $\mathcal O_{\widehat T}$-módulos
algebraizables. El Lema \sref{III.5.3.3} muestra entonces
que $\mathcal P$ y $\mathcal R$ son algebraizables, y se puede entonces aplicar (\hyperref[III.5.3.2-fr]{5.3.2}), que prueba finalmente que $\mathcal F$ es algebraizable.



\subsection{Aplicación: comparación entre morfismos de esquemas ordinarios y morfismos de esquemas formales. Esquemas formales algebraizables}
\label{subsection:III.5.4}
\label{subsection:III.5.4-fr}

\begin{theorem}[5.4.1]
\label{III.5.4.1-fr}
Sean $A$ un anillo ádico noetheriano, $\mathfrak J$ un ideal de definición de
$A$, $S=\operatorname{Spec}(A)$, $S'=\mathrm V(\mathfrak J)$. Sean
$u:X\to S$ un morfismo propio, $v:Y\to S$ un morfismo separado de tipo finito, y
sean $\widehat S$, $\widehat X$, $\widehat Y$ los completados de $S$, $X$, $Y$ a lo
largo de $S'$, $u^{-1}(S')$, $v^{-1}(S')$ respectivamente. Si, para todo
$S$-morfismo $f:X\to Y$, $\widehat f:\widehat X\to\widehat Y$ es la
prolongación de $f$ a los completados, la aplicación $f\leadsto\widehat f$ es una
biyección
\[
\operatorname{Hom}_S(X,Y)\xrightarrow{\ \sim\ }
\operatorname{Hom}_{\widehat S}(\widehat X,\widehat Y).
\]
\end{theorem}

\begin{proof}
Mostremos primero que $f\leadsto\widehat f$ es inyectiva. Supongamos en efecto que dos
$S$-morfismos $f$, $g$ de $X$ en $Y$ son tales que $\widehat f=\widehat g$.
Se sabe entonces (\textbf{I}, 10.9.4) que existe una vecindad abierta $V$ de
$X'=u^{-1}(S')$ en la que $f$ y $g$ coinciden. Ahora bien, como $u$ es una aplicación
cerrada, se tiene $V=X$ (\hyperref[III.5.1.3.1-fr]{5.1.3.1}), de donde $f=g$.

Probemos ahora que $f\leadsto\widehat f$ es sobreyectiva, y sea pues $h$ un
$\widehat S$-morfismo $\widehat X\to\widehat Y$. Sea $Z=X\times_S Y$, y
designemos por $p:Z\to X$ y $q:Z\to Y$ las proyecciones canónicas; $Z$ es de
tipo finito sobre $S$ (\textbf{I}, 6.3.4), luego noetheriano; designemos por $\widehat Z$ su
completado a lo largo de $Z'=p^{-1}(u^{-1}(S'))$; se sabe que $\widehat Z$ se identifica
canónicamente con $\widehat X\times_{\widehat S}\widehat Y$, identificándose las proyecciones
$\widehat Z\to\widehat X$ y $\widehat Z\to\widehat Y$ con las
prolongaciones $\widehat p$ y $\widehat q$ (\textbf{I}, 10.9.7). Como $Y$ es separado sobre
$S$, $\widehat Y$ es separado sobre $\widehat S$ (\textbf{I}, 10.15.7), luego el morfismo
grafo
$\Gamma_h=(1_{\widehat X},h):\widehat X\to\widehat Z$
es una inmersión cerrada (\textbf{I}, 10.15.4). Sean $\mathfrak T$ el subpreesquema
formal cerrado de $\widehat Z$ asociado a esta inmersión, y
$j:\mathfrak T\to\widehat Z$ la inclusión canónica, de modo que
$\Gamma_h=j\circ w$, donde $w:\widehat X\to\mathfrak T$ es un isomorfismo
(\textbf{I}, 10.14.3) cuyo isomorfismo inverso es $\widehat p\circ j$; además,
$\mathfrak T$ es evidentemente propio sobre $\widehat S$, puesto que $\widehat X$ lo es;
se concluye (\hyperref[III.5.1.8-fr]{5.1.8}) que existe un subpreesquema
cerrado $T$ de $Z$ tal que
$\mathfrak T=\widehat T=T_{/(T\cap Z')}$, y que $j=\widehat i$, donde $i$ es
la inclusión canónica $T\to Z$ (\textbf{I}, 10.14.7). Entonces $p\circ i:T\to X$ es un
isomorfismo, pues lo mismo ocurre con
$(p\circ i)^{\widehat{\,}}=\widehat p\circ\widehat i$ por hipótesis, y basta
aplicar
\oldpage[III]{157}
(4.6.8) observando como antes que $S$ es la única vecindad de $S'$ en $S$.
Sea $g:X\to T$ el isomorfismo inverso de $p\circ i$, y pongamos
$f=q\circ i\circ g$, que es un morfismo $X\to Y$ cuyo grafo, por definición, es
$\Gamma_f=i\circ g$. Como $\widehat g$ es el isomorfismo inverso de
$(p\circ i)^{\widehat{\,}}=w$, se tiene
$(\Gamma_f)^{\widehat{\,}}=\widehat i\circ\widehat g=j\circ w=\Gamma_h$.
Pero se sabe que $(\Gamma_f)^{\widehat{\,}}=\Gamma_{\widehat f}$ (\textbf{I}, 10.9.8), de donde
finalmente $h=\widehat f$, lo que termina la demostración.
\end{proof}

Se puede pues decir, en el lenguaje de las categorías, que el funtor
$X\leadsto\widehat X$ es \emph{plenamente fiel} (\textbf{0}, 8.1.6) de la categoría de los
esquemas propios sobre $\operatorname{Spec}(A)$ en la categoría de los esquemas formales
propios sobre $\operatorname{Spf}(A)$, para todo anillo ádico noetheriano $A$;
establece por tanto una equivalencia entre la primera de estas categorías y una
\emph{subcategoría} de la segunda; los objetos de esta última se llamarán
\emph{esquemas formales algebraizables}. Para tal esquema $\mathfrak X$, existe
un esquema usual $X$, propio sobre $\operatorname{Spec}(A)$, determinado salvo isomorfismo
único, tal que $\mathfrak X$ es isomorfo a $\widehat X$.

\begin{scholium}[5.4.2]
\label{III.5.4.2-fr}
Con las notaciones de (\hyperref[III.5.4.1-fr]{5.4.1}), pongamos
$S_n=\operatorname{Spec}(A/\mathfrak J^{n+1})$,
$X_n=X\times_S S_n$, $Y_n=Y\times_S S_n$. Resulta de
(\hyperref[III.5.4.1-fr]{5.4.1}) y de (\textbf{I}, 10.12.3) que darse un
$S$-morfismo $f:X\to Y$ \emph{equivale} a dar un sistema inductivo
ádico $(S_n)$ \sref[I]{I.10.12.2} de
$S_n$-morfismos $f_n:X_n\to Y_n$.
\end{scholium}

\begin{remark}[5.4.3]
\label{III.5.4.3}
\label{III.5.4.3-fr}
Contrariamente a lo que podría sugerir el teorema de existencia
\sref{III.5.1.6}, existen esquemas formales propios sobre $\operatorname{Spf}(A)$
que \emph{no son algebraizables}, así como existen espacios analíticos
compactos que no proceden de variedades algebraicas complejas. Más
adelante encontraremos tales esquemas en la «teoría de moduli»,
que, cuando el cuerpo base es $\mathbb C$, trata precisamente de las
variaciones infinitesimales de la estructura compleja de una variedad
algebraica completa; es bien sabido que tales variaciones pueden dar
lugar a variedades analíticas que no son algebraicas.
\end{remark}

\begin{proposition}[5.4.4]
\label{III.5.4.4-fr}
Sean $A$ un anillo ádico noetheriano, $\mathfrak S=\operatorname{Spf}(A)$,
$g:\mathfrak X\to\mathfrak S$, $h:\mathfrak Y\to\mathfrak S$ dos morfismos
propios de esquemas formales, $f:\mathfrak X\to\mathfrak Y$ un
$\mathfrak S$-morfismo. Si $f$ es finito y si $\mathfrak Y$ es algebraizable,
entonces $\mathfrak X$ es algebraizable.
\end{proposition}

\begin{proof}
Se observará que las hipótesis sobre $g$ y $h$ implican ya que $f$ es propio
(3.4.1), y para que $f$ sea finito, basta que para todo $y\in\mathfrak Y$ la
fibra $f^{-1}(y)$ sea \emph{finita} (\hyperref[III.4.8.11-fr]{4.8.11}). La hipótesis
implica que $\mathcal B=f_*(\mathcal O_{\mathfrak X})$ es una
$\mathcal O_{\mathfrak Y}$-Álgebra coherente
(\hyperref[III.4.8.6-fr]{4.8.6}), luego resulta del teorema de existencia que, si
$\mathfrak Y=\widehat Y$ y $h=\widehat w$, donde
$w:Y\to\operatorname{Spec}(A)$ es un morfismo propio de esquemas usuales,
existe una $\mathcal O_Y$-Álgebra coherente $\mathcal C$ tal que
$\mathcal B=\widehat{\mathcal C}$. Sean $X=\operatorname{Spec}(\mathcal C)$, y
$u:X\to Y$ el morfismo estructural; entonces resulta inmediatamente de la definición de
$\mathfrak X$ a partir de $\mathcal B$
(\hyperref[III.4.8.7-fr]{4.8.7}) que $\mathfrak X$ es canónicamente isomorfo a
$\widehat X$ y que $f$ se identifica con $\widehat u$ (basta para verlo
considerar el caso en que $Y$ es afín).

Se observará que (\hyperref[III.5.1.8-fr]{5.1.8}) es un caso particular de
(\hyperref[III.5.4.4-fr]{5.4.4}).
\end{proof}

\begin{theorem}[5.4.5]
\label{III.5.4.5-fr}
Sean $A$ un anillo ádico noetheriano, $\mathfrak J$ un ideal de definición de
$A$, $S=\operatorname{Spec}(A)$,
$\mathfrak S=\widehat S=\operatorname{Spf}(A)$,
$f:\mathfrak X\to\mathfrak S$ un morfismo propio de esquemas formales. Se pone
$S_k=\operatorname{Spec}(A/\mathfrak J^{k+1})$,
$X_k=\mathfrak X\times_{\mathfrak S}S_k$, y para todo
$\mathcal O_{\mathfrak X}$-Módulo $\mathcal F$,
$\mathcal F_k=\mathcal F\otimes_{\mathcal O_{\mathfrak X}}\mathcal O_{X_k}
=\mathcal F/\mathfrak J^{k+1}\mathcal F$.
\oldpage[III]{158}
Sea $\mathcal L$ un $\mathcal O_X$-Módulo invertible, y supongamos
que $\mathcal L_0=\mathcal L/\mathfrak J\mathcal L$ es un
$\mathcal O_{X_0}$-Módulo amplio. Entonces $\mathfrak X$ es algebraizable, y si
$X$ es un $S$-esquema propio tal que $\mathfrak X$ es isomorfo a
$\widehat X$, existe un $\mathcal O_X$-Módulo amplio $\mathcal M$ tal que
$\mathcal L$ es isomorfo a $\widehat{\mathcal M}$ \textup{(lo que implica que
$X$ es \emph{proyectivo} sobre $S$)}.
\end{theorem}

\begin{proof}
Apliquemos (\hyperref[III.5.2.3-fr]{5.2.3}) a
$\mathcal F=\mathcal O_{\mathfrak X}$: existe pues un entero $n_0$ tal que
para $n\geq n_0$ el homomorfismo canónico
$\Gamma(\mathfrak X,\mathcal L^{\otimes n})\to
\Gamma(X_0,\mathcal L_0^{\otimes n})$ es \emph{sobreyectivo}. Se puede suponer
$n\geq n_0$ suficientemente grande para que $\mathcal L_0^{\otimes n}$ sea \emph{muy
amplio} para $S_0$ (\textbf{II}, 4.5.10). Como el morfismo $f_0:X_0\to S_0$ es propio,
$\Gamma(X_0,\mathcal L_0^{\otimes n})$ es un $A$-módulo de tipo finito (3.2.1),
luego existe un sub-$A$-módulo de tipo finito $E$ de
$\Gamma(\mathfrak X,\mathcal L^{\otimes n})$ cuya imagen en
$\Gamma(X_0,\mathcal L_0^{\otimes n})$ es este último módulo entero.
Dicho esto, para todo $k\geq0$, consideremos el homomorfismo
$u_k:E/\mathfrak J^{k+1}E\longrightarrow
\Gamma(X_k,\mathcal L_k^{\otimes n})$
deducido de la inclusión canónica
$E\to\Gamma(\mathfrak X,\mathcal L^{\otimes n})$. Observemos que
$(f_k)_*(\mathcal L_k^{\otimes n})$ es cuasicoherente, y como
$\Gamma(S_k,(f_k)_*(\mathcal L_k^{\otimes n}))
=\Gamma(X_k,\mathcal L_k^{\otimes n})$, $u_k$ define un homomorfismo
$\widetilde u_k:(E/\mathfrak J^{k+1}E)^\sim
\longrightarrow(f_k)_*(\mathcal L_k^{\otimes n})$,
y por tanto también un homomorfismo
$\widetilde u_k^\sharp:
f_k^*((E/\mathfrak J^{k+1}E)^\sim)\longrightarrow
\mathcal L_k^{\otimes n}$.
Por lo demás, si se pone
$\mathcal G_k=f_k^*((E/\mathfrak J^{k+1}E)^\sim)$, se tiene
$\mathcal G_0=\mathcal G_k/\mathfrak J\mathcal G_k$ (\textbf{I}, 9.1.5), luego
$\widetilde u_0^\sharp:\mathcal G_k/\mathfrak J\mathcal G_k
\to\mathcal L_k^{\otimes n}/\mathfrak J\mathcal L_k^{\otimes n}$ se deduce
de $\widetilde u_k^\sharp$ por paso a los cocientes. Ahora bien, por definición de
$E$, $\widetilde u_0^\sharp$ no es otro que el homomorfismo canónico
$\sigma:f_0^*((f_0)_*(\mathcal L_0^{\otimes n}))\longrightarrow
\mathcal L_0^{\otimes n}$,
y la hipótesis de que $\mathcal L_0^{\otimes n}$ es muy amplio implica que
$\widetilde u_0^\sharp$ es \emph{sobreyectivo} (\textbf{II}, 4.4.3); se deduce entonces del lema
de Nakayama que cada uno de los $\widetilde u_k^\sharp$ es también \emph{sobreyectivo}.
Cada uno de los $\widetilde u_k^\sharp$ define pues (\textbf{II}, 4.2.2) un
$S_k$-morfismo
$g_k:X_k\longrightarrow P_{\boldsymbol{k}}=\mathbf P(E/\mathfrak J^{k+1}E)$,
y como $P_h=P_k\times_{S_k}S_h$ para $h\leq k$ en virtud de (\textbf{II}, 4.1.3),
$(g_k)$ es un $(S_k)$-\emph{sistema inductivo ádico} (\textbf{I}, 10.12.2) en virtud de las
relaciones entre los $\widetilde u_k^\sharp$ y de (\textbf{II}, 4.2.10). Los $g_k$
definen pues un $\mathfrak S$-morfismo de esquemas formales
$g:\mathfrak X\to\mathfrak P$, donde $\mathfrak P$ es el límite inductivo del
sistema $(P_k)$, o también el completado $\widehat P$, donde $P=\mathbf P(E)$.
Además, la hipótesis de que $\mathcal L_0^{\otimes n}$ es muy amplio implica que
$g_0$ es una \emph{inmersión cerrada} (\textbf{II}, 4.4.3); se concluye que $g$ es una
inmersión cerrada de esquemas formales (4.8.10), luego $\mathfrak X$ es
algebraizable (5.1.8). El hecho de que $\mathcal L$ sea isomorfo al completado
$\widehat{\mathcal M}$ de un $\mathcal O_X$-Módulo invertible resulta entonces del
teorema de existencia (5.1.6). Además, $\mathcal L^{\otimes n}$ es entonces el
completado de $\mathcal M^{\otimes n}$ (\textbf{I}, 10.8.10), y los homomorfismos
$\widetilde u_k^\sharp$ definen un homomorfismo bien determinado
$v:f^*(\widetilde E)\longrightarrow\mathcal M^{\otimes n}$
(5.1.7); por lo demás, como $\widetilde u_0^\sharp$ es sobreyectivo, lo mismo ocurre con
$\widehat v$ (\textbf{I}, 10.11.5), luego con $v$ (5.1.3); además, el morfismo
$r:X\to P$ definido por $v$ (\textbf{II}, 4.2.2) tiene por prolongación a los completados $g$,
y como $g$ es una inmersión cerrada, lo mismo ocurre con $r$, por (5.1.8)
y (5.4.1); se concluye que $\mathcal M^{\otimes n}$ es muy amplio
(\textbf{II}, 4.4.6) y $\mathcal M$ es amplio (\textbf{II}, 4.5.10).
\end{proof}

\begin{remark}[5.4.6]
\label{III.5.4.6}
\label{III.5.4.6-fr}
Sea $A$ un anillo ádico noetheriano, $\mathfrak S=\operatorname{Spf}(A)$, y sea
$f:\mathfrak X\to\mathfrak S$ un morfismo propio de esquemas
formales. Sea $\mathcal N$ el ideal coherente de
$\mathcal O_{\mathfrak X}$ tal que, para todo subconjunto abierto
formal afín $U$ de $\mathfrak X$,
$\Gamma(U,\mathcal N)$
es el nilradical de
$\Gamma(U,\mathcal O_{\mathfrak X})$. La existencia de este ideal
se deduce fácilmente de \sref[I]{I.10.10.2} y del hecho de que
todo homomorfismo de anillos $B\to C$ envía el nilradical de $B$ en
el de $C$. Sea $\mathfrak X'$ el subesquema formal cerrado de
$\mathfrak X$ definido por $\mathcal N$
\sref[I]{I.10.14.2}. Sería interesante saber si la
algebraizabilidad de $\mathfrak X'$ implica la algebraizabilidad del
propio $\mathfrak X$.

Sin duda se llegaría a una solución de este problema si se supiera
clasificar, por ejemplo mediante invariantes
\oldpage[III]{159}
de naturaleza cohomológica, las \emph{extensiones} de un haz estructural
$\mathcal O_{\mathfrak X}$ (para un preesquema usual o un preesquema formal) por un Ideal de cuadrado nulo, en otros términos los $\mathcal O_{\mathfrak X}$-álgebras $\mathcal A$ tales que $\mathcal O_{\mathfrak X}$ es isomorfo a
$\mathcal A/\mathcal J$, donde $\mathcal J$ es un
ideal de cuadrado nulo de $\mathcal A$.
\end{remark}


\subsection{Una descomposición de ciertos esquemas}
\label{subsection:III.5.5}
\label{subsection:III.5.5-fr}

\begin{proposition}[5.5.1]
\label{III.5.5.1-fr}
Sean $A$ un anillo ádico noetheriano, $\mathfrak J$ un ideal de definición
de $A$, $Y=\operatorname{Spec}(A)$. Sea $f:X\to Y$ un morfismo separado de
tipo finito; se pone $Y_0=\operatorname{Spec}(A/\mathfrak J)$,
$X_0=X\times_Y Y_0=f^{-1}(Y_0)$. Sea $Z_0$ una parte abierta de $X_0$,
propia sobre $Y_0$; existe entonces en $X$ una parte abierta y cerrada $Z$,
propia sobre $Y$ y tal que $Z\cap X_0=Z_0$.
\end{proposition}

\begin{proof}
Por hipótesis, hay un abierto $T$ de $X$ tal que $T\cap X_0=Z_0$; sea
$\widehat T$ el completado a lo largo de $Z_0$ del esquema inducido por $X$ sobre
el abierto $T$; al ser el soporte de $\mathcal O_{\widehat T}$ igual a $Z_0$, que es
propio sobre $Y_0$, $\widehat T$ es propio sobre
$\widehat Y=\operatorname{Spf}(A)$ (3.4.1). Resulta de (5.1.8) que existe
un subesquema cerrado $Z$ de $T$ \emph{propio} sobre $Y$, tal que, si $i:Z\to T$ es
la inclusión canónica, $\widehat i=\widehat Z\to\widehat T$ es un
isomorfismo (siendo $\widehat Z$ el completado de $Z$ a lo largo de $Z_0$). Se
concluye (4.6.8) que existe en $T$ una vecindad abierta $V$ de $Z_0$ tal que
la restricción $i^{-1}(V)\to V$ de $i$ es un isomorfismo. Pero
$i^{-1}(V)$ es una vecindad de $Z_0$ en $Z$, luego es necesariamente
idéntico a $Z$ (5.1.3.1). Se concluye que $Z$ es abierto en $T$, luego
en $X$, lo que termina la demostración.
\end{proof}

\begin{corollary}[5.5.2]
\label{III.5.5.2}
\label{III.5.5.2-fr}
Si $X_0$ es propio sobre $Y_0$, entonces $X$ es la unión de dos
subconjuntos abiertos disjuntos $Z$ y $Z'$ tales que $Z$ es propio sobre
$Y$ y contiene a $X_0$. Además, todo subconjunto cerrado $P$ de $X$
que es propio sobre $Y$ está contenido en $Z$.
\end{corollary}

\begin{proof}
La última afirmación resulta de que $P\cap Z'$ es cerrado en $P$ y, por tanto, propio
sobre $Y$; si $P\cap Z'$ no fuera vacío, $f(P\cap Z')$ sería un subconjunto cerrado
no vacío de $Y$ y, por tanto, cortaría a $Y_0$ (5.1.3.1), lo que contradice la
definición de $Z$.
\end{proof}

\begin{flushright}
\emph{(Continuará.)}
\end{flushright}

% La página impresa 504 está en blanco.
\clearpage
\oldpage[III]{160}
\thispagestyle{empty}
\null
\clearpage
\thispagestyle{plain}

\oldpage[III]{161}
\section*{BIBLIOGRAFÍA \textit{(continuación)}}

\begin{flushleft}
[27] P. \textsc{Gabriel}, \emph{Des catégories abéliennes}, tesis, París (1961).

\medskip
[28] J. E. \textsc{Roos}, Sur les foncteurs dérivés de $\varprojlim$.
Applications, \emph{C. R. Acad. Sci. Paris}, vol. CCLII, 1961, pp. 3702-3704.

\medskip
[29] H. \textsc{Cartan}, Séminaire de l'École Normale Supérieure,
13.er año (1960-1961), exposición n.º 11.
\end{flushleft}

\clearpage
\oldpage[III]{162}
\section*{ÍNDICE DE NOTACIONES}

\begin{flushleft}
$\boldsymbol{\mathit{Ens}}$: \textbf 0, 8.1.1.

$h_X(Y)$, $h_X(u)$ ($X$, $Y$ objetos de una categoría $\boldsymbol{C}$, $u$ un morfismo
de $\boldsymbol{C}$): \textbf 0, 8.1.1.

$h_w(u)$ ($u$ y $w$ morfismos de $\boldsymbol{C}$): \textbf 0, 8.1.2.

$Z_k(E_r^{pq})$, $B_k(E_r^{pq})$ ($k\geq r$): \textbf 0, 11.1.1.

$Z_\infty(E_2^{pq})$, $B_\infty(E_2^{pq})$: \textbf 0, 11.1.1.

$K^\bullet$, $K_\bullet$ (complejos): \textbf 0, 11.2.1.

$K^{\bullet\bullet}$, $K_{\bullet\bullet}$ (bicomplejos): \textbf 0, 11.2.1.

$B^i(K^\bullet)$, $Z^i(K^\bullet)$, $H^i(K^\bullet)$,
$H^\bullet(K^\bullet)$: \textbf 0, 11.2.1.

$B_i(K_\bullet)$, $Z_i(K_\bullet)$, $H_i(K_\bullet)$,
$H_\bullet(K_\bullet)$: \textbf 0, 11.2.1.

$T(K^\bullet)$, $T(K_\bullet)$ ($T$ un funtor): \textbf 0, 11.2.1.

$E(K^\bullet)$ ($K^\bullet$ un complejo filtrado): \textbf 0, 11.2.2.

$K^{i,\bullet}$, $K^{\bullet,j}$, $Z_{\mathrm{II}}^p(K^{i,\bullet})$,
$B_{\mathrm{II}}^p(K^{i,\bullet})$, $H_{\mathrm{II}}^p(K^{i,\bullet})$,
$Z_{\mathrm I}^p(K^{\bullet,j})$, $B_{\mathrm I}^p(K^{\bullet,j})$,
$H_{\mathrm I}^p(K^{\bullet,j})$: \textbf 0, 11.3.1.

$H_{\mathrm{II}}^p(K^{\bullet\bullet})$,
$Z_{\mathrm I}^q(H_{\mathrm{II}}^p(K^{\bullet\bullet}))$,
$B_{\mathrm I}^q(H_{\mathrm{II}}^p(K^{\bullet\bullet}))$,
$H_{\mathrm I}^q(H_{\mathrm{II}}^p(K^{\bullet\bullet}))$,
$H^n(K^{\bullet\bullet})$: \textbf 0, 11.3.1.

$F_{\mathrm I}^p(K^{\bullet\bullet})$,
$F_{\mathrm{II}}^p(K^{\bullet\bullet})$, ${}'E(K^{\bullet\bullet})$,
${}''E(K^{\bullet\bullet})$: \textbf 0, 11.3.2.

$K_{i,\bullet}$, $K_{\bullet,j}$, $H_p^{\mathrm{II}}(K_{i,\bullet})$,
$H_p^{\mathrm I}(K_{\bullet,j})$, $H_p^{\mathrm{II}}(K_{\bullet\bullet})$,
$H_p^{\mathrm I}(K_{\bullet\bullet})$,
$H_q^{\mathrm I}(H_p^{\mathrm{II}}(K_{\bullet\bullet}))$,
$H_q^{\mathrm{II}}(H_p^{\mathrm I}(K_{\bullet\bullet}))$,
$H_n(K_{\bullet\bullet})$: \textbf 0, 11.3.5.

$\mathbf R^\bullet T(K^\bullet)$ ($T$ un funtor covariante aditivo): \textbf 0, 11.4.3.

$\mathbf R^\bullet T(K^\bullet,K'^{\bullet})$ ($T$ un bifuntor covariante
aditivo): \textbf 0, 11.4.6.

$\mathbf L_\bullet T(K_\bullet)$ ($T$ un funtor covariante aditivo): \textbf 0, 11.6.2.

$\mathbf L_\bullet T(K_\bullet,K'_\bullet)$ ($T$ un bifuntor covariante
aditivo): \textbf 0, 11.6.6.

$\mathbf L_\bullet T(K_{\bullet\bullet})$ ($T$ un funtor covariante aditivo):
\textbf 0, 11.7.2.

$|\sigma|$, $\Sigma(A)$, $\mathrm C_\bullet(A)$, $\mathrm L_\bullet(A)$
($A$ un conjunto finito, $\sigma$ un símplex de $A$): \textbf 0, 11.8.1.

$\mathrm C^\bullet(A,B;\mathscr S)$, $\mathrm L^\bullet(A,B;\mathscr S)$
($A$, $B$ conjuntos finitos, $\mathscr S$ un sistema de coeficientes): \textbf 0, 11.8.4.

$\mathrm P^\bullet(A,B,\mathscr S)$ ($A$, $B$ conjuntos finitos, $\mathscr S$ un
sistema de coeficientes): \textbf 0, 11.8.6.

$\mathrm C^\bullet(A,B;\mathscr S^\bullet)$,
$\mathrm L^\bullet(A,B;\mathscr S^\bullet)$ ($A$, $B$ conjuntos finitos,
$\mathscr S^\bullet$ un complejo de sistemas de coeficientes): \textbf 0, 11.8.8.

$\mathrm P^\bullet(A,B;\mathscr S^\bullet)$ ($A$, $B$ conjuntos finitos,
$\mathscr S^\bullet$ un complejo de sistemas de coeficientes): \textbf 0, 11.8.10.

$\chi(M^\bullet)$ ($M^\bullet$ un complejo finito de módulos de longitud finita):
\textbf 0, 11.10.1.

$C^\bullet(\mathfrak U,\mathcal F)$ ($\mathcal F$ un haz de grupos abelianos
sobre $X$, $\mathfrak U$ un recubrimiento abierto de $X$): \textbf 0, 12.1.2.

$\mathrm R^pf_*(\mathcal F)$ ($f:X\to Y$ un morfismo de espacios anillados,
$\mathcal F$ un haz de $\mathcal O_X$-módulos): \textbf 0, 12.2.1.

$\mathscr H^p(f,\mathcal K^\bullet)$,
$\mathscr H_f^p(\mathcal K^\bullet)$, ${}'\mathcal E(f,\mathcal K^\bullet)$,
${}''\mathcal E(f,\mathcal K^\bullet)$ ($f:X\to Y$ un morfismo de espacios anillados,
$\mathcal K^\bullet$ un complejo de $\mathcal O_X$-módulos): \textbf 0, 12.4.1.

$\mathbf H^\bullet(X,\mathcal K^\bullet)$,
$\mathbf H^\bullet(V,\mathcal K^\bullet)$ ($X$ un espacio anillado,
$\mathcal K^\bullet$ un complejo de $\mathcal O_X$-módulos, $V$ abierto en $X$):
\textbf 0, 12.4.2.

$\mathbf{gr}^\bullet(\mathbf A)$ ($\mathbf A$ un sistema proyectivo que satisface (ML)):
\textbf 0, 13.4.2.

$E(A)$ ($A$ un objeto filtrado): \textbf 0, 13.6.4.

$K_\bullet(\mathbf f)$, $i_{\mathbf f}$ ($\mathbf f$ una familia finita de
elementos de un anillo): \textbf{III}, 1.1.1.

$K^\bullet(\mathbf f,M)$, $K_\bullet(\mathbf f,M)$ ($\mathbf f$ una familia
finita de elementos de un anillo $A$, $M$ un $A$-módulo): \textbf{III}, 1.1.2.

$H^\bullet(\mathbf f,M)$, $H_\bullet(\mathbf f,M)$ ($\mathbf f$ una familia
finita de elementos de un anillo $A$, $M$ un $A$-módulo): \textbf{III}, 1.1.3.

$C^\bullet((\mathbf f),M)$, $H^\bullet((\mathbf f),M)$ ($\mathbf f$ una
familia finita de elementos de un anillo $A$, $M$ un $A$-módulo): \textbf{III}, 1.1.6.

$H^\bullet(U,\mathcal F(*))$,
$H^\bullet(\mathfrak U,\mathcal F(*))$ ($\mathcal F$ un
$\mathcal O_X$-módulo cuasicoherente, $U$ abierto en $X$, $\mathfrak U$ un recubrimiento abierto de $X$):
\textbf{III}, 2.1.1.
\end{flushleft}

\clearpage
\oldpage[III]{163}
\section*{ÍNDICE TERMINOLÓGICO}

\begin{flushleft}
Término de convergencia de una sucesión espectral: \textbf 0, 11.1.1.

Algebraizable ($\mathcal O_{\widehat X}$-módulo): \textbf{III}, 5.2.1.

Algebraizable (esquema formal): \textbf{III}, 5.4.2.

Analíticamente íntegro (anillo): \textbf{III}, 4.3.6.

Aplicación cuasicompacta: \textbf 0, 9.1.1.

Aumentación de una resolución: \textbf 0, 11.4.1.

Característica de Euler--Poincaré de un complejo: \textbf 0, 11.10.1.

Característica de Euler--Poincaré de un $\mathcal O_X$-módulo coherente
($X$ proyectivo sobre $\operatorname{Spec}(A)$, $A$ un anillo artiniano): \textbf{III}, 2.5.1.

Cocadena bialternada: \textbf 0, 11.8.4.

Complejo definido por un bicomplejo: \textbf 0, 11.3.1.

Complejo de álgebra exterior: \textbf{III}, 1.1.1.

Condición (ML): \textbf 0, 13.1.2.

Construible (subconjunto, conjunto): \textbf 0, 9.1.2.

Construible (función): \textbf 0, 9.3.1.

Producto cup: \textbf 0, 12.1.2.

Dihomomorfismo para estructuras algebraicas sobre categorías: \textbf 0, 8.2.1.

Exacto (subconjunto) en una categoría abeliana: \textbf{III}, 3.1.1.

Esencialmente constante (sistema proyectivo): \textbf 0, 13.4.2.

Factorización de Stein: \textbf{III}, 4.3.3.

Codiscreta, coseparada, discreta, exhaustiva, finita, separada
(filtración): \textbf 0, 11.1.3.

Final (objeto) de una categoría: \textbf 0, 8.1.10.

Finito (morfismo) de preesquemas formales: \textbf{III}, 4.8.2.

Finito (preesquema formal) sobre $\mathfrak Y$,
$\mathfrak Y$-finito (preesquema formal): \textbf{III}, 4.8.2.

Funtor covariante canónico $\boldsymbol{C}\to\boldsymbol{\mathit{Hom}}(\boldsymbol{C}^0,\boldsymbol{\mathit{Ens}})$: \textbf 0, 8.1.2.

Género aritmético: \textbf{III}, 2.5.1.

Geométricamente conexa (fibra): \textbf{III}, 4.3.4.

Homomorfismo de estructuras algebraicas sobre categorías: \textbf 0, 8.2.1.

Hipercohomología de un funtor respecto de un complejo $K^\bullet$:
\textbf 0, 11.4.3.

Hipercohomología de un bifuntor respecto de dos complejos
$K^\bullet$, $K'^{\bullet}$: \textbf 0, 11.4.6.

Hiperhomología de un funtor respecto de un complejo $K_\bullet$:
\textbf 0, 11.6.2.

Hiperhomología de un bifuntor respecto de dos complejos $K_\bullet$,
$K'_\bullet$: \textbf 0, 11.6.6.

Hiperhomología de un funtor respecto de un bicomplejo: \textbf 0, 11.7.2.

Localmente construible (conjunto): \textbf 0, 9.1.11.

Ley de composición externa (interna) sobre los objetos de una categoría: \textbf 0, 8.2.1.

\oldpage[III]{164}

$S$-$\boldsymbol{C}$-módulo filtrado, $\mathbf{gr}^\bullet(S)$-$\boldsymbol{C}$-módulo graduado,
$\mathbf{gr}^\bullet(S)$-$\boldsymbol{C}$-módulo bigraduado: \textbf 0, 13.6.6.

Morfismo de sucesiones espectrales: \textbf 0, 11.1.2.

Número geométrico de componentes conexas de una fibra: \textbf{III}, 4.3.4.

Objeto grupo en $\boldsymbol{C}$, $\boldsymbol{C}$-grupo, $\boldsymbol{C}$-anillo, $\boldsymbol{C}$-módulo: \textbf 0, 8.2.3.

Objeto de bordes, objeto de cobordes, objeto de cociclos, objeto de
ciclos: \textbf 0, 11.2.1.

Objeto de imágenes universales: \textbf 0, 13.1.1.

Objeto graduado asociado a un sistema proyectivo que satisface (ML): \textbf 0, 13.4.2.

Objeto graduado acotado inferiormente (superiormente): \textbf 0, 11.2.1.

Subconjunto propio (sobre $\mathfrak Y$) de un preesquema formal: \textbf{III}, 3.4.1.

Plena (subcategoría): \textbf 0, 8.1.5.

Plenamente fiel (funtor): \textbf 0, 8.1.5.

Polinomio de Hilbert: \textbf{III}, 2.5.3.

Propio (morfismo) de preesquemas formales: \textbf{III}, 3.4.1.

Representable (funtor): \textbf 0, 8.1.8.

Cohomológica, derecha, izquierda, homológica, inyectiva, libre, plana, proyectiva
(resolución): \textbf 0, 11.4.1.

Resolución de Cartan--Eilenberg derecha, inyectiva: \textbf 0, 11.4.2.

Resolución de Cartan--Eilenberg izquierda, proyectiva: \textbf 0, 11.6.1.

Retrocompacto (conjunto): \textbf 0, 9.1.1.

Estricto (sistema proyectivo): \textbf 0, 13.4.2.

Sucesión espectral en una categoría abeliana: \textbf 0, 11.1.1.

Débilmente convergente, regular, coregular, birregular (sucesión espectral):
\textbf 0, 11.1.3.

Sucesión espectral degenerada: \textbf 0, 11.1.6.

Sucesión espectral de un funtor relativa a un objeto filtrado: \textbf 0, 13.6.4.

Sucesiones espectrales de un bicomplejo: \textbf 0, 11.3.2 y 11.3.5.

Sucesiones espectrales de hipercohomología de un funtor respecto de un complejo:
\textbf 0, 11.4.3.

Sucesiones espectrales de hiperhomología de un funtor respecto de un complejo:
\textbf 0, 11.6.2.

Sucesiones espectrales de hiperhomología de un funtor respecto de un bicomplejo:
\textbf 0, 11.7.2.

Sistema de coeficientes: \textbf 0, 11.8.4.

Unirrama (anillo, punto): \textbf{III}, 4.3.6.

Universalmente abierto (morfismo): \textbf{III}, 4.3.9.
\end{flushleft}

\clearpage
\oldpage[III]{165}
\section*{ÍNDICE GENERAL}

\begin{flushleft}
\textsc{Capítulo 0.} --- \textbf{Preliminares} \textit{(continuación)} \dotfill 5

\medskip
\S~8. Funtores representables \dotfill 5

8.1. Funtores representables \dotfill 5

8.2. Estructuras algebraicas en categorías \dotfill 9

\medskip
\S~9. Conjuntos construibles \dotfill 12

9.1. Conjuntos construibles \dotfill 12

9.2. Subconjuntos construibles de espacios noetherianos \dotfill 14

9.3. Funciones construibles \dotfill 16

\medskip
\S~10. Suplemento sobre los módulos planos \dotfill 17

10.1. Relaciones entre módulos planos y módulos libres \dotfill 17

10.2. Criterios locales de planitud \dotfill 18

10.3. Existencia de extensiones planas de anillos locales \dotfill 20

\medskip
\S~11. Suplemento sobre álgebra homológica \dotfill 23

11.1. Recordatorio sobre las sucesiones espectrales \dotfill 23

11.2. La sucesión espectral de un complejo filtrado \dotfill 27

11.3. Las sucesiones espectrales de un bicomplejo \dotfill 29

11.4. Hipercohomología de un funtor respecto de un complejo $K^\bullet$
\dotfill 32

11.5. Paso al límite inductivo en hipercohomología \dotfill 35

11.6. Hiperhomología de un funtor respecto de un complejo $K_\bullet$
\dotfill 39

11.7. Hiperhomología de un funtor respecto de un bicomplejo
$K_{\bullet\bullet}$ \dotfill 41

11.8. Suplemento sobre la cohomología de los complejos simpliciales \dotfill 43

11.9. Un lema sobre complejos de tipo finito \dotfill 46

11.10. Característica de Euler--Poincaré de un complejo de módulos de longitud finita
\dotfill 48

\medskip
\S~12. Suplemento sobre la cohomología de haces \dotfill 49

12.1. Cohomología de haces de módulos sobre espacios anillados \dotfill 49

12.2. Imágenes directas superiores \dotfill 57

12.3. Suplementos sobre los funtores Ext de haces \dotfill 60

12.4. Hipercohomología del funtor de imagen directa \dotfill 62

\oldpage[III]{166}

\S~13. Límites proyectivos en álgebra homológica \dotfill 64

13.1. La condición de Mittag--Leffler \dotfill 64

13.2. La condición de Mittag--Leffler para grupos abelianos \dotfill 65

13.3. Aplicación: cohomología de un límite proyectivo de haces \dotfill 68

13.4. Condición de Mittag--Leffler y objetos graduados asociados a sistemas
proyectivos \dotfill 69

13.5. Límites proyectivos de sucesiones espectrales de complejos filtrados
\dotfill 71

13.6. Sucesión espectral de un funtor relativa a un objeto dotado de una
filtración finita \dotfill 73

13.7. Funtores derivados de un límite proyectivo de argumentos \dotfill 75

\medskip
\textsc{Capítulo III.} --- \textbf{Estudio cohomológico de los haces coherentes}
\dotfill 81

\medskip
\S~1. Cohomología de los esquemas afines \dotfill 82

1.1. Recordatorio sobre el complejo de álgebra exterior \dotfill 82

1.2. Cohomología de \v Cech de un recubrimiento abierto \dotfill 85

1.3. Cohomología de un esquema afín \dotfill 88

1.4. Aplicación a la cohomología de preesquemas cualesquiera \dotfill 89

\medskip
\S~2. Estudio cohomológico de los morfismos proyectivos \dotfill 95

2.1. Cálculos explícitos de ciertos grupos de cohomología \dotfill 95

2.2. El teorema fundamental de los morfismos proyectivos \dotfill 100

2.3. Aplicación a haces graduados de álgebras y módulos \dotfill 102

2.4. Una generalización del teorema fundamental \dotfill 107

2.5. Característica de Euler--Poincaré y polinomio de Hilbert \dotfill 109

2.6. Aplicación: criterios de amplitud \dotfill 111

\medskip
\S~3. Teorema de finitud para los morfismos propios \dotfill 115

3.1. El lema de dévissage \dotfill 115

3.2. El teorema de finitud: caso de los esquemas ordinarios \dotfill 116

3.3. Generalización del teorema de finitud (esquemas ordinarios) \dotfill 118

3.4. Teorema de finitud: caso de los esquemas formales \dotfill 119

\medskip
\S~4. El teorema fundamental de los morfismos propios. Aplicaciones \dotfill 122

4.1. El teorema fundamental \dotfill 122

4.2. Casos particulares y variantes \dotfill 129

4.3. Teorema de conexión de Zariski \dotfill 130

4.4. «Teorema principal» de Zariski \dotfill 135

4.5. Compleciones de módulos de homomorfismos \dotfill 138

4.6. Relaciones entre morfismos formales y morfismos ordinarios \dotfill 139

4.7. Un criterio de amplitud \dotfill 145

4.8. Morfismos finitos de preesquemas formales \dotfill 146

\oldpage[III]{167}

\S~5. Un teorema de existencia para los haces algebraicos coherentes \dotfill 149

5.1. Enunciado del teorema \dotfill 149

5.2. Demostración del teorema de existencia: caso proyectivo y cuasiproyectivo
\dotfill 151

5.3. Demostración del teorema de existencia: caso general \dotfill 154

5.4. Aplicación: comparación entre morfismos de esquemas ordinarios y
morfismos de esquemas formales. Esquemas formales algebraizables \dotfill 156

5.5. Una descomposición de ciertos esquemas \dotfill 159

\medskip
\textsc{Bibliografía (continuación)} \dotfill 161

\textsc{Índice de notaciones} \dotfill 162

\textsc{Índice terminológico} \dotfill 163

\begin{flushright}
\emph{Recibido el 28 de junio de 1961.}
\end{flushright}
\end{flushleft}

% END tree/repo/editions/III1_DIPLOMATIC/ega3/ega3-5.tex


\clearpage
\phantomsection
\addcontentsline{toc}{section}{Referencias consolidadas de la edición}
\section*{Referencias consolidadas de la edición}
\begin{thebibliography}{Roo61}

\bibitem[Car]{III-29}
H.~Cartan, \emph{S\'eminaire de l'\'ecole normale sup\'erieure, 13th year
  (1960--61), expos\'e no 11}.

\bibitem[Roo61]{III-28}
J.~E. Roos, \emph{Sur les foncteurs d\'eriv\'es de $\varprojlim$.
  applications}, C. R. Acad. Sci. Paris \textbf{CCLII} (1961), 3702--3704.

\bibitem[SZ60]{I-13}
P.~Samuel and O.~Zariski, \emph{Commutative algebra, 2 vol.}, New York (Van
  Nostrand), 1958--60.

\bibitem[Zar51]{I-20}
O.~Zariski, \emph{Theory and applications of holomorphic functions on algebraic
  varieties over arbitrary ground fields}, Mem. Amer. Math. Soc., no. 5, 1951.

\end{thebibliography}




\end{document}
